\documentclass[UTF8]{ctexbook}
\usepackage{graphicx}
\usepackage{amsmath}
\usepackage{geometry}
\usepackage{hyperref}
\usepackage{listings}
\usepackage{xcolor}
\usepackage{multirow,threeparttable,longtable,upgreek}
\usepackage{makecell}
\usepackage{amssymb}
\usepackage[many]{tcolorbox}
\usepackage{float}
\usepackage{pifont}
\newcommand{\◎}{\ensuremath{\bigcirc}}

\geometry{a4paper, scale=0.8}
\begin{document}
\section*{内容简介}

本书为“十四五”普通高等教育本科国家级规划教材，第一版荣获陕西省高等学校优秀教材一等奖，第二版荣获首届全国教材建设奖全国优秀教材二等奖。此次修订是依据教育部高等学校机械基础课程教学指导分委员会制定的《工程材料与机械制造基础课程教学基本要求》，并结合西北工业大学多年来的教学改革实践经验和本书使用过程中发现的问题而展开的。

本书集工程材料与机械制造工艺为一体，内容包括材料的种类与性能、材料的组织结构、金属热处理及表面改性、钢铁材料及其用途、非铁材料、失效及选材、铸造、压力加工、焊接、非金属及复合材料成形方法、毛坯成形方法选择及质量控制、现代成形技术及发展趋势、切削加工基础知识、零件表面的加工方法、机械零件的结构工艺性、机械加工工艺过程的基础知识、现代制造技术及发展趋势等内容。每章附有复习思考题，便于读者理解和巩固所学知识。

为便于教学和学生自学，本书配有工程材料与机械制造基础电子教案，书中附有二维码，可扫码观看知识点演示动画。

本书可作为普通高等学校机械类及近机械类专业的教材，也可供高等职业学校、高等专科学校、继续教育学校选用，亦可供有关工程技术人员参考。

\section*{第3版前言}

当前，新一轮科技革命与产业变革正加速演进，以科技创新为主导的新质生产力加速形成，这为制造强国建设注入强劲动能。但是，我国在部分核心技术领域仍面临严峻的国际挑战，高端装备、先进材料等领域的国产化替代亟待突破，这对掌握扎实基础理论与具有实践创新能力的工程人才提出了迫切需求。习近平总书记关于科技创新与教育强国建设的重要论述，为技术创新、产业升级和具有国际化视野的工程创新人才培养指明了方向，也为工程材料与机械制造相关课程的教学改革提供了根本遵循。

本次修订秉持“夯实基础、对接产业、赋能创新”的编写理念，结合编者多年来的创新教学实践经验及本书使用过程中发现的问题，对传统教学内容进行了系统性优化。

本书是国家精品课程、国家一流本科课程“工程材料与机械制造基础”配套教材；第一版为普通高等教育“十一五”国家级规划教材，荣获陕西省优秀教材一等奖；第二版荣获首届全国教材建设奖全国优秀教材二等奖。

本书集工程材料、材料成形基础和机械加工工艺基础三部分内容于一体，形成完整知识体系，有助于学生深入了解零件从材料选择、毛坯成形到加工成品的全过程及其内在联系，解决学时减少、知识量增多所带来的突出矛盾。本书可满足新形势下机械类和近机械类专业相应课程教学的需求，适应多学科知识交叉与渗透的发展趋势。

本次修订工作的要点包括：

(1)依据最新国家标准与行业标准，对全书内容进行了系统性核对与修订；

(2)更新和增加了部分章节的复习思考题；

(3)更新了部分插图，以加强图文呼应，增强表达效果；

(4)配套数字化资源（包括三维动画等），以帮助读者理解所学知识；

(5)增加了近年来工程材料成形与机械制造领域的前沿技术，介绍了人工智能在先进机械制造领域的应用；

(6)面向“总师型”人才培养需求，新增发动机叶片、飞机起落架等特色制造案例，完整覆盖材料选型、坯料成形到切削加工的全流程，在提升学生工程实践能力的同时，强化价值引领。


本书由国家级教学名师、西北工业大学齐乐华担任主编并统稿，西北工业大学罗俊、周计明、杨方任副主编。参与本次修订的人员有齐乐华、罗俊、周计明、付佳伟、田文龙、晁许江、赵志龙。

国家级教学名师、清华大学傅水根教授悉心审阅了本书，提出了许多宝贵的意见和建议。编者在此表示衷心感谢。

本书在编写过程中得到西北工业大学机械制造基础课程组教师和研究生的帮助，编者在此表示衷心感谢。

限于编者的水平，书中难免存在不妥之处，敬请广大同行及读者批评指正。邮箱：qilehua@nwpu.edu.cn。

\begin{flushright}
编者 \\
2025年11月30日
\end{flushright}

\tableofcontents

\chapter{材料的种类与性能}

\section{材料的种类}

材料是工业和科学技术的物质基础，是人类社会赖以生存和发展的重要条件，是衡量一个国家经济实力与技术水平的重要标志，它与信息、能源并列为现代化技术的三大支柱，而能源和信息的发展又依托于材料，因此世界各国均把对材料的研究开发放在突出的地位。

材料品种繁多，性能各异，而工程材料主要指用于机械工程和建筑工程等领域的材料。常见的工程材料按组成分类如图1-1所示。

\begin{figure}[H]
    \centering
    % Placeholder for Figure 1-1
    \caption{常见工程材料的分类}
\end{figure}

金属材料包括钢铁合金和非铁金属（又称有色金属）。由于金属材料具有良好的力学性能、物理性能、化学性能及工艺性能，能采用比较简便和经济的工艺方法制成零件，因此金属材料是目前应用最广泛的材料。

高分子材料包括塑料、橡胶等。其具有原料丰富、成本低、加工方便等优点，因此发展极其迅速，目前已在国民经济与日常生活的众多领域得到广泛应用。

无机非金属材料主要是陶瓷材料、水泥、玻璃和耐火材料。它具有不可燃烧性、高耐热性、高化学稳定性、耐老化性以及高硬度和良好的耐压性，且原料丰富，受到材料工作者和特殊行业的广泛关注。

复合材料是由基体材料（树脂、金属、陶瓷）和增强剂（颗粒、纤维、晶须）复合而成的。它既有组成材料的特性，又具有组成后的新特性，且其力学性能和功能可以根据使用需要进行设计、制造，所以自1940年玻璃钢问世以来，复合材料的应用领域一直在迅速扩大，品种、数量和质量也有了飞速发展。

其他材料主要包括超导材料、贮氢材料、形状记忆合金、非晶态材料、磁性材料与压电材料等，这些材料具有某种特殊物理性能、化学性能、生物性能以及某些功能之间可以相互转化的材料，它是在传统的金属、无机和有机材料基础上发展起来的新型材料，是材料高性能化、功能化和复合化的产物。

\section{材料的性能}

工程材料的性能分为使用性能和工艺性能。使用性能是指在服役条件下能保证安全可靠工作所必备的性能，包括材料的力学性能（也称机械性能）、物理性能和化学性能等。工程材料的使用性能决定了它的使用范围和寿命。对绝大多数工程材料而言，力学性能是最重要的使用性能。工艺性能是指材料的可加工性，包括锻造性能、铸造性能、焊接性能、热处理及切削加工性能等。在设计零件和选择工艺方法时，都要考虑材料的工艺性能。工艺性能将在后面有关章节中进行讨论。

\subsection{静载荷作用下材料的力学性能}

静载荷（简称静载）是指对试样进行缓慢加载。最常用的静载试验有拉伸、压缩、硬度、弯曲、扭转等，可利用不同的试验方法测得材料的各种力学性能指标，本书主要讨论工程领域应用最为广泛的强度、塑性和硬度指标。

1. 强度

强度是指在外力作用下材料抵抗变形和断裂的能力。强度有多种判据，工程上以抗拉强度和屈服强度最为常用。强度指标常通过拉伸试验测定。图1-2a为拉伸试样示意图，展示了退火低碳钢拉伸试样在拉伸前后的形貌变化，图1-2b和图1-2c所示分别为拉伸试验载荷（拉力）与变形量（伸长量）的变化图（拉伸曲线）及其应力-应变曲线。如果将拉（外）力 $F$ 除以试样的原始横截面积 $S_0$，则得到拉应力 $R$（单位横截面上的拉力），单位为 MPa；将伸长量 $\Delta L$ 除以试样的原始标距 $L_0$，则得到应变 $\varepsilon$（单位长度的伸长量）。根据 $R$ 和 $\varepsilon$，可以画出应力-应变曲线（图1-2c）。应力-应变曲线不受试样尺寸的影响，可以从曲线上直接读出材料的一些常规力学性能指标。静载拉伸下材料的力学性能指标主要有以下几个。



(1) 弹性和刚度

在 $R-\varepsilon$ 曲线（图1-2c）的开始段（$e$ 点以前）为直线，在该段产生的变形是可以恢复的（即中途卸除载荷，试样即恢复原状），称为弹性变形。工程上曾经使用 $\sigma_e$ 表示弹性极限（弹性变形阶段的最大应力），单位为 MPa。采用普通测量方法很难测出准确、唯一的弹性极限值。在实际工程中,弹性极限与条件屈服强度在本质上是一样的,均为材料开始产生微量塑性变形时的应力值,只不过所规定的微量塑性变形量不同。国家标准 GB/T 228.1—2021 规定弹性极限为塑性变形量为 $0.01\%$ 时的应力值,即 $R_{r0.01}$。

\begin{figure}[H]
	\centering
	% Placeholder for Figure 1-1
	\includegraphics[page=1, viewport={79.0 450.3 426.0 748.3}, clip]{12-07572-001FigRevised.pdf}
	\caption{常见工程材料的分类}
\end{figure}

材料在弹性变形范围内应力与应变的比值称为弹性模量,以 $E$ 表示,单位为 $\text{GPa}$,即
\[
E = \frac{R}{\varepsilon}
\]
弹性模量 $E$ 表征材料产生弹性变形的难易程度。弹性模量在工程上称为材料的刚度。材料在使用过程中,如刚度不足,会因发生过大的弹性变形而失效。弹性模量主要取决于材料内部原子间的作用力,如晶体材料的晶格类型、原子间距,各种强化手段对弹性模量的影响较小。

(2) 屈服强度 $R_e$

在 $R-\varepsilon$ 曲线(图 1-2c)中,当应力值超过 $e$ 点时,试样将产生塑性变形;当应力增至 $H$ 点时,试样开始产生明显的塑性变形,随后在曲线上出现水平波折线,表明即使外力不增加,试样仍继续塑性伸长,这种现象称为屈服。发生屈服所对应的应力值即屈服强度,用 $R_e$ 表示,单位为 $\text{MPa}$。屈服强度包括下屈服强度 $R_{eL}$ (在屈服期间的最低应力值)和上屈服强度 $R_{eH}$ (发生屈服而首次下降前的最高应力值)。

对于没有明显屈服点的材料,规定产生 $0.2\%$ 的塑性变形时的应力值为该材料的屈服强度,称为名义(条件)屈服强度,以 $R_{r0.2}$ 表示。工程中大多数零件都是在弹性范围内工作,如果产生过量的弹性变形就会使零件失效,所以屈服强度是零件设计和选材的主要依据之一。

(3) 抗拉强度 $R_m$

抗拉强度是试样拉断前所能承受的最大应力值,即试样所能承受的最大载荷除以原始横截面积,用 $R_m$ 表示,单位为 $\text{MPa}$,即
\[
R_m = \frac{F_m}{S_0}
\]
式中: $F_m$ ——试样所能承受的最大载荷, $\text{N}$。
$S_0$ ——试样的原始横截面积, $\text{mm}^2$。

抗拉强度表征的是材料抵抗断裂的能力,是设计和选材的主要依据之一。特别是对于低塑性或脆性材料,抗拉强度更应作为主要设计指标。

2. 塑性

材料断裂前产生塑性变形的能力称为塑性。塑性以材料断裂时的最大相对塑性变形量来表征。拉伸时用断后伸长率 $A$ 和断面收缩率 $Z$ 表示,两者均为量纲为一的量。

(1) 断后伸长率 $A$

断后伸长率 $A$ 表示试样断裂时的相对伸长量,即
\[
A = \frac{L_u - L_0}{L_0} \times 100\%
\]
式中: $L_0$ ——试样的原始标距, $\text{mm}$;
$L_u$ ——试样的断后标距, $\text{mm}$。

(2) 断面收缩率 $Z$

断面收缩率 $Z$ 表示试样断裂后截面的相对收缩量,即
\[
Z = \frac{S_0 - S_u}{S_0} \times 100\%
\]
式中: $S_0$ ——试样的原始横截面积, $\text{mm}^2$;
$S_u$ ——试样断裂处的最小横截面积, $\text{mm}^2$。

$A, Z$ 愈大,材料的塑性愈好。

3. 硬度

材料抵抗硬物压入其表面的能力称为硬度,也可将硬度看作材料抵抗局部塑性变形的能力。硬度愈高,表明材料抵抗局部塑性变形的能力愈大。

在材料制成的半成品和成品的质量检验中,硬度是标志产品质量的重要依据。常用的硬度有布氏硬度、洛氏硬度和维氏硬度等。

(1) 布氏硬度

用一定的载荷 $F$ 将直径为 $D$ 的碳化钨合金球压入试样表面,保持一定时间后卸除载荷,在读数显微镜下测量试样表面压痕直径 $d$,并根据所测直径查表,即可得到硬度值,用 $\text{HBW}$ 表示,单位为 $\text{MPa}$,一般不予标出。显然,材料越软,压痕直径越大,布氏硬度值越低;反之,则布氏硬度值越高。

布氏硬度试验的优点是压痕面积大,不受微小不均匀硬度的影响,试验数据稳定,重复性好,但不适用于成品零件和薄壁件的硬度检验。

(2) 洛氏硬度

洛氏硬度是将标准压头(锥顶角为 $120^\circ$ 的金刚石圆锥或一定直径的碳化钨合金球)压入试样表面,保持一定时间后卸除载荷,根据残余压痕深度确定试样的硬度,可直接由硬度计刻度盘读出。

国家标准规定,根据所加载荷和压头的不同,洛氏硬度有十五种标尺,其中以 $\text{HRA}$、$\text{HRB}$、$\text{HRC}$ 三种最为常用,尤以 $\text{HRC}$ 应用最多。

洛氏硬度的测定操作迅速、简便,压痕面积小,适用于成品检验。但由于接触面积小,当硬度不均匀时,数值波动较大,需多测几个点,取其平均值。

(3) 维氏硬度

维氏硬度是将正四棱锥体的金刚石压头压入试样表面,再测量压痕对角线长度,根据所加压力和对角线平均长度查表得到硬度值,用符号 $\text{HV}$ 表示。

维氏硬度压痕深度浅,测量精确,硬度测量范围大,尤其能很好地测量薄试样的硬度。所加载荷较小的维氏硬度,又称为显微硬度(用 $\text{HM}$ 表示),可表示金相组织中不同相的硬度。缺点是需要在显微镜下测量压痕尺寸,工作效率较低。

\subsection{其他载荷作用下材料的力学性能}

1. 冲击韧性

材料抵抗冲击载荷的能力称为冲击韧性。

测定冲击韧性常用一次摆锤冲击试验来测定,摆锤冲断试样所做的冲击吸收功 $A_K$ 与试样缺口处横截面积 $S$ 的比值,即材料的冲击韧度,一般用 $a_K$ 表示,单位为 $\text{J/m}^2$。

实际上,在冲击载荷作用下的机械零件,很少是受大能量的一次冲击而发生破坏的,往往是经受小能量的多次冲击,因冲击损伤积累引起的裂纹扩展而发生断裂的。

2. 断裂韧性

材料抵抗裂纹(失稳)扩展的能力称为断裂韧性,其量化数值一般用断裂韧度 $K_{Ic}$ 表示,单位为 $\text{MPa} \cdot \text{m}^{1/2}$。

实际使用的材料中常常存在一定的缺陷,如夹杂物、气孔、微裂纹等,这些缺陷都可看作裂纹。在外力作用下,这些裂纹扩展的难易程度用断裂韧性表示。断裂韧性与材料本身的成分、组织和结构有关。

3. 疲劳强度

许多机器零件如轴、齿轮、弹簧、活塞杆等,均是在交变载荷下工作的,它们工作时所承受的应力一般低于材料的屈服强度。零件在这种交变动载荷作用下,经过较长时间的工作而发生断裂的现象称为疲劳。因此,疲劳是机械零件在循环或交变应力作用下,经过一段时间产生失效的现象。疲劳断裂往往无先兆,会突然断裂,因此危害很大。疲劳强度就是用来表征材料抵抗疲劳能力的。

疲劳强度是通过测定材料在交变载荷[钢的交变(循环)次数为 $10^6 \sim 10^7$、非金属的为 $10^7 \sim 10^8$]作用下而无断裂的最大应力得到的,用 $R_r$ 表示,单位为 $\text{MPa}$。

4. 磨损

机器运转时,任何零件在进行接触状态下的相对运动时都会产生摩擦,导致磨损,并最终失效。按磨损的破坏机理,磨损可分为黏着磨损、磨粒磨损和腐蚀磨损。

(1) 黏着磨损

黏着磨损又称为咬合磨损,其实质是接触面在接触压力作用下局部发生黏着(冷焊),在相对运动时黏着处分离,使接触面上有小颗粒被拉拽出来,反复黏着和分离,造成黏着磨损。

(2) 磨粒磨损

磨粒磨损是当摩擦副一方的硬度比另一方大得多,或者在接触面之间存在着硬质粒子时所产生的磨损,其特征是接触面上有明显的切削痕迹。

(3) 腐蚀磨损

腐蚀磨损是由于外界环境引起金属表面的腐蚀产物剥落,与金属表面之间的机械磨损(磨粒、黏着)相结合而出现的磨损。

\subsection{材料的高温性能和低温性能}

1. 高温性能

材料在长时间的恒温、恒应力作用下,发生缓慢塑性变形的现象称为蠕变。蠕变的一般规律是温度越高,工作应力越大,则蠕变的发展越快,产生断裂前的工作时间就越短。

蠕变的另一种表现形式是应力松弛。它是指承受弹性变形的零件,在工作过程中总变形量保持不变,但随时间的延长,工作应力自行逐渐衰减的现象,如高温紧固件因应力松弛使紧固失效。

2. 低温性能

随着温度的下降,多数材料会出现脆性增加的现象,严重时甚至发生脆断。通常,可通过材料的吸收冲击功 $A_K$ 与温度的变化关系来确定材料的韧、脆状态转化。当温度降低至某一值时, $A_K$ 值会急剧减小,使材料呈脆性状态。材料由韧性状态转变为脆性状态的温度 $T_K$ 称为冷脆转化温度。材料的 $T_K$ 低,表明其低温韧性好。

\subsection{材料的化学性能}

化学性能是指材料在室温或高温时抵抗各种介质化学侵蚀的能力。通常,将材料因化学侵蚀而损坏的现象称为腐蚀。非金属材料的耐腐蚀性远高于金属材料。金属腐蚀既容易造成一些隐蔽性和突发性的严重事故,也损失了大量的金属材料。金属腐蚀主要有化学腐蚀和电化学腐蚀,金属腐蚀的绝大多数是由电化学腐蚀引起的。为了提高金属的耐腐蚀能力,原则上应保证以下3点:一是尽可能使金属保持均匀的单相组织,即无电极电位差;二是尽量减小两极之间的电极电位差,并提高阴极的电极电位,以减缓腐蚀速度;三是尽量不与电解质溶液接触,减小甚至隔断腐蚀电流。

工程上常用的防腐蚀方法有以下几种。

1) 改变金属的化学成分,提高合金的耐腐蚀性,如不锈钢以及材料表面渗铬、渗铝处理等。

2) 通过覆盖法将金属同腐蚀介质隔离,如金属表面镀层、覆层和发蓝等。

3) 改善腐蚀环境,如干燥气体封存和密封包装等。

4) 阴极保护法,即将被保护的金属作为原电池的阴极,牺牲阳极金属,使阴极金属不遭受腐蚀的方法,或用外加电流法,保护阴极金属。

\section*{复习思考题}
 
1-1 简述金属材料的种类及其特点。

1-2 金属材料的使用性能包括哪些?

1-3 什么是金属的力学性能?它包括哪些主要指标?

1-4 一根直径为 $\phi 10 \text{ mm}$ 的钢棒,在拉伸断裂时直径变为 $\phi 8.5 \text{ mm}$,此钢的抗拉强度 $R_m = 450 \text{ MPa}$,问此棒能承受的最大载荷为多少?断面收缩率是多少?

1-5 简述洛氏硬度的测试原理。

1-6 简述冲击韧性的测试原理。

1-7 什么是蠕变和应力松弛?

1-8 什么是材料的高温性能?

1-9 功能材料有哪些种类?其特点是什么?

1-10 简述磨损的分类形式及含义。

1-11 金属腐蚀的方式主要有哪几种?金属防腐蚀的方法有哪些?

1-12 何谓疲劳现象?为什么说疲劳破坏比其他形式的破坏的危害性更大?

1-13 试比较布氏硬度和洛氏硬度试验的特点与应用范围。


\chapter{材料的组织结构}

\section{金属的晶体结构}

\subsection{金属的理想晶体结构}

固体物质按其原子(或分子)的聚集状态分为晶体和非晶体两大类。晶体是原子(或分子)在三维空间做有规律的周期性重复排列的固体,非晶体是由原子(或分子)无规则地堆砌而成的。常用的固态金属与合金都是晶体,而普通玻璃、沥青、松香等物质是非晶体。晶体具有固定的熔点、规则的几何外形和各向异性的特性;非晶体没有固定的熔点,且各向同性。

晶体中的原子是按一定规则排列的(图 2-1a)。为便于理解和描述,常用一些假想的连线将各原子的中心连接起来,把原子看作一个点,这样形成的几何图形称为晶格,如图 2-1b 所示。晶体是周期性重复排列的,通常取出晶格中的一个基本单元——晶胞(图 2-1c)来描述晶体构造。晶胞的各棱边长度称为晶格常数。在晶格中由一系列原子组成的平面称为晶面,而晶面又是由一行行的原子列组成的,晶格中各原子列的位向称为晶向。

\begin{figure}[H]
    \centering
    % Placeholder for Figure 2-1
    \includegraphics[page=1, viewport={94.5 303.3 396.5 419.3}, clip]{12-07572-001FigRevised.pdf}
    \caption{晶体结构示意图}
    \label{fig:2-1}
\end{figure}

工业上常用的金属中,除少数具有复杂的晶体结构外,绝大多数金属都具有比较简单的晶体结构,其中最常见的金属晶体结构(晶格类型)有体心立方晶格(bcc)、面心立方晶格(fcc)和密排六方晶格(hcp)3种类型,如图 2-2 所示。

(1) 体心立方晶格

体心立方晶格的晶胞如图 2-2a 所示,在立方体的8个顶角和立方体中心各有1个原子。具有这种晶格的金属有铬、钨、钼、钒、铌和 $912 \text{ }^\circ\text{C}$ 以下的铁等。

\begin{figure}[H]
    \centering
    % Placeholder for Figure 2-2
    \includegraphics[page=1, viewport={63.5 170.3 482.0 272.3}, clip]{12-07572-001FigRevised.pdf}
    \caption{典型晶格的晶胞示意图}
    \label{fig:2-2}
\end{figure}

(2) 面心立方晶格

面心立方晶格的晶胞如图 2-2b 所示,在立方体的8个顶角上和6个面的中心各有1个原子。具有这种晶格的金属有铝、铜、镍、铅、金、银和 $912 \sim 1394 \text{ }^\circ\text{C}$ 的铁等。

(3) 密排六方晶格

密排六方晶格的晶胞如图 2-2c 所示,在六棱柱的上、下六角形面的顶角和面的中心各有1个原子,在六棱柱中间还有3个原子。具有这种晶格的金属有镁、锌、镉和铍等。

\subsection{金属的实际晶体结构}

以上所讨论的金属的晶体结构是理想结构,是由原子排列位向或方式完全一致的晶格组成的,称为单晶体。而实际工程上使用的材料绝大多数是多晶体,是由很多个小的单晶体组成的,这些单晶体称为晶粒,每个晶粒的原子排列位向是不同的,如图 2-3 所示。研究结果还发现,即使在一个晶粒内,实际金属的结构与理想的状态也有差异。因此,在实际金属中或多或少地存在着偏离理想结构的微观区域,通常把这种偏离晶体理想结构的微观区域称为晶体缺陷。这些缺陷有点缺陷、线缺陷和面缺陷3种类型。

1. 点缺陷——空位和间隙原子

如图 2-4 所示,空位是指未被原子占据的晶格节点。间隙原子是指位于晶格间隙中的原子。它们主要是由于在结晶过程中原子堆积不完善、外来原子溶入等因素所造成的,或已形成的晶体在高温、冷变形加工、高能粒子轰击、氧化等作用下产生的。点缺陷的存在,使原子之间的作用力失去平衡,其周围的原子发生靠拢或撑开,使晶体结构的规律性遭到破坏,晶格发生歪扭,即晶格畸变。

\begin{figure}[H]
    \centering
    % Placeholder for Figure 2-3
    \includegraphics[page=2, viewport={93.0 588.8 243.5 736.8}, clip]{12-07572-001FigRevised.pdf}
    \caption{多晶体中不同位向晶粒示意图}
    \label{fig:2-3}
\end{figure}

\begin{figure}[H]
    \centering
    % Placeholder for Figure 2-4
    \includegraphics[page=2, viewport={262.5 587.3 396.5 744.8}, clip]{12-07572-001FigRevised.pdf}
    \caption{空位和间隙原子}
    \label{fig:2-4}
\end{figure}

2. 线缺陷——位错

位错是指在晶体中,某处有一列或若干列原子发生了某种有规律的错排现象。图 2-5 所示为刃型位错。

\begin{figure}[H]
    \centering
    % Placeholder for Figure 2-5
    \includegraphics[page=2, viewport={60.0 412.8 318.0 551.3}, clip]{12-07572-001FigRevised.pdf}
    \caption{刃型位错示意图}
    \label{fig:2-5}
\end{figure}

刃型位错可以描述为在一个完整晶体的某个晶面上,多出了半个原子面,这多余的半个原子面犹如刀刃一般地垂直切入(图 2-5a 中的 $EFGH$ 原子面),从而使晶体中某一晶面的上、下部分晶体产生了错排现象,故称之为刃型位错。刃边 $EF$ 为刃型位错线,其原子排列模型如图 2-5b 所示。

3. 面缺陷——晶界

晶界是位向不同、相邻晶粒之间的过渡层,如图 2-6 所示。它可看作由空位、间隙原子以及位错堆积起来的一种面层状晶体缺陷,也是各种杂质原子聚集的场所。

晶体缺陷的存在破坏了晶体的完整性,使晶格产生畸变,晶格能量增加,因而晶体缺陷相对于完整的晶体来说处于一种不稳定状态。它们在外界条件(温度、外力等)变化时会发生原子的扩散与迁移,从而引起金属某些性能发生变化。

一般情况下,晶体缺陷的存在虽然可以提高金属的强度,但是常常降低金属的耐腐蚀性能,所以可以通过腐蚀观察金属的各种缺陷。

\begin{figure}[H]
    \centering
    % Placeholder for Figure 2-6
    \includegraphics[page=2, viewport={343.0 414.8 485.0 540.3}, clip]{12-07572-001FigRevised.pdf}
    \caption{晶界示意图}
    \label{fig:2-6}
\end{figure}

\subsection{金属材料的结构特点}

金属材料主要由金属晶体组成,对纯金属而言,其结构主要是指不同类型晶体的显微组织形态和缺陷状态,包括晶粒的形状、大小,晶格的畸变、位错密度等。

机械工业中使用的金属材料主要是合金。合金是指由两种或两种以上的金属元素或金属元素与非金属元素组成的、具有金属特性的物质。组成合金最基本的独立物质称为组元。组元可以是金属元素、非金属元素或稳定的化合物。由两个组元组成的合金称为二元合金。根据组元数的多少,合金可分为二元合金、三元合金等。一系列相同组元按不同比例组成的合金称为合金系。

实践证明,在纯金属中加入适量的合金元素,可显著改变其性能。采用合金元素来改变金属性能的方法称为合金化。金属经合金化后,其性能发生明显变化的原因是其显微组织发生了明显的变化。组织是观察到的在金属及合金内部组成相的大小、方向、形状、分布及相互结合状态。仅有一种相组成的组织为单相组织,由两种或两种以上的相组成的组织为多相组织。如在铝中加入$11.7\%$硅构成的合金,其显微组织由两种基本组成物组成,即白色的基底上分布着黑色针状物,如图2-7所示。白色部分的化学成分和晶体结构是一样的,黑色针状部分是另一种化学成分和组织结构。通常,将合金中这种具有相同化学成分、相同晶体结构和相同物理或化学性能并与该系统的其余部分以界面相互隔开的均匀组成部分称作相。常见的合金中存在的相可以归纳为固溶体和金属化合物两大类。

1. 固溶体

固溶体是指溶质组元溶入溶剂晶格中而形成的单一的均匀固体。这与溶液的概念十分相似,只不过前者为固相,后者为液相。固溶体的特点是保持溶剂的晶格结构,但会引起溶剂晶格不同程度的畸变。
固溶体按照溶质原子在溶剂晶格中的位置可以分为置换固溶体和间隙固溶体,如图2-8所示。置换固溶体是溶质原子取代了溶剂晶格中某些节点上的原子,而间隙固溶体是溶质原子嵌入溶剂晶格的间隙中,不占据晶格节点位置。固溶体还可以按固溶度的大小分为有限固溶体和无限固溶体,按溶质原子在溶剂晶格中的分布特点分为无序固溶体和有序固溶体。

\begin{figure}[H]
    \centering
    % Placeholder for Figure 2-7
    \includegraphics[page=2, viewport={56.5 248.3 226.5 373.3}, clip]{12-07572-001FigRevised.pdf}
    \caption{铝硅合金的显微组织($400\times$)}
    \label{fig:2-7}
\end{figure}

\begin{figure}[H]
    \centering
    % Placeholder for Figure 2-8
    \includegraphics[page=2, viewport={263.0 236.8 511.0 380.8}, clip]{12-07572-001FigRevised.pdf}
    \caption{固溶体的类型}
    \label{fig:2-8}
\end{figure}

固溶体中溶质原子的溶入引起晶格畸变,使晶体处于高能状态,从而提高合金的强度和硬度,这种现象称为固溶强化。固溶强化使固溶体的强度、硬度比溶剂有所提高,但塑性和韧性下降。

2. 金属化合物

两组元A和B组成合金时,除了可形成以A或以B为基的固溶体外,还可能A和B相互作用生成化合物形成新相,这种新相一般可用化学式$\text{A}_m\text{B}_n$表示。它的晶格类型与两组元完全不同,性能差别也很大,一般都具有高熔点、高硬度,但很脆。合金中以固溶体为主,适量的金属化合物弥散分布可提高其强度、硬度及耐磨性能,所以常用金属化合物来强化合金。这种强化方式称第二相质点强化或弥散强化,它是各类合金钢及非铁金属的重要强化方法。

\section{非金属材料的结构}
非金属材料种类很多,目前在机械工程中常用的是高分子材料和陶瓷材料。

\subsection{陶瓷材料的结构特点}
1. 陶瓷材料的键合类型

陶瓷材料的结合键主要为离子键和共价键。例如$\text{Al}_2\text{O}_3$、$\text{MgO}$为离子键,金刚石、$\text{SiC}$为共价键。但陶瓷材料通常不是单一的键合类型,而是两种或两种以上的混合键。例如$\text{MgO}$晶体,离子键占$84\%$,共价键占$16\%$。混合键是陶瓷晶体的特点之一。
由于陶瓷材料具有键能高的结合键,所以陶瓷材料通常有熔点高、硬度高、耐腐蚀、塑性差等特性。例如$\text{Mg}$的熔点为$650\ ^{\circ}\text{C}$,而燃烧生成的$\text{MgO}$由于含离子键,其熔点高达$2\ 800\ ^{\circ}\text{C}$。又如金刚石是典型的共价键,其熔点高达$3\ 750\ ^{\circ}\text{C}$,是目前自然界中最坚硬的固体。如$\text{SiC}$、$\text{Si}_3\text{N}_4$等具有共价键的其他固体材料,均具有高熔点、高硬度等性质。

2. 陶瓷材料的组织

和金属材料一样,陶瓷材料的性能也是由其化学组成和结构所决定的。相对于金属材料而言,陶瓷材料的组成结构更复杂,一般由晶体相、玻璃相和气相组成,如图2-9所示。各相的组成、结构、数量、形状与分布都对陶瓷的性能有直接的影响。

\begin{figure}[H]
    \centering
    % Placeholder for Figure 2-9
    \includegraphics[page=2, viewport={14.0 95.3 184.5 206.8}, clip]{12-07572-001FigRevised.pdf}
    \caption{陶瓷显微组织示意图}
    \label{fig:2-9}
\end{figure}

(1) 晶体相

晶体相是一些化合物或以化合物为基的固溶体,是决定陶瓷材料物理、化学和力学性能的主要组成物。主要晶体相有氧化物和硅酸盐。陶瓷材料是多晶体,同金属一样,有晶粒和晶界。在一个晶粒内,也有线缺陷(位错)和点缺陷(空位和间隙原子),这些晶体缺陷的作用亦类同金属晶体中的缺陷。如果陶瓷的晶粒细小,晶界总面积大,则陶瓷材料的强度大;空位和间隙原子可加速陶瓷在烧结工艺过程中的扩散,并且也影响其物理性能。但是,陶瓷在常温下几乎没有塑性。

(2) 玻璃相

陶瓷制品在烧结过程中,作为主要原料的$\text{SiO}_2$已处在熔化状态,但在熔点附近$\text{SiO}_2$的黏度很大,原子迁移困难,所以当液态$\text{SiO}_2$冷却到熔点以下时,原子不能排列成长为有序(晶体)状态,而形成过冷液体。当过冷液体继续冷却到玻璃化转变温度时,凝固为非晶态的玻璃相。玻璃相的结构是由离子多面体构成的空间网络,呈不规则排列。
玻璃相在陶瓷组织中的作用是黏结分散的晶体相,降低烧结温度,抑制晶体长大和填充空隙等。玻璃相的熔点低、热稳定性差,使陶瓷容易在高温下产生蠕变,从而降低高温下的强度。所以,工业陶瓷须控制陶瓷组织中玻璃相的质量分数,一般陶瓷中玻璃相约占$30\%$。

(3) 气相

气相是指陶瓷组织中的气孔。气孔可以是封闭型的,也可以是开放型的(即气孔通向陶瓷的表面),可以分布在晶粒内(封闭型),也可分布在晶界上,甚至分布在玻璃相中。气孔在陶瓷组织中的比例约占$5\%$或更高。气孔会造成应力集中,使陶瓷容易开裂,降低强度。气孔还降低陶瓷抗电击穿能力,同时对光线有散射作用,故降低陶瓷的透明度。当陶瓷要求密度小、重量轻或者要求绝热性高时,应保留较多的气相。

\subsection{高分子材料的结构特点}
高分子材料的主要组分是高分子化合物。高分子化合物是由许多小分子(或称低分子)通过共价键连接起来形成的具有链状结构的大分子有机化合物,称为高聚物或聚合物,其结构称为大分子链。

高分子材料除以高聚物为主要组分外,还常包含各种添加剂,如塑料中的填料、增塑剂、固化剂等。虽然这些组分也会影响材料的性能,但相对而言,毕竟不像高聚物那样起关键作用。所以,高分子材料的结构主要是指高聚物的结构。

高聚物是由小分子化合物聚合而成的。可以聚合生成大分子链的小分子化合物称为单体。如聚乙烯是由乙烯聚合而成的,乙烯就是聚乙烯的单体,其聚合反应式为
\[ n(\text{CH}_2=\text{CH}_2) \xrightarrow{\text{加聚}} [\text{CH}_2-\text{CH}_2]_n \]
所以,单体是人工合成高聚物的原料。

需要说明的是,高聚物还可由两种或两种以上的单体共同聚合而成。但是,并非任何一种小分子化合物都可以作为单体。

大分子链的重复结构单元称为链节,如上述聚乙烯中的$[\text{CH}_2-\text{CH}_2]$即聚乙烯大分子链的链节。一个大分子链中链节的数量称为聚合度,上式中的$n$即聚乙烯大分子的聚合度。聚合度反映了大分子链的长短及相对分子质量的大小。

1. 大分子链的形态

高聚物的结构形式按其大分子链的几何形状,可分为线型结构和体型结构两种。

(1) 线型结构

线型结构是由许多链节连成一条长链,如图2-10a所示。通常卷曲成不规则的线团状,在拉伸时呈直线状。

\begin{figure}[H]
    \centering
    % Placeholder for Figure 2-10
    \includegraphics[page=2, viewport={197.5 95.3 532.0 197.3}, clip]{12-07572-001FigRevised.pdf}
    \caption{高聚物结构示意图}
    \label{fig:2-10}
\end{figure}

有一些高聚物的大分子链带有一些小的支链,如图2-10b所示,它也属于线型结构。具有线型结构的高聚物在加工成形时,分子链时而卷曲收缩,时而伸长,表现出良好的塑性和弹性。在适当溶剂中能溶解或溶胀,加热可软化或熔化,冷却后变硬,并可反复进行,故易于加工成形,可重复使用。一些合成纤维和热塑性塑料(如聚氯乙烯、聚苯乙烯等)就属此类结构。支链的存在使线型高聚物的性能钝化,如熔点升高、黏度增加等。

(2) 体型结构

体型结构是大分子链之间通过支链或化学键连接成一体的交联结构,在空间呈网状,如图2-10c所示。体型结构高聚物就是一个由化学键固结起来的不规则网状分子,所以具有较好的耐热性、尺寸稳定性和机械强度,但弹性、塑性低,脆性大,不能塑性加工,成形加工只能在网状结构形成之前进行,材料不能反复使用。热固性塑料(如酚醛塑料、环氧塑料等)和硫化橡胶等就属此类结构。

2. 大分子的聚集态结构

一般低分子材料有气态、液态和固态三种,高聚物由于分子特别大且分子间力也大,容易聚集为液态或固态,而无气态。按大分子几何排列是否有序,固态高聚物的结构分为无定型和结晶型两种,结晶型高聚物的分子排列规整有序,无定型高聚物的分子排列杂乱不规则。

结晶型高聚物由晶区(分子有规则紧密排列的区域)和非晶区(分子处于无序状态的区域)所组成,如图2-11所示。由于分子链很长,所以在每个部分都呈现规则排列是很困难的。在高聚物中晶区所占的质量分数(或体积分数)称为结晶度。结晶型高聚物的结晶度一般为$50\% \sim 80\%$。

无定型高聚物的结构,并非真正是大分子排列呈杂乱交缠状态,而是呈远程(远距离范围内)无序、近程(近距离范围内)有序状态。

晶态与非晶态影响高聚物的性能。结晶使分子排列紧密,分子间作用力增大,所以使高聚物的密度、强度、硬度、刚度、熔点、耐热性、耐化学性、抗液体及气体透过性有所提高,而依赖分子链运动的有关性能,如弹性、塑性和韧性较低。

3. 高聚物的物理、力学状态

高聚物在不同的温度下呈现出不同的物理状态,因而具有不同的性能,这对高聚物的成形加工和使用具有重要意义。图2-12为线型无定型高聚物的温度-变形曲线。由图可见,温度不同,线型无定型高聚物可处于玻璃态、高弹态或黏流态。

\begin{figure}[H]
    \centering
    % Placeholder for Figure 2-11
    \includegraphics[page=3, viewport={120.0 605.3 271.0 739.8}, clip]{12-07572-001FigRevised.pdf}
    \caption{高聚物的晶区与非晶区示意图}
    \label{fig:2-11}
\end{figure}

\begin{figure}[H]
    \centering
    % Placeholder for Figure 2-12
    \includegraphics[page=3, viewport={294.0 604.3 429.0 742.3}, clip]{12-07572-001FigRevised.pdf}
    \caption{线型无定型高聚物的温度-变形曲线}
    \label{fig:2-12}
\end{figure}

(1) 玻璃态

在温度低于$T_{\text{g}}$时,高聚物处于玻璃态,$T_{\text{g}}$称为玻璃化温度。在玻璃态时,高聚物的大分子链热运动处于停止状态。玻璃态下表现出的力学性能与低分子材料相似,具有一定的刚度,是塑料的应用状态。作为塑料使用的高聚物,其$T_{\text{g}}$应高于室温,且越高越好。

(2) 高弹态

当温度处于玻璃化温度$T_{\text{g}}$和黏流化温度$T_{\text{f}}$之间时,高聚物处于高弹态。这时,高聚物的分子链动能增加,由几个或几十个链节组成的链段可进行内旋转运动,但整个分子链并没有移动。处于高弹态的高聚物在受外力作用时,原卷曲链沿受力方向伸展,产生很大的弹性变形(变形量可达$100\% \sim 1\ 000\%$)。高弹态是橡胶的应用状态,故作为橡胶使用的高聚物,其$T_{\text{g}}$应低于室温,且越低越好。

(3) 黏流态

当温度升到黏流化温度$T_{\text{f}}$时,大分子链可自由运动,高聚物成流动的黏液,这种状态称为黏流态。黏流态是高聚物成形加工的工艺状态。

若高聚物中有部分结晶区域,则当温度升到$T_{\text{g}}$以上和结晶体的熔点以下时,非结晶区仍保持线型无定型高聚物的高弹态,而结晶区则由于分子链无法产生内旋运动,表现出较高的刚度和硬度,两者复合组成一种既韧又硬的皮革态。部分结晶高聚物的这种特性,为通过调整、控制结晶度来改变高聚物性能提供了可能。

4. 高分子材料的老化

高分子材料在热、光、化学、生物和辐射等作用下会产生老化现象,其结构和性能发生很大的变化,如硬化、脆化、发软、发黏等。老化现象的实质是高分子材料的主要组分——大分子链的结构通过交联或降解发生变化。例如,许多橡胶在空气中氧的作用下(特别是当氧以臭氧的形式存在时)发生进一步的交联,会变硬、变脆而失去弹性。紫外线的照射会加速氧化的进程。

降解对高分子材料结构和性能的影响更为突出。所谓降解,是指聚合物在长期储存或使用过程中,其聚合度由于热、光、氧化、水解、生物和力学等的作用而降低的一种化学反应。降解使高分子材料软化、发黏,需长久使用的工程塑料应尽量避免降解。但若合理利用也有其有利的一面,特别是可利用降解作为手段达到减少环境污染的目的。

\section{金属的结晶与细晶强化}
广义上讲,金属从液态变为固态的过程称为结晶。从原子排列的情况来看,结晶就是原子从一种排列状态(晶态或非晶态)变为另一种规则排列状态的过程。本书主要介绍纯金属由液态到固态的转变。

1. 结晶过程

结晶过程是从液体中形成一些称为结晶核心(晶核)的细小晶体开始的,然后已形成的晶核按各自不同的位向不断长大。同时,在液体中又产生新的结晶核心并逐渐长大,直至液体全部消失,形成由许多位向不同、外形不规则的晶粒所组成的多晶体。晶体的形核和长大过程如图2-13所示。

\begin{figure}[H]
    \centering
    % Placeholder for Figure 2-13
    \includegraphics[page=3, viewport={56.0 437.3 491.0 571.3}, clip]{12-07572-001FigRevised.pdf}
    \caption{结晶过程示意图}
    \label{fig:2-13}
\end{figure}

2. 结晶温度

金属结晶时,都存在着一个平衡结晶温度$T_{\text{m}}$。这时,液体中的原子结晶到晶体上的数目,等于晶体上的原子溶入液体中的数目。宏观上,此时既不结晶,也不熔化,液体和晶体处于动平衡状态。只有冷却到低于平衡结晶温度时才能有效地结晶。因此,实际结晶温度$T_1$总是低于平衡结晶温度$T_{\text{m}}$。两者之差($T_{\text{m}} - T_1$)称为过冷度$\Delta T$(图2-14)。过冷度的大小与冷却速度有关,冷却速度越快,过冷度越大。

金属的实际结晶温度可用热分析法加以测定。将熔化的金属以缓慢的速度进行冷却,同时记录下温度随时间的变化规律,绘出如图2-14所示的冷却曲线。金属结晶时放出的结晶潜热,补偿了冷却时向外散出的热量,冷却曲线上暂时出现水平线段,即温度保持不变的恒温现象。该温度即实际结晶温度$T_1$。当散热极其缓慢,即冷却速度极其缓慢时,实际结晶温度与平衡结晶温度趋于一致。

\begin{figure}[H]
    \centering
    % Placeholder for Figure 2-14
    \includegraphics[page=3, viewport={6.0 233.8 144.0 338.8}, clip]{12-07572-001FigRevised.pdf}
    \caption{金属的冷却曲线}
    \label{fig:2-14}
\end{figure}

3. 晶核的形成与细晶强化

液态金属内部的原子并不是完全无规则排列的,往往在局部存在有规则排列的原子团,它们时聚时散,在一定的过冷度下,原子团便成为晶核而长大。当然,结晶也可依附于液体中的杂质而进行,前者(由金属自身原子团形成晶核)称为自发形核,后者(依附外来固体杂质形成晶核)称为非自发形核。

晶核一旦形成,便开始长大,由规则的外形逐渐长出棱角,棱角处散热条件优于其他部位而优先长大,形成如图2-13c所示的树枝状晶体(简称枝晶),多个晶核一起长大,直至枝晶间液体消耗完毕,形成多颗晶粒。

结晶条件不同,晶核长大后形成的晶粒大小差别也很大。一般情况下,金属的强度和硬度、塑性和韧性都会随晶粒的细化而提高,这种现象称为细晶强化。因此,在生产中常采取提高冷却速度和变质处理等措施来细化晶粒以改善力学性能。提高冷却速度可增大过冷度,使晶核生成速度大于晶粒长大速度,因而使晶粒细化。但提高冷却速度受铸件的大小、形状的限制。变质处理是在液态金属中加少量变质剂(又称孕育剂)作为人工晶核,以增加晶核数,从而使晶粒细化。此外,在结晶过程中采用机械振动、超声波振动和电磁振动,也有细化晶粒的作用。

与纯金属相比,合金的结晶过程比较复杂,在合金结晶的过程中,液相和固相的成分都通过原子的扩散而不断变化。随温度的降低,晶核首先在含高熔点组元较多的区域形成,并通过吸收高熔点组元而长大,使其周围液相中高熔点组元减少。结晶完成后,如果原子的扩散不充分,在先结晶处高熔点组元的质量分数较高,而后结晶处高熔点组元的质量分数较小,在合金晶体内部造成成分的不均匀,这种成分不均匀现象称为偏析。对于大多数合金,结晶时极易出现树枝状晶体,先结晶的枝晶主轴上,高熔点组元较多,后结晶的枝干上,高熔点组元较少,这种树枝状结晶造成的晶体内化学成分不均匀的现象称为枝晶偏析或晶内偏析。为了消除这种偏析,可将合金重新加热到稍低于结晶温度,并长时间地保温,使原子扩散充分进行,从而使成分均匀一致,这种处理称为“扩散退火”。

\section{材料的同素异构现象}

\subsection{晶体的同素异构}
某些金属,例如铁、锰、钛、锡、钴等,凝固后在不同的温度下有着不同的晶格形式,这种金属在固态下由于温度的改变而发生晶格改变的现象称为同素异构转变。这一转变与液态金属的结晶过程相似,也包括晶核的形成和晶粒的长大,故又称为二次结晶或重结晶。

图2-15所示为铁的同素异构转变冷却曲线,铁在$1\ 538 \sim 1\ 394\ ^{\circ}\text{C}$时为体心立方晶格,称为$\delta\text{-Fe}$;在$1\ 394 \sim 912\ ^{\circ}\text{C}$时为面心立方晶格,称为$\gamma\text{-Fe}$;在$912\ ^{\circ}\text{C}$以下为体心立方晶格,称为$\alpha\text{-Fe}$。

\begin{figure}[H]
    \centering
    % Placeholder for Figure 2-15
    \includegraphics[page=3, viewport={146.0 238.3 313.0 377.8}, clip]{12-07572-001FigRevised.pdf}
    \caption{铁的同素异构转变冷却曲线}
    \label{fig:2-15}
\end{figure}

一些无机非金属多晶材料中,也存在着同素异构转变,有时也称为同质多晶转变,一个典型的例子就是石英。石英随着温度的变化而出现同质多晶转变,其过程如下:
\[
\begin{array}{cccccc}
\alpha\text{-石英} & \xrightarrow{870^{\circ}\text{C}} & \alpha\text{-磷石英} & \xrightarrow{1\ 470^{\circ}\text{C}} & \alpha\text{-方石英} & \xrightarrow{1\ 713^{\circ}\text{C}} \text{熔}\text{SiO}_2 \\
\updownarrow 573^{\circ}\text{C} & & \updownarrow 163^{\circ}\text{C} & & \updownarrow 180 \sim 270^{\circ}\text{C} & \updownarrow \text{急冷} \\
\beta\text{-石英} & & \beta\text{-磷石英} & & \beta\text{-方石英} & \text{石英玻璃} \\
& & \updownarrow 117^{\circ}\text{C} & & & \\
& & \gamma\text{-磷石英} & & &
\end{array}
\]

石英的同质多晶转变也是形核和长大的过程,通过硅氧四面体的重新排列和组合形成各种晶体结构,从而改变石英的性质。

\subsection{同分异构}
化学成分相同而分子结构中原子排列不同的现象称为同分异构。同分异构在有机物质中经常出现。在有机低分子物质中,丙醇和异丙醇、甲醚和乙醇就是同分异构体,它们的化学成分相同,但分子结构不同。它们的分子结构如下:

\begin{figure}[H]
    \centering
    % Placeholder for chemical structures
\end{figure}

同分异构在高聚物中也是普遍存在的。许多共聚物,其单体相同,而单体在大分子链中排列方式不同。如A、B两种单体共聚,可产生下列共聚物:

\begin{figure}[H]
	\centering
	% Placeholder for chemical structures
\end{figure}

同分异构对高分子材料的性能影响很大,例如聚丙烯($[\text{CH}_2-\text{CH}(\text{CH}_3)]_n$)的甲基($-\text{CH}_3$)在主链中排列规整性高,结晶度亦高,强度也高。而无规聚丙烯(甲基任意取向)强度低,基本无实用价值。所以,同分异构对高分子材料的开发和应用具有重要意义。

\section{铁碳合金相图}

\subsection{二元合金相图}
合金状态图是用图解的方法表示合金系中合金状态与温度和成分之间的关系。还可以说,用图解的方法表示不同成分、温度下合金中相的平衡关系,因而也称为相图。由于状态图是在极其缓慢的冷却条件下测定的,故又称为平衡图。
现以$\text{Cu-Ni}$合金为例,说明用热分析法建立二元合金相图的步骤。

1) 配制不同成分的$\text{Cu-Ni}$合金,例如:

\textcircled{1}合金I——纯$\text{Cu}$。

\textcircled{2}合金II——$80\%\text{Cu} + 20\%\text{Ni}$。

\textcircled{3}合金III——$60\%\text{Cu} + 40\%\text{Ni}$。

\textcircled{4}合金IV——$40\%\text{Cu} + 60\%\text{Ni}$。

\textcircled{5}合金V——$20\%\text{Cu} + 80\%\text{Ni}$。

\textcircled{6}合金VI——纯$\text{Ni}$。

配制的合金越多,则相图越精确。

2) 用热分析法测定合金的冷却曲线,并找出各冷却曲线上的临界点(相变点,即转折点和平台)的温度值(图2-16a)。

3) 画出温度-成分坐标,在相应成分垂线上标出临界点温度。

4) 将物理意义相同的临界点连接成曲线,即得 Cu-Ni 合金相图(图 2-16b)。

常见的二元合金相图有匀晶相图、共晶相图、共析相图和形成稳定化合物相图。

两组元在液态无限互溶,在固态也无限互溶的合金系称为匀晶系,它们的相图为匀晶相图。Cu-Ni 合金相图(图 2-16)即典型的匀晶相图。

两组元在液态下无限互溶,而在固态下仅有限溶解并发生共晶反应的合金系形成共晶相图,如图 2-17 所示的 Pb-Sn 共晶相图。共晶反应是指从某种成分固定的合金溶液中、在恒温(共晶温度)下同时结晶出两种成分和结构皆不相同的固相的反应,其反应式为

\[ \mathrm{L}_{E} \xrightarrow{\text { 共晶温度 }}\left(\alpha_{C}+\beta_{D}\right) \]

\begin{figure}[H]
\centering
\includegraphics[page=3, viewport={326.5 237.8 535.5 405.8}, clip]{12-07572-001FigRevised.pdf}
\caption{二元合金 Cu-Ni 匀晶相图}
\end{figure}

\begin{figure}[H]
\centering
\includegraphics[page=3, viewport={23.5 60.8 192.0 195.8}, clip]{12-07572-001FigRevised.pdf}
\caption{二元合金 Pb-Sn 共晶相图}
\end{figure}

\begin{figure}[H]
	\centering
	\includegraphics[page=3, viewport={201.0 62.8 373.5 202.3}, clip]{12-07572-001FigRevised.pdf}
	\caption{具有共析反应的二无合金相图}
\end{figure}

图2-18为具有共析反应的二元合金相图,共析反应与共晶反应非常类似,是由一种固相在恒温(共析温度)下同时转变成两种新的固相,其反应式为

\[ \gamma_{E} \xrightarrow{\text { 共析温度 }}\left(\alpha_{C}+\beta_{D}\right) \]

\begin{figure}[H]
\centering
\includegraphics[page=3, viewport={386.5 63.8 524.5 167.3}, clip]{12-07572-001FigRevised.pdf}
\caption{具有稳定化合物的相同}
\end{figure}

图2-19为具有稳定化合物的相图。稳定化合物是具有一定熔点,而且在熔点以下温度不发生分解的金属化合物,因此可以将稳定化合物作为一独立组元。在某些合金系中,常形成一种或几种稳定化合物。在图 2-19 中,化合物 $\mathrm{A}_{m} \mathrm{~B}_{n}$ 将整个相图分为两个相对独立的相图。

\subsection{Fe-Fe$_3$C 相图}

在二元合金中,铁碳合金是现代工业使用最广泛的合金。铁碳合金是由铁和碳两个基本组元组成的,根据碳的质量分数不同,可以分为碳钢和铸铁两大类,碳钢是指碳的质量分数为 $0.02 \% \sim 2.11 \%$ 的铁碳合金,铸铁是指碳的质量分数大于 $2.11 \%$ 的铁碳合金。

铁碳合金相图是研究平衡状态下铁碳合金成分、组织和性能之间的关系及其变化规律的重要工具。掌握铁碳合金相图对热加工工艺的制订及工艺废品原因的分析都有重要的指导意义。铁碳合金相图是用试验方法做出的温度-成分坐标图。当铁碳合金的碳的质量分数超过 $6.69 \%$ 时,合金太脆无法应用,所以人们研究铁碳合金相图时,主要研究简化后的 $\mathrm{Fe}-\mathrm{Fe}_{3} \mathrm{C}$ 相图(图 2-20)。


\begin{figure}[H]
\centering
\includegraphics[page=4, viewport={56.5 457.8 477.0 752.3}, clip]{12-07572-001FigRevised.pdf}
\caption{简化后的 $\mathrm{Fe}-\mathrm{Fe}_{3} \mathrm{C}$ 相图}
\end{figure}

1. 铁碳合金的基本组织

(1) 铁素体

铁素体是碳在 $\alpha-\mathrm{Fe}$ 中形成的间隙固溶体,用 $\mathrm{F}$ 或 $\alpha$ 表示。由于 $\alpha-\mathrm{Fe}$ 为体心立方结构,溶碳能力较差(在 $727{ }^{\circ} \mathrm{C}$ 时碳的质量分数最大为 $0.0218 \%$,随着温度的下降碳的质量分数逐渐减小,在室温时碳的质量分数为 $0.0008 \%$ ),因此铁素体的强度、硬度不高,但具有良好的塑性和韧性。其性能指标( $R_{\mathrm{m}}=250 \mathrm{MPa}, R_{\mathrm{e}}=120 \mathrm{MPa}$,硬度值为 $80 \mathrm{HBW}, A=50 \%, Z=85 \%, a_{\mathrm{K}}=3000 \mathrm{~kJ} / \mathrm{m}^{2}$ )几乎和纯铁相同。

(2) 奥氏体

奥氏体是碳在 $\gamma-\mathrm{Fe}$ 中形成的间隙固溶体,用 $\mathrm{A}$ 或 $\gamma$ 表示。由于 $\gamma-\mathrm{Fe}$ 为面心立方结构,溶碳能力较大(在 $1148{ }^{\circ} \mathrm{C}$ 时溶碳能力最大,碳的质量分数可达 $2.11 \%$;随着温度的下降溶碳能力逐渐减小,在 $727{ }^{\circ} \mathrm{C}$ 时碳的质量分数为 $0.77 \%$ ),因此奥氏体是高温组织,其力学性能与碳的质量分数及晶粒大小有关。一般来说,奥氏体的硬度较低( $170 \sim 220 \mathrm{HBW}$ )而塑性较高( $A=40 \% \sim 50 \%$ ),易于塑性成形。

(3) 渗碳体

渗碳体是铁和碳的金属化合物( $\mathrm{Fe}_{3} \mathrm{C}$ ),其碳的质量分数为 $6.69 \%$ 。渗碳体的熔点约为 $1227{ }^{\circ} \mathrm{C}$,其硬度很高( $>800 \mathrm{HV}$ ),而脆性极大,塑性和冲击韧度几乎等于零。渗碳体是碳钢中主要的强化相,其形状、数量与分布等对钢的性能有很大的影响。渗碳体又是一种亚稳定相,在一定的条件下会发生分解,形成石墨状的自由碳,即
\[ \mathrm{Fe}_{3} \mathrm{C} \rightleftharpoons 3 \mathrm{Fe}+\mathrm{C}(\text { 石墨 }) \]

(4) 珠光体

珠光体是铁素体与渗碳体的机械混合物,用 $\mathrm{P}$ 表示。珠光体是共析反应的产物,其碳的质量分数为 $0.77 \%$ 。由于渗碳体在混合物中起强化作用,因此珠光体具有良好的力学性能。

(5) 莱氏体

莱氏体是奥氏体和渗碳体的机械混合物,用 $\mathrm{L}_{\mathrm{d}}$ 表示。莱氏体是共晶反应的产物,其碳的质量分数为 $4.3 \%$ 。高温莱氏体冷却到 $727{ }^{\circ} \mathrm{C}$ 以下时,将转变为珠光体和渗碳体的机械混合物,称为低温莱氏体,用 $\mathrm{L}_{\mathrm{d}}^{\prime}$ 表示。由于莱氏体中含渗碳体较多,其力学性能与渗碳体相近,属脆性组织。

2. Fe-Fe$_3$C 相图的分析

相图中用字母标出的点都具有一定的特性(成分和温度),所以称为特性点,各主要特性点的含义如表 2-1 所示。相图中的各条线表示合金内部组织发生转变的界线,所以称为组织转变线,一些主要组织转变线的含义如表 2-2 所示。

\begin{table}[H]
\centering
\caption{Fe-Fe$_3$C 相图中的特性点及其含义}
\label{tab:2-1}
\begin{tabular}{|c|c|c|c|}
\hline
特性点 & 温度/$^{\circ}$C & 碳的质量分数/\% & 说明 \\
\hline
A & 1538 & 0 & 纯铁的熔点 \\
\hline
C & 1148 & 4.30 & 共晶点 \\
\hline
D & 1227 & 6.69 & 渗碳体的熔点 \\
\hline
E & 1148 & 2.11 & 碳在 $\gamma$-Fe 中的最大固溶度 \\
\hline
F & 1148 & 6.69 & 渗碳体的成分 \\
\hline
G & 912 & 0 & $\alpha$-Fe $\rightleftharpoons \gamma$-Fe 同素异构转变点 \\
\hline
K & 727 & 6.69 & 渗碳体的成分 \\
\hline
P & 727 & 0.0218 & 碳在 $\alpha$-Fe 中最大的固溶度 \\
\hline
S & 727 & 0.77 & 共析点 \\
\hline
Q & 室温 & 0.0008 & 碳在 $\alpha$-Fe 中的固溶度 \\
\hline
\end{tabular}
\end{table}

\begin{table}[H]
\centering
\caption{Fe-Fe$_3$C 相图中的组织转变线及其含义}
\label{tab:2-2}
\begin{tabular}{|c|p{370pt}|}
\hline
组织转变线 & 含义 \\
\hline
AC & 铁碳合金的液相线,液态合金开始结晶出奥氏体 \\
\hline
CD & 铁碳合金的液相线,液态合金开始结晶出渗碳体 \\
\hline
AE & 铁碳合金的固相线,即奥氏体的结晶终了线 \\
\hline
ECF & 铁碳合金的固相线,也是 $\mathrm{L}_{C} \longrightarrow \mathrm{A}_{E}+\mathrm{Fe}_{3} \mathrm{C}$ 共晶转变线 \\
\hline
GS & 奥氏体转变为铁素体的开始线(常称 $A_{3}$ 线) \\
\hline
GP & 奥氏体转变为铁素体的终了线 \\
\hline
ES & 碳在奥氏体中固溶线(常称 $A_{\mathrm{cm}}$ 线),开始析出二次渗碳体( $\mathrm{Fe}_{3} \mathrm{C}_{\text {II }}$ ) \\
\hline
PQ & 碳在铁素体中固溶线,开始析出三次渗碳体( $\mathrm{Fe}_{3} \mathrm{C}_{\text {III }}$ ) \\
\hline
PSK & $\mathrm{A}_{S} \longrightarrow \mathrm{F}_{P}+\mathrm{Fe}_{3} \mathrm{C}$ 共析转变线(常称 $A_{1}$ 线) \\
\hline
\end{tabular}
\end{table}
注:表格中各组织转变线的含义均指合金在缓慢冷却过程中的相变线,如果是加热过程,则相反。

3. 典型合金的结晶过程分析

(1) 铁碳合金的分类

铁碳合金相图上的各种合金,按其碳的质量分数及组织的不同,常分为以下 3 类。

1) 工业纯铁(碳的质量分数 $<0.0218 \%$ ),其显微组织为铁素体。

2) 碳钢(碳的质量分数为 $0.0218 \% \sim 2.11 \%$ ),其特点是高温固态组织为具有良好塑性的奥氏体,宜于锻造。根据室温组织的不同,其又可分为以下 3 种:

\textcircled{1}亚共析钢(碳的质量分数 $<0.77 \%$ ),组织是铁素体和珠光体;

\textcircled{2}共析钢(碳的质量分数为 $0.77 \%$ ),组织是珠光体;

\textcircled{3}过共析钢(碳的质量分数 $>0.77 \%$ ),组织是珠光体和二次渗碳体。

3) 白口铸铁(碳的质量分数为 $2.11 \% \sim 6.69 \%$ ),其特点是液态结晶时都有共晶转变,因而有较好的铸造性能。根据室温组织的不同,其又可分为以下 3 种:

\textcircled{1}亚共晶铸铁(碳的质量分数 $<4.3 \%$ ),组织是珠光体、二次渗碳体和莱氏体;

\textcircled{2}共晶铸铁(碳的质量分数为 $4.3 \%$ ),组织是莱氏体;

\textcircled{3}过共晶铸铁(碳的质量分数 $>4.3 \%$ ),组织是莱氏体和一次渗碳体( $\mathrm{Fe}_{3} \mathrm{C}_{\text {I }}$ )。

(2) 典型合金的结晶

现以几种典型合金(图 2-21)为例,分析其结晶过程和在室温下的显微组织。

1) 共析钢(图 2-21 中合金 I )

合金在温度 1$\sim$2 点之间按匀晶转变从钢液中结晶出奥氏体。奥氏体冷却至 $727{ }^{\circ} \mathrm{C}$ (3 点)时,将发生共析转变,即 $\mathrm{A}_{S} \longrightarrow \mathrm{P}\left(\mathrm{F}_{P}+\mathrm{Fe}_{3} \mathrm{C}\right)$,形成珠光体。珠光体中的渗碳体称为共析渗碳体。当温度由 $727{ }^{\circ} \mathrm{C}$ 继续下降时,铁素体沿固溶线 $P Q$ 改变成分,析出三次渗碳体。三次渗碳体常与共析渗碳体连在一起,不易分辨,且数量极少,可忽略不计。图 2-22 所示为共析钢的显微组织(珠光体)。

2) 亚共析钢(图 2-21 中合金 II )

合金在温度 1$\sim$2 点之间按匀晶转变形成奥氏体。温度由 2 点降至 3 点的过程是奥氏体的单相冷却过程。继续冷却至 3$\sim$4 点之间时,奥氏体结晶出铁素体,在此过程中,奥氏体成分沿 $G S$ 线变化,铁素体成分沿 $G P$ 线变化。当温度降到 $727{ }^{\circ} \mathrm{C}$,奥氏体的成分达到 $S$ 点成分(碳的质量分数为 $0.77 \%$ )时,则发生共析转变,即 $\mathrm{A}_{S} \longrightarrow \mathrm{P}\left(\mathrm{F}_{P}+\mathrm{Fe}_{3} \mathrm{C}\right)$,形成珠光体。此时,原先析出的铁素体保持不变。当继续冷却时,铁素体的碳的质量分数沿 $P Q$ 线下降,同时析出三次渗碳体,因析出量极少,一般可忽略不计。亚共析钢的室温组织由铁素体和珠光体组成,碳的质量分数越高,则珠光体越多,铁素体越少。亚共析钢的显微组织如图 2-23 所示。

3) 过共析钢(图 2-21 中合金 III )

合金在 1$\sim$2 点之间按匀晶过程转变为奥氏体。在 2$\sim$3 点之间为单相奥氏体的冷却过程。自 3 点开始,奥氏体的溶碳能力降低,从其中析出二次渗碳体,并沿其晶界呈网状分布。在 3$\sim$4 点之间,随着温度的降低,析出的二次渗碳体量不断增多。与此同时,奥氏体的碳的质量分数也逐渐沿 $E S$ 线降低。当冷却到 $727{ }^{\circ} \mathrm{C}$ (4 点)时,奥氏体的成分达到 $S$ 点成分,于是发生共析转变,即 $\mathrm{A}_{S} \longrightarrow \mathrm{P}\left(\mathrm{F}_{P}+\mathrm{Fe}_{3} \mathrm{C}\right)$,形成珠光体。4 点以下直至室温,合金组织变化不大。过共析钢的室温组织由珠光体和网状二次渗碳体组成,其显微组织如图 2-24 所示。

\begin{figure}[H]
\centering
\includegraphics[page=4, viewport={59.0 264.3 288.5 428.8}, clip]{12-07572-001FigRevised.pdf}
\caption{典型铁碳合金在 $\mathrm{Fe}-\mathrm{Fe}_{3} \mathrm{C}$ 相图中的位置}
\end{figure}

\begin{figure}[H]
\centering
\includegraphics[page=4, viewport={311.5 265.3 481.5 395.8}, clip]{12-07572-001FigRevised.pdf}
\caption{共析钢的显微组织图(400$\times$)}
\end{figure}

\begin{figure}[H]
\centering
\includegraphics[page=4, viewport={87.5 93.8 258.0 225.3}, clip]{12-07572-001FigRevised.pdf}
\caption{亚共析钢的显微组织图(400$\times$)}
\end{figure}

\begin{figure}[H]
\centering
\includegraphics[page=4, viewport={288.0 96.8 455.0 225.8}, clip]{12-07572-001FigRevised.pdf}
\caption{过共析钢的显微组织图(400$\times$)}
\end{figure}

可用同样的方法分析共晶铸铁、亚共晶铸铁及过共晶铸铁的结晶过程。它们的室温组织分别为莱氏体(图 2-25),珠光体、二次渗碳体和莱氏体(图 2-26),一次渗碳体和莱氏体(图 2-27)。

\begin{figure}[H]
\centering
\includegraphics[page=5, viewport={89.0 613.8 257.5 743.8}, clip]{12-07572-001FigRevised.pdf}
\caption{共晶铸铁显微组织图(400$\times$)}
\end{figure}

\begin{figure}[H]
\centering
\includegraphics[page=5, viewport={286.0 612.3 456.0 742.3}, clip]{12-07572-001FigRevised.pdf}
\caption{亚共晶铸铁显微组织图(400$\times$)}
\end{figure}

\begin{figure}[H]
\centering
\includegraphics[page=5, viewport={66.5 232.3 256.5 360.3}, clip]{12-07572-001FigRevised.pdf}
\caption{过共晶铸铁显微组织图(400$\times$)}
\end{figure}


4. 铁碳合金的组织和性能的变化规律

铁碳合金室温的成分、组织和性能的对应关系如图 2-28 所示。

(1) 相的变化规律

铁碳合金的室温组织均由铁素体与渗碳体两相组成(碳的质量分数低于 $0.0218 \%$ 的合金除外)。随着碳的质量分数的增加,铁素体的质量分数呈直线关系减少,当碳的质量分数达到 $6.69 \%$ 时降为零。与此同时,渗碳体的质量分数则由零直线增加至 $100 \%$,如图 2-28b 所示。

碳的质量分数的变化不仅引起铁素体和渗碳体相对量的变化,而且由于引起不同性质的结晶过程,使其出现不同的组织形态,发生不同的相互结合,因此造成不同的组织变化。

(2) 组织的变化

1) 组织组成物的变化

随着碳的质量分数的增加,组织变化顺序依次为
\[ \mathrm{F} \rightarrow \mathrm{F}+\mathrm{P} \rightarrow \mathrm{P} \rightarrow \mathrm{P}+\mathrm{Fe}_{3} \mathrm{C}_{\text {II }} \rightarrow \mathrm{P}+\mathrm{Fe}_{3} \mathrm{C}_{\text {II }}+\mathrm{L}_{\mathrm{d}}^{\prime} \rightarrow \mathrm{L}_{\mathrm{d}}^{\prime} \rightarrow \mathrm{L}_{\mathrm{d}}^{\prime}+\mathrm{Fe}_{3} \mathrm{C}_{\text {I }} \]
组织中各种组织组成物的相对数量由图 2-28c 中相应竖直高度来表示。

2) 组织形态的变化

同一种组织组成物或组成相,由于生成条件的不同,虽然本质相同,但形态差别却很大,对性能的影响也大不一样。

\textcircled{1}铁素体: 固溶体转变生成的单相铁素体为块状(等轴晶粒状);共析体中的铁素体则由于同渗碳体相互制约,主要呈交替片状。

\textcircled{2}渗碳体: 它的形态最复杂,钢铁组织的复杂化主要是其造成的。

一次渗碳体是从液体中直接析出的,呈长条状;二次渗碳体是从奥氏体中析出的,沿晶界呈网状;三次渗碳体是从铁素体中析出的,沿晶界呈小片或粒状;共晶渗碳体是同奥氏体相关形成的,在莱氏体中为连续的基体;共析渗碳体是同铁素体交互形成的,呈交替片状。由此可见,铁碳合金中这些组织的不同形态,决定了其性能变化的复杂性。

(3) 性能的变化

1) 力学性能

\textcircled{1}硬度:是一个与碳的质量分数有关的性能指标,对组织组成物或组成相的形态不十分敏感,其大小主要取决于组成相的硬度和相对数量。随着碳的质量分数的增加,硬度高的渗碳体增多,硬度低的铁素体减少,因此合金的硬度呈直线增高(图 2-28d),由完全为铁素体组织的 $80 \mathrm{HBW}$ 增大到完全为渗碳体的约 $68 \mathrm{HRC}$ 。

\textcircled{2}强度:是一种对组织组成物的形态很敏感的性能。

在工业纯铁中,随着碳的质量分数的增加,固溶强化或微量三次渗碳体的强化作用使强度提高。

在亚共析钢中,其组织为铁素体和珠光体的混合物。铁素体的强度低,珠光体的强度较高,随着珠光体的增加,强度提高(图 2-28d)。强度与组织的细密度有关,组织越细密,则强度越高,所以在亚共析钢中,随着碳的质量分数的增加,强度提高。

在过共析钢中,铁素体消失,而硬脆的二次渗碳体出现,合金强度增加变缓。在碳的质量分数约为 $0.90 \%$ 时,由于沿晶界形成的二次渗碳体网趋于完整,强度开始迅速下降,当碳的质量分数增至 $2.11 \%$ 时,组织中出现莱氏体,强度值降到很低。如果继续增加碳的质量分数,则由于基体变为连片的渗碳体,强度变化不大,但值很低,接近渗碳体的强度( $20 \sim 30 \mathrm{MPa}$ )。

\textcircled{3}塑性:铁碳合金中的渗碳体极脆,没有塑性,合金的塑性完全由铁素体来提供。所以,当碳的质量分数增加,铁素体减少时,合金的塑性不断降低。当基体变为渗碳体后,塑性便降低到接近于零值(图 2-28d)。

\textcircled{4}韧性 铁碳合金的冲击韧度对组织及其形态最为敏感。当碳的质量分数增加时,脆性的渗碳体增多,不利的形态更加严重,韧性下降很快,下降的趋势比塑性更急剧。

2) 工艺性能

\textcircled{1}铸造性能:良好的铸造性能要求铁碳合金熔点低、流动性好、收缩和偏析倾向小。共晶成分的铁碳合金熔点最低,结晶范围最小,具有良好的铸造性能。

\textcircled{2}可锻性:良好的可锻性要求铁碳合金塑性好、变形抗力小。在 $\mathrm{Fe}-\mathrm{Fe}_{3} \mathrm{C}$ 相图中,单相奥氏体区的铁碳合金最合适,其次为奥氏体和铁素体两相区的铁碳合金,而有渗碳体存在的两相区铁碳合金的塑性和韧性较差。

\begin{figure}[H]
\centering
\includegraphics[page=5, viewport={283.5 227.8 480.5 578.8}, clip]{12-07572-001FigRevised.pdf}
\caption{铁碳合金的成分、组织和性能的对应关系}
\end{figure}

\section*{复习思考题}
 
2-1 简述金属三种典型晶体结构的特点。

2-2 金属的实际晶体中存在哪些晶体缺陷? 它们对性能各有什么影响?

2-3 合金元素在金属中的存在形式有哪几种? 它们各自具有什么特性?

2-4 什么是固溶强化? 造成固溶强化的原因是什么?

2-5 简述高聚物大分子链的结构与形态,及其它们对高聚物性能的影响。

2-6 陶瓷的典型组织由哪几种相组成? 它们各自具有什么作用?

2-7 从结构入手比较金属、高聚物、陶瓷三种材料的优缺点。

2-8 简述纯金属的结晶过程。

2-9 金属结晶的基本规律是什么?

2-10 如果其他条件相同,试比较在下列铸造条件下,铸件晶粒的大小:
(1) 金属型浇注与砂型浇注;
(2) 铸成薄件与铸成厚件;
(3) 浇注时采用振动与不采用振动。

2-11 金属或合金的晶粒大小对其力学性能有何影响? 如何实现晶粒细化?

2-12 固溶体与化合物各有什么特点?

2-13 简述用热分析法建立二元合金相图的步骤。

2-14 何谓共晶反应、匀晶反应和共析反应? 试比较三种反应的异同点。

2-15 不同碳质量分数的奥氏体,为什么当其温度降至共析温度时都会发生共析反应?

2-16 已知 A(熔点为 $600{ }^{\circ} \mathrm{C}$ )与 B(熔点为 $500{ }^{\circ} \mathrm{C}$ )在液态无限互溶,在 $300{ }^{\circ} \mathrm{C}$ 固态时 A 溶于 B 的最大固溶度为 $30 \%$,室温时为 $10 \%$,但 B 不溶于 A;在 $300{ }^{\circ} \mathrm{C}$ 时含 $40 \% \mathrm{~B}$ 的液态合金发生共晶反应,现要求:
(1) 作出 A-B 合金相图。
(2) 分析 A 的质量分数分别为 $20 \% 、 45 \% 、 80 \%$ 合金的结晶过程。

2-17 奥氏体的形成过程分为哪几个阶段? 简述其形成原理。

2-18 什么是铁素体,它有何特点?

2-19 为什么铸造合金常选用接近共晶成分的合金? 为什么要进行压力加工的合金常选用单相固溶体成分的合金?

2-20 何谓 $\mathrm{F} 、 \mathrm{~A} 、 \mathrm{Fe}_{3} \mathrm{C} 、 \mathrm{P} 、 \mathrm{~L}_{\mathrm{d}}\left(\mathrm{L}_{\mathrm{d}}^{\prime}\right)$ ? 它们的结构、组织形态、力学性能有何特点?

2-21 碳素钢、铸铁两者的成分、组织和性能有何差别? 并说明原因。

2-22 分析碳的质量分数分别为 $0.20 \% 、 0.60 \% 、 0.80 \% 、 1.0 \%$ 的铁碳合金从液态缓冷至室温时的结晶过程和室温组织,指出这四种成分组织与性能的区别。

2-23 渗碳体有哪 5 种基本形态? 它们的来源和形态有何区别?

2-24 根据 $\text{Fe-Fe}_3\text{C}$ 相图，说明产生下列现象的原因。
(1) 碳的质量分数为 $1.0\%$ 的钢比碳的质量分数为 $0.5\%$ 的钢硬度高；
(2) 低温莱氏体的塑性比珠光体的塑性差；
(3) 捆扎物体一般用铁丝（镀锌的碳钢丝），而起重机起吊重物却用钢丝绳（用 60、65、70、75 等钢制成）；
(4) 一般要把钢材加热到高温（$1~000 \sim 1~250~^\circ\text{C}$）下进行热轧或锻造；
(5) 钢适宜于通过压力加工成形，而铸铁适宜于通过铸造成形。

2-25 为什么说 $\text{Fe-C}$ 合金相图是选材和制订热加工工艺的理论依据？


\chapter{金属热处理及表面改性}
金属热处理就是通过加热、保温和冷却改变金属内部组织或表面的组织，从而获得所需性能的工艺方法。常用的钢的热处理有退火、正火、淬火、回火及表面热处理等。在机械制造业中，热处理占有非常重要的地位，能使材料充分发挥其潜能，延长零件的使用寿命。例如，在各种机床设备中 $60\% \sim 70\%$ 的零部件需要热处理，重要机械中需要热处理的零件所占比例更高。

腐蚀和磨损是金属零件破坏的主要形式，为此发展出表面淬火、渗碳等表面热处理技术，以及电镀、热浸镀、气相沉积等涂层技术。以激光束、电子束和等离子束为代表的高能束表面改性技术，已发展成为表面工程技术，以改变金属材料表面层组织和结构，达到抗磨和耐蚀的目的或满足其他特殊性能要求。

\section{钢的热处理原理}

钢的热处理工艺虽然种类繁多，但大多采用一定的速度将钢加热到一定温度，保温一段时间，再以适当的速度冷却到室温。将上述过程建立如图 3-1 所示的温度-时间坐标图，便得到了热处理工艺曲线。通过改变工艺规范，可以获得不同的组织，以满足零件性能要求。

在 $\text{Fe-Fe}_3\text{C}$ 相图中，$A_1$、$A_3$ 和 $A_{\text{cm}}$ 线表示钢在缓慢加热和冷却时的相变温度，即平衡临界温度。在热处理中，加热和冷却速度不可能很慢，实际的临界温度和相图的平衡临界温度相比存在一定的滞后，即过热和过冷现象。通常，把实际加热时的临界温度用 $Ac_1$、$Ac_3$ 和 $Ac_{\text{cm}}$ 表示，把冷却时的临界温度用 $Ar_1$、$Ar_3$ 和 $Ar_{\text{cm}}$ 表示，如图 3-2 所示。

\subsection{钢在加热时的组织转变}
大多数热处理工艺是把钢加热到临界温度 $Ac_1$ 线以上，使其内部组织转变成均匀的奥氏体，这种加热转变过程称为奥氏体化。加热后得到的奥氏体组织状态（均匀性和晶粒大小）对热处理的最终质量影响很大。

当共析钢加热到 $Ac_1$ 线时，便发生珠光体向奥氏体的转变。如图 3-3 所示，奥氏体的形成过程可分为四个阶段：奥氏体晶核在铁素体和渗碳体界面上的形成（图 3-3a）；奥氏体晶核的长大（图 3-3b）；剩余渗碳体的溶解（图 3-3c）；奥氏体的均匀化（图 3-3d）。对于亚共析钢或过共析钢，当加热到 $Ac_1$ 线时，钢中的珠光体转变成奥氏体，其余的铁素体或二次渗碳体基本不变。随着温度升高，铁素体或二次渗碳体变成或溶入奥氏体，直至加热温度超过 $Ac_3$ 或 $Ac_{\text{cm}}$ 线，钢中的各种组织才会全部转变为奥氏体。

\begin{figure}[H]
\centering
\includegraphics[page=5, viewport={175.0 80.8 363.5 193.3}, clip]{12-07572-001FigRevised.pdf}
\caption{热处理温度-时间坐标图}
\end{figure}

\begin{figure}[H]
\centering
\includegraphics[page=6, viewport={150.0 551.8 347.5 743.8}, clip]{12-07572-001FigRevised.pdf}
\caption{钢加热（冷却）时$\text{Fe-Fe}_3\text{C}$ 相图上各临界点的实际位置}
\end{figure}
 
\begin{figure}[H]
\centering
\includegraphics[page=6, viewport={103.5 421.3 394.5 516.8}, clip]{12-07572-001FigRevised.pdf}
\caption{共析钢奥氏体形成过程示意图}
\end{figure}

奥氏体晶粒的大小对冷却钢的组织和性能有很大影响。金属零件加热获得细小、均匀的奥氏体，冷却后内部组织也较细小，力学性能高。奥氏体晶粒大小主要取决于加热温度和保温时间。适当的加热温度和保温时间可使奥氏体具有较高的形核率、较慢的晶核长大速度和均匀的成分，从而获得细小均匀的奥氏体晶粒；加热温度过高或高温下保温时间过长，均会产生粗大的奥氏体晶粒。因此，控制钢的加热温度和保温时间十分重要。

\subsection{钢在冷却时的组织转变}
冷却方式和冷却速度对奥氏体的组织转变及其钢的最终性能有决定性作用。钢的热处理工艺中有以下两种冷却方式。

1. 等温冷却

首先，将已奥氏体化的钢以较快的冷却速度冷却至 $A_1$ 线以下一定的温度，这时奥氏体尚未转变，成为过冷奥氏体；然后，保温使奥氏体在等温下进行组织转变；最后，奥氏体转变完成后冷却至室温。上述冷却过程即等温冷却，如热处理工艺中的等温淬火。

等温冷却方式对研究冷却过程中的组织转变较为方便。以共析钢为例，将一组共析钢薄片试样分别放入温度不同的恒温盐浴中，进行一系列不同过冷度的等温冷却试验，测出在不同等温温度下过冷奥氏体开始转变和转变终了的时间，先把相应的点画在温度-时间坐标图中，然后把表示开始转变时间的点和转变终了时间的点分别连接起来，即得到共析钢过冷奥氏体的等温转变曲线，如图 3-4 所示。由于其形状与字母“C”相似，故称为 C 曲线，又称为 TTT（tencperature-time transformation）曲线。它表示了一定成分的钢经奥氏体化后等温冷却转变的时间-温度-组织关系，是制订钢热处理工艺的重要依据。

\begin{figure}[H]
\centering
\includegraphics[page=6, viewport={99.5 105.3 447.0 390.3}, clip]{12-07572-001FigRevised.pdf}
\caption{共析钢过冷奥氏体等温转变曲线}
\end{figure}

在 C 曲线中，$A_1$ 线以上是奥氏体稳定存在区；在 $A_1$ 线以下，转变开始线 I 以左的区域是奥氏体不稳定存在区，称为过冷奥氏体区，此区中的过冷奥氏体要经一段孕育期才开始发生组织转变；在转变终了线 II 的右方是转变产物区；在两条曲线之间是转变过渡区，过冷奥氏体和转变产物同时存在；水平线 $M_{\text{s}}$ 为马氏体转变开始线，$M_{\text{f}}$ 线为马氏体转变终了温度线，在 $M_{\text{s}}$ 线与 $M_{\text{f}}$ 线之间为马氏体转变温度区。

孕育期的长短随过冷度而变化，孕育期越长，表明过冷奥氏体越稳定，反之，则越不稳定。对于共析钢，在 $550~^\circ\text{C}$ 左右等温时的孕育期最短，此处俗称鼻温。在 $A_1$ 线至鼻温之间，随着等温温度的下降，过冷度增大，孕育期逐渐变短，转变速度加快。在鼻温至 $M_{\text{s}}$ 线之间，随着等温温度的下降，原子扩散能力降低，孕育期逐渐变长，转变速度减慢。过冷奥氏体在鼻温处最不稳定。

共析钢过冷奥氏体等温转变产物可分为以下三种类型。

(1) 高温转变产物

在 $727 \sim 550~^\circ\text{C}$ 之间等温转变的产物为珠光体型组织，是由铁素体和渗碳体的层片组成的机械混合物。过冷度越大，层片越细小，钢的强度和硬度也越高。依照组织中层片的尺寸，又把在 $727 \sim 650~^\circ\text{C}$ 之间等温转变的组织称为粗片状珠光体；在 $650 \sim 600~^\circ\text{C}$ 之间等温转变的组织称为索氏体；在 $600 \sim 550~^\circ\text{C}$ 之间等温转变的组织称为托氏体。以上三种组织的显微组织如图 3-5 所示。

\begin{figure}[H]
\centering
\includegraphics[page=7, viewport={64.5 625.3 481.0 743.8}, clip]{12-07572-001FigRevised.pdf}
\caption{珠光体、索氏体、托氏体的显微组织图}
\end{figure}

(2) 中温转变产物

在 $550 \sim 230~^\circ\text{C}$ 之间等温转变的产物为贝氏体型组织，是由碳过饱和的铁素体和微小的渗碳体混合而成的。贝氏体一般具有比珠光体高的强度和硬度。根据转变温度和组织形态不同，贝氏体一般可分为上贝氏体和下贝氏体。上贝氏体为羽毛状（图 3-6a），由于韧性差，生产中很少采用。下贝氏体是在 $350 \sim 230~^\circ\text{C}$ 之间等温转变形成的，碳化物细小、分布均匀，在显微镜下呈黑色针状（图 3-6b），具有良好的韧性和塑性。

\begin{figure}[H]
\centering
\includegraphics[page=7, viewport={104.0 450.8 443.0 588.3}, clip]{12-07572-001FigRevised.pdf}
\caption{上贝氏体和下贝氏体显微组织图}
\end{figure}

(3) 低温转变产物

过冷奥氏体在 $230~^\circ\text{C}$ 开始转变成马氏体。随着等温转变温度的下降，马氏体逐渐增多。约在 $-50~^\circ\text{C}$ 过冷奥氏体停止向马氏体转变，未转变的奥氏体称为残余奥氏体，用 $A_{\text{r}}$ 表示。过冷奥氏体存在于 $M_{\text{s}}$ 与 $M_{\text{f}}$ 之间，转变后的组织为马氏体和残余奥氏体。

当奥氏体被迅速过冷到 $M_{\text{s}}$ 线以下时，$\gamma\text{-Fe}$ 会向 $\alpha\text{-Fe}$ 的晶格转变，由于碳原子无法扩散而留在 $\alpha\text{-Fe}$ 中，形成碳过饱和的 $\alpha\text{-Fe}$ 固溶体。过饱和的碳使 $\alpha\text{-Fe}$ 的晶格产生严重畸变，产生强烈的强化作用，因而马氏体具有很高的硬度（$62 \sim 65~\text{HRC}$）。马氏体转变能显著提高钢的强度和硬度。

马氏体按其碳的质量分数的不同，可分为高碳马氏体和低碳马氏体。碳的质量分数大于 $1\%$ 的马氏体是高碳马氏体（又称片状马氏体），呈凸透镜片状或针状，硬度很高，但塑性、韧性很差；碳的质量分数小于 $0.2\%$ 的马氏体是低碳马氏体（又称板条马氏体），呈板条状，强度、硬度高，塑性、韧性好；碳的质量分数为 $0.2\% \sim 1\%$ 的组织是两类马氏体的混合组织。马氏体的显微组织如图 3-7 所示。

\begin{figure}[H]
\centering
\includegraphics[page=7, viewport={102.0 272.8 443.0 411.8}, clip]{12-07572-001FigRevised.pdf}
\caption{马氏体显微组织图}
\end{figure}

2. 连续冷却

将已奥氏体化的钢在温度连续下降的过程中冷却，使其组织发生转变。这一冷却过程即连续冷却，如热处理生产中经常使用的工件在水中、油中和空气中的冷却。

同等温转变一样，奥氏体在连续冷却过程中也能发生珠光体转变、贝氏体转变和马氏体转变等，但是连续冷却过程要先后通过各个转变温度区，一个钢件可能先后发生几种转变，得到几种转变产物复合的不均匀组织。而且，冷却速度不同，可能发生的转变及其转变组织的相对量也不同，得到的组织及其性能也各异。因此，连续冷却转变比等温冷却转变更为复杂。图 3-8 为共析钢连续冷却转变曲线（实线）与等温冷却曲线（虚线）的对比。共析钢的连续冷却转变曲线 [ CCT（continuous-cooling transformation）曲线] 只有珠光体转变区和马氏体转变区，没有贝氏体转变。珠光体转变区由三条曲线构成：珠光体转变开始线（$P_{\text{s}}$ 线）、珠光体转变结束线（$P_{\text{z}}$ 线）和珠光体转变停止线（$K$ 线）。$v_{\text{K}}$ 是马氏体组织转变临界冷却速度。

由于等温冷却曲线比较容易测试，所以在热处理生产中常利用等温冷却曲线粗略估计连续冷却的转变产物与性能。

\begin{figure}[H]
\centering
\includegraphics[page=7, viewport={70.0 59.3 260.5 234.8}, clip]{12-07572-001FigRevised.pdf}
\caption{共析钢连续冷却转变曲线}
\end{figure}

\section{钢的普通热处理工艺}
\subsection{钢的退火与正火}
退火和正火是应用非常广泛的热处理工艺，主要用于铸件、锻件和焊件的预备热处理，目的在于消除冶金及热加工过程中产生的某些缺陷，改善组织及加工性能。对于某些性能要求不高的零件，亦可作为最终热处理。各种退火和正火的加热温度范围如图 3-9 所示。

钢退火和正火的目的是调整硬度以便进行切削加工（最佳切削硬度为 $170 \sim 250~\text{HBW}$）；消除残余内应力，以减少工件在淬火时产生变形或开裂的倾向；细化晶粒，改善组织，提高力学性能；为最终热处理（淬火/回火）作好组织准备。

1. 退火

将组织偏离平衡状态的金属或合金，先加热到适当温度，保持一定时间，然后缓慢冷却，以得到接近平衡状态组织的热处理工艺称为退火。工件常用的退火有下面几种。

(1) 完全退火

完全退火是将钢加热到 $Ac_3$ 线以上 $30 \sim 50~^\circ\text{C}$，保温一段时间，再缓慢冷却下来。在加热和保温过程中，室温下的珠光体和铁素体全部转变为奥氏体。

完全退火主要用于亚共析钢的锻件、轧件和铸件，它能细化晶粒、消除内应力、降低硬度和改善切削加工性。完全退火后的组织为珠光体加铁素体。

(2) 不完全退火

不完全退火是将钢加热到 $Ac_1$ 线以上 $20 \sim 30~^\circ\text{C}$，经适当保温后缓慢冷却。在加热和保温过程中，室温下的珠光体全部转变成奥氏体，而过剩相（渗碳体或铁素体）大部分保持不变。

不完全退火主要用于过共析钢。经不完全退火后，钢中的二次渗碳体和珠光体中的片状渗碳体变为球状渗碳体，故过共析钢的不完全退火又称为球化退火。

(3) 去应力退火

去应力退火是指将钢加热到 $Ac_1$ 线以下某一温度，经保温后缓慢冷却。去应力退火过程中不发生奥氏体相变，常用于消除铸件、锻件和焊件的内应力，以防止变形和开裂。因加热温度较低，故去应力退火又称为低温退火。

2. 正火

将钢加热到 $Ac_3$ 或 $Ac_{\text{cm}}$ 线以上 $30 \sim 50~^\circ\text{C}$，适当保温后在空气中冷却至室温，得到珠光体类型的组织，这种热处理工艺称为正火。与退火相比，正火处理的冷却速度较快，所得到的珠光体组织更细小，钢的强度和硬度均有所提高。

正火处理的冷却不占用设备，生产率高，所以低碳钢多用正火来代替完全退火。对于不太重要的零件，正火可以作为最终热处理。过共析钢经正火处理可减少二次渗碳体，消除网状渗碳体，便于进行球化退火。铸件经正火可改善性能，使粗大晶粒组织细化，且组织均匀。

\begin{figure}[H]
\centering
\includegraphics[page=7, viewport={278.0 58.8 473.5 235.8}, clip]{12-07572-001FigRevised.pdf}
\caption{各种退火和正火的加热温度范围}
\end{figure}

\subsection{钢的淬火与回火}
1. 淬火

将钢加热到 $Ac_3$ 或 $Ac_1$ 线以上 $30 \sim 50~^\circ\text{C}$，保温一段时间，再快速冷却，以获得马氏体、贝氏体等非平衡组织的热处理工艺称为淬火。

钢淬火的目的主要是提高钢的强度、硬度和耐磨性。钢淬火后再与适当的回火相配合，可获得所需要的力学性能。

淬火质量取决于淬火的加热温度和冷却方式。

(1) 淬火加热温度

确定钢的淬火加热温度时主要考虑的因素是钢的化学成分。

碳钢的淬火加热温度可利用 $\text{Fe-Fe}_3\text{C}$ 相图来选择。对于亚共析钢，适宜的淬火温度为 $Ac_3 + 30 \sim 50~^\circ\text{C}$，这样可获得均匀细小的马氏体组织。如果淬火加热温度过高，得到的粗大马氏体组织将使钢的塑性和韧性下降，淬火应力的增加还将引起工件较严重的变形甚至开裂；如果淬火加热温度过低，则淬火组织中出现铁素体，致使钢的强度不足，出现软点。对于过共析钢，适宜的淬火加热温度为 $Ac_1 + 30 \sim 50~^\circ\text{C}$，这样可获得均匀细小的马氏体和粒状渗碳体的混合组织。如果淬火温度过高，渗碳体溶解过多，得到粗片状马氏体，残余奥氏体增多，钢的硬度下降，变形、开裂倾向增大。

对于合金钢，其淬火的加热温度一般都比碳钢高，这是因为大多数合金元素会阻碍奥氏体的晶粒长大。提高淬火方式温度可使合金元素充分溶解和均匀化，取得较好的淬火效果。

(2) 淬火冷却

由共析钢过冷奥氏体等温转变曲线（图 3-8）可以看出，要得到马氏体，必须快速冷却以避免珠光体转变，但冷却速度过大会因内应力过大引起工件变形和开裂。淬火冷却方式要根据工件的化学成分及尺寸、形状和技术要求，选用合适的冷却介质。常用淬火冷却介质的冷却能力见表 3-1。

\begin{table}[H]
\centering
\caption{常用淬火冷却介质的冷却能力}
\begin{tabular}{|c|c|c|}
\hline
\multirow{2}{*}{淬火冷却介质} & \multicolumn{2}{c|}{冷却能力/($^\circ\text{C/s}$)} \\
\cline{2-3}
 & $650 \sim 550~^\circ\text{C}$ & $300 \sim 200~^\circ\text{C}$ \\
\hline
水($18~^\circ\text{C}$) & 600 & 270 \\
\hline
$10\% \text{NaCl}$ 水溶液($18~^\circ\text{C}$) & 1 100 & 300 \\
\hline
$10\% \text{NaOH}$ 水溶液($18~^\circ\text{C}$) & 1 200 & 300 \\
\hline
$10\% \text{Na}_2\text{CO}_3$ 水溶液($18~^\circ\text{C}$) & 800 & 270 \\
\hline
矿物机油 & 150 & 30 \\
\hline
菜籽油 & 200 & 35 \\
\hline
硝熔盐($200~^\circ\text{C}$) & 350 & 10 \\
\hline
\end{tabular}
\end{table}

碳钢的C曲线位置靠左,大截面的工件冷却较慢,应采用冷却能力较强的冷却介质,如水和盐水。合金钢的C曲线位置靠右,形状复杂的工件快速冷却易发生变形、开裂,应采用冷却能力较和缓的冷却介质,如矿物机油、水玻璃和聚乙烯醇合成的淬火剂等。

(3)钢的淬透性和淬硬性

钢的淬透性是指钢在淬火时获得马氏体的能力,在规定的淬火条件下该能力决定着钢的淬硬深度和硬度分布,通常用淬硬深度来表示。

对于截面尺寸较大的工件,加热后淬火时,工件表面冷却速度最大,而心部冷却速度最小。在工件截面上,凡是冷却速度大于临界冷却速度$v_K$的部分全部得到马氏体组织,而冷却速度低于$v_K$的部分将得到非马氏体组织,如图3-10所示。一般规定由工件表面到半马氏体区(即马氏体与珠光体型组织各占50\%的区域)的深度作为淬硬深度。半马氏体区的位置容易用金相显微镜观察和硬度仪测量硬度来确定。因此,淬透性也可以理解为钢在淬火后获得的淬硬深度大小,其实质是反映过冷奥氏体的稳定性。

\begin{figure}[H]
    \centering
    % Placeholder for Figure 3-10
    \includegraphics[page=8, viewport={52.5 486.8 299.5 678.3}, clip]{12-07572-001FigRevised.pdf}
    \caption{工件截面上不同过冷度及淬硬深度示意图}
\end{figure}

钢的淬透性与淬硬深度并不是同一概念,钢的淬透性是钢本身所固有的属性,它不受其他外部工艺条件的影响,而钢的淬硬深度虽然取决于钢的淬透性,但如果规定条件改变,淬透性强的钢材淬硬深度不一定大。例如,工件尺寸大、淬火冷却介质的冷却能力低,即便淬透性强的钢材也不容易获得较大的淬硬深度。

钢的淬硬性是指钢在淬火时所能获得一定硬度的能力是反映钢在淬火时的硬化能力。淬硬性强弱主要取决于马氏体的碳质量分数。马氏体的碳质量分数越高则淬硬性越强。在马氏体的碳质量分数相同的条件下,马氏体量越多,则淬硬性越强。

影响淬透性的决定因素是临界冷却速度$v_K$。临界冷却速度越小,钢的淬透性就越强。临界冷却速度是在C曲线上体现的,影响C曲线的因素也就是影响临界冷却速度的因素,所以化学成分与奥氏体化条件是影响淬透性的因素。化学成分中使C曲线右移的元素提升奥氏体的稳定性,使钢的临界冷却速度减小,其淬透性好;反之,淬透性差。因此,不同成分的钢有不同的淬透性。奥氏体化温度越高,保温时间越长,则晶粒越粗大,成分越均匀,因而过冷奥氏体越稳定。

C曲线越向右移,淬火临界冷却速度越小,钢的淬透性也越好。

2.回火

淬火后的钢先加热到$Ac_1$线以下某一温度,保温一定时间,然后冷却到室温的热处理工艺称为回火。

回火的目的在于降低淬火钢的脆性,减小或消除淬火后的内应力,防止工件的变形和开裂,稳定组织,调整力学性能,以满足使用要求。

根据加热温度的不同,回火可分为低温回火、中温回火和高温回火。

低温回火的加热温度为150$\sim$250$^\circ$C。低温回火后,钢的组织为马氏体,它是过饱和程度较小的$\alpha$固溶体。对于高碳钢,低温回火降低了其淬火应力和脆性,保持了高硬度和高耐磨性;对于低碳钢,低温回火降低使其获得了较高的塑性、韧性与较高强度的良好配合。低温回火广泛用于各类工具、模具、轴承、渗碳件、表面淬火件以及低碳合金结构钢工件。

中温回火的加热温度为250$\sim$500$^\circ$C。中温回火后,钢的组织为贝氏体,它是极细的球状渗碳体和铁素体的机械混合物。钢的内应力基本消除,弹性极限显著提高,但硬度有所降低。中温回火主要用于碳的质量分数为0.45\%$\sim$0.90\%的各类弹簧及某些要求高强度的结构钢工件。

高温回火的加热温度为500$\sim$650$^\circ$C。高温回火后,钢的组织为索氏体,它是较细的颗粒状渗碳体和铁素体的机械混合物。中碳钢经高温回火后内应力消除较完全,既有一定的强度和硬度,又有良好的塑性和韧性。生产中把淬火和高温回火相结合的热处理工艺称为调质。调质广泛应用于各种重要的机械零件,尤其是在交变载荷下工作的零件,如轴、齿轮、连杆和螺栓等。

\section{钢的表面热处理}

某些机器零件如凸轮、曲轴、齿轮等,要求表面具有高的硬度和耐磨性,而心部应具有足够的塑性和韧性。普通热处理难以实现这样的要求,因此采用表面热处理工艺。

常用的表面热处理工艺可分为两类:一类是只改变表面组织而不改变表面化学成分的表面淬火;另一类是同时改变表面化学成分和组织的表面化学热处理。

\subsection{表面淬火}

通过快速加热,零件表面很快达到淬火加热温度而内部还没有达到临界冷却温度时迅速冷却,使零件表面获得马氏体组织而心部仍保持塑性、韧性较好的原始组织的局部淬火方法称为表面淬火。

表面淬火的加热方式有火焰加热、感应加热、电接触加热、电解液加热等多种。应用较多的表面淬火有感应加热淬火和火焰加热淬火。

1.感应加热淬火

感应加热的基础是电磁感应、趋肤效应和热传导。感应加热淬火的工作原理是工件利用感应电流通过时产生的热效应使其表面迅速加热到$Ac_3$线以上80$\sim$150$^\circ$C,并快速喷水冷却,如图3-11所示。根据频率的不同,感应加热装置分为高频(60$\sim$300 kHz)、中频(1$\sim$10 kHz)及工频(50 Hz)三种。淬硬深度随频率降低而增加。常用高频感应加热淬火的淬硬深度为0.5$\sim$2.5 mm。工件淬火后获得细隐晶马氏体,表面硬度比一般淬火硬度高2$\sim$3 HRC,疲劳强度提高20\%$\sim$30\%。

感应加热淬火的加热速度快,生产率高,工件变形小,淬硬深度易于控制,淬火质量高,易于实现机械化、自动化,适宜于大批量生产。感应加热淬火一般用于中碳钢或中碳低合金钢,也可用于高碳工具钢或铸铁,但设备较复杂且昂贵,对批量小和不规则外形的零件不适宜。

\begin{figure}[H]
	\centering
	% Placeholder for Figure 3-11
	\includegraphics[page=8, viewport={322.0 485.3 489.0 746.3}, clip]{12-07572-001FigRevised.pdf}
	\caption{工件截面上不同过冷度及淬硬深度示意图}
\end{figure}

2.火焰加热淬火

火焰加热淬火利用氧-乙炔或其他可燃气体形成的高温火焰,喷射到工件表面上,使其迅速加热到淬火温度时立即喷水冷却,从而获得表面淬硬层的表面淬火方法。

火焰加热淬火的淬硬深度一般为1$\sim$6 mm,所需设备简单,适用于单件或小批生产的大型零件和需要局部表面淬火的零件。但其淬火质量不稳定,零件表面容易过热,生产效率低。

\subsection{化学热处理}

将工件放在一定的活性介质中加热保温,使介质中的活性原子渗入工件表层,从而改变表层的化学成分、组织和性能的工艺方法称为化学热处理。

化学热处理的目的是使工件心部具有足够的强度和韧性,而表面具有高的硬度和耐磨性;提高工件的疲劳强度和表面耐蚀性、耐热性等。根据渗入元素的不同,钢的化学热处理分为渗碳、渗氮、碳氮共渗、渗硫、渗硼、渗金属(铝、铬等)等,以渗碳、渗氮和碳氮共渗最为常用。

1.渗碳

根据渗碳介质的不同,渗碳方法可分为固体渗碳、液体渗碳和气体渗碳,以气体渗碳应用最广。工件在密封的加热炉中被加热至900$\sim$950$^\circ$C,向炉内通入渗碳气体(如煤气、天然气等)或滴入易高温裂解形成渗碳气氛的液体(如煤油、甲醇、丙酮等),以供给活性碳原子并在工件表面吸附、扩散,形成渗碳层。气体渗碳法的原理如图3-12所示。

\begin{figure}[H]
	\centering
	% Placeholder for Figure 3-12
	\includegraphics[page=8, viewport={16.0 305.3 150.5 453.3}, clip]{12-07572-001FigRevised.pdf}
	\caption{工件截面上不同过冷度及淬硬深度示意图}
\end{figure}

渗碳多用于低碳钢和低碳合金钢制成的齿轮、活塞销、轴类零件等重要零件。经过渗碳,工件的表层中碳的质量分数为0.85\%$\sim$1.05\%。再经过淬火、低温回火后,工件表层组织为细针状的高碳马氏体与渗碳体,其硬度、耐磨性及疲劳强度显著提高;内部组织为低碳马氏体或索氏体,具有一定的强度和良好的韧性。

2.渗氮

渗氮是向工件表面渗入氮原子,形成富氮的表面硬化层的化学热处理方法。渗氮的目的是提高工件表面硬度、耐磨性、耐蚀性和疲劳强度等性能。常用的渗氮方法有气体渗氮法和离子渗氮法。

气体渗氮法应用较广泛,它利用氨在500$\sim$600$^\circ$C加热分解出的活性氮原子渗入工件表面,通过扩散形成氮化层。适于渗氮的钢很多,如结构钢、工具钢、不锈钢等。当要求工件表面硬度高、抗磨损、抗疲劳、耐腐蚀,心部有良好的综合力学性能时,常选用含Cr、Mo、Al、Ti、V等元素的合金结构钢,最常用的钢种是38CrMoAl。渗氮前,一般需对工件进行调质或正火处理,其高温回火温度应比渗氮温度高。

与渗碳相比,渗氮后的表面硬化层薄($<$0.8 mm),但硬度高、耐磨性和耐腐蚀性好,且具有高的热硬性,因渗氮后无须再进行热处理,故工件变形小。渗氮常用于在交变载荷下工作的各种结构件,如各种高速传动精密齿轮、高精度机床主轴及阀门零件等。

3.碳氮共渗

碳氮共渗是将碳、氮原子同时渗入工件表面,形成碳氮共渗层的化学热处理方法。该方法兼有渗碳和渗氮两种方法的优点:工件表面硬度高,渗层较深,硬度变化平缓,具有良好的耐磨性和较小的表面脆性。工件碳氮共渗后,一般不需进行其他热处理和机械加工,可直接使用。碳氮共渗的工艺以中温气体碳氮共渗和低温气体碳氮共渗(又称气体氮碳共渗)应用最多。液体碳氮共渗有毒,污染环境,劳动条件差,已被限制使用。

中温气体碳氮共渗用煤油和氨气作为共渗剂,加热温度为820$\sim$860$^\circ$C,渗层厚度一般为0.5$\sim$0.8 mm,可直接进行淬火和低温回火。中温气体碳氮共渗常用于各种结构钢制造的机械零件,如齿轮、凸轮以及要求变形小的耐磨零件等。

低温气体碳氮共渗是在500$\sim$600$^\circ$C下进行的,共渗时长为1$\sim$6 h,氮和碳原子同时渗入工件表面,以渗氮为主,形成0.1$\sim$0.4 mm的碳氮共渗层,工件硬度为54$\sim$60 HRC。低温气体碳氮共渗普遍用于模具、量具以及要求耐磨的零件。

\section{金属材料的表面改性}

强化金属的方法很多,如前面介绍的固溶强化、弥散强化等。此外,通过表面改性也可强化金属材料。

磨损和腐蚀是机械零部件的两大主要失效形式,所造成的直接和间接经济损失十分惊人。表面改性技术赋予基体材料本身所不具备的高硬度、高耐磨性、耐蚀性和抗高温氧化性,目前已经发展成为表面工程技术。它主要包括薄膜技术、涂层技术、表面强化技术等,本节只对常用技术作简要介绍。

1.钢铁氧化、磷化与镀层钝化

生产中为了对工件表面进行有效的保护,使工件表面形成一层均匀致密的薄膜,以提高工件表面的防锈性能及减摩性能。

(1)氧化处理

氧化处理是把钢铁工件在沸腾的含浓碱和氧化剂(亚硝酸钠或硝酸钠)溶液中加热,使其表面形成一层厚0.5$\sim$1.5 $\upmu$m致密的氧化膜,既可以防止工件腐蚀和机械磨损,又可作为装饰性加工的工艺方法。由于氧化处理形成的氧化膜呈深黑蓝色,故又称发蓝处理或煮黑。氧化处理时对工件尺寸和表面粗糙度影响不大,并能在处理过程中消除内应力,因此在精密仪器、仪表,工具、模具等生产制造中得到了广泛应用。

(2)磷化处理

磷化处理是将钢铁工件浸入磷酸盐溶液中,在其表面形成一层5$\sim$15 $\upmu$m厚的磷化膜的工艺方法。磷化膜为多孔膜层,它能使基体表面层的吸附性、耐蚀性、减摩性得到改善。磷化膜与基体结合牢固,油漆、润滑油可以渗入到其孔隙中,故广泛用作钢铁工件油漆涂层的底层、冷变形加工过程中的减摩层和零件的防锈层。

(3)钝化处理

钝化处理是指经阳极氧化或化学氧化方法等处理后的金属零件,由活泼状态转变为不活泼状态的工艺方法,简称为钝化。钝化后的金属零件表面形成一层致密的氧化膜,可以起到防腐作用。

钝化往往作为电镀后的处理工艺,如对锌、镉、铜、银等金属镀层,经铬酸或重铬酸盐的处理,在其表面形成一层组织致密的氧化膜,使金属的腐蚀速度大为降低。如镀锌层钝化后,耐腐蚀能力一般可提高5倍以上。此外,钝化处理还能使镀层光亮美观,提高工件的装饰效果。

2.电火花表面强化和喷丸表面强化

(1)电火花表面强化

电火花表面强化是通过电火花放电的作用,把一种导电材料涂敷熔渗到另一种导电材料的表面,从而改变后者表面物理和化学性能的工艺方法。

电火花表面强化的原理如图3-13所示。在电极与工件之间接上直流电源或交流电源,振动器的作用使电极与工件之间的放电间隙频繁发生变化,电极与工件之间不断产生火花放电,从而使电极的材料黏结、覆盖在工件上,在工件表面形成强化层。

\begin{figure}[H]
\centering
\includegraphics[page=8, viewport={163.5 306.8 322.5 384.3}, clip]{12-07572-001FigRevised.pdf}
\caption{电火花表面强化的原理示意图}
\end{figure}

电火花表面强化具有设备简单、操作容易、成本低等优点,可用于模具、刀具及机械零件的表面强化和磨损部位的修补。例如,把硬质合金等材料涂敷在各类碳钢模具、刀具、量具及机械零件的表面,可以提高其表面硬度,增加耐磨性、耐蚀性,可提高使用寿命一至数倍。

(2)喷丸表面强化

喷丸表面强化是将大量高速运动的弹丸喷射到工件表面,使其表面产生强烈的塑性变形,从而获得一定厚度的强化层的工艺方法。

在这个过程中,弹丸犹如无数的小锤反复锤击工件表面,表面由于塑性变形产生一定厚度的冷作硬化层(也称喷丸表面强化层),这种表面强化层具有与内层材料完全不同的组织结构和应力状态。高密度的位错相互缠结阻碍塑性变形,提高了工件材料的屈服强度,工件表面的疲劳强度也大大提高;表面残余压应力可部分抵消引起工件疲劳破坏的循环拉应力,延缓疲劳裂纹的扩展,从而显著地提高了工件的抗疲劳性能。喷丸表面强化广泛应用于机械制造和航空制造等领域,如齿轮、连杆、飞机起落架和涡轮发动机中的关键承力件等都要经过喷丸表面强化处理。

3.离子注入材料表面改性技术

离子注入材料表面改性技术是20世纪70年代初逐渐发展起来的一种表面改性方法。它是把工件(金属、合金、陶瓷等)放在离子注入机中,在几十至几百千伏的电压下,把所需元素的离子注入工件表面的一种工艺。图3-14为离子注入机的结构示意图。作为离子流,目前用得较多的非金属有N、B、C;耐蚀抗磨合金元素有Ti、Cr、Ni;固体润滑元素有S、Mo、Sn、In;还有耐高温元素Y及稀土元素等。

\begin{figure}[H]
    \centering
    % Placeholder for Figure 3-14
    \includegraphics[page=8, viewport={337.5 304.3 530.0 407.8}, clip]{12-07572-001FigRevised.pdf}
    \caption{离子注入机的结构示意图}
\end{figure}

离子注入材料表面改性技术已在工业上得到广泛应用,并已取得良好的经济效益。例如,离子注入材料表面改性技术可使某些刀具、模具的寿命延长数倍,还可使钛合金人工关节的磨损速率下降,使用寿命延长10$\sim$15年。离子注入材料表面改性技术使航空精密轴承的耐蚀性得到显著改善,从而延长其使用寿命。

\section*{复习思考题}

 
3-1 钢热处理的基本原理是什么?其目的和作用是什么?

3-2 指出$A_{cm}$、$Ac_{cm}$、$Ar_{cm}$各临界温度线的含义。

3-3 什么是连续冷却与等温冷却?两种冷却方式有何差异?试绘出共析钢过冷奥氏体的这两种冷却方式示意图,并说明图中各个区域、各条线的含义。

3-4 奥氏体在等温冷却转变时,按过冷度的不同可以获得哪些组织?

3-5 简述珠光体、索氏体、贝氏体、马氏体的形成条件及特点。

3-6 简述正火的主要应用场合。

3-7 退火的主要目的是什么?常用的退火方法有哪些?

3-8 完全退火与不完全退火在加热规范、组织转变和应用场合上有何不同?为什么亚共析钢一般不采用不完全退火,过共析钢不采用完全退火?

3-9 正火与退火的主要区别是什么?如何选用?

3-10 淬火的目的是什么?常用的淬火方法有哪几种?

3-11 钢获得马氏体组织的条件是什么?马氏体相变有何特点?

3-12 为什么要进行表面淬火?常用的表面淬火方法有哪些?各应用在什么场合?

3-13 淬透性与淬硬深度、淬硬性有哪些区别?影响淬透性因素有哪些?

3-14 回火的目的是什么?常用的回火操作有哪些?试指出各种回火操作得到的组织、性能及应用范围。

3-15 什么是调质?调质的主要目的是什么?钢在调质处理后是什么组织?

3-16 一般采用什么方法改善低碳钢、中碳钢和高碳钢的切削加工性能?并说明原因。

3-17 甲乙两厂生产同一种零件,材料均为45钢。甲厂采用正火处理,乙厂采用调质处理,测试硬度均达到要求。试分析两种产品的组织和性能差异。

3-18 化学热处理的基本过程是什么?常用的化学热处理方法有哪些?其目的分别是什么?

3-19 请简述几种常用的表面热处理工艺。

3-20 渗碳、渗氮的目的是什么?试述其应用范围。

3-21 机床摩擦片用于传动或主轴刹车,要求耐磨性好。材料选用15钢渗碳后淬火,要求渗碳层0.4$\sim$0.5 mm,淬火、回火后硬度40$\sim$45 HRC。加工摩擦片的工艺路线如下:
下料$\rightarrow$机加工$\rightarrow$渗碳后淬火及回火$\rightarrow$机加工$\rightarrow$回火。
试说明其中各热处理工序的目的。

3-22 金属的氧化处理、磷化处理和钝化处理的目的各是什么?试述它们在工业上的实际应用情况。

3-23 电火花表面强化和喷丸表面强化的目的和基本原理是什么?

3-24 离子注入材料表面改性技术的基本原理是什么?试说明其应用情况。



\chapter{钢铁材料及其用途}

在机器制造中用得最多的材料是钢和铸铁。钢按照化学成分的差异可分为碳素钢(简称碳钢)和合金钢两大类。与合金钢相比,碳钢冶炼简便,价格低廉,在工程制造中占有很大的比例;合金钢是在碳钢的基础上添加某些合金元素以满足特殊性能要求,如不锈钢中含有铬(Cr)、耐热钢中含有钨(W)等。铸铁中碳的质量分数一般为2.5\%$\sim$4.0\%,具有优良的铸造性能和良好的耐磨性、减振性,广泛应用于机床和汽车等制造行业。

\section{碳钢与合金钢}

\subsection{碳钢}

1.碳钢的分类

碳钢的分类方法很多,可按冶炼方法、碳的质量分数、质量和用途分类。

(1)按冶炼方法分类

按冶炼炉的不同有平炉钢、转炉钢和电炉钢;按脱氧方法又可分为沸腾钢(F)、镇静钢(Z)和特殊镇静钢(TZ)。

(2)按碳的质量分数分类

低碳钢:碳的质量分数小于0.25\%;

中碳钢:碳的质量分数为0.25\%$\sim$0.60\%;

高碳钢:碳的质量分数大于0.60\%。

(3)按质量分类

主要根据碳钢中杂质硫和磷的质量分数进行分类。

普通钢:S的质量分数$\leqslant$0.050\%,P的质量分数$\leqslant$0.045\%;

优质钢:S的质量分数和P的质量分数均$\leqslant$0.035\%;

高级优质钢:S的质量分数$\leqslant$0.020\%,P的质量分数$\leqslant$0.030\%;

特级优质钢:S的质量分数$\leqslant$0.015\%,P的质量分数$\leqslant$0.025\%。

(4)按用途分类

碳素结构钢包括建筑工程结构用钢与机械工程结构用钢。前者多为低碳钢,焊接性较好,如桥梁、船舶、建筑用钢等,在热轧状态下很少进行热处理。后者主要用于制造机械零件,要求有良好的力学性能,碳的质量分数一般为0.60\%以下,使用前通常要进行热处理。

碳素工具钢的碳的质量分数较高,经适当热处理后,具有较高的强度、硬度和耐磨性,可用于制造量具、模具和低速切削刀具以及机器上各种耐磨零件。

2. 碳钢的牌号

(1) 碳素结构钢

碳素结构钢的牌号由代表屈服强度的字母、其屈服强度的数值、质量等级符号和脱氧方法符号四部分按顺序组成,如 Q235AF。其牌号、化学成分、性能及用途举例见表 4-1。

\begin{table}[H]
\centering
\caption{碳素结构钢的牌号、化学成分、性能及用途举例}
\label{tab:4-1}
\small
\begin{tabular}{|c|c|c|c|c|c|c|c|c|c|c|c|c|}
\hline
\multirow{2}{*}{牌号} & \multirow{2}{*}{\makecell{统一\\数字\\代号}} & \multirow{2}{*}{\makecell{质量\\等级}} & \multicolumn{5}{c|}{化学成分 $w/\%$,不大于} & \multirow{2}{*}{\makecell{脱氧\\方法}} & \multicolumn{3}{c|}{力学性能} & \multirow{2}{*}{用途举例} \\
\cline{4-8}\cline{10-12}
 & & & C & Mn & Si & S & P & & \makecell{$R_{\text{eH}}$\\/MPa} & \makecell{$R_{\text{m}}$\\/MPa} & \makecell{$A$\\/\%} & \\
\hline
Q195 & U11952 & --- & 0.12 & 0.50 & 0.30 & 0.040 & 0.035 & F、Z & 195 & 315$\sim$430 & 33 & \multirow{7}{2.5cm}{受力不大的零件，如螺钉、螺母、垫圈等，焊接件、冲压件及桥梁建筑等金属结构件}  \\
\cline{1-12}
\multirow{2}{*}{Q215} & U12152 & A & \multirow{2}{*}{0.15} & \multirow{2}{*}{1.20} & \multirow{2}{*}{0.35} & 0.050 & \multirow{2}{*}{0.045} & \multirow{2}{*}{F、Z} & \multirow{2}{*}{215} & \multirow{2}{*}{335$\sim$450} & \multirow{2}{*}{31} & \\
\cline{2-3}\cline{7-7}
 & U12155 & B & & & & 0.045 & & & & & & \\
\cline{1-12}
\multirow{4}{*}{Q235} & U12352 & A & 0.22 & \multirow{4}{*}{1.40} & \multirow{4}{*}{0.35} & 0.050 & \multirow{2}{*}{0.045} & \multirow{2}{*}{F、Z} & \multirow{4}{*}{235} & \multirow{4}{*}{370$\sim$500} & \multirow{4}{*}{26} & \\
\cline{2-4}\cline{7-7}
 & U12355 & B & 0.20 & & & 0.045 & & & & & & \\
\cline{2-4}\cline{7-9}
 & U12358 & C & \multirow{2}{*}{0.17} & & & 0.040 & 0.040 & Z & & & & \\
\cline{2-3}\cline{7-9}
 & U12359 & D & & & & 0.035 & 0.035 & TZ & & & & \\
\hline
\multirow{4}{*}{Q275} & U12752 & A & 0.24 & \multirow{4}{*}{1.50} & \multirow{4}{*}{0.35} & 0.050 & 0.045 & F、Z & \multirow{4}{*}{275} & \multirow{4}{*}{410$\sim$540} & \multirow{4}{*}{22} & \multirow{4}{2.5cm}{承受中等载荷的零件,如销子、连杆等} \\
\cline{2-4}\cline{7-9}
 & U12755 & B & 0.22 & & & 0.045 & 0.045 & Z & & & & \\
\cline{2-4}\cline{7-9}
 & U12758 & C & \multirow{2}{*}{0.20} & & & 0.040 & 0.040 & Z & & & & \\
\cline{2-3}\cline{7-9}
 & U12759 & D & & & & 0.035 & 0.035 & TZ & & & & \\
\hline
\end{tabular}
    % 表格下方的注释
\begin{tablenotes}
	\item 注：1. 表中符号：Q——屈服强度“屈”字汉语拼音首字母；A,B,C,D——质量等级；F——沸腾钢“沸”字汉语拼音首字母；Z——镇静钢“镇”字汉语拼音首字母；TZ——特殊镇静钢“特镇”两个字汉语拼音首字母。在牌号中Z, TZ符号可以省略。
	\item 2. 拉伸试验值适用于钢板厚度或直径 $\le 16$ mm 的钢材。
	\item 3. Q195 的屈服强度值仅供参考，不作为交货条件。
\end{tablenotes}
\end{table}

(2) 优质碳素结构钢

优质碳素结构钢的牌号采用两位数字表示,这两位数字为该钢的平均碳的质量分数的百分数,如 45 钢表示碳的质量分数为 0.45\%的优质碳素结构钢。优质碳素结构钢的牌号、化学成分、性能及用途举例见表 4-2。

(3) 碳素工具钢

碳素工具钢的牌号由字母 T 和数字组成。T 表示碳素工具钢,数字表示平均碳的质量分数的千分数,如 T10 表示平均碳的质量分数为 1\%的碳素工具钢。这类钢数字后面加 A 表示高级优质钢,如 T10A。碳素工具钢的牌号、化学成分、性能及用途举例见表 4-3。

\begin{table}[H]
\centering
\caption{优质碳素结构钢的牌号、化学成分、性能及用途举例}
\label{tab:4-2}
\footnotesize\setlength{\tabcolsep}{3pt}
\begin{tabular}{|c|c|c|c|c|c|c|c|c|c|c|}
\hline
\multirow{2}{*}{牌号} & \multicolumn{3}{c|}{化学成分 $w/\%$} & \multicolumn{6}{c|}{力学性能(不小于)} & \multirow{2}{*}{用途举例} \\
\cline{2-10}
 & C & Si & Mn & $R_{\text{m}}$ & $R_{\text{eH}}$ & \makecell{$A$\\/\%} & \makecell{$Z$\\/\%} & \makecell{硬度/HBW\\(热轧)} & \makecell{$a_{\text{K}}/$\\($\text{J}/\text{m}^2$)} & \\
\hline
08 & 0.05$\sim$0.11 & 0.17$\sim$0.37 & 0.35$\sim$0.65 & 325 & 195 & 33 & 60 & 131 & & \multirow{2}{4cm}{受力不大但要求高韧性的冲压件、焊接件、紧固件如螺母等} \\
10 & 0.07$\sim$0.13 & 0.17$\sim$0.37 & 0.35$\sim$0.65 & 335 & 205 & 31 & 55 & 137 & & \\
\hline
20 & 0.17$\sim$0.23 & 0.17$\sim$0.37 & 0.35$\sim$0.65 & 410 & 245 & 25 & 55 & 156 & & \multirow{2}{4cm}{渗碳淬火后强度要求不高的耐磨零件,如凸轮、滑块、活塞销等} \\
25 & 0.22$\sim$0.29 & 0.17$\sim$0.37 & 0.50$\sim$0.80 & 450 & 275 & 23 & 50 & 170 & & \\
\hline
30 & 0.27$\sim$0.34 & 0.17$\sim$0.37 & 0.50$\sim$0.80 & 490 & 295 & 21 & 50 & 179 & & \multirow{2}{4cm}{受力较大的零件,如曲轴、连杆、表面淬火齿轮等} \\
35 & 0.32$\sim$0.39 & 0.17$\sim$0.37 & 0.50$\sim$0.80 & 530 & 315 & 20 & 45 & 197 & & \\
\hline
40 & 0.37$\sim$0.44 & 0.17$\sim$0.37 & 0.50$\sim$0.80 & 570 & 335 & 19 & 45 & 217 & $6\times 10^5$ & \multirow{2}{4cm}{齿轮、曲轴、连杆、联轴器、轴} \\
45 & 0.42$\sim$0.50 & 0.17$\sim$0.37 & 0.50$\sim$0.80 & 600 & 355 & 16 & 40 & 229 & $5\times 10^5$ & \\
\hline
60 & 0.57$\sim$0.65 & 0.17$\sim$0.37 & 0.50$\sim$0.80 & 675 & 400 & 12 & 35 & 255 & & \multirow{2}{4cm}{弹性极限或强度要求高的零件如弹簧、弹簧垫圈、钢丝绳等} \\
65 & 0.62$\sim$0.70 & 0.17$\sim$0.37 & 0.50$\sim$0.80 & 695 & 410 & 10 & 30 & 255 & & \\
\hline
\end{tabular}\end{table}

\begin{table}[H]
\centering
\caption{碳素工具钢的牌号、化学成分、性能及用途举例}
\label{tab:4-3}
\footnotesize\setlength{\tabcolsep}{3pt}
\begin{tabular}{|c|c|c|c|c|c|c|c|c|c|p{4cm}|}
\hline
\multirow{2}{*}{牌号} & \multicolumn{5}{c|}{化学成分 $w/\%$} & \multicolumn{2}{c|}{淬火} & \multicolumn{2}{c|}{回火} & \multirow{2}{*}{用途举例} \\
\cline{2-10}
 & C & Mn & Si & \makecell{S(不\\大于)} & \makecell{P(不\\大于)} & \makecell{加热\\温度/$^\circ$C} & \makecell{硬度\\/HRC} & \makecell{加热\\温度/$^\circ$C} & \makecell{硬度\\/HRC} & \\
\hline
T7 & \makecell{0.65$\sim$\\0.74} & \multirow{6}{*}{$\leqslant 0.40$} & \multirow{6}{*}{$\leqslant 0.35$} & \multirow{3}{*}{0.03} & \multirow{3}{*}{0.035} & \makecell{800$\sim$820\\水淬} & 62$\sim$64 & 200$\sim$250 & 55$\sim$60 & 承受冲击、振动的工具,如锤头、锯、钻头、木工用凿子 \\
\cline{1-2}\cline{7-11}
T8 & \makecell{0.75$\sim$\\0.84} & & & & & \makecell{780$\sim$800\\水淬} & 63$\sim$65 & 150$\sim$240 & 55$\sim$60 & 硬度、耐磨性要求较高的工具,如加工木材用的铣刀、木工工具等 \\
\cline{1-2}\cline{5-11}
T10 & \multirow{2}{*}{\makecell{0.95$\sim$\\1.04}} & & & \multirow{2}{*}{0.02} & \multirow{2}{*}{0.03} & \multirow{4}{*}{\makecell{760$\sim$780\\水淬}} & \multirow{2}{*}{63$\sim$65} & \multirow{2}{*}{200$\sim$250} & \multirow{2}{*}{60$\sim$64} & \multirow{2}{4cm}{要求耐磨、不受强烈振动的工具,如丝锥、板牙、锯条、刨刀等} \\
T10A & & & & & & & & & & \\
\cline{1-2}\cline{5-6}\cline{8-11}
T13 & \multirow{2}{*}{\makecell{1.25$\sim$\\1.35}} & & & 0.03 & 0.035 & & \multirow{2}{*}{63$\sim$65} & \multirow{2}{*}{150$\sim$270} & \multirow{2}{*}{62$\sim$64} & \multirow{2}{4cm}{要求高硬度但不受冲击的工具,如锉刀、量具、刮刀等} \\
T13A & & & & 0.02 & 0.03 & & & & & \\
\hline
\end{tabular}\end{table}

\subsection{合金钢}

合金钢的品种很多,按化学成分分类可将其分为低合金钢(合金元素的质量分数$<5\%$)、中合金钢(合金元素的质量分数为 $5\%\sim 10\%$)和高合金钢(合金元素的质量分数$>10\%$);按用途分类可分为合金结构钢(主要用于制造重要机械零件和工程构件的钢)、合金工具钢(主要用于制造重要的工具及模具的钢)和特殊性能钢(主要用于制造有特殊物理、化学、力学性能要求的钢)。

目前,在合金钢中经常加入的合金元素有锰(Mn)、硅(Si)、铬(Cr)、镍(Ni)、钼(Mo)、钨(W)、钒(V)、钛(Ti)、硼(B)、磷(P)和稀土元素等。这些合金元素可以改善钢的综合力学性能、淬透性、热稳定性和耐蚀性等性能。

1. 合金元素在钢中的作用

(1) 合金元素可以提高钢的力学性能

合金元素之所以能提高钢的力学性能,其主要因素以下几个方面。

1) 固溶强化铁素体

大多数合金元素均能不同程度地溶解于铁素体中,使钢的强度、硬度升高,塑性和韧性下降。有些合金元素,如 Mn、Cr、Ni 等,只要配比得当,不仅能强化铁素体,而且能提高钢的韧性,使之具有良好的综合力学性能。

2) 第二相强化

具有比 Fe 与 C 更大亲和力的合金元素,除固溶于铁素体之外,还形成合金渗碳体[如 $(\text{FeMn})_3\text{C}$、$(\text{FeCr})_3\text{C}$ 等]以及碳化物(如 TiC、VC、MoC、WC 等)。这些组织均具有较高的硬度和稳定性,有的硬度可达 $71\sim 75$ HRC,因此提高了钢的强度、硬度和耐磨性。

3) 细晶强化

强碳化物形成元素 Nb、Ti、V 及强氮化物形成元素 Al,可形成稳定性高的碳化物、氮化物粒子,阻碍奥氏体晶粒的长大,细化铁素体晶粒。细晶粒的钢具有较好的力学性能,特别能显著提高钢的韧性。

(2) 合金元素可以改善钢的热处理工艺性能

1) 提高钢的淬透性

除 Co 以外,所有溶入奥氏体中的合金元素都能提升过冷奥氏体的稳定性,不同程度地使 C 曲线右移,降低了钢的临界冷却速度,从而提高了钢的淬火速度。因此,在获得同样淬硬深度的情况下,钢在淬火时则可用冷却能力较低的淬火冷却介质,这就降低了淬火内应力,减少了工件的变形和开裂。这对形状复杂的工件来说尤为重要。

2) 提高钢的回火稳定性

合金元素对钢的回火过程具有很大的影响。一般来说,合金元素使淬火钢回火时马氏体不易分解,阻碍碳化物聚集长大,并使发生这些转变的温度升高。这就使钢的硬度随回火温度而下降的程度减慢,即增加了回火稳定性。

回火稳定性越高,钢在比较高的温度下越能保持高的硬度和耐磨性。这对于用来制造高速切削刀具和热加工模具的材料有重要意义。如在高速切削时,刀具温度很高,若刀具材料的回火稳定性高,则可提高刀具的耐热性和使用寿命。

3) 使钢获得特殊性能

当向钢中加入一定量的某些合金元素时,钢的组织和性能将发生某些特殊的变化,形成具有某些特殊性能的合金钢,例如不锈钢、耐热钢、耐磨钢、磁钢等。

2. 合金结构钢

合金结构钢的牌号采用两位数字 + 元素符号 + 数字表示。前面两位数字表示钢中平均碳质量分数的万分数;合金元素后面的数字表示该元素平均质量分数的百分数,小于 1.5\%不标出。高级优质钢在钢牌号后加“A”;特级优质钢在钢牌号后加“E”。例如,18Cr2Ni4WA 表示高级优质结构钢,其碳的平均质量分数为 0.18\%,铬的平均质量分数为 2\%,Ni 的平均质量分数为 4\%,钨的平均质量分数小于 1.5\%。常用合金结构钢的牌号、性能及用途举例见表 4-4。

\begin{table}[H]
\centering
\caption{常用合金结构钢的牌号、性能及用途举例}
\label{tab:4-4}
\small
\begin{tabular}{|c|c|c|c|c|c|c|c|}
\hline
\multirow{2}{*}{钢种类别} & \multirow{2}{*}{牌号} & \multicolumn{2}{c|}{热处理} & \multicolumn{3}{c|}{力学性能(不小于)} & \multirow{2}{*}{用途举例} \\
\cline{3-7}
 & & \makecell{淬火温度\\/$^\circ$C} & \makecell{回火温度\\/$^\circ$C} & \makecell{$R_{\text{m}}$\\/MPa} & \makecell{$R_{\text{eL}}$\\/MPa} & \makecell{$A$/\%} & \\
\hline
\multirow{2}{*}{\makecell{低合金\\结构钢}} & 16Mn & --- & --- & 510$\sim$660 & 345 & 22 & \multirow{2}{4.5cm}{压力容器,桥梁,船舶,机车车辆,高、中压压力容器等} \\
 & 15MnTi & & & 530$\sim$680 & 390 & 20 & \\
\hline
\multirow{2}{*}{合金渗碳钢} & 20Cr & 880 水、油淬 & 200 & 835 & 540 & 10 & \multirow{2}{4.5cm}{齿轮,活塞销,凸轮气门顶杆,汽车、拖拉机上变速箱齿轮等} \\
 & 20CrMnTi & 880 油淬 & 200 & 1 080 & 850 & 10 & \\
\hline
\multirow{2}{*}{合金调质钢} & 40Cr & \multirow{2}{*}{850 油淬} & 520 & 980 & 785 & 9 & \multirow{2}{4.5cm}{机床主轴、曲轴、连杆、重要轴、齿轮等} \\
 & 35CrMo & & 550 & 980 & 835 & 12 & \\
\hline
\multirow{2}{*}{合金弹簧钢} & 60Si2Mn & 870 油淬 & 440 & 1 570 & 1 375 & $5(A_{11.3})$ & \multirow{2}{4.5cm}{汽车、拖拉机、铁道车辆上的板簧、螺旋弹簧,承受大应力的各种尺寸螺旋弹簧等} \\
 & 50CrV & 850 油淬 & 500 & 1 280 & 1 130 & $10(A_{11.3})$\rule{0cm}{0.8cm} & \\
\hline
\end{tabular}
\begin{tablenotes}
	\item 注:1. 力学性能值测试试样为热轧低合金结构钢;	
	\item 2. 部分数据摘自 GB/T 3077—2015。
\end{tablenotes}
\end{table}



低合金结构钢一般用作工程结构用钢。这类钢的强度较低,塑、韧性和焊接性较好,价格低廉,一般在热轧态下使用,必要时进行正火处理以提高强度。

机械制造用合金结构钢按组织和性能特点可分为以下几种类型。

(1) 合金调质钢

合金结构钢用于制造机械零件时应有较高的强度和良好的塑性及韧性。碳的质量分数为 0.3\%$\sim$0.6\%的合金结构钢经淬火、高温回火后得到具有回火索氏体组织的合金调质钢,可较好地满足性能要求。为保证淬透性,钢中常添加 Cr、Mn、Ni、Mo 等合金元素。常用钢种有 40Cr、35CrMo、40CrNiMo 等。

(2) 表面硬化钢

有些机械零件除了要求有较高强度及良好的塑性和韧性外,还要求表面硬度高,耐磨性好,可选用表面硬化钢。按表面硬化方式的不同,表面硬化钢可分为以下几种。

1) 合金渗碳钢

这类钢的碳的质量分数低(0.15\%$\sim$0.25\%),可保证工件心部获得低碳马氏体,有较高强度和韧性,而表层经渗碳后,硬度高而耐磨。热处理方式为渗碳后淬火加低温回火。常用钢种有 15Cr、20Cr、20CrMnTi 等。

2) 合金渗氮钢

合金渗氮钢的碳的质量分数中等(0.3\%$\sim$0.5\%)。工件渗氮前经过调质,心部为回火索氏体,有良好的强度和韧性。渗氮后,工件表层为硬而耐磨、耐腐蚀的氮化层,心部组织与性能仍保持氮化前的状态。常用钢种有 38CrMoAl、30CrMnSi 等。

(3) 合金弹簧钢

这类钢有高的弹性极限和高的疲劳强度,碳的质量分数为 0.5\%$\sim$0.85\%,经淬火、中温回火后得到托氏体组织。常用钢种有 65Mn 和 50CrV 等。

3. 合金工具钢和特殊性能钢

(1) 合金工具钢

合金工具钢具有高的硬度和耐磨性以及足够的强度和韧性,主要用于制造刀具、模具、量具等工具。合金工具钢和特殊性能钢牌号形式与合金结构钢相似,只是碳的质量分数的表示方法不同。当碳的平均质量分数大于或等于 1.0\%时不予标出,碳的平均质量分数小于 1.0\%时,以其千分数表示。其合金元素的表示方法与合金结构钢的相同。常用合金工具钢的牌号、性能及用途举例见表 4-5。

\begin{table}[H]
\centering
\caption{常用合金工具钢的牌号、性能及用途举例}
\label{tab:4-5}
\small
\begin{tabular}{|c|c|c|c|c|c|}
\hline
\multirow{2}{*}{钢种类别} & \multirow{2}{*}{牌号} & \multirow{2}{*}{\makecell{退火交货状态的\\钢材硬度/HBW}} & \multicolumn{2}{c|}{试样淬火硬度} & \multirow{2}{*}{用途举例} \\
\cline{4-5}
 & & & 淬火温度/$^\circ$C & 硬度/HRC & \\
\hline
\multirow{4}{*}{\makecell{低合金\\工具钢}} & 9SiCr & $197\sim 241^{\text{\textcircled{1}}}$ & 820$\sim$860 & $\geqslant 62^{\text{\textcircled{2}}}$ & \multirow{4}{*}{\makecell{板牙、钻头、铰刀、形状\\复杂的冲模;各种量规、量\\块;中型热锻模等,量具、\\刀具等}} \\
 & Cr2 & $179\sim 229$ & 830$\sim$860 & $\geqslant 62$ & \\
 & 5CrMnMo & $197\sim 241$ & 830$\sim$860 & b & \\
 & CrWMn & $207\sim 255$ & 800$\sim$830 & 62 & \\
\hline
\multirow{2}{*}{\makecell{高速\\工具钢}} & W18Cr4V & $\leqslant 255$ & 1 250$\sim$1 280 & $\geqslant 63$ & \multirow{2}{*}{\makecell{铣刀、车刀、钻头、刨\\刀等}} \\
 & W6Mo5Cr4V2 & $\leqslant 255$ & 1 200$\sim$1 230 & $\geqslant 64$ & \\
\hline
\end{tabular}
\begin{tablenotes}
	\item 注：\textcircled{1} 根据需方要求,并在合同中注明,制造螺纹刃具用钢为 $187\sim 229$ HBW;
	\item \textcircled{2} 根据需方要求,并在合同中注明,可提供实测值。
\end{tablenotes}
\end{table}

合金工具钢按合金元素的质量分数可分为低合金工具钢和高合金工具钢。

1) 低合金工具钢

低合金工具钢是在碳素工具钢的基础上加入少量合金元素(合金元素的总质量分数低于 5\%)制成,热处理方式和组织与碳素工具钢相同。这类钢的碳的质量分数高,通常为 0.9\%$\sim$1.5\%。由于合金元素的加入,低合金工具钢淬透性和回火稳定性得以提高,比碳素工具钢有较好的热硬性、韧性和较小的变形、开裂倾向。但低合金工具钢的热硬性不高,最高工作温度不超过 300 $^\circ$C,主要用于制造低速切削的刀具和量规、量块等量具以及形状较为复杂的冷冲模、拉丝模等模具。常用钢种有 9SiCr 和 CrWMn 等。

2) 高合金工具钢

由于有较多合金元素的加入,高合金工具钢的淬透性、热硬性大大提高。按性能特点其可分为高速工具钢和高铬模具钢。

\textcircled{1} 高速工具钢。简称高速钢,其碳的质量分数为 0.7\%$\sim$1.5\%,含有较多的 W、Cr、Mo、V 等合金元素,总的质量分数超过 10\%。高速钢的牌号与一般合金工具钢稍有不同,当钢中碳的质量分数小于 1.0\%时,也不予标出。例如 W18Cr4V(读作钨 18 铬 4 钒),表示钢中碳的质量分数为 0.7\%$\sim$0.8\%,W、Cr 和 V 的质量分数分别为 18\%、4\%和小于 1.5\%。

高速钢最终热处理采用淬火加多次高温回火,组织为马氏体和碳化物。高速钢刀具在 600 $^\circ$C 切削温度下,仍保持 60 HRC 左右的高硬度,因此适于高速切削,所谓高速钢的名称也由此而来。高速钢刃磨性能良好,比一般低合金工具钢锋利,故又俗称为“锋钢”。

\textcircled{2} 高铬模具钢。铬在碳的质量分数高的过共析钢中,除以固溶体和合金渗碳体的形式存在外,还能形成弥散分布的碳化物,有效细化晶粒,均匀组织,提高淬透性和耐蚀性。由于具有硬度均匀,耐磨性好,疲劳强度高及尺寸稳定、易切削加工等特点,高铬模具钢广泛用于制作形状复杂的大截面模具等。常用钢种有 Cr12 和 Cr12MoV 等。

(2) 特殊性能钢

特殊性能钢是指不锈钢、耐热钢、耐磨钢等具有特殊的物理和化学性能的钢。

1) 不锈钢

具有抵抗大气或弱腐蚀介质侵蚀能力的钢称为不锈钢。实际上,不锈钢在大气中还是会发生腐蚀的,所谓“不锈”只是相对而言。目前常用的不锈钢按其组织状态不同可以分为马氏体不锈钢、铁素体不锈钢、奥氏体不锈钢、奥氏体-铁素体不锈钢和沉淀硬化型不锈钢五类,如 95Cr18、06Cr13Al、12Cr18Ni9、14Cr18Ni11Si4AlTi、07Cr17NiAl 等。

2) 耐热钢

在高温下具有良好化学稳定性和较高强度的钢称为耐热钢。按性能可分为热稳定钢和热强钢。按显微组织可分为马氏体型耐热钢、铁素体型耐热钢、奥氏体型耐热钢和沉淀硬化型耐热钢,如 20Cr13、06Cr13Al、06Cr18Ni11Ti 和 07Cr17Ni7Al 等。

适量的 Cr、Si、Al 等元素加入钢中后,在高温下与氧接触,均能生成结构致密的高熔点氧化物,覆盖于钢的表面,保护钢免受高温气体的继续氧化腐蚀;W、Mo、V、Cr、Ni、B 等元素,可使钢具有良好的组织稳定性,在高温下具有较高的强度。

3) 耐磨钢

耐磨钢一般是指在冲击载荷作用下能发生冲击硬化的高锰钢。它的主要成分:碳的质量分数为 1.0\%$\sim$1.3\%,锰的质量分数为 11\%$\sim$14\%。耐磨钢采用铸造成形,只有经热处理后获得全部奥氏体组织,才能呈现出良好的韧性和耐磨性。常用牌号如 ZGMn13、ZGMn13Cr2 等。

高锰钢广泛应用于制造承受较大冲击或压力的零部件,例如挖掘机的铲斗、坦克车履带等。此外,高锰钢在寒冷气候条件下不冷脆,适于高寒地区使用。

\section{铸铁}

\subsection{铸铁的成分和性能}

铸铁是碳的质量分数大于 2.11\%的铁碳合金。它还含有硅、锰、硫、磷等元素,有时还加入一定量的铬、钼、钒、铝等合金元素。与钢的化学成分相比较,铸铁含碳、硅和杂质元素硫、磷较多。工业上常用铸铁的碳质量分数一般为 2.5\%$\sim$4.0\%。

铸铁的抗拉强度、塑性和韧性虽然比钢差,又不能进行锻造,但它却具有良好的铸造性能、耐磨性、减振性等一系列优良性能,而且生产简便,价格低廉,所以在机械制造中得到广泛应用。据统计,按质量百分比计算,铸铁件在汽车、拖拉机中占 50\%$\sim$70\%,机床中占 60\%$\sim$90\%。

\subsection{铸铁的分类及其性能}

铸铁根据其试样断口色泽的不同,可以分为白口铸铁、灰口铸铁和麻口铸铁。白口铸铁中的碳绝大部分以渗碳体的形式存在,所以断口呈银白色。白口铸铁性能硬而脆,很难进行切削加工,所以工业上一般不用来制作机器零件,仅可用于要求耐磨而不受冲击的制件,如轧辊、磨球等,其主要用途是炼钢的原料。白口铸铁根据需要也可经过热处理制成可锻铸铁。工业中所用铸铁几乎都是灰口铸铁。麻口铸铁由于质地硬脆,所以工业生产中也很少应用。根据铸铁组织中石墨存在形态的差异,灰口铸铁又划分为灰铸铁、可锻铸铁、球墨铸铁和蠕墨铸铁等。还有一类是特殊性能的合金铸铁。

1. 灰铸铁

灰铸铁中的碳大部分或全部以片状石墨形式存在,断口呈暗灰色,是应用最广泛的一种铸铁,在铸铁件总产量中占 80\%以上。其微观组织是由钢的基体与片状石墨组成的。基体组织特征以及石墨的析出量与铸铁的石墨化程度有关。

(1) 铸铁的石墨化过程

铸铁的石墨化就是铸铁中碳原子析出和石墨形成的过程。在铁碳合金中,碳可以化合状态的渗碳体($\text{Fe}_3\text{C}$)和游离状态的石墨(常用 G 来表示)两种形式存在。渗碳体是亚稳定组织,石墨是一种稳定相;在长时间高温加热时,渗碳体会分解为铁和石墨,即 $\text{Fe}_3\text{C}\rightarrow 3\text{Fe}+\text{G}$。在铁碳合金结晶过程中,在缓慢冷却下,可以从液态金属中直接析出稳定的 G 相。

(2) 影响铸铁石墨化的因素

影响铸铁石墨化的因素很多,主要是化学成分和冷却速度。C、Si、Mn、S、P 对石墨化有不同的影响,C、Si、P 是促进石墨化的元素,Mn 和 S 是阻碍石墨化的元素。其中,铸铁中的 C 和 Si 是强烈促进石墨化的元素。铁液中 C 的浓度增加,可以加速石墨晶核的生成,促进石墨化。C、Si 的质量分数过低,铸铁中易出现白口组织,力学性能与铸造性能均较差,但 C、Si 的质量分数过高,铸铁中石墨数量多且粗大,力学性能也会下降。为了防止铸铁中形成白口组织或粗大的片状石墨,C、Si 的质量分数一般分别为 2.5\%$\sim$4.0\%和 1.0\%$\sim$3.0\%。

铸件的冷却速度越缓慢,即过冷度越小,越有利于石墨化过程的充分进行;反之,冷却速度越快,原子扩散能力减弱,越不利于石墨化的进行。在铸型材料和浇注温度等条件相同时,冷却速度与铸件壁厚有关。铸件壁越厚,冷却速度越慢。就同一铸件而言,由于各部位壁厚不同,冷却速度也不同,所得到的组织也不同。

(3) 灰铸铁的组织与性能。

灰铸铁按基体的不同可分为以下三种。

1) 珠光体灰铸铁是在珠光体基体上分布着细小均匀石墨的灰铸铁。其力学性能较高,用于制造较为重要的零件。

2) 珠光体-铁素体灰铸铁是在珠光体与铁素体的混合基体上分布着较粗大石墨片的灰铸铁。此种铸铁强度较低,但能满足一般零部件的要求,且铸造性能及切削加工性能较好,用途最广。

3) 铁素体灰铸铁是在铁素体基体上分布着粗大而且较多石墨片的灰铸铁。其力学性能较差,很少应用。

在实际铸铁铸件生产过程中,为了获得珠光体灰铸铁常采用孕育处理。孕育处理是在浇注前,在铁液中加入孕育剂的工艺方法。常用的孕育剂是硅的质量分数为 75\%的硅铁,加入量为铁液质量的 0.25\%$\sim$0.6\%。

三种灰铸铁显微组织如图 4-1 所示。这种基体显微组织的差异取决于石墨化过程。

\begin{figure}[H]
\centering
\includegraphics[page=8, viewport={62.5 153.3 480.5 272.8}, clip]{12-07572-001FigRevised.pdf}
\caption{灰铸铁显微组织图}
\label{fig:4-1}
\end{figure}


灰铸铁的性能,除与基体组织有关外,主要取决于石墨的形态、数量、大小和分布。石墨的力学性能较差,其强度、硬度、塑性均极低($R_{\text{m}}<20$ MPa、硬度约为 3 HBW,$A=0$),且密度小,在铸铁中占有很大的体积。所以,灰铸铁的组织因石墨的存在,如同在钢的基体上分布着无数的孔洞,片状石墨引起应力集中,致使其抗拉强度、塑性、韧性比钢差得多。石墨片越多、越粗大、分布越集中,对其性能的影响就越大。但在承受压力作用时石墨引起的不利影响较小,所以灰铸铁的抗压强度仍较高。

灰铸铁中石墨片虽然因割裂基体导致力学性能下降,但也带来一系列其他方面的优良性能,如良好的铸造性能、耐磨性、减振性、切削性能及低的缺口敏感性。

(4) 灰铸铁的牌号和用途

灰铸铁的牌号由“HT”加数字组成。其中“HT”是“灰铁”两字汉语拼音的首字母,后面的数字表示其最低抗拉强度值。常用灰铸铁的牌号、性能及用途举例见表 4-6。

\begin{table}[H]
\centering
\caption{常用铸铁的牌号、性能及用途举例}
\label{tab:cast_iron}
\begin{tabular}{|c|c|c|c|c|m{8cm}|}
\hline
分类 & 牌号 & \makecell{$R_m$ \\ /MPa} & \makecell{$A$ \\ /\%} & \makecell{硬度 \\ /HBW} & 性能及用途举例 \\
\hline
\multirow{4}{*}{灰铸铁} & HT100 & 80\textasciitilde130 & --- & 82\textasciitilde166 & 铸造性能好、工艺简便、减振性能优良，适于载荷很小，对摩擦、磨损无特殊要求的零件，如盖、外罩、油盘、手轮、支架、底板、重锤等 \\
\cline{2-6}
 & HT150 & 120\textasciitilde175 & --- & 105\textasciitilde205 & 性能特点与HT100基本相同，用于承受中等载荷的零件，如支柱、机座、箱体、法兰、泵体、阀体、轴承座、工作台、带轮等 \\
\cline{2-6}
 & HT200 & 160\textasciitilde220 & --- & 124\textasciitilde236 & \multirow{2}{=}{强度较高，耐热性、耐磨性较好，减振性良好；铸\\ 造性能好，但铸件需进行人工时效，适用于承受较大载荷或较为重要的零件，如气缸体、气缸盖、活塞、刹车轮、联轴器盘、油缸、泵体、阀体、齿轮、机座、机床床身及立柱等} \\
\cline{2-5}
 & HT250 & 200\textasciitilde270 & --- & 150\textasciitilde262 \rule{0pt}{60pt}&  \\
\hline
\multirow{3}{*}{可锻铸铁} & \makecell{KTH300-06 \\ （黑心 \\ 可锻铸铁）} & 300 & 6 & $\leqslant 150$ & 管道的弯头、接头、三通、中压阀门 \\
\cline{2-6}
 & \makecell{KTB350-04 \\ （白心 \\ 可锻铸铁）} & 350 & 4 & 230 & 适用于制作厚度小于15 mm的薄壁铸件和焊后不需要进行热处理的零件 \\
\cline{2-6}
 & \makecell{KTZ450-06 \\ （珠光体 \\ 可锻铸铁）} & 450 & 6 & 150\textasciitilde200 & 曲轴、凸轮轴、连杆、齿轮、摇臂、活塞环、轴套、犁刀、耙片、万向接头、棘轮、扳手、传动链条、矿车轮等 \\
\hline
\multirow{4}{*}{球墨铸铁} & QT400-18 & 400 & 18 & 120\textasciitilde175 & 泵、阀体、受压容器、受冲击零件等 \\
\cline{2-6}
 & QT450-10 & 450 & 10 & 160\textasciitilde210 & 壳、箱体零件，要求韧性零件等 \\
\cline{2-6}
 & QT500-7 & 500 & 7 & 170\textasciitilde230 & 机器底座、齿轮、支架等 \\
\cline{2-6}
 & QT800-2 & 800 & 2 & 245\textasciitilde335 & 曲轴、凸轮、齿轮等高强度零件等 \\
\hline
\multirow{3}{*}{蠕墨铸铁} & RuT500 & 500 & 0.5 & \makecell{220\textasciitilde260 \\ (HBW)} & 高负荷内燃机缸体、气缸套 \\
\cline{2-6}
 & RuT450 & 450 & 1.0 & \makecell{200\textasciitilde250 \\ (HBW)} & 汽车内燃机缸体与缸盖、气缸套、载重卡车制动盘、泵体和液压件、玻璃模具、活塞环 \\
\cline{2-6}
 & RuT300 & 300 & 2.0 & \makecell{140\textasciitilde210 \\ (HBW)} & 排气歧管，大功率船用、机车、汽车和固定式内燃机缸盖，增压器壳体，纺织机，农机零件 \\
\hline
\multirow{2}{*}{合金铸铁} & MTPCuTi15 & $\geqslant 150$ & --- & 170\textasciitilde229 & 机床导轨、气缸套、活塞环、凸轮轴等 \\
\cline{2-6}
 & RTCr & 200 & --- & 189\textasciitilde288 & 炉条、高炉支梁式水箱、金属型玻璃模 \\
\hline
\end{tabular}\end{table}

2. 可锻铸铁

可锻铸铁是用碳、硅的质量分数较低的铁液浇注成白口铸铁件，再经过长时间的高温退火，使渗碳体分解成团絮状的石墨而制成的。可锻铸铁仅意味着比灰铸铁具有较好的塑性、韧性，实际上并不能进行锻造。

由于在退火过程中共析反应时的冷却速度不同，可锻铸铁的基体组织可分为铁素体（黑心）和珠光体（白心）两种，其显微组织如图4-2所示。前者具有较高的塑性、韧性；后者则强度、硬度较高，耐磨性好。

\begin{figure}[H]
\centering
% Placeholder for images
\includegraphics[page=9, viewport={103.0 606.3 441.0 744.8}, clip]{12-07572-001FigRevised.pdf}
\caption{可锻铸铁显微组织图}
\end{figure}

由于团絮状的石墨对基体金属的割裂作用和引起的应力集中比片状石墨小，所以可锻铸铁比灰铸铁有较高的强度、塑性和韧性。但可锻铸铁退火工艺复杂，生产周期长，成本较高，目前主要用于制造一些形状较复杂而在工作中又承受冲击振动的薄壁小型零件。可锻铸铁的牌号、性能和用途见表4-6。其中“KT”是“可铁”两字汉语拼音的第一个字母，“H”表示黑心可锻铸铁，“B”表示白心可锻铸铁，“Z”表示珠光体基体，后面两组数字分别表示最低抗拉强度和最低伸长率。

3. 球墨铸铁

在浇注前向铁液中加入适量的球化剂和孕育剂，使碳呈球状析出，所获得的铸铁称为球墨铸铁。常用的球化剂为镁或稀土镁合金（镁和稀土的质量分数小于10\%，其余为硅和铁），球化剂的加入量一般为铁液质量的1.3\%\textasciitilde1.8\%。

根据基体组织的不同，球墨铸铁主要有铁素体球墨铸铁和珠光体球墨铸铁，其显微组织如图4-3所示。前者具有较高的塑性，后者则强度较高。

球状石墨的形态比团絮状石墨更为圆整，因而对基体金属的不利影响更小。所以，球墨铸铁的强度、塑性、韧性远远超过灰铸铁，甚至优于可锻铸铁，同时还具有良好的铸造性能、耐磨性和切削加工性等优点。此外，球墨铸铁还可以通过热处理来提高力学性能。因此，球墨铸铁在机械制造中应用广泛，用来制造各种受力复杂、负荷较大和耐磨的重要零件，在许多场合成功地替代了铸钢件、锻钢件。

球墨铸铁的牌号、性能和用途举例见表4-6。其中“QT”是“球铁”汉语拼音的第一个字母，后面两组数字分别表示最低抗拉强度和最低伸长率。

\begin{figure}[H]
\centering
% Placeholder for images
\includegraphics[page=9, viewport={103.5 432.8 440.5 569.3}, clip]{12-07572-001FigRevised.pdf}
\caption{球墨铸铁显微组织图}
\end{figure}

4. 蠕墨铸铁

蠕墨铸铁是一种新型高强度铸铁，石墨形状短、厚，端部圆滑，呈蠕虫状，其显微组织如图4-4所示。蠕墨铸铁保留了灰铸铁优良的工艺性能，力学性能介于相同基体组织的灰铸铁与球墨铸铁之间。蠕墨铸铁一般不进行热处理。

蠕墨铸铁的牌号用汉语拼音缩写“RuT”加一组最低抗拉强度的数字表示。蠕墨铸铁的牌号、性能和用途举例见表4-6。

\begin{figure}[H]
\centering
% Placeholder for image
\includegraphics[page=9, viewport={95.0 233.8 257.5 362.8}, clip]{12-07572-001FigRevised.pdf}
\caption{蠕墨铸铁显微组织$\times 200$}
\end{figure}

5. 合金铸铁

为了满足工业上对铸铁性能的特殊要求，如耐磨性、耐热性和耐蚀性等，向铸铁中加入某些合金元素，从而获得具有特殊性能的合金铸铁。

（1）耐磨铸铁

1）耐磨铸铁（MT）

耐磨铸铁是在铸铁基体中加入Cr、Mo、Cu等少量合金元素，从而获得的具有高耐磨性的铸铁，主要用于机床导轨、汽车发动机气缸套、活塞环等耐磨零件。

2）冷硬铸铁（LT）

冷硬铸铁是在铸铁表面通过激冷处理形成一层白口层，使表面具有高硬度和高耐磨性的铸铁。其主要用于制作轧辊、凸轮轴等零件。

（2）耐热铸铁（RT）

耐热铸铁是在铸铁基体中加入Al、Si、Cr等合金元素，得到高耐热性的铸铁。其主要用于制作炉底、换热器、坩埚和热处理炉内的运输链条等零件。

（3）耐蚀铸铁（ST）

耐蚀铸铁是在铸铁基体中加入Al、Si、Cr等合金元素，使其表面形成一层连续致密的保护膜，从而获得的具有高耐蚀能力的铸铁。其用于制作在腐蚀介质中工作的零件，如化工设备的管道、阀门、泵体、反应釜和盛储器等。

\section*{复习思考题}

 
4-1 碳钢按质量分数和用途可分为几类？试举例说明。

4-2 分析说明优质碳素结构钢的碳质量分数差异会造成较大的性能差异的原因。

4-3 合金钢中经常加入的合金元素有哪些？这些合金元素对钢的力学性能和热处理工艺有何影响？

4-4 什么是合金调质钢？其主要用途是什么？其所含主要合金元素有哪些？

4-5 优质碳素钢的牌号是如何规定的？45钢的化学成分、抗拉强度、屈服强度分别是多少？

4-6 什么是合金渗碳钢和合金弹簧钢？常用的钢种有哪些？

4-7 试述不锈钢、耐热钢以及耐磨钢的主要合金元素。

4-8 简述高速钢的特点。

4-9 碳素工具钢的牌号是如何规定的？碳素工具钢的主要用途是什么？

4-10 指出下列合金钢的类别、用途和各合金元素的质量分数及主要作用。

（1）40CrNiMo；（2）60Si2Mn；（3）9SiCr；（4）Cr12MoV；（5）37CrNi3

4-11 与钢相比，铸铁具有哪些特点？

4-12 简述铸铁中主要包含的元素。

4-13 灰口铸铁可分哪几类？影响其组织和性能的因素有哪些？

4-14 试述灰口铸铁、球墨铸铁、可锻铸铁的主要用途，并说明石墨形态的差异会带来性能差异的原因？

4-15 简述白口铸铁的特点。

4-16 三块形状和颜色完全相同的铁碳合金，分别是15钢、60钢和白口铸铁，用什么简便方法可以迅速区分它们？

4-17 灰铸铁中石墨是以什么形态存在于钢的基体中？对铸铁的使用性能有何影响？

4-18 简述可锻铸铁基体组织的种类、特点及其形成原因。

4-19 什么是石墨化过程说明？影响石墨化的主要因素。

4-20 灰铸铁按基体组织的差异可分为哪三种类型？性能有何差异？

4-21 说明下列铸铁牌号的含义。

（1）HT200；（2）HT250；（3）KTH350-10；（4）KTZ550-04；（5）KTB400-05；（6）QT700-2

4-22 如何使铸铁中石墨以团絮状形态存在？可锻铸铁是否“可锻”？

4-23 如何通过“球化处理”使铸铁中石墨以球状形态存在？球墨铸铁常用的热处理工艺有哪些？

4-24 什么是合金铸铁？耐蚀铸铁主要添加哪些合金元素？

4-25 特殊性能钢主要有哪些？一般是通过添加什么元素实现其特殊性能的？


\chapter{非铁材料}

金属材料分为钢铁金属（钢和铸铁，又称黑色金属）和非铁金属（钢铁以外的所有金属，又称有色金属）两大类。非铁金属包括Al、Cu、Ni、Ti、Co、Zn、Mg等，贵金属（铂族金属），难熔金属（W、Mo、Ta、Nb），稀土金属以及由它们组成的合金。非铁材料是指除钢铁金属以外的其他材料，主要包括非铁金属及其合金、高分子材料、陶瓷材料、复合材料和其他新材料等。

\section{非铁金属及其合金}

\subsection{铝及铝合金}

纯铝熔点为660 $^\circ$C，面心立方晶格，无同素异构转变。纯铝呈银白色，密度为2.7 g/cm$^3$，具有良好的导电性、导热性，塑性好，但强度不高，不宜用来制作承力结构件，主要用来制造电线，电缆，强度要求不高的器皿、用具以及配制各种铝合金等。工业纯铝的牌号用1$\times\times\times$系列表示。牌号的最后两位数字表示最低铝百分含量。当最低铝百分含量精确到0.01\%时，牌号的最后两位数字就是最低铝百分含量中小数点后面的两位。牌号第二位的字母表示原始纯铝的改型情况。如果第二位字母为A，则表示为原始纯铝；如果是B\textasciitilde Y的其他字母，则表示为原始纯铝的改型，与原始纯铝相比，其元素含量略有不同。例如，1A99中99表示铝的百分含量为99.99\%。

铝合金按其加工方法可分为变形铝合金和铸造铝合金。图5-1是铝合金的分类示意图。$D$点以左的铝合金加热后均可形成单相固溶体，具有良好的塑性，适合进行变形加工，故称为变形铝合金。

$D$点以右的铝合金，由于其组织中含有低熔点的共晶体，流动性好，适于铸造，故称为铸造铝合金。

1. 变形铝合金

变形铝合金通常经不同的变形加工方式生产成各种半成品，如板、棒、管、线、型材及锻件等。铝合金根据其特性，可分为防锈铝合金、硬铝合金、超硬铝合金、锻铝合金四类。

防锈铝合金主要是Al-Mg和Al-Mn合金，其特点是耐蚀性好，易于加工成形和焊接，但其强度较低，不能通过热处理强化；硬铝合金主要有Al-Cu-Mg和Al-Cu-Mn合金，其特点是具有极强的时效硬化能力，强度高，耐蚀性和焊接性较差；超硬铝合金主要是Al-Cu-Mg-Zn合金，在硬铝合金的基础上加锌而成，强度高于硬铝合金，故称超硬铝合金；锻铝合金主要是Al-Mg-Si-Cu和Al-Cu-Mg-Ni-Fe合金，其特性是具有良好的热塑性、铸造性能和较高的力学性能。常用变形铝合金的牌号、化学成分、力学性能及用途见表5-1。

\begin{figure}[H]
\centering
% Placeholder for Figure 5-1
\includegraphics[page=9, viewport={282.0 235.3 453.5 397.3}, clip]{12-07572-001FigRevised.pdf}
\caption{铝合金的分类示意图}
\end{figure}

\begin{table}[H]
\centering
\caption{常用变形铝合金的牌号、化学成分、力学性能及用途}
\label{tab:alum_alloys}
\footnotesize\setlength{\tabcolsep}{3pt}
\begin{tabular}{|c|c|c|c|c|c|c|c|c|c|c|c|p{3cm}|}
\hline
\multirow{2}{*}{类别} & \multicolumn{2}{c|}{牌号} & \multicolumn{5}{c|}{化学成分 $w$/\%} & \multirow{2}{*}{\makecell{热处 \\ 理状态}} & \multicolumn{3}{c|}{力学性能} & \multirow{2}{*}{用途} \\
\cline{2-8}\cline{10-12}
 & 新国标 & 旧国标 & Cu & Mg & Mn & Zn & Al & & \makecell{$R_m$ \\ /MPa} & \makecell{$A$ \\ /\%} & \makecell{硬度 \\ /HBS} & \\
\hline
\multirow{3}{*}{\makecell{防锈 \\ 铝合金}} & 5A05 & LF5 & 0.1 & 4.8\textasciitilde5.5 & 0.3\textasciitilde0.6 & & \multirow{10}{*}{余量} & \multirow{3}{*}{M} & 270 & 23 & 70 & \multirow{3}{*}{\makecell{中等载荷零件、铆钉 \\ 以及焊接油箱、油管等}} \\
\cline{2-7}\cline{10-12}
 & 5A02 & LF2 & 0.1 & 2.0\textasciitilde2.8 & 0.15\textasciitilde0.4 & & & & 195 & 25 & 47 & \\
\cline{2-7}\cline{10-12}
 & 3A21 & LF21 & 0.2 & 0.05 & 1.0\textasciitilde1.6 & 0.10 & & & 110 & 30 & 28 & \\
\hline
\multirow{3}{*}{\makecell{硬铝 \\ 合金}} & 2A01 & LY1 & 2.2\textasciitilde3.0 & 0.2\textasciitilde0.5 & 0.2 & 0.10 & & \multirow{3}{*}{CZ} & 300 & 24 & 70 & \makecell{中等强度和工作温 \\ 度$\leqslant 100$ $^\circ$C的铆钉材料} \\
\cline{2-7}\cline{10-13}
 & 2A11 & LY11 & 3.8\textasciitilde4.8 & 0.4\textasciitilde0.8 & 0.4\textasciitilde0.8 & 0.30 & & & 420 & 18 & 100 & \makecell{中等强度结构件和零 \\ 件，如骨架、螺旋桨叶 \\ 片、铆钉等} \\
\cline{2-7}\cline{10-13}
 & 2A12 & LY12 & 3.8\textasciitilde4.9 & 1.2\textasciitilde1.8 & 0.3\textasciitilde0.9 & 0.30 & & & 470 & 20 & 105 & \makecell{高强度构件以及150 $^\circ$C \\ 以下工作的零件，如骨架、 \\ 梁等} \\
\hline
\multirow{2}{*}{\makecell{超硬 \\ 铝合金}} & 7A04 & LC4 & 1.4\textasciitilde2.0 & 1.8\textasciitilde2.8 & 0.2\textasciitilde0.6 & 5\textasciitilde7 & & \multirow{2}{*}{CS} & 230 & 17 & 60 & \multirow{2}{*}{\makecell{结构中主要受力件，如 \\ 飞机大梁、桁架、加强框 \\ 及起落架等}} \\
\cline{2-7}\cline{10-12}
 & 7A09 & LC9 & 1.2\textasciitilde2.0 & 2.0\textasciitilde3.0 & 0.15 & 5.1\textasciitilde6.1 & & & 570 & 11 & 150 \rule{0pt}{20pt} & \\
\hline
\multirow{3}{*}{\makecell{锻铝 \\ 合金}} & 2A50 & LD5 & 1.8\textasciitilde2.6 & 0.4\textasciitilde0.8 & 0.4\textasciitilde0.8 & 0.30 & & \multirow{3}{*}{CS} & 420 & 13 & 105 & \makecell{形状复杂和中等强度 \\ 的锻件和模锻件等} \\
\cline{2-7}\cline{10-13}
 & 2A70 & LD7 & 1.9\textasciitilde2.5 & 1.4\textasciitilde1.8 & 0.20 & 0.30 & & & 440 & 12 & 120 & \makecell{高温下工作的复杂锻 \\ 件和结构件、内燃机活 \\ 塞等} \\
\cline{2-7}\cline{10-13}
 & 2A14 & LD10 & 3.9\textasciitilde4.8 & 0.4\textasciitilde0.8 & 0.4\textasciitilde1.0 & 0.30 & & & 480 & 19 & 135 & \makecell{形状简单和高载荷的 \\ 锻件和模锻件} \\
\hline
\end{tabular}
\begin{tablenotes}
	\item 注：1. 新国标指GB/T 3190—2020；
	\item 2. 热处理代号：M表示退火；CZ表示淬火+自然时效；CS表示淬火+人工时效。
\end{tablenotes}
\end{table}

2. 铸造铝合金

铸造铝合金的优点是密度小，比强度较高，且具有良好的耐蚀性和铸造性能。铸造铝合金按成分不同可分为Al-Si系合金、Al-Cu系合金、Al-Mg系合金、Al-Zn系合金和Al-Re系合金。

Al-Si系合金又称硅铝明，硅的质量分数一般为10\%\textasciitilde13\%，这类合金具有优良的铸造性能，加入Cu、Mg、Mn等元素可使其强化，通过热处理可进一步提高其力学性能；Al-Cu系合金具有较高的强度和耐热性，但铸造性能差；Al-Mg系合金属于高强度和高耐腐蚀合金；Al-Zn系合金具有自淬火效应，铸造后经变质处理和时效处理后强度较高，但耐蚀性差，热裂倾向大；Al-Re系合金以Al-Si系合金为基础再加入适量稀土元素，因而具有良好的铸造性能和耐热性能。

铸造铝合金的牌号用“ZAl”表示“铸铝”，其后标注合金元素及其质量分数（百分含量）；若平均质量分数小于0.5\%则不予标出；牌号后标注A表示优质合金，如ZAlSi5Cu1MgA。铸造铝合金代号由“ZL”+三位数字表示，第一位数字代表合金系，1\textasciitilde4依次代表Al-Si、Al-Cu、Al-Mg和Al-Zn合金系，后两位数字为顺序号，如ZL101表示铸造Al-Si合金系。常用铸造铝合金的牌号、化学成分、力学性能及用途见表5-2。

\begin{table}[H]
\centering
\caption{常用铸造铝合金的牌号、化学成分、力学性能及用途}
\label{tab:cast_alum}
\footnotesize\setlength{\tabcolsep}{3pt}
\begin{tabular}{|c|c|c|c|c|c|c|c|c|c|c|c|c|p{3cm}|}
\hline
\multirow{2}{*}{类别} & \multirow{2}{*}{牌号} & \multirow{2}{*}{代号} & \multicolumn{5}{c|}{化学成分 $w$/\%} & \multirow{2}{*}{\makecell{铸造 \\ 方法}} & \multirow{2}{*}{\makecell{热 \\ 处理}} & \multicolumn{3}{c|}{\makecell{力学性能 \\ （不小于）}} & \multirow{2}{*}{用途} \\
\cline{4-8}\cline{11-13}
 & & & Si & Cu & Mg & Mn & Zn & & & \makecell{$R_m$ \\ /MPa} & \makecell{$A$ \\ /\%} & \makecell{硬度 \\ /HBW} & \\
\hline
\multirow{2}{*}{\makecell{铝硅 \\ 合金}} & ZAlSi7Mg & ZL101 & 6.5\textasciitilde7.5 & & 0.25\textasciitilde0.45 & & & J、JB & T5 & 205 & 2 & 60 & \makecell{形状复杂的 \\ 零件，如飞机、 \\ 仪器零件、抽水 \\ 机壳体等} \\
\cline{2-14}
 & ZAlSi9Mg & ZL104 & 8.0\textasciitilde10.5 & & 0.17\textasciitilde0.3 & & & J & T1 & 200 & 1.5 & 70 & \makecell{形状复杂、工 \\ 作温度在200 $^\circ$C \\ 以下的零件，如 \\ 电动机壳体、气 \\ 缸体等} \\
\hline
\makecell{铝铜 \\ 合金} & ZAlSiCu4 & ZL203 & & 4.0\textasciitilde5.0 & & & & J & T5 & 225 & 3 & 70 & \makecell{中等载荷和 \\ 形状复杂的零 \\ 件；工作温度不 \\ 超过200 $^\circ$C且要 \\ 求切削性能好 \\ 的小零件} \\
\hline
\makecell{铝镁 \\ 合金} & ZAlMg5Si & ZL303 & 0.8\textasciitilde1.3 & & 4.5\textasciitilde5.5 & 0.1\textasciitilde0.4 & & \makecell{S、 \\ J、 \\ R、 \\ K} & T5 & 143 & 1 & 55 & \makecell{腐蚀介质作用 \\ 下的中等载荷零 \\ 件；在严寒以及工 \\ 作温度不超过 \\ 200 $^\circ$C的零件，如 \\ 海轮配件及各种 \\ 壳体等} \\
\hline
\makecell{铝锌 \\ 合金} & ZAlZn11Si7 & ZL401 & 6.0\textasciitilde8.0 & & 0.1\textasciitilde0.3 & & 9.0\textasciitilde13.0 & J & T1 & 245 & 1.5 & 90 & \makecell{结构形状复杂 \\ 的零件；压力铸造 \\ 零件；工作温度不 \\ 超过200 $^\circ$C的零 \\ 件} \\
\hline
\end{tabular}
\begin{tablenotes}
	\item 注：1. 铸造方法代号：B表示变质处理；J表示金属型铸造；K表示壳型铸造；R表示熔模铸造；S表示砂型铸造；
	\item 2. 热处理代号：T1表示人工时效；T5表示淬火+部分时效。
\end{tablenotes}
\end{table}

\subsection{铜及铜合金}

纯铜的密度为8.94 g/cm$^3$，熔点为1083 $^\circ$C，面心立方晶格，无同素异构转变，无磁性，具有良好的导电性、导热性（仅次于银），有较高的塑性和耐腐蚀性，但强度、硬度低，不能通过热处理强化。故工业上常通过添加合金元素来改善其性能。纯铜中的杂质主要有Pb、Bi、O、S、P等，按纯度可分为T1、T1.5、T2、T3四种。“T”为铜的汉语拼音的第一个字母，其后数字越大，纯度越低。纯铜广泛用于制造电线、电缆、电刷、铜管、铜棒以及作为配制铜合金的原料。铜合金按化学成分可分为黄铜、青铜和白铜三类。

1. 黄铜

以Zn为唯一或主要合金元素的铜合金称为黄铜。按其化学成分的不同，黄铜分为普通黄铜和特殊黄铜两类。锌的质量分数为30\%\textasciitilde32\%时，黄铜组织为面心立方晶格的$\alpha$固溶体；锌的质量分数大于32\%（不超过45\%\textasciitilde47\%）时，黄铜为$\alpha+\beta$两相黄铜。普通黄铜的牌号以“H”+数字表示。“H”为“黄”字汉语拼音的第一个字母，数字表示铜的质量分数，如H80即80\%铜和20\%锌的普通黄铜。

特殊黄铜又称为复杂黄铜，是在铜锌合金中再加入其他合金元素，除主加元素Zn外，常加入的其他合金元素有Pb、Al、Mn、Sn、Fe、Ni、Si等，又分别称为铅黄铜、铝黄铜、锰黄铜等。特殊黄铜的牌号用“H”+主加元素的化学符号+铜的质量分数+主加元素的质量分数表示。如HPb59-1表示铜的质量分数为59\%，铅的质量分数为1\%，其余为锌。铸造用黄铜在牌号“H”前加“Z”（“铸”字汉语拼音的第一个字母),如 ZHA167-2.5 表示铜的质量分数为 67\%,铝的质量分数为 2.5\% 的铸造铝黄铜。

常用黄铜的牌号、化学成分、力学性能及用途见表 5-3。

\begin{table}[H]
\centering
\caption{常用黄铜的牌号、化学成分、力学性能及用途}
\label{tab:brass_properties}
\footnotesize\setlength{\tabcolsep}{2pt}
\begin{tabular}{|c|c|c|c|c|c|c|c|c|c|c|p{3cm}|}
\hline
\multirow{2}{*}{类别} & \multirow{2}{*}{牌号} & \multicolumn{5}{c|}{化学成分 $w/\%$} & \multirow{2}{*}{状态} & \multicolumn{3}{c|}{力学性能} & \multirow{2}{3cm}{用途} \\
\cline{3-7}\cline{9-11}
 & & Cu & Pb & Si & Al & 其他 & & $R_m$/MPa & $A$/\% & 硬度/HBW & \\
\hline
\multirow{3}{*}{普通黄铜} & H90 & 89.0~91.0 & 0.05 & & & Fe:0.05 & M & $\geqslant 220$ & $\geqslant 42$ & 40~70 & 双金属片、供水和排水管、装饰品等 \\
\cline{2-12}
 & H68 & 67.0~70.0 & 0.03 & & & Fe:0.10 & M & $\geqslant 280$ & $\geqslant 43$ & 50~80 & 散热器外壳、弹壳、导管、冷凝器管等 \\
\cline{2-12}
 & H62 & 60.5~63.5 & 0.08 & & & Fe:0.15 & M & $\geqslant 300$ & $\geqslant 43$ & 55~85 & 铆钉、垫圈、弹簧、螺钉等 \\
\hline
\multirow{5}{*}{特殊黄铜} & 铅黄铜 & HPb59-1 & 57.0~60.0 & 0.8~1.9 & & Fe:0.50 & M & $\geqslant 340$ & $\geqslant 35$ & 70~100 & 销、螺钉等冲压或加工件 \\
\cline{2-12}
 & 铝黄铜 & HAl59-3-2 & 57.0~60.0 & 0.10 & & 2.5~5.5 & Fe:0.50 & & & & 在常温下工作的高强度、耐腐蚀零件 \\
\cline{2-12}
 & 锰黄铜 & HMn58-2 & 57.0~60.0 & 0.1 & & & Fe:1.0 & $\geqslant 395$ & $\geqslant 29$ & & 船舶和弱电流用零件 \\
\cline{2-12}
 & 铸造硅黄铜 & HSi80-3 & 79.0~81.0 & 0.1 & 2.5~4.0 & & Fe:0.60 & $\geqslant 295$ & $\geqslant 28$ & & 耐腐蚀介质条件下工作的零件 \\
\cline{2-12}
 & 铸造铝黄铜 & HAl66-6-3-2 & 64.0~68.0 & 0.5 & & 6.0~7.0 & Fe: 2.0~4.0 & $\geqslant 735$ & $\geqslant 8$ & & 压紧螺母、重型螺杆、轴承、衬套等 \\
\hline
\end{tabular}
\begin{tablenotes}
	\item 注:合金状态 M 表示软状态。
\end{tablenotes}
\end{table}

2. 青铜

除了以 Zn 或 Ni 为主加元素的铜合金外,其余铜合金统称为青铜。青铜的牌号以“Q”(“青”字汉语拼音的第一个字母)为首,其后标注主要的合金元素及其质量分数。铸造用青铜,在牌号“Q”前冠以“Z”字,例如 ZQSn10 表示锡的质量分数为 10\%的铸造锡青铜。青铜分为普通青铜(锡青铜)与特殊青铜(如铝青铜、铅青铜、铍青铜、硅青铜等无锡青铜)两类。

锡青铜以 Sn 为主加元素:当 Sn 的质量分数小于 5\%时,合金的退火组织为 $\alpha$ 固溶体;当 Sn 的质量分数超过 6\%时,合金组织出现硬而脆的 $\delta$ 相;当 Sn 的质量分数超过 10\%时,只能用于铸造。工业用锡青铜的 Sn 质量分数为 3\%~14\%,另外还含有少量的 Zn、Pb、P、Ni 等元素。锡青铜具有良好的减摩性、抗磁性及低温韧性,但其耐酸腐蚀能力较差。

特殊青铜是不含 Sn 的青铜,根据主加元素的不同,分为铝青铜、铍青铜、硅青铜等。

铝青铜中 Al 的质量分数为 5\%~10\%,化学稳定性高,耐蚀、耐磨,且工艺性好,强度和塑性较高,主要用于在海水或高温下工作的高强度耐磨零件;铍青铜中 Be 的质量分数为 1.7\%~2.5\%,可进行固溶和时效强化处理,强度高,弹性极限、疲劳强度高,耐磨性、耐蚀性、导电性、导热性好,同时还具有抗磁、受冲击时不产生火花等特殊性能,主要用于精密仪器中的弹性元件及电动机中的防爆部件;硅青铜以 Si 为主加元素,具有比锡青铜高的力学性能,铸造性能、冷加工和压力加工性能良好。硅青铜中加入 Ni 可显著提高强度和耐磨性,主要用于长距离架空电话线和输电线等。常用青铜的牌号、化学成分、力学性能及用途见表 5-4。

\begin{table}[H]
\centering
\caption{常用青铜的牌号、化学成分、力学性能及用途}
\label{tab:bronze_properties}
\footnotesize\setlength{\tabcolsep}{1pt}
\begin{tabular}{|c|c|c|c|c|c|m{1.5cm}|c|c|c|c|m{1.6cm}|}
\hline
\multicolumn{2}{|c|}{\multirow{2}{*}{类别}} & \multirow{2}{*}{牌号} & \multicolumn{4}{c|}{化学成分 $w/\%$} & \multirow{2}{*}{状态} & \multicolumn{3}{c|}{力学性能 (不小于)} & \multirow{2}{1.2cm}{用途} \\
\cline{4-7}\cline{9-11}
 & & & Sn & Pb & Al & 其他 & & $R_m$/MPa & $A$/\% & 硬度/HBW & \\
\hline
\multirow{4}{*}{\makecell{普通\\锡青\\铜}} & \multirow{2}{*}{\makecell{铸造\\锡青\\铜}} & ZCuSn10P1 & 9.0$\sim$11.5 & & & P:0.8$\sim$1.1 & S、R & 220 & 3 & 80 & 轴承、齿轮、轴套等 \\
\cline{3-12}
 & & ZCuPb20Sn5 & 4.0$\sim$6.0 & 18.0$\sim$23.0 & & & S & 150 & 5 & 45 & 中速、中载轴承、轴套等 \\
\cline{2-12}
 & \multirow{2}{*}{锡青铜} & QSn4$\sim$3 & 3.5$\sim$4.5 & 0.02 & 0.002 & Zn:2.7$\sim$3.3 & M & 290 & 40 & & 弹性元件,耐磨、抗磁元件 \\
\cline{3-12}
 & & QSn6.5$\sim$0.1 & 6.0$\sim$7.0 & 0.02 & 0.002 & Zn:0.3 & M & 315 & 40 & & 接触片、弹簧、耐磨件等 \\
\hline
\multirow{3}{*}{特殊青铜} & 铝青铜 & QAl9-4 & 0.1 & 0.01 & 8.0$\sim$10.0 & Zn:1.0 Mn:0.5 & Y & 635 & & & 耐磨零件及在蒸汽、海水中工作的高强度零件 \\
\cline{2-12}
 & 铍青铜 & QBe2 & 0.005 & 0.15 & 0.15 & Si:0.15 & & & & & 重要弹簧、弹性元件、轴承等 \\
\cline{2-12}
 & 硅青铜 & QSi3-1 & 0.25 & 0.03 & & Si:2.7$\sim$ 3.5 Mn:1.0$\sim$1.5 Zn:0.5 & M & 370 & 45 & & 弹簧以及在腐蚀介质中工作的零件 \\
\hline
\end{tabular}\end{table}

\subsection{钛及钛合金}

钛在地壳中的蕴藏量仅次于铝、铁、镁,居金属元素的第四位。纯钛为银白色金属,其熔点为 1 667 $^\circ$C,密度为 4.5 g/cm$^3$,室温下为密排六方结构(晶格)。钛具有同素异构转变,低于 882.5 $^\circ$C 为密排六方晶格,称为 $\alpha$-Ti,高于 882.5 $^\circ$C 为体心立方晶格,称为 $\beta$-Ti。

钛合金的主要优点是:强度高,其抗拉强度可达 1 500 MPa,可与超高强度钢媲美,而较低的密度使其成为比强度最高的常用工程材料之一;热强度高,可在 500 $^\circ$C 以上的温度下工作,其工作温度可望提高到 800 $^\circ$C;断裂韧性较高,优于铝合金和一些结构钢;耐蚀性高,钛在大气、淡水、海水、硝酸、稀硫酸等腐蚀介质中的耐腐蚀性优于不锈钢。钛及钛合金作为一种高耐蚀性材料,已在航空、化工、造船及医疗等行业得到广泛应用,特别是近年来在激光立体成形(3D 打印)大型钛合金零件中得到应用。

钛合金的主要缺点是:工艺性差,导热系数小、摩擦因数大,故切削性差;弹性模量小,变形时回弹大,冷变形困难;硬度低,不耐磨,不能作抗磨结构件。另外,钛合金成本高,故其应用受到限制。钛合金按其退火组织可分为 $\alpha$ 型钛合金(TA)、$\beta$ 型钛合金(TB)、$\alpha+\beta$ 型钛合金(TC)。

1. $\alpha$ 型钛合金

这类合金的退火组织为单相 $\alpha$ 固溶体,压力加工性较差,多采用热压加工成形。这类合金不能进行热处理强化,只能进行退火处理,室温强度不高。绝大多数 $\alpha$ 单相钛合金中均含 5\% Al 与 2.5\% Sn,其作用是固溶强化。

2. $\beta$ 型钛合金

钛合金加入 V 或 Mo 后其组织在室温下全部为 $\beta$ 相,通过淬火可得到介稳定 $\beta$ 相的钛合金,这类合金可进行热处理强化(淬火、时效),故室温强度较高。其主要用途为高强度紧固件、杆类件及其他航空配件。

3. $\alpha+\beta$ 型钛合金

加入适当的稳定 $\beta$ 相合金元素,可形成室温下的 $\alpha+\beta$ 组织,兼有 $\alpha$ 型及 $\beta$ 型两类钛合金的优点。高度合金化的 $\alpha+\beta$ 型钛合金可通过热处理获得。在航空、航天领域,钛合金应用最广泛的是多用途的 $\alpha+\beta$ 型 Ti-6Al-4V 合金。

常用钛合金的牌号、化学成分、力学性能及用途见表 5-5 中。

\begin{table}[H]
	\centering
	\caption{常用钛合金的牌号、化学成分、力学性能及用途}
	\label{tab:titanium_properties}
	%\footnotesize\setlength{\tabcolsep}{2pt}
	\begin{tabular}{|p{1cm}|p{0.8cm}|p{2cm}|p{1cm}|p{1cm}|p{1cm}|p{1cm}|p{1cm}|p{1cm}|m{3.5cm}|}
		\hline
		\multirow{2}{*}{类别} & \multirow{2}{*}{代号} & \multirow{2}{*}{化学成分} & \multirow{2}{*}{\makecell{热处理\\状态}} & \multicolumn{2}{c|}{室温力学性能} & \multicolumn{3}{c|}{高温力学性能} & \multirow{2}{*}{用途} \\
		\cline{5-9}
		& & & & $R_m$ MPa & $A$\newline \% & 试验温度 $^\circ$C & $R_m$ MPa & $R_{100h}$ MPa & \\
		\hline
		\multirow{3}{*}{纯钛} & TA1 & Ti(杂质极微) & 退火 & 370 & 20 & & & & \multirow{3}{3.5cm}{在 350$^\circ$C 以下工作、强度要求不高的零件} \\
		\cline{2-9}
		& TA2 & Ti(杂质微) & 退火 & 440 & 18 & & & & \\
		\cline{2-9}
		& TA3 & Ti(杂质微) & 退火 & 540 & 15 & & & & \\
		\hline
		\multirow{2}{1cm}{$\alpha$ 型钛合金} & TA5 & Ti-4Al-0.005B & 退火 & 685 & 13 & & & & \multirow{2}{3.5cm}{在 500 $^\circ$C 以下工作的零件;导弹燃料罐、飞机的涡轮机匣等} \\
		\cline{2-9}
		& TA6 & Ti-5Al & 退火 & 685 & 10 & 350 & 420 & 390 & \\
		\hline
		\multirow{2}{1cm}{$\beta$ 型钛合金} & \multirow{2}{*}{TB2} & \multirow{2}{2cm}{Ti-5Mo-5V-8Cr-3Al} & 淬火 & $\leqslant 980$ & 18 & & & & \multirow{2}{3.5cm}{在 350 $^\circ$C 以下工作的零件,压气机叶片、轴、轮盘等重载荷旋转件;飞机构件等} \\
		\cline{4-9}
		& & & 淬火+时效 & 1 370 & 7 & & & & \\
		\hline
		\multirow{4}{1cm}{$\alpha+\beta$ 型钛合金} & TC1 & Ti-2Al-1.5Mn & 退火 & 585 & 15 & 350 & 345 & 325 & \multirow{4}{3.5cm}{在 400 $^\circ$C 以下工作的零件,有一定高温强度的发动机零件,低温用部件} \\
		\cline{2-9}
		& TC2 & Ti-4Al-1.5Mn & 退火 & 685 & 12 & 350 & 420 & 390 & \\
		\cline{2-9}
		& TC3 & Ti-5Al-4V & 退火 & 800 & 10 & & & & \\
		\cline{2-9}
		& TC4 & Ti-6Al-4V & 退火 & 895 & 10 & 400 & 600 & 560 & \\
		\hline
\end{tabular}
\begin{tablenotes}
	\item 注:参考国标 GB/T 2965—2023。
\end{tablenotes}
\end{table}

\section{高分子材料}

高分子材料包括塑料、合成纤维、橡胶和胶黏剂等。本节主要介绍工程上常用的高分子材料。

1. 工程塑料

塑料是一种以有机合成树脂为主要成分的高分子材料,它通常可在加热、加压条件下塑造成形,故称为塑料。塑料的种类很多,按照热性能分为热塑性塑料和热固性塑料两大类。前者的特点是加热时软化,可塑造成形,冷却后变硬,此过程可反复进行。后者的特点是初加热时软化,可塑造成形,但固化之后再加热既不再软化,也不溶于溶剂。常用塑料的主要特性和用途见表 5-6。

2. 橡胶

橡胶是优良的高弹性材料,同时还具有良好的耐磨、隔声和绝缘等性能,是重要的工业原料之一,广泛应用于各个领域。表 5-7 列出了常用橡胶的性能和用途。

\begin{table}[H]
\centering
\caption{常用塑料的主要特性和用途}
\label{tab:plastic_properties}
\small
\begin{tabular}{|c|c|c|c|p{3.5cm}|p{3.5cm}|}
\hline
塑料名称 & 代号 & 使用温度/$^\circ$C & 抗拉强度/MPa & 主要特性 & 用途 \\
\hline
聚乙烯 & PE & -70~100 & 8~36 & 优良的耐磨性,尤其是高频绝缘性 & 一般机械构件、电缆包覆、耐蚀涂层等 \\
\hline
聚丙烯 & PP & -35~121 & 40~49 & 密度小、力学性能较高,耐热性好,耐蚀性优良,高频绝缘性良好,不受湿度影响,低温易老化 & 一般机器零件,高频绝缘电缆、电缆包覆等 \\
\hline
聚氯乙烯 & PVC & -15~55 & 30~60 & 优良的耐蚀性,可改性。硬聚氯乙烯强度高;软聚氯乙烯强度低,伸长率大,易老化;泡沫聚氯乙烯质轻 & 硬聚氯乙烯用作耐蚀件、化工机械零件,软聚氯乙烯用作薄膜、密封件衬底等 \\
\hline
聚苯乙烯 & PS & -30~75 & $\geqslant 60$ & 优良的电绝缘性,尤其是高频绝缘性,无色透明,着色性好,质脆,不耐苯、汽油等有机溶剂,可改性 & 绝缘件、透明件、装饰件、泡沫聚苯乙烯包装和铸型等 \\
\hline
有机玻璃 & PMMA & -60~100 & 42~50 & 透光性、着色性好,表面硬度不高,易擦伤,可改性 & 透明件、装饰件等 \\
\hline
聚酰胺(尼龙) & PA & $<100$ & 45~90 & 坚韧,耐磨,耐疲劳,耐油,耐水,抗霉菌,无毒,吸水性大。尼龙:弹性好,冲击强度高;芳香尼龙:耐热 & 一般机械零件,减摩、耐磨传动件;铸型尼龙可制作大型减摩、耐磨传动件等 \\
\hline
ABS 塑料 & ABS & -60~100 & 21~63 & 综合性能良好,较高的强度和冲击韧度,耐热,表面硬度高,尺寸稳定,耐化学性及电性能良好,易于成形和机械加工 & 一般机械零件及传动件,壳体、管道、储槽、仪表盘、轿车车身等 \\
\hline
聚甲醛 & POM & -40~100 & 60~75 & 良好的综合力学性能,强度、刚性、冲击韧性、疲劳强度、蠕变强度等均较高;减摩、耐磨性良好,吸水性好,尺寸稳定性好 & 减摩、耐磨传动件,绝缘、耐蚀件及化工容器等 \\
\hline
聚四氟乙烯 & F-4 & -180~260 & 21~28 & 耐所有化学药品(包括王水)的腐蚀,摩擦因数低,不黏,不吸水,流动性好,不能注射成形 & 耐腐蚀件,减摩、耐磨件,密封件,绝缘件等 \\
\hline
聚砜 & PSE & -65~150 & $\sim 70$ & 强度高,冲击强度大,在水、湿空气或高温下具有良好绝缘性,不耐芳香烃及卤代烃 & 高强度耐热件、绝缘件及传动件、高频印制电路板等 \\
\hline
酚醛塑料 & PF & $<140$ & 21~56 & 优良的耐热、绝缘、化学稳定及尺寸稳定性,抗蠕变性能优良,因填料不同性能有差异 & 一般机械零件,绝缘件、耐蚀件、水润滑轴承等 \\
\hline
环氧塑料 & EP & -80~155 & 56~70 & 强度较高,电绝缘性优良,化学稳定性好,耐有机溶剂性好,因填料不同性能有差异 & 塑料膜,电气、电子元件及线圈的灌封与固定,修复机件等 \\
\hline
\end{tabular}\end{table}

\begin{table}[H]
\centering
\caption{常用橡胶的性能及用途}
\label{tab:rubber_properties}
\footnotesize\setlength{\tabcolsep}{2pt}
\begin{tabular}{|p{1.8cm}|p{1.5cm}|p{1.5cm}|p{1.2cm}|p{1.8cm}|p{1.5cm}|p{1.5cm}|p{1.8cm}|p{1.8cm}|p{1.2cm}|}
\hline
\multirow{2}{1.8cm}{名称} & \multicolumn{5}{c|}{通用橡胶} & \multicolumn{4}{c|}{特种橡胶} \\
\cline{2-10}
 & 天然橡胶 & 丁苯橡胶 & 氯丁橡胶 & 顺丁橡胶 & 乙丙橡胶 & 丁腈橡胶 & 硅橡胶 & 氟橡胶 & 聚硫橡胶 \\
\hline
抗拉强度/MPa & 25~35 & 15~20 & 25~27 & 17~21 & 10~25 & 15~30 & 4~10 & 20~22 & 20~35 \\
\hline
伸长率/\% & 650~900 & 500~800 & 800~1000 & 650~800 & 400~800 & 300~800 & 50~500 & 100~500 & 300~800 \\
\hline
使用温度/$^\circ$C & -50~120 & -50~140 & -35~130 & -120~170 & -65~170 & -35~175 & -90~300 & -50~300 & -55~115 \\
\hline
耐磨性 & 中 & 好 & 中 & 中 & 中 & 中 & 差 & 中 & 好 \\
\hline
耐油性 & 差 & 差 & 好 & 中 & & 好 & 差 & 中 & 好 \\
\hline
耐碱性 & & & 好 & 好 & & & & 好 & 差 \\
\hline
耐老化 & & & 好 & 好 & 好 & & & 好 & \\
\hline
绝缘性 & 好 & & & 好 & 好 & 差 & 好 & & \\
\hline
成本 & & & 高 & & & & 高 & 高 & \\
\hline
特殊性能 & 高强度、绝缘、防振 & 耐碱、耐燃 & 耐酸、耐碱、耐燃 & 耐酸、耐碱、气密、绝缘、防振 & 耐水、绝缘、耐臭氧 & 耐油、耐水、气密 & 耐热、绝缘、无毒 & 耐油、耐碱、耐热 & 高强度、耐磨 \\
\hline
用途 & 通用橡胶制品、轮胎 & 轮胎、传送带、密封垫圈 & 传送带、制动器、汽车门窗嵌条 & 内胎、化工容器衬里、减振制品、嵌缝胶条、密封设备中的耐热输送带 & 汽车配件、散热管、电绝缘件 & 耐油垫圈、油管、油封、印刷胶辊 & 航空工业和电子设备中的橡胶制品、加热管道用的隔热材料 & 主要用于航空、火箭、导弹等方面的高级密封件、高真空橡胶件 & 实心胎胶辊、耐磨件 \\
\hline
\end{tabular}\end{table}

3. 胶黏剂

胶黏剂是以富含黏性的物质为基料,再加入各种添加剂后组成的。借助胶黏剂实现的连接称为胶接。胶接的主要优点是劳动强度低、工艺操作简单、生产效率高,缺点是胶接强度低。常用胶黏剂的性能及用途见表 5-8。

\begin{table}[H]
\centering
\caption{常用胶黏剂的性能及用途}
\label{tab:adhesive_properties}
\begin{tabular}{|p{3cm}|p{4.5cm}|p{6.5cm}|}
\hline
胶黏剂名称 & 基本性能 & 用途 \\
\hline
环氧树脂胶黏剂 & 黏接力强;工艺性能好,能配制成多种黏度的胶种,固化时间、温度使用期间可调节,且稳定好;胶层机械强度高,绝缘性能好,体积收缩小 & 用于金属、玻璃、瓷器、橡胶、木材、塑料本身及相互之间的胶接,有“万能胶”之称 \\
\hline
酚醛-丁腈胶黏剂 & 黏接力强,尤其是有较高的剥离强度;使用温度范围广,一般可在 -50~120 $^\circ$C 使用;耐油性好 & 用于金属、瓷器、玻璃、塑料、层压板等胶接。常用于刹车片、飞机油箱的整体密封、各种机械修理 \\
\hline
聚醋酸乙烯酯乳液胶(白胶) & 对木材有良好的黏接力 & 用于木材、皮革、纸张、纤维等的胶接 \\
\hline
聚氨酯胶黏剂 & 耐寒性好,几乎超过其他所有胶种;弹性好,强度也高;耐疲劳性好 & 用于金属、橡胶、织物、塑料等胶接 \\
\hline
丙烯酸酯胶黏剂 & 室温下能快速固化,且固化速度可调;黏接强度(包括剥离、剪切、冲击等综合强度)高;黏接范围广,特别适用异种材料的黏接;使用方便 & 用于钢、铝、ABS 塑料、聚氯乙烯、有机玻璃、玻璃、木材、混凝土、橡胶等的胶接 \\
\hline
\end{tabular}
\end{table}

\section{陶瓷材料}

\subsection{陶瓷材料的性能及分类}

所谓陶瓷是指以天然硅酸盐(黏土、石英、长石等)或人工合成化合物(氮化物、氧化物、碳化物等)为原料,经过制粉、配料、成形、高温烧结而成的无机非金属材料。工业陶瓷可分为普通陶瓷和特种陶瓷两大类。普通陶瓷主要用于工艺品、餐具等日常生活用品。特种陶瓷具有特殊的力学、物理或化学性能,主要用于机械、电子、宇航、医学工程等尖端科学技术领域。

陶瓷材料通常具有强度高、硬度高、化学和热稳定性好,且耐高温、耐腐蚀等特点。由于陶瓷的键合特点,陶瓷材料还具有绝缘、绝热的性能。但陶瓷的最大缺点是塑性很差,多数陶瓷材料在常温下没有塑性,因此陶瓷材料应用面临的主要问题是“增韧”,工程实践中常见的增韧途径包括相变增韧、颗粒增韧、纤维增韧等。

\subsection{常用特种陶瓷}

1. 氧化物陶瓷

常用的氧化物陶瓷包括 Al$_2$O$_3$、BeO、ZrO$_2$、MgO、CaO、ThO$_2$、UO$_2$ 等,其熔点大多为 2 000 $^\circ$C 以上,烧结温度为 1 800 $^\circ$C 左右。

(1) 氧化铝(Al$_2$O$_3$)陶瓷

氧化铝是高熔点(熔点达 2 050 $^\circ$C)氧化物中被研究和应用最成熟、最多的一种。其原料来源丰富,约占地壳总质量的 25\%,价格低廉。Al$_2$O$_3$ 中的 O$^{2-}$ 为密排六方结构(晶格),Al$^{3+}$ 占据间隙位置。如果其含有不同的杂质原子,则为红宝石或蓝宝石。氧化铝陶瓷按氧化铝的质量分数可分为 75 瓷、95 瓷和 99 瓷。由于氧化铝熔点高,热强度大,抗氧化,故耐热性好。氧化铝常温硬度仅次于立方氮化硼,而且具有高的电阻率和低的热导率,所以除用作高温结构材料、绝缘和绝热材料以及模具材料外,重要的应用就是制作切削刀具。

氧化铝陶瓷刀具的主要成分有两种类型:一类是在 Al$_2$O$_3$ 中加入 1\% 左右的 MgO;另一类是在 Al$_2$O$_3$ 中加入较多(约 30\%)的金属碳化物,如 TiC、WC 和 Mo$_2$C 等。添加物可提高刀具的韧性和抗弯强度。氧化铝陶瓷刀具切削温度可达 1 000 $^\circ$C 以上,适用于切削难切削材料,如淬火钢、冷硬铸铁等。

(2) 氧化铍(BeO)陶瓷

氧化铍陶瓷除了具备普通陶瓷的特性外,其最大的特点是导热性极好,因而具有很高的热稳定性。虽然其强度不高,但抗热冲击性较高。氧化铍陶瓷具有消散高能辐射的能力强、热中子阻尼系数大等特性,所以经常用于制造坩埚,还可用作真空陶瓷和原子反应堆陶瓷等。另外,气体激光管、晶体管散热片和集成电路的基片和外壳等也多用该种陶瓷。

(3) 氧化锆(ZrO$_2$)陶瓷

氧化锆陶瓷的熔点为 2 700 $^\circ$C 以上,能耐 2 300 $^\circ$C 的高温,其推荐使用温度为 2 000 ~ 2 500 $^\circ$C。由于它还能抗熔融金属的侵蚀,所以多用作铂、锆等金属的冶炼坩埚和 1 800 $^\circ$C 以上的发热体及炉子、核反应堆绝热材料等。氧化锆作添加剂可大大提高其他陶瓷材料的强度和韧性。氧化锆增韧氧化铝陶瓷材料的强度达 1 200 MPa、断裂韧度为 15 MN/m$^{3/2}$,分别比原氧化铝提高了 3 倍和近 3 倍。氧化锆陶瓷可替代金属制造模具、拉丝模、泵叶轮等,还可制造汽车零件,如凸轮、推杆、连杆等。氧化锆陶瓷制成的剪刀既不生锈,也不导电。

(4) 氧化镁/氧化钙(MgO/CaO)陶瓷

氧化镁/氧化钙陶瓷通常是通过加热白云石(镁或钙的碳酸盐)矿石,除去 CO$_2$ 而制成的。其特点是能抗各种金属碱性渣的作用,因而常用作炉衬的耐火砖。但这种陶瓷的缺点是热稳定性差,MgO 在高温下易挥发,而 CaO 易在空气中受潮变质。

(5) 氧化钍/氧化铀(ThO$_2$/UO$_2$)陶瓷

这是具有放射性的一类陶瓷,具有极高的熔点和密度,多用于制造熔化铑、铂、银和其他金属的坩埚及核反应堆中的放热元件等,氧化钍陶瓷还可用于制造电炉构件。

2. 碳化物陶瓷

最常见的碳化物陶瓷包括 SiC、BC、GeC、Mo$_2$C、NbC、TiC、WC、TaC、VC、ZrC 等。该类陶瓷的突出特点是具有很高的熔点、硬度(接近于金刚石)和耐磨性(特别是在侵蚀性介质中),缺点是

耐高温能力差(工作温度为$900\sim1~000~^\circ\mathrm{C}$)。

(1)碳化硅(SiC)陶瓷

碳化硅陶瓷的密度为$3.2~\mathrm{g/cm^3}$,抗弯强度和抗压强度分别为$200\sim250~\mathrm{MPa}$和$1~000\sim1~500~\mathrm{MPa}$,硬度为莫氏9.2。该材料热导率很高,而热膨胀系数很小,但在$900\sim1~300~^\circ\mathrm{C}$时会缓慢氧化。碳化硅陶瓷是应用最为广泛的碳化物陶瓷。

碳化硅陶瓷通常用作加热元件、石墨表面保护层以及砂轮和磨料等。将用有机黏接剂黏接的碳化硅陶瓷在$1~700~^\circ\mathrm{C}$热压成形,可得到高强度、高密度、高耐磨性和高抗化学侵蚀的耐火材料。

(2)碳化硼(BC)陶瓷

碳化硼陶瓷的硬度极高,抗磨粒磨损能力很强;熔点高达$2~450~^\circ\mathrm{C}$左右,但在高温下会快速氧化,并且与热或熔融钢铁合金发生反应,因此其使用温度限定在$980~^\circ\mathrm{C}$以下。其主要用途是作磨料,有时用作超硬质工具材料。

3.氮化物陶瓷

氮化硅($\mathrm{Si_3N_4}$)和氮化硼(BN)是最常见的氮化物陶瓷。

(1)氮化硅($\mathrm{Si_3N_4}$)陶瓷

氮化硅陶瓷是键能高而稳定的共价键晶体,晶体内无正离子,也无自由电子。氧化铝陶瓷硬度高,摩擦因数低,有自润滑作用,是优良的耐磨材料,在$1~400~^\circ\mathrm{C}$以下,热强度和化学稳定性高,热膨胀系数小,而且抗热冲击,也是优良的高温结构材料。

氧化硅陶瓷的应用与制造方法与其的晶体类型有关。用反应烧结法得到$\alpha\text{-}\mathrm{Si_3N_4}$,用热压烧结法得到$\beta\text{-}\mathrm{Si_3N_4}$。用热压烧结法制备的陶瓷,其密度和抗弯强度要比用反应烧结法制备的陶瓷高得多,这是由于热压烧结陶瓷气相数量少。但热压烧结法只适用于形状简单的零件。

近年来,在$\mathrm{Si_3N_4}$中添加一定比例(如50\%)的$\mathrm{Al_2O_3}$,可实现常压下的烧结,得到和用热压烧结法性能近似的陶瓷,而且抗氧化性更高。这种添加$\mathrm{Al_2O_3}$的氧化硅陶瓷又称塞伦陶瓷,可制造柴油机气缸、活塞及燃气轮机转子叶片等零件。

(2)氮化硼(BN)

氮化硼有六方氮化硼和立方氮化硼两种。六方氮化硼结构与石墨相似,故又称为“白石墨”,具有良好的耐热性,导热系数与不锈钢相当,热稳定性好;在$2~000~^\circ\mathrm{C}$仍然是绝缘体,是理想的高温绝缘材料和散热材料;与石墨相似,六方氮化硼硬度低,有自润滑性。氮化硼陶瓷具有高温绝缘性,故常用来制造热电偶套管和半导体散热绝缘零件;它还具有优良的热稳定性和化学稳定性,因此也被用作熔炼半导体GaAs、GaP的单晶坩埚及一般冶金的高温容器和管道。六方氮化硼硬度低,可进行车、铣、刨等切削加工,制品精度可达$1/1~000~\mathrm{mm}$,又因为具有自润滑性,故可用来制造高温轴承和玻璃制品成形模具。

立方氮化硼晶体的硬度与金刚石相近,是优良的耐磨材料,常用作磨料和金属切削刀具。

4.硼化物陶瓷

常见的硼化物陶瓷包括硼化铬、硼化钼、硼化钛、硼化钨和硼化锆等。其特点是高硬度,同时具有较好的耐化学侵蚀能力。其熔点为$1~800\sim2~500~^\circ\mathrm{C}$。比起碳化物陶瓷,硼化物陶瓷具有较高的抗高温氧化性能,可在高达$1~400~^\circ\mathrm{C}$的温度下工作。硼化物陶瓷主要用于制作高温轴承、内燃机喷嘴、各种高温器件、处理熔融非金属的器件等。此外,还可用作电触点材料。

\section{复合材料}

\subsection{复合材料的分类及性能}

所谓复合材料,是指两种或两种以上的物理、化学性质不同的物质经一定方法得到的一种新的多相固体材料。复合材料的主要优点是能根据人们的要求来改善材料的使用性能。

按基体材料不同,复合材料可分为金属基复合材料、树脂基复合材料和陶瓷基复合材料三种。按其结构特点,复合材料又可分为颗粒增强复合材料、纤维增强复合材料、层合复合材料和骨架复合材料等,其中纤维增强复合材料应用最广。

复合材料和其他材料相比具有比强度和比模量高,抗疲劳性能和抗断裂性能好,减摩、耐磨、减振性能优良等特点。金属基复合材料具有高韧性和抗热冲击性能。另外,复合材料还具有高耐辐射性能、耐蠕变性能以及特殊的光、电、磁等性能。

\subsection{常用复合材料}

1.颗粒增强复合材料

颗粒增强复合材料是由一种或多种材料的颗粒均匀分散在基体内所形成的材料。由陶瓷颗粒与金属结合的颗粒增强复合材料称为金属陶瓷。材料中的陶瓷为氧化物、碳化物、硼化物和氮化物,起强化作用。目前,应用较多的是碳化物基和氧化物基金属陶瓷。

(1)碳化物基金属陶瓷

碳化物基金属陶瓷又称为硬质合金,是由难熔的碳化物粉末和Co等金属黏结剂混合加压成形,再经烧结而成的材料。其硬度可达93 HRA,热硬性可达$1~000~^\circ\mathrm{C}$,耐磨性好,与高速钢相比,切削速度可提高4$\sim$7倍,刀具寿命提高5$\sim$8倍,可切削淬硬钢和奥氏体钢。硬质合金主要有钨钴类硬质合金(P类)、钨钴钛类硬质合金(M类)、通用硬质合金(K类)和钢结硬质合金等几类。其中,钢结硬质合金是一种新型工模具材料,是以WC、TiC、VC为硬化相,以高速钢或铬钢粉作为黏结剂,用烧结法制成的合金。

(2)氧化物基金属陶瓷

氧化铝基金属陶瓷是应用最多的氧化物基金属陶瓷,常用Cr作黏结剂,具有高强度、高硬度、高耐磨性和高热硬性,耐腐蚀,但韧性和热稳定性较差,主要用作工具材料。

2.纤维增强复合材料

(1)玻璃纤维复合材料

这种由玻璃纤维与热固性树脂或热塑性树脂复合的材料通常称为玻璃钢,其主要成分是$\mathrm{SiO_2}$,具有强度高、工艺性好的特点,广泛应用于国民经济各行业中。玻璃钢可用来制造游艇、舰艇,各种车辆的车身及配件,各种耐腐蚀的管道、阀门、储罐、防护罩以及轴承、法兰盘、齿轮、螺钉等各种机械零件。玻璃钢的缺点是刚度差、易变形、耐热性差(使用温度不超过$200~^\circ\mathrm{C}$)、导热性差、易蠕变等。玻璃钢分为热塑性玻璃钢和热固性玻璃钢两种。

1)热塑性玻璃钢是由玻璃纤维与热塑性树脂复合而成的玻璃钢,具有高的强度、韧性和耐热性,化学稳定性好。应用较多的热塑性树脂有尼龙、聚烯烃类、聚苯乙烯类、热塑性聚酯和聚碳酸酯等五种,以尼龙的增强效果最好。

2)热固性玻璃钢是由玻璃纤维与热固性树脂复合而成的玻璃钢。其具有重量轻、比强度高、耐腐蚀性能好、介电性能优越、成形性能良好等优点,但刚度较差,耐热性不高。

(2)碳纤维增强复合材料

碳纤维是以人造纤维或天然纤维为原料,在隔绝空气的条件下经高温碳化而成。碳纤维比玻璃纤维具有更高的强度和弹性模量,并且在$2~000~^\circ\mathrm{C}$以上的高温下,其强度和弹性模量基本上保持不变,在$-180~^\circ\mathrm{C}$以下的低温下也不变脆。此外,碳纤维还具有密度小、疲劳强度高、冲击韧度高、耐水性好、化学稳定性高、摩擦因数小和导热性好等特点。因此,碳纤维是理想的增强材料,可用来增强塑料、碳、金属和陶瓷等。

1)碳纤维-树脂复合材料是碳纤维与热固性树脂结合而成的材料,可用来制作航天飞行器的外层材料,人造卫星和火箭的机架、壳体、天线及各种机器的齿轮、轴承等。

2)碳纤维-碳复合材料是以碳或石墨作基体的新型特种工程材料,能承受极高的温度和加热速度,可用作烧蚀防热材料。在航空、航天中,该材料用来制作导弹鼻锥、飞船的前缘、超声速飞机的制动装置等。

3)碳纤维-金属复合材料是在碳纤维表面镀上一层金属,再与金属基体复合而成的材料。如用碳纤维和铝锡合金制成的复合材料,是一种减摩性能比铝锡合金更优越、强度更高的轴承材料。

(3)其他纤维增强复合材料

其他纤维增强复合材料如硼纤维与树脂复合材料、晶须增强复合材料、金属纤维复合材料等。晶须是一种自由长大的金属或陶瓷型针状单晶纤维,直径在$30~\mathrm{\mu m}$以下,长度约为几毫米,其强度极高,因此,以其作为增强材料制备的复合材料具有优良的性能。但晶须制造成本高,目前多用于尖端工业或用作玻璃钢制品的局部辅助增强材料。

3.层合复合材料

层合复合材料是由两层或两层以上不同材料结合而成的。用层合法增强的复合材料可使其强度、刚度、耐蚀性、绝热性、隔热性、隔声性等性能分别得到改善。常见的层合复合材料有以下几种。

(1)双金属复合材料

双金属复合材料是用压力加工、铸造、热压、焊接、喷涂等方法将两种不同金属复合在一起的材料。如化工设备上用于代替全钛材料的包钛钢、不锈钢-普通钢复合钢板等。

(2)塑料覆层材料

近年来,国际上盛行的彩色涂层钢板就是以锌板或酸洗后的冷热轧带钢为基底,表面涂覆聚氯乙烯、聚四氟乙烯、环氧树脂等配制成多种色彩的有机涂料,从而制成的塑料覆层材料。这种材料常用于化工和食品工业。

(3)夹层玻璃

如两层玻璃板夹一层聚乙烯醇缩丁醛可制成安全玻璃。

(4)SF性三层复合材料

它是一种以钢板为基体,烧结铜网或多孔性青铜为中间层,塑料为表面层的自润滑材料。常用于表面层的塑料有聚四氟乙烯(如SF-1型)和聚甲醛(如SF-2型)。这种材料能承受的最高应力为$140~\mathrm{MPa}$,使用温度为$-195\sim270~^\circ\mathrm{C}$,混入20\%铅粉后是制造无油润滑轴承的好材料。目前,这种材料已用于汽车、矿山机械、化工机械等方面。

(5)金属塑料复合

这种材料以层状的金属与塑料相复合,具有金属的力学、物理性能和塑料的表面耐摩擦、耐磨损性能。如塑料-青铜-钢形成的复合材料,在钢与塑料之间以青铜网为媒介,可以使三者获得更可靠的结合力。

\section{其他材料}

\subsection{超导材料}

材料的电阻随温度降低而减小并最终出现零电阻的现象称为超导电现象。超导体的主要特性为:在临界温度$T_c$以下,超导体电阻为零,完全导电;在超导态下,超导体内没有磁力线通过,磁场强度恒为零;对处于$T_c$以下的超导体施加足够强的磁场时,其超导电性即消失;超导体在无外磁场条件下通电,当电流密度超过一定值后,即失去超导电性而恢复正常态。

目前,常压下具有超导电性的元素已发现20多种,其中铌(Nb)、铅(Pb)、镧(La)、钒(V)和锡(Sn)等的临界温度$T_c$较高,铌的$T_c$最高,为9.2 K,难以实际应用。超导合金是超导材料中强度最高、应力应变变小、磁场强度低的超导体,广泛应用的有Nb-Zr系合金和Ti-Nb系合金。超导陶瓷的出现,使超导体的$T_c$获得重大突破。例如,我国科学家用稀土钇(Y)取代La-Ba-Cu-O系合金中的镧(La),使$T_c > 90~\mathrm{K}$。

超导材料具有非常广泛的应用领域,如可用于以节能为目标的发电机、输电电缆、储能等电力系统的超导化,核聚变、磁流体发电等新能源装备,也可用于采用约瑟夫森器件的大型超导计算机以及核磁共振设备等医疗器械领域等。

\subsection{贮氢材料}

许多金属(或合金)可固溶氢气形成含氢的固溶体($\mathrm{MH}x$),固溶体的固溶度$[H]_\mathrm{M}$与平衡氢压$p_{\mathrm{H_2}}$的平方成正比。在一定温度和压力条件下,固溶相($\mathrm{MH}x$)与氢反应生成金属氢化物,贮氢合金正是靠其与氢起化学反应生成金属氢化物来贮氢的。目前,贮氢材料主要有镁系贮氢合金、稀土系贮氢合金和钛系贮氢合金。

1.镁系贮氢合金

镁系贮氢合金是以镁为基础、通过与Ni、Y、Cu等元素合金化而成的,具有储氢量高、主要元素镁资源丰富、成本低、重量轻等优势。其主要类型包括Mg-Ni等二元合金及Mg-Ni-Y、Mg-Cu-N等三元/多元合金。镁系贮氢合金已在移动式与固定式储氢场景中实现商业化应用,是氢能存储领域极具应用潜力的关键材料之一。

2.稀土系储氢合金

稀土系贮氢合金以$\mathrm{LaNi_5}$为典型代表,可在室温下快速活化并实现可逆吸放氢,具有平衡压低、动力学性能好、抗杂质能力强等优点,但镧价格较高,其大规模应用受到限制。工业上通常采用混合稀土(Mm)替代纯镧,并以Co、Mn、Al等元素部分替代Ni,开发多元稀土合金,在保持储氢性能的同时有效降低成本。

3.钛系贮氢合金

钛系贮氢合金主要有钛铁系合金(TiFe和$\mathrm{TiFe_2}$)和钛锰系合金($\mathrm{TiMn_{1.5}}$)。TiFe系合金室温下释氢压力不到$1~\mathrm{MPa}$,且价格低。缺点是活化困难,抗杂质气体能力差,且在反复吸释氢后性能下降。钛锰系合金中的$\mathrm{TiMn_{1.5}}$贮氢性能最佳,在室温下即可活化,但随Ti的质量分数增加室温下释氢量下降。

目前,非晶态贮氢合金由于贮氢量大而受到人们的重视,是未来的研究方向。贮氢合金可用作贮运氢的容器、氢能汽车的氢燃料贮存器。此外,还可用于工业生产中分离与回收氢,制造高纯度氢气和氢化物电极等方面。

\subsection{形状记忆合金}

在研究Ni-Ti合金时发现,原来弯曲的合金丝被拉直后,当温度升高一定值时,它又恢复到原来弯曲的形状。人们把这种现象称为形状记忆效应(shape memory effect, SME),具有形状记忆效应的金属称为形状记忆合金(SMA)。大部分形状记忆合金的形状记忆机理是热弹性马氏体相变。马氏体相变具有可逆性,即把马氏体(低温相)以足够快的速度加热,可以不经分解直接转变为高温相(母相)。

已发现的形状记忆合金种类很多,可以分为Ni-Ti系合金、Cu系合金、Fe系合金三大类。另外,近年发现一些聚合物和陶瓷材料也具有形状记忆功能,但其形状记忆原理与合金的不同。目前,已实用化的形状记忆材料只有Ti-Ni系合金和Cu系形状记忆合金。

形状记忆合金在工程上的应用很多,如紧固件、连接件、密封垫等。另外,也可以用于一些控制元件,如一些与温度有关的传感器及自动控制器,还可制作各种智能仿生机械。医学上使用的形状记忆合金主要是Ti-Ni系形状记忆合金,这种材料可以作为移植材料埋入人体,如用作固定折断骨架的销、固定接骨的接骨板,还可用于心血管介入治疗和口腔科相关治疗,例如,可用于制作血管支架,牙齿矫正线等。

\subsection{非晶态合金}

熔融的合金以极高的速度冷却,凝固后结构呈玻璃态,这种形态的合金称为非晶态合金,俗称金属玻璃。非晶态合金与金属相比,成分相同,但结构不同,从而引起二者在性能上的差异。非晶态合金具有超耐蚀性、高磁导率、恒弹性、高强韧性、低热膨胀系数和磁致伸缩效应等许多优异性能。

非晶态是一种新的、特殊的物质状态,至今尚无统一的严格定义。一般来说,它具有以下基本特征:不具有长程有序晶体结构;具有高达$10^{13}~\mathrm{Pa\cdot s}$以上的黏滞系数;在某一窄的温度区内能够发生明显的结构变化。原则上,所有的金属熔体都可通过急冷制成非晶体。也就是说,只要冷却速度足够快,使熔体中原子来不及做规则排列就完成凝固过程,即可形成非晶态金属。非晶态合金的主要缺点表现在两方面:一是采用急冷法制备材料,其厚度受到限制;二是热力学性能不稳定,受热有晶化倾向。解决上述缺点的办法主要是采取表面非晶化及微晶化。非晶态金属在电子工业中也得到广泛应用,如制作漏电开关、磁头材料等。

\subsection{磁性材料}

磁性材料可分为软磁材料、硬磁材料、磁记录材料及一些特殊用途的磁性材料等。

1.软磁材料

它在较低的磁场中被磁化而呈强磁性,但在磁场去除后磁性基本消失。这类材料被用作变压器、继电器、电感器的铁心,发电机与电动机转子和定子以及磁路中的磁轭材料等。软磁材料还包括耐磨高导磁材料、矩磁材料、恒导磁材料、磁温度补偿材料和磁致伸缩材料。典型的软磁材料有纯铁、Fe-Si合金(硅钢)、Ni-Fe合金、Fe-Co合金、Mn-Zn铁氧体、Ni-Zn铁氧体和Mg-Zn铁氧体。另外,非晶态合金也是很好的软磁材料。

2.硬磁材料

硬磁材料也称为永磁材料,是指材料被外磁场磁化以后,去掉外磁场仍然保持着较强剩磁的材料。它也是人类最早发现和应用的磁性材料。常用的永磁材料有以下三种。

(1)变形永磁钢

变形永磁钢的牌号有2J63、2J64、2J65、2J67等。这类合金有高的矫顽力和剩磁,用于制造电气仪表中的永磁体,如继电器、同步电动机的磁钢等。

(2)铝镍钴永磁合金

铝镍钴永磁合金的牌号由合金主成分的元素符号即AlNiCo和阿拉伯数字两部分组成,牌号首字母为“S”表示烧结永磁合金,无“S”表示铸造永磁合金。阿拉伯数字分别表示材料的主要磁性能最大磁能积($\mathrm{kJ/m}$)和材料的磁极化强度矫顽力($\mathrm{kA/m}$)的十分之一。例如,采用烧结法生产、标称最大磁能积为$33~\mathrm{kJ/m}$、磁极化强度矫顽力为$150\mathrm{kA/m}$的铝镍钴永磁合金牌号表示为S AlNiCo 33/15。这类合金磁性能好,可用铸造和烧结方法制造,广泛用于控制系统中,如磁滞电动机转子、继电器、指示器、温度表以及电磁仪表的铁心等。

(3)稀土永磁材料

稀土永磁材料分为稀土钴永磁材料和Nd-Fe-B系合金。Nd-Fe-B系合金典型成分为$\mathrm{Nd_{15}Fe_{77}B_8}$,硬磁化相为金属间化合物$\mathrm{Nd_2Fe_{14}B}$。Nd-Fe-$\beta$系合金主要用于电动机制造,具有体积小、重量轻、比功率大、效率高的优势。此外,电声器件中的传声器、高频扬声器和立体声耳机,磁流体密封器、磁水器、测量仪器、磁力器、磁传感器等都是Nd-Fe-$\beta$系合金的主要应用领域。在医疗方面,核磁共振成像仪中,也要用到Nd-Fe-$\beta$系合金。

\subsection{光纤材料}

光导纤维(简称光纤)是光通信的传输材料。这种通信电路不是用一般的电线和电缆,而是采用与头发丝一样细的透明玻璃纤维制成的光缆。其基本原理是把声音变为电信号,由发光元件(如GaP)将电信号变为光信号,由光导纤维传向远方,再由接收元件(CdS、ZnSe)恢复为电信号,使受话机发出声音。

利用光导纤维进行远距离通信的效率非常高,比普通电缆通信的效率高$10^{10}$倍以上。同时,其通信容量大、质量轻、耐腐蚀、不怕电子干扰,而且保密性好、施工方便、成本低,可节约大量有色金属。例如,$1~000~\mathrm{km}$长的同轴电缆,大约需要5万吨铜,而采用光导纤维则只需要几十千克石英玻璃就可拉成$1~000~\mathrm{km}$的光导纤维。

光导纤维是超纯石英掺杂磷、锗等元素制成的,其光损已接近理论值($0.2~\mathrm{dB/km}$)。目前,正在发展中的多组分玻璃纤维,熔点低、易生产、损耗小(理论值为$0.01\sim0.001~\mathrm{dB/km}$),但其强度较低。目前,普遍应用的还有塑料光纤,虽然其每千米的损耗高($100\sim200~\mathrm{dB/km}$),但由于纤维质地柔软、价格便宜,因此在短距离传输中仍有十分广泛的应用。

\subsection{压电材料}

当对$\alpha$石英晶体在某些特定方向上加力时,在力方向的垂直平面上出现正、负束缚电荷,这种现象称为压电效应。钛酸钡($\mathrm{BaTiO_3}$)是第一个被发现可以制成陶瓷的铁电体,其熔点为$1~618~^\circ\mathrm{C}$,室温下为铁电体,晶体呈四方结构,在$120~^\circ\mathrm{C}$时转变为立方晶相,此时铁电性消失。

压电材料中研究得比较早的是石英晶体,它的机电性能稳定,没有内耗,在频率稳定器、扩音器、电话、钟表等领域里都有广泛应用。此外,酒石酸钾钠、磷酸二氢铵、钽酸锂、碘酸锂、铌酸锂等晶体也都是比较好的压电材料。其中,使用最多的是铌酸锂($\mathrm{LiNbO_3}$)。铌酸锂属畸变的钙钛矿结构,密度为$4.64~\mathrm{g/cm^3}$,熔点为$1~253~^\circ\mathrm{C}$。

压电材料作为体材料,已在机电转换领域得到广泛应用。基于逆压电效应,其在高压驱动电场下能够产生高强度超声波,并以此作为动力被广泛用于超声清洗、乳化、焊接、打孔、粉碎及分散等设备中的机电换能器。凭借其优良的机电性能和高化学稳定性,压电材料也常用于制造电声器件,例如拾音器、扬声器和蜂鸣器等。此外,在压力计、流量计、计数器等多种测量仪表中,压电材料也扮演着重要角色。水声换能器则是压电材料的另一项重要应用。

\subsection{生态金属材料}

生态金属材料是指具有满意的使用性能和优良的环境协调性,或者是能够改善环境的金属材料。所谓环境协调性,是指对资源和能源消耗少、对环境污染小和再生循环利用率高。

金属材料生态设计采取的主要方法有以下两种。

(1)采用通用合金

从提高金属材料的再生循环性能出发,金属制品的全部零部件由单一合金系来制造最为理想,而且合金系中组元越少,合金的再生循环性能越好。

(2)采用简单合金

简单合金在成分设计上应有以下特点:

1)合金组元、规格简单,再生循环过程中易于分选;

2)原则上不添加目前尚不能精炼脱除的元素;

3)尽量不使用环境协调性不好的合金元素。

从这种设计思想出发,研制简单合金时应遵循两个原则:一是在维持合金高性能的前提下,尽量减少合金组元数;二是获取合金高性能时,以控制显微组织作为加入合金元素的替代方法。

\subsection{纳米材料}

1.纳米材料的特性

凡是至少在一维方向上的线度为$1\sim100~\mathrm{nm}$($1~\mathrm{nm}=10^{-9}~\mathrm{m}$)的单元和由这种纳米单元作基本结构单元的材料均称为纳米材料。纳米单元包括纳米膜、纳米线(晶须、大分子链)和纳米颗粒。纳米微粒由于尺寸小、表面大和表面活性高而具有小尺寸效应、表面效应和量子尺寸效应。

2.纳米材料的应用

(1)在环境科学中的应用

光催化氧化法是一种高效的深度氧化技术,通过纳米级光催化剂的作用,将光能转化为化学反应所需的能量,使催化剂周围的氧气及水分子被激发成极具氧化力的自由基或负离子,可以将水体中的烃类、卤代烃、羧酸、表面活性剂、染料、含氮有机物、有机磷杀虫剂等较快地完全氧化为$\mathrm{CO_2}$和$\mathrm{H_2O}$等无害物质,被广泛应用于水体中有机污染物的除毒、脱色、去臭。例如,利用纳米$\mathrm{TiO_2}$催化剂可高效降解或去除水体中的有机物、重金属离子和菌类。

(2)在复合材料中的应用

层状硅酸盐与聚合物形成的纳米复合材料,由于其纳米尺寸效应和强的界面黏结,具有高耐热性、高强度、高模量和低的膨胀系数,而密度仅为一般复合材料的65\%$\sim$75\%,因此可作为新型高性能工程塑料广泛应用于航空、汽车等行业。

(3)在涂料中的应用

将用纳米铁粉、镍粉、铁氧体粉末改性的有机涂料涂在飞机、导弹、军舰等武器上,可使该装备具有隐身性能。

(4)磁性纳米材料

磁性纳米材料的典型应用是磁流体和磁记录材料。目前,磁流体大多数以$\mathrm{Fe_3O_4}$纳米微粒为磁性粒子;磁性纳米微粒由于具有尺寸小和矫顽力高的特性,用其制作的磁记录材料可以提高声噪比,改善图像质量。

(5)超微粒传感器

超微粒一般具有吸收红外线等特点,而且表面积巨大、表面活性高,对周围环境(温度、气氛、光、湿度等)敏感,可用于制作敏感度高的超小型、低能耗和多功能传感器。

(6)在生物和医学上的应用

纳米微粒的尺寸一般比生物体内的细胞小得多,可用于细胞分离、细胞染色。纳米微粒溶液还可用于制作特殊药物或新型抗体进行局部定向治疗等。

\subsection{梯度功能材料}

梯度功能材料(functional gradient materials, FGM)是采用全新的设计理念而开发的新型功能材料,它将两种或多种材料复合使其成分和结构呈连续梯度变化。其设计思想是使材料的构成要素(包括组成、结构、结合形式等)从一侧向另一侧呈现连续性变化,从而得到单一或复合功能渐变的非均质材料。在复杂环境下使用时,梯度功能材料比均质材料具有更大优势。

梯度功能材料的性能取决于其体系组分选择及内部结构的合理设计,而且需要采取有效的制备技术来保证材料的设计。目前,制备梯度功能材料的主要方法有化学气相沉积法(chemical vapor deposition, CVD)、自蔓延高温合成法(self-propagation high-temperature synthesis, SHS)、电化学法(electrochemistry, EC)、物理气相沉积法(physical vapor deposition, PVD)、等离子喷涂法(plasma spraying, PS)、粉末冶金法(powder metallurgy, PM)、激光熔覆法(laser cladding, LC)等。

梯度功能材料具有比传统复合材料更为优异的性能,在航天工业、能源工业、电子工业、光学

材料、化学工程和生物医学工程等领域具有重要的应用前景。以美国航天飞机为例,目前唯一的再用型火箭发动机其再用次数设计目标为100次,但实际只能使用20~30次,如果改用具有良好隔热性能与缓和热应力型的梯度功能材料,有望达到并超过原有的设计目标。此外,由羟基磷灰石陶瓷与钛合金组成的梯度功能材料可用于仿生活性人工关节和牙齿,金属-陶瓷结合的梯度功能材料可用于核反应堆材料,消除热传递及热膨胀引起的应力,从而解决界面问题。

随着应用领域的逐步扩大,梯度功能材料的组成也由金属-陶瓷的单一组合发展成为金属-合金、非金属-非金属、非金属-陶瓷、高分子膜(I)-高分子膜(II)等多种组合,应用前景十分广阔。

\section*{复习思考题}

 
5-1 铝合金是如何分类的?各类铝合金可通过哪些途径进行强化?

5-2 变形铝合金包括哪几类铝合金?用2A01作铆钉应在何状态下进行铆接?在何时得到强化?

5-3 简述铝合金的分类方法、牌号、性能和用途。

5-4 简述铜合金的分类方法、牌号、性能和用途。

5-5 铜合金与纯铜相比有何特点与优势?

5-6 与纯铜相比,黄铜与青铜这两种铜合金各有什么特点?有哪些用途?

5-7 钛合金分为几类?钛合金的性能特点与应用是什么?

5-8 简述热塑性塑料与热固性塑料的区别。

5-9 与金属材料相比,工程塑料的主要性能特点是什么?

5-10 简述制约高分子材料大量应用的因素,以及进一步提高其性能的措施。

5-11 某厂使用库存数年的尼龙绳吊具时,在承载能力远大于吊装应力时却发生断裂事故,试分析其原因。

5-12 对于下面列出的四种构件,在给出的材料中选择最合适的。
构件:(1)电源插座;(2)飞机窗玻璃;(3)化工管道;(4)齿轮。
材料:(1)酚醛树脂;(2)尼龙;(3)聚氯乙烯;(4)聚甲基丙烯酸甲酯

5-13 简述常用特种陶瓷的种类、性能特点及应用。

5-14 陶瓷材料的力学性能有何特性?如何进一步改善其特性?

5-15 常用的复合材料有哪些?

5-16 复合材料的性能有什么特点?

5-17 简述氧化物基金属陶瓷的性能与特点。

5-18 常用的增强纤维有哪些?试比较它们的性能特点。

5-19 与纤维增强塑料相比,纤维增强金属在性能上有何特点?举例说明它们的用途。

5-20 比较玻璃钢、碳纤维增强塑料、硼纤维增强塑料、碳化硅纤维增强塑料和凯芙拉(Kevlar)纤维增强塑料的性能特点,举例说明它们的用途。

5-21 纳米材料具有什么样的性能特点?

5-22 未来材料有何特点?其发展方向是什么?


\chapter{失效及选材}

\section{零件的失效}

所谓零件失效,主要指由于某种原因导致零件的尺寸、形状或其组织与性能发生变化而不能圆满地完成指定的功能。零件的失效一般包括以下几种情况。

1) 零件完全破坏,不能继续工作;

2) 虽能继续工作,但不能完全起到预期的作用。

3) 零件严重损伤,继续工作不安全。

\subsection{失效的主要因素}

机械产品失效的原因很多,错综复杂,主要有设计、材质、工艺、安装、使用等方面的因素。

1. 设计因素

零件的结构形状、尺寸设计不合理易引起失效。例如,结构上存在尖角、尖锐缺口或圆角过渡过小,产生应力集中引起失效;安全系数小,未达到实际承载能力等。

2. 材质因素

选材不当,所选用的材料性能未达到使用要求,或材质较差(夹杂物、杂质元素超标等)。

3. 工艺因素

零件在锻造过程中的夹层、冷热裂纹,焊接过程中的未焊透、偏析、冷热裂纹,铸造过程中的夹杂、疏松,机械加工过程的尺寸偏差和表面粗糙度不合适,热处理工艺产生的裂纹等工艺因素都有可能引起零件失效。

4. 安装使用不当

零件安装不良,操作失误、过载使用、维修保养不当等,均可导致零件在使用中失效。

零件失效的原因是多种多样的,实际情况也往往错综复杂,失效分析就是寻找构件断裂、变形、磨损、腐蚀等失效现象的特征和规律,并从中找出损坏的成因。

\subsection{失效形式}

根据零件损坏的特点,所承受载荷的形式及外界条件,可将失效分为下列几种类型。

1. 变形失效

变形失效包括弹性变形失效、塑性变形失效和蠕变变形失效。其特点是非突发性失效,一般不会造成灾难性事故。但塑性变形失效和蠕变变形失效有时也可能造成灾难性事故,应引起充分重视。

2. 断裂失效

断裂失效包括以下几种形式。

(1) 延性断裂失效

延性断裂失效的特点是断裂前有一定程度的塑性变形,一般是非灾难性的,用电子显微镜观察断口时,到处可见韧窝断裂形貌,观察断口附近金相(显微)组织,可见到有明显塑性变形层组织。

(2) 脆性断裂失效

脆性断裂失效在断裂前无明显的塑性变形,是突发性的断裂。电子显微镜下它的特征为河流花样或冰糖状形貌,如解理断裂和沿晶界断裂。

(3) 疲劳断裂失效

疲劳的最终断裂是瞬时的,因此危害性较大,甚至会造成重大事故。用电子显微镜观察断口时,在疲劳扩展区可看到疲劳特征的条纹。疲劳断裂占工程上断裂失效大多数,占比为80\%以上。

(4) 蠕变断裂失效

在高温缓慢变形过程中发生的断裂属于蠕变断裂失效。蠕变的最终断裂也是瞬时的。在工程中常见的多属于高温、低应力的沿晶界蠕变断裂。

3. 腐蚀失效

金属与周围介质之间发生化学或电化学作用而造成的破坏属于腐蚀失效。其中,应力腐蚀、氢脆和腐蚀疲劳等是突发性失效,而点腐蚀、缝隙腐蚀等局部腐蚀和大部分均匀腐蚀失效不是突发性的,而是逐渐进展的。腐蚀失效的特点是失效形式众多,机理复杂,占金属材料失效事故中的比率较大。

4. 磨损失效

凡相互接触并作相对运动的物体,由于机械作用所造成的材料位移及分离的破坏形式称为磨损。磨损失效所造成的后果一般不像断裂失效和腐蚀失效那样严重,然而一些灾难性的事故也来自磨损失效。磨损失效主要有黏着磨损、磨粒磨损和腐蚀磨损等几种失效形式。

\section{选材的一般原则}

在零件设计时,不仅要考虑材料的性能能够适应零件的工作条件,使零件经久耐用,而且还要求材料有较好的加工工艺性和经济性。只有满足这三方面的要求,所选择的材料才能视为合理材料。

\subsection{使用性能原则}

使用性能是保证零件完成规定功能的必要条件。对使用性能的要求,是在分析零件工作条件和失效形式的基础上提出来的。零件的工作条件包括以下三个方面。

1. 受力状况

受力状况主要是指所受载荷的类型(如动载、静载、循环载荷或单调载荷等)和大小、载荷的形式(如拉伸、压缩、弯曲或扭转等)以及载荷的特点(如均布载荷或集中载荷等)。表6-1列出了几种常见机械零件的工作条件、失效形式及其要求的力学性能。

\begin{table}[H]
\centering
\caption{几种常见机械零件的工作条件、失效形式及其要求的力学性能}
\footnotesize\setlength{\tabcolsep}{3pt}
\begin{tabular}{|p{1.2cm}|p{2cm}|p{2cm}|p{2.5cm}|p{4cm}|p{4cm}|}
\hline
\multirow{2}{1.2cm}{零件} & \multicolumn{3}{c|}{工作条件} & \multirow{2}{4cm}{常见失效形式} & \multirow{2}{4cm}{力学性能指标} \\
\cline{2-4}
 & 变形方式 & 载荷性质 & 其他 & & \\
\hline
紧固螺栓 & 拉、剪 & 静 & & 变形、断裂 & 强度、塑性 \\
\hline
传动轴 & 弯、扭 & 循环、冲击 & 轴颈处摩擦、振动 & 疲劳破坏、变形、轴颈处磨损 & 综合力学性能 \\
\hline
齿轮 & 压、弯 & 循环、冲击 & 强烈摩擦、振动 & 磨损、疲劳麻点、齿断裂 & 表面有高硬度及高的疲劳极限,心部有较高强度及韧性 \\
\hline
弹簧 & 扭(螺旋簧) & 循环、冲击 & 振动 & 弹性丧失、疲劳破坏 & 弹性极限、屈强比、疲劳强度 \\
\hline
油泵柱塞副 & 压 & 循环、冲击 & 摩擦、油的腐蚀 & 磨损 & 硬度、抗压强度 \\
\hline
冷作模具 & 复杂组合变形 & 循环、冲击 & 强烈摩擦 & 磨损、脆断 & 高硬度、高强度、足够的韧性 \\
\hline
压铸型 & 复杂组合变形 & 循环、冲击 & 高温、摩擦、金属液腐蚀 & 热疲劳、磨损、脆断 & 高温强度、抗热疲劳性、足够韧性与热硬性 \\
\hline
\end{tabular}\end{table}

2. 环境状况

环境状况主要是指温度特性(如低温、常温、高温或变温等)以及介质情况(如有无腐蚀或摩擦作用)等。

3. 特殊要求

特殊要求主要是指对导电性、磁性、热膨胀、密度、外观等的要求。零件的失效形式如前所述,主要包括变形、断裂和表面损伤三个方面。

通过对零件工作条件和失效形式的全面分析,确定零件对使用性能的要求,利用使用性能与实验室性能的相应关系,将使用性能转化为力学性能指标,如强度、韧性或耐磨性等,这是选材最关键的步骤,也是最困难的一步。之后,根据零件的几何形状、尺寸及工作中所承受的载荷,计算出零件中的应力分布,再由工作应力、使用寿命或安全性与力学性能指标的关系,确定对力学性能指标要求的具体数值。

\subsection{工艺性原则}

对于任何机械零件,在根据使用要求选择合适的材料后,还要通过一定的加工方法或加工工艺制造出来。金属材料的加工(成形)方法有铸造、压力加工、焊接、机械加工等。陶瓷材料可通过粉末压制烧结成形,有的还需进行磨削加工、热加工;高分子材料利用有机物原料,通过热压、注塑、热挤压等方法成形,有的不需进行切削加工、焊接等加工过程。

材料的加工方法和加工工艺性,不仅影响机械零件的外观,还影响零件的性能,甚至影响生产率和成本。因此,选用的材料应具有良好的工艺性,至少要有可行的工艺性。常用加工方法的良好工艺性表现为以下几个方面。

1) 铸造合金应有高的流动性,小的缩松、缩孔、偏析和吸气性倾向。

2) 塑性加工材料应有高的塑性和低的变形抗力。

3) 切削加工材料应有小的切削力,切屑处理容易,对刀具的磨损小等。

4) 热处理要求材料过热敏感性小,氧化和脱碳倾向小,淬透性高,变形和开裂倾向小等。

应当指出,材料在不同的状态具有不同的工艺性,而且某种工艺性好,不等于其他工艺性也好。例如,2Cr13等马氏体不锈钢退火后切削加工性尚好,但焊接时容易开裂。奥氏体不锈钢塑性好,但切削加工性差。镁合金和有些钛合金冷变形性差,而在加热状态下则有良好的变形加工能力。

\subsection{经济性原则}

在满足零件性能要求的前提下,选材应使其总成本(包括材料和加工费用等)尽可能低。

1) 材料选用应考虑我国资源,例如尽可能选择那些以锰、硅、钼、稀土等元素完全或部分代替镍、铬等稀缺元素的合金。

2) 应考虑国内生产和供应情况,品种不宜过多。

3) 应考虑选用节省材料和加工成本的工艺方法,如精铸和精锻等。

\section{典型零件选材分析}

下面通过典型零件的选材举例,将选材中的有关问题加以分析,以便加深对选材方法的认识。

\subsection{齿轮类零件的选材与工艺}

1. 齿轮的工作条件、失效形式及性能要求

(1) 工作条件和失效形式

齿轮是应用广泛的机械零件,主要起传递扭矩、变速或改变传力方向的作用。其工作条件如下:

1) 传递扭矩时齿根部承受较大的交变弯曲应力;

2) 齿啮合时齿面承受较大的接触压应力并受强烈的摩擦和磨损;

3) 换挡、起动、制动或啮合不均匀时承受一定冲击力。

齿轮的失效形式主要是齿根的折断(包括疲劳断裂和冲击过载断裂)和齿面损伤(包括接触疲劳麻点剥落和过度磨损)。

(2) 性能要求

根据齿轮的工作条件和失效形式,齿轮材料应具有如下性能。

1) 高的抗弯疲劳强度,以防齿轮疲劳断裂。

2) 足够高的齿心强度和韧性,以防齿轮过载断裂。

3) 足够高的齿面接触疲劳强度和硬度及好的耐磨性,以防齿面损伤。

4) 较好的工艺性能,以便于制造和热处理等。

2. 齿轮的选材及热处理

下面以机床齿轮为例来分析齿轮的选材及热处理工艺。

机床齿轮工作平稳,无强烈冲击,负荷不大,转速中等,对齿轮心部强度和韧性的要求不高,一般选用40钢或45钢制造。经正火或调质后再经高频感应加热淬火,齿面硬度可达52 HRC左右,齿心硬度为220~250 HBW,完全可以满足性能要求。对于一部分性能要求较高的齿轮,可用中碳低合金钢(如40Cr、40MnB、45Mn2等)制造,齿面硬度可提高到58 HRC左右,心部强度和韧性也有所提高。

机床齿轮的加工工艺路线如下:

下料—锻造—正火—粗加工—调质—半精加工—高频感应加热淬火+低温回火—精磨—成品

正火可使组织均匀化,消除锻造应力,调整硬度,改善切削加工性,对于一般齿轮,正火也可作为高频感应加热淬火前的最后热处理工序;调质可使齿轮具有较高的综合力学性能,提高齿心的强度和韧性,使齿轮能承受较大的弯曲应力和冲击载荷,并减小淬火变形;高频感应加热淬火可提高齿轮表面硬度和耐磨性,提高齿面接触疲劳强度;低温回火可在不降低表面硬度的情况下消除淬火应力、防止产生磨削裂纹和提高齿轮抗冲击的能力。

\subsection{轴类零件的选材与工艺}

轴类零件是机械传动中广泛使用的重要结构件。机床主轴、花键轴、丝杠,内燃机的曲轴、连杆和汽车传动轴、半轴等均属于轴类零件。其主要作用是支承传动零件并传递扭矩。

1. 轴的工作条件、失效形式及性能要求

(1) 工作条件和失效形式

轴的工作条件如下:

1) 承受交变扭矩载荷、交变弯曲载荷或拉-压载荷;

2) 局部(轴颈、花键等处)承受摩擦和磨损;

3) 特殊条件下受高温或介质作用。

轴的失效形式主要是疲劳断裂和轴颈处磨损,有时会发生冲击过载断裂,个别情况下也会发生塑性变形或腐蚀失效。

(2) 性能要求

根据轴的工作条件及失效形式,轴的材料应具有如下性能。

1) 高的疲劳强度,以防轴疲劳断裂。

2) 优良的综合力学性能,即较高的屈服强度和抗拉强度、较好的韧性,以防塑性变形及在过载或冲击载荷作用下的折断和扭断。

3) 局部承受摩擦的部位具有高硬度和耐磨性,以防磨损失效。

4) 在特殊条件下工作时应具有特殊性能,如蠕变抗力、耐腐蚀性等。

2. 轴的选材及热处理

下面以汽轮机主轴为例分析轴的选材及热处理工艺。

汽轮机主轴尺寸大,工作载荷大,受弯矩、扭矩及离心力和温度的联合作用,工作条件恶劣。其失效形式主要是蠕变变形和由内部缺陷(如白点、夹杂物、焊接裂纹等)引起的低应力脆断或疲劳断裂和应力腐蚀断裂。因此,除要求材料有较高的强度及足够的塑性和韧性外,还要求其锻件中不出现较大的夹杂物、氢引起的微裂纹(白点)及焊接裂纹。需要根据汽轮机功率和主轴工作温度选用不同的材料。对于工作温度在450 $^\circ$C以下的主轴,不必考虑高温强度。若汽轮机功率较小(<12 000 kW),主轴尺寸较小,可选用45钢;若汽轮机功率较大(>12 000 kW),主轴尺寸较大,必须选用35CrMo钢,以提高淬透性。对于工作温度在500 $^\circ$C以上的主轴,由于汽轮机功率大,可达125 000 kW以上,要求较高的高温强度,选用合金结构钢制造。一般高中压主轴选用25CrMoV或27Cr2MoV钢,低压主轴选用15CrMoV或17CrMoV钢。对于燃气轮机主轴,由于工作温度更高,要求材料具有更高的高温强度,一般选用合金结构钢20Cr3MoWV(工作温度<540 $^\circ$C)、铁基高温合金Cr14Ni26MoTi(GH2132)(工作温度<650 $^\circ$C)和Cr14Ni35MoWTiAl(GH2135)(工作温度<680 $^\circ$C)。

3. 汽轮机主轴制造工艺分析

汽轮机主轴的加工工艺路线如下:

钢锭—锻造—第一次正火—去氢处理—第二次正火—高温回火—机械加工—成品

第一次正火可消除锻造应力,使组织均匀;去氢处理(又称为去氢退火)的目的是使氢自锻件中扩散出去,防止产生白点;第二次正火是为了获得细片状珠光体,提高高温强度;高温回火可消除正火应力,并使合金元素V、Ti充分进入碳化物中,而使Mo充分溶入铁素体中,进一步提高其高温强度。

\subsection{飞机起落架的选材与工艺}

1. 起落架的工作条件及性能要求

起落架是飞机起飞、降落的关键性部件,承受着飞机的全部重量。其性能直接关系飞机的使用与安全性。飞机起飞、降落滑行时,由于跑道不可能绝对平整,所以起落架承受着强烈的振动和交变载荷。每一次起飞和降落,在承受飞机重量条件下,起落架的振动和冲击等载荷都将经历一次循环。因此,起落架失效形式主要是疲劳破坏。根据起落架的工作条件,它必须具有以下性能。

1) 高的强度,以承受飞机载荷。

2) 高的疲劳强度,以抵抗疲劳破坏。

3) 高的冲击韧度,以抵抗冲击载荷。

2. 起落架的选材分析

起落架的支柱是主要受力部件,是经模锻、机械加工的上接头和下筒焊接而成的。因此,材料的选择在考虑强度要求的同时,还要考虑工艺性。

根据计算,上接头和下筒所选用的材料应具有很高的抗拉强度,可选用超高强度钢。查阅相关标准可知,30CrMnSiNi2A、30CrMnSi以及40CrNiMo的抗拉强度都可以达到要求,但后两种钢的韧性和塑性不及30CrMnSiNi2A。从淬透性方面来分析,30CrMnSiNi2A的淬硬深度在三种钢中最大,可满足起落架支柱淬透性要求。

起落架支柱是焊件,所选用钢材应有良好的焊接性。30CrMnSi的焊接性略高于30CrMnSiNi2A,而40CrNiMo的焊接性差。30CrMnSiNi2A钢可用电弧焊焊接,若用氩弧焊,则更易焊接,可以满足焊接生产要求。

就切削加工性而言,三种钢中30CrMnSi退火状态最好。但30CrMnSiNi2A在不完全退火或正火后可切削加工。从力学性能和工艺性方面综合考虑,应选用30CrMnSiNi2A。

3. 起落架制造的工艺分析

起落架的加工工艺路线如下:

模锻—正火+低温回火—机械加工—焊接—去应力退火(600~650 $^\circ$C)—等温淬火(900 $^\circ$C加热、保温后于180~230 $^\circ$C等温60 min)—低温回火(250~300 $^\circ$C)—探伤—磷化

正火+低温回火可以改善模锻后的组织,同时改善切削加工性;去应力退火可以消除焊接应力,避免焊接裂纹;等温淬火(180~230 $^\circ$C)+低温回火可以获得贝氏体和马氏体的混合组织;再通过低温回火进一步消除应力,提高韧性。

\subsection{航空发动机压气机叶片选材与工艺}

压气机叶片是航空发动机的核心部件之一,其性能直接决定了发动机的效率、可靠性和推重比。

1. 压气机叶片的工作条件、失效形式及性能要求

(1) 工作条件和失效形式

压气机叶片,尤其是高压压气机后几级的叶片,工作在极其严苛的环境中,不仅需要在高压气动环境中耐受高周疲劳载荷,而且还需抵抗介质腐蚀的影响。叶片的主要失效形式为蠕变变形、疲劳断裂、冲击损伤、腐蚀及磨损等。

(2) 性能要求

根据其工作条件,航空发动机压气机叶片的材料应具有以下性能。

1) 高温强度与蠕变抗力,能够在高温环境下保持高机械强度并抵抗蠕变变形。

2) 高疲劳强度,具备良好的抗高周与低周疲劳性能。

3) 抗氧化和耐腐蚀性能,能够有效抵御高温燃气及多种介质的侵蚀。

2. 压气机叶片的选材分析

根据叶片所处位置(低压前端、高压后端等)不同,压气机叶片主要使用三类材料:低压压气机和高压压气机前几级叶片主要采用钛合金,可在500~600 $^\circ$C下长期稳定工作,但钛合金冶炼和加工难度很大;高压压气机后几级叶片主要用镍基高温合金,在600~750 $^\circ$C高温下,仍能保持极高的强度和蠕变抗力,但因镍基高温合金密度大,硬度高,加工极其困难;对于低压压气机前端的大尺寸风扇叶片和部分宽弦叶片主要采用树脂基复合材料,减重效果显著,疲劳寿命远超金属材料,但树脂基复合材料的耐温能力有限,长期工作温度一般不超过$200\sim250\text{ }^\circ\text{C}$,无法用于压气机中后段。

航空发动机压气机叶片的选材是一个典型的“性能-重量-成本-可靠性”多维权衡过程,是实现高推重比航空发动机的关键技术决策之一。

3. 压气机叶片制造工艺分析

压气机叶片的加工工艺路线如下:

模锻—机械加工—固溶处理($980\sim1\ 020\text{ }^\circ\text{C}$)—时效处理($540\sim580\text{ }^\circ\text{C}$)—退火—表面处理—探伤

固溶处理旨在使合金元素充分溶解,以提高材料的强度和韧性;时效处理则通过低温析出强化相,进一步提升材料的硬度和耐磨性;表面处理包括阳极氧化、电镀、化学镀和物理气相沉积等方法,可显著增强叶片的耐磨性与耐腐蚀性能。

\section*{复习思考题}

 
6-1 零件的失效形式主要有哪些?分析零件失效的主要目的是什么?

6-2 引起零件失效的基本因素是什么?

6-3 选择零件材料应遵循哪些原则?在利用手册上的力学性能数据时应注意哪些问题?

6-4 从选材的角度如何理解“好钢要用在刀刃上”?
6-5 简要叙述齿轮材料应具备的性能。

6-6 试论述常用力学性能指标在选材中的意义。

6-7 断裂是材料破坏的重要形式之一,与材料的组分、服役条件等因素密切相关,试从材料组分、当时船体的服役条件等方面分析造成“泰坦尼克号”邮轮沉船的可能原因。

6-8 零件的使用要求包括哪些?以车床主轴为例说明其使用要求。

6-9 试确定下列齿轮的材料,并确定热处理方式。
(1)承受冲击的高速重载齿轮;
(2)小模数仪表用无润滑小齿轮;
(3)农机用受力小、无润滑大型直齿圆柱齿轮。

6-10 分析断裂韧度在选材中的意义。

6-11 今有一储存液化气的压力容器,工作温度为$-196\text{ }^\circ\text{C}$,试回答下列问题,并说明理由。
(1)低温压力容器要求材料具有哪些力学性能?
(2)在下列材料中选择何种材料较为合适?
低温合金高强度钢;奥氏体不锈钢;变形铝合金;黄铜;钛合金;工程塑料。

6-12 试叙述飞机起落架、航空发动机压气机叶片应具备的基本性能。


\chapter{铸造}

铸造是将熔融金属浇注、压射或吸入铸型型腔,冷却凝固后获得一定形状和性能的零件或毛坯的金属成形工艺。它是金属材料液态成形的一种重要方法。

铸造使金属一次成形,工艺灵活性大,各种成分、尺寸、形状和重量的铸件几乎都能适应,且成本低。铸造适用于形状复杂、特别是内腔复杂的毛坯或零件;对于不宜锻压生产和焊接的材料,铸造生产方法具有特殊的优势。因此,铸造在机器制造业中应用极其广泛。金属铸造成形的原理和方法,还被广泛借鉴,用于高分子材料、陶瓷材料及复合材料的成形。

铸造生产也存在着不足之处。铸造组织的晶粒比较粗大,内部常有缩孔、缩松、气孔、砂眼等铸造缺陷,因而铸件的力学性能一般不如锻件;铸造生产工序繁多,工艺过程较难控制,致使铸件的废品率较高;铸造工人的工作条件较差、劳动强度比较大。

随着科学技术的不断进步,铸造技术也获得了飞速的发展。现代铸造技术是集计算机技术(如计算机凝固模拟、应力计算、计算机辅助铸造工艺设计、计算机控制熔炼等)、信息技术、自动控制技术、真空技术、电磁技术、激光技术、新材料技术、现代管理技术与传统铸造技术之大成,形成优质、高效、低耗、清洁、灵活的铸造生产的系统工程。这些现代技术的应用使铸件的表面精度、内在质量和力学性能都有显著提高,使铸造生产率及铸件的工艺出品率大大提高,也使工人的劳动强度减小,劳动条件大为改善。目前,随着全球可持续发展战略的实施,铸造生产正朝着绿色化、高度专业化、智能化、网络化、虚拟化与集约化生产和铸件的高性能化、精密化、轻薄化方向发展。

\section{液态金属成形理论基础}

在液态金属成形过程中,液态金属的充型及其收缩是影响成形工艺及铸件质量的两个主要问题,工艺参数及工艺方案(如熔炼和浇注温度、浇注系统、冒口位置及尺寸等)的确定和铸造缺陷(如冷隔、浇不足、缩松、缩孔、变形、应力、裂纹等)的产生均与两者有关。

\subsection{液态合金的充型能力}

1. 合金的流动性

合金的流动性是指液态合金填充铸型的能力,即流动能力。液态合金具有良好的流动性,不仅易于获得形状复杂、轮廓清晰的薄壁铸件,而且有利于气体和夹杂物在凝固过程中向液面上浮和排出,有利于合金的补缩。合金流动性不好,则会产生冷隔、浇不足、气孔、夹渣及缩孔等铸造缺陷。因此,合金的流动性是衡量合金铸造性能优劣的主要指标之一。

合金的流动性通常用浇注螺旋形流动性试样的方法来测量。先将液态合金在相同的浇注温度或相同的过热度条件下,浇注成如图7-1所示的试样,然后比较各种合金浇注的试样的长度。浇注的试样越长,合金的流动性越好。

\begin{figure}[H]
    \centering
    % Placeholder for Figure 7-1
    \includegraphics[page=9, viewport={179.5 34.3 364.0 198.3}, clip]{12-07572-001FigRevised.pdf}
    \caption{测定合金流动性的螺旋形试样}
    \label{fig:7-1}
\end{figure}

2. 影响合金流动性的因素

影响合金流动性的主要因素有合金的成分、物理性质、温度、不溶杂质和气体等。

(1) 合金的成分

合金的流动性主要取决于合金的成分。纯金属和共晶成分的合金在恒定温度下凝固,已凝固层和未凝固层之间界面分明、光滑,对未凝固液体的流动阻力小,因而流动性好。

具有宽凝固温度范围的合金凝固时,在铸件断面上既有发达的树枝晶,又有与未凝固液体合金相混杂的两相区。初生的树枝晶阻碍剩余液体合金的流动,因而合金的流动性差。合金的凝固温度范围越宽,其流动性也越差。图7-2为合金结晶特性对流动性的影响示意图。

\begin{figure}[H]
    \centering
    % Placeholder for Figure 7-2
    \includegraphics[page=10, viewport={89.5 675.3 409.5 746.3}, clip]{12-07572-001FigRevised.pdf}
    \caption{合金结晶特性对流动性的影响示意图}
    \label{fig:7-2}
\end{figure}

在相同过热度的条件下,铁碳合金的流动性与其碳的质量分数的关系如图7-3所示。从图可以看出,纯铁的流动性好,随碳的质量分数的增加,合金的凝固温度范围增大,流动性也随之下降。在亚共晶铸铁中,越靠近共晶成分,合金的凝固温度范围越小,其流动性越好;共晶成分铸铁在恒温下凝固,流动性最好。

因此,从合金流动性的角度考虑,在铸造生产中,应尽量选择共晶成分、近共晶成分或凝固温度范围小的合金作为铸造合金。

(2) 合金的物理性质

与合金流动性有关的物理性质有比热容、密度、导热系数、结晶潜热和黏度等。液态合金的比热容和密度越大、导热系数越小,凝固时结晶潜热释放得越多,都能使合金较长时间的保持液态,因而流动性越好;液态合金的黏度越小,流动时的内摩擦力也就越小,流动性自然越好。

\begin{figure}[H]
    \centering
    % Placeholder for Figure 7-3
    \includegraphics[page=10, viewport={157.0 476.8 387.0 643.3}, clip]{12-07572-001FigRevised.pdf}
    \caption{铁碳合金的流动性与其碳的质量分数的关系}
    \label{fig:7-3}
\end{figure}

(3) 液态合金的温度

在一定温度范围内,液态合金的流动性随其温度的升高而直线上升。但液态合金的温度过高,会造成液态合金的氧化、吸气非常严重,易使铸件产生气孔、夹渣、黏砂、缩松、缩孔等铸造缺陷。因此,液态合金的浇注温度必须合理。

3. 液态合金的充型能力及影响因素

液态合金充满铸型型腔,获得形状完整、轮廓清晰的铸件的能力,称为液态合金的充型能力。充型能力不足,铸件易形成浇不足、冷隔等缺陷。液态合金的充型能力主要取决于合金本身的流动性和各种工艺因素。对于流动性较差的合金,可通过改善工艺条件来提高其充型能力。影响液态合金充型能力的工艺因素主要有以下几个方面。

(1) 铸型条件

在液态金属充型时,凡是增加液态金属的流动阻力,降低其流动速度,提高其冷却能力的因素,均会降低液态合金的充型能力。

铸型条件包括铸型的蓄热能力、铸型的排气、铸型温度和铸件结构等方面。铸型的蓄热能力越大,铸型对液态合金的冷却能力越强,其充型能力越差,如液态合金在金属型中比砂型中的充型能力差;如果铸型的排气能力差,型腔中气体的压力增大,则阻碍液态金属充型;铸型温度越高,铸型对液态金属的冷却能力越差,因而提高了其充型能力;铸件结构越复杂,铸件壁厚越薄,液态金属充型越困难。

(2) 浇注条件

浇注条件包括浇注系统的结构、充型压力和浇注温度。浇注系统越复杂,液态合金流动的阻力越大,其充型能力越差;液态合金所受的静压力越大,其充型能力就越好,如在砂型铸造中,常用加高直浇道等工艺措施提高金属的静压力,而在压力铸造和低压铸造等特种铸造中,液态合金在压力下充型,能有效地提高其充型能力;浇注温度越高,合金的流动性越好。提高浇注温度能显著提高液态合金的充型能力。实际生产中提高液态合金的充型能力主要是通过提高浇注温度来实现的。但对铸件质量而言,并非浇注温度越高越好,应在保证充型能力的前提下,采用较低的浇注温度。对铸铁件,可采用“高温出炉、低温浇注”的方法,高温出炉能使铁液中一些难熔的固体质点熔化,铁液中的未熔质点和气体在浇包中的镇静阶段有机会上浮而除去。在保证铁液具有足够流动性的条件下,应选择尽可能低的浇注温度。

\subsection{铸造合金的收缩}

1. 收缩的概念

合金从液态冷却至室温的过程中,其体积或尺寸缩小的现象称为收缩。收缩是铸造合金本身的物理性质。它是铸件产生缩孔、缩松、热应力、变形及裂纹等铸造缺陷的基本原因。

液态金属注入铸型以后,从浇注温度冷却到室温要经历三个互相联系的收缩阶段。

(1) 液态收缩

液态收缩是指液态金属由浇注温度冷却到凝固开始温度(液相线温度)之间的收缩。在此阶段,金属处于液态,体积的缩小仅表现为型腔内液面的降低。

(2) 凝固收缩

凝固收缩是指从凝固开始温度到凝固终了温度(固相线温度)之间的收缩。合金的凝固温度范围越大,则凝固收缩越大。液态收缩和凝固收缩使金属液体积缩小,一般表现为型腔内液面降低,因此常用单位体积收缩量(即体收缩率)来表示。液态收缩和凝固收缩是缩孔和缩松形成的基本原因。

(3) 固态收缩

固态收缩是指合金从凝固终了温度冷却到室温之间的收缩,这是处于固态下的收缩。该阶段收缩不仅表现为合金体积的缩减,还直接表现为铸件的外形尺寸的减小,因此常用单位长度上的收缩量(即线收缩率)表示。固态收缩是产生铸件应力、变形和裂纹的基本原因。

铸件收缩的大小主要取决于合金成分、浇注温度、铸件结构和铸型条件。常用铸造合金中,灰铸铁的体收缩率约为$7\%$,线收缩率为$0.7\%\sim1.0\%$;铸钢的体收缩率为$12\%$,线收缩率为$1.5\%\sim2.0\%$。

2. 铸件中的缩孔和缩松

在铸件的凝固过程中,合金的液态收缩和凝固收缩使铸件的最后凝固部位出现孔洞,容积较大而集中的孔洞称为缩孔,细小而分散的孔洞称为缩松。铸件中不论存在何种形态的孔洞,都会减小铸件的有效受力面积,产生应力集中,使其承载能力和气密性等使用性能下降,因此缩孔和缩松是铸件的重要缺陷,必须设法防止。

(1) 缩孔

缩孔通常隐藏在铸件上部或最后凝固部位,经机械加工后可暴露出来。有时,缩孔产生在铸件的上表面上,呈明显凹坑。缩孔的外形特征近于倒锥形,内表面不光滑。

铸件缩孔的形成过程如图7-4所示。假定合金在恒温下凝固或凝固温度范围很窄,合金由表及里逐层凝固。液态金属填满铸型(图7-4a)以后,由于铸型的吸热及不断向外散热,靠近型腔表面的金属温度很快就降低到凝固温度,凝固成一层外壳(图7-4b)。温度继续下降,外壳不断加厚。同时,内部的剩余液体由于本身的液态收缩和补充凝固层的凝固收缩,其体积减小。液态收缩和凝固收缩造成的体积缩减逐渐积累,在重力的作用下液面与顶面相脱离(图7-4c)。如此进行下去,外壳不断加厚,液面不断下降。待合金完全凝固,就在铸件中形成了缩孔(图7-4d)。已经产生缩孔的铸件从凝固终了温度冷却到室温,因固态收缩使其外形尺寸略有缩小(图7-4e)。由此可知,铸件中的缩孔是由于合金的液态收缩和凝固收缩而产生的。

\begin{figure}[H]
    \centering
    % Placeholder for Figure 7-4
    \includegraphics[page=10, viewport={93.5 335.8 404.5 446.3}, clip]{12-07572-001FigRevised.pdf}
    \caption{铸件缩孔的形成过程}
    \label{fig:7-4}
\end{figure}

(2) 缩松

缩松多分布于铸件的轴线区域、内浇道附近甚至厚大铸件的整个断面,分布面广,难于控制,对铸件的力学性能影响很大,是铸件最危险的缺陷之一。

缩松的形成过程如图7-5所示。具有较宽凝固温度范围的合金,在铸件的断面上温度梯度又较小的条件下凝固时,液态金属最后在心部较宽的区域内同时凝固(图7-5a),初生的树枝晶把液体分隔成许多小的封闭区(图7-5b)。这些小封闭区液态金属的收缩得不到外界金属液的补充,便形成细小、分散的孔洞(图7-5c),即缩松。缩松形成的基本原因和缩孔一样,也是合金的液态收缩和凝固收缩所致。

\begin{figure}[H]
    \centering
    % Placeholder for Figure 7-5
    \includegraphics[page=10, viewport={39.5 185.8 264.5 286.8}, clip]{12-07572-001FigRevised.pdf}
    \caption{(圆柱形铸件)缩松的形成过程}
    \label{fig:7-5}
\end{figure}

由以上分析可得出以下结论。

1) 纯金属及共晶成分的合金在恒温下凝固,其铸件通常逐层由表向里凝固,倾向于形成集中缩孔;凝固温度范围宽的合金,其铸件通常在截面上较宽的区域内同时凝固,易形成缩松。缩孔比缩松易于检查和修补,也便于采取工艺措施来防止。因此,从收缩的角度考虑,在生产中应尽量选择共晶成分、近共晶成分或凝固温度范围小的合金作为铸造合金。

2) 对于给定成分的铸件,在一定的浇注条件下,缩孔和缩松的总容积是一定值。适当增大铸件的冷却速度可以促进缩松向缩孔转化。例如,在砂型铸造中,湿型比干型对铸件的激冷能力强,使铸件的凝固区域变窄,缩松量减少,而缩孔体积增加;在金属型铸造中,铸型的激冷能力更大,缩松的量显著减小。

3) 合金(如铸钢、白口铸铁、锡青铜等)的液态收缩和凝固收缩越大,铸件的缩孔体积越大。

4) 铸造合金的浇注温度越高,液态收缩越大,缩孔的体积也越大。

5) 缩孔和缩松总是存在于铸件的最后凝固部位。如果铸件壁厚不均匀,在厚壁处易出现缩孔或缩松。

3. 防止铸件产生缩孔的方法

虽然收缩是铸造合金的物理本性,但铸件中的缩孔并不是不可避免的。在进行铸造工艺设计时,只要采取一定的工艺措施,就能有效地防止缩孔的产生。

在实际生产中,通常采用顺序凝固原则,并设法使分散的缩松转化为集中的缩孔,再使集中的缩孔转移到冒口中(图7-6),最后将冒口去除,即可获得完好的铸件。通过设置冒口和冷铁,铸件从远离冒口的地方开始凝固并逐渐向冒口推进,冒口最后凝固。

为了实现铸件的顺序凝固,应合理地选择内浇道在铸件上的引入位置,开设冒口并放置冷铁(图7-7)。生产上为提高液态金属的补缩效果,常采用发热保温冒口。

\begin{figure}[H]
    \centering
    % Placeholder for Figure 7-6
    \includegraphics[page=10, viewport={330.0 183.3 502.5 305.8}, clip]{12-07572-001FigRevised.pdf}
    \caption{冒口补缩示意图}
    \label{fig:7-6}
\end{figure}

\begin{figure}[H]
    \centering
    % Placeholder for Figure 7-7
    \includegraphics[page=10, viewport={159.5 20.3 380.5 151.3}, clip]{12-07572-001FigRevised.pdf}
    \caption{铸件顺序凝固示意图}
    \label{fig:7-7}
\end{figure}

4. 缩孔位置的确定

准确地估计铸件上缩孔可能产生的位置是合理安置冒口和冷铁的前提。在生产中,确定缩孔位置的常用方法有凝固等温线法、内切圆法和计算机凝固模拟法等。

计算机凝固模拟法是在计算机上利用凝固模拟软件对实际铸件的凝固过程进行模拟计算以确定缩孔位置的方法。该方法可以较为准确地预测铸件最后凝固的部位和可能产生缩孔的位置,并能优化铸造工艺,提高铸件的工艺出品率。目前,该方法在实际生产中已得到应用。

\subsection{铸造内应力及铸件的变形、裂纹}

铸件凝固后在冷却至室温的过程中还会继续收缩,有些合金甚至会发生固态相变而引起收缩或膨胀,这些收缩或膨胀如果受到阻碍或因铸件各部分互相牵制,都将使铸件内部产生应力。内应力是铸件产生变形及裂纹的主要原因。

1. 热应力

铸件在凝固和其后的冷却过程中,因壁厚不均,各部分冷却速度不同,便会造成同一时刻各部分收缩量不同,因此在铸件内彼此相互制约而产生应力。

金属在冷却过程中,从凝固终了温度到再结晶温度阶段处于塑性状态,此时在较小的外力下就会产生塑性变形,变形后应力可自行消除。低于再结晶温度的金属处于弹性状态,受力时产生弹性变形,变形后应力仍然存在。

下面用如图7-8所示的框形铸件来分析热应力的产生过程。I杆比II杆的直径大。凝固开始时,I、II两杆均处于塑性状态,冷却速度虽不同,但不产生应力。继续冷却,冷却速度大的II杆已进入弹性状态,而I杆仍处于塑性状态,此时,因II杆冷却快,收缩大于I杆,必然压缩I杆,所以II杆受拉,I杆受压(图7-8b)。I杆在应力作用下,发生微量塑性变形而被压短,内应力消失(图7-8c)。进一步冷却,I杆处于弹性状态,产生较大的固态收缩,此时II杆处于更低温度,其收缩已很小或收缩已趋停止,将阻碍I杆的收缩。结果,I杆受拉伸,II杆受压缩(图7-8d),直到室温,形成了内应力。当I杆所受的内拉应力超过强度极限,就会被拉断(图7-8e)。

\begin{figure}[H]
    \centering
    % Placeholder for Figure 7-8
    \includegraphics[page=11, viewport={47.5 653.8 450.0 739.3}, clip]{12-07572-001FigRevised.pdf}
    \caption{铸造热应力与变形}
    \label{fig:7-8}
\end{figure}

综上所述,固态收缩使铸件厚壁或心部受拉伸,薄壁或表层受压缩。合金固态收缩率越大,铸件壁厚差别越大,形状越复杂,所产生的热应力越大。

目前,对于铸件残余应力,不仅能进行定性分析(分析其应力状态),还能利用有限元法或有限差分法对其进行定量模拟计算,求得铸件不同温度下的应力场。

2. 机械应力

铸件收缩受到铸型、型芯及浇注系统的机械阻碍而产生的应力称为机械阻碍应力,简称为机械应力。

铸型或型芯退让性良好,则机械应力小。机械应力在铸件落砂之后可自行消除。但是,机械应力在铸型中能与热应力共同起作用,增加了铸件产生裂纹的可能性。

铸造内应力的存在,将引起铸件变形和裂纹等缺陷。

3. 铸件的变形及其防止

如果铸件存在内应力,则处于一种不稳定状态。铸件厚的部分受拉应力,薄的部分受压应力。如果内应力超过合金的屈服强度,则铸件本身总是力图通过变形来减缓内应力,因此细而长或大又薄的铸件易发生变形。

如图7-9所示,车床床身的导轨部分因较厚而受拉应力,床壁部分因较薄而受压应力,于是导轨会产生向下的挠曲变形。

\begin{figure}[H]
    \centering
    % Placeholder for Figure 7-9
    \includegraphics[page=11, viewport={37.5 425.8 315.0 542.3}, clip]{12-07572-001FigRevised.pdf}
    \caption{车床床身导轨的挠曲变形}
    \label{fig:7-9}
\end{figure}

为了防止铸件变形,设计时应使铸件各部分壁厚尽可能均匀或形状对称。在铸造工艺上可采取同时凝固原则。所谓同时凝固原则,就是采取工艺措施保证铸件结构上各部分之间没有温差或温差尽量小,使各部分能同时凝固,如图7-10所示。采用该原则,可使铸件内应力较小,不易产生变形和裂纹。但在铸件中心区域往往有缩松,组织不够致密。同时凝固原则主要用于凝固收缩小的合金(如灰铸铁)以及壁厚均匀、结晶温度范围宽而对铸件的致密性要求不高的铸件等。此外,还可在制模时采用反变形法(将模样制成与铸件变形方向相反的形状,此时应较精确地计算铸件的变形量);有时也在薄壁处附加工艺筋。

\begin{figure}[H]
	\centering
	% Placeholder for Figure 7-10
	\includegraphics[page=11, viewport={341.0 425.8 502.0 628.3}, clip]{12-07572-001FigRevised.pdf}
	\caption{车床床身导轨的挠曲变形}
	\label{fig:7-10}
\end{figure}

实践证明,尽管铸件冷却时发生部分变形,但内应力仍未彻底消除。经过机械加工后内应力会重新分布,铸件仍会发生变形,影响零件的精度。因此,对某些重要、精密的铸件,如车床床身等,必须采取去应力退火或自然时效等方法,将残余应力消除,必要时还可在粗加工后先进行去应力退火或人工时效,然后再进行精加工,以确保零件的精度。

4. 铸件的裂纹及防止

如果铸造内应力超过合金的强度极限,则会产生裂纹。裂纹分为热裂和冷裂两种。

(1) 热裂

热裂是在凝固后期高温下形成的,主要是由于收缩受到机械阻碍作用而产生的。它具有裂纹短、形状曲折、缝隙宽、断面严重氧化、无金属光泽、裂纹沿晶界产生和发展等特征,常见于铸钢和铝合金铸件中。

防止热裂的主要措施除合理设计铸件结构外,还应合理选用型砂或芯砂的黏结剂,以改善其退让性;大的型芯可采用中空结构或内部填以焦炭;严格限制铸钢和铸铁中硫的质量分数;选用收缩率小的合金等。

(2) 冷裂

冷裂是在较低温度下形成的,常出现在铸件受拉伸部位,特别是有应力集中的地方。其裂缝细小,呈连续直线状,缝内干净,有时呈轻微氧化色。

壁厚差别大,形状复杂或大而薄的铸件易产生冷裂。因此,凡是能减少铸造内应力或降低合金脆性的因素,都能防止冷裂的形成。同时,在铸钢和铸铁中要严格控制合金中磷的质量分数。

\section{砂型铸造}

\subsection{砂型铸造的生产过程及特点}

砂型铸造是目前应用最广泛的铸造方法,其生产过程如图7-11所示。首先根据零件图设计出铸件图及模样图,制出模样及其他工装设备,并用模样、砂箱等和配制好的型砂制成相应的砂型,然后把熔炼好的合金液浇入型腔。待合金液在型腔内凝固冷却后,把砂型破坏,取出铸件。最后清除铸件上附着的型砂及浇冒口系统,经过检验,便可获得所需要的铸件。砂型铸造的生产过程是周期性循环的。

\begin{figure}[H]
    \centering
    % Placeholder for Figure 7-11
    \includegraphics[page=11, viewport={62.0 241.8 476.5 395.8}, clip]{12-07572-001FigRevised.pdf}
    \caption{砂型铸造的生产过程}
\end{figure}

在单件、小批或生产批量不大(但有一定的年产量)的现代铸造车间中,通常采用机械化的型砂处理及输送系统,而造型则采用手工造型、机器造型或手工结合机器造型,铸型、合金液及铸件的搬运、浇注则采用起吊设备来完成,生产效率较低。

在大批量生产的机械化铸造车间中,生产过程是按流水作业连续进行的。型砂的处理及输送、造型、制芯、合箱(合型)、浇注、落砂清理以及砂箱、铸型、合金液及铸件的输送等绝大部分工作都是由机器来完成的。机械化铸造车间的造型-浇注-落砂生产线是由铸型输送机或传输轨道将造型机和铸造辅助设备(如翻箱机、合箱机、压铁机、分箱机、落砂机、捅箱机、台面清扫机、浇注机、过渡小车、冷却装置等)有机地结合在一起而组成的(图7-12)。整个铸造过程协调、连续地循环进行,生产率显著提高。需要指出的是,在机器造型流水线上,不能生产厚壁和大型铸件,且只能采用模板进行两箱造型,因此铸件外形受到一定的限制。又因砂型铸造的工艺繁琐,许多操作(如安装型芯)仍然离不开手工操作,所以砂型铸造的机械化程度至今仍受到一定的限制。

\begin{figure}[H]
    \centering
    % Placeholder for Figure 7-12
    \includegraphics[page=12, viewport={60.0 521.8 475.5 746.3}, clip]{12-07572-001FigRevised.pdf}
    \caption{砂型铸造造型线}
\end{figure}

砂型铸造的缺点是一个砂型只能使用一次,要耗费大量的造型工时。而且,每生产1t合格的铸件,需使用4\textasciitilde5t型砂,处理大量型砂的工作十分繁重。此外,车间中有砂尘污染,劳动条件较差。但砂型铸造适合于在各种生产条件下生产各种合金的铸件,所以仍是当今普遍应用的铸造方法。

\subsection{砂型铸造工艺简介}

砂型铸造工艺包括造型、熔炼与浇注、落砂与清理等工序。

造型是指利用型砂、模样与砂箱等原材料及工艺装备制造铸型的工艺过程。该过程通常分为手工造型和机器造型两大类。

手工造型主要包括整模、分模、挖砂、活块、三箱、刮板、假箱造型以及地坑造型等方法。这类造型方式操作灵活、适应性强、生产准备周期短,但生产率低、劳动强度大,铸件质量较差,因此主要用于单件或小批量生产。

机器造型是指用机器完成全部或至少完成紧砂和起模操作的造型工序。机器辅助紧砂方法主要包括:振压紧实,适用于中、小型铸件;抛砂紧实,适用于大、中型铸件;射砂紧实,主要用于制芯,亦可用于造型。机器辅助起模方法通常有顶箱起模(7-13a)、漏模起模(7-13b)和翻转起模(7-13c)。顶箱起模适用于结构简单、型腔较浅的铸型;漏模起模适用于结构相对简单但型腔较深的铸型;翻转起模则用于形状复杂、型腔较深的铸型。机器造型可大大提高生产率和铸件的尺寸精度,降低表面粗糙度值,减少加工余量,并改善工人的劳动条件,正日益广泛应用于铸件的大批量生产中。

\begin{figure}[H]
    \centering
    % Placeholder for Figure 7-13
    \includegraphics[page=12, viewport={15.5 218.3 344.0 496.3}, clip]{12-07572-001FigRevised.pdf}
    \caption{机器辅助起模方法}
\end{figure}

型芯主要用于形成铸件的内腔及尺寸较大的孔。最常用的制芯方法是用芯盒制芯。为防止铸件产生黏砂、夹砂及砂眼等缺陷,提高铸件表面质量,通常在砂型和型芯表面涂刷铸造涂料。

熔炼提供化学成分和温度都合格的液态金属。浇注是将液态金属从浇包注入铸型的操作。浇注时,浇注温度应尽可能低些,以减少气体的溶解量及液态收缩量,从而减少气孔、缩孔等铸件缺陷。但液态金属出炉的温度应尽可能高些,以利于熔渣上浮,从而便于清渣和减少夹杂物类铸件缺陷。

落砂是用手工或机械使铸件与型砂、砂箱分开的操作。落砂时间过早,可能导致灰铸铁铸件表层产生白口铸铁组织,难以进行切削加工;落砂时间过晚,则可能由于收缩应力使铸件产生裂纹。因此,浇注后应及时落砂。

清理是落砂后清除铸件表面的黏砂、型砂以及多余金属(包括浇冒口、氧化皮)的过程。落砂后应及时清理铸件。清理后的铸件应根据其技术要求仔细检验,判断铸件是否合格。技术条件允许焊补的铸件缺陷应焊补。合格的铸件应进行去应力退火或自然时效。变形的铸件应进行矫正。

\subsection{铸造工艺图}

铸造必须首先根据零件的结构特点、技术要求、生产批量和生产条件等进行铸造工艺设计,然后绘制铸造工艺图。

1. 浇注位置的选择

铸件的浇注位置是指浇注时铸件在铸型内所处的空间位置。它对铸件质量、造型方法、砂箱尺寸、加工余量等都有很大影响。选择浇注位置时应以保证铸件质量为主,一般应遵循以下几个原则。

1) 应将铸件上质量要求高的表面或主要加工面放在铸型的下面。如果做不到这点,也应将该表面置于铸型的侧面或倾斜位置进行浇注。图7-14表示锥齿轮的两种不同的浇注位置,应将要求较高并需机械加工的轮齿置于铸型的下面。图7-15所示为卷扬筒的浇注位置,应将铸件的主要加工面放在铸型侧面,而将次要且面积较小的凸缘放在上面。

\begin{figure}[H]
    \centering
    % Placeholder for Figure 7-14
    \includegraphics[page=12, viewport={364.5 219.8 533.5 337.3}, clip]{12-07572-001FigRevised.pdf}
    \caption{锥齿轮的浇注位置}
\end{figure}

\begin{figure}[H]
    \centering
    % Placeholder for Figure 7-15
    \includegraphics[page=12, viewport={11.5 20.3 113.5 202.3}, clip]{12-07572-001FigRevised.pdf}
    \caption{卷扬筒的浇注位置}
\end{figure}

2) 对于一些需要补缩的铸件,应把截面较厚的部分放在铸型的上部或侧面。这样便于在铸件的厚壁处放置冒口,造成合理的凝固顺序,有利于铸件补缩。

3) 对于具有大面积的薄壁铸件,浇注时应使薄壁部分位于铸型下部,同时尽量使薄壁立放或倾斜浇注,这样有利于金属的充填。图7-16所示为箱盖的两种浇注位置。其中,图7-16b是合理的,它将铸件大面积的薄壁部分放在铸型的下面,使这部分能在较高的金属液压力下充满铸型,防止浇不足或冷隔等缺陷产生。

\begin{figure}[H]
    \centering
    % Placeholder for Figure 7-16
    \includegraphics[page=12, viewport={127.5 22.8 338.5 167.3}, clip]{12-07572-001FigRevised.pdf}
    \caption{箱盖的浇注位置}
\end{figure}

4) 对于具有大平面的铸件,应将铸件的大平面放在铸型的下面。例如,在浇注带有筋条的平板时,应选如图7-17所示的浇注位置,这样可使铸件的大平面不容易产生夹砂等缺陷。

\begin{figure}[H]
    \centering
    % Placeholder for Figure 7-17
    \includegraphics[page=12, viewport={352.5 22.3 532.0 154.3}, clip]{12-07572-001FigRevised.pdf}
    \caption{平板铸件的浇注位置}
\end{figure}

2. 铸型分型面的选择

分型面是指两半铸型相互接触的表面。分型面的选择合理与否,对铸件质量及制模、造型、制芯、合箱或清理等工序的复杂程度有很大影响。在选择铸型分型面时应遵循如下原则。

1) 分型面应选在铸件的最大截面上,并力求采用平面。这样可使模样顺利取出,简化造型工艺,不用或少用挖砂造型或假箱造型,如图7-18所示。

\begin{figure}[H]
    \centering
    % Placeholder for Figure 7-18
    \includegraphics[page=13, viewport={136.5 653.8 406.5 740.3}, clip]{12-07572-001FigRevised.pdf}
    \caption{分型面应选在铸件最大截面上}
\end{figure}

2) 应尽量减少分型面的数量,并尽量做到只有一个分型面,以便采用工艺简便的两箱造型或采用机器造型,避免三箱造型。有时可用型芯来减少分型面,如图7-19所示。

3) 应尽可能减少活块和型芯的数量,注意减少砂箱高度。这样可简化制模及造型工艺,便于起模和修型。

4) 尽量把铸件的大部分或全部放在一个砂箱内,并使铸件的重要加工面、工作面、加工基准面及主要型芯位于下型内。这样便于型芯的安放和检验,还可使上型的高度降低,便于合箱,并可保证铸件的尺寸精度,防止错箱。图7-20所示为管子堵头铸件分型面的选择,其中图7-20a是正确的,它将铸件全部放在下型,避免错箱,保证铸件质量。

\begin{figure}[H]
    \centering
    % Placeholder for Figure 7-19
    \includegraphics[page=13, viewport={69.0 540.3 257.0 619.3}, clip]{12-07572-001FigRevised.pdf}
    \caption{用型芯减少绳轮铸件分型面}
\end{figure}

\begin{figure}[H]
    \centering
    % Placeholder for Figure 7-20
    \includegraphics[page=13, viewport={268.0 538.8 471.5 626.8}, clip]{12-07572-001FigRevised.pdf}
    \caption{管子堵头铸件分型面的选择}
\end{figure}

3. 主要工艺参数的确定

铸造生产的工艺方案确定以后,还应根据产品的形状、尺寸和技术要求,选定好各种铸造工艺参数,以保证铸件的形状和尺寸等符合要求。

铸造工艺参数是由金属种类和铸造方法等特点决定的。其内容包括铸造收缩率、机械加工余量、起模斜度、铸造圆角和芯头尺寸等,有时还要考虑工艺补正量。

(1) 铸造收缩率

由于合金的收缩,铸件的实际尺寸要比模样的尺寸小。为确保铸件的尺寸,必须根据合金收缩率放大模样尺寸。合金的收缩率受多种因素的影响,不同成分合金的收缩率相差较大。通常灰铸铁的收缩率为$0.7\% \sim 1.0\%$,铸钢的收缩率为$1.6\% \sim 2.0\%$,非铁合金为$1.0\% \sim 1.5\%$。

(2) 机械加工余量

机械加工余量是指在铸件加工表面上留出的、准备切去的金属层厚度。机械加工余量取决于铸件生产批量、合金的种类、铸件的大小、加工面与基准面的距离以及加工面在浇注时的位置等。机器造型铸件精度高,余量小;手工造型误差大,余量应加大。灰铸铁件表面平整,加工余量小;铸钢件表面粗糙,加工余量应加大。铸铁件的机械加工余量通常取在$3 \sim 15\text{ mm}$之间,具体可参阅GB/T 42124.3—2025《产品几何技术规范(GPS) 模制件的尺寸和几何公差 第3部分:铸件尺寸公差、几何公差与机械加工余量》。

(3) 起模斜度

起模斜度是指为方便起模和避免损坏砂型,在模样、芯盒的出模方向留有的斜度。起模斜度应留在铸件垂直于分型面要加工的表面上,其大小取决于立壁的高度、造型方法、模样材料等因素,通常为$15' \sim 3^{\circ}$。如图7-21所示,图中$\alpha, \beta, \gamma$表示不同侧面的起模斜度,$a, b, c$表示铸件不同表面上的机械加工余量。

\begin{figure}[H]
    \centering
    % Placeholder for Figure 7-21
    \includegraphics[page=13, viewport={15.5 363.3 169.0 493.8}, clip]{12-07572-001FigRevised.pdf}
    \caption{起模斜度与机械加工余量}
\end{figure}

(4) 铸造圆角

为便于造型,避免铸件在尖角处产生裂纹和应力集中以及在浇注时造成冲砂、砂眼和黏砂等缺陷,提高铸件强度,应将模样上的尖角做成圆角。因此,铸件最明显的特征结构就是铸造圆角。

铸件的内圆角半径可取铸件转角处两壁平均壁厚的$1/4$左右,外圆角半径比内圆角半径大转角处两壁壁厚之和的$\frac{1}{3} \sim \frac{1}{2}$,如图7-22所示。

\begin{figure}[H]
    \centering
    % Placeholder for Figure 7-22
    \includegraphics[page=13, viewport={189.5 367.3 296.0 450.3}, clip]{12-07572-001FigRevised.pdf}
    \caption{铸造圆角}
\end{figure}

(5) 芯头尺寸

芯头是指伸出铸件以外不与金属接触的型芯部分。它主要用于定位、支承和固定型芯,使型芯在铸型中有准确的位置。要使型芯工作可靠,就须使芯头有合适的尺寸,如芯头长度、芯头斜度和芯头装配间隙等,如图7-23所示。

\begin{figure}[H]
    \centering
    % Placeholder for Figure 7-23
    \includegraphics[page=13, viewport={322.0 361.3 531.0 507.8}, clip]{12-07572-001FigRevised.pdf}
    \caption{芯头}
\end{figure}

(6) 最小铸出孔及槽

零件上的孔、槽、台阶等是否铸出来,应从铸件质量及节约方面综合考虑。一般来说,较大的孔、槽等应铸出来,以便节约金属和加工工时,同时还可避免铸件的局部过厚所造成的热节,提高铸件质量。较小的孔、槽或者铸件壁很厚,则不宜铸孔,直接进行加工反而更方便。有些特殊要求的孔,如弯曲孔,无法实行机械加工,则必须铸出。表7-1为铸件的最小铸出孔的孔径。

\begin{table}[H]
    \centering
    \caption{铸件的最小铸出孔的孔径}
    \begin{tabular}{|c|c|c|}
    \hline
    \multirow{2}{*}{生产批量} & \multicolumn{2}{c|}{最小铸出孔直径/mm} \\ \cline{2-3} 
     & 灰铸铁件 & 铸钢件 \\ \hline
    大量生产 & 12\textasciitilde15 & --- \\ \hline
    成批生产 & 15\textasciitilde30 & 30\textasciitilde50 \\ \hline
    单件、小批生产 & 30\textasciitilde50 & 50 \\ \hline
    \end{tabular}
\end{table}

4. 铸造工艺图的绘制

铸造工艺图是在零件图上用各种工艺符号表示出铸造工艺方案的图形,是在对铸件进行工艺分析的基础上,确定铸件的浇注位置,分型面,型芯的数量、形状及其固定方法,加工余量,起模斜度,收缩率,浇注系统,冒口,冷铁的尺寸及布置,砂箱的形状及尺寸等。铸造工艺图是指导模样(芯盒)设计、生产准备、铸型制造和铸件检验的基本技术文件。完整的铸造工艺图一般包括铸件(毛坯)图、模样(芯盒)图和铸型装配图(砂型合箱图)。支座的铸造工艺图如图7-24所示。

\begin{figure}[H]
    \centering
    % Placeholder for Figure 7-24
    \includegraphics[page=13, viewport={94.5 167.8 450.0 331.3}, clip]{12-07572-001FigRevised.pdf}
    \caption{支座的铸造工艺图}
\end{figure}

\section{特种铸造}

砂型铸造具有适应性强、生产准备简单等优点,被广泛用于制造业,但是砂型铸造生产的铸件的尺寸精度较低,表面粗糙、内在质量较差,生产过程较复杂,不易实现机械化,工人的劳动条件差、劳动强度大。为克服砂型铸造的这些缺点,人们在砂型铸造的基础上,通过改变铸型材料(如金属型、磁型、陶瓷型铸造)、模样材料(如熔模铸造、实型铸造)、浇注方法(如离心铸造、压力铸造)、金属液充填铸型的形式或铸件凝固的条件(如压力铸造、低压铸造、真空吸铸)等,又创造了许多其他的铸造方法。通常,把这些不同于普通砂型铸造的铸造方法统称为特种铸造。常用的特种铸造方法有熔模铸造、金属型铸造、压力铸造、离心铸造、低压铸造、陶瓷型铸造等。近年来,特种铸造在我国发展非常迅速,尤其在非铁金属的铸造生产中占有重要的地位。

\subsection{熔模铸造}

熔模铸造又称为精密铸造,是用蜡料制成模样,再在蜡模表面涂覆多层耐火材料,待硬化干燥后,将蜡模熔去,而获得具有与蜡模形状相应空腔的型壳,再经焙烧后进行浇注而获得铸件的一种方法。

1. 熔模铸造的工艺过程

熔模铸造包括压制蜡模、结壳、脱蜡、焙烧、造型与浇注等工艺过程。

(1) 压制蜡模

制造蜡模的材料有石蜡、蜂蜡、硬脂酸和松香等,最常用的是质量分数各为50\%的石蜡和硬脂酸的混合料。压制时,将蜡料加热至糊状后,在0.2~0.3 MPa的压力下将蜡料压入到压型内,待蜡料冷却凝固后便可从压型内取出,修除分型面上的毛刺,即可得到单个蜡模。为了提高生产效率,常需将单个蜡模黏焊在预制好的蜡质直浇道模样上,制成蜡模组。压制蜡模过程如图7-25a~f所示。

(2) 结壳

先将蜡模浸挂一层用水玻璃或硅溶胶和石英粉配成的涂料,再向其表面撒一层硅砂,并硬化涂层。如此反复3~7次,直至结成5~10 mm的硬壳为止,如图7-25g所示。

(3) 脱蜡

将结壳后的型壳放入85~95 $^\circ$C的水中(或放在高压釜中,通入0.2~0.5 MPa压力的水蒸气),使蜡模熔化而脱出,型壳则形成了铸型空腔,如图7-25g所示。

(4) 焙烧

将铸型在850~950 $^\circ$C的温度下焙烧,使其中所含的残余挥发物进一步排出。

(5) 造型与浇注

浇注时将型壳置于铁箱中,周围用干砂填紧,以避免型壳开裂。为提高液态合金的填充能力,防止浇不足缺陷,常在焙烧后趁热(600~700 $^\circ$C)进行浇注,如图7-25h所示。


\begin{figure}[H]
\centering
\includegraphics[page=14, viewport={112.0 586.3 427.5 742.3}, clip]{12-07572-001FigRevised.pdf}
\caption{熔模铸造的工艺过程}
\end{figure}

2. 熔模铸造的特点及应用范围

熔模铸造件的精度及表面质量高(尺寸公差等级为IT14~IT11,表面粗糙度$Ra$值可达12.5~1.6 $\mu$m);能够铸造各种合金铸件;生产批量不受限制,从单件、小批到大批量生产均可;熔模铸件的形状可比较复杂,铸件上可铸出的最小孔径为0.5 mm,铸件的最小壁厚为0.3 mm。铸件的质量不宜太大,一般不超过25 kg,目前生产的最大熔模铸件为80 kg左右。
熔模铸造工艺过程较复杂,且不易控制,使用和消耗的材料较贵,因而适用于生产形状复杂、精度要求较高,或难以进行机械加工的小型零件,如涡轮发动机叶片和叶轮、高速钢切削刀具等。

\subsection{金属型铸造}

将金属液浇入用金属制成的铸型型腔中,以获得铸件的方法称为金属型铸造。金属型可反复多次使用,大大提高了生产效率。

1. 金属型分类及其特点

根据分型面的位置不同,金属型的结构可分为整体式、竖直分型式、水平分型式和复合分型式。其中,竖直分型式和复合分型式(图7-26)使用方便,应用最广。

用金属型代替砂型,克服了砂型的许多缺点,但也会带来一系列新问题。如金属型导热快,没有退让性,所以铸件易产生冷隔、浇不足、裂纹等缺陷,灰铸铁件常产生白口组织。为了确保铸件质量和延长铸型寿命,金属型铸造必须采取下列措施。

1) 浇注前预热金属型,以防止浇不足、冷隔、裂纹等缺陷。金属型铸造合理的工作温度一般为:非铁金属铸件为100~250 $^\circ$C,铸铁件为250~350 $^\circ$C。

2) 加强金属型的排气。

3) 金属型的型腔应喷刷涂料,使金属液和铸型隔开,以延长金属型的寿命和获得表面光洁的铸件。

4) 应尽早开型取出铸件。

5) 合理设计铸件壁厚。为防止铸铁产生白口铸铁组织,其壁厚不宜过薄(一般应大于15 mm),并控制铁液合适的碳、硅总的质量分数或采用孕育处理。


\begin{figure}[H]
\centering
\includegraphics[page=14, viewport={11.5 410.3 232.0 558.8}, clip]{12-07572-001FigRevised.pdf}
\caption{垂直分型式和复合分型式金属型}
\end{figure}

2. 金属型铸造的特点及应用范围

与砂型铸造比较,金属型铸造有如下特点。

1) 实现了“一型多铸”,节约了大量工时和型砂,提高了劳动生产率,改善了劳动条件。

2) 铸件的力学性能高,如铝合金金属型铸件比砂型铸件的抗拉强度可平均提高20\%,同时耐腐蚀性和硬度也显著提高。

3) 铸件的精度较高(尺寸公差等级IT16~IT12,表面粗糙度$Ra$值为12.5~6.3 $\mu$m),可少加工或不加工,提高了金属材料的利用率,减少了机械加工费用。

4) 金属型的制造成本高、周期长,铸型透气性差、无退让性,易使铸件产生冷隔、浇不足、裂纹等铸造缺陷;受铸型的限制,金属型铸件合金的熔点不宜太高,重量也不易太大;金属型铸造必须采用机械化或自动化装置,否则劳动条件反而更加恶劣,故金属型铸造的应用范围受到了较大限制。
金属型铸造主要适用于大批量生产非铁金属铸件,如发动机零件,汽车、拖拉机、内燃机、摩托车的铝活塞、气缸体、缸盖,油泵壳体以及铜合金轴瓦等。金属型铸造有时也可用来制造形状较简单的小型可锻铸铁件或铸钢件。

\subsection{压力铸造}

在高压作用下,液态或半固态金属以较高的速度填充铸型型腔,并在压力作用下凝固而获得铸件的方法称为压力铸造。高压和高速充型是压力铸造的两大特点。

1. 压力铸造的工艺过程

压力铸造在压铸机上进行。压铸机按压射部分的特征分为热压室式和冷压室式两大类。热压室式压铸机上装有贮存液态金属的坩埚,压室浸在液态金属中,因此只能压铸低熔点合金,应用较少。目前广泛应用的是冷压室式压铸机,金属的熔炼设备不在压铸机上。卧式冷压室式压铸机的工作过程如图7-27所示。


\begin{figure}[H]
\centering
\includegraphics[page=14, viewport={240.5 408.8 490.0 531.8}, clip]{12-07572-001FigRevised.pdf}
\caption{卧式冷压室式压铸机的工作过程}
\end{figure}

压力铸造所用的铸型称为压铸型,用耐热钢制成。压铸型与竖直分型式金属型相似,一半固定在压铸机上,称为定型;另一半可水平移动,称为动型。压力铸造时,首先动型和定型以很大的合型力(一般为500~15 000 MPa)合型,随后将液态金属注入压室(图7-27a),压射活塞向前推进,将金属液高速压入型腔,并使金属液在高压下凝固(图7-27b)。开型后,用顶杆将铸件与余料一起顶出(图7-27c),完成压力铸造。

2. 压力铸造的特点及应用范围

与砂型铸造相比较,压力铸造有如下优点。

1) 铸件的尺寸精度高,一般尺寸公差等级可达到IT13~IT11,表面粗糙度$Ra$值可达3.2~0.8 $\mu$m,有时达0.4 $\mu$m,一般可不经机械加工直接使用。

2) 铸件的强度和表面硬度高。因为液态金属在压力下结晶,冷却速度又较快,所以压铸件的组织致密、晶粒较细,其抗拉强度可比砂型铸件提高25\%~30\%,但伸长率有所下降。

3) 可压铸形状复杂的薄壁铸件,如铝合金压铸件的最小壁厚可为0.5 mm,最小铸出孔直径可为0.7 mm。

4) 压铸件中可嵌铸其他材料(如钢、铁、铜合金、钻石等)的零件,以节省贵重材料和机械加工工时。有时嵌铸还可以代替部件的装配过程。

5) 生产效率高,一般班产600~700件,是所有铸造方法中生产效率最高的方法。

压力铸造虽然是实现金属零件少、无切屑加工的有效方法,但也存在着如下不足之处。

1) 设备投资大,制作压铸型的成本高。

2) 压铸高熔点合金(如钢、铸铁等)时,压铸型的寿命低,因而限制了压力铸造的应用范围。

3) 由于液态金属高速充型,液流会包裹住大量空气,最后以气孔的形式留在压铸件中。因此,压铸件不能进行多余量的机械加工,以免气孔暴露,削弱铸件的使用性能。有气孔的压铸件也不能进行热处理,因为在高温时,气孔内气体膨胀会使铸件表面鼓泡。

压力铸造是目前应用较广泛的一种铸造方法,主要适用于低熔点的锌、铝、镁及铜等非铁金属中小型铸件的大批量生产,如用来生产发动机气缸体、气缸盖、变速箱体、发动机罩、仪表和照相机壳体及支架、管接头等。

\subsection{离心铸造}

将液态金属浇入高速旋转的铸型中,使其在离心力的作用下填充铸型并凝固成形的铸造方法称为离心铸造。

离心铸造的铸型有金属型和砂型两种。目前,广泛应用的是金属型离心铸造。

离心铸造在离心铸造机上进行。根据铸型旋转轴在空间的位置,离心铸造机分为立式离心铸造机和卧式离心铸造机两类,绕竖直轴旋转的为立式离心铸造机,绕水平轴旋转的为卧式离心铸造机,离心铸造过程如图7-28所示。


\begin{figure}[H]
\centering
\includegraphics[page=14, viewport={127.5 248.3 413.0 378.8}, clip]{12-07572-001FigRevised.pdf}
\caption{离心铸造过程示意图}
\end{figure}

与砂型铸造相比较,离心铸造有如下特点。

1) 工艺过程简单。铸造中空筒类、管类零件时,省去了型芯、浇注系统和冒口,节约金属和其他原材料。

2) 离心铸造使液态金属在离心力作用下充型并凝固,其中密度较小的气体、夹渣等均集中于铸件内表面,而金属则从外向内呈方向性结晶,因而铸件组织致密,无缩孔、气孔、夹渣等缺陷,力学性能较好。

3) 便于铸造“双金属”铸件,如制造钢套内壁挂衬巴氏合金的滑动轴承,既可达到滑动轴承的使用要求。又可节约较贵的滑动轴承合金材料。

在离心铸造中,铸造合金的种类几乎不受限制。目前,已有高度机械化、自动化的离心铸造机,有年产量达十万吨的机械化离心铸管厂。

离心铸造的缺点是铸件的内表面质量差,孔的尺寸不易控制,但这并不妨碍其作为一般管道的使用要求。对于内孔待加工的机器零件,则可采用加大内孔加工余量的方法来解决。

目前,离心铸造已广泛用于大批量生产灰铸铁及球墨铸铁管、气缸套及滑动轴承等中空件,也可采用熔模离心铸造浇注刀具、齿轮等成形铸件。

\subsection{其他铸造方法简介}

1. 反压铸造

反压铸造又称为压差铸造,它是使液体金属在压差作用下,充填到预先有一定压力的铸型型腔内,进行结晶、凝固而获得铸件的一种工艺方法。

反压铸造按压差产生的方式不同,可分为增压法和减压法两种,其铸造工艺原理如图7-29所示。增压法是在型腔中预先建立一定压力$p_1$的基础上,向坩埚内充入气体使其压力升至$p_2$,借助正压差推动金属液沿升液管平稳上升并充满型腔,随后维持压力,完成凝固。减压法则先使坩埚压力与型腔压力相等(均为$p_3$),再通过排气降低型腔压力至$p_4$,利用负压差将金属液“吸入”型腔,从而实现无冲击、低卷气的充型过程。这两种方法均有利于获得组织致密、表面光洁、壁厚均匀的高质量铸件。


\begin{figure}[H]
\centering
\includegraphics[page=14, viewport={95.5 64.3 455.0 220.3}, clip]{12-07572-001FigRevised.pdf}
\caption{反压铸造工艺原理}
\end{figure}

反压铸造虽然和真空吸铸一样,金属液是在压差作用下沿升液管上升充型的,但反压铸造在充型过程中型腔内始终有较大的反压作用,且铸件的结晶、凝固又类似于压力铸造,是在额定的结晶压力下完成的,这使反压铸造既有真空吸铸的优点,又有压力铸造的优点。

反压铸造的充型压力和充型速度可以控制,能实现平稳充型,避免液流氧化、卷入气体和冲刷型壁;铸件成形性好、表面粗糙度$Ra$值小,用反压铸造可浇注出壁厚为0.5 mm左右、形状复杂的薄壁件;铸件晶粒细小均匀,组织致密,力学性能高;能实现可控气氛浇注;金属液利用率高,劳动条件好。

2. 挤压铸造

挤压铸造是对定量浇入铸型型腔中的液态金属施加较大的机械压力,使其成形、凝固而获得零件毛坯的一种工艺方法,又称为液态模锻。

挤压铸造的工艺过程如图7-30所示。


\begin{figure}[H]
\centering
\includegraphics[page=15, viewport={165.0 565.3 384.0 752.3}, clip]{12-07572-001FigRevised.pdf}
\caption{挤压铸造的工艺过程}
\end{figure}

按加压时型腔中金属的状态,挤压铸造可分为液态成形和半固态成形两种。液态成形是将金属液浇入铸型后立即挤压成形。半固态成形是将金属液浇入铸型后,待一定时间使其结晶成半固态时再加压成形,而且压力一直作用至金属凝固完毕。

挤压铸造是介于铸造与锻造之间的一种新工艺方法,它兼有二者的优点。挤压铸造金属液直接浇入型腔中而不经过浇注系统,吸气少,铸件可进行热处理。挤压铸造没有浇注系统,金属液在压力作用下充型、结晶和凝固,补缩效果好,晶粒较细,组织致密均匀,力学性能较好。挤压铸造的模具结构较简单,加工费用较低,寿命较长,金属的利用率较高。挤压铸件的力学性能虽稍低于锻件,但只要工艺正确,其力学性能可接近或达到锻件的水平,且各方向性能均匀。挤压铸件的成形是通过封闭型腔里的液态金属在压力作用下结晶、凝固实现的,而锻件是固态金属在压力作用下成形的,因而挤压铸造所需的压力比锻造小得多,所需设备的功率也比锻造小。挤压铸件的尺寸精度比锻件高,甚至比熔模铸件高,加工余量小。因此,挤压铸造适用于多种合金材料,不仅可用于铸造铝合金、锌合金、铜合金、铸铁、铸钢以及部分变形合金,也可用于铝-铁双金属、碳纤维-铝等复合材料。挤压铸造的工艺过程较简单,节省能源,容易实现机械化和自动化,生产效率高,一般比金属型铸造高1~2倍。

3. 实型铸造

实型铸造又称为汽化模铸造或消失模铸造。它是采用聚苯乙烯泡沫塑料模样代替普通模样,将刷过涂料的模样放入可抽真空的特制砂箱,填干砂后振动紧实、抽真空,不用取出模样就可浇入金属液。在高温液体金属的热作用下,泡沫塑料模样汽化、燃烧而消失,金属液取代原来泡沫塑料模样所占据的空间位置,冷却凝固后即可获得所需要的铸件。实型铸造的工艺过程如图7-31所示。


\begin{figure}[H]
\centering
\includegraphics[page=15, viewport={21.0 318.3 328.5 548.3}, clip]{12-07572-001FigRevised.pdf}
\caption{实型铸造的工艺过程}
\end{figure}

与砂型铸造相比,实型铸造的铸件尺寸精度较高;由于模样无分型面,不存在分型起模等问题,因而增大了设计铸造零件的自由度;在多数情况下不用型芯,省去芯盒制造、芯砂配制、芯骨准备、型芯的制造及烘干等工序,简化了生产工序,缩短了生产周期,提高了生产效率(造型效率可提高2~5倍);冒口的金属利用率高。采用实型铸造可节省大量木材,所用的泡沫塑料模样的成本一般只为木模样的1/3左右。采用无黏结剂干砂真空实型铸造,可节省大量的型砂黏结剂,型砂可以回收再利用,型砂处理简单,所需的设备少,可节省投资60\%~80\%。

4. 磁型铸造

磁型铸造是在实型铸造的基础上发明的一种新的铸造方法。它是用聚苯乙烯泡沫塑料制成汽化模,在其表面上刷涂料,放进特制的砂箱内,填入磁丸(又称铁丸)并微振紧实,将砂箱放在磁型机里后通电,使磁丸相互吸引,形成强度好、透气性高的铸型。浇注时,汽化模在液态金属的热作用下汽化消失,金属液取代了汽化模原来的位置,待金属凝固冷却后解除磁场,磁丸恢复原来的松散状态,从而可方便地取出铸件。图7-32所示为磁型铸造原理示意图。

\begin{figure}[H]
\centering
\includegraphics[page=15, viewport={347.5 321.8 523.5 443.3}, clip]{12-07572-001FigRevised.pdf}
\caption{磁型铸造原理示意图}
\end{figure}

磁型铸造中用作汽化模的材料与实型铸造基本相同,所用的磁丸在未通电时与无黏结剂的干砂一样,具有良好的流动性,可充填到砂箱的各个部位。磁型铸造采用磁丸作为造型材料,无需黏结剂,磁丸的颗粒均匀,流动性及透气性好,不含水分,发气性小,铸件产生气孔、夹砂等缺陷少。磁型的冷却速度比砂型大3倍以上,但无一般金属型的激冷作用,改善了铸件的凝固条件,细化了金属组织,提高了铸件的力学性能。

磁型铸造的生产面积较小,设备投资少,原材料消耗少,金属的利用率高。磁型铸件的清理工作少,不需配制、运输型砂,落砂容易。埋模、浇注、落砂等工序紧凑,可集中采用吸尘排烟设备,减少污染,改善了劳动条件。

5. 气冲造型

气冲造型是近来发展起来的一种新的造型方法。它是将储存在压力罐内的压缩空气突然释放出来,作用在砂箱里松散的型砂上,使其紧实成形,或利用可燃气体燃烧爆炸产生的冲击波使型砂紧实成形。气冲造型的工艺过程如图7-33所示,具体工序如下。


\begin{figure}[H]
\centering
\includegraphics[page=15, viewport={71.0 125.8 420.5 294.8}, clip]{12-07572-001FigRevised.pdf}
\caption{气冲造型的工艺过程}
\end{figure}

1) 向砂箱及辅助砂框里加入松散的型砂,并向压力罐里充入压力为0.3~0.7 MPa的压缩空气(图7-33a)。

2) 砂箱及辅助砂框上升与空气冲击装置(压力罐)接触并压紧(图7-33b)。

3) 打开冲击阀,压力罐内的压缩空气迅速地进入型砂上面冲击型砂,使型砂紧实(图7-33c)。

4) 关闭冲击阀,排出砂箱的残留压缩空气,将模板、砂箱及辅助砂框分别下降不同的高度,并回程起模(图7-33d)。

气冲造型的紧实过程简单,速度快,型砂在没有振击与压实的条件下被紧实。与其他造型方法相比,气冲造型铸型的紧实度高且均匀;铸件的吃砂量小,砂箱的利用率高;铸件的尺寸精度高,表面粗糙度值较小;生产效率高,每次冲击紧实的时间不足0.1 s,制造一个铸型只要0.5 s左右;气冲造型机无振击机构,无连续的噪声,在冲击紧实的瞬间,噪声虽达100 dB(A)左右,但一般都在86 dB(A)以下,改善了劳动条件。由于气冲造型机没有振动部分,结构较简单,运动部件少,机器的安装容易,维修工作量少。

\subsection{铸造技术的发展趋势}

随着科学技术的迅速发展和全球可持续发展战略的实施,现代铸造技术正朝着清洁化、专业化、高效化、智能化、数字化、网络化、集成化和铸件的高性能化、精确化、轻质薄壁化的方向发展。

(1) 实现绿色集约化铸造,保证可持续发展

铸造应不断应用新技术、新工艺、新材料、新设备,实现低消耗、无废弃物铸造,使铸造生产的

环境宜人化,生产的产品清洁化。

(2) 实现专业化铸造生产

规模经营、技术先进、管理良好是现代企业的基本要素,专业化生产、规模经营有利于提高技术水平,实现机械化、自动化、智能化生产,提高产品质量,降低生产成本,提高效益。

(3) 不断提高铸件的外观质量和内在质量,实现铸件的高性能化、精确化、轻质薄壁化

应大力发展和应用树脂砂、冷芯盒、温芯盒法及壳型(芯)法以及涂料技术;采用特种铸造和现代铸造方法(如熔模铸造、压力铸造、低压铸造、半固态铸造、实型铸造和快速铸造等),以实现铸件外观质量的精确化及近净成形。大力发展和应用铸铁感应电炉熔炼工艺、先进的铁液脱硫技术和过滤技术;采用氩气净化、钙线喂入工艺、氩氧脱碳(argon oxygen decarburization, AOD)精炼工艺和真空氧脱碳(vacuum oxygen decarburization, VOD)精炼工艺等钢液精炼技术,并应加强金属基复合材料的开发和应用等,从而大幅度改善铸件的内在质量,提高铸件的性能,实现铸件的轻量化和薄壁化。

(4) 加强计算机和机器人(机械手)的应用,实现铸造生产的集成化、智能化、数字化、网络化和虚拟化

计算机技术几乎被应用到铸造生产的各个环节中。计算机辅助设计、模拟与仿真技术已向集成化、智能化、网络化、虚拟化和绿色化方向发展。以三维造型为基础的CAD/CAE/CAPP/CAM的集成,尤其进一步与快速成形(rapid prototyping manufacturing, RPM)技术的集成,可以构成一个闭环快速产品开发系统,在并行工程(concurrent engineering, CE)环境下,能对产品设计进行快速评价、修改和仿真,以响应市场大规模个性化生产的需要。铸造CAE技术的高度集成,包括多物理场(温度场、流动场、应力应变场等)的集成、多合金材质的集成及多铸造方式的集成;铸造生产智能化是基于铸造CAE技术、铸造专家系统及信息处理系统来实现的,能实现铸造缺陷的自动预报、自动反馈与自动响应,完成最佳铸造工艺的智能决策;网络化是利用互联网技术,实现计算机协同工作,达到远程设计和制造的目的。计算机模拟与仿真技术已实现微观组织模拟(从毫米、微米到纳米尺度),为深入研究材料的结构、组成及其各物理、化学过程中宏观、微观变化机制奠定了基础,因此可以材料成分、结构及制备参数的最佳组合进行材料设计。计算机检测与控制主要用于型砂性能及处理过程的在线检测与控制、熔炼过程的在线检测与控制、金属液质量的炉前快速检测与监控、铸件成形过程的检测与控制及铸件成品质量的检验等方面。随着人工智能技术水平的提高,机器人(机械手)在工作环境恶劣的铸造生产中的应用也会越来越多。

\section{常用合金铸件生产特点}

\subsection{铸铁件生产}

1. 灰铸铁

(1) 灰铸铁的生产特点

目前,大多数灰铸铁用冲天炉熔炼。冲天炉炉料由金属炉料、燃料(焦炭、天然气)和熔剂(石灰石、氟石)组成。金属炉料包括高炉铸造生铁、回炉铁(废旧铸件、浇冒口等)、废钢和铁合金(硅铁、锰铁等)。近年来,已有不少工厂用工频感应炉来熔炼灰铸铁,可获得洁净、高温、成分准确的优质铁液。

灰铸铁件主要用砂型铸造,因其铸造性能优良,便于制出薄而复杂的铸件,一般不需设置冒口和冷铁,铸造工艺简化;又因其浇注温度较低,故中、小型铸件多采用经济简便的湿型铸造。

(2) 灰铸铁的孕育处理

为了改善灰铸铁的组织,提高灰铸铁的强度和其他性能,生产中常进行孕育处理。

孕育处理就是在浇注前往铁液中加入孕育剂,使石墨细化,基体组织细密(珠光体基体)。生产中常用的孕育剂是硅的质量分数为75\%的硅铁,加入量为铁液质量的0.25\%$\sim$0.6\%。

孕育处理的灰铸铁生产要求比灰铸铁高。为了保证较好的孕育效果,必须先熔炼出低碳、低硅的铁液(一般$w_{\mathrm{C}}=2.7\% \sim 3.3\%$、$w_{\mathrm{Si}}=1.0\% \sim 2.0\%$)。另外,由于孕育处理过程中铁液温度要降低,故铁液应有较高的出炉温度(1450$\sim$1470$^{\circ}$C)。

2. 球墨铸铁

(1) 球墨铸铁的化学成分

球墨铸铁的化学成分与灰铸铁基本相同,但要求严格:高碳($w_{\mathrm{C}}=3.6\% \sim 4.0\%$)和高硅,以改善铸造性能和球化效果;低硫($w_{\mathrm{S}}<0.06\%$),硫会增加球化剂损耗,严重影响球化效果;低磷,磷会降低塑性、韧性和强度,增加冷脆性。

(2) 球化处理和孕育处理

球化处理和熔炼优质铁液同为生产球墨铸铁件的关键环节。球化剂的作用是使石墨呈球状析出。常用的球化剂是稀土镁合金,其球化能力强,球化效果好,与铁液反应平稳,工艺过程简便。

孕育剂的作用主要是进一步促进铸铁石墨球化,防止球化衰退和球化元素所造成的白口倾向。同时,孕育处理还可使石墨圆整、细化,改善球墨铸铁的力学性能。常用孕育剂是硅的质量分数为75\%的硅铁,加入量为铁液质量的0.4\%$\sim$1.0\%。

球化处理普遍采用冲入法,如图7-34所示。冲入法是在浇包底部修成“堤坝”,把球化剂置于坝内,盖以硅铁粉或铁屑,以防球化剂上浮,延缓反应速度,使处理过程充分进行。球化剂的加入量为铁液质量的1.3\%$\sim$1.8\%(含硫高时取上限)。处理时,先冲入铁液总量的1/3$\sim$1/2,待反应完毕后,加稻草灰扒渣,再加入其余铁液,经孕育处理、炉前检验合格后即可浇注。

\begin{figure}[H]
    \centering
    \includegraphics[page=15, viewport={117.5 17.3 427.5 98.3}, clip]{12-07572-001FigRevised.pdf}
    \caption{冲入法球化处理示意图}
    \label{fig:7-34}
\end{figure}

(3) 球墨铸铁的铸造工艺特点

球墨铸铁碳的质量分数高,接近共晶成分,其流动性与灰铸铁相近,可生产壁厚为3$\sim$4 mm的铸件。球墨铸铁在浇注后的一段时间,内、外壁与中心几乎同时凝固,造成凝固后期外壳强度低,同时球状石墨析出时膨胀力很大,若铸型刚度不够,易造成铸件内部金属液的不足而产生缩孔和缩松。因此,球量铸铁铸造时常常增设冒口和冷铁,采用顺序凝固,同时使用干型或快干型水玻璃等措施增大铸型刚度,以防止上述缺陷的产生。

由于铁液中的MgS与型砂中的水分反应,生成H$_2$S气体,易产生皮下气孔,所以应严格控制型砂中水分和铁液中硫的质量分数。

另外,由于铁液经球化和孕育处理后温度要降低50$\sim$100$^{\circ}$C,为防止冷隔和浇不足,出炉的铁液温度必须高达1400$^{\circ}$C以上。

3. 可锻铸铁

可锻铸铁是用低碳、低硅的铁液浇注出白口铸铁组织的铸件毛坯,再经高温(900$\sim$950$^{\circ}$C)石墨化(可锻化)退火使白口铸铁件中的渗碳体分解为团絮状石墨,从而得到由团絮状石墨和不同基体组织组成的铸铁。由于碳、硅的质量分数低,凝固时没有石墨析出,凝固收缩大,且可锻铸铁的熔点比灰铸铁高,结晶温度范围较宽,故其流动性差,所以易产生浇不足、冷隔、缩孔、缩松、裂纹等缺陷。因此,在工艺设计时应特别注意冒口及冷铁的位置,以增强补缩能力。同时,要求铁液出炉温度较高,一般不低于1360$^{\circ}$C。另外,由于白口铁铁液比灰铸铁液的含气量高,且黏度大,易产生皮下气孔,因此要求型砂的含水量低,并有足够的透气性。

4. 蠕墨铸铁

蠕墨铸铁件的生产与球墨铸铁件较为相似,但也具有不同的特点。

(1) 蠕墨铸铁的铸造性能

在充分蠕化的条件下,蠕墨铸铁的铸造性能与灰铸铁相近。它具有比灰铸铁更高的流动性(因除气和净化好),可浇注复杂铸件及薄壁铸件;其收缩性介于灰铸铁和球墨铸铁之间,倾向于形成集中缩孔;因其具有共晶成分或接近于共晶成分,故热裂倾向小;其有一定的塑性,不易产生冷裂纹。

(2) 蠕墨铸铁的铸造工艺特点

蠕墨铸铁件的生产过程与球墨铸铁件相似,主要包括熔炼铁液、蠕化-孕育处理和浇注等,但一般不进行热处理,而以铸态使用。为此,需特别重视其化学成分和蠕化-孕育效果。蠕墨铸铁的碳的质量分数也较高,一般为4.3\%$\sim$4.6\%,但以低碳、高硅为原则,以利于形成蠕虫状石墨。蠕化-孕育处理时,常用的蠕化剂是稀土硅铁蠕化剂(各元素的质量分数分别为稀土17\%$\sim$20\%, Si 40\%, Mg 0.5\%$\sim$1.0\%, Ca 2\%,其余为铁)、镁钛铈蠕化剂和稀土硅钙合金等,一般也采用冲入法,把蠕化剂埋入浇包底部凹坑内,用铁液冲熔和吸收。孕育剂采用硅的质量分数为75\%的硅铁,多用液流法进行孕育处理。

铸造蠕墨铸铁件时,也要注意防止蠕化-孕育的衰退现象,特别注意铁液中要有适宜的残留稀土量(稀土的质量分数为0.02\%$\sim$0.03\%),以保证蠕化效果。

\subsection{铸钢件生产}

1. 概述

铸钢是重要的铸造合金。按照化学成分,铸钢可分为铸造碳钢和铸造合金钢两大类。铸造碳钢应用最广,占总产量的80\%以上。

铸钢的综合力学性能高于各类铸铁,不仅强度高,且具有优良的塑性和韧性。此外,铸钢的焊接性好,可采用铸焊组合工艺制造重型零件。但铸钢件晶粒粗大、组织不均,且常存有残余内应力,致使铸件的强度,特别是塑性和韧性不如锻钢件高。因此,铸钢件必须进行热处理,一般采用正火或退火。正火后铸钢件的力学性能较退火铸钢件的好且成本低,所以一般采用正火。但正火铸钢件较退火铸钢件的内应力大,因此易产生裂纹或易硬化的铸钢件应进行退火。

2. 铸钢的熔炼

目前,在铸钢生产中应用最普遍的炼钢设备是三相电弧炉。近年来,感应电炉炼钢发展得很快。感应电炉炼钢的加热速度较快,氧化烧损较少。尤其采用酸性感应电炉不氧化法炼钢时,其熔炼过程基本上就是炉料的重熔过程,操作简便,因此广泛地应用于精密铸造生产。用于炼钢的感应电炉多为中频炉(500$\sim$1000 Hz),其容量多为0.25$\sim$3 t。

重型机械厂也用平炉作为炼钢设备,容量通常在100 t以下,适于铸造重型铸件。

3. 铸钢的铸造工艺特点

铸钢的熔点较高(约1500$^{\circ}$C)、流动性差、收缩率大(达2\%)。在熔炼过程中,易吸气和氧化;在浇注过程中易产生黏砂、浇不足、冷隔、缩孔、变形、裂纹、夹渣和气孔等缺陷。因此,在工艺上必须采取相应措施来防止上述铸造缺陷。

铸钢所用型(芯)砂需有良好的透气性、耐火度、强度和退让性。原砂要用颗粒大而均匀的硅砂,大铸件用人造硅砂。为防止黏砂,型腔表面要涂以石英粉或锆砂粉涂料。为减少气体来源、提高强度、改善填充条件,大件多采用干型或水玻璃快干型。

由于铸钢的流动性差、收缩率大,其浇注系统的形状较简单,截面积较大,并且应遵循顺序凝固原则进行工艺设计,配置较大的冒口和冷铁来防止缩孔的产生。对薄壁或易产生裂纹的铸钢件,则应采用同时凝固原则。

4. 铸钢的应用

铸钢主要用于制造一些形状复杂,用其他方法难以制造,且要求有较好力学性能的零件,如高压阀门壳体、轧钢机的机架、某些齿轮等。某些要求有极高耐磨性的零件,如碎石机颚板、挖掘机铲齿、坦克履带等,采用ZGMn13制造,由于材料切削性极差,只能用铸造生产。许多特大型零件,如大型发电机轴、水轮机叶轮、轧辊、水压机缸体等零件也只能采用铸钢件。航空发动机上的轴流转子、导风轮等零件,用沉淀硬化不锈钢制造。因零件形状复杂,材料切削性很差,因此采用熔模铸造生产毛坯,只需经少量磨削就可以装配使用。高精度的无余量精密铸件,只需抛光加工就可以使用。有些铸钢为了节省材料,减少切削加工量,降低生产成本,提高生产效率,就采用铸造方法生产,如齿轮、支架、联轴器、接头等。铸钢件可采用砂型铸造、熔模铸造等方法生产。

\section{铸件结构设计}

设计铸件时,不仅要使其结构能满足零件的使用要求,还应考虑铸造成形的可行性和经济性,使铸件结构尽可能简化和方便铸造生产。良好的铸件结构应与金属的铸造性能、铸件的铸造工艺相适应。

\subsection{铸造性能对铸件结构的要求}

铸件中很多缺陷的出现都是铸件结构设计时未考虑合金的铸造性能所致,如气孔、裂纹、缩孔等。设计时,应依据如下原则。

1. 铸件壁厚应适当

铸件壁厚影响金属液浇注后的冷却速度,每种合金都有其适宜的壁厚范围。如铸件壁厚设计合理,则既能保证铸件的力学性能,又能防止铸件产生某些缺陷。设计时,铸件壁厚应大于合金的最小壁厚,否则易产生浇不足、冷隔等缺陷。

不同合金的流动性不同,所能铸造的铸件的最小壁厚也不同。砂型铸造条件下不同合金的最小壁厚见表7-2。铸件的壁厚也不宜过大,否则铸件心部易产生晶粒粗大、缩孔、缩松、偏析等缺陷。为避免铸件具有厚大截面,又能充分发挥材料潜力,保证铸件的强度和刚度,可从铸件截面形状方面加以考虑,如选择T字形、工字形、槽形等截面形状,在铸件的薄弱部位还可设置加强筋来提高其强度和刚度。

\begin{table}[H]
\centering
\caption{砂型铸造条件下不同合金的最小壁厚}
\label{tab:7-2}
\small
\begin{tabular}{|c|c|c|c|c|c|c|}
\hline
铸件尺寸 & 铸钢 & 灰铸铁 & 球墨铸铁 & 可锻铸铁 & 铝合金 & 铜合金 \\
\hline
$<200 \times 200$ & 5$\sim$8 & 3$\sim$5 & 4$\sim$6 & 3$\sim$5 & 3$\sim$3.5 & 3$\sim$5 \\
\hline
$200 \times 200 \sim 500 \times 500$ & 10$\sim$12 & 4$\sim$10 & 8$\sim$12 & 6$\sim$8 & 4$\sim$6 & 6$\sim$8 \\
\hline
$>500 \times 500$ & 15$\sim$20 & 10$\sim$15 & 12$\sim$20 & --- & --- & --- \\
\hline
\end{tabular}\end{table}

2. 铸件壁厚应均匀

铸件的壁厚应尽可能均匀,尽量减小其差别。如铸件各部分的壁厚差别过大,则厚壁处易形成热节,产生缩孔、缩松、内应力、裂纹。壁厚均匀是指铸件具有冷却速度相近的壁厚。如内壁(隔墙)冷却慢应薄一点,外壁冷却快应厚一点。例如,图7-35a所示铸件的壁厚不均匀,在其厚大部位会形成许多小缩孔;图7-35b所示为改进后的结构设计,由于壁厚均匀,避免了缩孔缺陷的产生。

\begin{figure}[H]
    \centering
    \includegraphics[page=16, viewport={28.5 646.3 304.0 708.8}, clip]{12-07572-001FigRevised.pdf}
    \caption{铸件的壁厚}
    \label{fig:7-35}
\end{figure}

3. 铸件壁的连接应合理

铸件壁与壁之间应过渡连接,避免形成热节及产生较大的内应力,影响铸件质量。

(1) 结构圆角

以铸件为毛坯的零件结构应尽可能将壁间连接设计成结构圆角,以免局部金属聚集产生缩孔、应力集中等缺陷。圆角半径的大小应与铸件壁厚相适应。例如,图7-36a所示结构直角处内接圆较大,表明金属在这里聚集,可能产生缩孔;直角处易产生应力集中,可能产生裂纹。而图7-36b所示结构圆角处则没有发生金属聚集及应力集中的现象。

\begin{figure}[H]
    \centering
    \includegraphics[page=16, viewport={322.5 644.8 521.5 731.8}, clip]{12-07572-001FigRevised.pdf}
    \caption{铸件的结构圆角}
    \label{fig:7-36}
\end{figure}

(2) 过渡连接

铸件各部分之间的连接都要逐步过渡,避免锐角和交叉连接。例如,图7-37a所示筋的交叉连接在交叉处有金属聚集,可能形成缩孔。小型铸件的筋应设计成图7-37b所示的交错连接,大型铸件的筋应设计成图7-37c所示的环状连接,以改善金属的聚集程度。

\begin{figure}[H]
    \centering
    \includegraphics[page=16, viewport={48.0 466.8 347.0 573.3}, clip]{12-07572-001FigRevised.pdf}
    \caption{铸件壁的连接形式}
    \label{fig:7-37}
\end{figure}

铸件壁间连接应避免形成锐角。例如,图7-38a所示锐角连接容易形成金属聚集;图7-38b和图7-38c所示大角度连接改善了金属的分布情况。

铸件的薄、厚壁之间的连接可采用圆角过渡、倾斜过渡、复合过渡等形式,如图7-39所示。过渡连接可防止因壁间突然变化而产生内应力、变形和裂纹。

\begin{figure}[H]
    \centering
    \includegraphics[page=16, viewport={368.0 468.3 500.5 613.3}, clip]{12-07572-001FigRevised.pdf}
    \caption{铸件壁之间的连接}
    \label{fig:7-38}
\end{figure}

\begin{figure}[H]
    \centering
    \includegraphics[page=16, viewport={62.0 330.3 237.0 435.8}, clip]{12-07572-001FigRevised.pdf}
    \caption{铸件不等厚壁之间的连接}
    \label{fig:7-39}
\end{figure}

4. 铸件结构应能减少变形

壁厚均匀的细长铸件和面积较大的平板类铸件容易产生变形,通常设计成对称结构或增设加强筋,以防止变形。例如,图7-40a所示工字梁铸件常设计成对称结构;图7-40b所示平板铸件底面上增设了加强筋。

5. 铸件结构应能减缓收缩受阻

铸件在冷却过程中,固态收缩受阻是产生内应力、变形和裂纹的根本原因。铸件结构设计应尽可能使其各部分能自由收缩,减缓收缩受阻。图7-41a所示为弯曲的轮辐设计,可使轮辐在冷却时产生一定的自由收缩;图7-41b所示为奇数轮辐设计,同样可使其在冷却过程中产生一定的自由收缩,从而减小铸造内应力。

\begin{figure}[H]
    \centering
    \includegraphics[page=16, viewport={256.5 327.3 480.5 436.8}, clip]{12-07572-001FigRevised.pdf}
    \caption{防止变形的铸件结构}
    \label{fig:7-40}
\end{figure}

\begin{figure}[H]
    \centering
    \includegraphics[page=16, viewport={35.0 47.8 191.0 139.8}, clip]{12-07572-001FigRevised.pdf}
    \caption{轮辐的设计}
    \label{fig:7-41}
\end{figure}

\subsection{铸造工艺对铸件结构的要求}

铸件结构应能简化铸造工艺,便于机械化生产,提高铸件质量,降低废品率。为此,铸件结构应尽量简单合理,避免不必要的复杂结构。

1. 铸件的外形设计应便于起模,尽量避免外部侧凹结构,使分型面少而平直

铸件结构设计时应考虑铸造工艺方便,铸件外形应尽量简单,尽量避免外部侧凹结构,凸台、筋条的设计应便于起模,并尽可能使铸件具有一个简单的平直分型面,如图7-42、图7-43和图7-44所示。

\begin{figure}[H]
    \centering
    \includegraphics[page=16, viewport={216.0 48.3 513.0 288.8}, clip]{12-07572-001FigRevised.pdf}
    \caption{避免侧凹结构}
    \label{fig:7-42}
\end{figure}

\begin{figure}[H]
    \centering
    \includegraphics[page=17, viewport={110.5 558.3 434.0 749.3}, clip]{12-07572-001FigRevised.pdf}
    \caption{凸台设计应避免活块造型}
    \label{fig:7-43}
\end{figure}

\begin{figure}[H]
    \centering
    \includegraphics[page=17, viewport={50.0 416.8 278.0 523.3}, clip]{12-07572-001FigRevised.pdf}
    \caption{分型面应平直}
    \label{fig:7-44}
\end{figure}

2. 铸件内腔设计要合理、简单

铸件内腔设计应能减少型芯的数量(图7-45和图7-46),简化型芯形状,并有利于型芯的定位、固定、排气和清理(图7-47和图7-48),防止产生偏芯、气孔等缺陷。

\begin{figure}[H]
    \centering
    \includegraphics[page=17, viewport={295.5 419.3 497.5 503.3}, clip]{12-07572-001FigRevised.pdf}
    \caption{以砂垛代替型芯}
    \label{fig:7-45}
\end{figure}

\begin{figure}[H]
    \centering
    \includegraphics[page=17, viewport={92.0 227.3 256.0 387.8}, clip]{12-07572-001FigRevised.pdf}
    \caption{采用开式结构省去型芯}
    \label{fig:7-46}
\end{figure}

\begin{figure}[H]
    \centering
    \begin{minipage}{0.45\textwidth}
        \centering
        \includegraphics[page=17, viewport={281.0 229.3 450.5 298.8}, clip]{12-07572-001FigRevised.pdf}
        \caption{轴承支架}
    \end{minipage}
    \begin{minipage}{0.45\textwidth}
        \centering
        \includegraphics[page=17, viewport={80.5 105.8 273.5 192.3}, clip]{12-07572-001FigRevised.pdf}
        \caption{增设工艺孔}
    \end{minipage}
\end{figure}

3. 结构斜度

铸件上垂直于分型面的不加工表面上所加的斜度称为结构斜度。一般铸件上垂直于分型面的不加工表面均应设计结构斜度，如图 7-49 所示。

\begin{figure}[H]
    \centering
    % Placeholder for Figure 7-49
    \includegraphics[page=17, viewport={290.0 106.3 461.5 169.8}, clip]{12-07572-001FigRevised.pdf}
    \caption{铸件的结构斜度}
\end{figure}

设计结构斜度不仅使起模方便，而且零件也更加美观；具有结构斜度的内腔常常可以采用砂垛代替型芯。零件上垂直于分型面的不加工表面的高度越低，结构斜度应越大。

\section*{复习思考题}

 
7-1 什么是液态合金的充型能力？它对铸件质量有何影响？与合金的流动性有何关系？化学成分对合金的流动性有何影响？

7-2 既然提高浇注温度能提高合金液的充型能力，为何又要防止浇注温度过高？

7-3 液态合金浇入铸型后要经历哪几个收缩阶段？对铸件质量各有何影响？铸造模样尺寸与哪个收缩阶段密切相关？

7-4 铸件中的缩孔和缩松是如何形成的？其形成的根本原因何在？如何防止铸件的缩孔和缩松？从工艺上看哪种更难防止？为什么？

7-5 生产上经常采用哪些方法来确定缩孔的位置？

7-6 何谓顺序凝固原则？何谓同时凝固原则？从工艺上如何实现这两种凝固原则？它们各适用于什么场合？

7-7 铸造内应力分为哪几类？热应力是如何形成的？在铸件不同部位的应力状态如何？

7-8 铸件变形的原因何在？如何防止铸件的变形？

7-9 球墨铸铁在化学成分、力学性能和铸造工艺方面有何特点？

7-10 分析如图 7-50 所示的槽形铸件的热应力形成过程，标出最终热应力状态，并用虚线画出铸件的变形方向。

7-11 简述砂型铸造的生产过程，并说明其优缺点。

7-12 什么是浇注位置？浇注位置选择的一般性原则是什么？

7-13 什么是分型面？分型面选择的一般原则是什么？

7-14 铸造工艺图包括哪些内容？试确定图 7-51 所示的机床床身铸件的分型面和浇注位置，并说明原因。

\begin{figure}[H]
    \centering
    % Placeholder for Figure 7-50
    \includegraphics[page=18, viewport={128.0 643.8 302.5 697.8}, clip]{12-07572-001FigRevised.pdf}
    \caption{题 7-10 图}
\end{figure}
\begin{figure}[H]
    \centering
    % Placeholder for Figure 7-51
    \includegraphics[page=18, viewport={316.5 648.8 417.0 746.8}, clip]{12-07572-001FigRevised.pdf}
    \caption{题 7-14 图}
\end{figure}

7-15 图 7-52 所示的 3 种铸件应采用何种手工造型方法？试确定它们的分型面和浇注位置。

\begin{figure}[H]
    \centering
    % Placeholder for Figure 7-52
    \includegraphics[page=18, viewport={38.5 444.3 382.0 614.3}, clip]{12-07572-001FigRevised.pdf}
    \caption{题 7-15 图}
\end{figure}

7-16 熔模铸造有何优缺点？和实型铸造相比有哪些不同？

7-17 金属型铸造、压力铸造、反压铸造和挤压铸造有何不同？

7-18 离心铸造有何优缺点？主要适用于哪些场合？

7-19 在大批量生产时，铝活塞、机床床身、大直径铸铁管、汽轮机叶片和薄壁波导管最适宜采用哪一种铸造方法？

7-20 为什么铸钢件的铸造性能比铸铁件差？

7-21 简要论述铸造技术的发展趋势。

7-22 铸件结构应力求简单，以简化铸造工艺，提高铸造质量与生产效率，因此从铸件结构设计考虑，应从哪些方面入手来提高其铸造工艺性？

7-23 什么是铸件的最小壁厚？为何要规定铸件的最小壁厚？是否铸件壁厚越大越好？

7-24 为什么要设计铸件的结构圆角？图 7-53 所示铸件的结构设计是否合理？如不合理，在不改变分型面和浇注位置的前提下加以修改。

\begin{figure}[H]
    \centering
    % Placeholder for Figure 7-53
    \includegraphics[page=18, viewport={408.0 446.8 505.5 540.3}, clip]{12-07572-001FigRevised.pdf}
    \caption{题 7-24 图}
\end{figure}

7-25 指出图 7-54 所示铸件结构设计的不合理之处，并加以改正。

\begin{figure}[H]
    \centering
    % Placeholder for Figure 7-54
    \includegraphics[page=18, viewport={86.5 73.3 456.5 418.3}, clip]{12-07572-001FigRevised.pdf}
    \caption{题 7-25 图}
\end{figure}

7-26 为什么要设计铸件的结构斜度？铸件的结构斜度与起模斜度有何异同？


\chapter{压力加工}

压力加工也称为塑性加工，它是利用金属坯料在外力作用下产生塑性变形，从而获得所需形状、尺寸和性能的毛坯或零件的加工方法。

压力加工方法主要有轧制、挤压、拉拔、自由锻、模锻和板料冲压等（表 8-1）。前三种方法以生产原材料为主，后三种方法以生产毛坯为主。

\begin{longtable}{|c|m{3cm}|m{4.3cm}|m{2.2cm}|m{4cm}|}	
	\caption{压力加工方法的工作原理、应用与发展趋势}\\
	\hline
	名称 & 示意图 & 加工方法示例 & 工作原理 & 应用与发展趋势 \\
	\hline
	轧制 & \includegraphics[page=29, viewport={229.5 327.3 294.0 388.8}, clip]{12-07572-001FigRevised.pdf} & \includegraphics[page=29, viewport={289.5 219.3 390.5 292.8}, clip]{12-07572-001FigRevised.pdf} & 使金属坯料通过一对回转轧辊间的空隙而产生连续成形 & 应用：生产钢板、无缝钢管及各种型钢（如圆钢、方钢、扁钢、角钢、槽钢、工字钢、钢轨等）；\newline 发展趋势：高速轧制，精密轧制，轧锻组合等 \\
	\hline
	挤压 & \includegraphics[page=29, viewport={316.5 323.8 386.0 389.3}, clip]{12-07572-001FigRevised.pdf} & \includegraphics[page=29, viewport={404.5 217.8 533.0 273.8}, clip]{12-07572-001FigRevised.pdf} & 使金属坯料从挤压模的模孔中挤出而成形 & 应用：生产低碳钢、非铁金属及其合金的型材、管件和零件；\newline 发展趋势：高速精密挤压，挤锻结合 \\
	\hline
	拉拔 & \includegraphics[page=29, viewport={407.0 325.8 480.0 385.8}, clip]{12-07572-001FigRevised.pdf} & \includegraphics[page=29, viewport={8.0 110.8 132.0 187.3}, clip]{12-07572-001FigRevised.pdf} & 将金属坯料从拉拔模的模孔中拉出而成形 & 应用：可拉制直径仅为 0.02 mm 的金属丝和薄壁管，也可用于提高轧制型材和管材的精度和表面质量；\newline 发展趋势：高尺寸精度，表面粗糙度值小 \\
	\hline
	自由锻 & \includegraphics[page=29, viewport={11.0 219.3 74.5 274.8}, clip]{12-07572-001FigRevised.pdf} & \includegraphics[page=29, viewport={147.0 112.8 263.0 186.3}, clip]{12-07572-001FigRevised.pdf} & 将金属坯料放在上、下砧间受冲击力或压力而成形 & 应用：单件、小批生产力学性能高、形状简单的零件毛坯，是制造大型锻件的唯一方法；\newline 发展趋势：锻件大型化，操作机械化，液压机代替大锻锤 \\
	\hline
	模锻 & \includegraphics[page=29, viewport={96.5 220.3 168.0 278.8}, clip]{12-07572-001FigRevised.pdf} & \includegraphics[page=29, viewport={277.0 111.3 394.5 180.8}, clip]{12-07572-001FigRevised.pdf} & 将金属坯料放在模锻模膛内受冲击力或压力而成形 & 应用：成批、大量生产形状较复杂的中、小型模锻件；\newline 发展趋势：少、无切屑精密化（如精密模锻） \\
	\hline
	板料冲压 & \includegraphics[page=29, viewport={188.5 219.8 270.0 270.3}, clip]{12-07572-001FigRevised.pdf} & \includegraphics[page=29, viewport={415.5 110.8 535.0 187.8}, clip]{12-07572-001FigRevised.pdf} & 利用冲模，使金属板料产生分离或成形 & 应用：成批、大量生产形状复杂的薄板件，仪器、仪表件，中空零件或汽车覆盖件等；\newline 发展趋势：自动化，精密化，非传统成形工艺的发展 \\
	\hline
\end{longtable}

压力加工产品质量高，性能好，缺陷少。金属材料经塑性变形后，可以消除铸造组织的内部缺陷，得到细晶粒结构，组织致密，并能获得一定的锻造流线组织，使其力学性能提高。因而，承受重载的零件一般都采用锻件作毛坯。压力加工可以节省原材料，由于压力加工是依靠塑性变形重新分配坯料体积而进行成形的，与切削加工方法相比，可以减少零件制造过程中的金属消耗。一般压力加工方法的材料利用率可达 $60\% \sim 70\%$，先进压力加工方法现已达到 $85\% \sim 90\%$。另外，多数压力加工方法容易实现机械化和自动化，生产效率高。线材轧制速度可达 $100 \text{ m/s}$，在 $120\ 000 \text{ kN}$ 压力机上每小时可生产汽车发动机曲轴 90 件。

压力加工与铸造相比，成本较高，成形较困难，由于是在固态下成形，无法获得截面形状（特别是内腔）复杂的产品。

压力加工在工业生产中具有重要地位，在机械、冶金、船舶、航空航天、电子、仪器仪表等行业以及国防工业中得到广泛应用。

\section{压力加工理论基础}

金属压力加工主要是通过工件（坯料）的塑性变形来实现的。金属材料经塑性变形后，其内部组织会发生很大变化，性能得以提高。压力加工时，必须对金属材料施加外力，使之产生塑性变形，若设备吨位不足便达不到预期的变形程度，故外力是坯料转化为产品的外界条件。与此同时，还必须保证坯料产生足够的塑性变形量而不破裂，即要求材料具有良好的塑性，塑性是坯料转化为产品的内因。因而，必须深入了解塑性变形的实质及其对金属组织与性能的影响以及其他相关理论。

\subsection{金属的塑性变形与再结晶}

1. 金属塑性变形的实质

金属受外力作用时，其内部会产生应力，并发生相应的变形，当内应力超过材料的屈服强度时，在外力去除以后金属所保留的变形称为塑性变形。

（1）单晶体的塑性变形

单晶体的塑性变形主要通过滑移和孪生两种方式进行。

1）滑移

金属塑性变形最常见的方式就是滑移。即在切应力作用下，晶体的一部分相对于另一部分沿一定的晶面（滑移面）和晶向（滑移方向）发生的相对位移（图 8-1）。这种位移在应力去除后不能恢复，从而形成金属的宏观塑性变形。图 8-1a 表示晶格变形前的未受力状态；图 8-1b 表示晶格受到切应力作用产生弹性切向变形的情况；图 8-1c 表示切应力逐渐升高至其屈服强度时，晶格产生切向弹、塑性变形；图 8-1d 表示塑性变形后的晶格状态。

\begin{figure}[H]
    \centering
    % Placeholder for Figure 8-1
    \includegraphics[page=19, viewport={146.5 673.8 397.0 748.3}, clip]{12-07572-001FigRevised.pdf}
    \caption{单晶体滑移变形示意图}
\end{figure}

实际上，金属在切应力作用下产生的滑移，并非如图 8-1 所示的滑移面上、下两部分作刚性滑动，而是通过晶体内的位错运动来实现的（图 8-2a）。由于实际晶体存在位错，其在切应力作用下，只有位错中心的少量原子作微量位移（图 8-2b），所以滑移所需要的切应力比理论值低很多。

\begin{figure}[H]
    \centering
    % Placeholder for Figure 8-2
    \includegraphics[page=19, viewport={80.5 540.8 422.5 650.8}, clip]{12-07572-001FigRevised.pdf}
    \caption{位错运动引起滑移变形示意图}
\end{figure}

2）孪生

孪生是在切应力作用下，晶体的一部分沿着孪生面相对于另一部分发生转动的结果。转动后，以孪生面为界面形成镜像对称（图 8-3）。当滑移变形受到限制时，塑性变形可能以孪生方式进行，孪生变形量虽然很小，但是由于转动后改变了晶格的位向，因而有利于滑移的继续进行。

（2）多晶体的塑性变形

多晶体是由许多微小的原子呈不同位向排列的单个晶粒组合而成的。大多数金属材料是由多晶体组成的，其变形主要方式也是滑移和孪生。

\begin{figure}[H]
    \centering
    % Placeholder for Figure 8-3
    \includegraphics[page=19, viewport={74.0 427.8 359.0 512.8}, clip]{12-07572-001FigRevised.pdf}
    \caption{孪生变形示意图}
    $I-I$—生面
\end{figure}

多晶体的塑性变形包括各个单晶体的塑性变形（即晶内变形）和各个晶粒之间的变形（即晶间变形）。一方面，由于多晶体的晶粒取向不同，彼此之间在变形过程中有约束作用，滑移时各晶粒间相互协调配合转动；另一方面，晶界的存在也会对塑性变形产生影响。因此，多晶体的塑性变形比单晶体要复杂得多，并且需要施加更大的外力。

晶界的存在对多晶体变形造成很大障碍，相邻晶粒的位向差越大，其影响也越大。图 8-4 所示为两个晶粒在拉伸前后的变形情况，在远离晶界处的变形很明显，但晶界附近变形却很小，说明晶界处的变形抗力（强度）远较晶粒本身的高。晶界越多，多晶体塑性变形抗力越大，细晶粒多晶体的强度比粗晶粒多晶体要高。因此，除有特殊要求，生产中获得细晶组织是提高产品性能的重要措施之一。

\begin{figure}[H]
    \centering
    % Placeholder for Figure 8-4
    \includegraphics[page=19, viewport={376.0 430.8 469.5 498.3}, clip]{12-07572-001FigRevised.pdf}
    \caption{多晶体试样的拉伸变形}
\end{figure}

2. 塑性变形对金属组织与性能的影响

金属在常温下经塑性变形后，其内部组织将发生变化：晶粒沿变形最大方向伸长；晶格发生畸变，产生内应力；晶粒之间产生碎晶。

随着金属冷变形（再结晶温度以下的变形）程度的增加，材料的力学性能也会随之发生变化，金属的强度、硬度增高，塑性和韧性下降（图 8-5）。这种现象称为冷变形强化，又称为加工硬化。这是因为随变形量增大，晶格畸变越严重，位错的交互作用使塑性变形的阻力增大，使滑移难于进行。冷变形强化具有强化金属的作用，如大型发电机的护环就是利用冷变形强化来提高其强度的。

\begin{figure}[H]
    \centering
    % Placeholder for Figure 8-5
    \includegraphics[page=19, viewport={14.0 186.8 230.0 326.3}, clip]{12-07572-001FigRevised.pdf}
    \caption{冷变形程度对低碳钢力学性能的影响}
\end{figure}

加工硬化是一种不稳定的组织状态。处于高势能的原子具有自发回复到稳定状态的趋势。但是，在常温下原子的活动能力较弱，几乎观察不到。当温度升高时，金属原子获得热能，热运动加剧，最后趋于较稳定的状态，变形金属的组织和性能也会发生一系列变化。随着加热温度的升高，变形金属相继会发生回复、再结晶和晶粒长大 3 个阶段的变化（图 8-6）。

\begin{figure}[H]
    \centering
    % Placeholder for Figure 8-6
    \includegraphics[page=19, viewport={230.5 181.3 534.0 394.8}, clip]{12-07572-001FigRevised.pdf}
    \caption{冷变形金属的回复和再结晶示意图}
\end{figure}

（1）回复

温度升高到回复温度（$T_{\text{回}}$）时，原子回复到正常排列，晶格畸变基本消除。此时，金属的强度、硬度稍有降低，塑性、韧性略有提高，这一过程称为回复。
\[ T_{\text{回}} = (0.25 \sim 0.3) T_{\text{熔}} \]
式中：$T_{\text{回}}$——回复温度，K；
$T_{\text{熔}}$——金属的熔点，K。

（2）再结晶

当温度达到再结晶温度（$T_{\text{再}}$）时，变形金属的显微组织会发生显著变化。金属原子在高密度位错的晶粒边界或碎晶处形成晶核，并不断长大，按变形前的晶体结构形成新的均匀细小的等轴晶粒。金属的强度、硬度进一步降低，塑性、韧性显著提高，内应力完全消除，这一过程称为再结晶。
\[ T_{\text{再}} = 0.4 T_{\text{熔}} \]
式中：$T_{\text{再}}$——再结晶温度，K。

再结晶以后金属的加工硬化完全消除，并恢复良好的塑性以及冷变形前的金属耐蚀性、导电性和导磁性。

金属的加工硬化可以提高其强度和硬度，但也会给继续进行塑性变形带来困难，应加以消除。在实际生产中，如冷轧、冷压过程，常采用加热的方法使金属发生再结晶（称为再结晶退火），从而再次获得良好的塑性。

当金属在再结晶温度以上受力变形(称为热变形)时,加工硬化和再结晶同时存在。由于变形中的加工硬化随时被再结晶过程所消除,所以热变形后无加工硬化现象。

金属在热变形时能够以较小的功完成较大的变形,并可以压合铸造组织中的气孔、缩松、微裂纹等缺陷,使金属的致密度提高,获得具有较好力学性能的再结晶组织。因此,金属的塑性加工生产多采用热变形来进行。

(3) 晶粒长大

再结晶以后,如继续升高加热温度或延长保温时间,再结晶产生的细晶粒又会逐渐长大。晶粒长大是一种自发过程,通过大晶粒吞并小晶粒和晶界的迁移来实现。

晶粒长大会相应降低金属的力学性能,所以应尽量避免。

晶粒长大主要受以下两个方面因素的影响。

1) 加热温度与保温时间

温度越高,晶粒长大越快;保温时间越长,晶粒越粗大。

2) 变形程度

变形程度很小时几乎不发生再结晶;当变形程度达到 $2\% \sim 10\%$ 时,再结晶晶粒特别粗大(此变形程度称为临界变形度);超过临界变形度后,随变形量增大,再结晶晶粒细化;当变形程度超过 $90\%$ 以后,某些金属(如铁)中又会出现晶粒再次粗化的现象。

\subsection{金属的纤维组织及锻造比}

热变形可以使金属的致密度提高,细化组织,从而提高其力学性能,但同时会使金属出现流线组织,使其性能呈现异向性。

在热变形过程中,材料内部的夹杂物及其他非基体物质沿塑性变形方向所形成的流线组织,称为纤维(流线)组织。

纤维组织的明显程度与锻造比有关。锻造比通常用拔长时的变形程度来衡量,即
\[
Y = \frac{S_0}{S}
\]
式中: $Y$——锻造比; $S_0$——拔长前坯料的横截面积, $\mathrm{m}^2$; $S$——拔长后坯料的横截面积, $\mathrm{m}^2$。

锻造比的大小影响金属的力学性能和锻件质量。通常情况下,增加锻造比有利于改善金属的组织与性能,但其过大也无益。一般来说, $Y=2 \sim 5$ 时,变形金属中开始形成纤维组织,纵向(顺纤维方向)的强度、塑性和韧性增高,横向(垂直纤维方向)同类性能下降,力学性能出现各向异性; $Y>5$ 时,钢料的组织细密化程度已接近极限,力学性能不再提高,各向异性则进一步增加。因此,选择合适的锻造比十分重要。

纤维组织的稳定性很高,不会因热处理而改变,采用其他方法也无法消除,只能通过合理的锻造方法来改变纤维组织在零件中的分布方向和形状。因而,在设计和制造零件时,必须考虑纤维组织的合理分布,充分发挥其纵向性能高的优势,限制横向性能差的劣势。设计原则是使零件工作时承受的最大正应力方向与纤维方向一致,最大切应力方向与纤维方向垂直,并尽可能使纤维方向沿零件的轮廓分布而不被切断。图 8-7 所示为几种不同加工方法生产齿轮时纤维组织分布的比较。其中,图 8-7d 所示齿轮的纤维组织分布合理,使用寿命最高,材料消耗最少;图 8-7c 所示齿轮次之;图 8-7a 所示齿轮是采用轧制棒料经切削加工而成的,受力时齿根处产生的正应力垂直于纤维方向,性能最差;图 8-7b 所示的齿轮,其顺纤维方向的齿根处正应力与纤维方向重合,质量好,但竖直方向的质量较差。

\begin{figure}[H]
    \centering
    % Placeholder for Figure 8-7
    \includegraphics[page=20, viewport={67.0 548.8 319.0 680.3}, clip]{12-07572-001FigRevised.pdf}
    \caption{几种不同加工方法生产齿轮时纤维组织的分布比较}
\end{figure}

\subsection{金属的可锻性及影响因素}

金属的可锻性又称为锻造性能,是衡量材料压力加工难易程度的工艺性能。它包括塑性和变形抗力两个因素。塑性高,变形抗力小,则可锻性好;反之,可锻性差。

影响金属可锻性的因素主要包括金属的本质和变形条件两个方面。

1. 金属的本质

(1) 化学成分的影响

一般来说,纯金属的可锻性优于其合金的可锻性;合金元素的质量分数越大,成分越复杂,则金属的可锻性越差;碳钢中碳的质量分数增加使其可锻性降低。因此,低碳钢的可锻性比高碳钢好,碳钢的可锻性比合金钢好,低合金钢的可锻性比高合金钢好。

(2) 组织结构的影响

同样成分的金属在形成不同的组织结构时,其可锻性有很大差别。固溶体组织具有良好的可锻性,因而锻造时常常将钢加热至奥氏体状态;合金中化合物的增加会使其可锻性迅速下降;金属在单相状态下的可锻性优于多相状态,因为多相状态下各相的塑性不同,变形不均匀会引起内应力,甚至开裂。细晶组织金属的可锻性优于粗晶组织金属,因而锻造时要控制加热温度,避免晶粒长大。

2. 变形条件

(1) 变形温度

变形温度对材料的塑性和变形抗力影响很大。一般而言,随着温度的升高,原子的动能增加,原子间的吸引力削弱,减小了滑移阻力,从而使塑性提高、变形抗力减小,改善了金属的可锻性。热变形的变形抗力通常只有冷变形的 $1/15 \sim 1/10$,故在生产中得到广泛应用。

金属的加热应控制在一定的温度范围内,否则会产生过热、过烧以及脱碳和严重氧化等加热缺陷。过热是指由于加热温度过高或高温下保温时间过长而引起晶粒粗大的现象。过热组织的力学性能和可锻性均不好,可通过正火使晶粒细化,恢复其可锻性。过烧是指加热温度过高、接近金属的熔点时,使晶界出现氧化或熔化的现象。过烧组织的晶粒非常粗大,且晶界的氧化破坏了晶粒间的结合,使金属完全失去可锻性,是一种不可挽回的加热缺陷。

锻造时,必须合理地控制锻造温度范围,即始锻温度(开始锻造的温度)与终锻温度(停止锻造的温度)之间的间隔。在锻造过程中,随着温度的降低,工件材料的变形能力下降,变形抗力增大,下降至终锻温度时,必须停止锻造。继续锻造需重新加热,以保证材料具有足够的塑性以防止锻裂。但终锻温度不宜太高,否则无法充分利用有利的变形条件,且增加了加热火次,还容易使锻件在冷却后得到粗晶组织。

确定锻造温度范围的主要理论依据是合金相图。碳钢的始锻温度应在固相线 $ae$ 以下 $150 \sim 250\ ^{\circ}\mathrm{C}$,终锻温度约为 $800\ ^{\circ}\mathrm{C}$(图 8-8)。亚共析钢的终锻温度虽处于两相区,但仍具有足够的塑性和较小的变形抗力;对于过共析钢,在两相区停锻,是为了击碎沿晶界分布的网状二次渗碳体。

(2) 变形速度

变形速度是指单位时间内的变形程度。变形速度与可锻性的关系如图 8-9 所示。变形速度有一个临界值 $C$。低于临界值 $C$ 时,随变形速度增加,金属的变形抗力增加,塑性减小。这是由于金属的再结晶过程来不及消除金属变形所产生的加工硬化现象,残余的加工硬化作用逐渐积累使可锻性变差。当高于临界值 $C$ 时,塑性变形产生的热效应(消耗于金属塑性变形的能量一部分转化为热能,使金属的温度升高)加快了再结晶过程,使金属的塑性提高,变形抗力减小,可锻性得以改善。高速锻造便是利用这一原理来改善金属可锻性的。

\begin{figure}[H]
    \centering
    % Placeholder for Figure 8-8
    \includegraphics[page=20, viewport={340.0 552.3 475.5 746.3}, clip]{12-07572-001FigRevised.pdf}
    \caption{碳钢的锻造温度范围}
\end{figure}

\begin{figure}[H]
	\centering
	% Placeholder for Figure 8-9
	\includegraphics[page=20, viewport={47.0 388.8 194.0 514.3}, clip]{12-07572-001FigRevised.pdf}
	\caption{变形速度与可锻性的关系}
\end{figure}

但是,在普通锻压设备上锻压金属的变形速度均不可能高于临界值 $C$。因此,对于塑性差的材料(如高合金钢)或大型锻件,宜采用较小的变形速度(如在压力机上成形),以防锻裂坯料。

(3) 应力状态

变形方法不同,在金属中产生的应力状态也不同。即使同一种变形方式,金属内部不同位置的应力状态也可能不同。例如,金属在挤压时三向受压(图 8-10a),表现出较高的塑性和较大的变形抗力;拉拔时两向受压,一向受拉(图 8-10b),与挤压相比,表现出较低的塑性和较小的变形抗力;(平砧)镦粗时(图 8-10c),坯料内部处于三向压应力状态,但侧表面层在水平方向却处于拉应力状态,因而工件侧表面上容易产生垂直方向的裂纹。

\begin{figure}[H]
    \centering
    % Placeholder for Figure 8-10
    \includegraphics[page=20, viewport={210.5 389.3 491.5 477.8}, clip]{12-07572-001FigRevised.pdf}
    \caption{金属变形时的应力状态}
\end{figure}

金属三向受压时的塑性最好,出现拉应力则使塑性变差。这是因为压应力阻碍了微裂纹的产生和扩展,而金属处于拉应力状态时,内部缺陷处会产生应力集中,使缺陷易于扩展从而导致金属破坏。因此,选择变形方法时,对于塑性好的金属,变形时出现拉应力是有利的,可减少变形时的能量消耗;而对于塑性差的金属材料,应避免在拉应力状态下成形,尽量采用在三向压应力下变形。随着变形抗力的增大,加工设备的吨位需要相应增加。

坯料的表面状况对材料的塑性也有影响,特别在冷变形时尤为显著。坯料表面粗糙或有划痕、微裂纹和粗大夹杂物等,都会在变形过程中产生应力集中而引起开裂。故加工前应对坯料进行清理和消除缺陷。

\section{常用锻造方法}

\subsection{自由锻}

自由锻是指用简单的通用性工具,或在锻造设备的上、下砧之间,利用锤头的冲击力或上砧的静压力直接使坯料变形而获得锻件的方法。

常用自由锻设备有空气锤、蒸汽-空气自由锻锤和水压机。

空气锤由自身携带的电动机直接驱动,锤击能量小,只能锻造 $100\ \mathrm{kg}$ 以下的小型锻件。蒸汽-空气自由锻锤以蒸汽或压缩空气作为驱动力,以实现锤头的连续打击动作。蒸汽-空气自由锻锤可以锻造中型或较大型锻件。锻锤的吨位以落下部分(活塞、锤杆和上砧)的质量来表示。水压机主要由立柱、横梁、工作缸、回程缸和操作系统所组成(图 8-11)。水压机工作时,依靠工作缸通入高压水($20 \sim 40\ \mathrm{MPa}$)推动工作柱塞带动活动横梁和上砧向下运动,使坯料产生变形。回程时,高压水通入回程缸,通过回程柱塞和拉杆将活动横梁拉起。水压机的吨位以所能产生的最大压力来表示。水压机工作时的变形速度较慢,有利于改善坯料的可锻性,一般用于碳钢、合金钢等大型锻件的单件、小批生产。

自由锻所用工具和设备简单,通用性好,工艺灵活,成本低。锻件质量可以从数十克到二三百吨,对于大型锻件如轧辊、发电机转子、主轴、汽轮机叶轮、大型多拐曲轴等,大多采用自由锻方法成形。因此,自由锻在重型机械制造中占有重要地位。但是,自由锻锻件精度低,加工余量大,生产率低,故主要用于单件、小批生产。

\begin{figure}[H]
    \centering
    % Placeholder for Figure 8-11
    \includegraphics[page=20, viewport={99.5 37.3 440.0 353.3}, clip]{12-07572-001FigRevised.pdf}
    \caption{水压机的结构示意图}
\end{figure}

1. 自由锻基本工序

自由锻工序分为基本工序、辅助工序和修整工序。基本工序有镦粗、拔长、冲孔、弯曲、切割、错移和扭转;辅助工序有压钳口、倒棱和压痕等;修整工序有矫正、滚圆、平整等。表 8-2 列出了常用自由锻基本工序的定义及应用。

\begin{longtable}{|c|m{1cm}|m{1.5cm}|m{8.5cm}|m{1.5cm}|}
\caption{常用自由锻基本工序的定义及应用} \\
\hline
\multicolumn{2}{|c|}{工序名称} & 定义 & 图例 & 应用 \\
\hline
		镦粗 & 平砧镦粗(图a); 局部镦粗(图b) & 镦粗:使坯料的高度减小,横截面积增大; 平砧镦粗:坯粒在上下平砧间进行的镦粗; 局部镦粗:对坯料上某一部分进行镦粗 & \includegraphics[page=30, viewport={110.5 653.8 250.5 746.3}, clip]{12-07572-001FigRevised.pdf} & 制造盘类锻件,如齿轮坯、圆盘等; 作为冲孔前的准备工序; 增大锻造比 \\
		\hline
		拔长 & 平砧拔长(图a); 芯棒拔长(图b) & 拔长:使坯料的横截面积减小、长度增加; 平砧拔长:坯料在上下平砧间进行的拔长; 芯棒拔长:减小空心坯料的外径和壁厚,增加其长度 & \includegraphics[page=30, viewport={268.0 652.8 429.5 742.3}, clip]{12-07572-001FigRevised.pdf} & 制造长而截面小的锻件,如轴、连杆、曲轴等; 制造长轴类空心件、圆环类零件,如炮筒、圆环、套筒等 \\
		\hline
		冲孔 & 实心冲子冲孔(图a); 空心冲子冲孔(图b) 芯轴扩孔(图c) & 冲孔:在坯料上冲出通孔或不通孔; 芯轴扩孔:减小空心坯料的壁厚,增加其内径和外径 &\includegraphics[page=30, viewport={15.5 385.8 259.0 624.3}, clip]{12-07572-001FigRevised.pdf} & 制造空心锻件,如齿轮毛坯、圆环、套筒等; 锻件质量要求高的大型工件,可用空心冲孔去掉质量较低的铸锭中心部分 \\
		\hline
\end{longtable}

2. 自由锻工艺规程的制订

自由锻工艺规程是指导锻件生产、管理和质量检验的依据,其主要内容和步骤如下。

(1) 绘制锻件图

锻件图是以零件图为基础,考虑锻造工艺特点绘制而成的。

为了简化锻件形状、便于锻造而增加的一部分金属称为敷料(余块)。当零件上带有难以直接锻出的凹槽、台阶、凸肩、小孔时,均需添加敷料(图 8-12a)。

锻件上凡需切削加工的表面,应留有加工余量。零件的基本尺寸加上加工余量即锻件的基本尺寸。锻件的实际尺寸与基本尺寸之间所允许的偏差,称为锻件公差。加工余量与锻件公差通常根据有关手册和实际生产条件确定。

\begin{figure}[H]
    \centering
    % Placeholder for Figure 8-12
    \includegraphics[page=21, viewport={29.5 553.8 284.0 743.3}, clip]{12-07572-001FigRevised.pdf}
    \caption{锻件图的表示方法}
\end{figure}

确定了加工余量、公差和敷料后,便可绘出锻件图。锻件图的外形用粗实线表示,零件的外形用双点画线表示。锻件的基本尺寸与公差标注在尺寸线上面,零件的尺寸标注在尺寸线下面的括号内,如图 8-12b 所示。

(2) 选择锻造工序

确定锻造工序的依据是锻件的形状、尺寸、技术要求 and 生产数量等。自由锻件的分类及基本工序方案见表 8-3。

\begin{longtable}[h]{|c|c|c|m{4cm}|m{3cm}|}
    
    \caption{自由锻件分类及基本工序方案}\\
   
        \hline
        序号 & 类别 & 图例 & 基本工序方案 & 实例 \\
        \hline
        1 & 饼块类 & \includegraphics[page=30, viewport={271.5 384.8 337.5 442.3}, clip]{12-07572-001FigRevised.pdf} & 镦粗或局部镦粗 & 圆盘、齿轮、模块、锤头等 \\
        \hline
        2 & 轴杆类 &\includegraphics[page=30, viewport={351.5 385.3 427.5 436.8}, clip]{12-07572-001FigRevised.pdf} & 拔长 \newline 镦粗—拔长(增大锻造比) \newline 局部镦粗—拔长(截面相差较大的阶梯轴) & 传动轴、主轴、连杆类零件 \\
        \hline
        3 & 空心类 & \includegraphics[page=30, viewport={453.0 385.3 528.0 422.8}, clip]{12-07572-001FigRevised.pdf} & 镦粗—冲孔 \newline 镦粗—冲孔—扩孔 \newline 镦粗—冲孔—芯棒拔长 & 圆环、法兰、齿圈、套筒、空心轴等 \\
        \hline
        4 & 弯曲类 &\includegraphics[page=30, viewport={17.5 268.8 91.0 309.8}, clip]{12-07572-001FigRevised.pdf} & 轴杆类锻件工序—弯曲 & 吊钩、弯杆、轴瓦盖等 \\
        \hline
        5 & 曲轴类 &\includegraphics[page=30, viewport={112.5 261.3 244.5 312.3}, clip]{12-07572-001FigRevised.pdf} & 拔长—错移(单拐曲轴) \newline 拔长—错移—扭转 & 曲轴、偏心轴等 \\
        \hline
   \end{longtable}

(3) 确定坯料质量和尺寸

坯料有铸锭和型材两种,前者用于大、中型锻件,后者用于中、小型锻件。

坯料的质量 $m_{\text{坯}}(\mathrm{kg})$ 为锻件的质量 $m_{\text{锻}}(\mathrm{kg})$ 与锻造时的各种损耗质量 $m_{\text{损}}(\mathrm{kg})$(如加热时的烧损质量、冲孔时芯料的质量和锻造过程中切除的料头质量等)之和,即
\[
m_{\text{坯}} = m_{\text{锻}} + m_{\text{损}}
\]
坯料尺寸依据锻造工序和变形程度(锻造比)来确定。

采用拔长方法锻造时,
\[
S_{\text{坯}} \geqslant Y S_{\text{锻}}
\]
式中: $S_{\text{坯}}$——坯料最大横截面积, $\mathrm{m}^2$; $S_{\text{锻}}$——锻件最大横截面积, $\mathrm{m}^2$; $Y$——拔长时的锻造比。

对于碳钢钢锭, $Y$ 一般不小于 $2.5 \sim 3$;对于合金结构钢钢锭, $Y$ 一般不小于 $3 \sim 4$。

平砧镦粗时,为避免产生纵向弯曲现象和下料困难,坯料的高径比($H_0/D_0$)应大于 $1.25$,但不得超过 $2.5$。

由于坯料的质量已知,可先计算出坯料的体积,再确定坯料的截面尺寸(直径或边长),最后确定坯料的长度。

(4) 选择锻造设备

锻造设备的选择,应根据坯料的种类、质量以及锻造基本工序、设备的锻造能力等因素,并结合工厂现有设备条件综合确定。

表 8-4 为盘类典型锻件齿轮坯的自由锻工艺过程示例,包括了锻造工艺规程的内容,可用于指导自由锻生产。

\begin{longtable}[h]{|m{2cm}|m{3cm}|m{2.2cm}|m{3cm}|}
    
    \caption{齿轮坯的自由锻工艺过程}\\
    
        \hline
        锻件名称 & 齿轮坯 & 工艺类别 & 自由锻 \\
        \hline
        材料 & 18CrMnTi & 设备 & $65\ \mathrm{kg}$ 空气锤 \\
        \hline
        加热火次 & 1 & 锻造温度范围 & $1150 \sim 800\ ^{\circ}\mathrm{C}$ \\
        \hline
        \multicolumn{2}{|c|}{锻件图} & \multicolumn{2}{c|}{坯料图} \\
        \hline
        \multicolumn{2}{|c|}{\includegraphics[page=30, viewport={266.5 262.3 390.5 352.3}, clip]{12-07572-001FigRevised.pdf}} & \multicolumn{2}{c|}{\includegraphics[page=30, viewport={409.0 264.8 525.0 325.8}, clip]{12-07572-001FigRevised.pdf}} \\
        \hline
        序号 & 工序名称 & \multicolumn{2}{c|}{工序简图}\\
        \hline
        1 & 镦粗 & \multicolumn{2}{c|}{\includegraphics[page=30, viewport={8.0 76.8 72.5 139.3}, clip]{12-07572-001FigRevised.pdf}} \\
        \hline
        2 & 冲孔 & \multicolumn{2}{c|}{\includegraphics[page=30, viewport={95.0 79.8 140.0 136.3}, clip]{12-07572-001FigRevised.pdf}} \\
        \hline
        3 & 修整外圆 & \multicolumn{2}{c|}{\includegraphics[page=30, viewport={150.0 76.3 214.0 135.8}, clip]{12-07572-001FigRevised.pdf}} \\
        \hline
        4 & 修整平面 & \multicolumn{2}{c|}{\includegraphics[page=30, viewport={224.5 77.3 297.0 133.3}, clip]{12-07572-001FigRevised.pdf}} \\
        \hline
    \end{longtable}

\subsection{模锻}

模锻是金属在外力作用下产生塑性变形并充满模腔而获得锻件的方法。常用模锻设备有模锻锤、热模锻压力机、平锻机和摩擦压力机等。

与自由锻相比，模锻件尺寸精度高，机械加工余量小，锻件的纤维组织分布更为合理，可进一步提高零件的使用寿命。模锻生产率高，操作简单，容易实现机械化和自动化。但设备投资大，锻模成本高，生产准备周期长，且模锻件的质量受到模锻设备吨位的限制，因而适用于中小型锻件（一般$<150\text{ kg}$）的成批和大量生产。典型模锻件如图8-13所示。

按所用设备类型的不同，模锻可分为锤上模锻、压力机上模锻、其他设备模锻等。

1. 锤上模锻

锤上模锻所用的设备有蒸汽-空气模锻锤、无砧座锤（对击锤）和高速锤等。一般工厂主要采用蒸汽-空气模锻锤。其工作原理与蒸汽-空气自由锻锤基本相同。但由于模锻时受力大，锻件精度要求高，故模锻设备需刚性好，导向精度高。锤上模锻的锤头与导轨之间的间隙比自由锻小，以保证上模、下模对准；机架直接与砧座相连，以便提高打击刚度和冲击效率。模锻锤的吨位一般为$1\sim 16\text{ t}$。

(1) 锻模结构

锻模由带有燕尾的上模和下模组成（图8-14），在上模和下模内均制有相应的模腔。燕尾和紧固楔铁分别安装在锤头和模座上；键槽与键配合，起定位作用，防止锻模前后移动；起重孔是为了安装锻模方便而设置的。

\begin{figure}[H]
    \centering
    \includegraphics[page=21, viewport={303.0 552.3 519.0 701.8}, clip]{12-07572-001FigRevised.pdf}
    \caption{典型模锻件}
    \label{fig:8-13}
\end{figure}

\begin{figure}[H]
    \centering
    % Labels: 紧固楔铁, 锤头, 上模, 模腔, 飞边槽, 分模面, 下模, 模座, 锁扣, 飞边槽, 模锻模腔, 制坯模腔, 制坯模腔, 起重孔, 燕尾
    \includegraphics[page=21, viewport={118.0 397.3 424.5 518.8}, clip]{12-07572-001FigRevised.pdf}
    \caption{锻模结构}
    \label{fig:8-14}
\end{figure}

(2) 锻模模腔的分类

锻模模腔按其功用的不同，分为制坯模腔和模锻模腔两大类（表8-5）。模锻件的复杂程度不同，其所需模腔的数量也不等，锻模可设计为单腔锻模（图8-14a）或多腔锻模（图8-14b），一个锻模一般不超过7个模腔。图8-15所示为弯曲连杆的模锻过程，所使用的锻模为多腔锻模。

\begin{longtable}{|c|m{1cm}|c|m{9.5cm}|m{1.8cm}|}
\caption{锻模模腔分类及用途}\\
\hline
类别 & 性质 & 模腔名称 & 简图 & 用途 \\
\hline
\multirow{4}{*}{\makecell[c]{制\\坯\\模\\腔}} & \multirow{4}{1cm}{改变坯料横截面积和形状，使金属合理分布并基本接近锻件形状，以利于金属顺利充满模锻模腔} & 拔长模腔 & \includegraphics[page=30, viewport={305.5 76.8 533.5 236.8}, clip]{12-07572-001FigRevised.pdf} & 减小坯料某部分的横截面积，增加其长度，兼有去除氧化皮的作用。主要用于长轴类锻件制坯 \\
\cline{3-5}
& & 滚压模腔 & \includegraphics[page=31, viewport={39.0 587.3 308.5 747.3}, clip]{12-07572-001FigRevised.pdf} & 减小坯料某部分的横截面积，增大另一部分的横截面积，使坯料沿轴线的形状更接近锻件。主要用于某些变截面长轴类锻件的制坯 \\
\cline{3-5}
& & 弯曲模腔 &\includegraphics[page=31, viewport={334.5 585.8 504.0 723.8}, clip]{12-07572-001FigRevised.pdf} & 改变坯料轴线形状，以符合锻件水平投影形状。主要用于具有弯曲轴线的锻件的制坯 \\
\cline{3-5}
& & 切断模腔 & \includegraphics[page=31, viewport={172.0 416.8 372.0 554.8}, clip]{12-07572-001FigRevised.pdf} & 当一块坯料锻造两个或多个锻件时，将已锻好的锻件从坯料上切下 \\
\hline
\multirow{2}{*}{\makecell[c]{模\\锻\\模\\腔}} & \multirow{2}{1cm}{使坯料进一步变形，直至最终形成锻件} & 预锻模腔 & \multirow{2}{3.5cm}{\includegraphics[page=31, viewport={43.5 227.3 211.5 384.8}, clip]{12-07572-001FigRevised.pdf}} & 获得与终锻锻件相近的形状，以利于锻件在终锻模腔中清晰成形，提高锻件质量，并减小终锻模腔的磨损，延长其使用寿命。无飞边槽。主要用于形状复杂的锻件 \\
\cline{3-3}\cline{5-5}
& & 终锻模腔 & & 最终获得所需形状和尺寸的锻件。有飞边槽，其作用是增加坯料成形时所受到的三向压应力作用，促使金属充满模腔和容纳多余金属 \\
\hline
\end{longtable}

\begin{figure}[H]
    \centering
    % Labels: 滚压模腔, 终锻模腔, 预锻模腔, 弯曲模腔, 拔长模腔, 原始料坯, 拔长, 滚压, 弯曲, 预锻, 终锻, 毛边, 锻件
    \includegraphics[page=21, viewport={16.5 92.3 269.5 369.3}, clip]{12-07572-001FigRevised.pdf}
    \caption{弯曲连杆的模锻过程}
    \label{fig:8-15}
\end{figure}

(3) 模锻件图的绘制

模锻件图是确定模锻工艺、设计和制造锻模以及检验锻件的主要依据。绘制模锻件图时，应主要考虑以下问题。


1)选择分模面
    
    分模面是指上、下（锻）模在锻件上的分界面。确定分模面一般按以下原则。
 
     \textcircled{1}应保证锻件从模腔中顺利取出，故分模面一般应选取在锻件最大尺寸的截面上。
     
     \textcircled{2}应使分模面处上、下模模腔外形一致，以便能及时发现错模。
     
     \textcircled{3}应使模腔浅而宽，以利于金属充满模腔。
     
     \textcircled{4}应保证锻件上所加敷料最少。

    根据上述原则，图8-16中的$d-d$作分模面最为合适。
    
2)确定加工余量、公差、余块和连皮
    
    模锻件的加工余量一般为$1\sim 4\text{ mm}$，公差一般取$\pm 0.3\sim 3\text{ mm}$。具体可查阅相关手册确定。模锻件均为批量生产，应尽量减少或不加余块，直径小于$30\text{ mm}$的孔一般不予锻出。
    
    模锻时不能直接锻出通孔，在该部位留有一层厚度为$\delta$的较薄金属，称为连皮（图8-17），在锻造后与飞边一同切除。

\begin{figure}[H]
    \centering
    % Labels: a-a面, b-b面, c-c面, d-d面
    \includegraphics[page=21, viewport={329.5 95.8 528.5 222.3}, clip]{12-07572-001FigRevised.pdf}
    \caption{分模面的选择}
    \label{fig:8-16}
\end{figure}



3)确定模锻斜度和圆角半径
    
    模锻件平行于锤击方向的侧面，应设计成一定斜度，以便顺利取出锻件。外斜度$\alpha$（锻件外壁上的斜度）一般取$5^\circ \sim 10^\circ$，内斜度$\beta$（锻件内壁上的斜度）一般取$7^\circ \sim 15^\circ$。
    
    模锻件所有转角处均应设计成圆角，以便使金属易于在模腔内流动，保持金属纤维组织的连续性，提高锻件质量和模具寿命。一般外圆角半径$r$取$1.5\sim 12\text{ mm}$，内圆角半径$R$取$(3\sim 4)r$。
    
    图8-18为齿轮坯模锻件图，图中双点画线为零件的轮廓外形。
    
\begin{figure}[H]
    \centering
    \includegraphics[page=22, viewport={79.0 607.8 238.5 686.8}, clip]{12-07572-001FigRevised.pdf}
    \caption{模锻斜度、圆角半径和连皮}
    \label{fig:8-17}
\end{figure}

\begin{figure}[H]
    \centering
    \includegraphics[page=22, viewport={255.5 608.3 472.5 746.3}, clip]{12-07572-001FigRevised.pdf}
    \caption{齿轮坯模锻件图}
    \label{fig:8-18}
\end{figure}


(4) 变形工步的选择

首先根据锻件图的复杂程度确定变形工步，再根据已确定的工步设计制坯模腔、预锻模腔和终锻模腔。表8-6为锤上模锻件分类和变形工步示例。

\begin{longtable}[h]{|c|c|c|}

\caption{锤上模锻件分类和变形工步示例}\\



\hline
模锻件分类 & 变形工步示例 & 主要变形工步 \\
\hline
盘类 &\includegraphics[page=31, viewport={227.0 224.8 344.0 273.3}, clip]{12-07572-001FigRevised.pdf} & 镦粗（预锻）、终锻 \\
\hline
直轴类 & \includegraphics[page=31, viewport={365.0 223.3 495.0 318.3}, clip]{12-07572-001FigRevised.pdf} & 拔长、滚压（预锻）、终锻 \\
\hline
弯轴类 &\includegraphics[page=31, viewport={64.5 106.3 188.5 187.8}, clip]{12-07572-001FigRevised.pdf} & 拔长、滚压、弯曲（预锻）、终锻 \\
\hline
叉类 &\includegraphics[page=31, viewport={206.5 107.3 320.0 188.8}, clip]{12-07572-001FigRevised.pdf} & 拔长、滚压、预锻、终锻 \\
\hline
枝芽类 &\includegraphics[page=31, viewport={338.0 107.3 476.0 196.3}, clip]{12-07572-001FigRevised.pdf} & 拔长、滚压、成形（预锻）、终锻 \\
\hline

\end{longtable}

坯料在锻模内制成锻件后，还需要经过一系列修整工序。如对于带有飞边及连皮的锻件，需要在压力机上利用切边模和冲孔模将其切除（图8-19），以保证和提高锻件质量。

目前，锤上模锻是我国应用最多的一种模锻方法，但由于其工作时振动和噪声大，蒸汽效率低，难以实现较高程度的操作机械化；完成一个变形工步需要多次锤击，生产率仍不是很高，故在大批量生产中有逐渐被压力机上模锻所取代的趋势。

2. 压力机上模锻

(1) 曲柄压力机上模锻

曲柄压力机的结构如图8-20所示。电动机通过带轮和齿轮副的传动，带动曲柄连杆机构运动，从而使滑块作上下往复运动。锻模分别安装在滑块下端和工作台上。

与锤上模锻相比，曲柄压力机上模锻主要有以下优点。

1)变形力为静压力，坯料的变形速度较低，这对于成形低塑性材料较为有利，如可在曲柄压力机上成形耐热合金和镁合金等。

2)锻造时滑块行程不变，坯料变形在一次行程内完成，生产率高。

3)滑块运动精度高，并设有上下顶出装置，能使锻件自动脱模，便于实现机械化和自动化。


曲柄压力机模锻的缺点是滑块行程和压力不能随意调节，不宜进行拔长、滚压等操作；设备复杂、费用高，适用于大批量生产。

(2) 平锻机上模锻

平锻机相当于卧式曲柄压力机（图8-21）。它没有工作台，锻模由固定（凹）模、活动（凹）模和凸模三部分组成，具有两个相互垂直的分模面。当活动（凹）模与固定（凹）模合模时，便夹紧坯料，主滑块带动凸模进行模锻成形。

\begin{figure}[H]
    \centering
    \includegraphics[page=22, viewport={10.5 374.8 193.5 469.3}, clip]{12-07572-001FigRevised.pdf}
    \caption{切边模和冲孔模}
    \label{fig:8-19}
\end{figure}

\begin{figure}[H]
    \centering
    \includegraphics[page=22, viewport={211.0 372.8 533.0 576.8}, clip]{12-07572-001FigRevised.pdf}
    \caption{曲柄压力机的结构}
    \label{fig:8-20}
\end{figure}

\begin{figure}[H]
    \centering
    \includegraphics[page=22, viewport={82.5 198.8 416.5 343.8}, clip]{12-07572-001FigRevised.pdf}
    \caption{平锻机示意图}
    \label{fig:8-21}
\end{figure}

平锻机上模锻主要有以下特点。

1)坯料多是棒料和管材，可锻造出曲柄压力机模锻所不能锻造的长杆类锻件，并能锻出通孔（图8-22）。


\begin{figure}[H]
    \centering
    \includegraphics[page=22, viewport={155.5 82.8 391.0 167.3}, clip]{12-07572-001FigRevised.pdf}
    \caption{平锻机模锻件}
    \label{fig:8-22}
\end{figure}



2)锻模有两个分模面，可以锻出其他设备上无法成形的侧面带有凸台和凹槽的锻件。锻件无飞边、精度高。


平锻机模锻也是一种高效率、高质量、容易实现机械化和自动化的模锻方法。但平锻机造价高，投资大，仅适用于大批量生产。

(3) 摩擦压力机上模锻

摩擦压力机（图8-23）是靠飞轮旋转所积蓄的能量转化为金属的变形能而进行锻造的。电动机经带轮、摩擦盘、飞轮和螺杆带动滑块作上下往复运动，操纵机构控制左、右摩擦盘分别与飞轮接触，利用摩擦力改变飞轮转向。

\begin{figure}[H]
    \centering
    \includegraphics[page=23, viewport={72.5 541.3 467.5 752.3}, clip]{12-07572-001FigRevised.pdf}
    \caption{摩擦压力机}
    \label{fig:8-23}
\end{figure}

摩擦压力机的行程速度介于模锻锤和曲柄压力机之间，滑块行程和打击能量均可自由调节，坯料在一个模腔内可以多次锤击，能够完成镦粗、成形、弯曲、预锻等成形工序和矫正、修整等后续工序。

摩擦压力机构造简单，投资费用少，工艺适应性广，但传动效率低，一般只能进行单模腔模锻，对于形状复杂的锻件，需要在自由锻设备或其他设备上制坯。摩擦压力机上模锻广泛用于中、小批生产的中、小型模锻件，以及某些低塑性合金锻件的生产。

3. 其他设备模锻

(1) 胎模锻

胎模锻是在自由锻设备上使用胎模来生产模锻件的方法。通常，用自由锻方法使坯料初步成形，再在胎模内终锻成形。

胎模的结构形式很多，常用胎模结构如图8-24所示。扣模主要用于非回转体锻件的局部或整体成形；筒模主要用于锻造法兰盘、齿轮坯等回转体盘类零件；合模由上、下模两部分组成，主要用于锻造形状较复杂的非回转体锻件。

胎模锻的特点介于自由锻与锤上模锻之间，比自由锻生产率高，锻件质量较好，锻模简单，生产准备周期短，广泛用于中、小批的小型锻件的生产。

(2) 精密模锻

精密模锻是指在普通锻造设备上锻造高精度锻件的方法。其主要工艺特点是使用两套不

\begin{figure}[H]
    \centering
    % Placeholder for Figure 8-24
    \includegraphics[page=23, viewport={87.0 298.3 466.0 513.3}, clip]{12-07572-001FigRevised.pdf}
    \caption{常用胎模的结构}
    \label{fig:8-24}
\end{figure}

同精度的锻模。先使用普通锻模锻造,留有 $0.1 \sim 1.2 \text{ mm}$ 的精锻余量,然后切下飞边并进行酸洗,再使用高精度锻模直接锻造出所要求的产品零件。例如,在摩擦压力机和曲柄压力机上精锻锥齿轮、汽轮机叶片等形状复杂的零件。在精密模锻过程中,要采用无氧化和少氧化的加热方法。

提高锻件精度的另一条途径是在中温(碳钢的始锻温度一般为 $600 \sim 875 \text{ }^\circ\text{C}$)或室温下进行精密锻造。但这种锻造方法只能锻造小型钢锻件或非铁金属锻件。精密模锻件的精度较高,一般不需切削加工或只进行少量切削加工便可投入使用,是一种先进的锻造方法。但由于模具制造复杂,对坯料尺寸和加热质量要求较高,精密模锻只适宜于大批量生产。

\subsection{锻造方法的选择}

生产同一种锻件可以选用不同的锻造方法。例如,汽车发动机的连杆,可以在模锻锤上和各种压力机上模锻,也可以采用胎模锻造。即使在同一种锻压设备上,也可采用不同的工艺方案。如在模锻锤上模锻连杆,可以用拔长、滚压制坯,再进行预锻、终锻;也可以先在其他设备上制坯,之后在模锻锤上终锻。因此,需对各种工艺方案进行比较分析,在满足性能和质量要求的前提下,应选用生产成本低,生产效率高的方案。

锻件的成本由材料、模具、人员工资、直接和间接的管理费用、能耗费用等组成。各部分比例随工艺方案的合理性和管理水平而异,工艺方案的确定应以生产批量为重要依据(图 8-25),并综合考虑各种因素,合理选择锻造方法(表 8-7)。

\begin{figure}[H]
    \centering
    % Placeholder for Figure 8-25
    \includegraphics[page=23, viewport={170.0 121.8 369.5 271.8}, clip]{12-07572-001FigRevised.pdf}
    \caption{生产批量对锻件成本的影响}
    \label{fig:8-25}
\end{figure}

\begin{longtable}[h]{|m{1cm}|m{1cm}|m{2cm}|m{1.5cm}|m{2cm}|m{0.5cm}|m{2cm}|m{0.5cm}|m{0.5cm}|m{0.5cm}|m{1cm}|}
    
    \caption{常用锻造方法的特点和应用比较}
   
\footnotesize\setlength{\tabcolsep}{2pt}\\
   
    \hline
    加工方法 & 具体方法 & 使用设备 & 锻造力性质 & 应用范围 & 生产率 & 模具特点 & 模具寿命 & 机械化与自动化 & 劳动条件 & 对环境影响 \\
    \hline
    \multirow{3}{*}{自由锻} &  & 空气锤 & 冲击力 & \multirow{3}{2.5cm}{小型锻件,单件、小批生产中型锻件,单件、小批生产大型锻件} & \multirow{3}{*}{低} & \multirow{3}{*}{无模具} & \multirow{3}{*}{} & \multirow{3}{*}{难} & \multirow{3}{*}{差} & \multirow{3}{*}{\makecell[l]{振动和\\噪声大}} \\
    &  & 蒸汽-空气自由锻锤 & 冲击力 & & & & & & & \\
    &  & 水压机 & 静压力 & & & & & & & \\
    \hline
    \multirow{5}{*}{模锻} & 锤上模锻 & \makecell[l]{蒸汽-空气模\\锻锤\\无砧座锤\\高速锤} & 冲击力 & \makecell[l]{中、小型锻\\件,大批量生\\产各种类型模\\锻件} & 高 & \makecell[l]{锻模固定在\\锤头和模座\\上,模膛复杂,\\造价高} & 中 & 较难 & 较差 & \makecell[l]{振动和\\噪声大} \\
    \cline{2-11}
    & 曲柄压力机上模锻 & \makecell[l]{热模锻压力机\\曲柄压力机} & 静压力 & \makecell[l]{中、小型锻\\件,大批量生\\产,不适宜进\\行拔长和滚压\\工序} & 高 & \makecell[l]{组合模具,\\有导柱、导套\\和顶出装置} & 较高 & 易 & 好 & 较小 \\
    \cline{2-11}
    & 平锻机上模锻 & 平锻机 & 静压力 & \makecell[l]{中、小型锻\\件,大批量生\\产,适合锻造\\法兰轴和带孔\\的模锻件} & 高 & \makecell[l]{三块模组\\成,有两个分\\模面,可锻出\\侧面带凸台、\\凹槽的锻件} & 较高 & 较易 & 较好 & 较小 \\
    \cline{2-11}
    & 摩擦压力机上模锻 & 摩擦压力机 & \makecell[l]{介于\\冲击力\\与静压力\\之间} & \makecell[l]{小型锻件,\\中、小批生产,\\可进行精密\\模锻} & 较高 & 一般为单模膛 & 较高 & 较易 & 好 & 较小 \\
    \cline{2-11}
    & 胎模锻 & \makecell[l]{空气锤\\蒸汽-空气自\\由锻锤} & 冲击力 & \makecell[l]{中、小型锻\\件,中、小批\\生产} & 较高 & \makecell[l]{模具简单,\\且不固定在设\\备上,取换\\方便} & 较低 & 较易 & 差 & \makecell[l]{振动和\\噪声大} \\
    \hline
    \end{longtable}

\section{板料冲压}

板料冲压是利用冲模在压力机上对板料施加压力使其变形或分离,从而获得一定形状、尺寸的零件的加工方法。板料冲压通常在常温下进行,又称为冷冲压,只有当板厚大于 $8 \sim 10 \text{ mm}$ 时,才采用热冲压。

冲压加工的应用范围广泛,既适用于金属材料,也适用于非金属材料;既可加工仪表上的小型制件,也可加工汽车覆盖件等大型制件,在汽车、拖拉机、电气、航空、仪表及日常生活用品等制造行业中均占有重要地位。

板料冲压具有下列特点。

1) 生产率高,操作简单,便于实现机械化和自动化。

2) 产品的尺寸精度和表面质量较高,互换性好,一般不需进一步加工。

3) 可冲制形状复杂的零件,废料少,材料利用率高。

但是,板料冲压的冲模制造复杂,成本高,只有在大批量生产的条件下才能显示出优越性。

冲压设备主要有剪床和冲床两大类。剪床(亦称为剪板机)的用途是将板料按要求切成一定宽度的条料,供下一步冲压使用。冲床(亦称为曲柄压力机)是冲压成形的基本设备,可用于切断、落料、冲孔、弯曲、拉深和其他冲压工序。

常用小型冲床的结构如图 8-26 所示。电动机通过减速机构带动曲柄连杆机构运动,从而使固定在滑块上的上模作上下往复运动,与下模配合,完成各种冲压工序。大批量生产时,常采用多工位自动冲床,生产率很高。

\begin{figure}[H]
    \centering
    % Placeholder for Figure 8-26
    \includegraphics[page=24, viewport={90.5 515.8 452.5 750.8}, clip]{12-07572-001FigRevised.pdf}
    \caption{常用小型冲床的结构}
    \label{fig:8-26}
\end{figure}

\subsection{板料冲压的基本工序}

板料冲压的基本工序可分为分离工序和成形(或变形)工序两大类。

分离工序是使冲压件与板料沿所要求的轮廓线相分离的工序,如剪切、落料和冲孔等;成形工序是使板料产生塑性变形而不破裂的工序,如弯曲、拉深、卷边、翻边和胀形等(表 8-8)。

\begin{longtable}{|l|l|m{2.5cm}|m{6.5cm}|m{3.5cm}|}
    \caption{冲压基本工序分类}\\
    \hline
    \multicolumn{2}{|c|}{工件名称} & 定义 & 简图 & 应用举例 \\
    \hline
    \multirow{3}{*}{分离工序} & 剪切 & 利用剪床或冲模,沿不封闭的曲线或直线切断 &\includegraphics[page=31, viewport={14.5 17.8 150.5 81.8}, clip]{12-07572-001FigRevised.pdf}& 用于下料或加工形状简单的平板零件,如冲制变压器的矽钢芯片 \\
    \cline{2-5}
    & 落料 & 利用冲模沿封闭轮廓曲线或直线将板料(坯料)分离,冲下部分是成品,余下部分为废料 &\includegraphics[page=31, viewport={169.0 20.8 340.5 89.8}, clip]{12-07572-001FigRevised.pdf} & 用于下料或直接冲出工件,如汽水瓶扳头、垫片等 \\
    \cline{2-5}
    & 冲孔 & 利用冲模沿封闭轮廓曲线或直线将板料分离,冲下部分是废料,余下部分为成品 & \includegraphics[page=31, viewport={357.5 20.3 524.0 87.3}, clip]{12-07572-001FigRevised.pdf} & 用于中间工序或冲制带孔零件,如冲制垫圈孔、电气箱百叶窗等 \\
    \hline
    \multirow{5}{*}{成形工序} & 弯曲 & 利用冲模或折弯机,将平直的板料弯成一定的形状 & \includegraphics[page=32, viewport={42.0 681.3 210.0 740.3}, clip]{12-07572-001FigRevised.pdf} & 用于生产板材、角钢料和各种板料箱柜的边框等 \\
    \cline{2-5}
    & 拉深 & 利用冲模将板料加工成中空形状,壁厚基本不变或局部变薄 & \includegraphics[page=32, viewport={224.5 685.3 322.5 751.3}, clip]{12-07572-001FigRevised.pdf} & 用于生产各种金属日用品(如碗、锅、盆、易拉罐身)和汽车油箱等 \\
    \cline{2-5}
    & 翻边 & 利用冲模在带孔工件上用扩孔的方法获得凸缘或把边缘按曲线或圆弧弯成竖直的边缘 & \includegraphics[page=32, viewport={347.0 680.8 501.5 745.8}, clip]{12-07572-001FigRevised.pdf} & 用于增加冲制件的强度或美观性,如汽车覆盖件的周边一般应进行翻边 \\
    \cline{2-5}
    & 卷边 & 利用冲模或旋压法,将工件竖直的边缘翻卷 &\includegraphics[page=32, viewport={87.0 572.8 269.5 653.8}, clip]{12-07572-001FigRevised.pdf} & 用于增加冲制件的强度或美观性,如金属锅具、盆的开口边缘一般应进行卷边 \\
    \hline
    & 胀形 & 利用冲模或内旋压法,使中空坯料或管坯沿径向胀形成所需形状 & \includegraphics[page=32, viewport={294.5 571.8 457.0 652.3}, clip]{12-07572-001FigRevised.pdf} & 用于制造各种形状中部较大的容器、管接头等,如球形管接头、军用水壶等 \\
\hline    
\end{longtable}

1. 冲裁

冲裁是落料和冲孔工序的统称。

落料、冲孔所用的冲模结构以及板料的变形过程均相同,但作用不同。落料是为了制取工件的外形,故冲下的部分为工件,带周边的为废料;冲孔则相反,是要制取工件的内孔,故冲下的部分为废料,带孔的部分为工件(图 8-27)。

冲裁时板料的变形过程如图 8-28 所示。当冲头接触板料向下运动时,板料首先产生弹性变形,继而进入塑性变形阶段;随凸模继续向下运动,变形程度增大;金属的加工硬化现象及位于凸模与凹模刃口处金属产生的应力集中使板料出现微裂纹,裂纹逐渐扩展,相互汇合,从而使板料与工件相分离。

\begin{figure}[H]
    \centering
    % Placeholder for Figure 8-27
    \includegraphics[page=24, viewport={4.5 387.8 137.0 484.3}, clip]{12-07572-001FigRevised.pdf}
    \caption{冲孔与落料示意图}
    \label{fig:8-27}
\end{figure}

板料分离后所形成的断口区域包括塌角、光亮带、剪裂带和毛刺等四部分。其中,光亮带尺寸准确,表面质量好,其余部分则使断口表面质量下降。

\begin{figure}[H]
    \centering
    % Placeholder for Figure 8-28
    \includegraphics[page=24, viewport={145.5 386.8 537.5 485.8}, clip]{12-07572-001FigRevised.pdf}
    \caption{金属板料的冲裁过程及断面特征}
    \label{fig:8-28}
\end{figure}

冲裁件断面质量的优劣,与冲模间隙、刃口锋利程度和材料排样方式密切相关。为了顺利完成冲裁过程,保证冲裁件的断面质量,要求凸模、凹模具有锋利的刃口以及合理的模具间隙 $z$。间隙过大或过小,均会影响冲裁件断面质量,甚至损坏冲模。模具间隙主要取决于板厚和冲裁件的精度要求,一般取值为板厚 $t$ 的 $5\% \sim 10\%$。

冲裁件的排样方式(即在板料或条料上的布置方法)对材料的利用率、生产成本和产品质量均有较大影响。采用有接边排样(图 8-29a),冲裁件质量较高,模具寿命也较长,但材料利用率较低;采用少无接边排样(图 8-29b)可减少废料,降低成本,但冲裁件尺寸精度不高,故应依据实际情况合理选取。生产中通常采用有接边排样。

由于冲裁时板料产生的裂纹并非在与板面垂直的方向,而是呈一定角度,故设计同样尺寸的落料模具和冲孔模具时,其刃口尺寸是不同的。工件上孔的尺寸取决于凸模尺寸,外形尺寸取决于凹模尺寸。所以,冲孔时凸模刃口尺寸应等于孔的尺寸,凹模刃口尺寸为凸模尺寸加上模具间隙值 $z$;落料时,凹模刃口尺寸应等于工件的外形尺寸,凸模刃口尺寸为凹模尺寸减去模具间隙值 $z$。

普通冲裁件只能满足一般产品的精度要求,远不能满足钟表、照相机、电子仪器等精密器械的要求。对于高精度冲裁件,需采用精密冲裁工艺。精冲技术的基本要素是精冲设备、精冲工艺、精冲模具、精冲材料和精冲润滑等。

精密冲裁方法很多,如带圆角模精冲法、负间隙精冲法和强力压边精冲法等。图 8-30 为目前应用较为普遍的强力压边精冲法示意图。它依靠 V 形压边环、极小间隙($z$ 一般为 $1\%t$)的模具、刃口略带小圆角的凹模和反压力顶杆等,得到精密冲裁件。

\begin{figure}[H]
    \centering
    % Placeholder for Figure 8-29
    \includegraphics[page=24, viewport={113.0 215.8 271.0 359.3}, clip]{12-07572-001FigRevised.pdf}
    \caption{材料排样方式}
    \label{fig:8-29}
\end{figure}

\begin{figure}[H]
    \centering
    % Placeholder for Figure 8-30
    \includegraphics[page=24, viewport={288.0 215.8 431.5 348.3}, clip]{12-07572-001FigRevised.pdf}
    \caption{强力压边精冲法示意图}
    \label{fig:8-30}
\end{figure}

2. 弯曲

弯曲是将板料、型材或管材弯成一定角度或圆弧的工序(图 8-31)。

板料弯曲时,内侧金属受切向压应力,产生压缩变形;外层金属受切向拉应力,产生伸长变形,表现为如图 8-32 所示板料上产生的矩形网格的变化。当拉应力超过材料的抗拉强度时,板料就会破裂。坯料厚度 $t$ 越大,弯曲半径 $r$ 越小,材料所受的内应力就越大,越容易弯裂。因此,必须控制最小弯曲半径 $r_{\min}$,通常取 $r_{\min} \geqslant (0.25 \sim 1)t$。材料塑性好时取下限。

\begin{figure}[H]
    \centering
    % Placeholder for Figure 8-31
    \includegraphics[page=24, viewport={76.5 73.3 313.0 181.3}, clip]{12-07572-001FigRevised.pdf}
    \caption{弯曲件示意图}
    \label{fig:8-31}
\end{figure}

\begin{figure}[H]
    \centering
    % Placeholder for Figure 8-32
    \includegraphics[page=24, viewport={326.0 72.8 468.0 185.3}, clip]{12-07572-001FigRevised.pdf}
    \caption{弯曲示意图}
    \label{fig:8-32}
\end{figure}

弯曲时还应尽可能使弯曲曲线与坯料纤维方向垂直(图 8-33),即使材料所受的拉应力与纤维方向一致,否则容易产生破裂。在双向弯曲时,应使弯曲曲线与纤维方向成 $45^\circ$ 角。在材料排样时应加以注意。

\begin{figure}[H]
    \centering
    % Placeholder for Figure 8-33
    \includegraphics[page=25, viewport={7.5 593.3 286.5 743.8}, clip]{12-07572-001FigRevised.pdf}
    \caption{弯曲方向对弯曲件质量的影响}
    \label{fig:8-33}
\end{figure}

弯曲时还应使弯曲件的毛刺区位于内侧。若位于外侧,圆角部位受拉应力而产生应力集中,容易引起该部位破裂。

在弯曲过程中,当载荷卸除后,工件的弯曲角度由于弹性变形的影响会略有增大,此现象称为回弹现象。增大的角度称为回弹角,一般为 $0^\circ \sim 10^\circ$。设计模具时应该考虑它的影响。例如,使模具角度比产品件角度小一个回弹角,以便在弯曲后得到准确的弯曲角度。另外,设计弯曲件时应加强其变形部位的刚性,如在弯曲部位设置加强筋等(图 8-34)。

\begin{figure}[H]
    \centering
    % Placeholder for Figure 8-34
    \includegraphics[page=25, viewport={297.0 593.8 538.0 701.8}, clip]{12-07572-001FigRevised.pdf}
    \caption{在弯曲部位设置加强筋}
    \label{fig:8-34}
\end{figure}

3. 拉深

拉深是将板料变形为中空形状零件的工序,可以生产筒形、锥形、球形、方盒形以及其他非规则形状的零件(图 8-35)。

\begin{figure}[H]
    \centering
    % Placeholder for Figure 8-35
    \includegraphics[page=25, viewport={67.5 371.8 297.5 484.8}, clip]{12-07572-001FigRevised.pdf}
    \caption{拉深件示意图}
    \label{fig:8-35}
\end{figure}

拉深过程如图 8-36 所示。在凸模作用下,板料被拉入凸模和凹模的间隙中,形成中空零件。板料在凸缘和凸模圆角部位变形最大。凸缘部分在圆周切线方向受压应力,压应力过大时会起皱(图 8-37),坯料厚度越小,拉深深度 $h$ 越大,越容易起皱。为防止起皱,可采用有压板(压力圈)拉深(图 8-38)。板料在凸模圆角部位承受拉应力,变薄严重,容易在此处拉裂。为防止拉裂,拉深模具的凸、凹模必须具有一定的圆角,圆角半径 $r_{\text{凸}} \leqslant r_{\text{凹}} = (5 \sim 10)t$;且模具间隙 $z$ 应稍大于板厚 $t$,一般 $z = (1.1 \sim 1.2)t$。

当筒形件直径 $d$ 与坯料直径 $D$ 相差较大时,不能一次拉深至产品尺寸,而应进行多次拉深,并在中间穿插进行退火处理,以消除前几次拉深变形所产生的加工硬化现象。

\begin{figure}[H]
    \centering
    % Placeholder for Figure 8-36
    \includegraphics[page=25, viewport={315.5 373.3 478.5 557.3}, clip]{12-07572-001FigRevised.pdf}
    \caption{拉深过程}
    \label{fig:8-36}
\end{figure}

\begin{figure}[H]
    \centering
    % Placeholder for Figure 8-37
    \includegraphics[page=25, viewport={82.0 250.8 231.5 342.3}, clip]{12-07572-001FigRevised.pdf}
    \caption{拉深废品}
    \label{fig:8-37}
\end{figure}

\begin{figure}[H]
    \centering
    % Placeholder for Figure 8-38
    \includegraphics[page=25, viewport={255.5 255.3 472.0 326.3}, clip]{12-07572-001FigRevised.pdf}
    \caption{用压边圈拉深}
    \label{fig:8-38}
\end{figure}

\subsection{冲压模具}

冲压模具(冲模)按冲床的每一次冲程所完成工序的多少可划分为简单模、连续模和复合模三大类。

简单模在冲床的一次冲程内只能完成一道工序,又称为单工序模。简单模的工作部分由凸模和凹模所组成,其采用导料板和限位销来控制板料的送进方向和送进量;依靠导柱与导套的精密配合来保证凸模准确进入凹模,进行冲裁工作,如图 8-39 所示。简单模的结构简单,成本低,但生产率较低,主要用于简单冲裁件的批量生产。

连续模在冲床的一次冲程内,在模具的不同位置上可以同时完成两道以上的工序,又称为级进模。图 8-40 所示为落料、冲孔连续模。左侧为落料模,右侧为冲孔模,条料送进时,先冲孔,后落料,而且在同一冲程内完成。连续模生产率高,易于实现自动化,但结构复杂,成本也相应增高,适用于大批量生产精度要求不高的中、小型零件。

复合模在冲床的一次冲程内,在模具的同一位置上可以同时完成两道以上的工序。图 8-41 所示为生产中经常采用的落料、拉深复合模。它的特点是有一个凸凹模,其外圆为落料凸模,内孔为拉深凹模。当凸凹模下降时,首先与落料凹模配合进行落料,然后与拉深凸模配合进行拉深。这样在一个冲程内,同一位置上便可完成落料和拉深两道工序。复合模生产率高,零件加工精度高,但模具制造复杂,成本高,适用于大批量生产。

\begin{figure}[H]
    \centering
    \includegraphics[page=25, viewport={136.0 59.3 358.0 214.8}, clip]{12-07572-001FigRevised.pdf}
    \caption{简单模}
    \label{fig:8-39}
\end{figure}

\begin{figure}[H]
    \centering
    \includegraphics[page=26, viewport={126.0 601.3 417.5 746.8}, clip]{12-07572-001FigRevised.pdf}
    \caption{连续模}
    \label{fig:8-40}
\end{figure}

\begin{figure}[H]
    \centering
    \includegraphics[page=26, viewport={99.5 347.8 445.0 569.8}, clip]{12-07572-001FigRevised.pdf}
    \caption{复合模}
    \label{fig:8-41}
\end{figure}

\section{其他压力加工方法简介}
随着经济发展和需求的变化,压力加工领域的新设备、新工艺、新技术层出不穷。已从过去的单一材料成形扩展为复合材料成形;由金属材料成形扩展到陶瓷、塑料等非金属材料成形;从常态成形扩展到超塑成形;从普通压力机上成形扩展到高速成形和蠕变成形等。本节就相关内容只做简要介绍。

\subsection{超塑成形}
超塑成形自20世纪70年代以来得到迅速发展和应用。超塑性实际上是材料在特定条件下所表现出的异常高伸长率的能力,即在低变形速率($\varepsilon=10^{-2}\sim 10^{-4}\text{s}^{-1}$)、一定的变形温度(约为其熔点的一半)以及一定的晶粒度(一般为$0.2\sim 5\text{ \textmu m}$)下,其伸长率($A$)超过$100\%$的特性。如钢的伸长率超过$500\%$、纯钛的伸长率超过$300\%$、锌铝合金的伸长率超过$1000\%$。

超塑性金属的变形特性近似于高温玻璃或高温聚合物,在比常规变形低得多的载荷下,可以成形出高质量、高精度的薄壁、薄腹板、高筋件和其他复杂件。超塑成形特别适宜于变形力大,塑性低,在常规成形条件下较难成形的金属材料(如钛合金、镁合金、镍合金、合金钢等)。近年来,还发现金属间化合物、复合材料和陶瓷经细晶处理后也有超塑性,从而为这些高性能、难加工材料的成形开辟了新途径。

常用超塑成形工艺有以下几种。

    1. 超塑性模锻
    
    超塑性模锻与普通模锻的主要区别在于工艺参数不同(表8-9)和具有一套能够使模具和变形材料在成形过程中保持恒温的加热装置(通常采用感应加热和电阻加热)。如采用普通热模锻成形高温合金及钛合金时,机械加工损耗量高达$80\%$,但采用超塑性模锻,其损耗可降低一半以上。
    
    \begin{longtable}[]{|l|c|c|}
        
        \caption{普通模锻和超塑性模锻工艺参数}\\
        
        
            \hline
            模锻工艺参数 & 普通模锻 & 超塑性模锻 \\
            \hline
            毛坯加热温度/$^\circ\text{C}$ & 940 & 940 \\
            模具加热温度/$^\circ\text{C}$ & 480 & 940 \\
            变形速度/($\text{mm/s}$) & $12.7\sim 42.3$ & 0.025 \\
            平均单位压力/($\text{N/mm}^2$) & $50.0\sim 58.3$ & 11.7 \\
            模锻工步次数 & 4 & 1 \\
            \hline
        
    \end{longtable}

    2. 超塑性无模拉拔
    
    超塑性无模拉拔是利用感应线圈局部加热,使材料处于超塑性变形温度时而进行拉拔的工艺方法(图8-42),连续加热时,可生产等断面制品;断续加热并控制拉拔速度与感应线圈移动速度,可生产不等断面制品。超塑性无模拉拔的最大断面收缩率可达$83\%$,加工精度可达$\pm 0.013\text{ mm}$,主要用于管材、棒材的二次成形及断面形状简单的制品。


\begin{figure}[H]
    \centering
    \includegraphics[page=26, viewport={76.5 208.3 259.5 295.8}, clip]{12-07572-001FigRevised.pdf}
    \caption{超塑性无模拉拔示意图}
    \label{fig:8-42}
\end{figure}

\begin{figure}[H]
    \centering
    \includegraphics[page=26, viewport={277.0 208.8 423.5 316.8}, clip]{12-07572-001FigRevised.pdf}
    \caption{超塑性气压胀形工装示意图}
    \label{fig:8-43}
\end{figure}

2. 超塑性气压胀形
    
    超塑性气压胀形是利用凹模或凸模的形状,先把板料和模具加热到预定温度,然后向模具内通入压缩空气,使板料紧贴在凹模或凸模上,从而获得所需制件(图8-43)的工艺方法。该工艺方法主要用于钛合金、铝合金和双相不锈钢薄板(板厚一般为$0.4\sim 4\text{ mm}$)的成形。

    3. 超塑性胀形与扩散连接复合工艺
    
    超塑性胀形与扩散连接复合(SPF/DB)工艺是先将板坯胀形至所需形状,而后通过局部扩散连接使其结合在一起的工艺方法。对于多层板结构,需先进行扩散连接,而后气胀成夹层结构(图8-44)。该工艺可生产外形复杂的结构件,并能简化装配工序,主要用于成形钛合金与铝合金的夹层结构件,已成功地用于飞机和卫星一类航空航天器结构件的制造上。


\begin{figure}[H]
    \centering
    \includegraphics[page=26, viewport={132.5 35.8 409.0 181.3}, clip]{12-07572-001FigRevised.pdf}
    \caption{夹层结构 SPF/DB 过程}
    \label{fig:8-44}
\end{figure}

除金属材料的超塑性外,目前国际上对陶瓷材料及金属基复合材料的超塑性研究也有很大进展。例如,在金属基超塑性材料中加入SiC纤维后,也可达到超塑性气压胀形的要求。

\subsection{粉末锻造}
粉末锻造是将粉末烧结的预成形坯经加热后,在闭式锻模中锻造成零件的工艺方法,具体工艺过程如图8-45所示,即先制备合格的粉末,再成形、烧结、精密热锻,然后进行热处理。粉末锻造制品具有尺寸精度高、组织结构均匀、无成分偏析等特点,可以锻造难变形的高温铸造合金,在许多领域中得到应用。粉末锻造生产的高速钢车刀、钻头、工具等的寿命比普通方法生产要提高$1\sim 2$倍。钛合金的粉末锻造产品在国防领域,尤其在航空航天和高性能装备中已得到重要应用,已研制出Ti-6Al-4V合金喷气发动机叶片、燃气轮机导向叶片等。另外,粉末锻造还在汽车制造业中得到广泛应用。例如,汽车发动机中的齿轮和连杆,其动平衡性能要求高,材质要求均布,最适宜采用粉末锻造生产。

\begin{figure}[H]
    \centering
    \includegraphics[page=27, viewport={60.0 603.3 327.0 748.3}, clip]{12-07572-001FigRevised.pdf}
    \caption{粉末锻造示意图}
    \label{fig:8-45}
\end{figure}

常用的粉末锻造方法有粉末热锻、烧结锻造、锻造烧结和粉末冷锻。近年来,新的工艺方法不断涌现,如松装锻造法、球团锻造法、粉末热等静压法、粉末准等静压法,喷雾锻造法、粉末等温锻、粉末热挤压、粉末摆动碾压、粉末连续挤压、粉末轧制和粉末爆炸成形等。金属粉末轧制原理如图8-46所示,粉末经过轧机轧辊之间的压力作用而制成一定厚度的带材。连续轧制生产带材的工艺过程如图8-47所示,在连续生产时,轧制带材在烧结以后还可以进一步轧制或进行多次退火轧制。粉末喷雾锻造工艺过程如图8-48所示,首先采用高速氮气喷射金属液流,使雾化的粉末直接沉积到预成形模具中,沉积的预成形坯密度很高,相对密度可达$99\%$;然后将预成形坯从雾化室中取出,加热至锻造温度进行锻造,再送至切边压力机切边即获得成品锻件。该方法适合于生产大型锻件,还可进行喷射轧制、喷射挤压等。

\begin{figure}[H]
    \centering
    \includegraphics[page=27, viewport={346.5 603.8 486.0 729.8}, clip]{12-07572-001FigRevised.pdf}
    \caption{金属粉末轧制原理图}
    \label{fig:8-46}
\end{figure}

\begin{figure}[H]
    \centering
    \includegraphics[page=27, viewport={6.5 273.8 285.5 419.3}, clip]{12-07572-001FigRevised.pdf}
    \caption{连续轧制生产带材示意图}
    \label{fig:8-47}
\end{figure}

\subsection{回转成形}
回转成形过程是局部变形的连续累积过程,因生产率高、设备吨位小,应用范围逐步扩大。表8-10为几种较为成熟的回转成形新工艺的工作原理、特点及应用。

\begin{longtable}{|l|m{5cm}|m{4cm}|m{4cm}|}
    \caption{几种较为成熟的回转成形新工艺的工作原理、特点及应用}\\
        \hline
        名称 & 简图 & 工作原理 & 特点及应用 \\
        \hline
        辊锻 & \includegraphics[page=32, viewport={68.0 441.3 185.0 532.3}, clip]{12-07572-001FigRevised.pdf} & 坯料纵向通过辊锻机上的一对装有圆弧形模块且相对旋转的轧辊时,受压变形而形成锻件 & 设备吨位小,尺寸稳定,精度、效率高,材料消耗少,劳动条件好。适于长杆类件的大批量生产(如汽轮机叶片等),也可为其他模锻件制坯 \\
        \hline
        楔横轧 &\includegraphics[page=32, viewport={212.5 441.8 337.5 535.8}, clip]{12-07572-001FigRevised.pdf} & 圆柱形坯料在两个同向旋转、带有楔形模具的轧辊作用下旋转并受压变形,当轧辊转一周时,形成一个锻件 & 生产率高,模具寿命高,节省材料,易于实现机械化操作,劳动条件好。适于轴类零件的大批量生产(如汽车、摩托车上的轴类件),也可用于其他轴类模锻件的制坯 \\
        \hline
        斜轧 & \includegraphics[page=32, viewport={364.5 439.3 475.5 539.3}, clip]{12-07572-001FigRevised.pdf} & 圆柱形坯料在两个同向旋转且轴心线呈一定角度的轧辊作用下,在旋转的同时作直线运动,在不同孔形的轧辊作用下,局部连续成形为所需毛坯或零件 & 生产率高,模具寿命高,节省材料,易于实现机械化操作,劳动条件好。广泛用于生产钢球、轴承滚子、麻花钻头和空心轴等零件或毛坯 \\
        \hline
        摆动辗压 & \includegraphics[page=32, viewport={77.5 338.3 174.0 423.3}, clip]{12-07572-001FigRevised.pdf} & 坯料在有摆角的上模旋转挤压下,连续局部变形,高度减小,直径增大,形成盘状或局部盘状锻件 & 设备吨位小,产品质量高,劳动条件好。适于制造薄盘形锻件(如铣刀片、汽车半轴等) \\
        \hline
        旋压 & \includegraphics[page=32, viewport={196.0 335.8 325.5 418.3}, clip]{12-07572-001FigRevised.pdf} & 坯料随芯模旋转(或旋压工具绕坯料与芯模旋转),旋压工具相对芯模进给,坯料受压产生连续逐点变形,完成工件加工 & 变形力小,模具费用低,可批量生产筒形、卷边等回转体工件,以及形状复杂或高强度难变形材料(如薄壁食品罐等) \\
        \hline
\end{longtable}

\subsection{静液挤压}
静液挤压(图8-49)是一种采用液体静压力对工件实施挤压的塑性成形方法。挤压时挤压冲头不直接接触坯料,而是向液体施加压力($2000\sim 3000\text{ MPa}$),经液体将压力传递给坯料,在液体压力的作用下,金属坯料通过成形模挤出制件。由于静液挤压对坯料施加的是等静压力,有利于材料的塑性变形,所以可用于低塑性材料的挤压成形,如铍、铬、钼、钨等金属及其合金的成形,并可提高一次挤压的变形量。此外,静液挤压力也较其他挤压方法小$10\%\sim 50\%$,对于常用材料可以采用大变形量一次挤压成线材和型材。目前,静液挤压已用于挤制螺旋齿轮及麻花钻等形状复杂的零件。

\begin{figure}[H]
    \centering
    \includegraphics[page=27, viewport={289.0 273.3 496.5 573.8}, clip]{12-07572-001FigRevised.pdf}
    \caption{粉末喷雾锻造工艺过程}
    \label{fig:8-48}
\end{figure}

\begin{figure}[H]
    \centering
    \includegraphics[page=27, viewport={79.5 50.8 278.5 249.3}, clip]{12-07572-001FigRevised.pdf}
    \caption{静液挤压示意图}
    \label{fig:8-49}
\end{figure}

\subsection{压力加工技术的发展趋势}
压力加工技术发展的总趋势是交叉-综合化、数字-智能化、洁净-高效化和柔性-集成化,具体可以概括为以下几个方面。


(1)采用柔性成形工艺及发展锻压生产的柔性加工系统
    
柔性加工是指适应产品多变而工装变化很少的方法。如高内压成形(internal high pressure forming, IHPF)、黏性介质压力成形(Viscous pressure forming, VPF)、无模液压成形、数控增量成形及多点成形等。柔性加工系统(fleibel manufacturing system, FMS)包括自动化加工系统、物料输送系统及计算机信息系统三大部分,可以自动完成从零件设计、工艺制订、物料输送、零件加工直至工件输出的全过程。锻压生产的柔性加工系统,特别是板材成形的柔性加工系统技术已相当成熟。美国、日本、德国、意大利等国已生产出成套的板材柔性加工设备,目前已向计算机集成制造系统(computer integrated manufacturing system, CIMS)方向发展。

(2)增加设备的柔性和发展低能耗成形技术
    
以汽车覆盖件拉深为例,传统方法是采用刚性压边圈,很难成形深度较大的覆盖件,如果压边力在成形过程中可以调节,则可生产各种复杂拉深件。目前,国外已开发出多方位联动(或分别动作)可调压力机,为新产品开发带来诸多便利。

(3)实现精密塑性成形并与其他工艺交叉运用
    
精密塑性成形可分为净成形及近净成形,具有节约材料和节省后续机加工量的优点,与其他工艺交叉运用更具有广泛的应用前景。

(4)构建数字化塑性成形技术体系,并注重数字化技术与传统加工方式相结合
    
目前,计算机辅助技术在塑性成形领域的应用不断深入,使制件质量提高,制造周期下降。构建数字化塑性成形技术体系,就是要在塑性成形全过程(塑性成形产品设计、分析和制造过程)中融合数字化技术,并以系统工程为理论基础的技术体系,实现优质、高效、低耗、洁净生产。


\section{锻件与冲压件结构设计}
以锻件为毛坯的零件,其结构应适应于金属可锻性和锻造工艺的要求。由于锻造是在固态下成形的,锻件的复杂程度远不如铸件。设计锻件结构时,应考虑锻造成形的可行性和经济性。

\subsection{自由锻件的设计原则}
自由锻件的形状一般都较为简单。由于自由锻是在平砧上用简单工具锻造而成的,其形状、尺寸主要靠工人的操作技术水平来保证,不可能锻造出形状复杂的锻件。因而,自由锻件应设计得尽量简单。自由锻件的结构工艺性见表8-11。

\begin{longtable}{|m{2cm}|m{4.5cm}|m{4.5cm}|}
    
    \caption{自由锻件的结构工艺性}\\
    
     \hline
       
        设计原则 & 不合理结构 & 合理结构 \\
        \hline
        尽量避免锥面或斜面 & \includegraphics[page=32, viewport={344.5 337.3 466.5 412.8}, clip]{12-07572-001FigRevised.pdf} & \includegraphics[page=32, viewport={68.5 202.3 194.0 271.8}, clip]{12-07572-001FigRevised.pdf} \\
        \hline
        避免圆柱面与圆柱面、圆柱面与棱柱面相交 & \includegraphics[page=32, viewport={215.0 202.3 334.0 304.8}, clip]{12-07572-001FigRevised.pdf} & \includegraphics[page=32, viewport={358.5 202.3 475.5 305.8}, clip]{12-07572-001FigRevised.pdf} \\
        \hline
        避免椭圆形、工字形及其他非规则斜面或外形 & \includegraphics[page=32, viewport={65.5 70.3 184.5 123.3}, clip]{12-07572-001FigRevised.pdf} & \includegraphics[page=32, viewport={204.0 69.8 316.0 123.8}, clip]{12-07572-001FigRevised.pdf} \\
        \hline
        避免加强筋或凸台等结构 & \includegraphics[page=32, viewport={336.0 69.8 478.0 170.3}, clip]{12-07572-001FigRevised.pdf} & \includegraphics[page=33, viewport={64.5 630.8 205.0 737.3}, clip]{12-07572-001FigRevised.pdf} \\
        \hline
        对横截面尺寸相差较大和形状复杂的零件,可采用分体锻造,再采用焊接或机械连接组合为整体 & \includegraphics[page=33, viewport={226.5 629.3 345.5 747.3}, clip]{12-07572-001FigRevised.pdf} & \includegraphics[page=33, viewport={365.5 629.8 484.0 747.8}, clip]{12-07572-001FigRevised.pdf} \\
        \hline
    
\end{longtable}

\subsection{模锻件的设计原则}
模锻件是依靠锻模模膛控制成形的,其形状可以比自由锻件的形状复杂得多,允许有复杂的曲面,细小的花纹和文字,可以带筋和小凸台,但必须保证模锻件能从模膛中顺利取出。零件上较小的孔,很小的沟槽,尤其是与锤击方向平行的内、外表面上局部凹入部分是无法锻制出来的。

另外,模锻时金属的流动阻力较大,设计时必须考虑使金属易于充满模膛。模锻件的结构工艺性见表8-12。

\begin{longtable}[h]{|m{0.25\textwidth}|m{0.3\textwidth}|m{0.3\textwidth}|}

\caption{模锻件的结构工艺性}\\

\hline
设计原则 & 不合理结构 & 合理结构 \\
\hline
零件形状力求简单,避免带有长而复杂的分枝和多向弯曲等复杂形状 & \includegraphics[page=33, viewport={10.0 528.8 108.5 588.3}, clip]{12-07572-001FigRevised.pdf} & \includegraphics[page=33, viewport={27.0 427.3 148.0 469.8}, clip]{12-07572-001FigRevised.pdf} \\
\hline
零件形状尽可能对称,以使锻模和设备受力均匀,延长其使用寿命 & \includegraphics[page=33, viewport={132.0 528.3 245.0 577.8}, clip]{12-07572-001FigRevised.pdf} & \includegraphics[page=33, viewport={163.0 431.3 286.0 475.3}, clip]{12-07572-001FigRevised.pdf} \\
\hline
零件上与分模面垂直的表面尽可能避免凹槽和孔,以便于取出锻件 & \includegraphics[page=33, viewport={266.5 530.8 354.5 600.8}, clip]{12-07572-001FigRevised.pdf} & \includegraphics[page=33, viewport={305.0 429.8 398.5 502.3}, clip]{12-07572-001FigRevised.pdf} \\
\hline
避免薄壁、高筋、深孔和直径过大的凸缘,以减少金属充模阻力 & \includegraphics[page=33, viewport={368.5 532.3 453.0 590.8}, clip]{12-07572-001FigRevised.pdf} & \includegraphics[page=33, viewport={422.0 429.8 502.5 495.3}, clip]{12-07572-001FigRevised.pdf} \\
\hline
对于复杂零件可以采用锻-焊组合结构,以简化模锻工艺和降低废品率 & \includegraphics[page=33, viewport={472.5 529.3 539.5 582.3}, clip]{12-07572-001FigRevised.pdf} & \includegraphics[page=33, viewport={98.0 282.8 241.5 338.3}, clip]{12-07572-001FigRevised.pdf} \\
\hline
\end{longtable}

\subsection{冲压件的设计原则}
1. 冲压件的材料选择

冲压件材料选择应满足冲压基本工序对材料性能和表面质量的要求。例如,成形零件所用的材料应具有良好的塑性,而对于平板冲裁件的塑性则无过高要求。

金属材料的冲压成形性能与其化学成分、组织性能、材料表面质量及板料厚度公差等有关,具体可查阅有关技术资料。

2. 冲压件的形状

一般来讲,冲压件结构形状应力求简单、对称,以使模具受力均衡,材料变形均匀,便于保证冲压件质量。

对于普通冲裁件,为保证模具强度和冲裁件质量,应避免细长槽与悬壁结构;对于弯曲件可采用压制加强筋来提高刚性(图8-50);对于带孔弯曲件,为避免孔的变形,孔的位置应如图8-51所示,图中$L$应大于$(1.5 \sim 2)s$;对于拉深件,为防止拉裂与起皱,应尽量设计成轴对称零件,并尽量减小变形高度。

\begin{figure}[H]
\centering
\includegraphics[page=27, viewport={298.5 50.3 461.0 135.8}, clip]{12-07572-001FigRevised.pdf}% Placeholder for Fig 8-50
\caption{压制加强筋提高刚性}
\end{figure}
\begin{figure}[H]
\centering
\includegraphics[page=28, viewport={51.0 601.8 176.5 684.3}, clip]{12-07572-001FigRevised.pdf} % Placeholder for Fig 8-51
\caption{带孔弯曲件}
\end{figure}

对于形状过于复杂的冲压件,可采用冲-焊组合结构(图8-52);也可采用冲口结构来代替组合零件,以降低生产成本。图8-53所示的组合件,原工艺为三构件铆接或焊接而成,现采用冲口和弯曲工艺制成整体零件,可节省材料,简化工艺。

\begin{figure}[H]
\centering
\includegraphics[page=28, viewport={197.0 603.8 348.0 692.3}, clip]{12-07572-001FigRevised.pdf} % Placeholder for Fig 8-52
\caption{冲压-焊接组合零件}
\end{figure}
\begin{figure}[H]
\centering
\includegraphics[page=28, viewport={365.0 603.8 498.5 744.3}, clip]{12-07572-001FigRevised.pdf} % Placeholder for Fig 8-53
\caption{冲口工艺举例}
\end{figure}

3. 冲压件的尺寸

冲压件的转角一般应设计成圆角,圆角半径应大于冲压基本工序所允许的最小圆角半径;设计冲裁件要保证落料和冲孔所允许的最小尺寸、孔与孔和孔与零件边缘的距离;设计弯曲件时要注意弯曲件的直边长、孔与弯曲半径中心的最小距离;对于拉深件应注意拉深件的深度和拉深件凸缘尺寸等。具体要求可查阅相关手册。另外,在强度、刚度允许的条件下,应尽可能采用较薄的材料制作零件,以减少材料的消耗。

\section*{复习思考题}

 
8-1 常用的金属压力加工方法有哪些?各有何特点?

8-2 为什么弹簧钢丝一般采用冷拉、冷卷成形?在拉制过程中,为什么被拉过模孔而截面已缩小的钢丝,其截面不再缩小,也不会被拉断?

8-3 钢材在热变形过程中纤维组织是如何形成的?它的存在有何利弊?设计零件时如何合理利用纤维组织?

8-4 何谓冷变形和热变形?纯铅丝和纯铁丝反复折弯会发生什么现象?为什么?

8-5 何谓金属的可锻性?影响可锻性的因素有哪些?

8-6 为什么锻件的力学性能常优于铸件?

8-7 为什么重要的轴类锻件在锻造过程中均安排镦粗工序?

8-8 在如图8-54所示的两种砧铁上拔长时,其效果有何不同?

\begin{figure}[H]
\centering
\includegraphics[page=28, viewport={25.5 260.8 184.0 375.3}, clip]{12-07572-001FigRevised.pdf}
\caption{题8-8图}
\end{figure}

8-9 试述自由锻、胎模锻和模锻的特点及适用范围。

8-10 终锻模膛中设计飞边槽的作用是什么?

8-11 什么是分模面?分模面选择的一般原则是什么?

8-12 图8-55所示的零件采用锤上模锻制造,请选择最合适的分模面位置。

\begin{figure}[H]
\centering
\includegraphics[page=28, viewport={281.5 261.8 434.0 413.3}, clip]{12-07572-001FigRevised.pdf}
\caption{题8-12图}
\end{figure}

8-13 对于图8-56所示的零件拟采用自由锻制坯,试定性绘出锻件图,选择自由锻工序,并绘出变形简图。

\begin{figure}[H]
\centering
\includegraphics[page=28, viewport={117.0 118.8 422.5 231.3}, clip]{12-07572-001FigRevised.pdf}
\caption{题8-13图}
\end{figure}

8-14 图8-57所示模锻零件的设计有哪些不合理的地方?应如何改进?

\begin{figure}[H]
\centering
\includegraphics[page=29, viewport={149.5 642.3 396.0 746.8}, clip]{12-07572-001FigRevised.pdf}
\caption{题8-14图}
\end{figure}

8-15 图8-58所示三种零件若分别采用单件、小批和大批量生产时,应选用哪种锻造方法制造?试定性绘出锻件图。

\begin{figure}[H]
\centering
\includegraphics[page=29, viewport={52.0 504.8 482.5 614.8}, clip]{12-07572-001FigRevised.pdf}
\caption{题8-15图}
\end{figure}

8-16 下列零件应选用何种锻造方法制坯?

(1) 活口扳手(大批量);

(2) 大六角螺钉(成批);

(3) 铣床主轴(成批);

(4) 起重机吊钩(单件)。

8-17 简述板料冲压的特点和应用范围。

8-18 如何确定冲孔模和落料模的刃口尺寸?用$\phi 50 \text{ mm}$的冲孔模具来生产$\phi 50 \text{ mm}$的落料件能否保证冲压件精度?为什么?

8-19 板料拉深时,如何防止工件拉裂和边缘起皱?

8-20 冲模结构分为哪几类?各有何特点?

8-21 简述图8-59所示冲压件的生产过程。

8-22 超塑成形与普通塑性成形有何区别?常用的超塑成形方法有哪些?

8-23 简述摩擦压力机的工作原理。为什么在摩擦压力机上不宜进行多模膛锻造?

8-24 何谓精密模锻?它与普通模锻相比有何特点?

8-25 试分析通过压力加工实现材料强化的主要原因。

8-26 压力加工技术的发展趋势是什么?

\begin{figure}[H]
\centering
\includegraphics[page=29, viewport={65.0 327.3 201.0 465.8}, clip]{12-07572-001FigRevised.pdf}
\caption{题8-21图}
\end{figure}


\chapter{焊接}

焊接是利用局部加热或(和)加压等手段,使分离的金属通过原子的结合与扩散而形成永久性连接的工艺方法。根据实现金属原子间结合的方式不同,焊接方法可分为熔焊、压焊和钎焊3大类。

熔焊(又称为熔化焊)是利用热量使焊件接头局部熔化,再冷却结晶成一体的焊接方法。熔焊的加热温度较高,焊件容易变形。但接头表面的清洁程度要求不高,操作方便,适用于各种常用金属材料的焊接,应用较广。

压焊(又称为压力焊)是对焊件加热(或不加热)并施压,使其接头处紧密接触并产生塑性变形和再结晶,从而形成原子间结合的焊接方法。压焊只适用于塑性较好的金属材料的焊接。

钎焊是将低熔点的钎料熔化,填充到接头间隙,利用钎料与固态母材(焊件)的毛细作用和相互扩散实现连接的焊接方法。钎焊不仅适用于同种或异种金属的焊接,还广泛用于金属与玻璃、陶瓷等非金属材料的连接。

焊接工艺成形方便、适应性强、生产成本低(与铆接相比)。在制造大型、复杂结构的零件时,可与铸件、锻件和冲压件相结合,化大为小,化复杂为简单,再逐次装配焊接而成。采用相应的焊接方法,既能生产微型、大型和复杂的金属构件,也能生产气密性好的高温、高压设备和化工设备;既适应于单件、小批生产,也适应于大批量生产。另外,采用焊接结构能够按使用要求选用材料。在结构的不同部位,按强度、耐磨性、耐腐蚀性、耐高温等要求选用不同材料,如制造双金属构件或复合层容器,具有更好的经济性。因此,焊接在现代工业生产中具有十分重要的作用,在车辆、建筑、舰船、海洋钻探、航空航天、核电站及石油化工、电气、微电子等各现代工业领域均获得了广泛应用。

\section{焊接工程理论基础}

\subsection{焊接电弧与电弧焊机}

1. 焊接电弧

焊接电弧是在电极与工件之间的气体介质中产生的强烈、持久的气体放电现象,即在电极与工件之间局部气体介质中有大量电子流通过的导电现象。

焊接电弧如图9-1所示。引燃电弧后,弧柱中就充满了高温电离气体,并放出大量的热能和强烈的光。电弧的热量与焊接电流和电弧电压的乘积成正比。电流越大,电弧产生的总热量就越多。一般情况下,电弧热量在阳极区产生的较多,约占总热量的43\%;阴极区因放出大量的电子,消耗了一部分能量,所以产生的热量相对较少,约占36\%;其余21\%左右的热量是在弧柱区产生的。焊条电弧焊通常只有65\%~85\%的热量用于加热和熔化金属,其余的热量则散失在电弧周围和飞溅的金属滴中。

\begin{figure}[H]
\centering
\includegraphics[page=33, viewport={264.5 283.8 439.5 394.8}, clip]{12-07572-001FigRevised.pdf}
\caption{焊接电弧}
\end{figure}

电弧中阳极区和阴极区的温度因电极材料不同而有所不同。用钢焊条焊接钢材时,阳极区温度约为$2600 \text{ K}$,阴极区约为$2400 \text{ K}$,电弧中心区温度最高,可达$6000 \sim 8000 \text{ K}$。

由于电弧产生的热量在阳极区和阴极区有所不同,因此使用直流电源焊接其接线方法有正接和反接两种。

正接是将工件接到电源正极,焊条(或电极)接到负极;反接是将工件接到电源负极,焊条(或电极)接到正极,如图9-2所示。正接时工件的温度相对高些。

\begin{figure}[H]
\centering
\includegraphics[page=33, viewport={13.5 123.3 327.5 213.8}, clip]{12-07572-001FigRevised.pdf}
\caption{直流电源焊接的正接与反接}
\end{figure}

使用交流弧焊机焊接时,因为电极极性不断变化,两极加热温度都在$2500 \text{ K}$左右,因而不存在正接和反接问题。

2. 焊接电源

(1) 对焊接电源的要求

一般电气设备要求电源电压不随负载的变化而变化,但焊接电源却要求其电压随负载增大而迅速降低,即具有陡降的外特性,这样才能满足下列焊接要求。

1) 具有一定的空载电压(即焊接开始时的引弧电压),以满足引弧的需要,一般为$50 \sim 90 \text{ V}$。

2) 具有适当的短路电流,以保证焊接过程频繁短路时,电流不致无限增大而烧毁电源。短路电流一般不超过工作电流的$1.25 \sim 2$倍。

3) 电弧长度发生变化时,能保证电弧的稳定。电弧稳定燃烧时的电压降称为电弧电压,它与电弧长度(即焊条与工件之间的距离)有关。电弧长度愈大,电弧电压也愈大。电弧电压一般为$16 \sim 36 \text{ V}$。

此外,焊接电源还应能方便地调节焊接电流和焊接电压,以适应不同种类和厚度材料的焊接要求。

(2) 常用焊接设备的类型

手工电弧焊设备简称电弧焊机,其类型主要有交流弧焊机、直流弧焊机和交直流两用弧焊机。

交流弧焊机实质上是一个具有下降外特性的漏磁的降压变压器,可将工业电压($220 \text{ V}$或$380 \text{ V}$)降低到空载电压及工作电压。同时,能提供很大的焊接电流,并能在一定范围内进行调节。交流弧焊机结构简单,价廉,工作噪声小,使用和维修方便,应用较广泛。直流弧焊机包括直流弧焊发电机和弧焊整流器两大类。直流弧焊机焊接时电弧稳定,能适应各种焊条,接头质量较好,但结构复杂,价格较高,噪声大。交直流两用弧焊机常用作多用途弧焊机。

(3) 电弧焊机的选用

使用酸性焊条焊接一般低碳钢构件时,应优先考虑选用价格低廉、维修方便的交流弧焊机;使用碱性焊条焊接高压容器、高压管道等重要钢结构,或焊接合金钢、非铁金属、铸铁时,应选用直流弧焊机。

当购置能力有限而焊件材料的类型繁多时,可考虑选用通用性强的交直流两用弧焊机。在野外无电网供电时,应选用内燃机驱动的直流弧焊机。近年来,国内外竞相发展的逆变式弧焊整流器,具有效率高、重量轻、成本低、工艺性能好等优点,有逐渐取代其他焊接电源的趋势。

\subsection{焊接冶金过程与焊条}

1. 焊接冶金特点

在进行电弧焊时,被熔化的金属、熔渣、气体三者之间将发生复杂的冶金反应,如金属的氧化与还原、气体的溶解与析出、杂质的去除等。因此,焊接熔池可以看成一座微型冶金炉。但焊接冶金过程与一般的冶炼过程不同,主要有以下特点。

(1) 冶金温度高

在焊接碳素结构钢和普通低合金钢时,熔滴的平均温度约$2300 \text{ }^\circ\text{C}$,熔池也达$1600 \text{ }^\circ\text{C}$以上,远高于普通冶金温度,容易造成合金元素的烧损与蒸发。

(2) 冶金过程短

焊接时,由于焊接熔池体积很小(一般为$2 \sim 3 \text{ cm}^3$),冷却速度快(熔池周围是冷金属),液态停留时间很短(熔池从形成到凝固约$10 \text{ s}$),各种冶金反应无法达到平衡状态,反应进行得很不充分,焊缝中会出现化学成分不均匀的偏析现象。

(3) 冶金条件差

焊接熔池一般暴露在空气中,熔池周围的气体、铁锈、油污等在电弧的高温作用下,将分解成原子态的氢、氧、氮等,极易同金属元素产生化学反应。反应生成的氧化物、氮化物混入焊缝中,使焊缝的力学性能下降;液态金属氧化时生成$\text{FeO}$,其溶解于钢液中,冷凝时因溶解度减小而析出,杂质则滞留在焊缝里;$\text{FeO}$与钢中的$\text{C}$反应生成$\text{CO}$,易在焊缝中产生气孔;液态金属氮化时生成$\text{Fe}_4\text{N}$,其冷凝时呈针状的夹杂物分布在晶粒内,显著降低焊缝塑性和韧性;空气中水分分解成氢原子,使焊缝中产生气孔、裂纹等缺陷,会出现“氢脆”现象。

上述情况将严重影响焊接质量,因此必须采取有效措施来保护焊接区,防止周围有害气体侵入金属熔池;同时要控制焊缝金属的化学成分,向金属熔池中补充易烧损的合金元素;此外,还要进行脱氧和脱硫、脱磷,以减少焊接缺陷,获得优质焊接接头。

2. 焊条

(1) 焊条的组成及作用

焊条由金属焊芯和药皮两部分所组成。

焊芯的主要作用是作为电极和填充金属,用于引燃电弧和填充焊缝,其化学成分直接影响焊缝质量。焊芯通常用含碳、硫、磷较低的专用钢丝制成(表9-1)。焊芯的直径(即焊条直径)最小为$1.6 \text{ mm}$,最大为$8 \text{ mm}$,其中以$3.2 \sim 5 \text{ mm}$的焊条应用最广。

\begin{table}[H]
\centering
\caption{常用焊接用钢丝的牌号和化学成分}
\label{tab:9-1}
\small
\begin{tabular}{|l|c|c|c|c|c|c|c|l|}
\hline
\multirow{2}{*}{牌号} & \multicolumn{7}{c|}{化学成分 $w/\%$} & \multirow{2}{*}{用途} \\
\cline{2-8}
 & C & Mn & Si & Cr & Ni & S & P & \\
\hline
H08 & $\leqslant 0.10$ & $0.30 \sim 0.55$ & $\leqslant 0.03$ & $\leqslant 0.20$ & $\leqslant 0.30$ & $\leqslant 0.04$ & $\leqslant 0.04$ & 一般焊接结构 \\
\hline
H08A & $\leqslant 0.10$ & $0.30 \sim 0.55$ & $\leqslant 0.03$ & $\leqslant 0.20$ & $\leqslant 0.30$ & $\leqslant 0.03$ & $\leqslant 0.03$ & \multirow{2}{*}{重要焊接结构} \\
\cline{1-8}
H08E & $\leqslant 0.10$ & $0.30 \sim 0.55$ & $\leqslant 0.03$ & $\leqslant 0.20$ & $\leqslant 0.30$ & $\leqslant 0.02$ & $\leqslant 0.02$ & \\
\hline
H08MnA & $\leqslant 0.10$ & $0.80 \sim 1.10$ & $\leqslant 0.07$ & $\leqslant 0.20$ & $\leqslant 0.30$ & $\leqslant 0.03$ & $\leqslant 0.03$ & 埋弧焊焊丝 \\
\hline
H08Mn2SiA & $\leqslant 0.11$ & $1.80 \sim 2.10$ & $0.65 \sim 0.95$ & $\leqslant 0.20$ & $\leqslant 0.30$ & $\leqslant 0.03$ & $\leqslant 0.03$ & 二氧化碳气体保护焊焊丝 \\
\hline
\end{tabular}
\end{table}
注:焊丝牌号的含义:H是焊(Han)的汉语拼音的第一个字母,表示焊接用钢丝;H后面的数字表示碳的质量分数,A表示高级优质焊接钢丝(S、P的质量分数$\leqslant 0.03\%$);E表示特级优质焊接钢丝(S、P的质量分数$\leqslant 0.02\%$);化学元素后面的数字表示其质量分数($<1\%$不标出)。

焊接合金结构钢、不锈钢的焊条,应采用相应的合金结构钢、不锈钢的焊接钢丝作焊芯。

药皮的作用主要是稳弧、保护、脱氧、渗合金及改善焊接工艺性。焊条药皮中含有钾、钠等元素,能在较低电压下电离,容易引弧并使之稳定燃烧,从而改善焊条的工艺性能(如能减少焊接飞溅、使焊缝成形美观);药皮在高温下熔化,可产生保护熔渣及隔离气体,减少氧和氮侵入金属熔池;药皮中含有锰铁、硅铁等铁合金,在焊接冶金过程中起脱氧、去硫、渗合金等作用。

(2)焊条的分类与型号

焊条牌号是焊条行业统一的代号,焊条型号则是国家标准规定的代号。

国家标准(GB/T 5117—2012,GB/T 5118—2012等)将焊条按化学成分分为7大类,即碳钢焊条、低合金钢焊条、不锈钢焊条、堆焊焊条、铸铁焊条及焊丝、铜及铜合金焊条、铝及铝合金焊条等。其中应用最多的是碳钢焊条和低合金钢焊条。

焊条型号是按熔敷金属的力学性能、药皮类型、焊接位置、电流类型、化学成分和焊后状态划分的。以碳钢焊条为例,首字母E表示焊条;E后面的两位数字表示熔敷金属抗拉强度的最小值;第三位和第四位数字表示药皮类型、焊接位置和电流类型;其后的数字分别表示熔敷金属的化学成分及焊后状态。例如E5515-N5P,“E”表示焊条,“55”表示熔敷金属抗拉强度最小值为550 MPa,“15”表示药皮类型为碱性,适用于全位置焊接,采用直流反接,“N5”表示熔敷金属的化学成分,“P”表示焊后状态为热处理状态。

焊条牌号是一个汉语拼音字母加3位数字组成(字母+3位数字)。例如J422,“J”表示结构钢焊条,“42”表示焊缝金属抗拉强度不低于420 MPa;“2”表示药皮为氧化钛钙型,适用于直流或交流电源。其中,J422、J507和J506分别对应国标GB/T 5117—2012中的E4303、E5015和E5016。

根据焊条药皮性质的不同,结构钢焊条可以分为酸性焊条和碱性焊条两大类。酸性焊条药皮熔渣的主要成分是酸性氧化物(如$\text{SiO}_2$、$\text{TiO}_2$、$\text{Fe}_2\text{O}_3$等),氧化性较强,易烧损合金元素,但其电弧稳定,对焊件上的油污、铁锈不敏感,工艺性较好;酸性熔渣熔点低,流动性好,有利于脱渣和焊缝的形成,但难以有效清除熔池中硫、磷杂质,容易形成偏析,焊缝塑性和韧性稍差,渗合金作用弱,热裂倾向大。酸性焊条适合各种电源,常用于一般钢结构件的焊接。碱性焊条药皮熔渣的主要成分是碱性氧化物(如$\text{CaO}$、$\text{FeO}$、$\text{MnO}$、$\text{MgO}$、$\text{Na}_2\text{O}$等)和铁合金,氧化性弱,脱硫、脱磷能力强,抗裂性好,但对油污、水锈等敏感性较大,易产生气孔,工艺性较差。碱性焊条一般要求采用直流反接,主要用于压力容器等重要结构件的焊接。

药芯焊丝是一种新型焊接材料。药芯焊丝由外皮金属(钢管或钢带)包裹药粉而成,其中药粉相当于手工焊条的药皮。采用药芯焊丝焊接时,生产效率高,焊接质量好,焊接成本低。它继承了手工焊条成分可调的优点,并能实现连续、自动焊接,其熔敷效率高达手工电弧焊的4倍。与二氧化碳气体保护焊相比,药芯焊丝焊接具有可通过调整药芯添加物以适应各种类型钢材焊接的优势,同时克服了二氧化碳气体保护焊飞溅大、工艺性差的缺点,在国内外得到了较为广泛的应用。

(3)焊条的选用原则

选用焊条应首先根据焊件化学成分、力学性能、抗裂性、耐腐蚀性以及高温性能等要求,选用相应的焊条种类,再根据结构形状、受力情况、焊接设备条件和焊条售价来选定具体型号。具体选用时可考虑以下原则。

1)等强度原则

焊接低碳钢和低合金钢时,一般应使焊缝金属与母材等强度,即选用与母材同强度等级的焊条。但应注意,钢材是按屈服强度确定等级的,而结构钢焊条的等级是指焊缝金属抗拉强度的最低保证值。

2)同成分原则

焊接耐热钢、不锈钢等金属材料时,应使焊缝金属的化学成分与母材的化学成分相同或相近,即按母材化学成分选用相应成分的焊条。

3)抗裂纹原则

焊接刚度大、形状复杂、承受动载荷的焊接结构,应选用抗裂性好的碱性焊条,以免在焊接和使用过程中焊接接头产生裂纹。

4)抗气孔原则

受焊接工艺条件的限制,如对焊件接头部位的油污、铁锈等不易清理干净,应选用抗气孔能力强的酸性焊条,以免焊接过程中气体滞留于焊缝中,形成气孔。

5)低成本原则

在满足使用要求的前提下,应选用工艺性能好、成本低和效率高的焊条。

此外,应根据焊件的厚度、焊缝的位置等条件,选用不同直径的焊条。一般焊件越厚,选用的焊条直径越大。

\subsection{焊接接头的组织与性能}

熔焊使焊缝及其附近的母材经历了一个加热和冷却的热循环过程。由于温度分布不均匀,焊缝经历一次复杂的冶金过程,焊缝附近金属受到一次不同类型的热处理,因而会发生相应的组织和性能变化,并直接影响焊接质量。

焊接接头包括焊缝金属和焊接热影响区。

1.焊缝金属

焊接加热时,焊缝处的温度在液相线以上,母材与填充金属形成共同熔池,冷凝后成为铸态组织。在冷却过程中,液态金属从熔池壁向焊缝的中心方向结晶,形成柱状晶组织(图9-3)。焊缝金属的化学成分主要取决于焊芯金属的成分,但也受熔化母材的影响。焊条药皮在焊接过程中具有合金化作用,使焊缝金属的合金元素多于母材,故只要合理选择焊条和焊接规范,焊缝金属的强度一般不低于母材强度。

2.焊接热影响区

在焊接过程中,焊缝两侧金属因焊接热作用而产生组织和性能变化的区域称为焊接热影响区。低碳钢的热影响区分为熔合区、过热区、正火区和部分相变区(图9-4)。


\begin{figure}[H]
\centering
\includegraphics[page=33, viewport={354.0 122.8 526.5 251.3}, clip]{12-07572-001FigRevised.pdf}
\caption{焊缝的组织示意图}
\end{figure}

\begin{figure}[H]
\centering
\includegraphics[page=34, viewport={26.0 532.8 209.5 745.8}, clip]{12-07572-001FigRevised.pdf}
\caption{低碳钢焊接接头组织与性能变化示意图}
\end{figure}


熔合区位于焊缝与母材金属之间,部分金属熔化部分未熔化,又称为半熔化区。加热温度为$1490 \sim 1530\text{ }^\circ\text{C}$,处于固相线和液相线之间,此范围内熔合区的金属成分及组织极不均匀,强度下降,塑性很差,是产生裂纹及局部脆性破坏的发源地。

过热区紧靠着熔合区,加热温度为$1100 \sim 1490\text{ }^\circ\text{C}$。由于温度大大超过$Ac_3$,奥氏体晶粒急剧长大,形成过热组织,使塑性大大降低,冲击韧度下降$25\% \sim 75\%$。

正火区的加热温度为$850 \sim 1100\text{ }^\circ\text{C}$,属于正常的正火加热温度范围。冷却后得到均匀细小的铁素体和珠光体组织,其力学性能优于母材。

部分相变区加热温度为$727 \sim 850\text{ }^\circ\text{C}$。该区只有部分组织发生转变,冷却后组织不均匀,力学性能较差。

从图9-4所示的性能变化曲线可以看出,在焊接热影响区中,熔合区和过热区的性能最差,产生裂纹和局部破坏的倾向也最大。热影响区宽度增加会使焊缝金属的冷却速度减慢,晶粒变粗,并使焊接变形增大。因此,热影响区越窄越好。

热影响区的宽度主要取决于焊接方法和焊接参数。温度高、热量集中的焊接方法,其热影响区小。表9-2所示为用不同焊接方法焊接低碳钢时热影响区的平均尺寸。采用同样焊接方法,选用过大的焊接工艺参数(如电弧焊增大焊接电流、气焊选用过大焊嘴等)和减慢焊接速度,都会使热影响区变宽;反之亦然。

\begin{table}[H]
\centering
\caption{用不同焊接方法焊接低碳钢时热影响区的平均尺寸}
\label{tab:9-2}
\small\begin{tabular}{|l|c|c|c|c|c|c|}
\hline
\multirow{2}{*}{热影响区} & \multicolumn{6}{c|}{焊接方法} \\
\cline{2-7}
 & 气焊 & 手工电弧焊 & 埋弧焊 & 电渣焊 & 等离子弧焊 & 真空电子束焊 \\
\hline
过热区总宽度 & 21.0 & 2.2 & $0.8 \sim 1.2$ & $18 \sim 20$ & $\approx 0$ & $\approx 0$ \\
\hline
热影响区总宽度 & 27.0 & 6.0 & 2.5 & $25 \sim 30$ & $1.4 \sim 2.5$ & $0.05 \sim 0.75$ \\
\hline
\end{tabular}
\end{table}
选用合理的焊接方法和焊接工艺参数(如保证焊透的条件下提高焊速、减小焊接电流),可以缩小热影响区,但在焊接过程中无法消除热影响区。对于重要的焊件,常在焊后进行正火处理,以减弱热影响区的危害。在焊接碳的质量分数和合金元素的质量分数均较高的钢时,可采用焊前预热、焊后热处理等措施,以避免焊接接头的脆性断裂。

\subsection{焊接应力与变形}

1.焊接应力与变形的产生原因

焊接过程中,对焊件进行不均匀加热和冷却,是产生焊接应力和变形的根本原因。

焊接过程是一个极不平衡的热循环过程,即焊缝及其相邻区金属都要由室温被加热到很高的温度(焊缝金属已处于液态),再快速冷却下来。由于在这个热循环过程中,焊件各部分的温度不同,随后的冷却速度也各不相同,因而焊件各部位在热胀冷缩和塑性变形的影响下,必将产生内应力与变形,内应力过大时还会产生裂纹。

焊接应力的存在,对焊件质量、使用性能和焊后机械加工精度都有很大影响,甚至导致整个焊件断裂;对于接触腐蚀性介质的焊件(如容器),由于应力腐蚀现象加剧,焊件的使用寿命将降低,甚至产生应力腐蚀裂纹而报废;焊接变形不仅给装配工作带来很大困难,还会影响焊件的工作性能。变形量超过允许值时必须进行矫正,矫正无效时只能报废。因此,在设计和制造焊接结构时,应尽量减小焊接应力与变形。

2.焊接变形的基本形式

常见的焊接变形有收缩变形、角变形、弯曲变形、扭曲变形和波浪变形等5种形式(图9-5)。

\begin{figure}[H]
\centering
\includegraphics[page=34, viewport={229.0 532.3 523.5 680.8}, clip]{12-07572-001FigRevised.pdf}
\caption{焊接变形的基本形式}
\label{fig:9-5}
\end{figure}

收缩变形是由于焊缝金属沿纵向和横向的焊后收缩而引起的;角变形是由于焊缝截面上、下不对称,焊后沿横向上、下收缩不均匀而引起的;弯曲变形是由于焊缝布置不对称,焊缝较集中的一侧纵向收缩较大而引起的;扭曲变形常常是由于焊接顺序不合理而引起的;波浪变形则是由于薄板焊接后焊缝收缩时,产生较大的收缩应力,使焊件丧失稳定性而引起的。

3.减少和消除焊接应力与变形的措施

减少和消除焊接应力和变形的措施可从结构设计和工艺设计两方面入手。当材料塑性较好、结构刚度较小时,焊件的变形较大,焊接应力较小,此时应主要采取预防和矫正变形的措施,使焊件获得所需的形状和尺寸;当材料塑性较差,结构刚度较大时,焊接变形较小,焊接应力较大,此时应主要采取减小或消除应力的措施,以避免裂纹的产生。生产中常采取下列工艺措施。

(1)预留收缩变形量

根据理论值和经验值,在焊件备料及加工时预先考虑收缩余量,以便焊后工件达到所要求的形状和尺寸。

(2)反变形法

反变形法与防止铸件变形的反变形法原理相同。根据理论计算和经验,预先估计焊件变形的大小和方向,然后在焊接装配时给予一个方向相反、大小相等的人为变形,以抵消焊后产生的变形,使焊件得到正确形状(图9-6)。

\begin{figure}[H]
\centering
\includegraphics[page=34, viewport={94.0 397.8 454.0 505.3}, clip]{12-07572-001FigRevised.pdf}
\caption{防止变形的反变形法}
\label{fig:9-6}
\end{figure}

(3)刚性固定法

焊接时将焊件加以固定,焊后待焊件冷却到室温后再去掉刚性固定,可有效防止角变形和波浪变形,但会增大焊接应力。该方法只适用于焊接塑性较好的低碳钢,不能用于铸铁和淬硬倾向大的钢材,以免焊后断裂。

(4)选择合理的焊接顺序,尽量使焊缝自由收缩

拼焊图9-7所示的钢板时,应先焊错开的短焊缝,再焊直通长焊缝,即先将小块钢板拼焊成较大的钢板条,再将钢板条拼焊成形,使焊接时钢板能自由收缩,以防在焊缝交接处产生裂纹。对称截面梁的焊接顺序如图9-8所示,可有效地减少变形。

(5)锤击焊缝法

在焊缝的冷却过程中,用圆头小锤均匀迅速地锤击焊缝,使金属产生塑性伸长变形,抵消一部分焊接收缩变形,从而减小焊接应力和变形。

(6)加热“减应区”法

焊接前,在工件的适当部位(称为减应区)进行加热使之伸长,焊后冷却时,加热区与焊缝一起收缩,可大大减小焊接应力和变形。

(7)焊前预热和焊后缓冷

焊前预热的目的是减少焊缝区与焊件其他部分的温差,降低焊缝区的冷却速度,使焊件能较均匀地冷却下来,从而减少焊接应力与变形。预热焊件,可以局部加热或整体加热,预热温度一般为$100 \sim 600\text{ }^\circ\text{C}$。焊后缓冷也能起到同样的作用。但这种方法使工艺复杂化,只适用于塑性差、容易产生裂纹的材料,如高、中碳钢,铸铁和合金钢等。

(8)焊后去应力退火

可采用焊后去应力退火方法较彻底地消除焊接应力。此时,可将焊件加热至$500 \sim 650\text{ }^\circ\text{C}$,保温后缓慢冷却至室温。此外,亦可采用振动法消除焊接应力。

4.焊接变形的矫正

在焊接过程中,即使采用了上述措施,有时也会产生超过允许值的焊接变形,因此需要对焊接变形进行矫正。

(1)机械矫正

机械矫正是在焊后通过压力机、矫直机或采用辗压、锤击等方法矫正焊接变形的方法(图9-9)。该方法适用于矫正刚性较小、塑性较好和厚度不大的焊件。

(2)火焰加热矫正

火焰加热矫正利用火焰加热所产生的局部压缩塑性变形,来抵消焊件在该部分已产生的伸长变形。加热火焰通常选用氧-乙炔火焰,加热方式有点状加热、三角加热(图9-10)和条状加热,加热温度一般为$600 \sim 800\text{ }^\circ\text{C}$。该方法不需专用设备,简便、灵活,适用面广,但加热位置、加热面积和加热温度的选择需要借助一定的经验和焊接变形力学知识,否则不仅达不到目的,还会增大原有的变形。


\begin{figure}[H]
\centering
\includegraphics[page=34, viewport={93.5 254.8 277.5 366.8}, clip]{12-07572-001FigRevised.pdf}
\caption{拼接钢板的合理焊接顺序}
\end{figure}

\begin{figure}[H]
\centering
\includegraphics[page=34, viewport={285.5 254.3 449.0 351.8}, clip]{12-07572-001FigRevised.pdf}
\caption{对称截面梁的合理焊接顺序}
\end{figure}

\begin{figure}[H]
\centering
\includegraphics[page=34, viewport={7.0 88.3 167.5 227.3}, clip]{12-07572-001FigRevised.pdf}
\caption{机械矫正法}
\end{figure}

\begin{figure}[H]
\centering
\includegraphics[page=34, viewport={184.5 87.3 330.0 225.8}, clip]{12-07572-001FigRevised.pdf}
\caption{火焰加热矫正法}
\end{figure}


\section{常用焊接方法}

\subsection{电弧焊}

1.手工电弧焊

利用电弧热作为热源的熔焊方法称为电弧焊,简称弧焊。手工电弧焊是利用电弧热局部熔化焊件,并用手工操纵焊条进行焊接的电弧焊方法,是目前应用较为广泛的焊接方法之一。焊接时,焊条与工件之间产生电弧,电弧高温将焊件与焊条局部熔化形成共同熔池(图9-11),当焊条离开后,熔池迅速冷却,凝固形成焊缝,使分离的焊件牢固地连接成整体。

手工电弧焊设备简单,应用灵活、方便,适用面广,可焊接各种焊接位置(图9-12)和直缝、环缝及各种曲线焊缝,尤其适用于操作不便的场合和短小焊缝的焊接。但手工电弧焊对操作者的技能要求较高,生产率低,工作环境差,劳动强度大,且不适宜焊接钛等活泼金属、难熔金属及低熔点金属。


\begin{figure}[H]
\centering
\includegraphics[page=34, viewport={345.0 86.8 534.5 213.3}, clip]{12-07572-001FigRevised.pdf}
\caption{手工电弧焊焊接过程}
\end{figure}

\begin{figure}[H]
\centering
\includegraphics[page=35, viewport={115.5 659.8 415.5 746.3}, clip]{12-07572-001FigRevised.pdf}
\caption{手工电弧焊的焊接位置}
\end{figure}


2.埋弧自动焊

埋弧自动焊是利用连续送进的焊丝在焊剂层下产生电弧而自动焊接的方法(图9-13),是目前应用较广的自动焊接方法之一。它以连续送进的焊丝代替手工电弧焊的焊条,以颗粒状的焊剂代替焊条药皮。焊接时,电弧产生在焊丝和焊件之间,并在$40 \sim 60\text{ mm}$厚的焊剂下燃烧。在电弧高温作用下,焊件、焊丝和焊剂熔化形成熔池与熔渣,熔池和熔滴受熔渣和焊剂蒸汽的保护与空气隔绝。随着电弧向前移动,熔池在熔渣覆盖层下凝固成焊缝(图9-14)。

埋弧自动焊具有下列特点。

(1)生产率高

埋弧自动焊不存在焊条发热问题,允许采用大电流(可高达$1000\text{ A}$以上)焊接,熔敷效率大,焊速高;由于采用盘状焊丝连续送进,无须像手工电弧焊频繁更换焊条;由于在焊剂层下焊接,热能利用率高,对厚度在$25\text{ mm}$以下的焊件,可不开坡口一次焊成,节省了焊接辅助时间。埋弧自动焊比手工电弧焊生产率提高$5 \sim 10$倍甚至更高。


\begin{figure}[H]
\centering
\includegraphics[page=35, viewport={78.5 498.8 276.5 620.3}, clip]{12-07572-001FigRevised.pdf}
\caption{埋弧自动焊焊接过程}
\end{figure}

\begin{figure}[H]
\centering
\includegraphics[page=35, viewport={292.0 496.3 460.0 625.8}, clip]{12-07572-001FigRevised.pdf}
\caption{埋弧自动焊焊缝纵截面图}
\end{figure}


(2)焊缝质量高且稳定

电弧和熔池被封闭在液态熔渣中,保护效果好,焊接工艺参数自动控制,故焊接质量稳定,焊缝成形美观。

(3)劳动条件好

埋弧自动焊没有弧光,焊接烟尘较少,机械化操作减轻了劳动强度。

埋弧自动焊常用的焊丝有H08、H08A和H08MnA等。焊剂有熔炼焊剂和烧结焊剂两大类。埋弧自动焊适用于大批量生产,可焊接中、厚钢板($6 \sim 60\text{ mm}$)的直焊缝或大直径环缝,平焊位置的对接、搭接和T形接头,还可用于在基体金属表面堆焊耐磨、耐蚀合金。

3.气体保护焊

气体保护焊是用外加气体作为电弧介质并保护电弧和焊接区的电弧焊。常用焊接方法有氩弧焊和二氧化碳气体保护焊。

(1)氩弧焊

氩弧焊是以氩气作为保护气体的气体保护焊。氩气是一种惰性气体,它既不与焊缝金属发生化学反应,又不溶于液态金属,是理想的保护气体。

按使用电极不同,氩弧焊可分为不熔化极氩弧焊和熔化极氩弧焊。

不熔化极氩弧焊采用高熔点的铈钨棒作电极,故又称为钨极氩弧焊。焊接时,钨极不熔化,仅起引弧和维持电弧的作用,需另加焊丝作为填充金属(图9-15a)。在电弧高温作用下,填充金属与焊件熔融在一起形成焊接接头。该方法易于实现机械化、自动化焊接,但因通过电极的焊接电流不宜过大,通常用于焊接厚度为$6\text{ mm}$以下的薄板。

熔化极氩弧焊(图9-15b)以连续送进的金属焊丝作电极和填充金属,可采用较大的焊接电流。为使电弧稳定,通常采用直流反接。该方法适宜焊接$25\text{ mm}$以下的板材。

氩弧焊的保护效果好,电弧稳定,在氩气流的压缩作用下电弧热量集中。因此,氩弧焊的热影响区小,焊后变形小,焊缝外形光洁美观,无渣壳,便于实现机械化和自动化。但氩气成本高,氩弧焊设备复杂,目前主要用于铝、镁、钛及其合金和不锈钢的焊接,有时也用于合金钢的焊接。

(2)二氧化碳气体保护焊

二氧化碳气体保护焊是以$\text{CO}_2$作为保护气体的气体保护焊。其焊接装置(图9-16)与熔化极氩弧焊相似,也是以连续送进的金属焊丝为电极。$\text{CO}_2$气体从喷嘴中连续喷出,在电弧周围形成局部气体保护层。该方法宜采用自动或半自动方式进行焊接。

\begin{figure}[H]
	\centering
	\includegraphics[page=35, viewport={61.5 289.8 485.5 464.8}, clip]{12-07572-001FigRevised.pdf}
	\caption{氩弧焊示意图}
\end{figure}

由于 $\text{CO}_2$ 气体是氧化性气体，高温时可分解成 $\text{CO}$ 和氧原子，易造成合金元素烧损、焊缝吸氧，导致电弧稳定性差、金属飞溅等缺点，因而二氧化碳气体保护焊多配用含锰、硅元素的焊丝进行脱氧和渗合金，并使用直流电源，以使电弧稳定。

二氧化碳气体保护焊的成本低；$\text{CO}_2$ 电弧穿透力强，熔敷速度快，生产率比手工电弧焊高 $1\sim4$ 倍；抗锈力强，抗裂性能好；可用于低碳钢、低合金钢、耐热钢和不锈钢的焊接。

\begin{figure}[H]
	\centering
	\includegraphics[page=35, viewport={146.5 147.8 383.0 276.8}, clip]{12-07572-001FigRevised.pdf}
	\caption{二氧化碳气体保护焊示意图}
\end{figure}

\subsection{其他焊接方法}

1. 电阻焊

电阻焊是利用电流通过焊件接触面及邻近区域产生的电阻热，将接头接触处加热到塑性状态或局部熔化状态，在压力下形成牢固接头的焊接方法。

电阻焊的焊接电流很大，焊接时间很短（百分之一秒到几秒），生产率高，不用填充金属，不用保护，易于实现机械化、自动化，但对焊件厚度和接头形式有一定限制。

电阻焊的基本形式有点焊、缝焊和对焊，如图 9-17 所示。点焊是利用柱状电极加压通电，在搭接工件接触面之间形成一个个焊点的焊接方法。点焊适于焊接薄板壳体或型钢构件。缝焊过程与点焊相似，它用旋转的圆盘状滚动电极代替了柱状电极。焊接时，圆盘状电极压紧工件并转动（同时焊件向前移动），配合断续通电，即形成连续重叠的焊点。缝焊适于焊接有密封要求的薄壁容器。对焊是利用电阻热使两个工件在整个接触面上焊接起来的一种方法，它可分为电阻对焊和闪光对焊。对焊适于焊接杆状零件，如刀具、钢筋等。

2. 钎焊

钎焊是用比焊件熔点低的金属作钎料，将焊件和钎料加热到适当温度，焊件不熔化，钎料熔化后填满接头间隙，并与焊件相互扩散，冷凝后将焊件连接起来的焊接方法，如铜锌钎焊、锡铅钎焊等。

\begin{figure}[H]
	\centering
	\includegraphics[page=35, viewport={93.0 17.3 439.0 130.3}, clip]{12-07572-001FigRevised.pdf}
	\caption{电阻焊的基本形式}
\end{figure}

在钎焊过程中一般需使用钎剂，以去除工件表面的氧化膜和油污等杂质，保护母材接触面和钎料不受氧化，并增加钎料浸润性和毛细流动性。钎焊一般采用搭接接头和套件镶接，以弥补钎焊强度的不足。

根据钎料熔点高低，钎焊可分为硬钎焊（钎料熔点高于 $450\text{ }^\circ\text{C}$）和软钎焊（钎料熔点低于 $450\text{ }^\circ\text{C}$）。硬钎焊接头强度较高（大于 $200\text{ MPa}$），常用于结构件的焊接，如硬质合金刀片在其刀体上的焊接；软钎焊接头强度较低（小于 $70\text{ MPa}$），但导电性优异，在电子工业上的应用极为广泛，如各种生产线上应用的波峰焊（机）、再流焊（机）及气体保护助钎技术等。

气体保护助钎技术是在保护气体下进行的软钎焊。保护气体一般为纯惰性气体或混合气体，如 $\text{N}_2$、$\text{Ar}$、$\text{N}_2+\text{H}_2$ 及 $\text{Ar}+\text{H}_2$ 气体等。采用气体保护，可提高钎焊质量，并免去钎焊后的清洗工序。因此，气体保护助钎技术在印制电路板的组装中获得了广泛采用。

3. 电渣焊

电渣焊是利用电流通过液体熔渣所产生的电阻热进行熔焊的方法，其焊接过程如图 9-18 所示。将两焊件相距 $20\sim60\text{ mm}$ 垂直放置，焊接时，先在引弧板上撒上焊剂，焊丝伸入焊剂中，通电引弧，焊剂受热熔化，形成渣池。随即将焊丝插入渣池，电弧熄灭，转变为电渣。电流通过渣池产生的电阻热将不断送入的焊丝熔化，沉积于渣池下部，形成熔池。随焊丝送进，熔池液面升高，冷却滑块上移，熔池不断凝固形成焊缝。电渣焊可一次焊接较大厚度工件（通常用于板厚 $40\text{ mm}$ 以上的工件，最大厚度可达 $2\text{ m}$），生产效率比电弧焊高。另外，电渣焊焊缝金属较纯净，焊接时不开坡口，节省钢材和焊接材料，因此经济效益好。电渣焊还可实现“以焊代铸”“以焊代锻”，减轻结构质量。其缺点是焊接接头晶粒粗大，对于重要结构，可通过焊后热处理（正火）来细化晶粒，改善其力学性能。电渣焊主要用于碳钢、低合金钢、不锈钢等厚大工件的立焊位置焊接，以及铸-焊、锻-焊组合构件的焊接。

\begin{figure}[H]
	\centering
	\includegraphics[page=36, viewport={36.0 495.8 232.5 728.8}, clip]{12-07572-001FigRevised.pdf}
	\caption{电渣焊过程示意图}
\end{figure}


4. 摩擦焊

摩擦焊是利用焊件表面相互摩擦所产生的热，使端面达到热塑性状态时迅速顶锻的一种压焊方法。其特点是表面清理要求不高、焊接质量好、生产率高和易实现自动化等，尤其适于异种材料的焊接，如各种铝-铜过渡接头、铜-不锈钢接头、石油钻杆、电站锅炉蛇形管和阀门等。但设备投资大，工件必须有一个是回转体（作回转运动），不宜焊接摩擦因数小的材料和惰性材料。

5. 电子束焊

电子束焊是以电子束为热源的熔焊方法。在真空环境中，从炽热阴极发射的电子被高压静电场加速，并经磁场聚集成高能量密度的电子束，以极高的速度轰击焊件表面，由于电子运动受阻而被制动，遂将动能转化为热能而使焊件瞬间熔化，从而形成牢固的接头。其特点是焊接速度快、焊缝深而窄、热影响区和焊接变形很小，焊缝质量高，可以焊接其他焊接工艺难以焊接的形状复杂的焊件，能焊接特种金属和难熔金属，也适于异种金属及金属与非金属的焊接。

6. 等离子弧焊

等离子弧焊是利用压缩电弧（等离子弧）作热源的金属极气体保护焊。经强迫压缩后的电弧弧柱中的气体充分电离，形成高温、高能量的等离子弧，从而熔化金属并形成焊接接头。等离子弧焊能量密度大，弧柱温度极高，因而穿透力强，焊接质量好，热影响区小，焊接变形小。等离子弧焊可用于碳钢、合金钢、耐热钢、不锈钢、铜、镍、钛合金及其极薄板（$0.025\text{ mm}$）的焊接；在氩气保护下可焊接钨、钼、钽、铌、锆及其合金。

7. 激光焊

激光焊是以聚焦的激光束作为热源轰击焊件所产生的热量进行焊接的方法。其特点是焊缝窄，热影响区和变形小；激光在大气中能远距离传射到焊件上，不像电子束那样需要真空室，但穿透能力不及电子束。激光焊可进行同种金属或异种金属间的焊接，其中包括各种钢、铝、铜、银、钛、铌、镍、铬、钼以及难熔金属材料等，甚至还可焊接玻璃钢等非金属材料。

8. 爆炸焊

爆炸焊是利用炸药爆炸产生的巨大冲击力和热量使焊件的迅速碰撞而实现连接的焊接方法。任何具有足够强度和塑性并能承受工艺过程所要求的快速变形的金属，均可以进行爆炸焊。爆炸焊接头强度高，热影响区小，主要用于材料性能差异大而用其他方法难以焊接的场合，如铝-钢、钛-不锈钢、钽、锆等的焊接，也可用于制造复合板。爆炸焊无需专用设备，工件形状、尺寸不限，但以平板、圆柱、圆锥形为宜。

9. 扩散焊

扩散焊是在一定的温度和压力下使待焊表面相互接触，通过微观塑性变形，经较长时间的原子相互扩散而实现连接的焊接方法。扩散焊接头质量好，焊接变形小，可以焊接非金属材料和异种材料（如焊接同种或异种金属、陶瓷与金属等），还可制造多层复合材料，也可用于焊接不适于熔焊的材料（如钼、钨等）。

10. 磁力脉冲焊

磁力脉冲焊是利用被焊工件之间脉冲磁场相互作用而产生冲击作用来实现金属之间连接的焊接方法。其作用原理与爆炸焊相似，可用来焊接薄壁管材、异种金属如铜-铝、铝-不锈钢、铜-不锈钢、锆-不锈钢等。

11. 超声波焊

超声波焊是利用超声波的高频振荡能对焊件接头进行局部加热和表面清理，再施加压力实现焊接的一种压焊方法。因焊接过程中无电流流经焊件，也无火焰、电弧等热源作用，所以焊件表面无变形、无热影响区，表面无须严格清理，焊接质量好。超声波焊适于厚度小于 $0.5\text{ mm}$ 的工件焊接，尤其适于异种材料的焊接，但其功率小，应用受到一定限制。

\subsection{焊接技术的发展趋势}

1. 焊接机器人的开发与应用

工业发达国家在机器人研究方面取得了许多成果，机器人已广泛用于工业生产。而焊接机器人在机器人应用中占有很大分量，目前焊接机器人多为固定位置的手臂式机械。研究与开发智能机器人，特别是具有自动路径规划、自动校正轨迹、自动控制熔深的机器人是今后的重点研究方向。

2. 灵巧智能型焊接机

要真正满足工业需求和提高焊接质量，焊接机必须具有更高的智能。研究重点主要有：研究设计各种运动机构，以实现焊接电弧的运动，焊丝送进和焊炬姿态的调节；研究各种不同的传感器，以探测焊缝坡口形状与位置、焊接温度场等；研究传感器的信号处理方法及最佳控制方法，包括线性控制、非线性控制和模糊控制、神经网络控制等。

3. 计算机模拟技术与控制技术

研究焊接技术的模式应从传统的“理论-试验-生产”模式逐渐向“理论-计算机仿真-生产”模式转化。应建立焊接工程数据库，研究开发温度场、应力场、应变场、残余应力等的测试装置和方法，加快计算机辅助焊接过程控制的应用。在焊接工艺参数的控制中，引入神经网络和模糊控制，以实现焊接过程和质量的数字化、集成化和智能化控制。

4. 发展无损探伤技术及可靠性评估的方法及理论

重点研究包括：断裂失效准则的完善；焊件断裂失效的微观机理研究；可靠性评估的计算机软件及专家系统的开发研究；延寿金属技术的研究，如喷丸、锤击焊趾重熔、内应力消除、表面处理等。

5. 发展恶劣条件下的焊接技术

恶劣条件主要是指高温、放射性、水下及太空等条件。核能工业、海洋工业以及空间技术的发展，必然要求在放射性、水下及太空进行焊接或焊补。恶劣条件下的焊接技术研究内容主要包括：为太空焊接开发新的焊接工艺；为放射性及深水条件下的焊接研究远距离操纵的自动化焊接设备；为水下焊接建立深水焊接模拟中心；研究高气压的电弧及焊接质量。

6. 焊接装备的大型化和组体化

焊接操作机、配套焊件变位机械、搬运机械与传输辊道的组合联动，形成了制造系统、焊接生产线或焊接中心，如中重型厚壁容器焊接中心、集装箱外壳整体组装焊接中心、汽轮机导流隔板柔性制造系统、箱形梁焊接生产线、机车车厢总装组焊中心等。

7. 管理控制一体化

通过企业的局域网，利用 $\text{CAD/CAM/CAE/CAPP/PDM}$ 等仿真软件，将生产管理与制造系统实行全集成自动化控制，实现脱机编程，远程监控、诊断和检修。

\section{常用金属材料的焊接}

\subsection{材料的焊接性}

焊接性是指金属材料对焊接工艺的适应性，即金属材料在一定的焊接工艺条件下（焊接方法、焊接材料、焊接工艺参数和结构形式等），获得优质焊接接头的难易程度。它包括两方面的内容：一是结合性能，即在一定焊接工艺条件下，焊接接头产生焊接缺陷的敏感性；二是使用性能，即焊接接头在使用过程中的可靠性。

影响金属焊接性的因素很多，可归类为材料（化学成分、组织状态、力学性能等）、设计（结构形式）、工艺（焊接方法、焊接工艺参数等）及正常使用环境（工作温度、负荷条件、工作环境等）4 个方面。金属材料的焊接性可通过估算或试验确定。

焊接性主要取决于金属材料的化学成分、物理特性等。生产中，常根据钢材的化学成分估计其焊接性。钢的碳质量分数对其焊接性的影响最为显著，因而把钢中其他合金元素的质量分数对焊接性的影响折算成碳的影响。这些合金元素的折算量与碳的质量分数的总和称为碳当量 $C_{\text{E}}$。碳当量 $C_{\text{E}}$ 常作为评定钢材焊接性的一种参考指标。

国际焊接学会推荐的碳素结构钢和低合金结构钢的碳当量 $C_{\text{E}}$ 的经验公式为
\[
C_{\text{E}} = w_{\text{C}} + \frac{w_{\text{Mn}}}{6} + \frac{w_{\text{Ni}} + w_{\text{Cu}}}{15} + \frac{w_{\text{Cr}} + w_{\text{Mo}} + w_{\text{V}}}{5}
\]
式中，化学元素的质量分数均取其成分范围的上限。

实践表明，碳当量 $C_{\text{E}}$ 越高，钢材的焊接性越差。

当 $C_{\text{E}} < 0.4\%$ 时，钢材的塑性良好，焊接性良好。塑性变形能够减缓焊件中产生的应力，减少产生裂纹的倾向，焊接时，一般不需要采取特殊的工艺措施。

当 $C_{\text{E}} = 0.4\% \sim 0.6\%$ 时，钢材的塑性稍差，焊接性也较差。焊接时，需采取预热、缓冷等工艺措施，以防止裂纹产生。

当 $C_{\text{E}} > 0.6\%$ 时，钢材的塑性差，焊接性不好。焊前必须预热到较高温度，采取减少焊接应力和防止裂纹的严格工艺措施，焊后进行适当热处理，才能保证焊接质量。

上述碳当量法由于没有考虑焊接方法、结构刚度、焊件厚度、环境温度等因素的影响，故只能粗略估算焊接性。对钢材实际焊接性，应考虑焊接具体条件，通过焊接试验来确定。

\subsection{钢材的焊接}

1. 低碳钢的焊接

低碳钢的碳当量 $C_{\text{E}} < 0.4\%$，其塑性好，一般无淬硬倾向，焊接性良好，容易获得优质的焊接接头，几乎适用所有的焊接方法，通常不需要采取特殊的工艺措施。较厚的低碳钢焊件但在 $0\text{ }^\circ\text{C}$ 以下低温焊接时需预热；对厚度超过 $50\text{ mm}$ 的低碳钢焊件，焊后需进行热处理，以消除内应力；电渣焊后的低碳钢焊件需进行正火处理，以细化晶粒。

低碳钢焊接的常用焊接方法有手工电弧焊、埋弧自动焊、电渣焊、二氧化碳气体保护焊和电阻焊。手工电弧焊焊条一般选用 $\text{E4303 (J422)}$、$\text{E4313 (J421)}$ 或 $\text{E4320 (J424)}$。对于重要结构件，则应选用 $\text{E4315 (J427)}$、$\text{E4316 (J426)}$ 或 $\text{E5015 (J507)}$、$\text{E5016 (J506)}$ 等低氢型焊条。

2. 中、高碳钢的焊接

中碳钢的碳当量 $C_{\text{E}}$ 约为 $0.4\%$，其塑性较差，易产生淬硬组织和裂纹，焊接性较差，焊接时需要采取特殊的工艺措施。中碳钢的焊接通常选用抗裂性能好的低氢型焊条，焊前需进行预热。碳的质量分数小于 $0.45\%$ 的钢，其预热温度为 $150\sim250\text{ }^\circ\text{C}$，碳的质量分数较高或厚度较大的钢，其预热温度应提高至 $250\sim400\text{ }^\circ\text{C}$，采用细焊条、小电流、开坡口、多层焊，以减小母材的熔化深度，减缓热影响区的冷却速度；焊后缓冷，必要时应进行去应力退火。

高碳钢的碳当量 $C_{\text{E}} > 0.4\%$，焊接性差，一般不用于制造焊接结构，仅对损坏的机件进行焊补。高碳钢的焊接可采用手工电弧焊和气焊，其焊前预热温度为 $250\sim350\text{ }^\circ\text{C}$，刚度大的焊件在焊接过程中还应保持此温度，焊后应缓慢冷却。

3. 低合金结构钢的焊接

当低合金结构钢的 $R_{\text{m}} < 400\text{ MPa}$ 时，其碳当量 $C_{\text{E}} < 0.4\%$，焊接性良好，焊接时不需要采取特殊的工艺措施。但其在低温下焊接或厚度较大时，焊前应进行预热。

对于 $R_{\text{m}} > 400\text{ MPa}$ 的低合金结构钢，其碳当量 $C_{\text{E}} > 0.4\%$，焊接性较差，焊前一般要预热（预热温度 $> 150\text{ }^\circ\text{C}$），焊接时应调整焊接工艺参数来严格控制热影响区的冷却速度，焊后进行去应力退火。

目前，在焊接生产中广泛采用 $16\text{Mn}$ 钢，其焊接性接近低碳钢，但因锰的质量分数较高，在近缝区易出现晶粒粗大，冷却速度快时有淬硬现象，故在焊接厚度较大（$> 16\text{ mm}$）或低温（$< -5\text{ }^\circ\text{C}$）下施焊时，应进行温度为 $150\sim250\text{ }^\circ\text{C}$ 的预热。

各种牌号的普通低合金结构钢的焊接常采用手工电弧焊、埋弧自动焊和二氧化碳气体保护焊。

\subsection{铸铁的焊补}

1. 铸铁的焊接性

铸铁的碳质量分数高，含硫、磷等杂质较多，塑性差，焊接性也很差，主要原因有以下几点。

（1）焊接接头易产生白口铸铁组织

碳和硅是促进石墨化的元素，焊接时会大量烧损，焊后冷却速度又快，不利于石墨的析出，故焊接接头容易产生白口铸铁组织，给切削加工带来困难。

（2）易产生焊接裂纹

铸铁是脆性材料，焊接时容易产生白口和淬硬组织。当焊接应力超过铸铁的抗拉强度时，焊件就会在焊缝或近缝区产生裂纹，甚至完全开裂。

（3）易产生气孔和夹渣

铸铁中的碳、硅元素剧烈氧化，形成 $\text{CO}$、$\text{CO}_2$ 气体和硅酸盐熔渣，它们滞留在焊缝中会形成气孔和夹渣等缺陷。

铸铁不宜作为焊接结构材料，铸铁的焊接主要用于其局部损坏和铸造缺陷的焊补修复。

2. 铸铁的焊补方法

铸铁件的焊补通常采用气焊和手工电弧焊，要求不高时也可采用钎焊。按照焊前焊件是否预热，铸铁的焊补方法可分为热焊法和冷焊法。

（1）热焊法

热焊法是把工件整体或局部预热到 $600\sim700\text{ }^\circ\text{C}$ 再进行焊接的方法。在焊接过程中，焊件温度应不低于 $400\text{ }^\circ\text{C}$，焊后缓慢冷却或进行去应力退火。热焊法可防止焊件产生马氏体、白口铸铁组织和裂纹，焊补质量好，焊补处可进行机械加工。但热焊法需加热设备，成本较高，生产率低，劳动条件差。热焊法一般用于焊补形状复杂的重要铸件，如机床导轨和内燃机气缸体等。生产中常采用加热“减应区”法来提高焊补质量。

用气焊进行铸铁热焊比较方便，也可用涂有药皮的铸铁焊条进行手工电弧焊焊补。

（2）冷焊法

冷焊是焊前不预热或只进行 $400\text{ }^\circ\text{C}$ 以下的低温预热的焊接方法。冷焊法操作简便，劳动条件好，生产率高，但焊补质量不如热焊法，焊接处切削加工性较差，主要用于焊补要求不高的铸件及非加工表面。

常用的冷焊铸铁焊条有钢芯铸铁焊条、铸铁芯铸铁焊条、镍基铸铁焊条和铜基铸铁焊条等，其中镍基铸铁焊条、铜基铸铁焊条对防止白口和裂纹的效果较好。焊接时，一般采用细焊条、小电流、短弧、窄焊缝、分段焊（每段焊缝长度不超过 $50\text{ mm}$）及焊后锤击焊缝等措施，以减小热影响区的范围，防止热变形和热应力，避免裂纹和开裂。

\subsection{非铁金属的焊接}

1. 铜及铜合金的焊接

（1）铜及铜合金的焊接性

通常铜及其合金的焊接性很差，主要有如下原因。

1）裂纹倾向大

铜在高温下易氧化生成 $\text{Cu}_2\text{O}$，它与铜形成低熔点的脆性共晶体分布在晶界上，又因铜的热膨胀系数大，焊接时极易引起热裂纹。

2）气孔倾向大

液态铜溶氢能力强，凝固时其溶解度下降很快，来不及逸出的氢气存在于焊缝中形成气孔。

3）容易产生焊不透缺陷

铜的导热性很好，热容量大，焊接时热量极易散失。因此，焊件应焊前预热，焊接中要采用较大热量（如较大电流或火焰），否则不易焊透。但这样会增加热影响区宽度，降低焊缝的力学性能。

4）合金元素易氧化

铜合金中的合金元素，如黄铜中的锌、铝青铜中的铝等，比铜更易氧化，从而降低焊接接头的性能，并促使热裂纹、气孔、夹渣等缺陷的产生。

（2）铜及铜合金的焊接方法

铜及铜合金的焊接常用氩弧焊、气焊、手工电弧焊和钎焊等方法，以氩弧焊的焊接质量最好。选择焊接方法时应综合考虑焊件材料的成分、厚度、结构特点及使用性能等。焊件厚度小于 $5\text{ mm}$ 时可选用气焊，尤其是黄铜，其焊接目前主要采用气焊；紫铜和青铜的焊接采用氩弧焊较好，不仅能获得优质焊缝，还有利于减少焊接变形。

对于铜及其合金的焊接，焊接前应严格清理焊件，以减少氢的来源；焊前预热，以弥补热传导损失；焊后锤击焊缝及进行完全退火，以细化晶粒和提高焊接质量。

2. 铝及铝合金的焊接

（1）铝及铝合金的焊接性

通常，铝及铝合金的焊接性很差，其主要原因如下。

1）极易氧化

铝对氧的亲和力很强，生成的氧化膜组织致密，熔点很高，覆盖在熔池表面，阻碍相互熔合，易造成焊缝夹渣。

2）易产生气孔

液态铝能溶解大量的氢，而凝固过程中溶解度下降很快，来不及逸出的氢气残存在焊缝中形成气孔。

3）易产生裂纹

高温时铝的强度低，塑性差，且线膨胀系数大，因此焊件容易产生应力、变形和裂纹。

（2）铝及铝合金的焊接方法

铝及铝合金的焊接通常采用氩弧焊、电阻焊、钎焊和气焊等。氩弧焊的保护作用最好，没有熔渣，一般采用直流电源反接法。要求不高的焊件可采用气焊。

对于铝及其合金的焊接，焊接前需严格清理焊件表面的油污及杂质；焊后还应清除残留在焊件上的熔剂，以防腐蚀。

对于硬铝，由于其焊接性很差，到目前为止，用一般焊接方法包括氩弧焊，都不易获得满意的焊接接头，目前多采用铆接方法来连接。

\subsection{难熔金属及其合金的焊接}

难熔金属如钛、锆、钼、铌等在航空、航天、原子能、电子工业和民用工业中，均具有广泛的用途。但难熔金属的化学活性极强，采用一般焊接方法焊接时，其会强烈吸收氧、氢和氮等气体，使接头性能恶化。因此，其焊接性能很差，必须采用特种焊接方法。

难熔金属的焊接通常采用氩弧焊、等离子弧焊和电子束焊等焊接方法。

钛及钛合金要在高纯度的氩气保护下进行焊接。焊前通过真空退火使焊丝和母材脱气；焊接时需用氩气流对焊缝根部和尚未冷却到 $350\text{ }^\circ\text{C}$ 的焊缝区进行附加保护。此外，钛合金还有形成冷裂纹的倾向，重要部件应在可控氩气室中焊接，也可在充氩气室内进行。钛及其合金也可采用等离子弧焊和电子束焊进行焊接。

锆的焊接性与钛很接近，故其焊接工艺与钛相似。

钼、铌及其合金，与钛相比更容易在焊接时吸收氧氢和氮等气体，特别是氧气。当氧的质量分数超过 $0.01\%$ 时，其塑性会急剧下降。钼、铌及其合金应在可控氩气室中采用电弧焊焊接或在真空中采用电子束焊焊接。

\subsection{非金属材料的焊接}

陶瓷材料、高分子材料和复合材料以其有别于金属材料的明显特性在现代工业中起着重要作用，具有广泛的用途。其焊接方法也与金属材料明显不同。

1. 陶瓷的焊接

陶瓷的熔点高、硬度与强度高、不易变形，而热导率和电导率很低，因此其焊接性很差，用常规焊接方法难以焊接。目前，陶瓷材料常用的焊接方法包括钎焊、固相扩散焊、超声波压焊、摩擦压焊、过渡液相连接、微波连接、黏接、反应成形与反应烧结连接等方法，其中钎焊和固相扩散焊比较成熟。

钎焊是连接陶瓷与陶瓷或陶瓷与金属最常用的方法之一。陶瓷的钎焊可分为两步钎焊法和直接钎焊法。

两步钎焊法是将陶瓷表面预金属化后再进行普通钎焊的方法。常用的陶瓷预金属化的方法很多，有传统的氧化铝陶瓷封装的 Mo-Mn 法、物理气相沉积法、化学气相沉积法、离子注入法和热喷涂法等。例如，将活性元素 Ti 用离子注入法直接注入 $\mathrm{Al}_{2}\mathrm{O}_{3}$ 陶瓷表面，形成 $50 \sim 100 \mathrm{~nm}$ 深的预金属化层；又如用低压等离子弧喷涂的方法在 $\mathrm{Si}_{3}\mathrm{N}_{4}$ 陶瓷表面喷涂两层约 $200 \mathrm{~\mu m}$ 厚的铝等。经过预金属化的陶瓷表面可与金属钎料产生良好的浸润，提高了钎焊接头的力学性能。

直接钎焊法是在钎料中加入过渡金属 Ti、Zr、Hf、Nb、Ta 及 In 等活性元素形成活性钎料，如 Ag、Cu 或 Ag-Cu 钎料中常加质量分数为 $1 \% \sim 5 \%$ 的 Ti。这些活性元素与陶瓷通过化学反应形成由金属和陶瓷的复合物组成的反应层，该层可与熔化的金属钎料互相浸润，形成较好的钎焊接头。钎焊时应对活性元素进行保护，生产中常用真空或高纯气体来保护。目前，应用较多的陶瓷 $\mathrm{ZrO}_{2}$、$\mathrm{Al}_{2}\mathrm{O}_{3}$、$\mathrm{Si}_{3}\mathrm{N}_{4}$、$\mathrm{SiC}$ 和 AlN 等，均可由 Ag 基、Cu 基、Ag-Cu 基钎料来进行钎焊。

2. 高分子材料的焊接

高分子材料按其热行为可分为热固性高分子材料和热塑性高分子材料。

热固性高分子材料是在一定条件下加入固化剂通过交联固化反应，形成不熔不溶的三维网络结构，在过高的温度下只会发生碳化破坏。因此，固化后的热固性高分子材料无法进行熔化连接，只能采用黏接和机械连接。

热塑性高分子材料在加热时会变软或熔化，可进行熔化焊接。热塑性高分子材料常用的熔化焊接方法包括热气焊、热板焊、内藏加热元体（插入物）的电阻加热或感应加热焊、红外或激光焊、介电或微波加热焊、摩擦焊和超声波焊等。热气焊和热板焊是两种应用最为广泛的热塑性高分子材料焊接方法。热气焊采用热气流作为热源来熔化高分子材料，同时也可熔化填充棒料来形成焊接接头。热板焊焊接时将作为热源的热板置于待连接件间，热板直接与待连接件表面接触，同时将两个待连接件表面加热至软化，再迅速移出热板，对零件进行加压焊接。

3. 复合材料的焊接

复合材料通常可分为树脂基复合材料、金属基复合材料、陶瓷基复合材料和 C/C 复合材料。树脂基复合材料之间的焊接与同基体的高分子材料的焊接方法相同，可采用热气焊、热板焊、内置加热元体（插入物）的电阻加热或感应加热焊、红外或激光焊、介电或微波加热焊、摩擦焊和超声波焊等。树脂基复合材料与金属之间的焊接可采用内藏加热元体（插入物）的电阻加热焊工艺，如高性能的碳纤维增强热塑性树脂聚醚醚酮复合材料（C/PEEK）与铝合金的焊接，采用固结在两层 PEEK/S2 玻璃纤维预浸料中的 $0.18 \mathrm{~mm}$ 厚的不锈钢网作为电阻加热元件，放在 C/PEEK 和铝合金搭接接头的连接界面间。两层 PEEK/S2 预浸料作为不锈钢网的托架，并使不锈钢网与铝合金之间绝缘。加热元件的两端接到电阻焊设备上，通电加热并加足够压力即可形成焊接接头。

按焊接特征，金属基复合材料可分为非连续纤维增强金属基复合材料和连续纤维增强金属基复合材料。非连续纤维增强金属基复合材料的焊接可采用钨极惰性气体保护焊（tungsten inert gas welding, TIG）、扩散连接（中间层固态扩散连接和过渡液相扩散连接）、摩擦焊、电子束焊、电容放电焊接、电阻点焊、钎焊等；连续纤维增强金属基复合材料的焊接较前者困难，通常可采用激光焊、扩散连接、钎焊等。

C/C 复合材料之间的焊接可采用扩散连接（可用石墨、硼、钛、硼化物和碳化物等作中间层）、钎焊（可用 Si、Al、$\mathrm{Mg}_{2}\mathrm{Si}$ 和玻璃作为填充材料）等。C/C 复合材料与金属之间的焊接可采用钎焊，如 C/C 复合材料与铜之间的焊接可采用 Ag-Cu-Ti 为钎料的真空钎焊或用 Ti 的扩散钎焊。

陶瓷基复合材料的焊接工艺尚处于开发阶段，还不很成熟。其焊接可采用微波连接、过渡液相连接、反应成形法连接、原位或自蔓延高温合成法连接、无压固相反应连接等焊接方法，也可采用金属或玻璃作钎料的钎焊等。其中，钎焊工艺研究较多，相对成熟一些。

\section{焊接结构设计}

\subsection{焊接结构材料的选择}

在满足使用性能的前提下，应选用焊接性好的材料来制造焊接结构。一般来说，碳的质量分数 $<0.2 \%$ 的低碳钢和碳当量 $C_{\mathrm{E}}<0.4 \%$ 的低合金钢，都具有良好的焊接性，应优先选用。碳的质量分数 $>0.5 \%$ 的碳钢和碳当量 $C_{\mathrm{E}}>0.4 \%$ 的合金钢，焊接性不好，一般不宜选用，如必须采用，则应采取必要措施，以保证焊接质量。

在满足焊接性的条件下，应优先选用强度等级低的低合金结构钢。这样，既可减轻焊接结构重量，节约钢材，也可延长使用寿命。

异种材料的焊接因其焊接性不同，焊接时易产生较大的焊接应力及裂纹倾向，故应尽量不用，而应尽量选用同种材料进行焊接。但如对焊接结构使用性非常有利并有明显经济效益时，也可选用异种材料的焊接，但应采取必要的工艺措施以保证焊接质量。

设计焊接结构时，应尽量选用槽钢、角钢等各种型材，以减少焊缝数量，简化焊接工艺，同时也可增加结构件的强度和刚性。大批量生产形状复杂的薄壁焊接结构时，应尽量设计成冲压-焊接组合结构。

\subsection{焊接方法的选择}

应根据材料的焊接性、焊件厚度、生产率要求、各种焊接方法的适用范围和现场设备条件等综合考虑来选择焊接方法。

对于焊接性良好的低碳钢，可根据其焊件厚度、生产率等确定具体的焊接方法。如焊件为中等厚度（$10 \sim 20 \mathrm{~mm}$），可采用手工电弧焊、埋弧自动焊、气体保护焊等。氩弧焊成本较高，一般不选用。对于薄板轻型结构，无密封要求时，可优先采用生产率较高的点焊；无点焊设备时，可考虑气焊和手工电弧焊；若要求密封，可采用缝焊。对于厚度为 $35 \mathrm{~mm}$ 以上的重要结构，若条件许可应采用电渣焊。对于要求对接的棒材、管材、型材等材料的焊件，宜用电阻对焊或摩擦焊。

合金钢、不锈钢重要焊件的焊接，应采用氩弧焊以保证焊接质量。铝合金焊件的焊接，应优先选用氩弧焊，如无氩弧焊设备，也可考虑采用气焊。稀有金属或高熔点金属的特殊焊件的焊接，可采用等离子弧焊、真空电子束焊等焊接方法。

\subsection{焊接接头形式的设计}

焊接接头与坡口形式的选择，应根据焊接结构形状、尺寸、受力情况、强度要求、焊件厚度、焊接方法及坡口加工难易程度等因素综合决定。

1. 手工电弧焊焊接接头形式的选择

手工电弧焊焊接接头的基本形式有 4 种：对接接头（图 9-19a）、角接接头（图 9-19b）、T 形接头（图 9-19c）和搭接接头（图 9-19d）。

\begin{figure}[H]
    \centering
    % Placeholder for Figure 9-19
    \includegraphics[page=36, viewport={251.5 497.8 512.5 568.3}, clip]{12-07572-001FigRevised.pdf}
    \caption{手工电弧焊焊接接头的基本形式}
\end{figure}

对接接头受力较均匀，焊接质量易于保证，应用最广，应优先选用。角接接头和 T 形接头受力情况比对接接头复杂，但接头为直角或具有一定角度时必须采用这两种接头形式。搭接接头受力时，焊缝处易产生应力集中和附加弯矩，一般应避免选用。但因其不需开坡口，焊前装配方便，对受力不大的平面连接也可选用。

2. 手工电弧焊坡口形式的选择

除搭接接头外，其余接头在焊件较厚时均需开坡口。开坡口的目的是保证焊透。坡口的基本形式有 I 形（图 9-20a）、V 形（图 9-20b）、U 形（图 9-20c）和 X 形（图 9-20d）等，常采用砂轮打磨、切削加工、电弧加工和火焰切割等方法制成。

\begin{figure}[H]
    \centering
    % Placeholder for Figure 9-20
    \includegraphics[page=36, viewport={126.0 436.8 417.0 468.8}, clip]{12-07572-001FigRevised.pdf}
    \caption{手工电弧焊坡口的基本形式}
\end{figure}

I 形坡口主要用于厚度为 $1 \sim 6 \mathrm{~mm}$ 钢板的焊接；V 形坡口主要用于厚度为 $3 \sim 26 \mathrm{~mm}$ 钢板的焊接；U 形坡口主要用于厚度为 $20 \sim 60 \mathrm{~mm}$ 钢板的焊接；X 形坡口主要用于厚度为 $12 \sim 60 \mathrm{~mm}$ 钢板的焊接，需双面施焊。

3. 其他焊接方法接头形式

埋弧自动焊的接头形式与手工电弧焊基本相同。因其电流大，熔深大，故焊件厚度小于 $12 \mathrm{~mm}$ 时不可开坡口单面焊接；焊件厚度小于 $24 \mathrm{~mm}$ 时不开坡口双面焊接。电渣焊可采用对接、角接、T 形接头焊接，不需开坡口。点焊、缝焊多采用搭接接头。

焊接时应尽量避免厚薄相差很大的焊件焊接，以便获得优质焊接接头。必须采用时，在较厚焊件上应加工出过渡形式（图 9-21）。

\begin{figure}[H]
    \centering
    % Placeholder for Figure 9-21
    \includegraphics[page=36, viewport={112.0 316.8 426.5 399.3}, clip]{12-07572-001FigRevised.pdf}
    \caption{焊接时不同厚度焊件的过渡形式}
\end{figure}

\subsection{焊缝的布置}

焊缝布置一般应遵循以下原则。

（1）焊缝布置应便于焊接操作

焊缝设置必须留有足够的操作空间以满足焊接时运条的需要。手工电弧焊时，至少应使焊条能伸到待焊部位（图 9-22）。点焊与缝焊时，电极应能伸到待焊部位（图 9-23）。埋弧自动焊时，则应考虑施焊时接头处能存放焊剂（图 9-24）。

\begin{figure}[H]
    \centering
    % Placeholder for Figure 9-22
    \includegraphics[page=36, viewport={33.0 131.3 316.5 286.8}, clip]{12-07572-001FigRevised.pdf}
    \caption{手工电弧焊焊缝设计}
\end{figure}

\begin{figure}[H]
    \centering
    % Placeholder for Figure 9-23
    \includegraphics[page=36, viewport={335.5 131.3 501.0 259.8}, clip]{12-07572-001FigRevised.pdf}
    \caption{点焊与缝焊的焊缝设计}
\end{figure}

\begin{figure}[H]
    \centering
    % Placeholder for Figure 9-24
    \includegraphics[page=37, viewport={19.5 563.8 105.5 748.3}, clip]{12-07572-001FigRevised.pdf}
    \caption{埋弧自动焊焊缝的设计}
\end{figure}

（2）焊缝布置应有利于减少焊接应力与变形

设计焊接结构时，应尽量选用尺寸规格较大的板材、型材和管材，形状复杂的可采用冲压件和铸钢件，以减少焊缝数量，简化焊接工艺和提高结构的强度和刚度。同时，焊缝应尽可能对称布置（图 9-25），以减小变形。图 9-26 所示的箱形结构，可以由 4 块钢板拼焊而成（图 9-26a），也可以由两块槽钢（图 9-26b）或经弯曲的钢板（图 9-26c）拼焊而成，此时只需两条焊缝，且对称布置。

\begin{figure}[H]
    \centering
    % Placeholder for Figure 9-25
    \includegraphics[page=37, viewport={130.5 559.3 378.5 720.8}, clip]{12-07572-001FigRevised.pdf}
    \caption{焊缝的对称设计}
\end{figure}

\begin{figure}[H]
    \centering
    % Placeholder for Figure 9-26
    \includegraphics[page=37, viewport={392.5 565.3 521.0 714.8}, clip]{12-07572-001FigRevised.pdf}
    \caption{焊缝的合理选材与减少焊缝}
\end{figure}

焊缝的布置应避免密集、交叉或重叠（图 9-27）。因焊缝交叉或过分集中会造成接头部位产生过热，增大热影响区，使接头组织恶化，性能严重下降。两条焊缝间距一般要求大于 3 倍焊件厚度且不小于 $100 \mathrm{~mm}$。

\begin{figure}[H]
    \centering
    % Placeholder for Figure 9-27
    \includegraphics[page=37, viewport={28.5 379.8 341.0 532.8}, clip]{12-07572-001FigRevised.pdf}
    \caption{焊缝设计应避免密集、交叉或重叠}
\end{figure}

（3）焊缝布置应避开最大应力区和应力集中部位

焊接的热影响区很可能是焊接结构的薄弱环节。因此，焊缝应避开焊接结构上应力最大的部位（图 9-28）；另外，拐角处等应力集中部位也是焊接结构的薄弱部位，不应设计焊缝，如压力容器一般不用无折边封头，而应采用碟形封头和球形封头等（图 9-29）。

\begin{figure}[H]
    \centering
    % Placeholder for Figure 9-28
    \includegraphics[page=37, viewport={364.0 378.8 509.0 507.3}, clip]{12-07572-001FigRevised.pdf}
    \caption{焊缝应避开最大应力区}
\end{figure}

\begin{figure}[H]
    \centering
    % Placeholder for Figure 9-29
    \includegraphics[page=37, viewport={25.5 183.8 227.5 252.3}, clip]{12-07572-001FigRevised.pdf}
    \caption{焊缝应避开应力集中部位}
\end{figure}

（4）焊缝布置应避开或远离机械加工面

焊接时局部加热会引起焊件变形，设计焊缝时必须考虑。焊接结构上的加工面有两种不同情况：对焊接结构的位置精度要求较高时，一般应在焊后进行精加工，以免焊接变形影响加工精度；对焊接结构的位置精度要求不高时，可先进行机械加工，但焊缝位置与加工面要保持一定距离，以保证原有的加工面精度（图 9-30）。

\begin{figure}[H]
    \centering
    % Placeholder for Figure 9-30
    \includegraphics[page=37, viewport={251.0 187.3 511.0 343.8}, clip]{12-07572-001FigRevised.pdf}
    \caption{避开或远离机械加工面}
\end{figure}

为减少焊接工作量和焊接缺陷与变形，焊接时应尽量在平焊位置施焊，避免仰焊，减少横焊和立焊。尽量使全部焊件（至少主要焊件）能在焊前一次装配点固，简化焊件装配过程。

（5）便于焊接和检验

设计封闭容器时，要留工艺孔，如入孔、检验孔和通气孔，焊后再用其他方法封堵。

\section*{复习思考题}

 
9-1 什么是焊接电弧？它由哪几个区组成？各区产生的热量和温度如何？什么是正接和反接？

9-2 焊接与其他连接方法的本质区别是什么？

9-3 电焊条由哪几部分组成？各部分作用如何？其中为什么含有锰铁？在其他电弧焊工艺中，有无类似材料和作用？

9-4 什么是焊接热影响区？熔焊时低碳钢的焊接热影响区一般分为哪几个区？它们对焊接接头质量有何影响？可采用什么措施来减小或消除焊接热影响区？

9-5 什么是酸性焊条和碱性焊条？各有何特点？焊后焊接接头质量如何？各在何种场合使用？

9-6 焊条选用的基本原则有哪些？

9-7 产生焊接应力和变形的原因是什么？焊接应力是否一定要消除？如要消除，可采取哪些措施？

9-8 焊接变形的基本形式有哪些？预防和减小焊接变形常采用哪些措施？

9-9 简述常用的焊接变形矫正方法。

9-10 埋弧自动焊与手工电弧焊相比有何优点？

9-11 氩弧焊可分为哪几种？各有何特点？适合于焊接哪些材料？

9-12 电阻焊的基本形式有哪几种？各有何特点？

9-13 电渣焊有何特点？适用于什么场合？

9-14 钎焊和熔化焊的实质有何区别？钎焊分哪几类？各有何特点和应用？

9-15 下列制品应选用什么焊接方法？

自行车车架（钢）；自行车车圈（钢）；铝合金自行车车架；钢窗；家用液化气罐罐体；电子线路板；锅炉壳体；有缝钢管；不锈钢储罐；硬质合金刀片与刀杆。

9-16 何谓金属的焊接性？何谓碳当量？

9-17 低碳钢、中碳钢和高碳钢的焊接工艺有何不同？

9-18 低合金结构钢焊接时有哪些特点？为保证焊接接头质量应采取哪些措施？

9-19 铸铁焊补方法如何分类？各有何特点？常用的铸铁焊条有哪几种？

9-20 铜合金和铝合金常用的焊接方法有哪些？

9-21 陶瓷材料、高分子材料和复合材料的焊接性怎样？它们常用什么方法来焊接？

9-22 手工电弧焊接头的基本形式有哪几种？有何特点？焊接较厚板材时为什么要开坡口？

9-23 焊缝布置应该遵循什么原则？

9-24 指出图 9-31 所示焊接结构设计或焊接顺序的不合理之处，说明理由并加以修改。

\begin{figure}[H]
    \centering
    % Placeholder for Figure 9-31
    \includegraphics[page=38, viewport={27.0 544.8 286.5 744.3}, clip]{12-07572-001FigRevised.pdf}
    \caption{题 9-24 图}
\end{figure}

9-25 在图 9-32 中，图 9-32a 所示为不等厚度板的焊接，图 9-32b 所示为在有孔板上焊接圆管件，请指出图示结构中设计不合理之处，说明理由并改正之。

\begin{figure}[H]
    \centering
    % Placeholder for Figure 9-32
    \includegraphics[page=38, viewport={313.0 538.8 516.0 634.3}, clip]{12-07572-001FigRevised.pdf}
    \caption{题 9-25 图}
\end{figure}

9-26 在不改变图 9-33 中 A、B、C、D、E 和 F 六小块板料尺寸和不改变焊后面积 $3a \times 3a$ 的前提下，试将图中交叉焊缝连接改为 T 形连接，并标注其焊缝的焊接顺序。

\begin{figure}[H]
    \centering
    % Placeholder for Figure 9-33
    \includegraphics[page=38, viewport={5.0 328.8 131.5 457.8}, clip]{12-07572-001FigRevised.pdf}
    \caption{题 9-26 图}
\end{figure}

9-27 图 9-34 所示的焊接结构设计与焊缝位置是否合理？如不合理请加以改进。

\begin{figure}[H]
    \centering
    % Placeholder for Figure 9-34
    \includegraphics[page=38, viewport={143.0 331.3 530.5 511.3}, clip]{12-07572-001FigRevised.pdf}
    \caption{题 9-27}
\end{figure}


\chapter{非金属及复合材料成形方法}

随着材料科学技术及现代工业的发展，非金属及复合材料的应用范围越来越广。与金属材料一样，塑料、陶瓷、复合材料等都必须经过加工，才能得到具有一定形状、尺寸和使用性能的产品。本章简要介绍塑料、陶瓷及复合材料的常用成形方法。

\section{工程塑料件的成形}

工程塑料件的成形是将粉状、粒状、溶液或分散体等各种物态的聚合物料转变为所需形状制品的工艺过程。塑料的成形方法很多，在较低温度（400 $^\circ$C以下）条件下，一般采用挤出、注射、压制、浇注、吹塑、压延等成形方法。

\subsection{挤出成形}

挤出成形又称为挤塑成形。首先将粒状或粉状的热塑性塑料加入挤出机料斗中，在旋转着的挤出机螺杆的作用下，塑料沿螺旋槽向前输送，料筒内的塑料由于料筒的传热以及塑料与料筒及螺杆之间的剪切和摩擦热，逐渐熔融呈黏流态，然后在挤压系统作用下，塑料熔体通过具有一定形状的挤出模具（机头）口而连续成形为所需断面形状的型材。最通用的单螺杆式挤出成形如图10-1所示。

\begin{figure}[H]
    \centering
    \includegraphics[page=38, viewport={22.0 121.3 298.5 232.3}, clip]{12-07572-001FigRevised.pdf}
    \caption{单螺杆式挤出成形示意图}
\end{figure}

挤出成形塑料件组织均匀致密、尺寸稳定准确。挤出成形模具结构简单、制造维修方便，同时还能连续成形，生产效率高、成本低；除氟塑料外，所有热塑性塑料及部分热固性塑料均可采用挤出成形。该方法还可用于塑料的着色、造粒和共混改性等。

挤出成形加工的产品有各种管材、棒材、板材、电线电缆外皮和中空容器等。

\subsection{注射成形}

注射成形是将颗粒或粉状塑料置于注射成形机料斗内加热至流动状态，由注射柱塞（或螺旋杆）施加高压和较快速度，使已熔化的物料通过狭小的喷嘴注入闭合模具的模腔中，经冷却脱模，即可得到所需形状的塑料制品。其工艺过程包括成形前的准备、注射过程、制品的后处理等。注射成形主要用于热塑性塑料和流动性较好的热固性塑料，可以成形几何形状非常复杂、尺寸要求高及带各种嵌件的塑料制品，如无线电配件、电气零件、日常生活用品等。注射成形生产过程如图10-2所示。

\begin{figure}[H]
    \centering
    \includegraphics[page=38, viewport={315.0 126.3 510.5 298.3}, clip]{12-07572-001FigRevised.pdf}
    \caption{注射成形生产过程}
\end{figure}

\subsection{压延成形}

压延成形是生产高分子薄膜材料和片材的成形方法。压延过程是利用一对或数对相对旋转的加热滚筒，使加热塑化的热塑性塑料在滚筒间隙中被压延而连续形成一定厚度和宽度的薄型材料。所用设备为压延机。压延过程如图10-3所示。

压延成形可以生产大型塑料制品，特别是平板类制品，制品质量均匀、内应力小、尺寸稳定、耐热、强度较高。但此法难以生产形状复杂及带有金属嵌件的塑料制品。

\begin{figure}[H]
    \centering
    \includegraphics[page=39, viewport={81.5 608.3 245.5 704.3}, clip]{12-07572-001FigRevised.pdf}
    \caption{压延过程示意图}
\end{figure}

\section{陶瓷件成形}

陶瓷具有高强度、高硬度、耐磨损、耐高温、耐氧化等优点，陶瓷制品的生产过程主要包括配料、成形、烧结三个阶段。陶瓷经成形、烧结后还可根据需要进行磨削加工和抛光，甚至切削加工。本节着重介绍陶瓷的几种常用成形方法。

\subsection{干压成形}

干压成形类似金属模锻，是将粉料装入钢模内，通过冲头对粉末施加压力，压制成具有一定形状和尺寸的压坯的成形方法（图10-4）。该方法一般适用于形状简单、尺寸较小的制品。

\begin{figure}[H]
    \centering
    \includegraphics[page=39, viewport={263.5 612.3 462.0 745.3}, clip]{12-07572-001FigRevised.pdf}
    \caption{干压成形示意图}
\end{figure}

干压成形工艺简单、操作方便、周期短、效率高、便于自动化生产。此外，干压成形还具有坯体密度大、尺寸精确、收缩小、强度高等特点。但干压成形对大型坯体生产有困难，模具磨损大、加工复杂、成本高。

\subsection{等静压成形}

等静压成形是利用液体或橡胶等在各个方向传递压力相等的原理对坯体进行压制的。等静压成形可分为湿式等静压成形和干式等静压成形两种。

湿式等静压成形是将预压好的坯体（粉体）包封在具有弹性的塑料模、橡皮模或软模中，再置入压力容器之内。高压泵通过进液口将传压液体打入压力容器内，橡皮模内的坯体在各个方向受到同等大小的压力。传压液体可以是水，也可以是油等介质。压力可以在一定范围内调整。对于要求高的工件，进行胶套密封时还要做真空处理。湿式等静压成形如图10-5所示。

\begin{figure}[H]
    \centering
    \includegraphics[page=39, viewport={49.5 440.8 234.5 580.3}, clip]{12-07572-001FigRevised.pdf}
    \caption{湿式等静压成形示意图}
\end{figure}

等静压成形具有压力容易调整、坯体均匀致密、烧结收缩小、不易变形开裂等优点，但设备比较复杂，操作繁琐，生产效率低，目前仍只限于生产具有较高要求的电子元件及其他高性能材料。

\subsection{注浆成形}

注浆成形是将陶瓷颗粒悬浮于液体中成为料浆，再把料浆注入多孔模具中，通过模具的气孔把料浆中的液体吸出，而在模具内形成坯体的工艺方式。其工艺过程包括料浆准备、模具制备和料浆浇注（图10-6）三个阶段。注浆成形方法可制造形状复杂、大型薄壁制品。另外，金属铸造工艺的型芯、离心铸造、真空铸造、压力铸造等工艺方法被引入注浆成形，从而形成了离心注浆、真空注浆、压力注浆等方法。离心注浆适用于制造大型环状制品，且坯体壁厚均匀；真空注浆可有效地去除料浆中的气体；压力注浆可提高坯体的致密度，减少坯体中的残留水分，缩短成形时间，减少制品缺陷，是一种较先进的成形工艺。

\begin{figure}[H]
    \centering
    \includegraphics[page=39, viewport={261.0 442.3 501.5 545.8}, clip]{12-07572-001FigRevised.pdf}
    \caption{料浆浇注过程示意图}
\end{figure}

\subsection{热压成形}

热压成形是利用蜡类材料热熔冷固的特点，把粉料与熔化的蜡料黏结剂迅速搅和成具有流动性的料浆，在热压铸机中用压缩空气把热熔料浆注入金属模，冷却凝固后成形的工艺方法。这种成形方法操作简单，模具损耗小，可成形复杂制品，但坯体密度较低，生产周期长。

\subsection{注射成形}

注射成形是将粉料与有机黏结剂混合后，加热混炼，制成粒状粉料，用注射成形机在130$\sim$300 $^\circ$C温度下注射入金属模具中，冷却后形成坯体的工艺方法。坯体取出并经脱脂后，按常规工艺烧结。这种工艺成形简单，成本低，坯体密度均匀，适用于复杂零件的自动化大规模生产。

\section{复合材料成形}

复合材料的成形方法较多，本节在介绍通用制备方法的基础上，重点介绍树脂基复合材料、金属基复合材料、陶瓷基复合材料和C/C复合材料的成形方法。

\subsection{制备复合材料的通用方法}

1. 颗粒、晶须、短纤维增强复合材料

制备方法通常包括下列3个步骤。

（1）混合

基体材料熔化（溶化）为液态，采用搅拌方法均匀混入增强材料；或制成粉末，采用滚筒或球磨等方法混入增强材料，并均匀化。

（2）制坯

采用铸造、液态模锻、喷射、粉末热压等方法使复合材料凝固或固化，制备出复合材料坯体或零件。

（3）成形

根据需要，通过挤压、轧制、锻造、机械加工等二次加工，制备出性能、形状均满足要求的零件。

2. 纤维增强体增强复合材料

制备方法通常包括下列两个基本步骤。

（1）增强体预成形

按设计要求将增强纤维（或纤丝）排列成特定形状或模式。对长纤维（或纤丝），采用缠绕、织物铺层、三维编织等方法成形；对晶须或较短的纤维，采用磁力、静电、振荡、压延或悬浮法进行预处理，再用挤压等方法成形。

（2）复合

将基体材料与增强体复合，通常采用的复合方法有粉末冶金法、液态浸透法、化学气相沉积法等。

\subsection{金属基复合材料成形方法}

制备金属基复合（metal matrix composite, MMC）材料，关键在于获得基体金属与增强材料之间良好的浸润和合适的界面结合。下面介绍几种常用方法。

1. 液态金属浸润法

（1）常压铸造法

首先将经过预处理的纤维制成整体或局部形状的零件预成形坯，然后将预成形坯预热后放入浇注模，浇入液态金属，靠重力使金属渗入纤维预成形坯并凝固。此法适合较大规模生产。但复合材料制品易存在宏观或微观缺陷。

（2）液态金属搅拌法

首先将基体金属放入坩埚中熔化，插入旋转叶片，搅拌金属液，并逐步加入弥散增强材料，直至在熔体中均匀分布为止，然后进行脱气处理，注入模具中凝固成形。该方法设备较为简单，生产成本低，主要用于陶瓷颗粒增强金属基复合材料的制造。

（3）真空加压铸造法

首先在真空（或惰性气体）密闭容器中加热纤维预成形坯和熔化金属，然后将铸模的引流管插入熔融金属中，并通入惰性气体对金属液面施以一定压力，强制液态金属渗入预成形坯，冷却凝固后制成复合材料或制品（图10-7）。该方法适用于生产小型零件，但生产率较低。

\begin{figure}[H]
    \centering
    \includegraphics[page=39, viewport={53.5 203.8 212.0 324.3}, clip]{12-07572-001FigRevised.pdf}
    \caption{真空加压铸造装置示意图}
\end{figure}

（4）挤压铸造法

首先将增强材料放入配有黏结剂和纤维表面改性溶质的溶液中，充分搅拌；然后压滤、干燥（图10-8a$\sim$d），或通过直接抽吸、脱模和干燥（图10-8e$\sim$g），烧结成具有一定强度的预成形坯；最后将预热后的预成形坯放入固定在液压机上经预热的模具中，注入液态金属，加压，使金属渗透预成形坯，并在高压下凝固成形为复合材料制品（图10-9）。该成形方法可生产材质优良、加工余量小的制品，成本低，生产率高。

\begin{figure}[H]
    \centering
    \includegraphics[page=39, viewport={233.5 202.8 488.0 407.8}, clip]{12-07572-001FigRevised.pdf}
    \caption{短纤维预成形坯的两种制造过程}
\end{figure}

2. 扩散黏结法

扩散黏结法是在较长时间、较高温度和压力下，使固态金属基体与增强材料的接触界面通过原子间相互扩散黏结成形的工艺方法。制造时，首先把增强纤维用不同的方法，如等离子喷涂法、液态金属浸渍法、化学涂覆法等制成预成形坯，然后将预成形坯处理、清洗，按一定形状、尺寸和排列形式叠层封装，加热压制。压制过程可以在真空、惰性气体或大气环境中进行。常用的压制方法有热压法、热等静压法和热轧法。

3. 粉末冶金法

粉末冶金法是将金属粉末与陶瓷颗粒、晶须或短纤维，经均匀混合后放入模具中，而后进行高温、高压成形的工艺方法。该方法可直接制成零件，也可制坯进行二次成形。制得的材料致密度高，增强材料分布均匀。

4. 喷雾共沉积法

喷雾共沉积法是用于生产陶瓷颗粒增强金属基复合材料的一种新工艺（图10-10）。熔融金属从炉子底部的浇注孔流出，经雾化器被高速惰性气体流雾化，同时由气体携带陶瓷颗粒加入雾化流中使其混合、沉降，在金属滴尚未完全凝固前喷射在基板或特定模具上，并凝固成固态（复合材料）。这种方法成形的材料致密度高，陶瓷颗粒分布均匀，生产率高，还可直接生产不同规格的空心管、板、锻坯和挤压锭等。

\begin{figure}[H]
    \centering
    \includegraphics[page=40, viewport={87.5 566.3 265.0 730.8}, clip]{12-07572-001FigRevised.pdf}
    \caption{挤压铸造复合材料示意图}
\end{figure}

\begin{figure}[H]
    \centering
    \includegraphics[page=40, viewport={284.5 565.8 457.5 745.3}, clip]{12-07572-001FigRevised.pdf}
    \caption{喷雾共沉积法示意图}
\end{figure}

\subsection{树脂基复合材料成形方法}

树脂基复合（resin matrix composite, RMC）材料成形方法有手糊成形、热压罐成形、对模模压成形、纤维缠绕成形、拉挤成形、喷射成形以及注射成形等。下面介绍几种常用方法。

1. 热压罐成形

首先将预浸材料按一定排列顺序置于涂有分型剂（脱模剂）的模具上，铺放分离布和带孔的脱模薄膜，在脱模薄膜的上面铺放吸胶透气毡，再包覆耐高温的真空袋，并用密封胶条密封周边（图10-11）。然后，连续从真空袋内抽出空气并加热，使预浸材料的层间达到一定程度的真空度。当达到要求的温度后，向热压罐内充以压缩空气，给制品加压。由于无法直接观察到基体树脂的流变和固化行为，只能通过测定树脂在固化过程中的黏度、介电常数或反应热的变化，来确定加温和加压程序的实施。因而，该方法也被认为是一种带有“技巧性”的方法。

\begin{figure}[H]
    \centering
    \includegraphics[page=40, viewport={49.5 369.8 278.0 533.8}, clip]{12-07572-001FigRevised.pdf}
    \caption{热压罐成形示意图}
\end{figure}

2. 对模模压成形

对模模压成形是将模压料约束在两个模具型面之间形成制品形状，并加压使之固化的成形方法。根据所使用模压料形式和状态不同，对模模压成形可分为增强模压料模压、毡与预成形坯料模压、冷模压、树脂注入模压、泡沫蓄积模压等方法。在工业中占主导地位的是增强模压料模压和树脂注入（传输）模压。图10-12所示为厚片模压料制造流程。该方法复现性好、成形速度快，成形的制品质量高、尺寸精度高，适用大批量生产。

3. 热成形工艺

热成形工艺也称为预热坯料成形，与热固性树脂基复合材料的对模模压成形类似，是一种快速、大量成形热塑性复合材料制品的成形方法。复合材料制品的热成形工艺与纯塑料制品的热成形工艺不同，预浸料在模具内不能伸长，也不能变薄。模具闭合之前，预浸料要从夹持框架上松开，放置在下半模具上。模具闭合时，预浸料铺层边缘将向模具中滑移，并贴覆到模具型面上，预浸料厚度保持不变（图10-13）。

\begin{figure}[H]
    \centering
    \includegraphics[page=40, viewport={295.0 367.8 494.5 526.3}, clip]{12-07572-001FigRevised.pdf}
    \caption{厚片模压料制造流程}
\end{figure}

\begin{figure}[H]
    \centering
    \includegraphics[page=40, viewport={22.0 235.3 226.5 334.8}, clip]{12-07572-001FigRevised.pdf}
    \caption{热成形示意图}
\end{figure}

4. 缠绕成形

缠绕成形是制造具有回转体形状的复合材料制品的基本成形方法。首先将浸渍树脂的纤维按照要求的方向有规律、均匀地布满芯模表面，再送入固化炉固化，脱去芯模即可得到所需制品（图10-14）。该方法的基本设备是缠绕机、固化炉和芯模。

\begin{figure}[H]
    \centering
    \includegraphics[page=40, viewport={252.0 233.8 521.0 331.8}, clip]{12-07572-001FigRevised.pdf}
    \caption{缠绕成形示意图}
\end{figure}

缠绕成形可按设计要求确定缠绕方向、层数和数量，获得等强度结构。该方法机械化、自动化程度高，产品质量好。

5. 拉挤成形

拉挤成形是一种可连续制造恒定截面复合材料型材的工艺方法，与铝材的挤压成形或热塑

性塑料的挤出成形相似,可制造实心、空心以及各种复杂截面的制品,并且可以设计型材的性能,以满足各种工程和结构要求。如可在连续拉挤过程中,埋入金属件、木材或泡沫等。

拉挤成形有6个基本工艺环节(图10-15)。首先,将增强纤维送入树脂槽浸渍树脂;在牵引机构的牵引下,浸渍树脂的增强纤维在预成形模中按照产品形状预成形;随后,预成形坯进入固化模中精成形;热固性树脂基体在热的作用下固化成所需截面的型材;固化后的型材在牵引机构的牵引下,连续从热模具中脱出,在空气或水中冷却;最后,进入自动切割装置切成所需长度。

\begin{figure}[H]
\centering
   \includegraphics[page=40, viewport={109.5 91.8 425.5 204.3}, clip]{12-07572-001FigRevised.pdf}
\caption{拉挤成形工艺流程示意图}
\label{fig:10-15}
\end{figure}

\section{陶瓷基复合材料成形方法}

1. 浆体浸渗工艺

该方法是早期生产陶瓷基复合材料(ceramic matrix composite, CMC)最常用的一种方法。其工艺流程如图10-16所示。玻璃陶瓷基复合材料的浆体浸渍工艺如图10-17所示。纤维束(或纤维预制体)通常经过至少含有3种组分的泥浆混合物(基体粉末、水或乙醇、有机黏结剂)进行浸渍,再压制切断成单层薄片,按一定方式排列成层板,然后放入加热炉中烧去黏结剂,最后加压使之固化。

\begin{figure}[H]
\centering
    \includegraphics[page=41, viewport={36.5 570.3 220.5 750.3}, clip]{12-07572-001FigRevised.pdf}
    \\
\caption{浆体压制烧结工艺流程}
\label{fig:10-16}
\end{figure}

\begin{figure}[H]
\centering
   \includegraphics[page=41, viewport={243.5 570.8 503.5 723.3}, clip]{12-07572-001FigRevised.pdf}
\caption{玻璃陶瓷基复合材料的浆体浸渗工艺}
\label{fig:10-17}
\end{figure}

2. 气-液反应工艺

将熔融金属直接氧化制备陶瓷基复合材料的工艺方法,其商业名称为Lanxide工艺。该方法具有工艺简单、成本低廉、反应温度低、反应速度快等优点,且制品的形状及尺寸几乎不受限制,其性能还可由工艺调控。尽管其致命弱点是存在残余金属,使高温强度显著下降,但常温性能优越,所以成为陶瓷基复合材料制备中具有吸引力的方法之一。

3. 化学气相渗透法

化学气相渗透(chemical vapor infiltration, CVI)法的工艺过程如下:首先,将具有特定形状的纤维预制体置于沉积炉中,通入的气态前驱体通过扩散、对流等方式进入预制体内部,在一定温度下由于热激活而发生复杂的化学反应,生成固态的陶瓷类物质并以涂层的形式沉积于纤维丝表面;然后,随着沉积的继续进行,纤维表面的涂层越来越厚,纤维间的空隙越来越小,最终各涂层相互重叠,成为材料内的连续相,即陶瓷基体。CVI工艺简图如图10-18所示。

\begin{figure}[H]
	\centering
	\includegraphics[page=41, viewport={15.5 308.8 207.0 488.3}, clip]{12-07572-001FigRevised.pdf}
	\caption{CVI工艺简图}
	\label{fig:10-18}
\end{figure}

CVI工艺具有以下优点。

1)在无压和相对低温条件下进行,纤维类增强物的损伤较小,可制备出高性能(特别是高断裂韧性)的陶瓷基复合材料。

2)通过改变气态前驱体的种类、质量分数、沉积顺序、沉积工艺,可方便地对陶瓷基复合材料的界面、基体的组成与微观结构进行设计。

3)由于不需要加入烧结助剂,所得到的陶瓷基体在纯度和组成结构上优于常规方法制备的复合材料。

4)可成形形状复杂、纤维体积分数较高的陶瓷基复合材料,但成形周期长,成本高。

根据控制气体输送模式和反应温度不同,CVI法主要有等温CVI(isothermal chemical vapor infiltration, ICVI)法、等温强制流动CVI法、热梯度CVI法、强制流动热梯度CVI(force flow-thermal gradient chemical vapor infiltration, FCVI)法和脉冲CVI法。

\section{C/C复合材料成形方法}

C/C复合材料的增强剂为碳纤维及其织物,基体为碳或石墨,其工艺主要包括预成形体的成形和致密化两个过程。工艺流程如图10-19所示。

1. 预成形体的成形工艺

预成形体是指按产品形状和性能要求先把碳纤维成形为所需的结构,以便进一步进行C/C致密化工艺。短纤维增强的预成形体常采用压滤法、浇注法、喷涂法、热压法;连续纤维增强的预成形体可采用传统成形方法,如预浸布、层压、铺层、缠绕等方法,也可采用多向编织技术。

2. C/C的致密化工艺

C/C的致密化工艺过程就是基体碳形成的过程,实质是用高质量的碳填满碳纤维周围的空隙以获得结构、性能优良的C/C复合材料。最常用的致密化工艺包括树脂、沥青的液相浸渍工艺和碳氢化合物气体的化学气相渗透工艺两种基本工艺。

树脂浸渍工艺的典型流程是:先将预成形体置于浸渍罐中,在真空状态下用树脂浸没预成形体;再充气加压使树脂浸透预成形体;最后,将浸透树脂的预成形体放入固化罐内进行加压固化,随后在碳化炉中的保护气氛下进行碳化。

\begin{figure}[H]
\centering
    \includegraphics[page=41, viewport={225.0 312.3 518.5 550.8}, clip]{12-07572-001FigRevised.pdf}
    \\
\caption{C/C复合材料制造基本工艺流程}
\label{fig:10-19}
\end{figure}

\section*{复习思考题}

 
10-1 塑料的成形方法主要有哪些?其适用范围及其工艺过程如何?

10-2 橡胶的成形方法主要有哪些?其适用范围如何?

10-3 陶瓷的成形方法主要有哪些?工艺过程有何特点?

10-4 简述制备复合材料的通用方法。

10-5 陶瓷基复合材料主要有哪些成形方法?

10-6 金属基复合材料主要有哪些成形方法?

10-7 简述用挤压铸造法制备金属基复合材料的基本过程与特点。

10-8 树脂基复合材料有哪些类型?分别用何方法成形?

10-9 化学气相渗透法制备陶瓷基复合材料的基本流程及特点是什么?

10-10 简述C/C复合材料的制备工艺过程。


\chapter{毛坯成形方法选择及质量控制}

\section{毛坯成形方法选择}

为节省材料、提高产品生产效率和满足零件使用性能等方面的要求,机械加工对象常常是与零件形状已经比较接近的各种毛坯。毛坯主要有通过铸造方法成形的铸件和通过锻造方法成形的锻件,以及轧制的各种型材。铸造、锻造和轧制等毛坯制造方法有其各自的特点和适用性,应根据零件使用性、经济性和材料的成形性等方面的要求,做出合理地选择。

\subsection{毛坯分类}

常用机械零件毛坯依其形状大致分为轴(杆)类、机架箱体类和盘套类三类。

1. 轴(杆)类零件

轴(杆)类零件(图11-1)是重要的受力件和传动件,轴(杆)类零件的材料选择主要考虑强度,同时兼顾冲击韧度和表面耐磨性。轴(杆)类零件通常以锻件作毛坯,常用材料为中碳钢或中碳合金调质钢;等截面轴也可采用轧制型材直接加工,某些较复杂的轴(杆)类零件可采用锻-焊组合结构;具有部分空心结构的轴,可先分段下料锻造,然后拼焊;带有法兰端的轴,可将法兰和轴分别锻造并拼焊。轴(杆)类零件通常需经调质,以获得良好的综合性能,调质后还需进行高频感应加热淬火或渗氮处理,以提高表面硬度,增强耐磨性。

2. 机架箱体类零件

各类机床的机身、底座、支架、轴承座、齿轮箱及阀体、泵体等均属于机架箱体类零件(图11-2)。这类零件的特点是形状不规则,结构比较复杂,质量从几千克至数十吨,工作条件也相差很大。机身、底座、工作台和导轨面等基础零件,主要起支承和连接机床各部件的作用,以承受压应力、弯曲应力和摩擦力为主,要求具有较好的刚性、减振性和耐磨性,因此多数情况下选用灰铸铁进行铸造成形。对受力复杂或承受较大冲击载荷的和架箱体类零件,可选用中碳铸钢件或合金铸钢件。

个别机架箱体类零件还可采用铸钢-焊接结构。在下列情况下通常采用铸-焊结构毛坯。


1)重型机械的机身等大型零件,采用铸造时重量难以控制。

2)由于技术原因铸件的形状或重量受到限制。

3)须在短时间内制造出少量毛坯。


设计合理的焊接结构比铸件可靠性高,重量轻,且能满足刚性要求,但减振性差,易变形和产生应力集中。

\begin{figure}[H]
	\centering
	\includegraphics[page=41, viewport={8.5 142.3 156.0 278.3}, clip]{12-07572-001FigRevised.pdf}
	\caption{轴(杆)类零件}
	\label{fig:11-1}
\end{figure}

\begin{figure}[H]
\centering
    \includegraphics[page=41, viewport={155.5 143.3 358.0 266.3}, clip]{12-07572-001FigRevised.pdf}
\caption{机架箱体类零件}
\label{fig:11-2}
\end{figure}

3. 盘套类零件

齿轮、飞轮、带轮、轴承环、卡盘、法兰盘、联轴器等均可归类于盘套类零件(图11-3)。其大小、形状、工作条件各不相同,所选用的材料及毛坯也不同,应根据具体情况具体分析,灵活选用。例如,汽车发动机上的凸轮、齿轮,由于在低应力、低冲击载荷条件下工作,材料可选HT200,采用铸造成形。而对于手轮、带轮、端盖和卡盘等受力不大或结构复杂的零件,通常采用铸铁件。毛坯的生产方法不同,在结构、性能、质量、体积和生产成本上也会存在很大差异。表11-1给出了各种毛坯的基本特点,选择毛坯时可根据这些特点进行综合分析比较,选择合理方案。

\begin{figure}[H]
\centering
\includegraphics[page=41, viewport={369.5 139.8 535.0 227.3}, clip]{12-07572-001FigRevised.pdf}
\caption{盘套类零件}
\label{fig:11-3}
\end{figure}

\begin{table}[H]
\centering
\caption{各种毛坯的基本特点}
\footnotesize\setlength{\tabcolsep}{2pt}
\begin{tabular}{|p{0.8cm}|p{1.5cm}|p{1.8cm}|p{2.2cm}|p{1.2cm}|p{1.5cm}|p{1.2cm}|p{1.5cm}|p{1.2cm}|p{0.8cm}|p{1.2cm}|}
\hline
\multirow{2}{*}{种类} & \multirow{2}{*}{制造方法} & \multirow{2}{*}{形状复杂程度} & \multirow{2}{*}{内部组织特点} & \multirow{2}{*}{力学性能} & \multicolumn{2}{c|}{壁厚尺寸或质量} & \multirow{2}{*}{常用合金} & \multirow{2}{*}{表面质量} & \multirow{2}{*}{毛坯精度} & \multirow{2}{*}{生产批量} \\
\cline{6-7}
 & & & & & 最小 & 最大 & & & & \\
\hline
\multirow{3}{*}{铸件} & 砂型铸造 & 复杂,特别适宜于内腔复杂的铸件 & 组织疏松,易产生缩孔、缩松、气孔、砂眼等缺陷 & 一般较低 & 厚度$>3$ mm & 厚度为1 m,质量为100 t & 各种铸造合金 & 差 & 低 & 单件、小批 \\
\cline{2-10}
 & 金属型铸造 & 有限,受金属型限制形状不能太复杂 & 组织较为致密 & 一般 & 铸铝$>3$ mm 铸铁$>5$ mm & 100 kg (铝20 kg) & 非铁金属和灰口铸铁 & 较好,一般不需加工 & 中等 & 大批量 \\
\cline{2-10}
 & 压力铸造 & 高,只受铸型能否制造的限制 & 组织致密,有小气孔,不能热处理 & 较好 & 1 mm (锌合金0.5 mm) & 一般为10$\sim$16 kg & 锌、铝、镁、黄铜等非铁金属 & 好,一般不需加工 & 高 & 大批量 \\
\hline
\multirow{2}{*}{铸件} & 熔模铸造 & 高 & 晶粒较粗大 & 较好 & 0.7 mm & $<25$ kg (通常为5 kg) & 各类合金,难熔合金 & 好 & 高 & 成批、大量 \\
\cline{2-10}
 & 离心铸造 & 有限,多为回转体 & 组织致密,晶粒细小 & 较好 & 一般为3$\sim$5 mm & 200 kg & 非铁金属、灰口铸铁、铸钢 & 外表面好,内表面差 & 低 & 成批、大量 \\
\hline
\multirow{2}{*}{锻件} & 自由锻造 & 简单 & 形成流线组织 & 高 & 不能太薄 & 200 t & 碳钢、合金钢、非铁金属 & 较差 & 较低 & 单件、小批 \\
\cline{2-10}
 & 锤上模锻 & 有限 & 形成流线组织 & 高 & 不能太薄,质量可至几克 & 100 kg (通常$<25$ kg) & 碳钢、合金钢、非铁金属 & 一般 & 较好 & 中、大批 \\
\hline
冲压件 & 冲压 & 复杂 & & 好 & 0.08$\sim$0.13 mm & & 各种金属板材和塑料件 & 高 & 高 & 成批、大量 \\
\hline
型材 & 型材切割 & 简单 & & & & & 合金钢、非铁金属等 & 较差 & 较差 & 单件,中、大批 \\
\hline
拼接件 & 焊接和黏接 & 高 & 焊缝热影响区组织不均匀 & & & 不受限制 & 各类金属及复合材料 & & & 单件,中、大批 \\
\hline
\end{tabular}\end{table}

\subsection{毛坯成形方法选择原则}

毛坯成形方法的选择一般应遵循功能性原则、经济性原则和可行性原则。

1. 功能性原则

功能性原则是指所选毛坯材料的力学性能和工艺性能,应能满足零件使用性能的要求。例如,汽车和火车的壳体,多采用较薄的钢材经压力加工、铆接或焊接而成;机床床身和泵体等机架箱体类零件,一般采用铸造成形;而连杆、主轴等传动件则选用锻件。

对使用性能的要求,一般是在分析零件工作条件的基础上提出的。同样的零件在不同的工作条件下,其材料的选择和相应热处理状态以及成形方法都会存在差异。例如,在各类机械中常见的齿轮,其结构特点、功能基本相同,但实际工作条件却有很大差异,其毛坯材料的选择和制造方法也就不同。机床中的齿轮在较稳定的受力状态下工作,但有时会遇到冲击和振动,因此选用优质碳素结构钢,采用锻造工艺,零件整体正火或调质、齿面高频淬火;汽车中的齿轮,其工作条件远不如机床中的齿轮,要求具有更好的韧性,以免遇到冲击载荷、超载等情况引起事故,因此材料选用低碳合金钢,用锻造工艺方法、齿面先渗碳后淬火等工艺制造;低速机械中的齿轮,受力不大,不需要特殊的工艺方法就能保证其使用要求,可选用灰铸铁和铸造工艺方法制造;仪表齿轮的批量生产则采用压力铸造或精密冲裁等加工制造方法,海上作业的仪表齿轮多采用耐腐蚀、冷脆小的铜合金。

2. 经济性原则

经济性原则是指在满足使用要求的前提下,选择省工、省料、节约能源、成本低廉的方案。单件、小批生产时,应选择通用设备和精度、生产率均较低的成形工艺,如手工造型砂型铸造、自由锻等;大批量生产时,必须考虑先进、高生产率和少(无)切屑的制造方法,如精密铸造、精密模锻、高能率成形等,大批量生产的焊接结构应优先选用低合金结构钢和自动焊接方法。

毛坯的经济性还与其可靠性有关。可靠性高的毛坯表现为缺陷少、质量好、工艺出品率高,零件在正常工作条件下不易破坏,因而经济性也好。在单件、小批生产时,如采用钢材焊接组合结构来代替铸件,既可以提高可靠性,又能减轻重量。

3. 可行性原则

可行性原则是指根据零件形状或材料工艺性能及实际生产条件选择成形方法。如形状复杂的金属制件,特别是具有复杂内腔的复杂件,只能采用铸造而不能选用锻造方法成形,而形状相对简单的金属制件可选用压力加工、焊接成形工艺,亦可选用铸造成形工艺。在几种方案同时可行时,则应考虑本单位的具体生产能力,尽可能利用已有的条件,以便经济合理地成形出高质量毛坯或零件。

图11-4a所示的轴颈套,其毛坯既可以采用铸件(图11-4b),也可以采用模锻件(图11-4c)。当生产批量不大时,采用铸件较合理,因铸模的制造费用比锻模低,一般工厂都能制造。但在大批量生产,且工厂又拥有模锻设备的条件下,选用模锻件较合理,不仅能提高生产效率,还可以得到较高精度和力学性能的毛坯。



\begin{figure}[H]
\centering
\includegraphics[page=41, viewport={14.5 18.8 287.5 107.8}, clip]{12-07572-001FigRevised.pdf}
\caption{轴颈套毛坯的比较}
\end{figure}


\subsection{压气机叶片(毛坯)成形方法举例}

压气机叶片的作用是通过高速旋转逐级压缩并引导气流,为发动机燃烧室提供高压、稳定的空气,以确保燃油高效燃烧,最终产生强大的推进力。压气机叶片在高温、高速、强腐蚀和高疲劳环境下服役,毛坯需采用高性能合金精密成形,并经过严格热处理以确保其强度和精度。

压气机叶片尺寸取决于其在发动机中的位置和发动机尺寸。以高压压气机前级叶片为例,其长度一般为100$\sim$400 mm,宽度一般为50$\sim$150 mm,叶尖厚度仅为1$\sim$3 mm,属于薄壁、复杂空间曲面件;其展弦比大、曲率半径变化明显,兼具薄壁、变截面、扭转曲面和复杂型腔等特点,对组织均匀性和力学性能要求极高。

叶片毛坯一般选用30CrMnSiA高强度钢、镍基高温合金或钛合金等材料,采用模锻成形。毛坯加热时须严格控制加热温度与保温时间,防止晶粒粗化和过烧;模锻时要求金属流线沿叶身形状合理分布,以改善叶片的低倍组织和疲劳性能;模锻成形后,叶片需进行固溶处理与多级时效处理,以获得稳定的组织和弥散分布的强化相,从而提高叶片高温强度、抗蠕变性能和抗疲劳性能;部分合金叶片还需配合表面热处理或涂层工艺,以提升抗氧化和抗腐蚀能力。

压气机叶片(毛坯)成形工艺过程见表11.2。

\begin{table}[H]
\centering
\caption{压气机叶片(毛坯)成形工艺过程}
\small
\begin{tabular}{|c|c|p{4cm}|c|c|}
\hline
工序号 & 工序名称 & 工序内容 & 加工简图 & 设备 \\
\hline
5 & 下料 & 毛坯(GH2132合金毛坯棒料) &\includegraphics[page=42, viewport={327.0 427.8 421.5 484.3}, clip]{12-07572-001FigRevised.pdf} & 锯床 \\
\hline
10 & 车削 & ①端面切削倒角2 mm$\times$45$^{\circ}$ \newline ②表面喷涂高温玻璃润滑保护剂 &\includegraphics[page=42, viewport={19.0 287.8 208.5 374.8}, clip]{12-07572-001FigRevised.pdf} & 车床 \\
\hline
15 & 挤压 & ①将毛坯加热至1030$^{\circ}$C$\pm$10$^{\circ}$C,保温 \newline ②将毛坯倒角一侧向下放至挤压模具型腔中,挤压至要求的尺寸 &\includegraphics[page=42, viewport={227.5 286.8 432.0 375.3}, clip]{12-07572-001FigRevised.pdf} & 挤压机 \\
\hline
20 & 模锻 & ①将挤压件加热至1080$^{\circ}$C$\pm$10$^{\circ}$C,保温 \newline ②将挤压件放至模具型腔,模锻至要求的尺寸 & & 模锻机 \\
\hline
25 & 钳工 & 切除锻件飞边 & & \\
\hline
30 & 热处理 & 固溶和时效热处理 & & 热处理炉 \\
\hline
35 & 检验 & 理化检测、尺寸检测、入库 & & \\
\hline
\end{tabular}\end{table}

\section{毛坯成形质量及其质量控制}

毛坯的外在质量(如尺寸精度、表面粗糙度等)和内在质量(如晶粒大小、应力分布、孔洞和裂纹等)直接影响产品质量和使用性能。因此,应对毛坯的成形质量提出合理要求。对重要零件应制订严格的验收标准,而对于其他非工作面、非外露表面,在保证使用性能的前提下,可以要求低一些。

\subsection{常见毛坯成形缺陷及其预防措施}

1. 铸件常见缺陷及其预防措施

由于铸造工序较多,铸件容易产生各种缺陷,所以要对清理过的铸件进行严格的检验,及时发现缺陷,并且要对铸造缺陷进行认真分析,发现症结,降低废品率。值得注意的是,铸造缺陷的产生往往原因复杂,需经过工艺试验最终确定合理的工艺方案。表11-3列举了常见铸件缺陷的特征及其产生原因。

\begin{table}[H]
\centering
\caption{常见铸件缺陷的特征及其产生原因}
\begin{tabular}{|l|m{4cm}|m{4cm}|m{4cm}|}
\hline
类别 & 缺陷名称和特征 & 主要原因分析 \\
\hline
孔眼 & 气孔:铸件内部或表面有大小不等的孔眼,孔的内壁光滑,多呈圆形 \newline 缩孔:铸件厚截面处出现的形状不规则的孔眼,孔的内壁粗糙 \newline 砂眼:铸件的内部或表面有充满砂粒的孔眼,孔形不规则 \newline 渣眼:孔内充满熔渣,孔形不规则 & 
\includegraphics[page=42, viewport={448.0 287.8 531.0 397.3}, clip]{12-07572-001FigRevised.pdf} & 砂型舂得太紧或型砂透气性差,型砂太湿,起模、修型时刷水过多;浇注系统和冒口设计不当,浇注温度过高;型砂强度不够或局部没舂紧而掉砂,合箱时砂型局部挤坏而掉砂;浇注时没挡住熔渣,或者浇注温度太低,熔渣不易上浮 \\
\hline
表面缺陷 & 冷隔:铸件上有未完全融合的缝隙,接头处边缘光滑 \newline 夹砂:铸件表面有突起物,表面粗糙,中间有砂层   &    
\includegraphics[page=42, viewport={64.5 164.8 162.0 253.8}, clip]{12-07572-001FigRevised.pdf} & 浇注温度过低,浇注系统位置不当或太小;型砂湿压强度较低,内浇道过于集中,使局部砂型烘烤厉害 \\
\hline
形状尺寸不合格 & 偏芯:铸件局部形状和尺寸由于型芯位置偏移而变化 \newline 浇不足:铸件未浇满,外形不完整 \newline 错箱:铸件在分型面处错 &
\includegraphics[page=42, viewport={184.5 167.3 265.0 228.8}, clip]{12-07572-001FigRevised.pdf} & 型芯未固定好,浇注时被冲偏;浇注温度过低;浇注时金属液不够;合箱时上、下型未对准等 \\
\hline
裂纹 & 热裂:铸件开裂,该处表面氧化,呈蓝色 \newline 冷裂:裂纹处表面不氧化或轻微氧化 & \includegraphics[page=42, viewport={289.5 165.8 385.0 241.8}, clip]{12-07572-001FigRevised.pdf} & 铸件壁厚差别大;合金化学成分不当,收缩大;砂型(芯)退让性差;浇注系统开设不当,造成冷却及收缩不均,形成内应力等 \\
\hline
\end{tabular}\end{table}

2. 常见模锻件缺陷及其预防措施

在锻造生产中,由于模锻件生产批量较大,因此质量的控制就更为关键。模锻件在生产中,由于各种原因,会产生一些缺陷。所以,必须仔细分析模锻件产生缺陷的原因,从而采取切实有效的措施,降低模锻件的废品率,提高模锻件质量。

模锻件的缺陷种类较多,原因比较复杂。分析其产生的原因,应从下料、锻造成形过程、锻模设计和制造合理性、模锻件热处理以及清理等生产环节综合分析。表11-4列举了模锻件常见缺陷的主要特征及其预防措施。

\begin{longtable}[H]{|l|p{5cm}|p{7cm}|}
\caption{模锻件常见缺陷的主要特征及其预防措施}\\
\hline
缺陷名称 & 主要特征 & 预防措施 \\
\hline
过热 & 严重过热的模锻件,其组织为粗大晶粒,断口呈脆性断口特征 & ①避免加热温度过高和在高温加热时间太长; \newline ②一般碳钢可通过锻造和正火细化晶粒,改善性能 \\
\hline
过烧 & 严重过烧的毛坯,一击就裂。 \newline 过烧部分的切片,晶粒特别粗大,裂口间的表面呈浅蓝色,表明已经氧化。 \newline 过烧的铝合金毛坯,其表面呈黑色或暗黑色,并有鸡皮状的气泡 & ①避免加热温度太高和在高温加热时间太长; \newline ②毛坯在炉内的位置不要离喷嘴太近 \\
\hline
脱碳 & 毛坯或模锻件表层的碳被部分或全部烧掉而造成贫碳区,深度一般为 $0.01 \sim 0.6 \mathrm{~mm}$ & ①控制锻造和热处理的炉温和加热时间,尤其对高碳钢和硅的质量分数高、易于脱碳的钢; \newline ②控制炉气成分,不使水分过多 \\
\hline
欠压 & 模锻件在高度方向尺寸偏大 & ①选择足够吨位的模锻设备; \newline ②控制锤击轻重及锤击次数; \newline ③不使毛坯体积过大; \newline ④不使终锻温度过低; \newline ⑤选择合适的飞边槽,不使飞边阻力过大; \newline ⑥仔细加工锻模,模膛磨损或变形后及时维修 \\
\hline
局部未充满 & 模锻件的筋处、圆角半径处缺少金属,其尺寸不符合图样要求,或者明显的锻件轮廓不清晰 & ①选择足够的设备吨位; \newline ②毛坯应有足够的体积,制坯后毛坯断面及长度符合终锻要求; \newline ③加热至规定的锻造温度; \newline ④合理地设计模锻件,选择合适的斜度和圆角半径等; \newline ⑤选择具有足够阻力的飞边槽尺寸; \newline ⑥锻前把氧化皮清除干净; \newline ⑦仔细加工锻模,模膛磨损及变形后及时维修 \\
\hline
错差 & 模锻件沿分模面的上、下部分产生了错移 & ①调整好设备的导轨; \newline ②提高锻模加工精度,保证终锻模膛和检验角的相对位置准确; \newline ③把锻模装正及紧固好,及时对设备的装夹面(如锤头的燕尾支承面)进行维修; \newline ④尽量使终锻模膛中心与锻模中心重合。设法平衡错移力,设计导向装置; \\
\hline
冲裂 & 在冲孔的边缘有龟裂或裂纹 & ①检查原材料,保证材料质量; \newline ②冲孔时把毛坯加热到足够的温度; \newline ③一次冲孔的变形量不能太大; \newline ④把冲头预热到规定的温度; \newline ⑤冲头与凹模尺寸设计合理,配合适当 \\
\hline
表面凹穴 & 锻件表面不光滑,有凹穴 & ①不使用表面有严重凹穴的毛坯; \newline ②控制加热温度和加热时间,避免毛坯表面形成太多氧化皮; \newline ③终锻前把模膛中的氧化皮清除干净; \newline ④降低模膛表面粗糙度 \\
\hline
\end{longtable}

3. 焊件常见缺陷及其预防措施

焊接缺陷主要有焊接裂纹、未焊透、夹渣和气孔等。焊接缺陷直接影响焊接接头的质量和焊接结构的安全性,易产生应力集中造成焊件的脆性断裂,其中危害最大的是焊接裂纹和气孔。表11-4列举了常见焊接缺陷的产生原因及其防止措施。

\begin{table}[H]
\centering
\caption{常见焊接缺陷的产生原因及其防止措施}
\begin{tabular}{|l|m{4cm}|m{3cm}|m{3cm}|m{3cm}|}
\hline
缺陷名称 & 缺陷简图 & 缺陷特征 & 产生原因 & 防止措施 \\
\hline
烧穿 &\includegraphics[page=42, viewport={400.0 167.8 484.5 214.8}, clip]{12-07572-001FigRevised.pdf} & 液态金属从焊缝反面漏出而形成穿孔 & 坡口间隙太大;电流太大或焊速太慢;操作不当 & 确定合理的操作间隙;选择合适的焊接工艺参数,掌握正确的运条方法 \\
\hline
未焊透 &\includegraphics[page=42, viewport={54.0 66.3 168.0 133.8}, clip]{12-07572-001FigRevised.pdf} & 焊接时接头根部未完全焊透 & 焊接速度过快,焊接电流太小;坡口角度太小,间隙过窄;焊件坡口不干净 & 选择合适的焊接工艺参数;正确选用坡口形式、尺寸和间隙;加强清理,正确操作 \\
\hline
夹渣 &\includegraphics[page=42, viewport={183.0 67.3 253.5 118.3}, clip]{12-07572-001FigRevised.pdf} & 焊后残留在焊缝金属中的宏观非金属夹杂物 & 工件不洁;焊接电流过小,焊速太快;多层焊时各层熔渣未清除干净 & 多层焊时层层清渣;坡口清理干净;正确选择焊接工艺参数 \\
\hline
气孔 &\includegraphics[page=42, viewport={277.0 70.3 353.5 114.8}, clip]{12-07572-001FigRevised.pdf} & 焊接时,熔池中溶入的气体($\mathrm{H}_{2} 、 \mathrm{~N}_{2} 、 \mathrm{CO}$)在凝固时未能逸出,形成气孔 & 焊件表面不洁;焊条潮湿,焊接电流过小,焊速过快;焊件碳、硅的质量分数高 & 严格清除坡口上的水、锈、油;焊条按要求烘干;正确选择焊接工艺参数 \\
\hline
裂纹 &\includegraphics[page=42, viewport={370.5 67.8 493.0 135.3}, clip]{12-07572-001FigRevised.pdf} & 焊接过程中或焊接后,在焊接接头区域出现的金属局部破裂现象 & 熔池内含有较多的硫、磷或氢;结构刚度大,焊接应力过大,焊接顺序不当 & 焊前预热;限制原材料中的硫、磷的质量分数;选用低氢型焊条,严格对焊条烘干及清理焊件表面 \\
\hline
\end{tabular}\end{table}

\subsection{毛坯成形质量的检验}

1. 外观质量检验

(1) 尺寸和形状检验

通过目测,确保铸件没有因为造型错误、浇不足或由于清理不当而造成铸件较严重的形貌偏差。对照图样,借助检测夹具、卡规和样板检验铸件尺寸及其公差是否符合图样设计要求。

锻件外形尺寸可用卡尺、卡钳、游标卡尺等通用量具进行测量。成批或大量生产时,为了提高检测功效,可用专门量具如卡规、塞规、样板等进行检验。对于外形复杂的锻件,还可用专用仪器来检测。例如,对叶片型面尺寸采用电感量仪,一次可以准确地检测 $20 \sim 34$ 个测量点的尺寸公差。模锻件形状尺寸也可以通过划线方法来检验,但工作量大,一般只用于检查型槽制造精度或磨损程度。

焊件焊缝区的表面气孔、咬边、未焊透、裂纹等外观缺陷可用肉眼或借助大于20倍的放大镜进行检查。

(2) 表面质量检验

采用渗透法、化学腐蚀和磁力探伤等方法检测毛坯表面是否存在裂纹等缺陷。

1) 着色试验

先采用低黏度的染色液处理毛坯表面,再将毛坯彻底漂净并干燥。在裂纹等表面缺陷处存在残留染色液,刷悬浮白垩粉将显示出裂纹形貌。渗透法对毛坯表面的不规则很敏感,因此对光滑的毛坯表面更为可靠。

毛坯表面缺陷也可用荧光检验来显示。其原理是利用细小裂纹的毛细管作用,渗入荧光染色,在暗室中靠紫外线光源检测癖疵或裂纹。

2) 化学腐蚀

毛坯经喷丸处理往往在其表面生成塑性变形薄层,或由于热处理在铸件表面形成氧化皮,着色检验对这类毛坯上的掩盖层是无效的。因此,可采用酸浸溶液暴露薄层以下的表面缺陷。此法采用酸或酸-盐混合物溶液,先将毛坯浸在腐蚀溶液中,然后取出漂净并烘干,局部色斑和由腐蚀产生的无光泽表面的对比将使缺陷显露出来。必须注意,在试验完毕后应将过剩的酸中和并漂净。

3) 磁力探伤

磁力探伤用于检测由铁磁材料制成的毛坯的表层缺陷。将被检毛坯放在磁力探伤机的磁铁两极之间,被检毛坯中有磁力线通过。如果毛坯质地均匀,磁力线将均匀通过;但当毛坯上有裂纹、非铁磁性夹杂等缺陷时,由于这些缺陷的磁导率比铁磁材料低得多,故磁力线将绕过缺陷而发生弯曲现象。若在毛坯表面均匀地撒上铁粉,即可发现缺陷处存在“漏磁”现象,可根据缺陷处被吸附铁粉的形状、厚薄判断出缺陷的位置、大小和形状,如图11-5所示。

\begin{figure}[H]
\centering
\includegraphics[page=41, viewport={308.0 19.3 521.5 68.8}, clip]{12-07572-001FigRevised.pdf}
\caption{焊缝有缺陷时产生漏磁}
\end{figure}

2. 内部缺陷检验

内部缺陷的无损检测主要有射线照相法和超声检测两种方法。

(1) 射线照相法

将毛坯暴露在合适的放射源(如X射线或$\gamma$射线)的辐射下,如毛坯件内部存在孔洞将造成射线局部强度的差异,缺陷处的射线强度大,胶片感光就强,底片冲洗后,在缺陷部位明显地显示出黑色条纹和斑点。由此就可判断缺陷的位置、大小和类别,如图11-6和图11-7所示。

\begin{figure}[H]
\centering
\includegraphics[page=42, viewport={101.5 582.3 227.5 731.3}, clip]{12-07572-001FigRevised.pdf}
\caption{X射线透视示意图}
\end{figure}

\begin{figure}[H]
\centering
\includegraphics[page=42, viewport={241.0 581.3 438.5 722.8}, clip]{12-07572-001FigRevised.pdf}
\caption{X光底片的识别}
\end{figure}

(2) 超声检测

超声波是指频率在 $20 \mathrm{kHz} \sim 20 \mathrm{MHz}$ 范围内的声波。超声波具有很强的穿透能力,可以穿透几米甚至几十米的金属。超声波具有在固体介质中传输和被内表面(如裂纹、缩孔、气孔、夹渣等缺陷的表面)反射的性能,因此能较为准确地发现缺陷的位置及大小。但对缺陷的性质不易准确判断,必须配合其他方法或积累丰富经验加以推理比较才能做出结论。另外,对被检工件的表面精度要求较高,否则由于表面粗糙而产生的假反射会引起错误判断。该方法可以采用单探头(既是发生器又是接收器),也可采用分开的探头。

图11-8所示为超声波探伤示意图。用探头1和2从相互垂直方向探测工件内部纵向裂纹时,可得到的不同脉冲信号。当超声波的穿透方向与裂纹方向平行时,便不能显示裂纹的存在。


3. 成分和显微组织检验

铸件成分和性能通常采用同炉金属浇注试样进行试验的方法来确定。真正代表性的试样应与铸件同时浇注,否则会由于氧化或挥发损失而造成成分偏差。锻件和焊件材料的化学成分以冶炼炉前取样分析为准,但在原材料进厂时,需按技术标准进行复检,合格后才能投产。因此,一般锻件和焊件不进行化学成分检验,对重要的或可疑的锻件和焊件,可从工件上切下一些切屑,采用化学分析或光谱分析检验其化学成分。

毛坯的显微组织检验(高倍检验)是指在光学显微镜下检验其内部(或断口)组织状态与微观缺陷,主要是检验毛坯内部不同组织与夹杂物的状态和分布情况,以及毛坯表面缺陷和渗碳淬硬层等。

\begin{figure}[H]
\centering
\includegraphics[page=42, viewport={125.0 427.3 307.5 546.3}, clip]{12-07572-001FigRevised.pdf}
\caption{超声波探伤示意图}
\end{figure}

4. 毛坯力学性能检验

毛坯在热处理后,一般要做硬度检查。毛坯根据其重要程度及材料性质,还要进行强度、韧性、疲劳、高温蠕变等某一项或几项试验,从而确定其力学性能。

硬度检验,可以判断毛坯在机械加工时是否具有正常的切削加工性能,也可以发现锻件表面脱碳等情况。拉伸试验是金属材料力学性能试验中最基本的方法,它可以测定金属材料在单向静拉力作用下的屈服强度 $R_{e}(R_{r0.2})$、抗拉强度 $R_{m}$、伸长率 $A$ 及断面收缩率 $Z$ 等。冲击试验是测量材料的韧性和对缺口的敏感性。因此,只是对某些受冲击(振动)载荷或在高温、高速条件下工作的零件(如涡轮盘、涡轮叶片,汽车、拖拉机上的曲轴和连杆等)才做冲击试验。对某些在高温、高速、高压下工作的零件,还需进行高温持久试验,目的是检验其在给定温度与时间的特定条件下断裂持久程度。

5. 密封性检验

(1) 水压试验

将容器类毛坯注满水后用小锤轻击,观察是否渗水;再以1.5倍的工作压力进行水压试验,维持 $5 \mathrm{~min}$ 后用锤轻击,观察是否渗水。

(2) 气压试验

把压缩空气通入容器内,在焊缝及其他连接处外表涂肥皂水,观察是否起泡,以检查其密封性;或者在压力容器内通入混有10\%氨气的压缩空气,在焊缝及其他连接处外表上在5\%硝酸汞水溶液中浸过的纸条,若有渗漏,纸条上会留下墨斑。

\section*{复习思考题}

 
11-1 常用机械零件毛坯按照形状可以分为哪几类?

11-2 选择毛坯成形工艺的原则是什么?

11-3 零件的使用要求包括哪些?以机床主轴为例说明其使用要求。

11-4 在选择毛坯成形方法时,首先应该考虑的是什么?

11-5 为什么轴(杆)类零件一般采用锻件作毛坯,而箱体类零件多采用铸件作毛坯?

11-6 为什么齿轮大多是以锻件为毛坯,而带轮、飞轮则多用铸件为毛坯?

11-7 下列零件选用何种材料,采用什么成形方法制造毛坯比较合理?

(1) 形状复杂要求减振的大型机座;

(2) 大批量生产的重载中、小型齿轮;

(3) 薄壁杯状的低碳钢零件;

(4) 形状复杂的铝合金构件。

11-8 试为下列齿轮选择材料及成形方法:

(1) 承受冲击的高速重载齿轮, $\phi 200 \mathrm{~mm}$, 批量2万件;

(2) 不承受冲击的低速中载齿轮, $\phi 250 \mathrm{~mm}$, 批量50件;

(3) 小模数仪表用无油润滑小齿轮, $\phi 30 \mathrm{~mm}$, 批量3000万件;

(4) 卷扬机大型人字齿轮, $\phi 1500 \mathrm{~mm}$, 批量5件;

(5) 钟表用小模数传动齿轮, $\phi 15 \mathrm{~mm}$, 批量10万件。

11-9 砂型铸造的常见缺陷有哪些?其中哪些缺陷与合金流动性有关?

11-10 试分析形成气孔、缩孔等孔眼缺陷的原因。

11-11 常见模锻件的缺陷有哪些?

11-12 过热和过烧对锻件的质量有何影响?生产中采取何种措施消除过热所带来的不利影响?

11-13 冲压件主要缺陷是什么?在生产上如何避免?

11-14 常见焊接缺陷有哪些?其中对焊接接头性能危害最大的有哪几种?

11-15 毛坯常用的无损探伤方法有哪些?试述其工作原理及其适用性。

11-16 简述锻件的质量检验方法。

11-17 毛坯质量检验的内容包括哪些?

11-18 针对铸件的常见缺陷,分析需要何种质量检验方法。


\chapter{现代成形技术及发展趋势}

\section{快速成形技术的类型及应用}

\subsection{快速成形技术的基本原理}

快速成形技术[又称快速原型(rapid prototyping, RP)制造]是20世纪80年代初期出现的一项先进制造技术,是制造技术领域的一次重大突破。快速成形的基本原理是分层/堆积成形,通过逐层堆积材料而形成三维实体。快速成形技术将现代控制技术、CAD/CAM技术、精密伺服驱动技术、激光技术和先进材料成形技术等集于一体,突破了传统的加工模式,大大缩短了产品的生产周期。快速成形技术与反求工程技术相结合,成为快速开发新产品的有力工具。

快速成形的基本原理是由设计者首先在计算机上构造出所需产品的三维模型,用切片软件将其切成具有一定厚度(厚度越小,零件精度越高)的二维平面(轮廓曲线),而后经快速成形设备按此轮廓逐层叠加材料,直至形成三维产品。其工艺过程包括三维模型构造、三维模型的近似处理、三维模型的分层切片处理、一定厚度截面的加工成形、截面叠加形成制件和制件的后处理等过程(图12-1)。下面介绍几种应用较多的快速成形技术。

\begin{figure}[H]
\centering
\includegraphics[page=43, viewport={125.5 574.3 415.5 746.8}, clip]{12-07572-001FigRevised.pdf}
\caption{快速成形的工艺过程}
\end{figure}

\subsection{立体光固化技术}
立体光固化(stereo lithography, SL)技术是世界第一种快速成形技术。由查克・赫尔(Chahes Hull)在1982年发明。其基本原理为:将所设计零件的三维计算机成像数据转换成一系列很薄的模样截面数据,再在快速成形机上,用可控紫外线激光束,按计算机切片软件所得到的轮廓轨迹,对液态光敏树脂进行扫描固化,从而构成模样的一个薄截面轮廓。从零件的最底层截面开始,一次一层(一般为$0.076 \sim 0.381 \mathrm{~mm}$)连续进行,直至三维立体模样制成。而后将模样从树脂液中取出进行硬化处理,再打光、电镀、喷涂或着色即可。

图12-2所示为SL技术原理图。该技术主要特点是可成形任意复杂形状、成形精度高、仿真性强、材料利用率高、性能可靠、性价比高,可用于产品外形评估、功能实验、快速制造电极和各种快速经济模具。其缺点是在液态树脂固化过程中,模样收缩后引起的变形量较大,在型壳焙烧中,因树脂模膨胀易使型壳开裂,需要采用新型树脂制造尺寸精确的半空模样,且所需设备和光敏树脂价格昂贵,成本较高。


\begin{figure}[H]
    \centering
    % Placeholder for Figure 12-2
    \includegraphics[page=43, viewport={54.0 355.8 275.5 497.8}, clip]{12-07572-001FigRevised.pdf}
    \caption{SL技术原理图}
\end{figure}

\subsection{选择性激光烧结技术}
选择性激光烧结(selected laser sintering, SLS)技术是由美国得克萨斯大学开发的,1989年开始推广。SLS技术的工作原理为:在一个充满惰性气体的加工室中,先将一层很薄的可熔粉末沉淀到圆柱形容器底部的可上、下移动的板上,按CAD数据控制$\mathrm{CO}_{2}$激光束的运动轨迹,对可熔粉末材料进行扫描熔化,并调整激光束的强度对$0.125 \sim 0.25 \mathrm{~mm}$厚的粉末进行烧结,从而形成模样的截面形状。未烧结的粉末在制完模样后,可用刷子或压缩空气去除。图12-3所示为SLS技术原理图。

SLS技术所用制模材料包括蜡料、聚碳酸酯和尼龙。另外,高性能的热塑性塑料、陶瓷粉末和金属粉末也正在研究中。SLS技术适合制造那些采用机械加工方法难以成形的几何形状复杂的零件。其发展方向是用金属粉末和陶瓷粉末来直接制造工具、模具和铸造型壳。

\begin{figure}[H]
    \centering
    % Placeholder for Figure 12-3
    \includegraphics[page=43, viewport={298.5 357.3 489.5 542.3}, clip]{12-07572-001FigRevised.pdf}
    \caption{SLS技术原理图}
\end{figure}

\subsection{选择性激光熔化成形技术}
选择性激光熔化成形(selective laser melting, SLM)技术由德国弗劳恩霍夫(Fraunhofer)研究所于1995年首次提出,它是使用高功率激光光束来选择性熔化金属粉末,逐层制造金属零件的技术,即使用高能激光束直接熔化选择区域的金属或合金粉末,通过逐层熔化和堆积,最终制造出具有冶金结合、组织致密的金属零件。SLM技术结合了激光快速熔化和快速冷却凝固的优势,能够获得非平衡态过饱和固溶体以及均匀细小的金相组织,这些特性使得SLM技术制造的金属零件在组织致密度和力学性能方面可与铸锻件相当。图12-4所示为SLM技术原理图。

SLM技术是SLS技术的一种延伸,其区别在于SLM技术使用金属粉末代替SLS技术使用的高分子聚合物粉末,可直接形成致密的金属件。SLM技术突破了传统制造工艺的局限,成为高精度、复杂结构金属零件制造的重要手段,广泛应用于航空航天、医疗、汽车等高端制造领域。

\begin{figure}[H]
    \centering
    % Placeholder for Figure 12-4
    \includegraphics[page=43, viewport={57.0 170.8 302.5 284.3}, clip]{12-07572-001FigRevised.pdf}
    \caption{SLM技术原理图}
\end{figure}

\subsection{熔丝沉积成形技术}
熔丝沉积成形(fused deposition modeling, FDM)技术的原理为:通过计算机控制喷头运动,热塑性丝材在喷头中加热至熔融状态,从喷头的微细喷嘴挤出,并按二维切片薄层轨迹逐步堆积。同理,制造模样从底层开始,一层一层进行,由于热塑性树脂或蜡料冷却很快,便形成一个由二维薄层轮廓堆积并黏结成的立体模样,如图12-5所示。放大后的FDM喷头如图12-6所示。

\begin{figure}[H]
    \centering
    % Placeholder for Figure 12-5
    \includegraphics[page=43, viewport={322.5 176.3 500.5 323.8}, clip]{12-07572-001FigRevised.pdf}
    \caption{FDM技术原理图}
\end{figure}

\begin{figure}[H]
    \centering
    % Placeholder for Figure 12-6
    \includegraphics[page=44, viewport={36.0 588.3 226.5 745.3}, clip]{12-07572-001FigRevised.pdf}
    \caption{放大的FDM喷头}
\end{figure}

与其他快速原型制造技术相比,用FDM技术制模时,其模样上的突出部分无须支撑材料,制出的模样表面光洁,尺寸精度高,且消除了因层间黏结不良而形成的层间台阶、毛刺缺陷和分层问题。

\subsection{激光近净成形技术}
激光近净成形(laser engineering net shaping, LENS)技术,也称为激光立体成形技术,其技术原理为:以高功率激光为能量源,以金属粉末为原料,通过同轴喷嘴将原材料精确送入激光形成的熔池中,使其快速熔化并凝固,逐层沉积,最终形成金属零部件(图12-7)。该技术成形效率高,适合制造大型金属零部件;原材料范围广,可实现非均质和梯度材料制造;组织致密均匀,强度、塑性、抗疲劳性能好,为大型难加工金属零件的制造提供了快速、柔性、低成本制造新途径。

美国于2000年率先将这一先进技术实用化。包括先进飞机承力结构件,如钛合金支架、吊耳、框、梁等;航空发动机零件,如镍基高温合金单晶叶片;战术导弹、人造卫星;超声速飞行器的薄壁结构件,如导弹制导部外壳座、导弹姿态控制系统的铼燃烧室等在内的零件均可使用这一技术成形。美国通用电气公司正在研发世界上最大的激光增材制造设备,打印成形尺寸可达$1 \mathrm{~m}^{3}$。

我国大型金属零件激光立体成形技术水平已经进入世界前列。目前已制造出大型客机C919机头钛合金主风挡整体窗框,钛合金翼肋上、下缘条构件,钛合金飞机角盒,钛合金飞机座椅支座及腹鳍接头,歼20战斗机上重要承力部件等。

激光近净成形技术发展方向包括进一步提升设备加工效率和降低设备成本,与传统加工工艺复合,研发激光增材制造专用金属粉末,加强对激光增材制造过程中零件内应力演化规律、变形开裂行为、凝固组织形成规律以及内部缺陷形成机理等关键基础问题的研究等。

\begin{figure}[H]
    \centering
    % Placeholder for Figure 12-7
    \includegraphics[page=44, viewport={61.5 384.8 288.0 526.8}, clip]{12-07572-001FigRevised.pdf}
    \caption{LENS技术原理图}
\end{figure}

\subsection{均匀金属微滴喷射成形技术}
均匀金属微滴喷射成形(droplet-based manufacturing, DBM)在微小型金属零件成形领域极具发展潜力。该技术原理为(图12-8):在保护气氛中,金属微滴喷射器喷射出尺寸均匀的金属微滴,然后精准地控制这些均匀金属微滴在运动平台上进行逐点、逐层地堆积,使金属微滴迅速凝固并与基体或已凝固的金属层实现良好结合,并逐层成形出实体金属结构。由于金属微滴直径较小,均匀性好,故可实现较高的尺寸精度,又因为金属微滴冷却与凝固速度较快,可得到细小均匀的内部微观组织,能有效提高成形制件的力学性能。

\begin{figure}[H]
    \centering
    % Placeholder for Figure 12-8
    \includegraphics[page=44, viewport={291.5 382.3 478.0 556.3}, clip]{12-07572-001FigRevised.pdf}
    \caption{DBM技术原理图}
\end{figure}

均匀金属微滴喷射成形技术可用于复杂微小金属件、微电子器件、多材质件、智能器件等的快速成形,目前已成形出微小金属薄壁件、金属蜂窝件、微型扑翼飞行器支架等微小型零件,有望应用于太空金属零件在轨制造、立体电路3D打印、异形微小换热结构制造等领域。

目前,均匀金属微滴喷射成形技术正在向着高温、高精度方向发展,打印的材料从铝、镁等熔点较低的有色金属合金发展到铜、铁等熔点较高的合金以及碳纳米管、石墨烯等功能材料,微滴直径也从几百微米逐渐减小至微米级甚至亚微米级。

快速成形技术发展至今,已从最初只能成形零件“原型”(需进行后处理)发展为可以直接制造复杂(金属)零件,即可以自动、直接、快速、精确地将设计思想转变为具有一定功能的原型(构件)或直接制造零件,从而为零件原型制作、功能构件或功能器件制作、新设计思想校验等方面提供了一种高效低成本的实现手段。目前,快速原型制造技术多被称为3D打印或增材制造技术,近年来得到国内外各界特别是工业界的高度重视,其应用范围也扩大至航空、航天、民用等各个领域。

\section{快速铸造与快速精密铸造技术}
快速铸造技术是将快速成形技术与传统铸造技术相结合的一种高效成形工艺,其基本原理是:利用快速成形技术直接或间接制造铸造用消失模、模样、模板、铸型、型芯或型壳等,再结合传统铸造工艺,快捷地制造金属零件。快速铸造技术可以可使铸造生产实现短周期、多品种、低费用和高精度。

快速精密铸造技术是将快速成形技术与精密铸造工艺相结合的一种先进成形工艺,其基本原理是:用快速成形技术制得的原型(如用SL技术制造的光固化树脂原型)代替熔模铸造中的蜡模,在其上涂挂耐火型壳,高温焙烧使树脂原型燃烧去除,最后造型浇注,从而达到缩短生产周期(只需几小时)和降低成本之目的。快速精密铸造的工艺流程见图12-9。

\begin{figure}[H]
    \centering
    % Placeholder for Figure 12-9
    \includegraphics[page=44, viewport={119.0 285.3 425.5 353.3}, clip]{12-07572-001FigRevised.pdf}
    \caption{快速精密铸造工艺流程图}
\end{figure}

用快速成形技术实现快速精密铸造主要有以下方法。

(1) 基于SL原型快速制造零件

用SL原型模代替熔模精密铸造中的蜡模,在SL原型模上直接涂挂耐火浆料(多层),待耐火浆料固化后,再焙烧除去SL原型模,得到耐火材料壳层作为铸造型壳,浇注液态金属从而得到金属件(其工艺与熔模铸造工艺相同)。此方法适用于中等复杂程度的中、小型铸件。

(2) 基于SLS原型生产金属零件

采用陶瓷粉末或包覆黏结剂的陶瓷粉末或覆膜砂作为成形材料,按照铸型CAD模型(包括浇注系统等工艺信息)的轮廓信息精确控制激光束在造型材料粉末层进行扫描,直接烧结成铸造用型壳,或使包覆在陶瓷粉末或覆膜砂表面的黏结剂熔化黏结,逐步堆积得到铸型的型壳,清理出型腔内未烧结的松散粉末,就可用于浇注金属零件。铸型和型芯可分别制造,再装配成完整铸型,也可一体化制造,减少下芯装配带来的误差。此方法适用于中、小型复杂铸件。

(3) 基于FDM原型生产金属零件

采用石蜡和塑料等低熔点材料的FDM原型,以FDM原型代替熔模铸造中的蜡模。用FDM原型制造金属零件的方法与用SL原型生产金属零件的方法相同。此方法适用于中等复杂程度的中、小型铸件。

(4) 基于3DP工艺的型壳直接快速制造工艺

3DP(three dimensional printing)工艺又称为三维印刷,它采用多通道喷头逐点喷洒黏结剂来黏结粉末材料的方法制造原型。根据3DP原理开发的型壳直接快速制造工艺,以陶瓷粉末为造型材料,以硅溶胶为黏结剂。由于陶瓷粉末颗粒尺寸在$75 \mu \mathrm{m} \sim 150 \mu \mathrm{m}$之间,所以型壳直接快速制造工艺造型的铸型表面质量较高,但强度较低,必须经过焙烧之后才能用于浇注金属。对于大型铸件的铸型则需价格高昂、体积庞大的加热设备,所以该工艺不适合大、中型铸件的生产。

(5) 无模铸型快速制造工艺

无模铸型快速制造(patternless casting manufacturing, PCM)工艺是将快速原型理论引入树脂砂造型工艺中,采用轮廓扫描喷射固化工艺,实现了无模铸型的快速制造。首先从零件CAD模型得到铸型CAD模型,由铸型CAD数据得到分层截面轮廓数据,再以层面信息产生控制信息。在计算机的控制下,分别喷射树脂和固化剂的两个喷头在每一层铺好压实的型砂上分别精确地喷射黏结剂和催化剂。喷射出的黏结剂与催化剂发生胶连反应而被固化在一起,其他地方型砂仍为颗粒态干砂。固化完一层后,再黏结下一层,直至形成三维实体。在其内表面涂敷或浸渍涂料之后便可用于浇注金属。该工艺实现了砂型、型芯同时成形的铸型一体化制造,可以方便地制造含自由曲面的铸型和组合零件等,突破了传统工艺的许多约束,具有较高的柔性,尤其适合于制造单件、小批的大、中型铸件。

\section{高能率成形技术}
高能率成形技术是一种在极短时间内释放高能量而使金属变形的成形方法。高能率成形的历史可追溯到一百多年前,但由于成本太高及当时工业发展的局限,该工艺并未得到应用。随着航空及导弹技术的发展,高能率成形技术才进入生产实践中。

高能率成形主要包括利用高压气体使活塞高速运动来产生动能的高速成形、利用火药产生化学能的爆炸成形、利用电能的电液成形以及利用磁场力的电磁成形。

\subsection{高能率成形的特点}
与常规成形方法相比,高能率成形具有以下特点。

(1) 模具简单

高能率成形仅用凹模就可以实现。因此,节省模具材料,缩短模具制造周期,降低模具成本。

(2) 零件精度高,表面质量好

高能率成形时,零件以很高的速度贴模,在零件和模具之间产生很大的冲击力,这不但有利于提高零件的贴模性,而且可有效地减小零件的弹复现象。高能率成形时,毛坯的变形不是由于刚性凸模的作用,而是在液体、气体等传力介质作用下实现的(电磁成形则无需传力介质)。因此,毛坯表面不受损伤,而且可提高变形的均匀性。

(3) 可提高材料的塑性变形能力

与常规成形方法相比,高能率成形可提高材料的塑性变形能力。因此,对于塑性差的成形材料,高能率成形是一种较理想的工艺方法。

(4) 成本低

用常规成形方法需多道工序才能成形的零件,采用高能率成形方法可在一道工序中完成。因此,可有效地缩短生产周期,降低成本。

\subsection{高能率成形的类型}
1. 爆炸成形

爆炸成形是利用爆炸物质在爆炸瞬间释放出巨大的化学能对金属毛坯进行加工的高能率成形方法。爆炸成形装置简单,操作容易,能加工工件的尺寸一般不受设备能力限制,在试制或小批生产大型工件时经济效益尤为显著。

爆炸成形主要用于板材的拉深、胀形、矫形等,还可用于爆炸焊接、表面强化、管件结构的装配、粉末压制等。

爆炸成形时,爆炸物质的化学能在极短时间内转化为周围介质(空气或水)中的高压冲击波,并以脉冲波的形式作用于毛坯,使其产生塑性变形。冲击波对毛坯的作用时间为微秒级,仅占毛坯变形时间的一小部分。图12-10为爆炸拉深示意图,图12-11为爆炸胀形示意图。

\begin{figure}[H]
    \centering
    % Placeholder for Figure 12-10
    \includegraphics[page=44, viewport={121.0 93.3 431.5 252.3}, clip]{12-07572-001FigRevised.pdf}
    \caption{爆炸拉深示意图}
\end{figure}

药包起爆后,爆炸物质以极高的传爆速度在极短的时间内完成爆轰过程。位于爆炸中心周围的介质,在爆炸过程中产生的高温和高压气体的骤然作用下,形成了向四周急速扩散的高压冲击波。当冲击波与成形毛坯接触时,由于冲击波压力大大超过毛坯塑性变形抗力,毛坯开始运动并以很大的加速度积累运动能量。冲击波压力很快降低,当其降低至毛坯的塑性变形抗力时,毛坯位移速度达最大值。这时,毛坯所获得的动能,使其在冲击波压力低于毛坯变形抗力和在冲击波停止作用以后仍能继续变形,并以一定的速度贴模,从而完成成形过程。

\begin{figure}[H]
    \centering
    % Placeholder for Figure 12-11
    \includegraphics[page=45, viewport={140.0 604.8 398.5 745.3}, clip]{12-07572-001FigRevised.pdf}
    \caption{爆炸胀形示意图}
\end{figure}

2. 电液成形

电液成形是利用液体中强电流脉冲放电所产生的强大冲击波对金属进行加工的一种高能率成形方法。与爆炸成形相比,电液成形时能量容易控制,成形过程稳定,操作方便,生产率高,便于组织生产。但由于受到设备容量限制,电液成形还只限于中、小型零件的加工。电液成形可进行拉深、胀形、翻边、冲裁、校形等。

电液成形原理如图12-12所示。该装置主要由两部分组成:充电回路和放电回路。充电回路主要由升压变压器、整流器和充电电阻组成;放电回路主要由电容器、辅助间隙及电极组成。交流电由升压变压器和整流器后变为$2040 \mathrm{~V}$的高压直流电,并向电容器充电。当充电电压达到一定数值时,辅助间隙被击穿,高压电加在由两个电极板形成的主间隙上,将其击穿并放电,形成的强大冲击电流(达$3 \times 10^{4} \mathrm{~A}$以上)在介质(水)中引起冲击波及液流冲击,使金属毛坯成形。

\begin{figure}[H]
    \centering
    % Placeholder for Figure 12-12
    \includegraphics[page=45, viewport={116.5 419.8 425.5 571.3}, clip]{12-07572-001FigRevised.pdf}
    \caption{电液成形原理图}
\end{figure}

3. 电磁成形

电磁成形是利用脉冲磁场对金属进行压力加工的高能率成形方法(图12-13)。电磁成形无须传压介质,可以在真空或高温条件下成形,能量易于控制,成形过程稳定,再现性强,生产率高,易于实现机械化。电磁成形适用于板材、管材的胀形、缩口、翻边、压印、剪切、装配、连接等,尤其适于管子、管接头的连接装配,目前在生产中得到推广应用。

\begin{figure}[H]
    \centering
    % Placeholder for Figure 12-13
    \includegraphics[page=45, viewport={58.0 145.3 274.5 242.3}, clip]{12-07572-001FigRevised.pdf}
    \caption{电磁成形装置示意图}
\end{figure}

\section{半固态成形技术}
\subsection{半固态成形的概念}
20世纪70年代初美国麻省理工学院的弗莱明斯(M. C. Flemigs)和斯宾塞(D. B. Spencer)发现,处于固-液相区间的合金经过连续搅拌后呈现出低的表观黏度,此时在结晶过程中形成的树枝晶被粒状晶代替。这种浆料很容易变形,只要加很小的力就可以充填复杂的型腔,从而开发出一种新的金属成形方法——金属半固态成形。金属半固态成形(semi-solid forming of metals)的原理为:在金属凝固过程中,对其施以剧烈的搅拌或扰动,或改变金属的热状态,或加入晶粒细化剂,或进行快速凝固,即改变初生固相的形核和长大过程,得到一种液态金属母液中均匀地悬浮着一定球状初生固相的固-液混合浆料,利用这种固-液混合浆料直接进行加工成形,或先将固-液混合浆料完全凝固成坯料,根据需要将坯料切分,再将切分的坯料重新加热至固液两相区,用这种半固态坯料进行成形加工。金属半固态成形具有能消除气孔、缩孔,提高零件的力学性能及模具寿命等优点。半固态金属易于搬运和输送,为连续高效自动化生产创造了条件。在节省能源、保护环境方面,金属半固态成形也较传统的铸造方法更为优越。目前,半固态成形技术已应用于生产。

\subsection{半固态金属浆料的制备}
半固态成形技术采用机械或电磁搅拌的方法,可以得到固体组分的颗粒大小为$50 \sim 100 \mu \mathrm{m}$的半固态浆料。图12-14所示为机械搅拌式半固态金属浆料制备装置。对于铝、铜合金和铸铁,该装置可实现固相率为0.5的浆料的连续生产。电磁搅拌法与机械搅拌法相比,减少了搅拌器对浆料的污染,但在制备高固相率的浆料时,搅拌速度会急剧降低,表观黏度迅速增加,使浆料排出困难。图12-15所示为半固态金属浆料电磁搅拌装置。该装置中的4对磁极

以 $0 \sim 3~000~\mathrm{r/min}$ 的速度旋转。为了使浆料产生三维运动，磁铁与旋转中心轴之间有 $10^{\circ}$ 的偏转角，呈螺旋形放置。采用该装置已制造出 ZL101 铝合金为基体，加入平均颗粒尺寸为 $29~\mu\mathrm{m}$ 的 $\mathrm{SiC}$ 颗粒的复合材料锭。

\begin{figure}[H]
    \centering
    % Placeholder for Figure 12-14
    \includegraphics[page=45, viewport={298.5 153.3 495.0 393.8}, clip]{12-07572-001FigRevised.pdf}
    \caption{机械搅拌式半固态金属浆料制备装置}
    \label{fig:12-14}
\end{figure}

\begin{figure}[H]
    \centering
    % Placeholder for Figure 12-15
    \includegraphics[page=46, viewport={61.0 463.3 209.0 743.8}, clip]{12-07572-001FigRevised.pdf}
    \caption{半固态金属浆料电磁搅拌装置}
    \label{fig:12-15}
\end{figure}

\subsection{半固态金属的成形与应用}

金属半固态成形是介于铸造和锻造之间的一种工艺过程，适用于很多常规的成形方法。通常，根据采用的成形设备对其命名，这些设备包括改进的压铸设备、注射成形设备、连续铸造设备和模锻设备等。在研究和应用中，铸造设备在金属半固态成形中占有较大比例，因而金属半固态成形多称为半固态铸造。目前，已经对铝、镁、锌、铜合金及钢、铸铁、镍基超耐热合金、复合材料的半固态成形进行过许多试验研究。生产中实际应用的合金还是直接取自现有的铸造或锻造合金系列，例如，ZAISi7Mn 铸造铝合金，2618、7075 锻造铝合金。应用最多的铸造铝合金为铝硅合金。

根据半固态金属浆料在进入模具内腔之前处理方法的不同，金属半固态成形可分为流变成形（rheoforming）和触变成形（thixo-forming）两大类。流变成形是将获得的半固态金属浆料直接成形；触变成形是将半固态金属浆料首先制成锭料，生产时将定量的锭坯重新加热至半固态后再成形。

1. 流变铸造

流变铸造（rheoforming or rheocast）是将金属液从液相到固相冷却过程中进行强烈搅动，在一定固相率下，直接将所得到的半固态金属浆液压铸或挤压成形（图 12-16），例如可将用电磁搅拌方法制备的半固态合金浆液直接送入压铸机的压射室中成形。该方法生产的铝合金铸件的力学性能较挤压铸件高，与半固态触变铸件的性能相当。但因半固态金属浆液的保存和输送难度较大，故实际投入应用的不多。

\begin{figure}[H]
    \centering
    % Placeholder for Figure 12-16
    \includegraphics[page=46, viewport={230.0 458.3 476.0 730.8}, clip]{12-07572-001FigRevised.pdf}
    \caption{触变压铸成形工艺示意图}
    \label{fig:12-16}
\end{figure}

2. 触变压铸成形

半固态金属的触变压铸成形（thixoforming or thixocast）是在一般液态金属压铸的基础上发展起来的一种新型压铸工艺，其主要包含三个工艺流程：半固态金属原始坯料的制备、原始金属坯料的半固态重熔加热和半固态坯料的触变压铸成形。首先，根据压铸件的质量或尺寸，将已经完全凝固的半固态金属坯料切割成一定的大小，将切割的半固态金属坯料放入加热装置内进行快速半固态重熔加热，并控制坯料的固相率或液相率；然后，将半固态金属坯料送入压铸机的压射室，进行压射成形，并进行适当的保压；最后，卸压开型，取出铸件，清理型腔和喷刷涂料。

由于该方法对坯料的加热、输送易于实现自动化，是目前半固态铸造的主要工艺方法。

3. 触变锻造成形

半固态金属锻造与半固态金属触变压铸实质上并无明显差别，其主要不同之处在于半固态金属在锻造设备上加工成形。锻造半固态金属可以在较低的压力下进行，这使得一些传统锻造无法成形的形状复杂构件可以用半固态金属锻造方法来生产。半固态金属锻件的生产流程为：将合金液冷却至半固态，用电磁搅拌装置搅拌后在水平连铸机上铸成坯料，切断的坯料感应加热至半固态（固相率约为 0.5），而后在压力机上锻造成形，并进行适当的保压，再卸压开型，取出锻件（图 12-17）。材料的加热、运送、夹持和锻造均实现了自动化。触变锻造成形减少了机械加工，减轻了锻件质量，成本大大降低。

\begin{figure}[H]
    \centering
    % Placeholder for Figure 12-17
    \includegraphics[page=46, viewport={78.5 311.3 460.5 429.8}, clip]{12-07572-001FigRevised.pdf}
    \caption{触变锻造成形示意图}
    \label{fig:12-17}
\end{figure}

4. 压射成形

压射成形（injection molding）是将普通压铸与注塑成形这两种工艺结合在一起的成形工艺。该工艺直接把熔化的金属液而不是处理后半固态浆液冷却至适宜的温度，并在一定的工艺条件下压射入型腔成形，取消了熔化设备。该工艺可进行镁合金的压射成形，压射成形设备如图 12-18 所示。颗粒状的 AZ91D 镁合金通过加料器加入到多段控温的圆筒中，为防止氧化从加料器处通入氩气，圆筒内装有可前后运动及旋转的螺旋搅拌器。圆筒用感应与电阻两种方式加热。转动的螺旋搅拌器将加热至半固态的原材料向前输送，材料在混合的同时受剪切力的作用，当一定数量的半固态镁合金进入螺旋搅拌器前方的储存室后，螺旋搅拌器即以预定的速度向前运动，将金属浆料压射入模腔，压射完成后，螺旋搅拌器向后回复到原位。AZ91D 镁合金的压射温度为 $580~^{\circ}\mathrm{C}$，较普通压铸低 $70 \sim 80~^{\circ}\mathrm{C}$，此时金属浆料的固相率为 0.3。压射成形设备由计算机控制，生产的零件尺寸精确，性能也较压铸件更为优越。

\begin{figure}[H]
    \centering
    % Placeholder for Figure 12-18
    \includegraphics[page=46, viewport={137.5 159.8 400.5 279.3}, clip]{12-07572-001FigRevised.pdf}
    \caption{压射成形设备示意图}
    \label{fig:12-18}
\end{figure}

5. 流变轧制成形

流变轧制成形就是将半固态金属浆料直接进行轧制变形，可连续制备金属薄带。图 12-19 为流变轧制成形过程示意图。其工艺过程是：采用电磁搅拌器对浇入搅拌室的合金液进行连续变温搅拌，控制搅拌室中合金液的冷却速度，使其处于液相线温度和固相线温度之间的时间足够长，以便进行较充分的电磁搅拌，获得球状或近球状初生固相的半固态钢铁浆料，通过中间塞杆及提升机构将半固态浆料定量输送至空心水冷双辊轧机的辊缝，对半固态浆料直接进行轧制成形。

\begin{figure}[H]
    \centering
    % Placeholder for Figure 12-19
    \includegraphics[page=47, viewport={54.5 571.3 347.5 743.3}, clip]{12-07572-001FigRevised.pdf}
    \caption{流变轧制成形过程示意图}
    \label{fig:12-19}
\end{figure}

\section{焊接新技术}

近几十年来，各种高效节能的焊接方法、焊接生产的机械化与自动化、焊接机器人等新技术，给制造业带来巨大的变革。

\subsection{焊接机器人和智能化}

焊接机器人是 20 世纪 70 年代开始发展的一种新型自动化焊接设备，现已成为焊接自动化的重要发展方向。目前，焊接机器人已广泛用于车身装焊。

1. 焊接机器人的组成

焊接机器人主要包括机器人和焊接设备两部分。机器人由机器人本体和控制柜（硬件及软件）组成。焊接装备主要由焊接电源（包括其控制系统）、送丝机（电弧焊）、焊枪（钳）等部分组成。智能机器人还应有传感系统，如激光或摄像传感器及其控制装置等。

2. 焊接机器人的主要结构

焊接机器人本体的机械结构主要有两种形式：一种为侧置式（摆式）结构，一种为平行四边形结构。侧置式结构的主要优点是上、下臂的活动范围大，使机器人的工作空间几乎能达一个球体，如图 12-18 所示。因此，这种机器人可倒挂在机架上工作，以节省占地面积，方便地面物件的流动。但是，这种侧置式机械结构为悬臂结构，降低了机器人的刚度，一般适用于负载较小的机器人，用于电弧焊、切割或喷涂。平行四边形机器人的上臂是通过一根拉杆驱动的，拉杆与下臂组成一个平行四边形的两条边，故而得名。早期开发的平行四边形机器人工作空间比较小（局限于机器人的前部），难以倒挂工作。但 20 世纪 80 年代后期以来开发的新型平行四边形机器人，已能把工作空间扩大到机器人的顶部、背部及底部，又没有侧置式机器人的刚度问题，从而得到普遍的重视。这种结构不仅适合于轻型机器人，也适合于重型机器人。近年来，点焊机器人（负载 $100 \sim 150~\mathrm{kg}$）大多选用平行四边形结构的机器人。

\begin{figure}[H]
    \centering
    % Placeholder for Figure 12-20
    \includegraphics[page=47, viewport={47.0 285.8 483.5 542.3}, clip]{12-07572-001FigRevised.pdf}
    \caption{侧置式焊接机器人}
    \label{fig:12-20}
\end{figure}

新一代焊接机器人正朝着智能化方向发展，能自动检测材料的厚度、工件形状、焊缝轨迹和位置、坡口的尺寸和形式、对缝的间隙等，并自动设定焊接工艺参数、电弧焊枪运动点位或轨迹、填丝或送丝速度、焊钳摆动方式等；亦可实时检测是否形成所要求的焊点或焊缝，是否有内部或外部焊接缺陷及排除等情况。智能机器人的关键在于计算机硬件和软件功能的完善和发展，以及各种高功能、高可靠性的传感器的研制。目前，许多国家都在从事这方面的研制工作，带有视觉、听觉、触觉功能的智能焊接机器人已逐步应用于生产。

\subsection{提高焊接生产率}

提高焊接生产率是推动焊接技术发展的重要驱动力。其途径有以下两个方面。

1. 提高焊接熔敷效率

手工电弧焊直接采用铁粉焊条、重力焊条等工艺，埋弧自动焊中采用多焊丝、热焊丝均可提高焊接熔敷效率，且效果显著。例如，三丝埋弧焊，其工艺参数分别为 $2~200~\mathrm{A} \times 33~\mathrm{V}$、$1~400~\mathrm{A} \times 40~\mathrm{V}$、$1~100~\mathrm{A} \times 45~\mathrm{V}$，采用较小坡口截面，背面用挡板或衬垫，$50 \sim 60~\mathrm{mm}$ 的钢板可一次焊透成形，焊速达到 $0.4~\mathrm{m/min}$ 以上，其熔敷效率是手工电弧焊的 100 倍以上。

2. 减少坡口截面及熔敷金属量

近 10 年来，焊接技术研究最突出的成就是窄间隙焊接。它是以气体保护焊为基础，利用单丝、双丝、三丝进行焊接。无论接头厚度如何，均可采用对接接头形式。如焊接厚度为 $50 \sim 30~\mathrm{mm}$，间隙为 $13~\mathrm{mm}$ 左右的钢板，因所需熔敷金属量大为降低，从而大大提高了生产率。窄间隙焊接的主要技术关键是保证焊件两侧熔透和保证电弧中心自动跟踪并处于坡口中心线上。为解决这两个问题，世界各国开发出多种不同方案，出现了种类多样的窄间隙焊接法。电子束焊、激光焊及等离子弧焊可采用对接接头，且不需开坡口，因而得到了较为广泛的应用。

\section{现代成形技术发展趋势}

随着金属间化合物、复合材料、超导材料及其他各类新型功能材料的不断涌现，传统成形技术正面临新的挑战。与新材料的制备与合成相对应，新的成形方法正成为材料加工研究开发的一个重要领域，材料制备和加工一体化、数智化已成为新的发展趋势。根据未来可能面临的挑战和机遇，材料成形加工技术将会出现如下新特征：精密特征，成形精度向净成形（即近无余量成形）的方向发展；优质特征，成形质量向近无缺陷方向发展；快速特征，成形过程的快速化；复合特征，成形方法向复合方向发展；绿色特征，成形加工生产向清洁生产方向发展；数智化特征，成形加工全流程向数据与模型驱动的自主优化决策方向发展。

近年来出现了很多新的精确成形加工技术。例如，在精确铸造成形方面，汽车工业中 cosworth 铸造、消失模铸造及高真空压铸已成为新一代汽车薄壁、高质量铝合金缸体铸件的主要精确铸造成形方法。半固态铸造精确成形技术由于熔体在压力下充型、凝固，从而使铸件的表面及内部质量大大提高。连轧连铸是连续铸造和连续轧制复合的一种短流程成形技术，可使由钢液至成卷的时间从传统生产的 $5~\mathrm{h}$ 降至 $15 \sim 30~\mathrm{min}$。喷射铸造将铸造与粉末冶金工艺相复合，通过将液态金属气体雾化成微细颗粒，然后喷射在一定形状的收集器上制成半成品金属。3D 打印砂型/砂芯技术通过喷射树脂固化砂粒直接制造复杂整体砂型，突破了传统工艺的限制，极大缩短了产品开发周期，满足铸件中小批量及柔性生产。金属板料增量成形工艺基于数字控制，采用预先编制的控制程序对板材进行逐点成形而获得零件，因不需要专用模具，成形极限较高，适用于航空航天、汽车工业中小批量、多品种、形状复杂的板材零件加工。智能微铸锻铣技术将 3D 打印与锻压工艺同步复合，突破了传统铸、锻、焊、铣分离制造模式对大型设备依赖度高、流程长、污染重以及复杂构件难以整体成形的局限，解决了 3D 打印难以获得锻态组织的难题。

新材料的合成与制造，往往利用极端条件作为必要手段，如超高压、超高温、超高真空、极低温、超高速冷却及超高纯等。例如，电磁成形是一种既可控形又可控性的材料成形方法。激光、电子束成形技术种类多样，广泛应用于电子元器件的精密微焊接，汽车和船舶制造中的焊接、切割和成形等。激光-电弧复合焊融合了激光焊的深熔能力与电弧焊的优良桥接性，在船舶、管道等厚板焊接中表现突出。纳米材料是现代材料科学的一个重要发展发方向，作为新型的结构功能材料，其未来应用很大程度上取决于纳米粉末零件成形技术的发展，以保证纳米材料微结构的稳定性和性能的发挥。

随着计算机技术的快速发展，基于智能材料成形技术的模拟仿真已成为材料科学与工程领域的研究热点，高性能、高保真和高效率则是模拟仿真的发展目标。铸造及锻造过程的宏观、模拟在工程中已获得广泛应用，多尺度、跨尺度、分子动力学模拟，特别是微观组织模拟是近年来研究的新热点。通过计算机模拟，可深入研究材料的结构、组成及各种物理、化学过程的宏观、微观变化机制，并由材料成分、结构及制备参数的最佳组合进行材料设计。目前，基于高性能计算的材料多尺度模拟已实现从微观组织到宏观性能的精准预测。依托数字孪生、物联网与人工智能，自适应化实现了成形过程的实时感知、分析、决策与调控，模拟仿真驱动的智能成形已成为工艺开发的标准工具，保障产品质量一致可靠，标志着协同智能制造成为现实。

材料成形技术的发展体现了新材料、新工艺与人工智能深度融合的系统性变革。可以预见，随着人工智能与绿色制造等前沿技术的持续突破，必将推动材料科学技术进入一个全新的发展阶段。

\section*{复习思考题}

 
12-1 什么是快速成形技术？

12-2 快速成形技术有哪几种类型？它们有何异同点？

12-3 哪一种快速成形技术使铸件适时生产成为可能？为什么？

12-4 SL 技术在哪些方面得到应用？

12-5 用快速成形技术实现快速精密铸造的主要方法有哪些？

12-6 什么是高能率成形技术？主要有哪些类型？

12-7 高能率成形技术具有哪些特点？

12-8 爆炸成形主要用于哪类零件的生产？

12-9 什么是半固态成形技术？常用的半固态成形技术有哪些？各有何特点？

12-10 生产中常用什么方法制备半固态金属？

12-11 简述触变压铸的生产过程。

12-12 流变成形和触变成形的区别是什么？

12-13 简述压射成形的生产过程及特点。

12-14 焊接机器人主要由哪些部分组成？

12-15 简述焊接机器人的主要结构及特点。

12-16 提高焊接生产率的主要途径是什么？

12-17 简述现代成形技术发展趋势。

\chapter{切削加工基础知识}

\section{切削加工的分类与特点}

切削加工是利用切削刀具从工件毛坯上切除多余的材料，以获得具有一定形状、尺寸、精度和表面粗糙度的零件的加工方法。切削加工可分为机械加工和钳工两部分。机械加工通过人工操纵机床来实现零件的加工，其主要方法有车削、钻削、镗削、铣削、刨削、磨削等，所用的机床分别有车床、钻床、镗床、铣床、刨床、磨床等。钳工一般是在钳工工作台上通过手持工具进行的，其主要方法有划线、錾削、锯割、锉削、刮削、研磨、钻孔、扩孔、铰孔、锪孔、攻螺纹、套螺纹、机械装配和设备维修等。虽然钳工中的某些工作已实现机械化，但是在机器装配和维修等工艺过程中，钳工比机械加工更为灵活、方便和经济，并容易保证产品的质量，所以钳工是切削加工中不可缺少的一部分。

切削加工主要有以下特点。

1)切削加工的尺寸公差等级可达 IT12 $\sim$ IT3，甚至更高；表面粗糙度 $Ra$ 值为 $25 \sim 0.008~\mu\mathrm{m}$。

2)切削加工一般不受零件形状和尺寸的限制，只要能在机床上实现装夹，大多可进行切削加工，且可实现常见型面，如外圆、内圆、锥面、平面、螺纹、齿形及空间曲面等的加工。切削加工零件的质量范围大，重的可达数百吨，轻的只有几克。

3)在常规条件下，切削加工的生产率一般高于其他加工方法。

4)刀具和工件必须具有一定的强度和刚度，同时刀具材料的硬度必须大于工件材料的硬度。


随着科学技术和现代工业日新月异的飞速发展，切削加工技术必然要与计算机、自动化、系统论、控制论及人工智能、计算机辅助设计与制造（CAD/CAM）、计算机集成制造系统（CIMS）等高新技术及理论相融合，向着精密化、柔性化、智能化方向发展。

\section{切削运动与切削要素}

\subsection{机床的切削运动}

在切削加工过程中，刀具对工件的切削作用是通过刀具与工件之间的相对运动和相互作用来实现的。刀具与工件之间的相对运动称为切削运动，它包括主运动和进给运动。图 13-1 所示为加工不同（零件）表面时的切削运动。

\begin{figure}[H]
    \centering
    % Placeholder for Figure 13-1
    \includegraphics[page=47, viewport={107.5 87.3 435.0 260.3}, clip]{12-07572-001FigRevised.pdf}
    \caption{加工不同(零件)表面时的切削运动}
    \label{fig:13-1}
\end{figure}

1) 主运动(图中 I)是切下切屑所需的最基本的运动。与进给运动相比,它的速度高、消耗机床功率多。主运动一般只有一个。

2) 进给运动(图中 II)是多余材料不断被投入切削,从而加工出完整表面所需的运动。进给运动可以有一个或几个。

主运动和进给运动的适当配合,就可对工件的不同表面进行加工。

机床的切削运动除了主运动和进给运动以外,还有吃刀、退刀、让刀等辅助运动,普通机床的辅助运动大多通过手动来完成。

在切削加工过程中,工件上形成三种表面,如图 13-2 所示。

\begin{figure}[H]
    \centering
    % Placeholder for Figure 13-2
    \includegraphics[page=48, viewport={11.5 502.3 276.5 638.8}, clip]{12-07572-001FigRevised.pdf}
    \caption{车削时在工件上形成的表面}
    \label{fig:13-2}
\end{figure}

1) 已加工表面是工件上经刀具切削后形成的表面。

2) 过渡表面是工件上正在被切除的表面。

3) 待加工表面是工件上等待切除的表面。

\subsection{切削用量}

切削用量是衡量切削运动大小的参数,包括切削速度、进给量和背吃刀量(又称为切削深度),它们是切削过程中不可缺少的因素,称为切削用量三要素。

1. 切削速度 $v_c$

主运动的线速度称为切削速度,即切削刃上选定点相对工件沿主运动方向单位时间内移动的距离。当主运动为转动时,切削速度为
\begin{equation}
v_c = \frac{\pi d n}{1000 \times 60}
\end{equation}
式中: $v_c$——切削速度, m/s;
$d_w$——切削刃上选定点处工件(待加工表面)或刀具直径, mm;
$n$——工件或刀具转速, r/min。

当主运动为往复直线运动时(如刨削运动),以平均速度为切削运动的速度
\begin{equation}
v_c = \frac{2 L n_r}{1000 \times 60}
\end{equation}
式中: $L$——往复运动的行程长度, mm;
$n_r$——主运动每分钟的往复次数, str/min。

2. 进给量 $f$

在单位时间内,刀具在进给方向上相对工件的位移量称为进给量。不同的加工方法,由于所用刀具和切削运动形式不同,进给量的表述和度量方式也不同。

车削时,进给量指工件每转一周,刀具沿进给方向所移动的距离,以 $f$ 表示,单位为 mm/r;在牛头刨床上刨削时,进给量指刀具每往复运动一次,工件所移动的距离,单位为 mm/str;铣削时,因铣刀为多齿刀具,还规定了每齿进给量 $f_z$,即铣刀每转过一个齿,工件沿进给方向所移动的距离,单位是 mm/z。

3. 背吃刀量(又称切削深度) $a_p$

通过切削刃上的选定点,垂直于进给运动方向上测量的主切削刃切入工件的深度,称为背吃刀量。车外圆时,可以用已加工表面和待加工表面之间的垂直距离计算,即
\begin{equation}
a_p = \frac{d_w - d_m}{2}
\end{equation}
式中: $a_p$——背吃刀量, mm;
$d_w$——工件待加工表面直径, mm;
$d_m$——工件已加工表面直径, mm。

\subsection{切削加工的阶段划分}

为了加工出具有一定形状、尺寸和精度的合格零件,保证切削加工质量,工件表面的加工余量往往并不是一次切除,而是逐步减少背吃刀量分阶段加工的。加工工艺过程按照加工性质和目的的不同可以划分为粗加工、半精加工、精加工、精密加工、超精密加工几个阶段。表 13-1 为加工阶段的划分,列出了各加工阶段的目的、尺寸公差等级范围、表面粗糙度及相应加工方法。

工艺过程阶段的划分不是绝对的,对于大多数零件而言,只需经过粗加工、半精加工和精加工三个阶段,但对于精度要求很高的零件还需要进行精密加工甚至超精密加工。对于批量较小、形状简单及毛坯精度较高且加工精度要求低的零件,也可以不必划分加工阶段。

\begin{table}[H]
\centering
\caption{加工阶段的划分}
\label{tab:13-1}
\small
\begin{tabular}{|p{1.5cm}|p{4cm}|p{2.5cm}|p{2.5cm}|p{3cm}|}
\hline
阶段名称 & 目的 & 尺寸公差等级范围 & $Ra$ 值范围/$\mu$m & 相应加工方法 \\
\hline
粗加工 & 快速从毛坯上切除多余材料,使其接近零件的形状和尺寸 & IT12$\sim$IT11 & 25$\sim$12.5 & 粗车、粗铣、粗镗、粗刨、钻孔等 \\
\hline
半精加工 & 进一步提高加工精度和降低表面粗糙度,并留下合适的加工余量,为精加工做准备 & IT10$\sim$IT9 & 6.3$\sim$3.2 & 半精车、半精铣、半精镗、半精刨扩孔等 \\
\hline
精加工 & 一般零件的表面达到技术要求,或为要求很高的主要表面的精密加工做准备 & IT8$\sim$IT6 & 1.6$\sim$0.2 & 精车、精铣、精刨、精镗、粗磨、粗拉、粗铰等 \\
\hline
精密加工 & 在精加工的基础上进一步提高精度和降低表面粗糙度 & IT5$\sim$IT3 & 0.1$\sim$0.008 & 研磨、珩磨、超精加工、抛光等 \\
\hline
超精密加工 & 更高级的亚微米加工和纳米加工 & IT3 以上 & 0.012 或更低 & 金刚石刀具切削、超精密研磨和抛光等 \\
\hline
\end{tabular}\end{table}

\section{刀具材料与刀具构造}

在金属切削过程中,刀具切削部分要受到高温、剧烈摩擦和很大的切削力、冲击力的作用,因此刀具必须选用合适的材料、合理的角度及适当的结构。

\subsection{刀具材料}

1. 刀具材料应具备的性能

(1) 高的硬度

刀具材料的硬度必须高于工件材料的硬度。在常温下,刀具材料的硬度一般应在 60 HRC 以上。

(2) 足够的强度和韧性

刀具材料必须有足够的强度和韧性,以便承受切削力和切削时产生的振动,防止刀具脆性断裂和崩刃。

(3) 高的耐磨性

高的耐磨性即抵抗磨损的能力。一般情况下,刀具材料硬度越高,耐磨性越好。

(4) 高的耐热性

高的耐热性是指刀具材料在高温下仍能保持硬度、强度、韧性和耐磨等性能。

(5) 良好的工艺性和经济性

为便于刀具本身的制造,刀具材料应具有好的切削性能、磨削性能、焊接性能及热处理性能等,而且要追求高的性能价格比。

2. 刀具材料

刀具材料主要根据工件材料、刀具形状和类型以及加工要求等进行选择。目前,在切削加工中常用的刀具材料有碳素工具钢、合金工具钢、高速钢及硬质合金等。其中,国产硬质合金一般分为两大类:一类是由 WC 和 Co 组成的钨钴类(YG 类),韧性较好,但耐磨性较差,适用于加工铸铁等脆性材料;另一类是 WC、TiC 和 Co 组成的钨钛钴类(YT 类),耐热性、耐磨性较好,但韧性较差,一般适用于加工钢件。在现行标准(GB/T 18376.1—2008)中,硬质合金牌号分为 P、M、K、N、S、H 六大类,可满足不同使用要求。超硬刀具材料目前应用较多的有陶瓷、立方氮化硼、金刚石。刀具材料的种类、性能和用途见表 13-2。

\begin{table}[H]
\centering
\caption{刀具材料的种类、性能和用途}
\label{tab:13-2}
\footnotesize\setlength{\tabcolsep}{3pt}
\begin{tabular}{|p{1.5cm}|p{4cm}|p{1.5cm}|p{1.5cm}|p{1.5cm}|p{2.5cm}|p{2.5cm}|}
\hline
种类 & 常用牌号 & 硬度 & 抗弯强度 $R_b$/GPa & 红硬性 /$^\circ$C & 工艺性能 & 用途 \\
\hline
优质碳素工具钢 & T8A$\sim$T10A, T12A, T13A & $\geqslant$62 HRC & 2.16 & 200 & 可冷、热加工成形,刃磨性能好 & 手动工具,如锉刀、锯条等 \\
\hline
合金工具钢 & 9SiCr, CrWMn & $\geqslant$62 HRC & 2.35 & 250$\sim$300 & 可冷、热加工成形,刃磨性能好,热处理变形小 & 用于低速成形刀具,如丝锥、板牙、铰刀 \\
\hline
高速钢 & W18Cr4V, W6Mo5Cr4V2 & $\geqslant$63 HRC & 1.96$\sim$4.41 & 550$\sim$600 & 可冷、热加工成形,刃磨性能好,热处理变形小 & 中速及形状复杂的刀具,如钻头、铣刀等 \\
\hline
硬质合金 & \makecell{P01, P10, P20, P30, P40 \\ M01, M10, M20, M30, M40 \\ K01, K10, K20, K30, K40 \\ N01, N10, N20, N30 \\ S01, S10, S20, S30 \\ H01, H10, H20, H30} & 88.5$\sim$92.3 HRA & 0.7$\sim$1.8 & 800$\sim$1000 & 粉末冶金成形,多镶片使用,性能较脆 & 用于高速切削刀具,如车刀、刨刀、铣刀 \\
\hline
涂层刀具 & TiC, TiN, TiN-TiC & 3200 HV & 1.08$\sim$2.16 & 1100 & 刀具材料表面的硬度和耐磨性大为提高 & 用于高速切削刀具,如车刀、刨刀、铣刀,但切削速度可提高 30\%左右,同等速度下,寿命提高 2$\sim$5 倍多 \\
\hline
陶瓷 & AM, AMT, SG4, AT6 & 93$\sim$94 HRA & 0.4$\sim$0.785 & 1200 & 硬度高于硬质合金,脆性大于硬质合金 & 精加工优于硬质合金,可加工淬火钢 \\
\hline
立方氮化硼(CBN) & FN, LBN-Y & 7300$\sim$9000 HV & & 1300$\sim$1500 & 硬度高于陶瓷,性能较脆 & 切削加工优于陶瓷,可加工淬火钢 \\
\hline
人造金刚石 & & 10000 HV 左右 & & 600 & 硬度高于 CBN,性能较脆 & 用于非铁金属精密加工,不宜切削钢铁合金 \\
\hline
\end{tabular}\end{table}

为了提高高速钢的硬度和耐磨性,常采用如下措施。

(1) 在高速钢中增添新的元素

在通用高速钢的成分中,提高碳的质量分数,形成高碳高速钢;提高碳化钒的质量分数,形成高钒高速钢;增添钴、铝等元素形成钴高速钢、铝高速钢。这些高速钢也被称为高性能高速钢,其加工性能更好,刀具耐用度为通用高速钢刀具的 1.5$\sim$3 倍。

(2) 改进刀具制造的工艺方法

用粉末冶金法制造的高速钢称为粉末冶金高速钢,它可以有效解决一般熔炼高速钢在铸锭中易产生粗大的碳化物共晶偏析问题。粉末冶金高速钢得到的是细小均匀的结晶组织,使这种钢具有高的硬度和强度,并减小了热处理变形,适用于制造高精度刀具,切削各种难加工材料。

(3) 进行涂层处理

采用物理气相沉积方法对高速钢刀具表面进行 TiC(灰色)或 TiN(金色)等涂层处理后,其表面硬度可达 80 HRC 以上,刀具耐用度可提高 2$\sim$5 倍,切削力和切削温度可降低约 25\%。适用于钻头、铣刀、拉刀及齿轮刀具等。

为了克服常用硬质合金强度和韧性低、脆性大、易崩刃的缺点,常采用如下措施。

(1) 调整化学成分

在钨钛钴类硬质合金中加入适量的碳化钽(TaC)、碳化铌(NbC),可使其具有硬度高、耐热性好和强度高、韧性好等特点。这类硬质合金也被称为通用硬质合金,可用于难加工材料的切削,也可加工钢材、铸铁和非铁金属。

(2) 采用涂层刀片

在韧性较好的硬质合金(如 K 类)基体表面,涂敷一层硬度和耐磨性更高的、厚度只有几微米的 TiC 或 TiN,使之成为既有高硬度和耐磨性的表面,又有强而硬的基体的刀具材料。TiC、TiN 复合涂层以及 Al$_2$O$_3$ 涂层,其使用效果更佳。

\subsection{刀具结构}

切削刀具种类很多,如车刀、钻头、刨刀、铣刀等。它们的几何形状各异,复杂程度不同。其中,车刀是最常用、最简单和最基本的切削刀具,各种复杂刀具都可以看作是以车刀为基本形态演变而成的。

1. 车刀的结构形式及组成

车刀的结构形式有整体式、焊接式、机夹式、机夹可转位式等几种,如图 13-3 所示。整体式车刀结构简单,但对贵重的刀具材料消耗较大;焊接式车刀刚性好,灵活性较大,可以根据加工要求,方便地磨出所需要的角度,但硬质合金刀片经过高温焊接和刃磨后,易产生内应力和裂纹,使其切削性能下降,对提高生产率不利;机夹式车刀可避免焊接高温带来的缺陷,提高刀具切削性能,并能使刀柄多次使用;机夹可转位式车刀所用的硬质合金刀片具有多个切削刃,当一个切削刃用钝后,不需要重磨,只要松开夹紧元件,将刀片转换一个位置,重新夹紧,即能继续使用,可获得较大的经济效益。

车刀由切削部分(刀体)和夹持部分(刀柄)组成。切削部分又由前面、主后面、副后面、主切削刃、副切削刃和刀尖(简称三面、二刃、一尖)组成,如图 13-3a 所示。

\begin{figure}[H]
    \centering
    % Placeholder for Figure 13-3
    \includegraphics[page=48, viewport={289.0 502.8 535.0 746.3}, clip]{12-07572-001FigRevised.pdf}
    \caption{车刀的结构形式}
    \label{fig:13-3}
\end{figure}

	\[
\text{切削部分} \begin{cases}
	\text{三面} & \begin{cases}
		\text{前面：切屑流出所经过的表面} \\
		\text{主后面：与工件过渡表面相对的表面} \\
		\text{副后面：与工件已加工表面相对的表面}
	\end{cases} \\[6pt] % 调整行间距，避免拥挤
	\text{二刃} & \begin{cases}
		\text{主切削刃:} \text{刀具前面与主后面的交线，承担主要的切削工作} \\
		\text{副切削刃:} \text{刀具前面与副后面的交线，一般仅承担微量的切削}
	\end{cases} \\[6pt]
	\text{一尖} & \text{刀尖：主、副切削刃的交点。为强化刀尖，常磨成圆弧形或直线形的过渡刃}
\end{cases}
\]

2. 车刀的标注角度

为了确定上述表面和切削刃的空间位置,先介绍三个相互垂直的辅助平面,如图 13-4 所示。

1) 切削平面是通过主切削刃上某一点并与工件过渡表面相切的平面。

2) 基面是通过主切削刃上某一点并与该点切削速度方向相垂直的平面。

3) 主剖面是通过主切削刃上某一点并与主切削刃在基面上的投影相垂直的平面。

车刀的标注角度是指在刀具图样上标注的角度,也称为刃磨角度。主要的标注角度有前角 $\gamma_o$、后角 $\alpha_o$、主偏角 $\kappa_r$、副偏角 $\kappa'_r$ 和刃倾角 $\lambda_s$,如图 13-5 所示。

\begin{figure}[H]
    \centering
    % Placeholder for Figure 13-4
    \includegraphics[page=48, viewport={11.5 319.3 216.5 472.3}, clip]{12-07572-001FigRevised.pdf}
    \caption{确定刀具角度的辅助平面}
    \label{fig:13-4}
\end{figure}

1) 前角 $\gamma_o$。在主剖面中测量,是前刀面与基面之间的夹角,有正负之分,如图 13-5b 所示。

2) 后角 $\alpha_o$。在主剖面中测量,是主后刀面与切削平面之间的夹角。

3) 主偏角 $\kappa_r$。在基面中测量,是主切削刃在基面上的投影与进给运动方向之间的夹角。

\begin{figure}[H]
    \centering
    % Placeholder for Figure 13-5
    \includegraphics[page=48, viewport={226.5 321.8 536.5 467.3}, clip]{12-07572-001FigRevised.pdf}
    \caption{车刀的主要角度}
    \label{fig:13-5}
\end{figure}

4) 副偏角 $\kappa'_r$。在基面中测量,是副切削刃在基面上的投影与进给运动反方向之间的夹角。

5) 刃倾角 $\lambda_s$。在切削平面中测量,是主切削刃与基面之间的夹角,有正负之分。

3. 车刀的工作角度

上述车刀角度是在静止参考系中,假定车刀刀尖和工件回转轴线等高、刀柄中心线垂直于进给方向,且不考虑进给运动对坐标平面空间位置的影响等条件下标注的角度。在实际切削过程中,这些条件往往会改变,致使刀具切削时的几何角度不等于上述标注的角度。刀具在切削过程中的实际切削角度称为工作角度。

安装车刀时,刀尖如果高于或低于工件回转轴线,则切削平面和基面的位置将发生变化,如图 13-6 所示。当刀尖高于工件回转轴线时,前角增大,后角减小(图 13-6a);反之,前角减小,后角增大(图 13-6c)。

如果车刀刀柄中心线与进给方向不垂直,车刀的主、副偏角将发生变化,如图 13-7 所示。刀柄右偏,则主偏角增大,副偏角减小(图 13-7a);反之,主偏角减小,副偏角增大(图 13-7c)。

可见,刀具安装的正确与否,对切削是否顺利,工件加工表面是否光洁等有较大的影响。如果安装不当,就不能发挥它应有的作用。

\begin{figure}[H]
    \centering
    % Placeholder for Figure 13-6
    \includegraphics[page=48, viewport={110.0 196.8 429.0 293.8}, clip]{12-07572-001FigRevised.pdf}
    \caption{外圆车刀安装对前角和后角的影响}
    \label{fig:13-6}
\end{figure}

\begin{figure}[H]
    \centering
    % Placeholder for Figure 13-7
    \includegraphics[page=48, viewport={114.0 76.3 430.0 165.8}, clip]{12-07572-001FigRevised.pdf}
    \caption{车刀安装偏斜对主偏角和副偏角的影响}
    \label{fig:13-7}
\end{figure}

进给运动对工作角度也有一定的影响。切削过程中由于进给运动的存在,当纵向走刀车外圆时,工件上形成的加工表面实际上是一个螺旋面,使实际的切削平面和基面都要偏转一个螺旋面的升角 $\psi$,从而引起工作前角 $\gamma_{oe}$ 加大,工作后角 $\alpha_{oe}$ 减小,如图 13-8 所示。车外圆(或内孔)时,进给量 $f$ 小($\psi$ 小),进给运动对工作角度的影响不大,可忽略不计。但当车螺纹时,特别是车削螺距较大($\psi$ 大)的螺纹时,一定要考虑螺旋升角 $\psi$ 对工作前角和工作后角的影响。

\begin{figure}[H]
    \centering
    % Placeholder for Figure 13-8
    \includegraphics[page=49, viewport={138.0 636.8 403.0 743.3}, clip]{12-07572-001FigRevised.pdf}
    \caption{进给运动对 $\gamma_{oe}$ 和 $\alpha_{oe}$ 的影响}
    \label{fig:13-8}
\end{figure}

\section{刀具切削过程及其物理现象}

金属切削过程中的许多物理现象,例如切削力、切削热、刀具磨损、加工表面质量等,都是以切屑的形成过程为依据的,而实际中的许多问题,如振动、断屑等,都同切削过程紧密相关。因此,研究金属切削过程,对于发展切削加工技术,保证加工质量,降低加工成本,提高生产率,具有重要的意义。

\subsection{切削过程及切屑类型}

切削过程就是利用刀具从工件上切下切屑的过程。金属切削过程其实是一种挤压变形。在这一变形过程中会产生许多现象，有弹性变形、塑性变形、切削力、切削热、刀具磨损以及加工表面质量的变化等。

1. 切屑形成过程

对塑性金属以缓慢的速度进行切削时，切屑形成的过程如图 13-9a 所示。当刀具逐渐向工件推进时，切削层金属在始滑移面 $OA$ 以左发生弹性变形，越靠近 $OA$ 面，弹性变形越大。在 $OA$ 面上，应力达到材料的屈服强度 $R_e$ 时，会发生塑性变形，产生滑移现象。随着刀具的继续推进，原来处于始滑移面上的金属不断向刀具靠拢，应力和变形也继续加大。在终滑移面 $OE$ 上，应力和变形达到最大值。此时，切削层金属越过 $OE$ 面，沿剪切面与工件母体分离，即切屑沿着刀具前面排出，完成切离阶段。经过塑性变形的金属，其晶粒沿大致相同的方向伸长。

由此可见，塑性金属的切削过程是挤压和切削变形的过程，经历了弹性变形、塑性变形、挤裂和切离四个阶段。

\begin{figure}[H]
\centering
    \includegraphics[page=49, viewport={104.5 506.8 435.0 601.8}, clip]{12-07572-001FigRevised.pdf}
\caption{切屑的形成过程}
\label{fig:13-9}
\end{figure}

切削塑性金属材料时，在刀具与工件接触的区域产生三个变形区，如图 13-9b 所示。$OA$ 与 $OE$ 之间的区域 I 为第一变形区，或称基本变形区，此区域是切削层金属产生剪切滑移和大量塑性变形的区域，常用来说明切削过程的变形情况。切屑与刀具前面摩擦的区域 II 称为第二变形区，或称摩擦变形区，此区域的状况对积屑瘤的形成和刀具前面的磨损有很大影响。工件已加工表面与刀具后面接触的区域 III 称为第三变形区，或称加工表面变形区，此区域的状况对工件表面的加工硬化、残余应力和刀具后面的磨损有很大影响。

2. 切屑类型

在金属切削的过程中，所切削的工件材料力学性能各异、刀具的前角以及采用的切削用量不同，从而会对切削过程产生不同的影响，形成不同的切屑种类。从切屑形成机理出发，根据切屑产生基本变形过程中的特征以及变形程度的不同，一般把切屑分为如图 13-10 所示的几种类型。


(1)带状切屑

带状切屑（图 13-10a）是最常见的一种切屑，外观呈延绵的长带状，底层光滑，外表面呈毛茸状，无明显裂纹。一般在加工塑性金属材料，使用较大的切削速度、较小的进给量和较大的刀具前角时，形成带状切屑。形成带状切屑时，切削变形小，切削力较平稳，加工表面粗糙度值小，但是切屑连续容易产生缠绕，划伤已加工表面，因此必须采取有效的断屑、排屑措施。

\begin{figure}[H]
\centering
        \includegraphics[page=49, viewport={119.0 376.8 415.0 476.3}, clip]{12-07572-001FigRevised.pdf}
\caption{切屑类型}
\label{fig:13-10}
\end{figure}

(2)节状切屑

当粗加工中等硬度的材料，采用较大的进给量、较低的切削速度，刀具前角很小时，剪切面上的应力超过材料的抗剪强度，并且裂纹扩展，在整个剪切面上产生破裂，以致形成节状的分离切屑，这种切屑称为节状切屑（图 13-10b）。形成节状切屑时，产生的切削振动较大，工件表面质量较差，因此应当对切削过程进行控制，使节状切屑转变为带状切屑，以提高加工质量。可以采取高速切削和采用较大的前角、较小的进给量，以及采用切削液以减小切削变形等措施，从而改变切屑的类型。

(3)崩碎切屑

在加工铸铁、黄铜等脆性材料时，切削层金属发生弹性变形后，一般不经过塑性变形就发生挤裂或脆断，突然崩落，形成不规则的细粒状碎片，这种切屑称为崩碎切屑（图 13-10c）。当切削厚度越大，刀具前角越小时，越容易形成这种切屑。形成崩碎切屑时，切削力变化大，容易产生振动、冲击，切削热和切削力集中在切削刃和刀尖附近，使刀具的寿命降低，加工表面粗糙，加工质量变差。因此，在生产中应当采取减小切削厚度、适当提高切削速度等措施，使切屑转化为针状或片状，力求避免产生崩碎切屑。


上述情况说明切屑类型随切削过程的具体条件不同而改变。在生产中，可以根据具体条件采取不同措施得到所需的切屑类型，以保证切削加工的顺利进行。

\subsection{积屑瘤和表面变形强化}

1. 积屑瘤

在切削铝合金、钢等塑性材料时，当切削速度不高而且形成带状切屑的情况下，常有一些来自工件或切屑的金属黏结在刀具前面靠近切削刃的附近，形成硬度很高（一般为工件硬度的 2\textasciitilde3 倍）的楔块，称作积屑瘤。

在切削过程中，切屑沿刀具前面流出，在一定的温度和压力的作用下，与前面接触的切屑底层受到很大的摩擦阻力，致使这一层金属的流速减慢，形成一层很薄的滞流层。当刀具前面对滞流层的摩擦阻力大于切屑或工件材料的结合力时，就会有一部分金属黏附在刀具上，形成积屑瘤（图 13-11）。积屑瘤形成以后会不断长大，长到一定高度后因不能承受切削力而破裂脱落。因此，积屑瘤是一个反复长大、脱落的动态形成过程。

\begin{figure}[H]
\centering
    \includegraphics[page=49, viewport={26.5 192.8 343.5 332.8}, clip]{12-07572-001FigRevised.pdf}
\caption{积屑瘤}
\label{fig:13-11}
\end{figure}

积屑瘤的形成必须具备两个条件：一是切削塑性金属材料；二是采用中等切削速度（$v_c = 5 \sim 60 \text{ m/min}$）切削。而当采用低速（$v_c < 5 \text{ m/min}$）和高速（$v_c > 60 \text{ m/min}$）时，切屑底层和刀具前面之间的摩擦因数较小，一般不会产生积屑瘤。

积屑瘤的存在有利也有弊。积屑瘤的不稳定使切削深度和厚度不断发生变化，影响加工精度，并引起切削力的变化，产生振动、冲击，而且脱落的积屑瘤碎片黏附在已加工表面上，使已加工表面粗糙。因此，在精加工过程中，应采取改变切削速度、选用适当的切削液等措施避免积屑瘤的产生。由于积屑瘤的硬度高于刀具，可以代替切削刃进行切削，起到保护刀具的作用，同时积屑瘤的存在，增大了刀具的实际工作前角，使切削过程轻快，因此粗加工时积屑瘤的存在是有利的。

2. 表面变形强化和残余应力

切削金属材料时，工件已加工表面表层的强度、硬度明显提高而塑性下降的现象称为表面变形强化。在切削过程中，切削层金属产生变形、分离的同时，切削力会扩展到切削层以下，使即将成为已加工表面的表层金属产生一定的塑性变形。塑性变形越大，表面强化现象越严重。表面变形强化可以提高工件表面的耐磨性和疲劳强度，但也会加剧刀具的磨损，并给后续工序带来不便。在切削加工中，可以通过控制工件表层金属塑性变形的大小，适当控制表面变形强化的程度。

残余应力是指在外力消失后，残留在工件材料内部而总体又保持平衡的内应力。切削过程中，在金属的塑性变形以及切削力、切削热等因素的综合作用下，在已加工表面的表层金属内产生一定的残余应力。残余应力一般与表面变形强化同时出现，而残余应力的存在容易引起工件的变形，影响加工精度的稳定性。因此，在切削加工过程中，应采用减小金属塑性变形、降低切削力和切削温度等方法，尽量减小残余应力，以保证加工表面的质量。

\subsection{切削力和切削功率}

1. 切削力

切削加工时，刀具切入工件，使被加工材料发生变形并成为切屑所需的力称为切削力。切削力来源于切屑作用于刀具前面及工件作用于后面上的弹性变形和塑性变形抗力，以及切屑和刀具前面之间、工件已加工表面和后面之间的摩擦力。这些力构成了作用在刀具上的总切削力，如图 13-12 所示。


\begin{figure}[H]
\centering
\includegraphics[page=49, viewport={355.0 192.8 510.5 350.8}, clip]{12-07572-001FigRevised.pdf}
\caption{切削力的产生}
\end{figure}
    
\begin{figure}[H]
\centering
\includegraphics[page=49, viewport={32.0 24.3 171.5 164.8}, clip]{12-07572-001FigRevised.pdf}
\caption{切削力的分解}
\end{figure}


总切削力 $F$ 是一个空间矢量，很难直接测量。为了分析 $F$ 对加工的影响，在进行机床、刀具设计和工艺分析时，常将 $F$ 分解为三个互相垂直的切削分力。切削力的分解如图 13-13 所示。


1)主切削力（切向力）$F_z$ 是总切削力 $F$ 在切削速度方向上的分力，一般占总切削力的 80\% \textasciitilde 90\%，是计算机床传动、主运动系统零件和刀具强度、硬度的主要依据。主切削力过大，会使刀具崩刃或使机床发生闷车现象。

2)进给抗力（轴向力）$F_x$ 是总切削力 $F$ 在进给方向上的分力，它所消耗的功率一般只占机床总功率的 1\% \textasciitilde 5\%，是设计和验算进给机构的依据。

3)背向分力（径向力）$F_y$ 是总切削力 $F$ 在背吃刀量上的分力。切削时，在这一方向上的速度为零，因此 $F_y$ 不做功。但是，$F_y$ 作用于工件刚性最弱的方向，容易使工件产生弯曲变形，影响加工精度。


三个切削分力与总切削力存在以下关系：
\[ F^2 = F_x^2 + F_y^2 + F_z^2 \]

影响切削力的因素很多，主要有以下几方面。


(1)工件材料

    若工件材料的强度、硬度、冲击韧度及塑性越大，则切削时的变形抗力越大，切削力就越大；加工硬化程度越大，切削力也会越大。

(2)切削用量

    切削用量中，背吃刀量和进给量是影响切削力的主要因素。当 $a_p$ 和 $f$ 增大时，切削面积 $A_0$ 增大，因而切削力也会明显地增大。实验表明，当背吃刀量增加一倍时，切削力也增加一倍；进给量增加一倍时，切削力增大 75\% 左右。

(3)刀具角度

    刀具角度中，前角对切削力的影响最大。加工塑性材料时，前角越大，切削变形越小，切削力就越小。主偏角对 $F_x$ 和 $F_y$ 两个分力影响较大，当增大主偏角时，$F_x$ 增大而 $F_y$ 减小。

(4)切削液

    在切削加工时喷注切削液，可以减小摩擦阻力而使切削力降低。


2. 切削功率

切削功率是三个切削分力消耗功率的总和。由于背向分力 $F_y$ 方向的运动速度为零，故不消耗功率。进给抗力 $F_x$ 消耗的功率也很小（占 $P_m$ 的 1\% \textasciitilde 5\%），一般忽略不计。故切削功率 $P_m$ 可以简化为
\[ P_m = 10^{-3} F_z v_c \]
式中：$P_m$——切削功率，kW；
$F_z$——主切削力，N；
$v_c$——切削速度，m/s。

机床电动机功率为
\[ P_e \geqslant P_m / \eta_m \]
式中：$P_e$——机床电动机功率，kW；
$\eta_m$——机床传动效率，一般取 0.75 \textasciitilde 0.85，新机床取大值，旧机床取小值。

\subsection{切削热和切削温度}

1. 切削热

在切削过程中，因变形和摩擦所消耗的功的绝大部分变成了热能，产生出大量的热，即切削热。切削热来源于三个变形区：由于切削层金属弹性变形和塑性变形而产生的热（I 变形区）；由于切屑与刀具前面摩擦而产生的热（II 变形区）；由于工件与刀具后面摩擦而产生的热（III 变形区）。切削热通过切屑、工件、刀具以及周围的介质向外传散，如图 13-14 所示。

各部分传散切削热的比例根据加工方式不同而异。用高速钢车刀车削钢材时，切屑带走的切削热为 50\% \textasciitilde 86\%，工件带走的为 10\% \textasciitilde 40\%，车刀带走的为 3\% \textasciitilde 9\%，1\% 左右的传入空气；以上述条件钻削钢件时，切屑带走的切削热为 28\%，工件带走的为 14.5\%，钻头带走的为 52.5\%，5\% 传入周围介质。


\begin{figure}[H]
\centering
\includegraphics[page=49, viewport={182.0 18.8 358.0 155.3}, clip]{12-07572-001FigRevised.pdf}
\caption{切削热的来源与传散}
\end{figure}


切削热传入刀具，会引起刀具温度升高（高速切削时，刀体温度可达 1000 $^\circ$C 以上），使刀具材料的性能发生变化，加速刀具的磨损；切削热传入工件，可引起工件变形，从而产生形状和尺寸误差，降低加工精度；切削热传入机床、夹具，也会对加工精度和表面质量产生影响。

因此，在切削加工中，应设法减少切削热，改善散热条件，以减小高温对刀具和工件的不良影响。

2. 切削温度

切削温度是指刀具前面与切屑接触区域的平均温度。切削温度的高低取决于切削热的产生与传散情况。产生的切削热越多，传出得越慢，切削温度越高。切削温度主要受工件材料、切削用量、刀具角度和冷却条件等因素的影响。

切削过程中，凡是增大切削力和切削功率的因素都会使切削温度上升，而有利于切削热传出的因素都会降低切削温度。

实验表明，在切削用量中，切削速度增大一倍，切削温度一般增加 20\% \textasciitilde 33\%；进给量增大一倍，切削温度大约升高 10\%；背吃刀量增大一倍，切削温度大约只升高 3\%。从降低切削温度出发，相同的金属切除率，应优先采用大的背吃刀量，较小的进给量，最后再确定合适的切削速度为宜。总之，在切削加工中，可通过合理选择切削用量及刀具角度、提高工件材料和刀具材料的热导率、充分浇注切削液等措施，使切削温度降低。

\subsection{刀具的磨损和耐用度}

1. 刀具的磨损

刀具在切削过程中，一直处于切削热和摩擦力的作用下，会逐渐产生磨损并变钝。刀具的磨损可以分为初期磨损（I）、正常磨损（II）、急剧磨损（III）三个阶段，如图 13-15 所示。一般在正常磨损的后期、急剧磨损之前，应换刀重磨为好，否则会增加切削力和切削热，降低加工精度。对于可重磨的刀具，经过刃磨之后恢复锋利，可以继续使用。在经过使用$\to$磨钝$\to$刃磨若干个循环之后，刀具会无法使用而报废。这样，从最开始使用到完全报废，刀具实际切削时间的总和称为刀具的寿命，用 $t$ 表示。


\begin{figure}[H]
\centering
\includegraphics[page=49, viewport={374.0 23.3 503.0 108.3}, clip]{12-07572-001FigRevised.pdf}
\caption{刀具的磨损过程}
\end{figure}


刀具的磨损速度与刀具材料、刀具的几何形状、工件的材料、是否使用切削液等因素有关，可以选用耐热性好的材料或加大刀具前角，以减少刀具的磨损。刀具的磨损极限，是以刀具后面的磨损带宽度 $VB$ 来衡量的。但在实际生产中 $VB$ 不可能经常测量，一般根据刀具进行切削的时间来确定刀具是否磨钝。

2. 刀具耐用度

刀具耐用度是指刀具由开始切削至磨钝为止实际进行切削的时间，用 $T$ 表示。

凡影响刀具磨损的因素也必然影响刀具耐用度，主要因素有工件材料、刀具材料及几何角度、切削用量以及是否使用切削液等。切削用量对刀具耐用度的影响，以切削速度的影响为最大。一般切削速度越大，刀具耐用度越小；耐热性好的刀具材料，不易磨损，刀具耐用度大；适当加大刀具前角，使用切削液，都可减少刀具的磨损，提高刀具耐用度。

刀具耐用度与刀具寿命的关系为
\[ t = nT \]
式中：$t$——刀具寿命，min；
$n$——刀具刃磨次数；
$T$——刀具耐用度，min。

\section{加工质量和生产率}

\subsection{加工质量}

零件的加工质量直接影响着产品的使用性能和使用寿命，主要包括加工精度和表面质量两个方面。

1. 加工精度

零件的加工精度是指零件加工后的实际几何参数与理想几何参数的符合程度。这些几何参数包括零件加工表面的尺寸、形状、各表面的相互位置及表面粗糙度。若实际值越接近理想值，加工精度就越高。

加工精度在数值上一般通过加工误差值的大小来表示。在加工过程中，机床、夹具、刀具及工件组成的工艺系统中的各个环节的误差都会影响加工误差。因此，在加工同一批工件时，每个工件的误差值也是不相同的。公差则限制了误差的数值。当工件上某一参数的误差超过公差的范围时，则此工件为不合格件。给定的公差数值越小，加工精度越高，加工的难度也就越大。加工精度又分为尺寸精度、形状精度和位置精度。


(1)尺寸精度

    尺寸精度是指零件实际加工的尺寸与设计给定的尺寸相符合的程度，它由尺寸公差控制。国家标准（GB/T 1800.1—2020、GB/T 1804—2000）规定尺寸公差等级分为 20 级，分别用 IT01、IT0、IT1、IT2、$\cdots$、IT18 表示。其中，IT01 精度最高，公差值最小；IT18 精度最低，公差值最大。

(2)几何精度

    几何精度是指零件表面实际几何要素与理想几何要素相符合的程度，由几何公差控制。国家标准（GB/T 1182—2018）规定几何公差分为 19 项。几何公差分类及符号见表 13-3。


\begin{table}[H]
    \centering
    \caption{几何公差分类及符号}
    \begin{tabular}{|c|c|c|c|c|c|}
        \hline
        公差类型 & 几何特征 & 符号 & 公差类型 & 几何特征 & 符号 \\
        \hline
        \multirow{6}{*}{形状公差} & 直线度 & \includegraphics[page=52, viewport={197.0 281.3 221.5 298.3}, clip]{12-07572-001FigRevised.pdf} & \multirow{6}{*}{位置公差} & 位置公差 & \includegraphics[page=52, viewport={284.0 232.8 303.0 250.8}, clip]{12-07572-001FigRevised.pdf}\\
        \cline{2-3}\cline{5-6}
        & 平面度 &\includegraphics[page=52, viewport={238.5 281.3 264.5 303.8}, clip]{12-07572-001FigRevised.pdf} & & 同心度（用于中心线） & \includegraphics[page=52, viewport={321.5 231.3 345.5 252.8}, clip]{12-07572-001FigRevised.pdf}\\
        \cline{2-3}\cline{5-6}
        & 圆度 & \includegraphics[page=52, viewport={283.0 285.3 303.0 303.3}, clip]{12-07572-001FigRevised.pdf} & & 同轴度（用于轴线） & \includegraphics[page=52, viewport={281.0 283.8 303.0 304.3}, clip]{12-07572-001FigRevised.pdf}\\
        \cline{2-3}\cline{5-6}
        & 圆柱度 & \includegraphics[page=52, viewport={323.5 285.3 345.0 304.3}, clip]{12-07572-001FigRevised.pdf} & & 对称度 & \includegraphics[page=52, viewport={363.0 231.8 383.5 251.8}, clip]{12-07572-001FigRevised.pdf}\\
        \cline{2-3}\cline{5-6}
        & 线轮廓度 & \includegraphics[page=52, viewport={71.5 232.3 96.5 248.3}, clip]{12-07572-001FigRevised.pdf} & & 线轮廓度 &\includegraphics[page=52, viewport={68.5 230.3 98.0 250.3}, clip]{12-07572-001FigRevised.pdf}\\
        \cline{2-3}\cline{5-6}
        & 面轮廓度 & \includegraphics[page=52, viewport={116.0 231.8 139.0 247.3}, clip]{12-07572-001FigRevised.pdf} & & 面轮廓度 & \includegraphics[page=52, viewport={114.5 231.8 139.5 251.3}, clip]{12-07572-001FigRevised.pdf}\\
        \hline
        \multirow{5}{*}{方向公差} & 平行度 &\includegraphics[page=52, viewport={157.5 233.3 178.5 253.3}, clip]{12-07572-001FigRevised.pdf} & \multirow{2}{*}{跳动公差} & 圆跳动 & \includegraphics[page=52, viewport={401.5 231.8 425.5 251.8}, clip]{12-07572-001FigRevised.pdf}\\
        \cline{2-3}\cline{5-6}
        & 垂直度 &\includegraphics[page=52, viewport={200.0 231.3 219.5 250.8}, clip]{12-07572-001FigRevised.pdf} & & 全跳动 & \includegraphics[page=52, viewport={442.5 228.8 473.0 255.3}, clip]{12-07572-001FigRevised.pdf}\\
        \cline{2-3}
        & 倾斜度 & \includegraphics[page=52, viewport={242.0 232.3 263.5 249.3}, clip]{12-07572-001FigRevised.pdf} & & & \\
        \cline{2-3}
        & 线轮廓度 & \includegraphics[page=52, viewport={73.5 231.8 96.5 249.3}, clip]{12-07572-001FigRevised.pdf} & & & \\
        \cline{2-3}
        & 面轮廓度 & \includegraphics[page=52, viewport={116.0 231.3 138.0 248.8}, clip]{12-07572-001FigRevised.pdf} & & & \\
        \hline
    \end{tabular}\end{table}

由于在加工过程中有各种因素影响加工精度，所以同一加工方法在不同的条件下所能达到的精度是不同的，甚至采用同一种方法，如果多费一些工时，精心操作，细心调整，并选择合适的切削参数进行加工，也能使加工精度得到较大的提高。但这样做会降低生产率，增加生产成本，因而是不经济的。所以，通常所说的某加工方法所达到的精度，是指在正常操作情况下所能达到的精度，称为经济精度。

影响零件加工精度的因素很多，主要包括采用了近似加工方法或形状近似的刀具等而造成的加工原理误差，机床、刀具及夹具误差，工件装夹误差，工艺系统变形误差，工件内应力误差等。

2. 表面质量

零件的表面质量是指零件在加工后表面层的状况，包括表面粗糙度、表面层加工硬化的程度及表面残余应力的性质和大小。零件的表面质量对零件的耐磨、耐腐蚀、耐疲劳等性能以及零件的使用寿命都有很大的影响。因此，对于高速、重载荷下工作的重要零件，除限制表面粗糙度外，还要控制其表层加工硬化的程度和深度，以及表层残余应力的性质（拉应力还是压应力）和大小。而对于一般的零件，则主要规定其表面粗糙度的数值范围。

零件的表面粗糙度是指零件加工表面上微小峰谷的高低不平度。它是由于切削过程中的振动、刀刃或磨粒摩擦而留下的加工痕迹。它与零件的耐磨性、配合性质、抗腐蚀性等有密切关系，会影响机器的使用性能、寿命与制造成本。为了保证零件的使用性能，应限制表面粗糙度的范围。国家标准（GB/T 1031—2009）规定了表面粗糙度的评定参数及其数值，其中最为常用的是轮廓算术平均偏差 $Ra$，单位为 $\mu$m。$Ra$ 值越小，零件的表面越光滑。

影响零件表面粗糙度的因素主要有切削残留面积、积屑瘤及工艺系统振动等。

在加工过程中，如果零件未达到规定的加工精度及表面粗糙度的要求，则会加大装配难度，并会影响其使用性能和寿命。但若对加工精度和表面粗糙度的规定超过了实际需要，则没有必要，也是不经济的。因此，选择加工精度和表面粗糙度的原则应在保证产品能达到技术要求的前提下，选用较低的等级。

在一般情况下，零件表面的尺寸精度要求越高，其形状和位置精度也就要求越高，表面粗糙度值越小。各种加工方法所能达到的尺寸公差等级及表面粗糙度见表 13-4。

\begin{table}[H]
    \centering
    \caption{各种加工方法所能达到的尺寸公差等级和表面粗糙度}
    \small
    \begin{tabular}{|c|c|c|c|c|c|}
        \hline
        \makecell{表面\\要求} & 加工方法 & 尺寸公差等级 & \makecell{表面粗糙度\\$Ra$ 值/$\mu$m} & 表面特征 & 应用举例 \\
        \hline
        不加工 & & IT16 \textasciitilde IT14 & & 消除毛刺 & 铸、锻件 \\
        \hline
        \multirow{3}{*}{粗加工} & \multirow{3}{*}{\makecell{粗车、粗铣、\\钻、粗镗、粗刨}} & IT13 \textasciitilde IT10 & 80 \textasciitilde 40 & 显见刀纹 & 底板、垫块 \\
        \cline{3-6}
        & & IT10 & 40 \textasciitilde 20 & 可见刀纹 & 螺钉不接合面 \\
        \cline{3-6}
        & & IT10 \textasciitilde IT8 & 20 \textasciitilde 10 & 微见刀纹 & 螺母不接合面 \\
        \hline
        \multirow{3}{*}{\makecell{半精\\加工}} & \multirow{3}{*}{\makecell{半精车、精车、\\精铣、精刨、粗磨}} & IT10 \textasciitilde IT8 & 10 \textasciitilde 5 & 可见刀痕 & 轴套不接合面 \\
        \cline{3-6}
        & & IT8 \textasciitilde IT7 & 5 \textasciitilde 2.5 & 微见刀痕 & \makecell{要求较高的轴套不接\\合面} \\
        \cline{3-6}
        & & IT8 \textasciitilde IT7 & 2.5 \textasciitilde 1.25 & 不见刀痕 & 一般轴套接合面 \\
        \hline
        \multirow{3}{*}{精加工} & \multirow{3}{*}{\makecell{精车、宽刃精\\刨、高速精铣、\\磨、铰、刮}} & IT8 \textasciitilde IT6 & 1.25 \textasciitilde 0.32 & 可辨加工痕迹的方向 & 要求较高的接合面 \\
        \cline{3-6}
        & & IT7 \textasciitilde IT6 & 0.63 \textasciitilde 0.32 & 微辨痕迹的方向 & 凸轮轴颈、轴承内孔 \\
        \cline{3-6}
        & & IT7 \textasciitilde IT6 & 0.32 \textasciitilde 0.16 & 不辨加工痕迹的方向 & 高速轴轴颈 \\
        \hline
        \multirow{4}{*}{\makecell{光整\\加工}} & \multirow{4}{*}{\makecell{精细磨、研磨、\\镜面磨、超精\\加工}} & IT7 \textasciitilde IT5 & 0.16 \textasciitilde 0.03 & 暗光泽面 & 阀面 \\
        \cline{3-6}
        & & IT6 \textasciitilde IT5 & 0.08 \textasciitilde 0.04 & 亮光泽面 & 滚珠轴承 \\
        \cline{3-6}
        & & IT6 \textasciitilde IT5 & 0.04 \textasciitilde 0.02 & 镜状光泽面 & 量规 \\
        \cline{3-6}
        & & & 0.02 \textasciitilde 0.01 & 雾状光泽面 & 量规 \\
        \cline{3-6}
        & & & $\leqslant 0.01$ & 镜面 & 量块 \\
        \hline
    \end{tabular}\end{table}

\subsection{生产率}
切削加工生产率 $R_0$ 是指单位时间内生产合格零件的数量，即
$$R_0 = 1/t_w$$
式中：$R_0$——切削加工生产率；
$t_w$——生产一个合格零件所需的时间，min/件。

在机床上加工一个合格零件所用的时间包括三个部分，即
$$t_w = t_m + t_c + t_0$$
式中：$t_m$——基本工艺时间，即加工一个零件所需的总切削时间，min；
$t_c$——辅助时间，即为了维持切削加工所消耗在各种辅助操作上的时间，如装卸工件、调整机床、空移刀具、刃磨刀具、安装及检验工件等时间，min；
$t_0$——其他时间，即熟悉工艺文件、润滑和擦拭机床、清扫切屑及工间休息时间，min。

因此，切削加工生产率 $R_0$ 可表示为
$$R_0 = \frac{1}{t_m + t_c + t_0}$$

由上式可知，提高切削加工的生产率，实际上就是设法减少零件加工的基本工艺时间、辅助时间及其他时间。

采用自动化生产过程以及使用先进的工具，如机夹可转位式刀具、气动夹具等，可以大大减少辅助时间。改进管理、妥善安排和调度生产可以减少其他时间的消耗。至于缩短零件基本工艺时间，则与切削用量的选择有着密切的关系。

\section{切削条件的合理选择}
\subsection{刀具角度及耐用度的选择}

1. 刀具角度的选择

(1) 前角 $\gamma_o$ 的选择

前角对切削的难易程度有很大影响。增大前角能使刀具锋利，减小切削变形、切削力和切削热。但前角过大，切削刃和刀尖的强度下降，刀具导热体积减小，影响刀具使用寿命。

前角大小的选择与工件材料、刀具材料、加工要求有关。工件材料的强度、硬度低，前角应选得大些，反之应选得小些；刀具材料韧性好（如高速钢），前角可选得大些，反之应选得小些（如硬质合金）；精加工时前角可选得大些，粗加工时应选得小些。通常，硬质合金车刀的前角 $\gamma_o$ 在 $-5^\circ \sim 20^\circ$ 的范围内选取。

(2) 后角 $\alpha_o$ 的选择

适当的后角可减小刀具后面与工件之间的摩擦和减少刀具后面的磨损。但后角过大会削弱切削刃强度，减小导热体积。

粗加工和承受冲击载荷的刀具，为了使切削刃有足够的强度，应取较小的后角，一般后角 $\alpha_o$ 为 $5^\circ \sim 7^\circ$；精加工时，为保证加工表面质量，应取较大的后角，一般后角 $\alpha_o$ 为 $8^\circ \sim 12^\circ$；高速钢刀具的后角可比同类型的硬质合金刀具稍大一些。

(3) 主偏角 $\kappa_r$ 的选择

主偏角的大小影响切削条件和刀具寿命，如图 13-16 所示。在进给量和背吃刀量相同的情况下，减小主偏角可以使切削刃参加切削的长度增加，切屑变薄，因而使切削刃单位长度上的切削负荷减轻。同时，加强了刀尖，增大了散热面积，使切削条件得到改善，刀具寿命提高。

主偏角的大小还影响切削分力的大小，如图 13-17 所示。在切削力同样大小的情况下，减小主偏角会使背向分力 $F_y$ 增大。当加工刚性较差的工件时，为避免工件变形和振动，应选用较大的主偏角。车刀常用的主偏角 $\kappa_r$ 有 $45^\circ$、$60^\circ$、$75^\circ$、$90^\circ$ 几种。

减小主偏角还可以减小已加工表面残留面积的高度，以减小工件的表面粗糙度 $Ra$ 值。

(4) 副偏角 $\kappa'_r$ 的选择

适当的副偏角可减小副切削刃和副后面与工件已加工表面之间的摩擦，防止切削时产生振动。副偏角的大小还影响表面粗糙度值的大小，如图 13-18 所示。切削时，在背吃刀量、进给量和主偏角相同的情况下，减小副偏角可以使残留面积减小，表面粗糙度值降低。

\begin{figure}[H]
\centering
\includegraphics[page=50, viewport={79.5 647.3 255.5 734.3}, clip]{12-07572-001FigRevised.pdf}
\caption{主偏角对切削宽度和厚度的影响}
\end{figure}

\begin{figure}[H]
\centering
\includegraphics[page=50, viewport={272.5 647.3 462.0 726.8}, clip]{12-07572-001FigRevised.pdf}
\caption{主偏角对背向分力的影响}
\end{figure}

\begin{figure}[H]
\centering
\includegraphics[page=50, viewport={94.0 505.3 445.5 613.3}, clip]{12-07572-001FigRevised.pdf}
\caption{副偏角对残留面积的影响}
\end{figure}

副偏角的大小主要根据表面粗糙度的要求来选取，一般副偏角 $\kappa'_r$ 为 $5^\circ \sim 15^\circ$。粗加工取较大值，精加工取较小值。

(5) 刃倾角 $\lambda_s$ 的选择

刃倾角主要影响主切削刃的强度和切屑流出的方向。刃倾角为正值时，切削刃强度降低，切屑流向待加工表面；刃倾角为负值时，切削刃强度增高，切屑流向已加工表面；刃倾角等于零时，切屑流向垂直于主切削刃方向。如图 13-19 所示。

\begin{figure}[H]
\centering
\includegraphics[page=50, viewport={56.5 328.3 255.0 449.3}, clip]{12-07572-001FigRevised.pdf}
\caption{刃倾角及其对切屑流向的影响}
\end{figure}

粗加工时，为增强主切削刃强度和散热能力，刃倾角常取负值；精加工时，为了不使切屑划伤工件已加工表面，刃倾角常取正值或零值；粗车、粗镗时，取 $\lambda_s = -5^\circ \sim 0^\circ$；精车、精镗时，取 $\lambda_s = 0^\circ \sim 5^\circ$；加工时有冲击载荷，取 $\lambda_s = -5^\circ \sim -15^\circ$。

随着科学技术的发展，各种新型刀具材料不断涌现，因此不但要掌握一般刀具的选用规律，还应该掌握各种新型刀具的特点和选用原则。

2. 刀具耐用度的选择

刀具耐用度规定得是否合适，直接影响着加工效率及加工成本。通过刀具磨损实验得到的刀具寿命方程式为
$$v_c T^m = C_0$$
式中：$v_c$——切削速度，m/s；
$T$——刀具耐用度，min；
$m$——指数，表示切削速度和刀具耐磨度之间的影响程度；
$C_0$——系数，与刀具、工件材料和切削条件有关。

上式在双对数坐标系中是一簇斜率为 $m$，截距为 $C_0$ 的直线，如图 13-20 所示。由图 13-20 可知，切削速度对刀具耐用度的影响很大。耐热性越低的刀具材料（斜率越小），切削速度对刀具耐用度的影响越大。对高速钢刀具而言，当切削速度升高一倍时，刀具耐用度将为原来的 1/32。

由此可知，若从加工效率上考虑，当刀具耐用度规定过高，允许采用的切削速度就低，使生产效率降低；当刀具耐用度规定过低，允许采用的切削速度就高，虽然减少了切削时间，但却因刀具耐用度的下降而大大增加了换刀和磨刀的辅助时间，使生产效率降低。若从加工成本上考虑，当刀具耐用度规定过高，切削速度必须降低，因此加工时间加长，使机床费用和人工费用增加，成本提高；当刀具耐用度规定过低，虽然允许的切削速度提高，但刀具消耗及与磨刀有关的成本均提高，成本也提高。为了合理解决耐用度及速度之间的关系，提出了最高生产率刀具耐用度 $T_p$ 和最低成本刀具耐用度 $T_c$ 两个概念。

\begin{figure}[H]
\centering
\includegraphics[page=50, viewport={274.0 326.8 485.5 476.8}, clip]{12-07572-001FigRevised.pdf}
\caption{各种刀具的寿命曲线比较}
\end{figure}

经计算得
$$T_p = \frac{1-m}{m} t_{ct}$$
$$T_c = \frac{1-m}{m} \left( t_{ct} + \frac{C_t}{M} \right)$$
式中：$t_{ct}$——换刀一次所消耗的时间，min；
$C_t$——磨刀费用（包括刀具成本及折旧费），元；
$M$——工时费（包括操作工人的工资、开机费及均摊到机床上的管理费和其他杂费），元。

比较以上两式，最高生产率刀具耐用度 $T_p$ 比最低成本刀具耐用度 $T_c$ 要低一些。一般情况下，多采用最低成本刀具耐用度 $T_c$。只有当生产任务紧迫或生产中出现不平衡的薄弱环节时，才采用最高生产率刀具耐用度 $T_p$。

最低成本刀具耐用度 $T_c$ 已由切削用量手册提供了参考数据，可根据具体生产条件进行选择。对形状简单、刃磨和装卸容易的刀具，耐用度可定得低些；对形状复杂，制造、刃磨成本高，换刀时间长的刀具，耐用度应定得高些。通常，硬质合金车刀 $T=60 \sim 90$ min；高速钢钻头 $T=80 \sim 120$ min；齿轮滚刀 $T=200 \sim 300$ min。根据刀具耐用度 $T$ 的数值由切削用量手册可查得相应的切削速度。

对于装刀、换刀和调刀比较复杂的自动线、多刀机床、组合机床等使用的加工刀具，其耐用度应选得高些。对于机夹可转位式刀具，由于换刀时间短，为了充分发挥其切削性能，提高生产效率，刀具耐用度可选得低些，一般取 $T=15 \sim 30$ min。

\subsection{切削用量的选择}
选择合理的切削用量对提高切削效率，保证必要的刀具耐用度和经济性以及加工质量，都有重要的意义。为了确定切削用量的选择原则，首先要了解它们对切削加工的影响。

(1) 对加工质量的影响

切削用量三要素中，背吃刀量对切削力影响最大，其次是进给量。进给量增大会使残留面积的高度显著增大（图 13-21），表面更加粗糙。切削速度增大时，切削力减小，并可减小或避免积屑瘤，有利于加工质量和表面质量的提高。

\begin{figure}[H]
\centering
\includegraphics[page=50, viewport={138.0 202.8 400.0 302.3}, clip]{12-07572-001FigRevised.pdf}
\caption{进给量对残留面积的影响}
\end{figure}

(2) 对基本工艺时间的影响

以车外圆为例，基本工艺时间的计算可简化为
$$t_m = \frac{K}{v_c f a_p}$$
式中：$K$——与毛坯直径及加工余量、车刀行程长度、走刀次数有关的系数。

由上式可知，切削用量三要素对基本工艺时间 $t_m$ 的影响是相同的。

(3) 对刀具耐用度的影响

用实验的方法（硬质合金车刀车削中碳钢）可以求出刀具耐用度与切削用量之间关系的经验公式
$$T = \frac{C_T}{v_c^5 f^{2.25} a_p^{0.75}} (\text{当 } f > 0.75 \text{ mm/r 时})$$
式中：$C_T$——与刀具耐用度有关的常数。

由上式知，切削用量中，切削速度对刀具耐用度的影响最大，进给量的影响次之，背吃刀量的影响最小。即当提高切削速度时，刀具耐用度下降的速度比增大同样倍数的进给量或背吃刀量时要快得多。由于刀具耐用度迅速下降，则增加了换刀或磨刀的次数（增加了辅助时间 $t_c$），从而影响生产率的提高。

综合切削用量三要素对刀具耐用度、生产率和加工质量的影响，选择切削用量的原则是在保证加工质量，降低成本和提高生产率的前提下，使 $a_p$、$f$、$v_c$ 的乘积最大，即使基本工艺时间最短。

粗加工时，首先要尽可能提高金属切除率（减少基本工艺时间），同时还要保证规定的刀具耐用度。在机床功率足够时，应尽可能选取大的背吃刀量，除留下精加工的余量外，最好一次走刀切除全部粗加工的余量。其次，根据工艺系统的刚度，按工艺装备及技术条件选择大的进给量 $f$。最后，再根据刀具耐用度的要求，针对不同的刀具材料和工件材料，选择合适的切削速度 $v_c$，以保证在一定刀具耐用度条件下达到最高生产率。

精加工时，首先应保证零件的加工精度和表面质量，同时还要保证必要的刀具耐用度和生产率。精加工往往采用逐渐减小背吃刀量的方法来逐步提高加工精度。为了既可保证加工质量，又可提高生产率，精加工时常采用专门的精加工刀具并选用较小的背吃刀量和进给量，以及较高的切削速度将工件加工到最终要求的精度。

选择切削用量时，还应考虑加工状况、工件材料及刀具材料等。如在断续切削时有较大的冲击力，切削用量应小一些；因硬质合金刀具的耐热性比高速钢好，故使用硬质合金刀具时，可提高切削用量；当工艺系统刚性较差时，易引起振动，应适当降低切削用量；切削非铁金属时，因材料的强度、硬度较低，切削力较小，一般可选用较大的切削用量；加工韧性好的金属（如紫铜）时，因切屑易黏附在刀刃上，故切削用量不宜过大。

切削用量的选取方法有计算法和查表法。对于大量生产的场合，切削用量按公式计算决定。但在大多数情况下，是根据给定的条件按金属切削用量手册或根据实践经验及实验数据来合理选定切削用量的。

\subsection{切削液的选择}
合理选择和使用切削液，对降低切削温度、减小刀具磨损、抑制积屑瘤的生成、提高加工质量等均能起到很好的效果。切削液具有冷却、润滑、洗涤、排屑和防锈的作用。

生产中常用的切削液可分为水溶液、乳化液、切削油三类。水溶液（主要成分是水，并加入防锈剂等添加剂）的冷却、排屑性能好，润滑性能差，常在磨削中使用。乳化液（将乳化油用水稀释而成）具有良好的冷却和清洗性能，也有一定的润滑性能，适用于粗加工及磨削。切削油（主要是矿物油，少数采用动植物油或混合油）的润滑性能好，防锈效果好，但冷却和清洗效果差，常用于精加工工序。

通常应根据加工性质、工件材料和刀具材料等来选择合适的切削液。

粗加工时，为了降低切削温度，一般应选用冷却作用较好的水溶液或低浓度的乳化液；精加工时，为了提高工件表面质量和减少刀具磨损，一般应选用润滑作用较好的切削油或高浓度的乳化液。磨削时，一般采用冷却、排屑性能好的水溶液或乳化液。

加工一般钢材时，通常选用乳化液或硫化切削油。加工铜合金等非铁金属时，一般不宜采用含硫化油的切削液，以免腐蚀工件；加工铸铁、青铜、黄铜等脆性材料时，切削温度不高，为避免崩碎切屑进入机床运动部件之间，一般不使用切削液；在低速精加工（如宽刀精刨、精铰、攻螺纹）时，为了提高工件的表面质量，可选用煤油。

高速钢刀具的耐热性较差，为了提高刀具的耐用度，一般要根据加工性质和工件材料选用合适的切削液。硬质合金刀具由于耐热性和耐磨性都较好，一般不用切削液。

\subsection{工件材料的切削加工性}
1. 工件材料切削加工性的概念

工件材料的切削加工性是指在一定的切削条件下，工件材料被切削加工的难易程度。

切削加工性的衡量指标是多方面的，可以从刀具使用寿命、切削力、切削温度、加工表面质量以及断屑性能等方面来衡量。但衡量一种材料切削加工性的好坏，还要看具体的加工要求和切削条件。例如，纯铁切除余量很容易，但获得光洁的表面比较难。所以，从加工难易程度上衡量，纯铁的加工性好；但从质量角度评定，则认为其切削加工性不好。同样，不锈钢在普通机床上加工并不困难，但在自动机床上加工难以断屑。

在生产和实验研究中，往往只取某一项指标来反映材料切削加工性的某一侧面。最常用的指标是一定刀具耐用度下的切削速度 $v_T$ 和相对加工性 $K_r$。
$v_T$ 的含义是指当刀具耐用度为 $T$ 时，切削某种材料所允许的最大切削速度。$v_T$ 越高，表示材料的切削加工性越好。通常取 $T=60$ min，则 $v_T$ 写作 $v_{60}$。

切削加工性的概念具有相对性。所谓某种材料的切削加工性好与坏，是相对于另一种材料而言的。在判别材料的切削加工性时，一般以切削正火状态 45 钢的 $v_{60}$ 作为基准，写作 $(v_{60})_j$，而把其他各种材料的 $v_{60}$ 同它相比，比值 $K_r$ 称为相对加工性，即
$$K_r = \frac{v_{60}}{(v_{60})_j}$$
常用材料的相对加工性 $K_r$ 分为八级，见表 13-5。凡 $K_r > 1$ 的材料，其加工性比 45 钢好；$K_r < 1$ 的材料，其加工性比 45 钢差。$K_r$ 实际上也反映了不同材料对刀具磨损和刀具耐用度的影响。

\begin{table}[H]
\centering
\caption{材料切削加工性等级}
\label{tab:13-5}
\small\setlength{\tabcolsep}{3pt}
\begin{tabular}{|c|c|c|c|p{4.5cm}|}
\hline
加工性等级 & \multicolumn{2}{c|}{名称及种类} & 相对加工性 $K_r$ & 代表性材料 \\
\hline
1 & 很容易切削材料 & 一般非铁金属 & $>3.0$ & 5-5-5 铜铅合金、9-4 铝铜合金、铝镁合金 \\
\hline
2 & \multirow{2}{*}{容易切削材料} & 易切削钢 & $2.5 \sim 3.0$ & 15Cr(退火, $R_m=380 \sim 450$ MPa); 自动机钢($R_m=400 \sim 500$ MPa); 30 钢(正火, $R_m=450 \sim 560$ MPa) \\
\cline{1-1}\cline{3-4}
3 & & 较易切削钢 & $1.6 \sim 2.5$ & \\
\hline
\multirow{2}{*}{4} & \multirow{2}{*}{普通材料} & 一般钢与铸铁 & $1.0 \sim 1.6$ & \begin{tabular}[t]{p{4.2cm}}45 钢、灰铸铁; \\ 2Cr13(调质, $R_m=850$ MPa); \\ 85 钢($R_m=900$ MPa)\end{tabular} \\
\cline{1-1}\cline{3-4}
5 & & 稍难切削材料 & $0.65 \sim 1.0$ & \\
\hline
6 & \multirow{3}{*}{难切削材料} & 较难切削材料 & $0.5 \sim 0.65$ & 45Cr(调质, $R_m=1050$ MPa); 65Mn(调质, $R_m=950 \sim 1000$ MPa); 50CrV(调质)、06Cr18Ni11Ti、某些钛合金、铸造镍基高温合金 \\
\cline{1-1}\cline{3-4}
7 & & 难切削材料 & $0.15 \sim 0.5$ & \\
\cline{1-1}\cline{3-4}
8 & & 很难切削材料 & $<0.15$ & \\
\hline
\end{tabular}
\end{table}

2. 改善工件材料切削加工性的途径

材料的切削加工性对生产率和表面质量有很大影响。因此，在满足工件使用要求的前提下，应尽量选用加工性较好的材料。

影响工件材料切削加工性的因素很多，主要有工件材料的物理力学性能（如导热系数、强度、塑性、韧性、硬度等）、化学成分及显微组织等。在实际生产中，可采取一些措施来改善材料的切削加工性。生产中常用的措施主要有以下两方面：

(1) 调整材料的化学成分

材料的化学成分直接影响其力学性能。如高碳钢强度和硬度较高，切削加工性较差；低碳钢塑性和韧性较高，切削加工性也较差；中碳钢的强度、硬度、塑性和韧性都居于高碳钢和低碳钢之间，故切削加工性较好。

在不影响工件对材料使用性能要求的前提下，可在钢中加入适量的硫、铅等元素，使之成为“易切钢”，即容易断屑、切削力小、刀具使用寿命长、加工表面质量好。

(2) 热处理改善显微组织

化学成分相同的材料，当其显微组织不同时，力学性能就不一样，其切削加工性也不同。因此，可通过对不同材料进行不同的热处理来改善其切削加工性。例如，低碳钢经正火处理或冷拔处理，其硬度略有提高，而塑性下降；白口铸铁硬度高、材质脆，加工困难，可在 $910 \sim 950$ $^\circ$C 经 $10 \sim 20$ h 退火或正火，降低硬度，变为可锻铸铁。这些热处理措施都能改善切削加工性。

\section{金属切削机床的基础知识}
\subsection{机床的分类与型号}
机床是机械加工系统的主要组成部分。机床的类型、规格繁多，以适应各种加工的需要。对机床进行分类并编制型号，以便于区别、管理和选用机床。

机床的分类方法很多，若按其使用性能来分，则可分为通用机床、专门化机床和专用机床；若按其的精度来分，则可分为普通机床、精密机床和高精度机床；若按其的自动化程度来分，则可分为一般机床、半自动机床和自动机床；若按其的质量来分，则可分为一般机床、大型机床和重型机床。

按机床的加工性质和所用刀具进行分类是最基本的机床分类方法。按照国家制订的机床型号编制方法，机床分为 11 类，即车床、钻床、镗床、磨床、齿轮加工机床、螺纹加工机床、刨床插床、拉床、铣床、切断机床及其他机床。

机床的型号由基本部分和辅助部分组成，中间用“/”隔开，读作“之”，前者需统一管理，后者纳入型号与否由企业自定，型号构成如图 13-22 所示。详细内容可参阅国家标准 GB/T 15375—2008。

CW6140B 的含义为最大车削直径为 400 mm 的经过第二次重大设计改进的万能车床。

MG1432A 的含义为最大磨削直径为 320 mm 的经过第一次重大设计改进的高精度万能外圆磨床。

XK5030 的含义为工作台台面宽度为 300 mm 的数控立式升降台铣床。

\begin{figure}[H]
\centering
\includegraphics[page=51, viewport={7.0 498.8 357.5 751.3}, clip]{12-07572-001FigRevised.pdf}
\caption{机床型号的含义}
\end{figure}

注 (1) 有“()”的代号或数字，当无内容时，则不表示。若有内容则不带括号。

(2) 有“$\bigcirc$”符号的，为大写的汉语拼音字母。

(3) 有“$\triangle$”符号的，为阿拉伯数字。

(4) 有“$\otimes$”符号的，为大写的汉语拼音字母，或阿拉伯数字，或两者兼有之。

\subsection{机床的机械传动}
机床的传动方式按传动机构的结构特点可分为机械传动、液压传动、电气传动和气压传动等。在机床的各种传动方式中，应用最多的是机械传动和液压传动。

机械传动方式因实现回转运动的结构简单、工作可靠、传动比准确、维修方便等，故主要用在机

床的回转运动中,也多用于机床的直线运动中。液压传动方式因易于实现自动化和无级变速,动作快速性好,吸振性能好,能实现系统的过载保护并能自行润滑等,故常在机床上的一些装置中采用。如往复主体运动传动装置:磨床及龙门刨床的工作台的运动、牛头刨床及插床的滑枕运动等。机床上的夹紧装置、变速操纵装置、工件和刀具装卸装置、工件输送装置等,均可采用液压传动来实现。采用液压传动有利于简化机床结构,提高机床自动化的程度。本书主要介绍机床的机械传动。

1. 机床上常用的传动副及传动关系

机械传动中,常用的传动元件有带与带轮、齿轮、蜗杆蜗轮、齿轮齿条和丝杠螺母等。每一对传动元件称为传动副,各种传动副具有不同的传动特点。常用传动副及其传动特点见表13-6。

\begin{longtable}{|m{1cm}|m{6cm}|m{2.5cm}|m{3cm}|m{3cm}|}
\caption{常用传动副及其传动特点}\\
\hline
传动形式 & 外形图 & 符号图 & 传动比 & 优缺点 \\
\hline
带传动 &\includegraphics[page=52, viewport={152.5 50.3 242.0 178.8}, clip]{12-07572-001FigRevised.pdf}&\includegraphics[page=52, viewport={262.0 51.8 329.5 143.8}, clip]{12-07572-001FigRevised.pdf} & $i_{I-II} = \frac{n_{II}}{n_I} = \frac{d_1}{d_2}\varepsilon$ \newline $\varepsilon$—带的滑动系数, \newline 一般取0.98 & 优点:传动平稳,中心距变化范围大;结构简单,制造、维修方便;过载时带打滑,起到安全装置作用。 \newline 缺点:外廓尺寸大,传动比不准确。摩擦损失大,传动效率低 \\
\hline
齿轮传动 & \includegraphics[page=53, viewport={153.0 650.3 286.5 746.8}, clip]{12-07572-001FigRevised.pdf} &\includegraphics[page=53, viewport={11.0 650.3 59.0 726.8}, clip]{12-07572-001FigRevised.pdf}& $i_{I-II} = \frac{n_{II}}{n_I} = \frac{z_1}{z_2}$ & 优点:传动比准确恒定;结构紧凑,工作可靠;可传递较大的扭矩且传动效率高,使用寿命长。 \newline 缺点:制造复杂,精度不高时传动不稳定,有噪声 \\
\hline
蜗杆蜗轮传动 &\includegraphics[page=52, viewport={6.0 51.8 139.0 202.3}, clip]{12-07572-001FigRevised.pdf} & \includegraphics[page=53, viewport={67.0 650.3 141.0 721.3}, clip]{12-07572-001FigRevised.pdf} & $i_{I-II} = \frac{n_{II}}{n_I} = \frac{z_1}{z_2}$ & 优点:可获得较大的传动比,传动准确;结构紧凑,承载能力大;传动平稳,无噪声。 \newline 缺点:传动效率低,摩擦产生的热量大,需良好的润滑条件 \\
\hline
齿轮齿条传动 &\includegraphics[page=52, viewport={350.0 49.8 529.5 143.8}, clip]{12-07572-001FigRevised.pdf} & \includegraphics[page=53, viewport={295.0 650.8 362.5 719.3}, clip]{12-07572-001FigRevised.pdf} & $S = n\pi d$ \newline $= n\pi m z$ & 优点:传动效率较高,结构紧凑。 \newline 缺点:当制造精度不高时,传动不够平稳 \\
\hline
丝杠螺母传动 & \includegraphics[page=53, viewport={373.0 650.3 460.0 706.3}, clip]{12-07572-001FigRevised.pdf} & \includegraphics[page=53, viewport={474.0 648.8 534.0 701.8}, clip]{12-07572-001FigRevised.pdf} & $S = n P_h$ & 优点:传动平稳,无噪声,可以达到高的传动精度。 \newline 缺点:传动效率较低 \\
\hline
\end{longtable}
注:$d$—带轮直径;$z$—齿轮齿数;$m$—齿轮模数;$n$—转速;$P_h$—丝杠导程。

2. 机床上常见的传动机构

变速机构用来改变从动件的旋转速度或移动速度。机床中常用塔轮、滑动齿轮、离合器及摆动齿轮等来实现变速。但无论是哪一种变速机构,都是通过改变传动比的大小,在主动轴转速不变时,从动轴得到各种不同的转速。表13-7列出了常用的四种变速机构。

\begin{longtable}{|m{2cm}|m{4.5cm}|m{2.5cm}|m{3cm}|m{3cm}|}
\caption{常用的四种变速机构}\\
\hline
传动形式 & 外形图 & 符号图 & 传动比 & 特点 \\
\hline
塔轮变速机构 & \includegraphics[page=53, viewport={10.5 486.3 103.0 597.8}, clip]{12-07572-001FigRevised.pdf} & \includegraphics[page=53, viewport={120.0 486.8 175.5 586.3}, clip]{12-07572-001FigRevised.pdf} & $i_{I-II} = \frac{n_{II}}{n_I}$ \newline 所以: \newline $i_1 = \frac{d_1}{d_4}$ \newline $i_2 = \frac{d_2}{d_5}$ \newline $i_3 = \frac{d_3}{d_6}$ & 运转平稳,结构简单,但需要在停止转动时用手来推带换挡,使用不方便 \\
\hline
滑动齿轮变速机构 & \includegraphics[page=53, viewport={190.5 486.3 324.0 618.8}, clip]{12-07572-001FigRevised.pdf} & \includegraphics[page=53, viewport={335.0 487.3 395.5 556.3}, clip]{12-07572-001FigRevised.pdf} & $i_{I-II} = \frac{n_{II}}{n_I}$ \newline 所以: \newline $i_1 = \frac{z_1}{z_2}$ \newline $i_2 = \frac{z_3}{z_4}$ \newline $i_3 = \frac{z_5}{z_6}$ & 结构紧凑,传动效率高,但不能在运转中变速 \\
\hline
离合器变速机构 & \includegraphics[page=53, viewport={410.5 489.3 530.0 601.8}, clip]{12-07572-001FigRevised.pdf} &\includegraphics[page=53, viewport={35.0 348.3 115.0 421.8}, clip]{12-07572-001FigRevised.pdf} & $i_{I-II} = \frac{n_{II}}{n_I}$ \newline 所以: \newline $i_1 = \frac{z_1}{z_2}$ (离合器M左移) \newline $i_2 = \frac{z_3}{z_4}$ (离合器M右移) & 变速时齿轮不需移动,因此可以采用斜齿轮使传动平稳。如果采用摩擦离合器,便可在运转中变速 \\
\hline
摆动齿轮变速机构 & \includegraphics[page=53, viewport={127.5 345.3 252.0 456.8}, clip]{12-07572-001FigRevised.pdf} &\includegraphics[page=53, viewport={272.0 345.8 337.0 442.3}, clip]{12-07572-001FigRevised.pdf} & $i_{I-II} = \frac{n_{II}}{n_I}$ \newline 所以: \newline $i_1 = \frac{z_1}{z_6}$ \newline $i_2 = \frac{z_2}{z_6}$ \newline $i_3 = \frac{z_3}{z_6}$ \newline $i_4 = \frac{z_4}{z_6}$ & 其外廓尺寸更小,结构刚度低,故传递力不宜过大 \\
\hline
\end{longtable}
注:$d$—带轮直径;$z$—齿轮齿数;$n$—转速。

换向机构用来改变机床运动部件的运动方向,机床上广泛采用由圆柱齿轮和锥齿轮组成的换向机构。表13-8列出了常用的三种换向机构。

\begin{table}[H]
\centering
\caption{常用三种换向机构}
\label{tab:13-8}
\small
\begin{tabular}{|m{2cm}|m{6.5cm}|m{2.5cm}|m{3cm}|m{3cm}|}
\hline
机构型式 & 符号图 & 传动路线 & 优缺点 \\
\hline
三星齿轮 & \includegraphics[page=53, viewport={353.0 347.3 505.0 452.3}, clip]{12-07572-001FigRevised.pdf} & 正转: \newline (a) $n_I \frac{z_1}{z_3} \frac{z_3}{z_4} = n_{II}$ \newline 反转: \newline (b) $n_I \frac{z_1}{z_2} \frac{z_2}{z_3} \frac{z_3}{z_4} = n_{II}$ & 优点:结构简单、紧凑,制造方便。 \newline 缺点:结构刚性差,只能传递小功率 \\
\hline
中间齿轮 & \includegraphics[page=53, viewport={74.0 215.3 263.0 314.8}, clip]{12-07572-001FigRevised.pdf} &正转: \newline (a) $n_I \frac{z_1}{z'_2} \frac{z'_2}{z_2} = n_{II}$ \newline (b) $n_I \frac{z_1}{z'_2} \frac{z'_2}{z_2} = n_{II}$ (离合器M左移) \newline 反转: \newline (a) $n_I \frac{z_3}{z_4} = n_{II}$ \newline (b) $n_I \frac{z_3}{z_4} = n_{II}$ (离合器M右移) & 优点:可传递较大的扭矩,结构稳固可靠,可以快速反转。 \newline 缺点:结构较大,制造成本较高 \\
\hline
锥齿轮 & \includegraphics[page=53, viewport={283.0 213.3 469.5 303.8}, clip]{12-07572-001FigRevised.pdf} &正转: \newline (a) $n_I \frac{z_1}{z_3} = n_{II}$ \newline (b) $n_I \frac{z_1}{z_3} = n_{II}$ (离合器M左移) \newline 反转: \newline (a) $n_I \frac{z_2}{z_3} = n_{II}$ \newline (b) $n_I \frac{z_2}{z_3} = n_{II}$ (离合器M右移) & 优点:可以改变两垂直轴之间的旋转方向。 \newline 缺点:制造较难 \\
\hline
\end{tabular}
\begin{tablenotes}
	\item 注:$z$—齿轮齿数;$n$—转速。
\end{tablenotes}
\end{table}

3. 机床传动链及其传动比

如果将基本传动方法中某些传动副按传动轴依次组合起来,就构成一个传动系统,这种传动系统也称为传动链。

如图13-23所示,运动自轴I输入,转速为$n_I$,经直径分别为$D_1$、$D_2$的两个带轮传至轴II,经圆柱齿轮1、2传至轴III,再经圆柱齿轮3、4传至轴IV,最后经蜗杆5及蜗轮6传至轴V,并把运动输出。此传动链的传动路线可用下面的方法来表达
\[ I \to \frac{D_1}{D_2}\varepsilon \to II \to \frac{z_1}{z_2} \to III \to \frac{z_3}{z_4} \to IV \to \frac{z_5}{z_6} \to V \]

传动链的总传动比等于传动链中的所有传动副的传动比的乘积。传动链总传动比为
\[ i_{I-V} = i_1 i_2 i_3 i_4 = \frac{n_V}{n_I} = \frac{D_1}{D_2}\varepsilon \frac{z_1}{z_2} \frac{z_3}{z_4} \frac{z_5}{z_6} \]

归纳成一般情况,若某传动链从输入轴I到输出轴K间共有$m$个传动副组成,则其运动链的传动比计算式可写成
\[ i_{I-K} = \frac{n_K}{n_I} = i_1 i_2 i_3 \cdots i_m \]

若已知主动轴I的转速$n_I$,由上式可求出输出轴K的转速,即
\[ n_K = n_I i_{I-K} = n_I i_1 i_2 i_3 \cdots i_m \]

利用此式,可确定出传动系统中任意一轴的转速。如图13-23所示,轴V的转速为
\[ n_V = n_I i_{I-V} = n_I \frac{D_1}{D_2}\varepsilon \frac{z_1}{z_2} \frac{z_3}{z_4} \frac{z_5}{z_6} \]

轴III的转速为
\[ n_{III} = n_I i_{I-III} = n_I \frac{D_1}{D_2}\varepsilon \frac{z_1}{z_2} \]

\begin{figure}[H]
	\centering
	\includegraphics[page=51, viewport={373.5 499.8 529.5 634.8}, clip]{12-07572-001FigRevised.pdf}
	\caption{传动链图例}
\end{figure}


4. 卧式车床的传动系统

车床的种类很多,有卧式车床、转塔车床、立式车床、自动和半自动车床等,其中卧式车床的通用性好,应用最为广泛。图13-24为C6132车床的传动系统图。

\begin{figure}[H]
\centering
\includegraphics[page=51, viewport={44.0 228.3 488.5 471.3}, clip]{12-07572-001FigRevised.pdf}
\caption{C6132车床的传动系统图}
\end{figure}

卧式车床的传动系统由两条传动链组成。一条是由电动机经变速箱、带轮、主轴箱到主轴,称为主运动链,其任务是将电动机的运动传到主轴,并使其获得各种不同的转速,以满足不同工件直径、工件和刀具材料以及不同加工工序的需要。另一条是由主轴经挂轮箱、进给箱、溜板箱到刀具,称为进给运动链,其任务是使刀架带着刀具实现机动的纵向进给、横向进给或车螺纹,以满足不同车削加工的需要。

(1) 主运动转动

1) 传动路线

如图13-24所示,电动机带动变速箱内的轴I转动,移动轴I上的双联齿轮,可使轴II获得不同的两种转速。移动轴III上的三联齿轮,可使变速箱输出轴III获得不同的六种转速。运动再经带轮$\phi 176$将轴III的六种转速传给主轴箱上的带轮$\phi 200$。带轮$\phi 200$和齿轮27固定在轴套IV上,并空套在主轴VI上。轴套IV的运动又分成两路传给主轴VI。一路是图示位置,运动经过齿轮27/63和17/58传给主轴,使主轴获得六种低速。另一路是通过将带有内齿轮的联轴器$M_1$向左移动,与齿轮27相啮合,同时齿轮63、27向左脱开,使主轴直接获得较高的六种转速。主运动传动链的传动路线表达式如下:

\[
\left(
\begin{matrix}
	\text{电动机} \\
	n = 1\,440\ \text{r/min} \\
	P = 4\ \text{kW}
\end{matrix}
\right)
\to \text{I} \to
\begin{Bmatrix}
	\frac{33}{22} \\
	\frac{19}{34}
\end{Bmatrix}
\to \text{II} \to
\begin{Bmatrix}
	\frac{34}{32} \\
	\frac{28}{39} \\
	\frac{22}{45}
\end{Bmatrix}
\to \text{III} \to \frac{\phi 176}{\phi 200} \varepsilon \to \text{IV} \to
\begin{Bmatrix}
	\overleftarrow{M_1} \\
	\frac{27}{63}
\end{Bmatrix}
\to \text{V} \to \frac{17}{58} \to \text{主轴 VI}
\]

2) 传动计算

根据传动路线表达式,主轴可得到$2 \times 3 \times 2 = 12$种转速。

根据传动链传动比的计算方法,主轴工作转速计算式为
\[ n_{VI} = n_I i_{I-VI} \]
即
\[ n_{\text{主}} = n_{\text{电}} i_{\text{变}} i_{\text{带}} i_{\text{主}} \]
式中:$n_{\text{电}}$——电动机转速,r/min;
$i_{\text{变}}$——变速箱总传动比;
$i_{\text{带}}$——带传动比;
$i_{\text{主}}$——主轴箱输入轴至主轴的传动比。

按图示齿轮的啮合位置,主轴的工作转速为

\[ n_{\text{主}} = \left[ 1440 \times \left( \frac{33}{22} \times \frac{34}{32} \right) \times \left( \frac{176}{200} \times 0.98 \right) \times \left( \frac{27}{63} \times \frac{17}{58} \right) \right] \text{ r/min} \approx 249 \text{ r/min} \]

当改变变速箱中滑移齿轮位置及主轴箱中离合器$M_1$的位置时,用同样方法可求出其余11种转速。其中:
\[ n_{\text{主最高}} = 1980 \text{ r/min}, \]
\[ n_{\text{主最低}} = 45 \text{ r/min}。 \]

主轴的反转是靠电动机反转实现的。

(2) 进给运动传动

1) 传动路线

如图13-24所示,主轴的运动通过滑移齿轮变向机构(齿轮$\frac{55}{35} \times \frac{35}{55}$或$\frac{55}{55}$),再经挂轮箱内的齿轮$\frac{29}{58}$及挂轮(交换齿轮)$\frac{a}{b} \times \frac{c}{d}$将运动传入进给箱中的轴XI,轴XI的运动又通过齿轮$\frac{27}{24}, \frac{30}{48}, \frac{26}{52}, \frac{21}{24}, \frac{27}{36}$中的任意一对将运动传到轴XII,再通过增倍机构,将运动传到轴XIII上。改变轴XIII上的滑移齿轮39的位置,即可分别与丝杠上的齿轮39或光杠上齿轮39相啮合,从而将运动传给丝杠或光杠。

运动经轴XV传出,则经丝杠($P_{h\text{丝}}=6$)与固定在溜板箱上的开合螺母配合,即可将丝杠的旋转运动变成溜板箱及刀架的移动,完成螺纹的加工。

运动从XIV轴传出,则经光杠传给溜板箱内的蜗杆蜗轮并传至轴XVI。当合上离合器$M_{\text{左}}$时,运动经齿轮$\frac{24}{60} \times \frac{25}{55}$传至轴XVIII,并带动固定在该轴上的小齿轮14转动,与齿轮14相啮合的齿条带动溜板箱及刀架作纵向进给。当合上离合器$M_{\text{右}}$时,运动就经齿轮$\frac{38}{47} \times \frac{47}{13}$传至横向进给丝杠($P_{h\text{横}}=4$),通过固定在横溜板上的螺母,使中滑板及刀架作横向进给。

综上所述,进给运动传动链的传动路线表达式如下:

	\begin{multline*}
	\text{主轴 VI} \to
	\underbrace{
		\begin{Bmatrix}
			\frac{55}{55} \\
			\frac{55}{35} \times \frac{35}{55}
		\end{Bmatrix}
	}_{\text{(变向机构)}}
	\to \text{VIII} \to \frac{29}{58} \to \text{IX} \to \frac{a}{b} \times \frac{c}{d}_{\text{(挂轮)}} \to \text{XI} \to 
	% 第二行开始
	\underbrace{
		\begin{Bmatrix}
			\frac{27}{24} \\
			\frac{30}{48} \\
			\frac{26}{52} \\
			\frac{21}{24} \\
			\frac{27}{36}
	\end{Bmatrix}}_{\text{(滑动齿轮)}}
	\to \text{XII} \to
	\underbrace{
		\begin{Bmatrix}
			\frac{39}{39} \times \frac{52}{26} \\
			\frac{26}{52} \times \frac{52}{26} \\
			\frac{39}{39} \times \frac{26}{52} \\
			\frac{26}{52} \times \frac{26}{52}
		\end{Bmatrix}
	}_{\text{(增倍齿轮)}}
	\to \text{XIII} \to \\
	% 第三行开始
	\begin{Bmatrix}
		\begin{aligned}  % 关键点：[l] 代表左对齐
			&\frac{39}{39} \to \text{丝杠} \to \text{纵向进给车螺纹} \\
			&\frac{39}{39} \to \text{光杠} \to \frac{2}{45} \to \text{XVI} \to 
			\begin{Bmatrix}
				\begin{aligned}
					&\frac{24}{60} \to \text{XVII} \to M_{\text{左}} \uparrow \to \frac{25}{55} \to \text{XVIII} \to \text{齿轮齿条} \to \text{刀架纵向自动进给} \\
					&M_{\text{右}} \uparrow \to \frac{38}{47} \times \frac{47}{13} \to \text{横向进给丝杠} \to \text{刀架横向自动进给}
				\end{aligned}
			\end{Bmatrix}
		\end{aligned}
	\end{Bmatrix}
\end{multline*}

2) 纵向和横向进给传动计算

由传动路线表达式可看出,刀架机动的纵向和横向进给,是经主轴至进给箱中轴XIII的传动路线实现的。因此,利用进给箱中传动链的不同传动路线,可获得加工时所需要的不同的纵向和横向进给量。

进给量的计算,可根据前述齿轮齿条及丝杠螺母传动时齿条或螺母的移动距离计算式求得。当主轴转一周时,刀具纵向或横向移动的距离便为进给量。计算式为

\[ f_{\text{纵}} = 1 \underbrace{(n_{\text{主}}=1) i_{\text{换}} i_{\text{挂}} i_{\text{进}} i_{\text{溜}}}_{n_{\text{主}}=1\text{时,小齿轮14的转速}} \pi z m \]
\[ f_{\text{横}} = 1 \underbrace{(n_{\text{主}}=1) i_{\text{换}} i_{\text{挂}} i_{\text{进}} i_{\text{溜}}}_{n_{\text{主}}=1\text{时,横向进给丝杠的转速}} P_{h\text{横}} \]
式中:$f_{\text{纵}}$——刀具纵向进给量,mm/r;
$f_{\text{横}}$——刀具横向进给量,mm/r;
$i_{\text{换}}$——换向机构传动比;
$i_{\text{挂}}$——挂轮箱总传动比;
$i_{\text{进}}$——进给箱总传动比;
$i_{\text{溜}}$——溜板箱总传动比;
$z, m$——分别为小齿轮的齿数($z=14$)及模数($m=2$);
$P_{h\text{横}}$——横向进给丝杠的导程($P_{h\text{横}}=4$)。

按图示啮合位置,如挂轮使用车公制螺纹的四个齿轮60、65、65、45,则当离合器$M_{\text{左}}$及$M_{\text{右}}$分别接合时,其纵向、横向进给量分别为

\[ f_{\text{纵}} = \left[ 1 \times \left( \frac{55}{35} \times \frac{35}{55} \right) \times \left( \frac{29}{58} \times \frac{60}{65} \times \frac{65}{45} \right) \times \left( \frac{26}{52} \times \frac{39}{39} \times \frac{26}{52} \times \frac{39}{39} \right) \times \left( \frac{2}{45} \times \frac{24}{60} \times \frac{25}{55} \right) \times \pi \times 2 \times 14 \right] \text{ mm/r} \]
\[ = 0.118 \text{ mm/r} \]
\[ f_{\text{横}} = \left[ 1 \times \left( \frac{55}{35} \times \frac{35}{55} \right) \times \left( \frac{29}{58} \times \frac{60}{65} \times \frac{65}{45} \right) \times \left( \frac{26}{52} \times \frac{39}{39} \times \frac{26}{52} \times \frac{39}{39} \right) \times \left( \frac{2}{45} \times \frac{38}{47} \times \frac{47}{13} \right) \times 4 \right] \text{ mm/r} \]
\[ = 0.087 \text{ mm/r} \]

(3) 螺纹加工

从上述传动路线可知,进给箱中的滑动齿轮和增倍机构,可使进给箱本身的变速级数达$5 \times 4 = 20$种,加之有7组不同传动比的挂轮可根据需要任意配换,故可以车出各种不同螺距的螺纹。

\subsection{机床的自动化}

提高机床的生产率仅从提高材料的切除率、缩短基本工艺时间等方面入手远远不够,还必须在保证加工精度、降低成本的前提下,有效减少辅助时间,这就要求提高机床的自动化程度。

所谓自动化就是以机械动作代替人力操作,自动完成特定的作业。对于应用于大批量生产的机床,以及对于加工工件形状复杂又要求精度高的机床,提高自动化程度非常必要。自动化生产不仅可提高生产效率,减少操作人员,节省能耗,降低产品成本,还可减少工人的劳动强度以及工人技术水平对加工质量的影响,稳定产品的质量。

机床的自动化分为单一品种大批量生产的自动化和多品种、小批生产的自动化两大类。品种和批量的不同,所采用的自动化手段也不同。

产品单一批量较大时,可采用专用设备、专用流水线和自动线等刚性自动化措施来提高生产效率,一般有自动机床、半自动机床、组合机床及自动生产线等。但是,一旦产品改变,其设备则很难适应。

对于多品种、小批生产的自动化问题,数控机床则提供了广阔前景。数控机床的通用性强,生产准备时间短,生产效率高,适应性强。当加工对象改变时,除了更换刀具外,只需更换数控程序,便可自动加工出所要求的新零件,而不必对机床做任何调整。因此,数控机床特别适宜于产量小、品种多、产品更新频繁、生产周期要求短以及零件形状复杂的场合。

数控技术可以应用于车床、钻床、镗床、铣床、磨床等各类机床中。目前,数控机床已发展到更先进的“加工中心”,由传统的专用硬件逻辑数控系统(NC)发展到计算数字控制系统(CNC),并已形成柔性单元、柔性系统、自动化工厂等成熟产业应用模式,正在朝着智能化、网络化、敏捷制造、虚拟制造方向发展。

\section*{复习思考题}

 
13-1 试述切削加工的分类、特点及发展方向。

13-2 试述合理选用切削用量三要素的基本原则。

13-3 表13-9中几种切削加工方法的主运动和进给运动方式是转动还是移动?分别是由刀具还是由工

件实现的?将答案填入表13-9中。

\begin{table}[H]
\centering
\caption{题13-3表}
\label{tab:13-9}
\begin{tabular}{|m{2cm}|m{2cm}|m{2cm}|m{2cm}|m{2cm}|m{2cm}|m{2cm}|}
\hline
运动方法 & 车床车外圆 & 车床钻孔 & 车床镗孔 & 钻床钻孔 & 龙门刨床刨平面 & 镗床镗孔 \\
\hline
主运动形式 & & & & & & \\
\hline
进给运动形式 & & & & & & \\
\hline
\hline
运动方法 & 牛头刨床刨平面 & 铣床铣平面 & 插床插键槽 & 拉床拉键槽 & 平面磨床磨平面 & 外圆磨床磨外圆 \\
\hline
主运动形式 & & & & & & \\
\hline
进给运动形式 & & & & & & \\
\hline
\end{tabular}\end{table}

13-4 车削外圆时,已知工件转速 $n=320\text{ r/min}$,进给速度 $v_{\text{f}}=64\text{ mm/min}$,工件待加工表面直径 $d_{\text{w}}=100\text{ mm}$,工件已加工表面直径 $d_{\text{m}}=94\text{ mm}$,求切削速度 $v_{\text{c}}$、进给量 $f$、背吃刀量 $a_{\text{p}}$。

13-5 车削外圆时,工件转速 $n=360\text{ r/min}$,切削速度 $v_{\text{c}}=150\text{ m/min}$,测得此时电动机功率 $P_{\text{e}}=3\text{ kW}$,设机床传动效率 $\eta_{\text{m}}=0.8$,试求工件直径 $d_{\text{w}}$、主切削力 $F_{\text{z}}$。

13-6 请说明常用高速钢而不采用硬质合金制造拉刀和齿轮刀具这类形状较复杂刀具的原因。

13-7 刀具材料应具备哪些性能?常用刀具材料有哪些?各有何优缺点?。

13-8 简述切屑的种类及其对表面质量的影响。

13-9 积屑瘤的形成原因及对切削加工的影响是什么?

13-10 已知下列车刀的主要角度,试画出它们切削部分的示意图。
(1) 外圆车刀: $\gamma_{\text{o}}=10^{\circ}$, $\alpha_{\text{o}}=8^{\circ}$, $\kappa_{\text{r}}=60^{\circ}$, $\kappa'_{\text{r}}=10^{\circ}$, $\lambda_{\text{s}}=4^{\circ}$;
(2) 端面车刀: $\gamma_{\text{o}}=15^{\circ}$, $\alpha_{\text{o}}=10^{\circ}$, $\kappa_{\text{r}}=45^{\circ}$, $\kappa'_{\text{r}}=30^{\circ}$, $\lambda_{\text{s}}=-5^{\circ}$;
(3) 切断刀: $\gamma_{\text{o}}=10^{\circ}$, $\alpha_{\text{o}}=6^{\circ}$, $\kappa_{\text{r}}=90^{\circ}$, $\kappa'_{\text{r}}=2^{\circ}$, $\lambda_{\text{s}}=0^{\circ}$。

13-11 弯头车刀刀头的几何形状如图13-25所示,试分别说明车外圆、车端面(由外向内中心供给)时的前角 $\gamma_{\text{o}}$、主后角 $\alpha_{\text{o}}$、主偏角 $\kappa_{\text{r}}$ 和副偏角 $\kappa'_{\text{r}}$。

13-12 切断刀如图13-26所示,试分别标出切削平面、基面、主剖面及 $\gamma_{\text{o}}$、$\alpha_{\text{o}}$、$\kappa_{\text{r}}$、$\kappa'_{\text{r}}$。

\begin{figure}[H]
\centering
% Placeholder for Figure 13-25
\includegraphics[page=51, viewport={96.5 44.3 227.5 151.8}, clip]{12-07572-001FigRevised.pdf}
\caption{题13-11图}
\label{fig:13-25}
\end{figure}

\begin{figure}[H]
\centering
% Placeholder for Figure 13-26
\includegraphics[page=51, viewport={249.5 46.3 442.0 200.3}, clip]{12-07572-001FigRevised.pdf}
\caption{题13-12图}
\label{fig:13-26}
\end{figure}

13-13 切削力是如何产生的?影响切削力的因素主要有哪些?

13-14 切削热是如何产生,又如何传散的?对切削加工有何影响?

13-15 从提高刀具耐用度的角度出发,粗加工时选择切削用量的顺序是什么?

13-16 刀具的正常磨损过程可分为几个阶段?各阶段的特点是什么?刀具使用时磨损应限制在哪一阶段?

13-17 如何评价材料切削加工性的好坏?改善材料切削加工性的途径有哪些?

13-18 单一产品大批量生产时,通常采用的自动化措施有哪些?

13-19 为什么采用数控机床对多品种、小批生产实现自动化非常有利?

13-20 如图13-27所示传动系统中的传动结构中,当轴I的转速 $n_{\text{I}}=100\text{ r/min}$ 时,螺母的移动速度 $v_{\text{f}}$ 是多少?

\begin{figure}[H]
\centering
% Placeholder for Figure 13-27
\includegraphics[page=52, viewport={131.0 566.8 396.0 748.3}, clip]{12-07572-001FigRevised.pdf}
\caption{题13-20图}
\label{fig:13-27}
\end{figure}

13-21 图13-28所示为某车床主轴的传动图,试:
(1) 写出主传动链;
(2) 确定主轴V转速的级数;
(3) 计算主轴最大转速。

\begin{figure}[H]
\centering
% Placeholder for Figure 13-28
\includegraphics[page=52, viewport={111.0 335.8 424.5 537.3}, clip]{12-07572-001FigRevised.pdf}
\caption{题13-21图}
\label{fig:13-28}
\end{figure}


\chapter{零件表面的加工方法}

\section{切削加工方法}

\subsection{车削加工}

车削加工是指在车床上利用车刀进行切削加工的一种方法。常见的车床有卧式车床、立式车床、转塔车床、自动车床、数控车床等。车削加工时,工件的旋转是主运动,刀具做直线进给运动。因此,车削加工主要用来进行各种回转面的加工,如外圆面、内圆面、锥面、螺纹和滚花面等。在卧式车床上可以完成的主要工作见表14-1。

\begin{table}[H]
\centering
\caption{卧式车床的主要工作}
\label{tab:14-1}
\begin{tabular}{|c|c|c|c|}
\hline
钻中心孔 & 钻孔 & 铰孔 & 攻螺纹 \\
\hline
\includegraphics[page=62, viewport={233.5 121.8 308.0 184.8}, clip]{12-07572-001FigRevised.pdf}
&\includegraphics[page=62, viewport={321.0 120.3 412.0 192.3}, clip]{12-07572-001FigRevised.pdf} &\includegraphics[page=62, viewport={430.5 118.8 535.5 197.3}, clip]{12-07572-001FigRevised.pdf} & \includegraphics[page=62, viewport={58.0 19.3 160.0 91.8}, clip]{12-07572-001FigRevised.pdf}\\
\hline
车外圆 & 镗孔 & 车端面 & 切断 \\
\hline
\includegraphics[page=62, viewport={181.5 18.3 278.0 89.8}, clip]{12-07572-001FigRevised.pdf}
& \includegraphics[page=62, viewport={292.0 22.3 362.0 72.8}, clip]{12-07572-001FigRevised.pdf}&\includegraphics[page=62, viewport={387.5 19.8 491.0 89.8}, clip]{12-07572-001FigRevised.pdf} &\includegraphics[page=63, viewport={22.0 659.8 120.0 737.8}, clip]{12-07572-001FigRevised.pdf} \\
\hline
车成形面 & 车锥面 & 滚花 & 车螺纹 \\
\hline
\includegraphics[page=63, viewport={134.5 664.3 195.5 717.3}, clip]{12-07572-001FigRevised.pdf}
&\includegraphics[page=63, viewport={225.0 666.3 297.5 721.3}, clip]{12-07572-001FigRevised.pdf} &\includegraphics[page=63, viewport={321.0 664.8 393.5 721.8}, clip]{12-07572-001FigRevised.pdf} &\includegraphics[page=63, viewport={414.5 661.8 524.5 725.3}, clip]{12-07572-001FigRevised.pdf} \\
\hline
\end{tabular}
\end{table}

不同的加工表面,需要用不同的车刀进行加工。常见的车刀有切断刀、左偏刀、右偏刀、外圆车刀、内孔车刀、端面车刀、螺纹车刀等。

1. 车削的工艺特点

(1) 易于保证各加工表面之间的位置精度 车削加工时,工件以主轴的回转中心或两顶尖的中心连线为轴线作回转运动,各个表面具有同一回转轴线,其相对位置精度易于保证。

(2) 切削过程平稳 与刨削和铣削相比,车削过程是连续的,没有大的冲击和振动。若切削用量一定时,其切削面积和切削力基本上不发生变化。因此,车削加工可以采用较大的切削用量进行高速切削或强力切削,提高生产效率。

(3) 适用于有色金属零件的精加工 由于有色金属材料本身的硬度低,塑性、韧性较高,如果采用砂轮磨削,从工件上脱离的磨屑容易堵塞砂轮,影响加工过程,难以得到光洁的加工表面。所以,有色金属零件的表面粗糙度 $Ra$ 要求较小时,不宜采用磨削,可以利用车削或精细车、铣刨来进行加工。

(4) 刀具简单 车刀是最简单的一种刀具,其制造、刃磨、安装方便,使工艺过程简单。

车削加工的尺寸公差等级为 IT8$\sim$IT7,表面粗糙度 $Ra$ 值为 $1.6\sim0.8\text{ \textmu m}$。

在车削细长轴时,应根据零件要求采用跟刀架或中心架,以提高工件刚度;同时车刀采用 $90^{\circ}$ 主偏角,以减少背向分力 $F_{\text{y}}$,提高零件的加工精度。

车削塑性材料时,往往需要加切削液;而加工脆性材料时,由于断屑容易,一般不采用切削液。

2. 车削的应用

(1) 车锥面:采用车床加工锥面的方法通常有宽刀法、小刀架转位法、偏移尾架法、靠模法四种(图14-1)。宽刀法:利用主刀刃与工件轴线间的夹角等于工件锥面斜角 $\alpha$ 的车刀径向进给实现锥面加工;小刀架转位法:将小刀架绕转盘转一个被切锥面的斜角 $\alpha$,通过转动小刀架手柄,使车刀沿工件母线移动,以实现锥面加工;偏移尾架法:将尾架顶尖偏移一个距离 $s$,使工件锥面母线平行于车刀纵向进给方向,以实现锥面加工;靠模法:将靠模板绕中心轴转动,调整成所需的锥面斜角 $\alpha$,沿靠模板移动的滑块与横向溜板连接,当纵向进给时,车刀平行于靠模板移动,以加工出锥角 $\alpha$ 的圆锥面。

\begin{figure}[H]
\centering
% Placeholder for Figure 14-1
\includegraphics[page=54, viewport={82.5 407.8 456.0 743.8}, clip]{12-07572-001FigRevised.pdf}
\caption{圆锥面的车削方法}
\label{fig:14-1}
\end{figure}

(2) 车成形面:对于具有回转形面的零件,根据零件的精度要求及生产批量的不同,可采用图14-2所示的车削方法。双手同时操作法:采用双手同时作纵向和横向进给运动,其合成运动使车刀的运动轨迹与工件母线相同;成形车刀法:利用刀具刃口与工件表面母线形状相吻合的成形车刀作横向进给运动,实现成形面加工。靠模法车成形面与靠模法车圆锥面的原理相同,靠模形状与工件表面母线形状相同。

\begin{figure}[H]
\centering
% Placeholder for Figure 14-2
\includegraphics[page=54, viewport={45.5 255.3 495.5 375.8}, clip]{12-07572-001FigRevised.pdf}
\caption{成形面的车削方法}
\label{fig:14-2}
\end{figure}

(3) 车螺纹:车螺纹时,螺纹的截面形状由车刀保证,车刀的形状必须与螺纹截面相吻合。螺纹的螺距由车床传动系统保证,为了得到准确的螺距,必须保证工件转一圈,刀具准确地移动一个导程。除了车螺纹外,常用的螺纹加工方法还有:攻螺纹,套螺纹,铣、磨和滚压螺纹等,它们各有不同的特点,应根据零件的结构、技术要求、加工批量、尺寸等因素来选择其加工方法。

\subsection{钻削与镗削加工}

1. 钻削加工

钻削是孔加工的一种基本方法。钻孔一般在钻床上进行,也可以在车床、镗床或铣床上进行。常用的钻床有台式钻床、立式钻床和摇臂钻床。在钻床上钻孔时,钻头(刀具)的旋转运动为主运动,同时钻头沿工件孔的轴心线方向的移动为进给运动。在钻床上可以实现钻孔、扩孔、铰孔、攻螺纹、锪孔、锪平面等加工。

(1) 钻孔

钻孔是利用钻头在实体材料上加工孔的一种方法。常用的钻头有扁钻、麻花钻、深孔钻等多种,其中以麻花钻应用最为普遍。麻花钻由夹持部分和工作部分组成,工作部分又由导向部分和切削部分组成。

麻花钻的导向部分(图14-3)有两条对称的螺旋槽,用来形成切削刃及前角,并起排屑和输送切削液的作用。在麻花钻工作部分外螺旋面上做出两条窄的棱带,其外径略带倒锥,前大后小,起减小摩擦面积并保持钻孔方向的作用。

\begin{figure}[H]
\centering
% Placeholder for Figure 14-3
\includegraphics[page=54, viewport={53.5 89.8 281.0 163.3}, clip]{12-07572-001FigRevised.pdf}
\caption{麻花钻的结构}
\label{fig:14-3}
\end{figure}

麻花钻的切削部分(图14-4)有两条主切削刃、两条副切削刃和一条横刃。切屑流过的两个螺旋槽表面为刀具前面,与工件切削表面(即孔底)相对的钻头顶端两曲面为主后面,与工件已加工表面(即孔壁)相对的两条棱带为副后面。前面与主刀面的交线为主切削刃,前面与副后面的交线为副切削刃,两个主后面的交线为横刃。

钻孔与车削外圆相比,工作条件要恶劣许多。一方面,麻花钻头在切削时其工作部分大多处于加工表面的包围之中,受力大,散热困难;另一方面,由于麻花钻的结构特点,其存在刚度和强度不足、容屑和排屑能力较差、导向和冷却润滑困难等诸多问题。钻孔具有以下工艺特点。

1) 钻头容易引偏:由于钻头存在横刃,钻孔时轴向抗力大;横刃较长又有较大负前角,使钻头很难定心;钻头比较细长,加之其具有两条宽而深的容屑槽,使得其刚性差;钻头仅有两条很窄的螺旋棱带与孔壁接触,使得其导向性较差。因此,钻头在开始切削时就容易引偏,切入以后易产生弯曲变形,致使钻头偏离原轴线。

在钻床上钻孔与在车床上钻孔,钻头偏斜对孔加工精度的影响不同。当钻头引偏时,前者孔的轴线也发生偏斜,但孔径无显著变化(图14-5a);后者孔的轴线无明显偏斜,但会引起孔径变化,常使孔出现锥形或腰鼓形等缺陷(图14-5b)。因此,钻小孔或深孔时应尽可能在车床上进行,以减小孔轴线的偏斜。

\begin{figure}[H]
\centering
% Placeholder for Figure 14-4
\includegraphics[page=54, viewport={303.5 89.8 483.5 223.8}, clip]{12-07572-001FigRevised.pdf}
\caption{麻花钻切削部分的组成}
\label{fig:14-4}
\end{figure}

\begin{figure}[H]
\centering
% Placeholder for Figure 14-5
\includegraphics[page=55, viewport={54.5 617.8 243.0 745.3}, clip]{12-07572-001FigRevised.pdf}
\caption{钻头引偏}
\label{fig:14-5}
\end{figure}

钻头引偏会导致加工后的孔出现孔轴线歪斜、孔径扩大和孔失圆等现象。在实际生产中常采用措施有:预钻锥形定心坑、采用钻套导向、对称修磨两个半锋角和两条主切削刃等,以进行钻前定位,钻时导向和径向力平衡,从而减小钻头引偏。

2) 排屑困难:钻孔时,由于切屑较宽,容屑尺寸又受到限制,因此在排屑过程中,往往与孔壁产生很大的摩擦和挤压,拉毛和刮伤已加工表面,从而大大降低孔壁质量。
为便于排屑,生产中常对麻花钻进行修磨(图14-6a),使横刃变短,横刃的前角增大,从而减少横刃的不利影响;开磨分屑槽(图14-6b),在加工塑性材料时,能使较宽的切屑分成几条,以便顺利排屑。

\begin{figure}[H]
\centering
% Placeholder for Figure 14-6
\includegraphics[page=55, viewport={265.0 613.8 483.0 720.3}, clip]{12-07572-001FigRevised.pdf}
\caption{麻花钻的修磨}
\label{fig:14-6}
\end{figure}

3) 切削热不易传散:由于钻削是一种半封闭式的切削,切削时产生大量的热量,且高温切屑不能及时排出,切削液难以注入切削区,切屑、刀具与工件之间摩擦又很大,因此切削温度较高,致使刀具磨损加剧,从而限制了钻削用量和生产效率的提高。
钻孔的应用非常广泛,钻床上的主要工作见表14-2。钻孔加工的尺寸公差等级一般为 IT12$\sim$IT11,表面粗糙度 $Ra$ 值为 $25\sim12.5\text{ \textmu m}$,主要用于粗加工。对于精度高、表面粗糙度值小的中、小直径孔,在钻削之后,经常采用扩孔和铰孔来进行半精加工和精加工。

\begin{table}[H]
\centering
\caption{钻床上的主要工作}
\label{tab:14-2}
\begin{tabular}{|c|c|c|c|}
\hline
钻孔 & 扩孔 & 铰孔 & 攻螺纹 \\
\hline
\includegraphics[page=63, viewport={117.5 558.3 190.0 629.8}, clip]{12-07572-001FigRevised.pdf}
&\includegraphics[page=63, viewport={202.0 556.3 265.5 631.8}, clip]{12-07572-001FigRevised.pdf} &\includegraphics[page=63, viewport={286.5 557.8 350.5 631.8}, clip]{12-07572-001FigRevised.pdf} &\includegraphics[page=63, viewport={363.5 553.8 425.0 632.3}, clip]{12-07572-001FigRevised.pdf} \\
\hline
锪锥孔 & 锪柱孔 & 反锪鱼眼坑 & 锪凸台 \\
\hline
\includegraphics[page=63, viewport={100.0 443.3 171.0 517.8}, clip]{12-07572-001FigRevised.pdf}
&\includegraphics[page=63, viewport={189.0 443.3 263.0 518.8}, clip]{12-07572-001FigRevised.pdf} &\includegraphics[page=63, viewport={284.0 440.3 353.0 520.8}, clip]{12-07572-001FigRevised.pdf} &\includegraphics[page=63, viewport={372.5 444.8 437.5 525.8}, clip]{12-07572-001FigRevised.pdf} \\
\hline
\end{tabular}
\end{table}

(2) 扩孔

扩孔是利用扩孔钻对已经钻出、铸造出或锻造出的孔进行扩大的加工。从图14-7可知,扩孔钻的容屑槽小而浅,使刀具的刚性提高;刀尖部位避免了横刃和由横刃引起的不良影响,改善了切削条件;扩孔钻的刀齿多(3$\sim$4个),棱带增多,导向作用好。因此,扩孔的加工质量比钻孔好,可以达到的尺寸公差等级为 IT10$\sim$IT9,表面粗糙度 $Ra$ 值为 $6.3\sim3.2\text{ \textmu m}$。扩孔可以作为精度要求不高的孔的最终加工工序。

(3) 铰孔

铰孔是采用铰刀对未淬硬的孔进行的一种精加工方法。铰刀的结构和切削条件比扩孔更为优越。如图14-8所示,铰刀的刀齿更多,刚性更好;铰孔时的切削速度低,铰孔余量小,排屑和冷却较好,产生的切削热较少,工件的受力和变形较小;铰刀具有修光部分,可以校准孔径、修光孔壁。因此,铰孔的加工质量更高,可以达到的尺寸公差等级为 IT8$\sim$IT7,表面粗糙度 $Ra$ 值为 $3.2\sim0.4\text{ \textmu m}$。铰孔主要用于中、小尺寸孔的精加工。

\begin{figure}[H]
\centering
% Placeholder for Figure 14-7
\includegraphics[page=55, viewport={79.5 400.3 261.0 566.8}, clip]{12-07572-001FigRevised.pdf}
\caption{扩孔钻的结构}
\label{fig:14-7}
\end{figure}

\begin{figure}[H]
\centering
% Placeholder for Figure 14-8
\includegraphics[page=55, viewport={287.5 397.3 460.5 582.8}, clip]{12-07572-001FigRevised.pdf}
\caption{铰刀及其刃部结构}
\label{fig:14-8}
\end{figure}

对于中、小尺寸孔的单件生产甚至大批量生产,钻$\to$扩$\to$铰是常用的加工工艺路线,一般多是采用市场上容易购买到的标准刀具进行加工。

2. 镗削加工

镗削的加工范围广泛,可以加工较大直径的内孔表面、内成形面、孔内环槽等,还可以加工平面、端面、螺纹等。镗孔是利用镗刀对已钻出、铸出或锻出的孔进行加工的工艺方法。镗孔可以在多种机床上进行,例如回转体零件上的孔大多在车床上加工。镗床上镗孔主要用于箱体、机架等大型和复杂零件上的孔或孔系的加工。

镗床按结构形式可分为立式镗床、卧式铣镗床、坐标镗床、专门化镗床等,应用最广泛的为卧式铣镗床。镗削的主运动是主轴的旋转,进给运动有主轴的轴向移动,主轴箱沿立柱的竖直移动,刀架沿平旋盘的径向移动,工作台的纵向、横向移动和圆周的转动等。

镗刀有单刃镗刀和多刃镗刀,它们的结构和工艺特点如下。

单刃镗刀(如图14-9)的结构与车刀类似,其结构简单,适用性广,可用于粗加工、半精加工及精加工。通过调整刀体的悬臂长度可以加工不同直径的孔。单刃镗刀镗孔,还可以校正原有孔的轴线歪斜或位置偏差,而扩孔和铰孔一般难以做到。单刃镗刀只有一个主切削刃参加工作,而且为了减少振动和变形,一般采用较小的切削用量,所以生产效率较低,多用于单件、小批生产。

\begin{figure}[H]
\centering
% Placeholder for Figure 14-9
\includegraphics[page=55, viewport={162.0 248.8 373.5 363.8}, clip]{12-07572-001FigRevised.pdf}
\caption{单刃镗刀}
\label{fig:14-9}
\end{figure}

多刃镗刀中最常用的是可调浮动镗刀,如图14-10所示。加工前,先根据所加工的孔径调节刀齿的径向尺寸。镗孔时,镗刀片在刀杆的长方形孔中并不紧固,在孔径方向可以自由浮动,依靠两个对称切削刃产生的径向切削力自动平衡。多刃镗刀镗孔的加工质量较高,表面粗糙度$Ra$值较小,生产率较高,但是不能校正原有孔的轴线歪斜或位置偏差,而且刀具复杂,成本比单刃镗刀高。由于以上特点,可调浮动镗刀主要用于批量生产和精加工箱体、机架等大型零件上的孔和孔系。

\begin{figure}[H]
\centering
\includegraphics[page=55, viewport={104.5 92.8 436.0 216.3}, clip]{12-07572-001FigRevised.pdf}
\caption{用浮动镗刀在铣镗床上镗孔的方法}
\label{fig:14-10}
\end{figure}

镗床具有以下工艺特点。

(1)镗床是加工机座、箱体、机架等外形复杂的大型零件的主要设备

在一些箱体上往往有一系列孔径较大、精度较高的孔,这些孔在一般机床上加工很困难,但在镗床上加工却很容易,并可方便地保证孔与孔之间、孔与基准之间的位置和尺寸精度要求。

(2)加工范围广泛

镗床是一种通用性强、功能多的通用机床。镗床既可以加工单个孔,又可以加工孔系;既可以加工小直径孔,又可以加工大直径孔;既可以加工通孔,又可以加工台阶孔及内环形槽。除此之外,镗床还可以进行部分铣削和车削加工。

(3)能获得较高的加工精度和较小的表面粗糙度值

普通镗床镗孔的尺寸公差等级可达IT7~IT6,表面粗糙度$Ra$值可达$0.8 \sim 0.2 \text{ \textmu m}$。如果采用金刚石镗刀或坐标镗床进行镗削加工,可以获得更高的加工精度和更小的表面粗糙度值。

(4)生产率较低

机床和刀具调整复杂,操作技术要求较高,特别是在单件、小批生产中。在大批量生产中常采用镗模,以提高生产率。

\subsection{刨削、插削和拉削加工}
1.刨削加工

刨削是在刨床上利用刨刀进行切削的工艺方法。它主要用于水平面、竖直面、斜面、直槽、V形槽及成形面等的加工,如图14-11所示。

常见的刨床类型有牛头刨床、龙门刨床和插床等。在牛头刨床上加工时,滑枕带动刨刀的往复直线运动为主运动,刨刀或工件随工作台作间歇进给运动。其最大刨削长度一般不超过$1000 \text{ mm}$,适于加工中、小型工件;在龙门刨床上加工时,工件随工作台的往复直线运动为主运动,刀架沿横梁或立柱作间歇的进给运动。由于刚性好,龙门刨床主要用来加工大型工件或同时加工多个中、小型工件,其加工精度和生产率均比牛头刨床高。

\begin{figure}[H]
\centering
\includegraphics[page=56, viewport={90.5 541.3 451.5 745.3}, clip]{12-07572-001FigRevised.pdf}
\caption{刨削的主要应用}
\label{fig:14-11}
\end{figure}

刨削具有以下工艺特点。

(1)机床与刀具简单,通用性好

刨床的结构简单,易于操作,所用的单刃刨刀与车刀基本相同,形状简单,制造、刃磨和安装都较方便。

(2)刨削精度低

刨削主运动为往复直线运动,刀具切入、切出时有冲击振动,影响了加工表面质量。精刨时加工的尺寸公差等级一般为IT8~IT7,表面粗糙度$Ra$值为$2.3 \sim 1.6 \text{ \textmu m}$。

(3)刨削的生产率较低

刨削时有空行程,且每一往复行程伴有两次冲击,从而限制了切削速度的提高。

刨削多用于单件、小批生产及修配工作中。

2.插削加工

插削是利用插刀在插床上进行加工的方法。插床可看作“立式牛头刨床”,滑枕带动插刀的上下直线往复运动为主运动,工件装夹在工作台上,工作台可以实现纵向、横向和圆周进给运动。

插削主要用在单件、小批生产中插削某些内表面,如方孔、长方孔、各种多边形孔及孔内键槽等,也可以加工某些零件外表面。

3.拉削加工

拉削是在拉床上用拉刀加工工件的加工方法。拉削是一种高生产率和高精度的加工方法。拉削时,拉刀的直线移动为主运动,进给运动是由拉刀的后一个刀齿较前一个刀齿递增的一个齿升量来实现的。

拉刀是一种多齿刀具,一把拉刀只能加工一种形状和尺寸规格的表面。各种拉刀的形状、尺寸虽然不同,但它们的组成部分大体一致。图14-12所示为常见拉刀的结构及组成。

\begin{figure}[H]
\centering
\includegraphics[page=56, viewport={100.5 182.8 441.5 515.8}, clip]{12-07572-001FigRevised.pdf}
\caption{常见拉刀的结构及组成}
\label{fig:14-12}
\end{figure}

拉削具有以下工艺特点。

(1)生产率高

拉刀是多刃刀具,一次行程可完成粗切、半精切、精切、校正和修光等工作。

(2)加工精度高、表面粗糙度值小

由于拉削的切削速度较低,切削过程平稳,避免了积屑瘤的出现,加之校准部分的作用,因此,可获得较好的加工质量。拉削一般可以达到的尺寸公差等级为IT8~IT7,表面粗糙度$Ra$值为$0.8 \sim 0.4 \text{ \textmu m}$。

(3)拉床结构简单,刀具寿命长

拉削的运动简单,只有一个主运动,如图14-13所示。拉刀的结构和形状复杂,精度和表面质量要求高,制造成本高。但拉削时速度低,刀具磨损慢,拉刀的寿命长。

拉削应用范围广,主要用来加工各种形状的通孔,如圆孔、方孔、多边形孔和内齿轮等,以及加工各种沟槽,如键槽、T形槽、燕尾槽等。外拉削可加工平面、成形面和外齿轮等,如图14-14所示。


\begin{figure}[H]
\centering
\includegraphics[page=57, viewport={42.0 488.8 261.0 649.3}, clip]{12-07572-001FigRevised.pdf}
\caption{拉孔示意图}
\end{figure}



\begin{figure}[H]
\centering
\includegraphics[page=57, viewport={280.0 490.3 500.5 749.3}, clip]{12-07572-001FigRevised.pdf}
\caption{拉削加工的应用举例}
\end{figure}


拉削加工主要用于大批量生产。对于单件、小批生产精度较高、形状复杂的成形面,若其他方法加工困难时,也可以采用拉削。拉削不能用于加工盲孔、深孔和阶梯孔等。

\subsection{铣削加工}
铣削加工是利用铣刀对工件进行切削加工的方法。铣削时,主运动是铣刀的回转运动,进给运动是工作台带动工件的直线运动或曲线运动。铣削可以用来加工平面、成形面、齿轮、沟槽(包括键槽、V形槽、燕尾槽、T形槽、圆弧槽、螺旋槽等),还可进行孔加工,如钻孔、扩孔等。常用的铣床有卧式铣床、立式铣床、万能铣床等。铣削可以达到的尺寸公差等级为IT8~IT7,表面粗糙度$Ra$值可达$2.3 \sim 1.6 \text{ \textmu m}$。

铣削加工的主要工作如图14-15所示。

1.铣削方式

(1)端铣与周铣

铣平面时,有端铣和周铣两种方法。端铣用铣刀端面上的齿进行铣削;周铣用铣刀圆周上的齿进行铣削,如图14-16所示。两种方法的加工质量均较高。端铣与周铣各具有以下特点。

1)端铣的生产率高于周铣

端铣时有较多的刀齿同时参加切削,工作过程更为平稳,端铣刀大多数镶有硬质合金刀块(条),刚性较好,可以采用大的铣削用量;而周铣用的圆柱铣刀多为高速钢制成,刀轴的刚性较差,使铣削用量受到很大的限制。


\begin{figure}[H]
\centering
\includegraphics[page=57, viewport={97.0 215.8 440.5 458.3}, clip]{12-07572-001FigRevised.pdf}
\caption{铣削加工的主要工作}
\end{figure}



\begin{figure}[H]
\centering
\includegraphics[page=57, viewport={159.0 73.3 376.5 179.3}, clip]{12-07572-001FigRevised.pdf}
\caption{端铣与周铣}
\end{figure}


2)端铣的加工质量比周铣好

端铣时可利用副切削刃对已加工表面进行修光,只要选取合适的副偏角,就可以减小表面粗糙度值;而周铣时只有圆周刃切削,已加工表面实际上由许多圆弧组成,使得表面粗糙度值较大。

3)周铣的适应性比端铣好

周铣可用多种铣刀铣削平面、沟槽、齿形、成形面等,适应性较强;而端铣只能加工平面。比较可知,端铣主要用于大平面的铣削,周铣多用于小平面、各种沟槽和成形面的铣削。

(2)逆铣与顺铣

用周铣铣平面又分为逆铣和顺铣两种方法。当铣刀和工件接触部分的旋转方向与工件的进给方向相反时,这样的铣削称为逆铣;当铣刀和工件接触部分的旋转方向与工件的进给方向相同时,这样的铣削称为顺铣,如图14-17所示。

\begin{figure}[H]
\centering
\includegraphics[page=58, viewport={97.5 513.3 447.5 749.3}, clip]{12-07572-001FigRevised.pdf}
\caption{逆铣与顺铣}
\label{fig:14-17}
\end{figure}

逆铣与顺铣各具有如下特点。

1)逆铣时,铣削厚度从零到最大。切削刃在开始切削时,要在工件已加工表面滑行一小段距离,使刀具磨损加剧,工件表面冷作硬化程度加重,加工表面质量下降;顺铣时,铣削厚度从最大到零,不存在逆铣时的滑行现象,刀具磨损小,工件表面冷作硬化程度轻。

2)逆铣时,工件所受的竖直分力$F_v$将工件向上抬起,不利于压紧工件,会引起振动;顺铣时,铣刀作用在工件上的竖直分力$F_v$有助于压紧工件,铣削过程比较平稳,可提高加工表面质量。

3)顺铣时,忽大忽小的水平分力$F_h$的方向与工作台的进给方向相同,由于工作台进给丝杠与固定螺母之间存在间隙,当水平分力$F_h$值较小时,丝杠与螺母之间的间隙位于右侧,而当水平分力$F_h$值足够大时,就会将工件连同丝杠一起向右拖动,使丝杠与螺母之间的间隙位于左侧。在这种情况下,水平分力$F_h$的大小变化会使工作台忽左忽右来回窜动,造成切削过程不平稳,严重时会打刀甚至损坏机床。逆铣时,水平分力$F_h$与进给方向相反,因此工作台进给丝杠与螺母之间在切削过程中始终保持一侧紧密接触,工作台不会窜动,切削过程平稳。

比较可知,当加工无硬皮的工件,且铣床工作台的进给丝杠和螺母具有消除间隙装置时,采用顺铣;反之,采用逆铣。

2.铣削加工的工艺特点

铣削加工的工艺特点如下。

(1)生产率较高

铣刀是典型的多齿刀具。铣削过程中,多个刀齿依次参加切削工作;主运动是回转运动,可以进行高速切削。

(2)铣削时容易产生振动

铣刀刀齿在切入、切出工件时会产生冲击,每个刀齿的切削厚度随刀齿的运动而发生变化,切削力也随之变化,使切削过程不平稳,容易产生振动,从而影响加工质量。同时,切削振动也限制了加工质量和生产率的进一步提高。

(3)刀齿散热条件较好

铣刀刀齿在切出工件的一段时间内,可以得到一定程度的冷却,有利于刀齿的散热。但是,刀齿在切入、切出工件时,不但受到冲击力,还受到热冲击,这将加速刀具的磨损,甚至使硬质合金刀具碎裂。

\subsection{磨削加工}
磨削是利用砂轮作为切削工具的一种精密加工方法。磨削加工时,砂轮的高速旋转运动为主运动,工件旋转运动为圆周进给,工作台带动工件一起作往复直线运动为纵向进给,当工件一次往复行程终了时,砂轮作周期性的横向进给。

磨削大多在磨床上进行。磨床的种类较多,常见的有外圆磨床、内圆磨床、平面磨床和工具磨床等,分别用于加工零件外圆面、内孔、平面及各种刀具的刃磨。

外圆磨削一般在外圆磨床或万能外圆磨床上进行,磨削方法如图14-18所示。


\begin{figure}[H]
\centering
\includegraphics[page=58, viewport={59.5 281.8 344.5 487.8}, clip]{12-07572-001FigRevised.pdf}
\caption{在外圆磨床上磨外圆}
\end{figure}


磨外圆的方法分为纵磨法和横磨法。纵磨法每次的磨削量很小,磨削余量是在多次往复行程中切除的。

纵磨法的磨削力小,产生热量少,散热较好,可以通过无横向进给光磨进行精磨,所以加工精度和表面质量较高。纵磨法可用一个砂轮加工不同长度的工件,但是生产效率较低,广泛用于单件、小批生产,尤其适合细长轴的磨削。

横磨法又称为切入磨法,工件不作纵向移动,而由砂轮作慢速连续的横向进给,直至磨去全部余量。

横磨法生产率高,适用于大批量生产,尤其是工件上的成形表面,只要将砂轮修整成形,就可直接磨出,较为简便。但是,横磨时工件与砂轮接触面积大,磨削力较大,发热量大,磨削温度高,工件易发生变形和烧伤,适于加工表面不太宽且刚性较好的零件。

孔的磨削可以在内圆磨床上进行,也可以在万能外圆磨床上进行。目前,应用的内圆磨床大多是卡盘式的,它可以加工圆柱孔、圆锥孔、成形内圆面等。纵磨圆柱孔时,工件安装在卡盘上,如图14-19所示,在其旋转的同时,沿轴向作往复直线运动。砂轮高速旋转,并在工件往复行程终了时作周期性的横向进给。

与磨外圆相似,内圆磨削也可以分为纵磨法和横磨法。

磨削具有以下工艺特点。

(1)加工精度高,表面粗糙度值小

由于磨粒的刃口半径小,能切削一层极薄的材料;又由于砂轮表面上的磨粒多,磨削速度高($30 \sim 35 \text{ m/s}$),同时参加切削的磨粒很多,在工件表面形成细小而致密的网纹磨痕,再加上磨床本身的精度高、液压传动平稳和微量进给机构的进给精度高,因此磨削加工的尺寸公差等级可达IT8~IT5,表面粗糙度$Ra$值可达$1.6 \sim 0.2 \text{ \textmu m}$。

\begin{figure}[H]
\centering
\includegraphics[page=58, viewport={362.5 282.3 484.0 377.3}, clip]{12-07572-001FigRevised.pdf}
\caption{磨圆柱孔}
\label{fig:14-19}
\end{figure}

(2)径向分力$F_y$大

磨削力(切削力)一般分解为进给抗力$F_x$、径向分力$F_y$和主切削力$F_z$。车削加工时,主切削力$F_z$最大。磨削加工时,磨削深度和磨粒的切削厚度都较小,因而$F_z$较小,$F_x$更小。因为磨削时砂轮与工件的接触宽度大,磨粒的切削能力较差,所以$F_y$较大。

(3)磨削温度高

由于具有较大负前角的磨粒在高压和高速下,对工件表面进行切削、刻划和滑擦作用,砂轮表面与工件表面之间的摩擦非常严重,消耗功率大,产生的切削热较多。而砂轮本身的导热性差,因此大量的磨削热在很短的时间内不易传出,使磨削区的温度很高,有时高达$800 \sim 1000 \text{ \textdegree C}$。很高的磨削温度容易烧伤工件表面。干磨淬火钢工件时,会使工件退火,硬度降低,因此磨削时必须向磨削区喷注大量的磨削液。

(4)砂轮具有自锐性

砂轮周边的磨粒在磨削力的作用下具有自行脱落的能力,这种能力被称为自锐性。砂轮的自锐性是其他刀具不具备的特性,该特性使砂轮总能以锋利的状态对工件进行持续的加工。

\subsection{精整加工}
1.研磨

研磨是在研(磨)具与工件之间施以研磨剂,研具在一定压力作用下与工件表面之间作复杂的相对运动,通过研磨剂的机械与化学作用,从工件表面上切除很薄的一层材料,从而达到很高的精度和很小的表面粗糙度值的一种精密加工方法。

研具的材料应比工件材料软,以使部分磨粒在研磨过程中能嵌入研具表面,发挥其切削作用。研具可以用铸铁、低碳钢、黄铜、塑料或硬木制造,最常用的是铸铁研具。

研磨剂由磨料、研磨液和辅助填料等混合而成,有液态、膏状和固态三种,以适应不同加工的需要。磨料的硬度高,主要起机械切削作用。机械切削是由游离分散的磨粒作自由滑动、滚动和冲击来完成的。常用的磨料有刚玉、碳化硅等。研磨液主要起冷却和润滑作用,并能使磨粒较均匀地分布在研具表面。常用的研磨液有煤油、汽油、机油等。辅助填料可以使金属表面产生极薄的、较软的化合物薄膜,以便工件表面凸峰容易被磨粒切除,提高研磨效率和表面质量。最常用的辅助填料是硬脂酸、油酸等化学活性物质。

研具与工件之间作复杂的相对运动,使每颗磨粒几乎都不会在工件表面上重复自己的轨迹。这就有可能保证均匀地切除工件表面上的凸峰,获得很小的表面粗糙度值。

图14-20所示为研磨外圆的方法及其所用的研具。工件安装在车床顶尖间或卡盘上,在加工工件表面上涂上研磨剂,将研具套上,工件低速旋转,手持研具作轴向往复移动。


\begin{figure}[H]
    \centering
    % Placeholder for Figure 14-20
    \includegraphics[page=58, viewport={78.5 141.8 463.5 248.8}, clip]{12-07572-001FigRevised.pdf}
    \caption{研磨外圆的方法及其研具}
\end{figure}

图14-21所示为研磨内圆的方法及所用的研具。内圆研具由锥度调节杆和开口锥套组成,锥套的锥孔与调节杆的外锥面相配。锥套上开口使其有一定的弹性,外圆上带有螺旋槽,用以排屑和储存研磨剂。


\begin{figure}[H]
    \centering
    % Placeholder for Figure 14-21
    \includegraphics[page=59, viewport={105.0 632.3 442.0 752.3}, clip]{12-07572-001FigRevised.pdf}
    \caption{研磨内圆的方法及其研具}
\end{figure}

研磨具有如下特点。

(1) 研磨应用广泛

常见的表面如平面、圆柱面、圆锥面、螺纹表面、齿轮齿面等,都可以用研磨进行精密加工。

(2) 研磨可以提高尺寸精度、形状精度,降低表面粗糙度

研磨是在精加工基础上进行的微量切削,尺寸公差等级可达IT5~IT3,表面粗糙度$Ra$值可达$0.1 \sim 0.008 \text{ \textmu m}$,但不能提高位置精度。

(3) 研磨的生产率低

研磨余量一般为$0.01 \sim 0.03 \text{ mm}$。

(4) 加工简单,质量高

加工方法简单,易保证质量,不需要复杂的高精度设备。

研磨不但可加工钢、铸铁、铜、铝及其合金,还可以加工硬质合金、半导体、陶瓷、玻璃、塑料等。在现代工业中,常采用研磨作为精密零件的最终加工。例如,在机械制造业中,用研磨的方法加工精密量块、量规、齿轮、钢球、喷油嘴等零件;在光学仪器制造业中,用研磨精加工镜头、棱镜等零件;在电子工业中,用研磨精加工石英晶体、半导体晶体、陶瓷元件等。

2. 珩磨

珩磨是利用带有油石条的珩磨头对工件进行精密加工的一种方法。图14-22为珩磨加工示意图。珩磨时,珩磨头上的油石以一定的压力压在被加工表面上,由机床主轴带动珩磨头旋转并沿轴向作往复运动(工件固定不动)。在相对运动的过程中,油石从工件表面切除一层极薄的金属,油石在工件表面上的切削轨迹是交叉而不重复的网纹(图14-22b),故可获得很高的精度和很小的表面粗糙度值。

\begin{figure}[H]
    \centering
    % Placeholder for Figure 14-22
    \includegraphics[page=59, viewport={112.0 338.8 434.5 604.3}, clip]{12-07572-001FigRevised.pdf}
    \caption{珩磨加工示意图}
\end{figure}

为了及时地排出切屑和切削热,降低切削温度和减小表面粗糙度值,珩磨过程中要充分浇注珩磨液。珩磨铸铁和钢件时,通常用煤油加少量(质量分数为$10\% \sim 20\%$)机油或锭子油作珩磨液;珩磨青铜等脆性材料时,可以用水剂珩磨液。

在大批、大量生产中,珩磨在专门的珩磨机上进行。在单件、小批生产中,常将立式钻床或卧式车床进行适当改装来完成珩磨加工。

珩磨具有如下工艺特点:

(1) 生产率较高

珩磨时多个油石同时工作,同时参加切削的磨粒较多,能较长时间保持磨粒刃口锋利。珩磨余量比研磨大,一般珩磨铸铁时为$0.02 \sim 0.15 \text{ mm}$,珩磨钢件时为$0.005 \sim 0.08 \text{ mm}$。

(2) 加工精度高

珩磨可提高孔的表面质量、尺寸和形状精度,但不能提高孔的位置精度。这是珩磨头与机床主轴的浮动连接所致。因此,在珩磨前孔的精加工中,必须保证其位置精度。珩磨的尺寸公差等级可达IT6~IT4,表面粗糙度$Ra$值可达$0.2 \sim 0.05 \text{ \textmu m}$。

(3) 珩磨表面耐磨损

由于已加工表面有交叉网纹,利于油膜形成,润滑性能好,故磨损慢。

(4) 不易加工非铁金属

珩磨实际上是一种特殊的磨削,为了避免磨屑堵塞油石孔隙,不易加工塑性较大的非铁金属零件。

(5) 珩磨头结构较复杂,使用过程中调整时间较长

珩磨主要用于孔的精密加工,加工范围很广,能加工直径为$5 \sim 500 \text{ mm}$或更大的孔,并且能加工深孔。

珩磨不仅在大批、大量生产中应用极为普遍,而且在单件、小批生产中应用也较广泛。对于某些零件的孔,珩磨已成为典型的精加工方法,例如飞机、汽车、拖拉机发动机的气缸、气缸套、连杆以及液压油缸、炮筒等。

3. 抛光

抛光是在高速旋转的抛光轮上涂以磨膏,对工件表面进行光整加工的方法。

抛光轮一般是用毛毡、橡胶、皮革、布或压制纸板制成的,磨膏由磨料(氧化铬、氧化铁等)和油酸、软脂酸等配制而成。

抛光时,将工件压于高速旋转的抛光轮上,在磨膏介质的作用下,金属表面产生的一层极薄的软膜,可以用比工件材料软的磨料切除,不会在工件表面留下划痕。高速摩擦使工件表面出现高温,表层材料被挤压而发生塑性流动,这样可填平表面原来的微观不平,获得很光亮、呈镜面状的表面。

抛光具有如下特点。

(1) 方法简便而经济

抛光一般不用特殊设备,工具和加工方法也比较简单,成本低。

(2) 容易对曲面进行加工

由于抛光轮是弹性的,能与曲面相吻合,容易实现曲面抛光,便于对模具型腔进行光整加工。

(3) 只能减小表面粗糙度值而不能保持或提高原加工精度

由于抛光轮与工件之间没有刚性的运动联系,抛光轮又有弹性,因此不能保证从工件表面均匀地切除材料,只是去掉前道工序所留下的痕迹,获得光亮的表面。经过抛光,表面粗糙度$Ra$值可达$0.1 \sim 0.012 \text{ \textmu m}$。

(4) 劳动条件较差

抛光多为手工操作,工作繁重,飞溅的磨粒、介质、微屑等污染环境,劳动条件较差。为改善劳动条件,可采用砂带磨床进行抛光,以代替用抛光轮的手工抛光。

抛光主要用于零件表面的装饰加工,或者用抛光消除前道工序的加工痕迹,以提高零件的疲劳强度,而不是以提高精度为目的。抛光可以加工外圆、孔、平面及各种成形面等。

\section{特种加工方法}

\subsection{特种加工概述}
在机械制造业中,切削加工方法占据着难以替代的主导地位。随着科学技术的不断发展以及新材料、新结构、新工艺的不断诞生,生产实践中出现了切削加工难以解决的加工问题,例如高硬度、高强度、高韧性、高脆性金属或非金属材料(硬质合金、耐热钢、不锈钢、宝石、石英等)的加工,复杂的特殊表面的加工以及特殊零件如细长零件、薄壁零件、弹性零件等的加工。鉴于此,人们发明创造了特种加工方法。

特种加工不是采用常规的刀具或磨具对工件进行切削加工的工艺方法,而是利用电能、光能、化学能、声能、磁能等物理、化学能量或几种复合形式直接施加在被加工的部位,从而使工件改变形状、去除材料、改变性能等。

与传统的切削加工相比,特种加工具有下列特点。

1) 工具材料的硬度可以大大低于工件材料的硬度。

2) 可直接利用电能、电化学能、声能或光能等能量对材料进行加工。

3) 加工过程中的机械力不明显。

4) 各种加工方法可以有选择地复合成新的工艺方法。

特种加工一般按能量形式和作用原理进行如下划分。

1) 电能与热能作用方式:电火花加工、电火花线切割加工、高能束加工。

2) 电能与化学能作用方式:电解加工、电铸加工。

3) 电化学能与机械能作用方式:电解磨削、电解珩磨。

4) 声能与机械能作用方式:超声加工。

5) 光能与热能作用方式:激光加工。

6) 电能与机械能作用方式:离子束加工。

7) 液流能与机械能作用方式:挤压珩磨和水磨料喷射切割。

本节主要介绍电火花加工、电解加工、超声(波)加工、激光加工的基本原理、加工特点及适用范围等。

\subsection{电火花加工}
电火花加工(electrical discharge machining, EDM)是在一定的液体介质中,利用正、负电极间脉冲放电时的电腐蚀现象对导电材料进行加工,从而使零件的尺寸、形状和表面质量达到技术要求的一种加工方法。在特种加工中,电火花加工的应用最为广泛。

1. 电火花加工原理

电火花加工原理如图14-23所示。在加工过程中,脉冲电源的一极接工具电极,另一极接工件电极。两极均浸泡在具有一定绝缘度的液体介质(常用煤油或矿物油)中。工具电极在自动进给调节装置的控制下,以保证工具与工件之间保持一定的放电间隙($0.01 \sim 0.05 \text{ mm}$)。

脉冲电压加到两极之间,便将两极相对间隙最小处或绝缘度最低处的液体介质电离为负电子和正离子而被击穿,形成放电通道。在此瞬间,形成通道的区域很小,放电时间极短,放电区域产生的瞬时高温使被加工材料熔化甚至蒸发,形成一个小凹坑。单个脉冲经过上述放电过程,完成了一次脉冲放电,从而在工件表面形成一个微小的凹坑。如此周而复始高频率地循环放电,工具电极不断地向工件进给,就将工具的形状复制在工件上,最终加工出所需要的零件表面。电火花过程如图14-24所示。

\begin{figure}[H]
    \centering
    % Placeholder for Figure 14-23
    \includegraphics[page=59, viewport={64.0 215.3 321.5 306.8}, clip]{12-07572-001FigRevised.pdf}
    \caption{电火花加工原理示意图}
\end{figure}

\begin{figure}[H]
    \centering
    % Placeholder for Figure 14-24
    \includegraphics[page=59, viewport={341.5 216.8 480.5 310.3}, clip]{12-07572-001FigRevised.pdf}
    \caption{电火花过程示意图}
\end{figure}

2. 电火花加工的特点

电火花加工具有如下特点。

1) 可以加工任何高强度、高硬度、高韧性、高脆性以及高纯度的导电材料,如不锈钢、铁合金、工业纯铁、淬火钢、硬质合金、导电陶瓷、立方氮化硼和人造聚晶金刚石等。

2) 加工时无明显的机械力,故适用于低刚度工件和微细结构的加工,特别适用于复杂的型孔和型腔加工。

3) 脉冲参数可以进行调节,可以在同一台机床上进行粗加工、半精加工和精加工。

4) 在一般情况下生产效率低于切削加工。为了提高生产率,常先采用切削加工进行粗加工,再进行电火花加工。

5) 放电过程有部分能量消耗在工具电极上,从而导致电极损耗,影响成形精度。

3. 电火花加工的应用

电火花加工应用最为广泛的是电火花成形加工和电火花线切割加工。

(1) 电火花成形加工

电火花成形加工是通过工具电极相对于工件作进给运动,将工件电极的形状和尺寸复制在工件上,从而加工出所需要的零件,如图14-25所示。

电火花成形加工主要用于加工各类锻模、压铸模、挤压模、塑料模和胶木模的型腔。这类型腔多为盲孔,内部形状复杂,各处深浅不同,加工较为困难。为了便于排出加工产物和冷却,以提高加工的稳定性,有时在工具电极中间开有冲油孔。

电火花成形还可以用于加工型孔、方孔、多边形孔、异形孔、螺旋孔等,如图14-26所示。

(2) 电火花线切割加工

电火花线切割加工(wire cut EDM, WEDM)是利用移动的细金属丝作工具电极,按预定的轨迹进行脉冲放电切割的加工方法,按线电极移动的速度大小分为高速走丝线切割和低速走丝线切割。我国普遍采用高速走丝线切割,近年来还在发展低速走丝线切割。高速走丝时,线电极是直径为$0.02 \sim 0.3 \text{ mm}$的高强度钼丝,钼丝往复运动的速度为$8 \sim 10 \text{ m/s}$。低速走丝时,多采用铜丝,线电极以小于$0.2 \text{ m/s}$的速度作单方向低速移动。电火花线切割的原理如图14-27所示。

\begin{figure}[H]
    \centering
    % Placeholder for Figure 14-25
    \includegraphics[page=59, viewport={124.0 58.3 410.5 185.3}, clip]{12-07572-001FigRevised.pdf}
    \caption{电火花型腔加工}
\end{figure}

\begin{figure}[H]
    \centering
    % Placeholder for Figure 14-26
    \includegraphics[page=60, viewport={128.5 641.3 413.0 749.3}, clip]{12-07572-001FigRevised.pdf}
    \caption{电火花穿孔加工}
\end{figure}

\begin{figure}[H]
    \centering
    % Placeholder for Figure 14-27
    \includegraphics[page=60, viewport={113.0 442.8 433.0 609.3}, clip]{12-07572-001FigRevised.pdf}
    \caption{电火花线切割原理}
\end{figure}

工作时,脉冲电源的一极接工件,另一极接缠绕金属丝的贮丝筒。如切割内封闭结构零件,电极丝须先穿过工件上预加工的工艺小孔,再经导轮由贮丝筒带动作正、反向的往复移动。工作台在水平面两个坐标方向按各自预定的控制程序,根据放电间隙状态作伺服进给移动,合成各种曲线轨迹,把工件“切割”成形。与此同时,工作液不断喷注在工件与钼丝之间,起绝缘、冷却和冲走屑末的作用。

线切割时,电极丝不断移动,其损耗很小,因而加工精度较高。其平均加工精度可达$0.01 \text{ mm}$,大大高于电火花成形加工。表面粗糙度$Ra$值可达$1.6 \text{ \textmu m}$或更小。

目前,电火花线切割广泛用于加工各种冲裁模(冲孔和落料用)、样板以及各种形状复杂的型孔、型面和窄缝等。

\subsection{电解加工}
电解加工(electrochemical machining, ECM)是电化学加工中的一种特种加工方式,已经广泛应用于涡轮、齿轮、异形孔等复杂型面、型孔的加工以及炮管内膛线的加工和去毛刺等工艺过程。

1. 电解加工的原理

电解加工是利用金属在电解液中产生阳极溶解的电化学原理对工件进行成形加工的一种工艺方法。电解加工原理如图14-28所示。加工时,工件接直流稳压电源正极,工具接负极,两极间保持$0.1 \sim 1 \text{ mm}$的间隙,具有一定压力($0.5 \sim 2.5 \text{ MPa}$)的电解液从两极间隙中高速($5 \sim 60 \text{ m/s}$)流过。加工过程中,工具阴极的凸出部分与工件阳极的电极间隙最小,此处的电流密度最大,单位时间内消耗的电量最多。因此,工件上与工具阴极凸起部位的对应处比其他部位溶解更快。随着工具阴极不断缓慢地向工件进给,工件则不断地按工具端部的型面溶解,电解产物则不断被高速流动的电解液带走,最终工具的形状就“复制”在工件上。

\begin{figure}[H]
    \centering
    % Placeholder for Figure 14-28
    \includegraphics[page=60, viewport={65.0 246.8 472.5 413.8}, clip]{12-07572-001FigRevised.pdf}
    \caption{电解加工原理}
\end{figure}

下面以铁在氯化钠电解液中进行电解加工为例,分析阴极、阳极发生的电化学反应。
水溶液:
\[ \text{H}_2\text{O} \rightleftharpoons \text{H}^+ + \text{OH}^- \]
阳极反应:
\[ \text{Fe} - 2\text{e} \longrightarrow \text{Fe}^{2+} \]
\[ \text{Fe}^{2+} + 2\text{OH}^- \rightleftharpoons \text{Fe(OH)}_2 \downarrow \]
阴极反应:
\[ 2\text{H}^+ + 2\text{e} \longrightarrow \text{H}_2 \uparrow \]

由此可以看出,在电解加工过程中,由于外电源的作用,阳极铁失去电子以$\text{Fe}^{2+}$的形式与水溶液中的负离子$\text{OH}^-$化合生成$\text{Fe(OH)}_2$而沉淀;阴极得到电子,与水溶液中的$\text{H}^+$结合而析出氢气。在电解过程中,工件阳极和水不断消耗,而工具阴极和氯化钠并不消耗。因此,在理想的情况下,工具阴极可长期使用。氯化钠电解液经不断过滤干净并经常补充适量的水,也可长期使用。

电解加工在专用的电解机床上进行,其中的直流稳压电源常采用低电压($6 \sim 24 \text{ V}$)和大电流($500 \sim 20\,000 \text{ A}$)。工具阴极材料常采用黄铜和不锈钢等。

2. 电解加工的特点

电解加工具有如下特点。

1) 可以加工高强度、高硬度和高韧性等难加工的金属材料,如淬火钢、硬质合金、不锈钢和耐热合金等。

2) 加工过程中无机械切削力和切削热,工件不会产生残余应力和变形,也没有飞边、毛刺,适合于加工易变形和薄壁类零件。

3) 平均加工精度可以达到$0.1 \text{ mm}$;表面粗糙度$Ra$值可达$1.6 \sim 0.2 \text{ \textmu m}$。

4) 加工过程中工具阴极在理论上不会损耗,可长期使用。

5) 生产率较高,为电火花加工的$5 \sim 10$倍。

6) 一次进给运动可加工出形状复杂的型腔与型面。

7) 电解加工的附属设备多,造价高,占地面积大,加工稳定性尚不够高。另外,电解液易腐蚀机床和污染环境。

3. 电解加工的应用

(1) 电解锻模型腔

由于电火花加工的精度容易控制,多数锻模的型腔采用电火花加工。但电火花加工的生产率较低,因此对精度要求不太高的矿山机械、汽车拖拉机所需锻模可采用电解加工。

(2) 电解整体叶轮

叶轮是喷气发动机、汽轮机中的关键零件,它的形状复杂,精度要求高,生产批量大。采用电解加工,不受材料硬度和韧性的限制,在一次行程中可加工出复杂的叶片型面,比机械加工有明显的优越性。

采用机械加工方法制造叶轮时,叶片毛坯是精密铸造的,经过机械加工和抛光,再分别镶入叶轮轮缘的榫槽中,最后焊接形成整体叶轮。这种方法加工量大,周期长,质量难以保证。电解加工整体叶轮时(图14-29),只要先将整体叶轮的毛坯加工好,便可用套料法进行加工。每加工完一个叶片,退出阴极,分度后再依次加工下一个叶片,这样不但可大大缩短加工周期,而且可保证叶轮的整体强度和质量。

\begin{figure}[H]
    \centering
    % Placeholder for Figure 14-29
    \includegraphics[page=60, viewport={107.5 53.3 427.0 214.3}, clip]{12-07572-001FigRevised.pdf}
    \caption{电解加工整体叶轮}
\end{figure}

(3)电解去毛刺

机械加工中常采用钳工方法去毛刺,工作量大,有时毛刺难以去除。采用电解加工,则可以提高工效,节省费用。

\subsection{超声加工}

人耳能感受到的声波频率为 $20 \sim 20\,000\text{ Hz}$。频率超过 $20\,000\text{ Hz}$ 的声波即超声波。利用超声波振动,不但能加工像淬火钢、硬质合金等硬脆的导电材料,而且更适合加工像玻璃、陶瓷、宝石和金刚石等硬脆非金属材料。

1. 超声加工原理

超声加工(ultrasonic machining, USM)是利用工具端面的超声频振动,通过磨料悬浮液加工硬脆材料的一种工艺方法。其加工原理如图 14-30 所示。

超声发生器产生的超声频电振荡,通过换能器转变为超声频的机械振动。变幅杆将振幅放大到 $0.01 \sim 0.15\text{ mm}$,再传给工具,并驱动工具端面作超声振动。在加工过程中,由于工具与工件间不断注入磨料悬浮液,当工具端面以超声频冲击磨料时,磨料再冲击工件,迫使加工区域内的工件材料不断被粉碎成很细的微粒脱落下来。此外,当工具端面以很大的加速度离开工件表面时,加工间隙中的工作液内由于负压和局部真空形成许多微空腔。当工具端面再以很大的加速度接近工件表面时,空腔闭合,从而形成可以强化加工过程的液压冲击波,这种现象称为超声空化。工具不断进给,加工持续进行,工具的形状便被复制在工件上,直到符合零件的尺寸要求。

因此,超声加工过程是磨粒在工具端面的超声振动下,以机械锤击和研抛为主,以超声空化为辅的综合作用过程。

\begin{figure}[H]
    \centering
    % Placeholder for Figure 14-30
    \includegraphics[page=61, viewport={54.5 550.3 270.0 688.3}, clip]{12-07572-001FigRevised.pdf}
    \caption{超声加工原理示意图}
\end{figure}

2. 超声加工的特点


1)特别适合加工各种硬脆材料,尤其是电火花加工等无法加工的不导电的非金属材料,如玻璃、陶瓷、玛瑙、宝石、金刚石以及锗和硅等。

2)切削力小,热影响小,适合加工薄壁、窄缝等低刚度工件。

3)加工精度高,加工表面质量好,尺寸精度可达 $0.02 \sim 0.01\text{ mm}$,表面粗糙度 $Ra$ 值可达 $0.8 \sim 0.1\text{ \textmu m}$。加工表面无组织改变、残余应力及烧伤等现象。

4)易加工复杂型面、型孔、型腔等形状。

5)工具和工件无须作复杂的相对运动,因此普通的超声加工设备结构较简单。


3. 超声加工的应用

(1)型孔和型腔加工

目前,超声加工主要用于加工硬脆材料的圆孔、异形孔和各种型腔,以及进行套料、雕刻和研抛等(图 14-31)。

(2)切割加工

锗、硅等半导体材料又硬又脆,用机械切割非常困难,采用超声切割则十分有效,如图 14-32 所示。

\begin{figure}[H]
    \centering
    % Placeholder for Figure 14-31
    \includegraphics[page=61, viewport={290.5 549.3 485.5 741.8}, clip]{12-07572-001FigRevised.pdf}
    \caption{型孔与型腔加工}
\end{figure}

\begin{figure}[H]
    \centering
    % Placeholder for Figure 14-32
    \includegraphics[page=61, viewport={31.0 367.8 204.0 514.3}, clip]{12-07572-001FigRevised.pdf}
    \caption{切割单晶硅片}
\end{figure}

(3)超声清洗

由于超声波在液体中会产生交变冲击波和超声空化现象,当这两种作用的强度达到一定值时,产生的微冲击就可以使被清洗物表面的污渍遭到破坏并脱落下来。目前,超波清洗不但用于机械零件或电子器件的清洗,也用于医疗器皿如生理盐水瓶、葡萄糖水瓶的清洗。利用超声振动去污原理,已生产出超声洗衣机。

(4)超声焊接

首先,被焊工件在超声频振动的作用下,表层的氧化膜被去除,露出新的本体表面。然后,被焊工件表层的分子在高速振动撞击下,摩擦生热并亲和从而焊接在一起。它不仅可以焊接表面易生成氧化物的铝制品及尼龙、塑料等高分子制品,而且可以使陶瓷等非金属材料在超声振动作用下挂上锡或银,从而改善这些材料的可焊接性。

\subsection{激光加工}

激光加工(laser beam machining, LBM)是利用光能经过透镜聚焦后达到很高的能量密度,依靠光热效应来加工各种材料的。它加工速度快、变形小,可以加工各种金属和非金属材料。目前,激光加工已较为广泛应用在切割、打孔、焊接、表面处理、切削加工、快速成形、电阻微调、基板划片以及信息存储等许多领域。

1. 激光加工原理

激光是一种经受激辐射产生的加强光。它的光强度高,方向性、相干性和单色性好,通过光学系统可将激光束聚焦成直径为几十微米到几微米的极小光斑,从而获得极高的能量密度 ($10^8 \sim 10^{10}\text{ W/cm}^2$)。当激光照射到工件表面,光能被工件吸收并迅速转化为热能,光斑区域的温度可达 $10\,000\text{ \textdegree C}$ 以上,使材料熔化甚至汽化。随着激光能量的不断吸收,材料凹坑内的金属蒸气迅速膨胀,压力突然增大,熔融物爆炸式地高速喷射出来,在工件内部形成方向性很强的冲击波。因此,激光加工是工件在光热效应下产生的高温熔融和冲击波的综合作用过程。

是固体激光器中激光的产生和工作原理示意图。激光的工作物质钇铝石榴石受到光泵的激发后,吸收具有特定波长的光,在一定条件下可导致工作物质中的亚稳态粒子数大于低能级粒子数,这种现象称为粒子数反转。此时,一旦有少量激发粒子产生受激辐射跃迁,造成光放大,再通过谐振腔内的全反射镜和部分反射镜的反馈作用产生振荡,此时由谐振腔的一端输出激光。激光通过透镜聚焦形成高能光束,照射在工件表面上,即可进行加工。


\begin{figure}[H]
    \centering
    % Placeholder for Figure 14-33
    \includegraphics[page=61, viewport={225.0 368.3 519.5 499.3}, clip]{12-07572-001FigRevised.pdf}
    \caption{固体激光器中激光的产生和工作原理示意图}
\end{figure}

2. 激光加工的特点


1)激光的功率密度可高达 $10^8 \sim 10^{10}\text{ W/cm}^2$,加工的热作用时间很短,热影响区小,几乎可以加工任何金属与非金属材料。

2)激光加工属于非接触加工,不需要工具,不存在工具损耗,无明显机械力,加工速度快,易于实现加工过程自动化。

3)激光可通过玻璃等透明材料进行加工,如对真空管内部进行焊接等。

4)激光可用于精密微细加工。平均加工精度可以达到 $0.01\text{ mm}$,表面粗糙度 $Ra$ 值可达 $0.4 \sim 0.1\text{ \textmu m}$。

5)激光加工不受电磁干扰。

6)激光除加工外,还可以进行焊接、热处理、表面强化等。


3. 激光加工的应用

(1)激光打孔

激光打孔主要应用于宝石、金刚石、玻璃、硬质合金、不锈钢等特殊材料或特殊零件上加工孔,如火箭发动机和柴油机的喷油嘴、化学纤维的喷丝板、钟表上的宝石轴承和聚晶金刚石拉丝模等零件上的微细孔加工。激光打孔的效率很高,如直径为 $0.12 \sim 0.18\text{ mm}$、深为 $0.6 \sim 1.2\text{ mm}$ 的宝石轴承孔,若工件自动传送,每分钟可加工数十件。在聚晶金刚石拉丝模坯料的中央加工直径为 $0.04\text{ mm}$ 的小孔,仅需十几秒。

(2)激光切割

激光既可以切割厚度达 $10\text{ mm}$ 以上的金属材料,也可以切割厚度达几十毫米的非金属材料。它还能透过玻璃切割真空管内的灯丝,这是任何机械加工所难以实现的。

(3)激光焊接

激光焊接一般无需焊料和焊剂,只需将工件的加工区域“热熔”在一起就可以。激光焊接过程迅速,热影响区小,焊缝质量高,既可以焊接同种材料,也可以焊接异种材料,还可以透过玻璃进行焊接。

(4)激光表面热处理

利用激光对工件材料表面进行扫射,使金属表层材料产生相变甚至融化,当激光束离开工件表面时,热量迅速向内部传导,表面冷却且硬化,从而可提高工件的耐磨性和疲劳强度。一般情况下,也可实现对铸铁、中碳钢甚至低碳钢等材料进行激光表面淬火。激光淬火层的淬硬深度一般为 $0.7 \sim 1.1\text{ mm}$。淬火层的硬度比常规淬火约高 $20\%$,产生的变形小,能解决低碳钢的表面淬火强化问题。

\section{零件表面加工方法选择}

\subsection{选择表面加工方法的依据}

组成零件的各种表面,不仅具有一定的形状和尺寸,同时还要求达到一定的技术要求。

工件表面的加工过程,就是获得符合要求的零件表面的过程。由于零件的结构特点、材料性能和表面加工要求的不同,所采用的加工方法也不一样。即使是同一精度要求,所采用的加工方法也是多种多样的。在选择某一表面的加工方法时,应遵循如下基本原则。


1)选择的加工方法要与加工表面的经济精度及表面粗糙度要求相适应。

2)选择的加工方法要与零件材料的性能、热处理要求及产品的生产类型相适应。

3)选择的加工方法要与零件的结构形状和尺寸要求相适应。

4)根据零件表面的具体要求,考虑各种加工方法的特点和应用,选用几种加工方法组合起来,完成精度较高的零件表面加工。

5)当零件表面的加工质量要求较高时,整个加工过程应分阶段进行。一般分为粗加工、半精加工和精加工三个阶段。表面质量要求更高时,还要进一步进行精整加工。


本节主要介绍螺纹及成形面的技术要求和加工方法。

\subsection{螺纹的加工}

螺纹是零件上常见的表面之一,它有多种形式,按用途的不同可分为如下两类。

(1)紧固螺纹

紧固螺纹用于零件间的固定连接,常用的有普通螺纹和管螺纹等,螺纹牙型多为三角形。对普通螺纹的主要要求是可旋入性和连接的可靠性;对管螺纹的主要要求是密封性和连接的可靠性。

(2)传动螺纹

传动螺纹用于传递动力、运动或位移,如丝杠和测微螺杆的螺纹等,其牙型多为梯形或锯齿形。对于传动螺纹的主要要求是传动准确、可靠,螺牙接触良好及耐磨等。

1. 螺纹的技术要求

螺纹和其他类型的表面一样,有一定的尺寸精度、几何精度和表面质量的要求。由于它们的用途和使用要求不同,技术要求也有所不同。

对于紧固螺纹和无传动精度要求的传动螺纹,一般只要求中径、外螺纹的大径、内螺纹的小径精度。

对于有传动精度要求或用于读数的螺纹,除要求中径和顶径的精度外,还要求螺距和牙型角的精度。为了保证传动或读数精度及耐磨性,对螺纹表面的粗糙度和硬度等也有较高的要求。

2. 螺纹加工方法的分析

螺纹的加工方法很多,可以在车床、钻床、螺纹铣床、螺纹磨床等机床上利用不同的工具进行加工。选择螺纹的加工方法要考虑的因素较多,主要有工件形状、螺纹牙型、螺纹的尺寸和精度、工件材料和热处理以及生产类型等。表 14-3 列出了常见螺纹加工方法所能达到的精度和表面粗糙度,可以作为选择螺纹加工方法的依据和参考。

\begin{table}[H]
    \centering
    \caption{常见螺纹加工方法所能达到的精度和表面粗糙度}
    \begin{tabular}{|c|c|c|}
        \hline
        加工方法 & 精度等级 & 表面粗糙度 $Ra$ 值/$\text{\textmu m}$ \\
        \hline
        攻螺纹 & $8 \sim 6$ & $6.3 \sim 1.6$ \\
        套螺纹 & $8 \sim 7$ & $3.2 \sim 1.6$ \\
        车削 & $8 \sim 4$ & $1.6 \sim 0.4$ \\
        铣刀铣削 & $8 \sim 6$ & $6.3 \sim 3.2$ \\
        磨削 & $6 \sim 4$ & $0.4 \sim 0.1$ \\
        研磨 & $4$ & $0.1$ \\
        滚压 & $8 \sim 4$ & $0.8 \sim 0.1$ \\
        \hline
    \end{tabular}
\end{table}

3. 常见的螺纹加工方法

(1)攻螺纹和套螺纹

攻螺纹和套螺纹是应用较广的螺纹加工方法。对于小尺寸的内螺纹,攻螺纹几乎是唯一有效的加工方法。单件、小批生产时,可以用丝锥手工攻螺纹;当批量较大时,则应在车床、钻床或攻螺纹机上用机用丝锥加工。套螺纹的螺纹直径一般不超过 $16\text{ mm}$,它既可以手工操作,也可以在机床上进行。

攻螺纹和套螺纹的加工精度较低,主要用于加工精度要求不高的普通螺纹。

(2)车螺纹

车螺纹是螺纹加工的基本方法,它可以使用通用设备,刀具简单(图 14-34),适应性广,可用来加工各种形状、尺寸及精度的内、外螺纹,特别适于加工尺寸较大的螺纹。但是,车螺纹的生产率较低,加工质量取决于工人的技术水平以及机床、刀具本身的精度,所以主要用于单件、小批生产。对于非淬硬精密丝杠,利用精密车床车削可以获得较高的精度和较小的表面粗糙度值,因此占有重要的地位。

车削螺纹所用的刀具为具有螺纹牙型廓形的成形车刀。当生产批量较大时,为了提高生产率,常采用螺纹梳刀(图 14-35)进行车削。螺纹梳刀实质上是一种多齿的螺纹车刀,只要一次走刀就能切出全部螺纹,生产率较高。但是,一般的螺纹梳刀加工精度不高,不能加工精密螺纹。此外,螺纹附近有轴肩的工件也不能用螺纹梳刀加工。

\begin{figure}[H]
    \centering
    % Placeholder for Figure 14-34
    \includegraphics[page=61, viewport={126.0 216.8 412.0 336.8}, clip]{12-07572-001FigRevised.pdf}
    \caption{几种常见的螺纹车刀}
\end{figure}

\begin{figure}[H]
    \centering
    % Placeholder for Figure 14-35
    \includegraphics[page=61, viewport={33.0 85.8 362.0 182.8}, clip]{12-07572-001FigRevised.pdf}
    \caption{螺纹梳刀}
\end{figure}

(3)铣螺纹

在成批和大量生产中,广泛采用铣削法加工螺纹。铣螺纹一般在专门的螺纹铣床上进行,根据所用铣刀的结构不同,可以分为如下两种方法。


1)用盘形螺纹铣刀铣削(图 14-36)。这种方法一般用于加工尺寸较大的传动螺纹。由于加工精度较低,该方法通常只作为粗加工,精加工可选用车削。

2)用梳形螺纹铣刀铣削(图 14-37)。一般用于加工螺距不大、短的普通(三角形)内、外螺纹。加工时,工件只需转一周多一点就可以切出全部螺纹,因此生产率较高。这种方法可以加工靠近轴肩或盲孔底部的螺纹,且不需要退刀槽,但其加工精度较低。


\begin{figure}[H]
    \centering
    % Placeholder for Figure 14-36
    \includegraphics[page=61, viewport={380.0 93.8 514.0 186.8}, clip]{12-07572-001FigRevised.pdf}
    \caption{盘形螺纹铣刀加工螺纹}
\end{figure}

\begin{figure}[H]
    \centering
    % Placeholder for Figure 14-37
    \includegraphics[page=62, viewport={7.0 621.8 165.0 730.8}, clip]{12-07572-001FigRevised.pdf}
    \caption{梳形螺纹铣刀加工螺纹}
\end{figure}

(4)磨螺纹

磨螺纹常用于淬硬螺纹的精加工,例如丝锥、螺纹量规、滚丝轮及精密螺杆上的螺纹。为了修正热处理引起的变形,提高加工精度,必须进行磨削。螺纹磨削一般在专门的螺纹磨床上进行。螺纹在磨削之前,可以用车、铣等方法进行预加工,而对于小尺寸的精密螺纹,也可以不经预加工而直接磨出。

根据所用砂轮的形状不同,外螺纹的磨削可以分为单线砂轮磨削(图 14-38)和多线砂轮磨削(图 14-39)。

\begin{figure}[H]
    \centering
    % Placeholder for Figure 14-38
    \includegraphics[page=62, viewport={190.0 621.3 376.5 742.3}, clip]{12-07572-001FigRevised.pdf}
    \caption{单线砂轮磨螺纹}
\end{figure}

\begin{figure}[H]
    \centering
    % Placeholder for Figure 14-39
    \includegraphics[page=62, viewport={395.0 625.3 539.5 744.3}, clip]{12-07572-001FigRevised.pdf}
    \caption{多线砂轮磨螺纹}
\end{figure}

用单线砂轮磨螺纹,砂轮的修整较方便,加工精度较高,并且可以加工较长的螺纹。而用多线砂轮磨螺纹,砂轮的修整比较困难,加工精度低于前者,且仅适于加工较短的螺纹。但是,用多线砂轮磨削时,工件转 $1\frac{1}{3} \sim 1\frac{1}{2}$ 转就可以完成磨削加工,生产效率较单线砂轮磨削高。

直径大于 $30\text{ mm}$ 的内螺纹,也可以用单线砂轮磨削。

\subsection{成形面的加工}

带有成形面的零件,机器上用得也相当多,如机床的手柄、内燃机凸轮轴上的凸轮、汽轮机的叶轮、锻模模膛等,如图 14-40 所示。

\begin{figure}[H]
    \centering
    % Placeholder for Figure 14-40
    \includegraphics[page=62, viewport={101.5 498.3 441.0 588.8}, clip]{12-07572-001FigRevised.pdf}
    \caption{各种常见的成形面}
\end{figure}

1. 成形面的技术要求

与其他表面类似,成形面的技术要求也包括尺寸精度、几何精度及表面质量等。但是,成形面往往是为了实现特定功能而专门设计的,因此其表面形状的要求是十分重要的。加工时,刀具的切削刃形状和切削运动,应首先满足表面形状的要求。

2. 成形面加工方法的分析

成形面的加工可以分别采用车削、铣削、刨削、拉削或磨削等方法。这些方法可以归纳为如下两种基本方式。

(1)用成形刀具加工

用成形刀具加工即用切削刃形状与工件廓形相符合的刀具直接加工出成形面。例如,用成形车刀车成形面(图 14-41a)、用成形刨刀刨成形面(图 14-41b)用成形铣刀铣成形面(图 14-41c)等。

\begin{figure}[H]
    \centering
    % Placeholder for Figure 14-41
    \includegraphics[page=62, viewport={12.0 312.8 315.0 416.8}, clip]{12-07572-001FigRevised.pdf}
    \caption{利用成形法加工成形面}
\end{figure}

用成形刀具加工成形面,机床的运动和结构比较简单,操作简便,但是刀具的制造和刃磨比较复杂(特别是成形铣刀和拉刀),成本较高。而且,这种方法的应用受工件成形面尺寸的限制,不宜用于加工刚度差而成形面较宽的工件。

(2)利用刀具和工件作特定的相对运动加工

用双手控制法车成形面(图 14-42)就是其中的一种。此外,还可以利用液压仿形装置或数控装置等来控制刀具与工件之间特定的相对运动来加工成形面。

利用刀具和工件作特定的相对运动来加工成形面,刀具比较简单,并且加工成形面的尺寸范围较大。但是,机床的运动和结构都较复杂,成本也高。

成形面的加工方法,应根据零件的尺寸、形状及生产批量等来选择。用靠模法车成形面(图 14-43),适应于零件形状较简单、生产批量大的情况。对大批、大量生产小型回转体零件上形状不太复杂的成形面,常用成形车刀在自动或半自动车床上加工;批量较小时,可用成形车刀在卧式车床上加工。对直槽和螺旋槽等,一般可用成形铣刀在万能铣床上加工。

\begin{figure}[H]
    \centering
    % Placeholder for Figure 14-42
    \includegraphics[page=62, viewport={340.5 312.8 532.5 464.8}, clip]{12-07572-001FigRevised.pdf}
    \caption{双手控制法车成形面}
\end{figure}

\begin{figure}[H]
    \centering
    % Placeholder for Figure 14-43
    \includegraphics[page=62, viewport={10.5 124.3 224.5 293.8}, clip]{12-07572-001FigRevised.pdf}
    \caption{用靠模法车成形面}
\end{figure}

大批、大量生产中，尺寸较大的成形面多采用仿形车床或仿形铣床加工；单件、小批生产时，可借助样板在卧式车床上加工，或者依据划线在铣床或刨床上加工，但这种方法加工的质量和效率较低。为了保证加工质量和提高生产效率，在单件、小批生产中可应用数控机床加工成形面。

大批、大量生产中，为了加工一定的成形面，常常专门设计和制造专用的刀具或专门化的机床，例如加工凸轮轴所用的凸轮轴车床、凸轮轴磨床等。

对于淬硬的成形面或精度高、表面粗糙度值小的成形面，其精加工则要采用磨削，甚至要用精整加工。

\section*{复习思考题}

 
14-1 简述车削的工艺特点。

14-2 常用的车床类型有哪些？在车床上能够完成哪些工作？

14-3 为什么车削应用广泛？车削能够达到的经济精度和表面粗糙度范围是多少？

14-4 在车床上钻孔，若钻头引偏对加工的孔有何影响？

14-5 简述改善钻削时的排屑条件的措施。

14-6 钻削的特点是什么？

14-7 一般情况下，为什么车削过程比铣削、刨削平稳？

14-8 铣平面时，为什么端铣比周铣优越？

14-9 何谓逆铣和顺铣？它们各有何特点？分别适用于何种场合？

14-10 简述铣削的工艺特点和应用。

14-11 刨削加工的工艺特点是什么？

14-12 试比较铰孔和镗孔的工艺特点和应用。

14-13 镗床上可加工哪些工件？镗削加工的工艺特点是什么？

14-14 镗削加工的切削运动与车削加工的切削运动有何不同？

14-15 镗刀有单刃镗刀和多刃镗刀，简述它们的结构和工艺特点。

14-16 拉削加工有何特点？其切削运动与一般机械加工方法有何不同？

14-17 磨削加工的精度和表面粗糙度一般可达到什么范围？

14-18 外圆表面磨削时，磨削用量包括哪几个要素？怎样合理选择磨削要素？

14-19 为什么珩磨不能提高孔的位置精度？

14-20 试述特种加工的分类方法以及电火花加工、电解加工、超声加工、激光加工的特点和应用范围。

14-21 电火花加工具有哪些特点？

14-22 列表比较车外圆、磨外圆和研磨外圆的设备、工具、主运动、进给运动、加工精度和表面粗糙度。

14-23 列表比较车削、铣削、磨削的加工要素、表面粗糙度、生产率、加工成本及适应性。


\chapter{机械零件的结构工艺性}

零件结构对零件的工艺过程影响很大。使用性能完全相同而结构不同的两个零件，其加工方法与制造成本却可能有很大的差别。因此，所设计的零件应具有良好的结构工艺性。

良好的结构工艺性是指根据使用要求所设计的零件结构，能在同样的生产条件下，采用高效率、低消耗和低成本的方法制造出来。

满足零件的使用要求是零件结构工艺性设计的前提。在此前提下，还应尽可能使毛坯生产、切削加工、热处理和装配调试等各个阶段都具有良好的结构工艺性。如果不能兼顾各阶段，也应做到保证主要，照顾次要。零件结构工艺性的优劣是相对的，与现有设备条件、生产类型、技术水平和工艺方法等密切相关。例如，非圆形通孔在批量小的情况下，结构工艺性不好，但在批量大的情况下，可采用拉削方法加工，则具有较好的结构工艺性；又如电火花加工的出现，使原来认为不易加工的复杂型腔、直径很小的孔等表面变得容易加工。

在整个制造过程中，零件的切削加工所消耗的时间及费用最多，因此零件结构的切削加工工艺性的好坏就显得特别重要。表 15-1 至表 15-3 以“便于零件安装、加工，提高切削效率，减少切削加工量和易于保证加工质量”为指导原则。

\section{便于安装、加工的机械零件结构}

表 15-1 给出了便于安装、加工的机械零件结构例子，通过合理与不合理结构的对比可以发现，合理的零件结构能够降低加工难度，减少工装投入，提高定位可靠性，保证加工质量。

\begin{longtable}{|m{1cm}|m{5cm}|m{5cm}|m{2cm}|}
\caption{便于安装、加工的机械零件结构举例}\\
\hline
\multirow{2}{1cm}{\centering 设计准则} & \multicolumn{2}{c|}{图例} & \multirow{2}{2cm}{\centering 说明} \\
\cline{2-3}
& \centering 不合理的结构 & \centering 合理的结构 & \\
\hline
\multirow{3}{1cm}{\centering 便于装夹、保证定位可靠} &\includegraphics[page=64, viewport={330.5 641.3 471.5 704.8}, clip]{12-07572-001FigRevised.pdf} &\includegraphics[page=64, viewport={33.0 533.3 180.5 587.8}, clip]{12-07572-001FigRevised.pdf} & 锥度心轴一般先车后磨，用顶尖、拨盘、卡箍装夹，应在心轴一端设计一圆柱表面，以便安装卡箍 \\
\cline{2-4}
& \includegraphics[page=64, viewport={203.0 534.8 349.5 603.8}, clip]{12-07572-001FigRevised.pdf} & \includegraphics[page=64, viewport={368.5 532.8 513.5 606.8}, clip]{12-07572-001FigRevised.pdf} & 为了安装方便，增加了工艺凸台B，精加工后，再把凸台B切除 \\
\cline{2-4}
&\includegraphics[page=64, viewport={109.0 382.3 256.0 490.3}, clip]{12-07572-001FigRevised.pdf} &\includegraphics[page=64, viewport={279.5 383.3 435.5 493.3}, clip]{12-07572-001FigRevised.pdf} & 马达端盖\\
\hline
\multirow{3}{1cm}{\centering 减小加工难度} & \includegraphics[page=64, viewport={342.0 20.8 446.0 115.3}, clip]{12-07572-001FigRevised.pdf} &\includegraphics[page=65, viewport={8.0 637.3 110.0 734.8}, clip]{12-07572-001FigRevised.pdf} & 箱体的同轴孔系应尽可能设计成无台阶的通孔。孔径应向一个方向递减，孔的端面应在同一平面上 \\
\cline{2-4}
& \includegraphics[page=65, viewport={134.5 637.3 218.5 683.8}, clip]{12-07572-001FigRevised.pdf} &\includegraphics[page=65, viewport={238.0 637.8 325.0 687.3}, clip]{12-07572-001FigRevised.pdf} & 复杂内孔表面采用组合件，可简化内部复杂面的加工，易于保证质量，并可简化刀具结构 \\
\cline{2-3}
& \includegraphics[page=65, viewport={345.5 637.3 433.0 690.3}, clip]{12-07572-001FigRevised.pdf} &\includegraphics[page=65, viewport={455.5 638.3 537.5 684.8}, clip]{12-07572-001FigRevised.pdf} & \\
\hline
\multirow{4}{1cm}{\centering 便于进刀和退刀} & \includegraphics[page=65, viewport={56.0 474.8 157.5 521.8}, clip]{12-07572-001FigRevised.pdf} &\includegraphics[page=65, viewport={179.5 475.3 273.0 577.3}, clip]{12-07572-001FigRevised.pdf} & 加工内、外螺纹时应留有退刀槽或保留足够的退刀长度 \\
\cline{2-4}
& \includegraphics[page=65, viewport={295.0 475.8 377.5 538.8}, clip]{12-07572-001FigRevised.pdf} &\includegraphics[page=65, viewport={399.5 473.3 482.5 602.3}, clip]{12-07572-001FigRevised.pdf} & \\
\cline{2-4}
 & \includegraphics[page=65, viewport={56.0 474.8 157.5 521.8}, clip]{12-07572-001FigRevised.pdf} &\includegraphics[page=65, viewport={179.5 475.3 273.0 577.3}, clip]{12-07572-001FigRevised.pdf} & 需要磨削的内、外圆，其根部应有砂轮越程槽 \\
\cline{2-3}
& \includegraphics[page=65, viewport={295.0 475.8 377.5 538.8}, clip]{12-07572-001FigRevised.pdf} &\includegraphics[page=65, viewport={399.5 473.3 482.5 602.3}, clip]{12-07572-001FigRevised.pdf} & \\
\hline
\multirow{2}{1cm}{\centering 力求几何形状简单} &\includegraphics[page=65, viewport={255.0 282.3 338.0 333.3}, clip]{12-07572-001FigRevised.pdf} & \includegraphics[page=65, viewport={352.0 281.3 423.5 333.3}, clip]{12-07572-001FigRevised.pdf} & 把阶梯孔改成简单的孔，便于加工 \\
\cline{2-4}
& \includegraphics[page=65, viewport={62.5 123.8 188.5 243.8}, clip]{12-07572-001FigRevised.pdf} & \includegraphics[page=65, viewport={207.0 123.8 334.0 217.8}, clip]{12-07572-001FigRevised.pdf} & 沟槽底部若是圆弧，铣刀直径必须与工件圆弧直径一致；槽底若为平面，则可选择任何直径的铣刀加工。槽底为平面，还可多件串联起来同时加工，提高生产率 \\
\hline
\end{longtable}


\section{减少切削加工量、降低成本的机械零件结构}

表 15-2 给出了减少切削加工量、降低成本的机械零件结构的例子，通过对合理与不合理结构的对比可以发现，合理结构能够减少原材料消耗，降低毛坯成本，便于使用标准刀具。

\begin{longtable}{|m{1.5cm}|m{4cm}|m{4cm}|m{4cm}|}
\caption{减少切削加工量、降低成本的机械零件结构举例}\\
\hline
\multirow{2}{1.5cm}{\centering 设计准则} & \multicolumn{2}{c|}{图例} & \multirow{2}{1.5cm}{\centering 说明} \\
\cline{2-3}
& \centering 不合理的结构 & \centering 合理的结构 & \\
\hline
\multirow{3}{1.5cm}{\centering 减少加工表面面积} &\includegraphics[page=65, viewport={352.0 121.3 472.5 174.8}, clip]{12-07572-001FigRevised.pdf} & \includegraphics[page=65, viewport={79.5 38.8 196.5 88.3}, clip]{12-07572-001FigRevised.pdf} & 当轴与盘套类零件相配合时，应保证配合部位的精度，非配合表面不必加工成高精度，这样不仅可减少精加工面积，且易于保证质量 \\
\cline{2-4}
& \includegraphics[page=65, viewport={214.0 35.3 327.0 92.3}, clip]{12-07572-001FigRevised.pdf} &\includegraphics[page=65, viewport={348.5 34.3 463.5 88.3}, clip]{12-07572-001FigRevised.pdf} & 当孔的长度与直径之比较大时，应保证与轴相配合部位的精度，以节省材料和降低工时 \\
\cline{2-4}
& \includegraphics[page=66, viewport={95.0 632.8 212.5 749.3}, clip]{12-07572-001FigRevised.pdf} &\includegraphics[page=66, viewport={227.5 634.8 349.5 747.8}, clip]{12-07572-001FigRevised.pdf}& 轴承座、箱体、机架等零件的底平面，应设计成中部成凹状的平面，以减少加工面积，并保证工作可靠 \\
\cline{2-3}
& \includegraphics[page=66, viewport={365.5 630.8 456.5 735.3}, clip]{12-07572-001FigRevised.pdf} &\includegraphics[page=66, viewport={83.0 497.8 172.0 599.3}, clip]{12-07572-001FigRevised.pdf}&\\
\hline
\multirow{3}{1.5cm}{\centering 便于采用标准刀具} &\includegraphics[page=66, viewport={190.5 499.8 318.0 554.8}, clip]{12-07572-001FigRevised.pdf} &\includegraphics[page=66, viewport={339.0 496.3 465.0 554.8}, clip]{12-07572-001FigRevised.pdf} & 孔应设计成标准直径，以便选用标准钻头等定径刀具；直径过渡应采用与钻头顶角相同的锥度 \\
\cline{2-4}
& \includegraphics[page=66, viewport={83.5 364.8 193.0 466.8}, clip]{12-07572-001FigRevised.pdf} &\includegraphics[page=66, viewport={209.5 365.8 329.5 471.8}, clip]{12-07572-001FigRevised.pdf} & 槽的形状（直角、圆角）和尺寸应与立铣刀形状相符 \\
\cline{2-4}
& \includegraphics[page=66, viewport={342.0 361.8 467.5 453.8}, clip]{12-07572-001FigRevised.pdf} & \includegraphics[page=66, viewport={70.5 199.3 201.5 332.3}, clip]{12-07572-001FigRevised.pdf} & 需铣削的凹面内圆角的直径应等于标准立铣刀直径，并且下凹表面越深，内圆角直径越大，以增强刀具刚性 \\
\hline
\end{longtable}

\section{减少加工辅助时间、减小工件变形的机械零件}

表 15-3 给出了减少加工辅助时间，减小工件变形的机械零件结构例子，通过合理与不合理结构的对比可以发现，合理结构能减少刀具种类、换刀和装夹次数，缩短工时，提高设备利用率；并能够增加工件刚度，减小工件变形，保证加工精度。

\begin{longtable}{|m{1.5cm}|m{5cm}|m{5cm}|m{3cm}|}
\caption{减少加工辅助时间、减小工件变形的零件结构举例}\\
\hline
\multirow{2}{1.5cm}{\centering 设计准则} & \multicolumn{2}{c|}{图例} & \multirow{2}{1.5cm}{\centering 说明} \\
\cline{2-3}
& \centering 不合理的结构 & \centering 合理的结构 & \\
\hline
\multirow{3}{1.5cm}{\centering 减少加工辅助时间} & \includegraphics[page=66, viewport={213.0 199.3 329.5 257.8}, clip]{12-07572-001FigRevised.pdf} &\includegraphics[page=66, viewport={355.0 200.8 465.0 253.3}, clip]{12-07572-001FigRevised.pdf} & 同一零件上结构相同的槽（键槽、刀槽），其宽度（包括内圆角半径）应尽可能一致，以减少刀具种类和换刀次数 \\
\cline{2-3}
 & \includegraphics[page=66, viewport={30.5 65.3 189.0 125.8}, clip]{12-07572-001FigRevised.pdf} &\includegraphics[page=66, viewport={204.5 67.8 350.0 122.3}, clip]{12-07572-001FigRevised.pdf} &\\
& \includegraphics[page=66, viewport={370.0 62.8 508.0 173.8}, clip]{12-07572-001FigRevised.pdf}  &\includegraphics[page=67, viewport={49.0 650.8 205.0 748.3}, clip]{12-07572-001FigRevised.pdf} & 箱体上的螺纹孔直径应尽量一致或减少种类，以便采用同一丝锥或减少丝锥规格 \\
\cline{2-4}
&\includegraphics[page=67, viewport={215.5 650.8 339.5 707.3}, clip]{12-07572-001FigRevised.pdf}  &\includegraphics[page=67, viewport={357.0 650.3 490.5 709.3}, clip]{12-07572-001FigRevised.pdf} & 键槽的尺寸、方位应尽量一致，便于在一次走刀中铣出各键槽 \\
\hline
\multirow{2}{1.5cm}{\centering 减少安装次数} &\includegraphics[page=67, viewport={89.0 557.3 197.0 614.3}, clip]{12-07572-001FigRevised.pdf} &\includegraphics[page=67, viewport={220.5 556.3 341.5 612.8}, clip]{12-07572-001FigRevised.pdf} & 需要磨削或精车的零件，应考虑在一次安装中完成，以保证加工精度和节省工时 \\
\cline{2-4}
&\includegraphics[page=67, viewport={362.5 557.8 456.0 620.3}, clip]{12-07572-001FigRevised.pdf} &\includegraphics[page=67, viewport={33.5 464.3 128.5 527.8}, clip]{12-07572-001FigRevised.pdf} & 原设计需要二次装夹，改进设计后一次装夹即可，且易于保证孔的同轴度 \\
\hline
\multirow{1}{1.5cm}{\centering 减少对刀次数} &\includegraphics[page=67, viewport={144.5 459.8 262.5 511.8}, clip]{12-07572-001FigRevised.pdf} & \includegraphics[page=67, viewport={280.0 460.8 399.0 505.8}, clip]{12-07572-001FigRevised.pdf} & 零件表面上的凸台，应尽可能布置在同一平面上，以便在一次对刀中加工出各凸台 \\
\hline
\multirow{1}{1.5cm}{\centering 考虑夹紧力} & \includegraphics[page=67, viewport={413.0 464.8 509.0 529.3}, clip]{12-07572-001FigRevised.pdf} &\includegraphics[page=67, viewport={11.5 343.8 102.0 414.3}, clip]{12-07572-001FigRevised.pdf} & 薄壁、套筒类零件夹紧时易变形，若一端加凸缘，可提高工件刚度 \\
\hline
\multirow{2}{1.5cm}{\centering 使零件有足够的刚度} &\includegraphics[page=67, viewport={124.5 344.3 209.0 416.3}, clip]{12-07572-001FigRevised.pdf} & \includegraphics[page=67, viewport={234.0 344.3 324.5 417.3}, clip]{12-07572-001FigRevised.pdf} & 设置加强筋，可提高零件刚度，减少刨削或铣削时工件的振动或变形 \\
\cline{2-3}
& \includegraphics[page=67, viewport={341.0 347.8 425.0 432.3}, clip]{12-07572-001FigRevised.pdf} &\includegraphics[page=67, viewport={440.5 349.8 530.0 433.8}, clip]{12-07572-001FigRevised.pdf} & \\
\hline
\end{longtable}

\section*{复习思考题}

 
15-1 什么是零件的结构工艺性?零件的结构设计在生产中有什么意义?

15-2 为什么说零件的结构工艺性具有相对性?

15-3 要使零件在切削加工过程中具有良好的结构工艺性,除满足使用性能要求外,还应遵守的基本原则是什么?

15-4 简述零件结构切削加工工艺性的重要性。

15-5 简述电动机端盖在便于安装、加工等方面的设计。

15-6 试举例说明零件必须设置退刀槽的结构。

15-7 试举例说明应该减少零件加工面积的结构。

15-8 试举例说明如何提高工件在加工时的刚度。

15-9 铣削牙嵌式离合器时,为什么铣四个齿不如铣五个齿省工?试绘图进行分析。

15-10 减少加工辅助时间的措施有哪些?

15-11 指出图15-1所示零件难以加工或无法加工的部位,并提出改进意见。


\begin{figure}[H]
    \centering
    % Placeholder for Figure 15-1 containing subfigures (a) through (h)
    \includegraphics[page=63, viewport={64.0 94.8 473.0 410.3}, clip]{12-07572-001FigRevised.pdf}
    \caption{题 15-11 图}
\end{figure}

15-12 判断图15-2所示零件结构工艺性的好坏,并简要说明理由。

\begin{figure}[H]
    \centering
    % Placeholder for Figure 15-2
    \includegraphics[page=64, viewport={57.5 640.8 302.5 745.8}, clip]{12-07572-001FigRevised.pdf}
    \caption{题 15-12 图}
\end{figure}


\chapter{机械加工工艺过程的基础知识}

\section{机械加工工艺过程的基本概念}

\subsection{生产过程和工艺过程}

1. 生产过程

制造机器时,由原材料制成各种零件并装配成机器的全过程称为生产过程。它包括原材料的运输和保管,生产准备工作,毛坯制造,零件加工和热处理,产品装配、调试、检验以及油漆和包装等。

2. 工艺过程与工艺规程

工艺过程是生产过程中的主要部分。在生产过程中,直接改变原材料或毛坯的形状、尺寸、性能以及相互位置关系,使之成为成品的过程,称为工艺过程。工艺过程主要包括毛坯的制造(铸造、锻造、冲压等)、热处理、机械加工和装配。因此,工艺过程可分为机械加工工艺过程、铸造工艺过程、锻造工艺过程、焊接工艺过程、热处理工艺过程、装配工艺过程等。

通常,把合理的工艺过程编写成技术文件(机械加工工艺过程卡片、机械加工工序卡片或典型零件工艺过程卡片)用于指导生产,这类文件称为工艺规程。

3. 机械加工工艺过程及其组成

机械加工工艺过程是指用机械加工的方法改变毛坯的形状、尺寸、相对位置和性能使其成为合格零件的全过程。该过程直接决定零件及产品的质量和性能,对产品的成本、生产周期都有较大影响,是整个工艺过程的重要组成部分。

组成机械加工工艺过程的基本单元是工序。工序又是由安装、工位、工步及走刀组成的。

(1) 工序

工序是一个(或一组)工人在一台机床(或一个工作场地)上,对一个(或一组)工件连续进行的那部分工艺过程。

划分工序的主要依据是工作地点是否改变以及加工是否连续。工序内容由被加工零件的复杂程度、加工要求及生产类型来决定,同样的加工内容,可以有不同的工序安排。图16-1所示阶梯轴的加工,对单件、小批生产和大批、大量生产时的工序安排分别见表16-1和表16-2。

\begin{figure}[H]
    \centering
    % Placeholder for Figure 16-1
    \includegraphics[page=67, viewport={34.0 185.3 279.5 310.3}, clip]{12-07572-001FigRevised.pdf}
    \caption{阶梯轴}
\end{figure}

(2) 安装

安装是工件经一次装夹所完成的那一部分工序。装夹是指将工件在机床或夹具中定位、夹紧的过程。工件在一道工序中,可能有一次或几次安装。如表16-1所示的工序1中,工件在一次装夹后还需要三次调头装夹,故该工序有四次安装。工件在加工过程中,应尽量减少装夹次数,以节省装夹时间,减少装夹误差。

\begin{table}[H]
    \centering
    \caption{阶梯轴工艺过程(单件、小批生产)}
    \begin{tabular}{|c|l|c|}
    \hline
    工序号 & 工序内容 & 设备 \\
    \hline
    1 & 车端面,钻中心孔,粗车各外圆,半精车各外圆,切槽,倒角 & 车床 \\
    \hline
    2 & 铣键槽,去毛刺 & 铣床 \\
    \hline
    3 & 磨各外圆 & 磨床 \\
    \hline
    \end{tabular}
\end{table}

(3) 工位

为了减少工件的安装次数,生产中常采用多工位夹具或多轴机床,使工件在一次装夹后经过若干个不同位置顺次进行加工。此时,工件在机床上占据的每一个加工位置上所完成的那部分工序即称为一个工位。如表16-2中的工序1的铣端面、钻中心孔就是两个工位。图16-2所示即利用移动工作台或移动夹具,在一次装夹中顺次完成铣端面、钻两中心孔的两工位加工的实例。

\begin{table}[H]
    \centering
    \caption{阶梯轴工艺过程(大批、大量生产)}
    \begin{tabular}{|c|l|c|}
    \hline
    工序号 & 工序内容 & 设备 \\
    \hline
    1 & 两边同时铣端面,钻中心孔 & 铣端面、钻中心孔机床 \\
    \hline
    2 & 粗车各外圆 & 车床 \\
    \hline
    3 & 半精车各外圆,切槽,倒角 & 车床 \\
    \hline
    4 & 铣键槽 & 铣床 \\
    \hline
    5 & 去毛刺 & 钳工台 \\
    \hline
    6 & 磨外圆 & 磨床 \\
    \hline
    \end{tabular}
\end{table}

\begin{figure}[H]
    \centering
    % Placeholder for Figure 16-2
    \includegraphics[page=67, viewport={302.5 185.8 508.5 300.8}, clip]{12-07572-001FigRevised.pdf}
    \caption{多工位加工}
\end{figure}

(4) 工步

在加工表面、切削刀具、切削速度和进给量都不变的情况下所完成的那部分工序,称为工步。变化其中的一个因素就是另一个工步。

为提高生产率,用几把车刀同时加工几个表面,这也可以看作一个工步,称为复合工步。

(5) 走刀

在一个工步内,若余量较大,需分几次切削。切削刀具在加工表面上切削一次所完成的那部分工序,称为一次走刀。一个工步包括一次或几次走刀。

\subsection{生产纲领与生产类型}

1. 生产纲领

产品的生产纲领是指包括备品率和废品率在内的该产品的年产量。产品的生产纲领一般按市场需求量与本企业的生产能力而定。生产纲领是划分生产类型的依据,对工厂的生产过程及管理有着决定性的影响。

2. 生产类型

在制订机械加工工艺过程中,工序的安排不仅与零件的技术要求有关,而且与生产类型有关。生产类型是企业(或车间、工段、班组、工作地)生产专业化程度的分类。根据产品零件的大小和生产纲领的不同,机械制造生产一般可以分为单件生产、成批生产和大量生产三种不同的生产类型。成批生产根据批量的大小和产品的特征,又可分为小批生产、中批生产和大批生产。大批生产和大量生产又统称为大批量生产。
表16-3按重型零件、中型零件和轻型零件的生产纲领划分出了不同生产类型。

\begin{table}[H]
    \centering
    \caption{生产类型的划分}
    \begin{tabular}{|c|c|c|c|}
    \hline
    \multirow{2}{*}{生产类型} & \multicolumn{3}{c|}{同一零件的生产纲领} \\
    \cline{2-4}
     & \makecell{重型零件 \\ (质量$>2000$ kg)} & \makecell{中型零件 \\ (质量$100 \sim 2000$ kg)} & \makecell{轻型零件 \\ (质量$<100$ kg)} \\
    \hline
    单件生产 & $<5$ & $<10$ & $<100$ \\
    \hline
    \multirow{3}{*}{成批生产} & $5 \sim 100$ & $10 \sim 200$ & $100 \sim 500$ \\
    \cline{2-4}
     & $100 \sim 300$ & $200 \sim 500$ & $500 \sim 5000$ \\
    \cline{2-4}
     & $300 \sim 1000$ & $500 \sim 5000$ & $5000 \sim 50000$ \\
    \hline
    大量生产 & $>1000$ & $>5000$ & $>50000$ \\
    \hline
    \end{tabular}\end{table}

在拟订零件的工艺过程时,由于生产类型不同,所采用的加工方法、机床设备、刀具、夹具、量具、毛坯以及对工人的技术要求等都有很大的不同。各种生产类型的工艺特征见表16-4。

\begin{table}[H]
    \centering
    \caption{各种生产类型的工艺特征}
    \begin{tabular}{|l|l|l|l|}
    \hline
    生产要素 & 单件生产 & 成批生产 & 大量生产 \\
    \hline
    机床设备 & 通用设备 & 通用和部分专用设备 & 广泛使用高效率专用设备 \\
    \hline
    夹具 & 通用夹具 & 广泛使用专用夹具 & 广泛使用高效率专用夹具 \\
    \hline
    刀具和量具 & 一般刀具,通用量具 & 部分采用专用刀具和量具 & 使用高效率专用刀具和量具 \\
    \hline
    毛坯 & 木模铸造,自由锻 & 部分采用金属型铸造和模锻 & 金属型铸造、模锻等 \\
    \hline
    工艺规程 & 工艺路线卡片 & 简单工艺规程 & 详细工艺规程 \\
    \hline
    对工人的要求 & 需要技术熟练的工人 & 需要技术比较熟练的工人 & \makecell{调整工要求技术熟练,操作工要 \\ 求熟练程度较低} \\
    \hline
    \end{tabular}\end{table}

\section{工件的安装与夹具的基本知识}

在机械加工前,必须将工件放置在机床或夹具上,使其占有一个正确的位置,这一工艺过程称为定位;在加工过程中,为了使工件能承受切削力,并保持其正确的位置,还必须把它压紧或夹牢,这一工艺过程称为夹紧。工件从定位到夹紧的全过程称为安装。工件的安装方式对零件的加工质量、生产率和制造成本都有较大的影响。但无论采用哪种安装方法,都必须使工件在机床或夹具上正确地定位,以便保证被加工面的精度。

\subsection{工件的定位}

任何一个不受约束的物体,在空间具有六个自由度,即沿三个互相垂直坐标轴的移动(用$\vec{x}$、$\vec{y}$、$\vec{z}$表示)和绕这三个坐标轴的转动(用$\widehat{x}$、$\widehat{y}$、$\widehat{z}$表示),如图16-3所示。因此,要使物体在空间占有确定的位置,就必须约束这六个自由度。

\begin{figure}[H]
    \centering
    % Placeholder for Figure 16-3
    \includegraphics[page=68, viewport={81.0 582.8 465.5 744.3}, clip]{12-07572-001FigRevised.pdf}
    \caption{物体的六个自由度}
\end{figure}

1. 工件的六点定位原理

在机械加工中,通过用一定规律分布的六个支承点来限制工件的六个自由度,使工件在机床或夹具中的位置完全确定,这种定位原理称为工件的六点定位原理。如图16-4所示,可以设想六个支承点分布在三个互相垂直的坐标平面内。其中,三个支承点在$xOy$平面上,限制$\widehat{x}$、$\widehat{y}$和$\vec{z}$三个自由度,两个支承点在$xOz$平面上,限制$\vec{y}$和$\widehat{z}$两个自由度,最后一个支承点在$yOz$平面上,限制了$\vec{x}$一个自由度。这样就限制了工件的六个自由度。

2. 六点定位原理的应用

在实际生产中,并不要求在任何情况下都要限制工件的六个自由度,一般要根据工件的加工要求来确定工件必须限制的自由度数。工件定位只要相应的限制那些对加工精度有影响的自由度即可,由此产生了下列各种定位情况。

(1) 完全定位

工件的六个自由度都限制了的定位称为完全定位。如图16-5所示,在铣床上铣削一批工件上的沟槽时,为了保证每次安装中工件的正确位置,保证三个加工尺寸$x$、$y$、$z$,就必须限制六个自由度。

\begin{figure}[H]
    \centering
    % Placeholder for Figure 16-4
   \includegraphics[page=68, viewport={62.0 410.3 265.0 564.3}, clip]{12-07572-001FigRevised.pdf}
    \caption{六点定位简图}
\end{figure}

\begin{figure}[H]
	\centering
	% Placeholder for Figure 16-5
	\includegraphics[page=68, viewport={284.5 410.8 476.5 560.8}, clip]{12-07572-001FigRevised.pdf}
	\caption{完全定位}
\end{figure}

(2) 不完全定位

没有完全限制六个自由度,但能保证加工要求的定位,称为不完全定位。图16-6a所示为铣削一批工件的台阶面,为保证两个加工尺寸$y$和$z$,只需限制$\vec{y}$、$\vec{z}$、$\widehat{x}$、$\widehat{y}$、$\widehat{z}$五个自由度即可;图16-6b所示为磨削一批工件的顶面,为保证加工尺寸$z$,仅需限制$\widehat{x}$、$\widehat{y}$、$\vec{z}$三个自由度。

(3) 欠定位

按照加工要求应该限制的自由度而未得到限制,致使工件定位不足,这种情况称为欠定位。很显然,欠定位不能保证加工要求,因此是不允许出现的。

(4) 过定位

若定位方案中的有些定位点重复限制了同一个自由度,这样的定位称为重复定位或过定位。如图16-7所示,加工连杆大头孔时,圆柱销限制了$\vec{x}$、$\vec{y}$、$\widehat{x}$、$\widehat{y}$四个自由度,支承板限制了$\vec{z}$、$\widehat{x}$、$\widehat{y}$三个自由度,止销限制$\widehat{z}$一个自由度,其中$\widehat{x}$、$\widehat{y}$被圆柱销和支承板两个定位元件重复限制,便出现了过定位。这时,若工件定位孔与端面的垂直度误差较大,而且圆柱销与孔的间隙又很小,当圆柱销刚度好时,定位情况将如图16-7b所示,定位后工件歪斜,端面只有一点接触,压紧后必然使连杆变形。当圆柱销刚度不足时,则由于夹紧力的作用,圆柱销产生倾斜,如图16-7c所示,工件也可能产生变形。两者都会引起被加工孔的位置误差,使连杆大、小头孔的轴线不平行。
一般情况下,应当尽量避免过定位,但是如果工件定位面和夹具定位元件的尺寸、形状和位置都做得比较准确、光整,则过定位不但对工件加工面的位置尺寸影响不大,反而可以增强加工时的刚性,这时过定位是允许的。

\begin{figure}[H]
    \centering
    % Placeholder for Figure 16-6 
    \includegraphics[page=68, viewport={13.5 210.3 266.0 325.3}, clip]{12-07572-001FigRevised.pdf}
    \caption{不完全定位}
\end{figure}

\begin{figure}[H]
	\centering
	% Placeholder for Figure 16-7
	\includegraphics[page=68, viewport={267.5 207.8 534.5 381.3}, clip]{12-07572-001FigRevised.pdf}
	\caption{连杆的过定位}
\end{figure}

\subsection{夹具的基本概念}

在切削加工中,为完成某道工序,用来正确、迅速安装工件的装置称为夹具。它对保证加工精度、提高生产效率和减轻工人劳动量有很大作用。

1. 夹具的种类

按使用范围的不同,机床夹具通常可以分为如下四类。

(1) 通用夹具

通用夹具一般指已经标准化、不需特殊调整就可以用于加工不同工件的夹具。如通用三爪自定心卡盘和四爪单动卡盘、机床用平口虎钳、万能分度头、磁力工作台等,这类夹具通常作为机

床附件由专业厂制造供应。它们的适应性较强,对于充分发挥机床的技术性能、扩大机床的使用范围起着重要作用,广泛应用于单件、小批生产。

(2)专用夹具

专用夹具是指为某一零件的加工而专门设计和制造的夹具,没有通用性。利用专用夹具加工工件,既可保证加工精度,又可提高生产效率。但当产品更换或工序内容变动后,专用夹具往往不能再使用。所以,专用夹具适用于产品固定、工艺相对稳定、批量大的零件的加工。

(3)可调夹具

可调夹具是指当加工完一种工件后,经过调整或更换个别元件,即可用于另一种工件的夹具。它主要用于形状相似、尺寸相近的工件,多用于中、小批生产。

(4)组合夹具

组合夹具是指在夹具零部件完全标准化的基础上,针对不同的工件对象和加工要求拼装组合而成的夹具。此类夹具在使用完毕后可拆散重新装成其他夹具,具有组装迅速、准备周期短、能反复使用等优点,适合于单件、小批、多品种生产。

根据夹紧力来源的不同,夹具又可分成手动夹具、气动夹具、电动夹具和液压夹具等。单件、小批生产中主要使用手动夹具,而大批和大量生产中则广泛采用气动、电动或液压夹具等。

此外,按使用机床的不同,夹具还可分为车床夹具、铣床夹具、钻床夹具、镗床夹具、齿轮机床夹具和其他机床夹具等类型。

2.夹具的主要组成部分

图16-8所示是在轴上钻孔所用的一种简单的专用夹具。钻孔时,工件以外圆面定位在夹具的长V形块上,以保证所钻孔的轴线与工件轴线垂直相交。轴的端面与夹具上的挡铁接触,以保证所钻孔的轴线与工件端面的距离。工件在夹具上定位之后,拧紧夹紧机构的螺杆,将工件夹牢,即可开始钻孔。钻孔时,利用钻套定位并引导钻头。

\begin{figure}[H]
\centering
   \includegraphics[page=68, viewport={161.5 18.3 382.0 184.8}, clip]{12-07572-001FigRevised.pdf}
\caption{在轴上钻孔的夹具}
\label{fig:16-8}
\end{figure}

虽然夹具的用途和种类各不相同,结构各异,但其主要组成与上例相似,可以概括为以下几个部分。

(1)定位元件

定位元件是夹具上与工件的定位基准接触,用来确定工件正确位置的零件。工件以平面定位时,用支承钉和支承板为定位元件。图16-9所示为支承钉和支承板的标准结构。支承钉有三种形式:平头式用于已加工过的平面定位;球头式用于粗糙不平的毛坯面的定位;带齿纹的形式可以增加摩擦力,但不易清除切屑,多用于侧面定位。

\begin{figure}[H]
\centering
   \includegraphics[page=69, viewport={92.5 589.8 447.5 738.3}, clip]{12-07572-001FigRevised.pdf}
\caption{支承钉和支承板的标准结构}
\label{fig:16-9}
\end{figure}

工件以外圆柱面定位时,用V形块和定位套筒作定位元件(图16-10)。

\begin{figure}[H]
\centering
   \includegraphics[page=69, viewport={9.5 461.3 294.5 556.8}, clip]{12-07572-001FigRevised.pdf}
\caption{外圆柱面定位用的定位元件}
\label{fig:16-10}
\end{figure}

工件以孔定位时,用定位心轴(图16-11)和定位销(图16-12)作定位元件。

\begin{figure}[H]
\centering
   \includegraphics[page=69, viewport={294.5 460.3 532.0 527.3}, clip]{12-07572-001FigRevised.pdf}
\caption{定位心轴}
\label{fig:16-11}
\end{figure}

\begin{figure}[H]
	\centering
	\includegraphics[page=69, viewport={9.0 282.8 125.5 393.3}, clip]{12-07572-001FigRevised.pdf}
	\caption{定位销}
	\label{fig:16-12}
\end{figure}

(2)夹紧机构

工件定位后,将其夹紧以承受切削力等作用的机构,称为夹紧机构。例如,图16-8所示夹具上的螺杆和框架等,就是夹紧机构中的一种。常用的夹紧机构还有螺钉压板和偏心压板等(图16-13)。

\begin{figure}[H]
\centering
   \includegraphics[page=69, viewport={141.0 284.8 524.0 426.8}, clip]{12-07572-001FigRevised.pdf}
\caption{夹紧机构}
\label{fig:16-13}
\end{figure}

(3)导向元件

导向元件是用来对刀及引导刀具进入正确加工位置的元件。例如图16-8所示夹具上的钻套。其他导向元件还有导向套、对刀件等。钻套和导向套主要用在钻床夹具(习惯上称为钻模)上,对刀件主要用在铣床夹具上。

(4)夹具体和其他部分

夹具体是夹具的基本零件,用来连接并固定定位元件、夹紧机构和导向元件等,使之成为一个整体,并将夹具安装在机床上。

根据加工工件的要求,有时还在夹具上设有分度机构、导向键、平衡铁和操作件等。

工件的加工精度在很大程度上决定于夹具的精度和结构,因此整个夹具及其零件都应具有足够的精度和刚度,并且结构要紧凑,形状要简单,装卸工件和清除切屑要方便等。

\section{基准及其选择原则}

\subsection{基准的概念}
在零件的设计和制造过程中,要确定一些点、线或面的位置,必须以一些指定的点、线或面作为依据,这些作为依据的点、线或面称为基准。按照基准的不同作用,常把基准分为设计基准和工艺基准两类。

1.设计基准

在零件图上用于标注尺寸和表面相互位置的基准称为设计基准。如图16-14所示,表面2和孔3的设计基准是表面1,孔4的设计基准是孔3的中心线。

2.工艺基准

工艺基准是在制造零件和装配机器的过程中所使用的基准。工艺基准又分为定位基准、测量基准和装配基准,它们分别用于工件加工时的定位、工件的测量检验和零件的装配。

\begin{figure}[H]
\centering
    \includegraphics[page=69, viewport={39.5 113.3 184.5 251.8}, clip]{12-07572-001FigRevised.pdf}
\caption{机体示意图}
\label{fig:16-14}
\end{figure}

(1)定位基准

工件在加工过程中,用于确定工件在机床或夹具上的正确位置的基准称为定位基准。如图16-15所示,阶梯轴用三爪自定心卡盘装夹,则大端外圆的轴线为径向的定位基准,A面为加工其余两端面时保证其轴向尺寸B、C的定位基准。

事实上,工件上作为定位基准的点或线,总是由具体表面来体现的,这个表面称为定位基准面。例如,图16-15所示大端外圆的轴线,并不具体存在,而是由大端外圆的表面来体现的,所以确切地说,此例中的大端外圆表面是径向的定位基准面。

(2)测量基准

用于测量已加工表面的尺寸及各表面之间位置精度的基准称为测量基准。如图16-16所示,在偏摆仪上利用锥度心轴检验齿轮坯外圆和两个端面相对孔轴线的圆跳动时,孔的轴线即测量基准。

\begin{figure}[H]
\centering
    \includegraphics[page=69, viewport={201.0 112.3 336.0 249.3}, clip]{12-07572-001FigRevised.pdf}
\caption{定位基准}
\label{fig:16-15}
\end{figure}

\begin{figure}[H]
	\centering
	\includegraphics[page=69, viewport={355.0 110.3 503.0 248.8}, clip]{12-07572-001FigRevised.pdf}
	\caption{齿轮坯圆跳动检验}
	\label{fig:16-16}
\end{figure}


(3)装配基准

在机器装配中,用于确定零件或部件在机器中正确位置的基准称为装配基准。

\subsection{定位基准的选择}
合理选择定位基准,对保证加工精度、安排加工顺序和提高加工生产率有着重要的作用。

定位基准分为粗基准和精基准两种。用毛坯表面作为定位基准的表面为粗基准,用加工过的表面作为定位基准的表面为精基准。

1.粗基准的选择原则

(1)不加工面原则

为了保证加工面与不加工面的位置要求,应选不加工面为粗基准。如图16-17所示,以不加工的外圆表面作为粗基准,既可在一次安装中把绝大部分要加工的表面加工出来,又能够保证外圆面与内孔同轴以及端面与孔轴线垂直。

如果零件上有几个不加工的表面,则应选择与加工表面相互位置精度要求高的表面作粗基准。

(2)重要面加工原则

为了保证某加工面加工余量均匀,应选择该表面为粗基准。例如,车床床身的导轨面(图16-18)要求耐磨性好,希望在加工时只切去少而均匀的一层余量,使其表层保留均匀一致的显微组织和力学性能。若先选择导轨面作粗基准,加工床腿的底平面(图16-18a),然后再以床腿的底平面为基准加工导轨面(图16-18b),则能达此目的。

\begin{figure}[H]
\centering
   \includegraphics[page=70, viewport={39.5 592.8 154.0 704.3}, clip]{12-07572-001FigRevised.pdf}
\caption{用不加工表面作粗基准}
\label{fig:16-17}
\end{figure}

\begin{figure}[H]
	\centering
	\includegraphics[page=70, viewport={170.0 596.8 376.5 748.3}, clip]{12-07572-001FigRevised.pdf}
	\caption{床身加工的粗基准}
	\label{fig:16-18}
\end{figure}

(3)余量最小表面原则

对于所有表面都要加工的零件,应选择余量和公差最小的表面作粗基准,以避免余量不足而造成废品。图16-19所示的自由锻件,其毛坯大圆柱面A的加工余量小,小圆柱面B的加工余量大。如果选取B面为粗基准,加工A面时,可能由于A面无足够的加工余量而使零件报废。故此时只能选择加工余量最小的A面为粗基准。

(4)平整光洁原则

应选取光洁、平整、面积足够大的表面为粗基准,要避开锻造飞边和铸造浇冒口、分型面等,以保证定位正确、夹紧可靠。

(5)使用一次原则

粗基准只能在第一道工序中使用一次,不应重复使用。因为粗基准表面粗糙,在每次安装中位置不可能一致,而使加工表面的位置超差。

以上粗基准的选择原则,具体使用时要全面考虑,灵活运用,保证主要方面的要求。

2.精基准的选择原则

选择精基准时,应着重考虑的是有利于保证各加工面之间的位置精度。其选择原则如下:

(1)基准重合原则

基准重合原则即尽可能选用设计基准作为定位基准,这样可以避免定位基准与设计基准不重合而引起的定基误差。

图16-20所示的轴承座,表面1、2已精加工完毕,现预加工孔3,要求孔3的轴线与设计基准面1之间的尺寸为$A_0^{+\delta A}$。如果按图16-20b所示,用面2作为定位精基准,而面2与面1之间的尺寸有公差$\delta B$。所以,加工一批零件时,在孔3轴线与面1之间尺寸A的误差中,除了因其他原因产生的加工误差外,还包括由于定位基准与设计基准不重合而引起的定基误差。这项误差值为$\varepsilon_{\text{定基}} = \delta B$。如图16-20c所示,如果用面1作为定基基准,则因基准一致,此时定基误差$\varepsilon_{\text{定基}} = 0$。


\begin{figure}[H]
\centering
   \includegraphics[page=70, viewport={393.0 592.8 499.5 734.8}, clip]{12-07572-001FigRevised.pdf}
\caption{保证各表面应有足够的加工余量}
\label{fig:16-19}
\end{figure}

\begin{figure}[H]
\centering
    \includegraphics[page=70, viewport={84.0 453.3 453.0 562.3}, clip]{12-07572-001FigRevised.pdf}
\caption{定位误差与定位基准选择的关系}
\label{fig:16-20}
\end{figure}

由上述分析可知,选择定位基准时,应尽量使它与设计基准重合,否则必然会因基准不重合而产生定基误差,增加了加工难度,造成零件尺寸超差甚至报废。

(2)基准同一原则

对位置精度要求较高的某些表面加工时,尽可能选用同一个定位基准,这样有利于保证各加工表面的位置精度。例如,加工较精密的阶梯轴时,往往以中心孔为定位基准车削其他各表面,在精加工之前还要修研中心孔,再以中心孔定位磨削各表面,这样有利于保证各表面的位置精度,如同轴度、垂直度等。

应当指出,基准同一原则常常会带来基准不重合的问题。在这种情况下,要针对具体问题进行认真分析,在满足设计要求的前提下,决定最终选择的精基准。

(3)自为基准原则

某些精加工工序,当要求加工余量小而均匀时,选择加工表面本身作为定位基准称为自为基准原则。图16-21所示的以自为基准原则磨削床身导轨面,首先用固定在磨头上的百分表找正工件上的导轨面,然后加工导轨面,保证导轨面余量均匀,以满足对导轨面的质量要求。用浮动镗刀、浮动铰刀和珩磨、拉孔等加工孔的方法,在无心外圆磨上磨削外圆,在滚丝机上滚压螺纹等,都是自为基准的例子。

\begin{figure}[H]
\centering
    \includegraphics[page=70, viewport={62.0 339.3 260.0 416.8}, clip]{12-07572-001FigRevised.pdf}
\caption{以自为基准原则磨削床身导轨面}
\label{fig:16-21}
\end{figure}

采用自为基准原则时,不能提高加工面的位置精度,只能提高加工面本身的精度。

(4)互为基准原则

当两个表面的相互位置精度要求很高,而表面自身的尺寸和形状精度又很高时,常采用互为基准原则反复加工来达到位置精度要求。例如,车床主轴轴颈和前锥孔是主轴的主要工作表面,它们之间的同轴度要求很高,因此常以轴颈为基准粗磨锥孔,再以锥孔定位粗、精磨轴颈,最后以轴颈定位精磨内锥孔。

此外,铣削箱体类零件的上、下表面,特别是加工薄壁易变形的零件,更需要应用这一原则反复进行加工,这样可以获得较高的位置精度。

(5)便于装夹原则

应选择精度较高、安装稳定可靠的表面作精基准,而且所选的基准应使夹具结构简单,安装和加工工件方便。

在实际工作中,定位基准的选择要完全符合上述所有的原则,有时是不可能的。这就需要结合实际情况,从解决主要问题着眼,对整个加工过程中的定位通盘考虑,确定出最合理的定位方案。

\section{工艺规程的制订}
规定产品或零部件制造过程和操作方法等的工艺文件称为工艺规程。它是生产准备、生产计划、生产组织、实际加工及技术检验等的重要技术文件,是进行生产活动的基础资料。根据生产过程中工艺性质的不同,又可以分为毛坯制造、机械加工、热处理及装配等不同的工艺规程。本节仅介绍拟订机械加工工艺规程的一些基本问题。

\subsection{制订工艺规程的内容和要求}
零件的工艺规程就是零件加工的方法和步骤。其内容包括排列加工工序(包括毛坯制造、切削加工、热处理和检验工序),确定各工序所用的机床、装夹方法、加工方法、测量方法、刀具、夹具、量具、加工余量、切削用量和工时定额等。将这些内容用工艺文件的形式表示出来,就是机械加工工艺规程,即通常所说的机械加工工艺卡片。

制订的零件加工工艺必须确保零件的全部技术要求,并使生产效率高、生产成本低、劳动生产条件好。在制订零件加工工艺时,尤其对较复杂零件的加工工艺,要根据客观生产条件,经过反复实践、反复修改,才能使之合理与完善。

\subsection{制订工艺规程时的一般步骤}
1.零件的工艺分析

首先要熟悉整个产品(如整台机器)的用途、性能和工作条件,结合装配图了解零件在产品中的位置、作用、装配关系以及其精度等技术要求对产品质量和使用性能的影响,然后从加工的角度对零件进行工艺分析。零件工艺分析的主要内容如下。

(1)检查零件的图样是否完整和正确

如视图是否足够、正确,所标注的尺寸、公差、表面粗糙度和技术要求等是否齐全、合理,并要分析零件主要表面的精度、表面质量和技术要求等在现有的生产条件下能否达到,以便采取适当的措施。

(2)审查零件材料的选择是否恰当

零件材料的选择应立足于国内,尽量采用我国资源丰富的材料,不要轻易选用贵重材料。另外,还要分析所选的材料会不会使工艺变得困难和复杂。

(3)审查零件结构的工艺性

零件的结构是否符合工艺性一般原则的要求,现有生产条件能否经济、高效、合格地加工出来。如果发现问题,应与有关设计人员共同研究,按规定程序对原图样进行必要的修改与补充。

2.毛坯的选择

机械加工的加工质量、生产效率和经济效益,在很大程度上取决于所选用的工件毛坯。常用的毛坯类型有型材、铸件、锻件、冲压件和焊件等。毛坯的选择应根据零件的作用、受力情况、生产批量及工厂现场条件综合考虑。例如,一般轴、套类零件都要承受动载荷,故应选用钢件型材作毛坯,对于某些要求高(一般有热处理要求)的零件,通常要经过锻造后,以锻件作为毛坯。对于箱体类零件,其特点是形状、结构复杂,大多数情况下承受静载荷以及起容纳作用,一般选择铸件作毛坯。若零件数量很少,如试制产品,则可以选择焊件作箱体零件的毛坯。这样,既可以节约切削加工的工时,还可以节约大量材料。

3.定位基准的选择

在加工过程中,合理确定定位基准面,对保证零件的技术要求和工序的安排有着决定性的影响。粗基准与精基准的选择应该遵循上一节所述的原则。一般在选择主要表面加工方法的同时就要确定其定位基准面。下面介绍三类典型零件的定位基准选择。

(1)阶梯轴类零件

常选择两端中心孔作为定位基面。如图16-22所示,采用双顶尖装夹,车削或磨削外圆、螺纹和轴肩端面,这样能较好地保证各外圆、螺纹的同轴度(或径向圆跳动)和轴肩对轴线的垂直度(或端面圆跳动)。在热处理后或磨削前,一般要研磨中心孔,以提高中心孔的定位精度。

\begin{figure}[H]
\centering
   \includegraphics[page=70, viewport={278.5 342.8 477.0 423.8}, clip]{12-07572-001FigRevised.pdf}
\caption{阶梯轴的加工}
\label{fig:16-22}
\end{figure}

(2)盘套类零件

一般以轴线部位的孔作为定位基面,采用心轴装夹,如图16-23a所示。以孔为定位基面车削或磨削其他表面能较好地保证各外圆和端面对孔轴线的圆跳动要求。值得注意的是,如果零件结构允许,常在一次装夹中完成孔及与其有关表面的精加工,这样不仅可获得较高的位置精度,而且加工十分方便,如图16-23b所示。

\begin{figure}[H]
\centering
    \includegraphics[page=70, viewport={65.5 172.3 476.0 321.8}, clip]{12-07572-001FigRevised.pdf}
\caption{盘套类零件定位精基准的选择}
\label{fig:16-23}
\end{figure}

(3)机架箱体类零件

该类零件形状、结构复杂,除有尺寸精度要求外,一般孔的轴线相对于底面(装配基准面)有位置精度要求。因此,箱体类零件多采用主要的装配基准面(一般为最大的底平面)作为定位精基准。

图16-24所示为变速箱壳体的定位简图。该箱体的装配基准面为底面 $P$，轴装在孔 $D$ 内。因此，可以在 $z$ 方向选用底平面 $P$ 作为定位精基准（夹具上可以用两条经过精加工的窄长支承块代替基准面），在 $y$ 方向选用侧平面 $M$ 作为定位精基准（夹具上可以用两个可调支承来实现），采用压板、螺栓等附件装夹即可进行加工。这样定位可以保证这两个加工面、孔的轴线与底面的平行度等位置精度要求。

\begin{figure}[H]
	\centering
	\includegraphics[page=70, viewport={13.0 21.3 242.0 148.3}, clip]{12-07572-001FigRevised.pdf}
	\caption{盘套类零件定位精基准的选择}
	\label{fig:16-24}
\end{figure}

4. 工艺路线的制订

制订工艺路线，就是把加工工件所需的各个工序按顺序合理地排列出来，这是制订零件工艺规程的核心，主要包括以下内容。

（1）确定加工方案

确定加工方案即根据零件每个加工表面（特别是主要表面）的技术要求，选择较合理的加工方案（或方法）。常见典型表面的加工方案（或方法），可参照第 14 章有关内容来确定。在确定加工方案（或方法）时，除了表面的技术要求外，还要考虑零件的生产类型、材料性能以及本单位现有的加工条件等。若是大批、大量生产，应采用高效率设备与专用设备，如使用各种类型的组合机床，加工平面和孔采用拉床，加工轴类零件采用多刀车床、多轴车床或仿形车床。对于中、小批生产，则应选用通用设备，或采用数控机床以提高加工质量和生产率。在选择加工方法和加工设备时，应进行方案比较，选择加工成本最低的方案。

（2）加工阶段的划分

零件的机械加工一般要经过粗加工、半精加工和精加工几个阶段才能完成，对于特别精密的零件还要进行精密加工。

粗加工阶段应高效率地切除各加工表面的大部分余量，并为半精加工提供定位基准。
半精加工阶段的任务是完成次要表面的加工，使之达到图样要求，并为主要表面的精加工做好准备（如消除热处理变形、准备修整精加工用的定位基准等）。

精加工阶段主要是保证零件的尺寸精度、形状和位置精度以及表面粗糙度。

精密加工阶段的任务是在改善零件的表面质量的同时，使加工表面的尺寸精度和形状精度达到图样要求。

划分加工阶段利于保证加工质量、合理使用设备、及时发现毛坯缺陷、避免浪费工时，同时还可避免精加工后的表面被碰伤。

划分加工阶段应结合具体生产情况，避免生搬硬套。例如，对于刚性好、加工精度要求又不高或生产批量较小的工件，则不一定划分加工阶段，以免增加工序，延长生产周期，导致产品成本提高；又如某些零件在精加工后装配时，才安排配钻之类的粗加工工序。在组合机床和自动机床上加工零件，也常常不划分加工阶段。

（3）工序集中与分散

选定加工方案和划分加工阶段后，就要确定工序的数目，即工序的集中或分散的问题。

工序集中是使每道工序所加工的表面数量尽量多，而使零件加工总的工序数目减少，所用机床和夹具的数量也相应减少。工序集中有利于保证零件各表面之间的相互位置精度，简化生产组织，节省辅助时间。工序分散则是减少每道工序的加工内容，增加总的工序数目，其所用设备较简单，对工人的技术水平要求较低，且生产准备工作量小，容易适应产品的变换。

在拟订工艺路线时，工序集中或分散的程度主要取决于生产类型、零件的结构特点及技术要求。在单件、小批生产时，由于使用通用机床、通用夹具和量具，工序安排通常尽可能集中；当产品固定且产量很大时，由于有条件采用各种专用机床和专用工装，则常常采用工序分散的原则；在重型机械制造厂，由于零件笨重，安装搬运困难，应尽可能实行工序集中。由于工序集中的优点较多，以及数控机床、柔性制造单元和柔性制造系统等的发展，现代生产多趋于工序集中。

（4）安排加工顺序
安排加工顺序即较合理地安排切削加工工序、划线工序、热处理工序、检验工序和其他辅助工序的先后次序。

1）切削加工工序的安排应遵循的原则

 \textcircled{1}基准先行。首先加工用作精基准的表面，以便为其他表面的加工提供可靠的基准表面。这是确定加工顺序的一个重要原则。

 \textcircled{2}先主后次。主要表面一般是指零件上的工作表面、装配基面等，它们的技术要求较高，加工工作量较大，应先安排加工。其他次要表面如非工作面、键槽、螺钉孔、螺纹孔等，一般可穿插在主要表面加工工序之间或稍后进行加工，但应安排在主要表面之后精加工之前。

 \textcircled{3}先粗后精。粗加工时切削力大，切削热多，工件变形大。此外，被加工零件的内应力会由于表面切掉一层金属而重新分布，也会使工件产生变形，先粗后精有利于保证加工精度和提高生产率。

 \textcircled{4} 先面后孔。底座、箱体之类零件的加工，应先加工平面后加工孔，因为这些零件上平面的轮廓尺寸较大，用平面定位比较稳定可靠，易于保证孔与平面之间的位置精度。

2）划线工序的安排

形状较复杂的铸件、锻件和焊件等，在单件、小批生产中，为了给安装和加工提供依据，一般在切削加工之前须安排划线工序。有时为了加工的需要，在切削加工工序之间，可能还要进行第二次或多次划线。但是在大批、大量生产中，由于采用专用夹具等，可免去划线工序。

3）热处理工序的安排

根据热处理工序的性质和作用不同，一般可以分为预备热处理、时效处理和最终热处理。

 \textcircled{1}预备热处理是指为改善金属的组织和切削加工性而进行的热处理，如退火、正火等，一般安排在切削加工之前。调质也可以作为预备热处理，但若是以提高材料的力学性能为主要目的，则应放在粗加工之后、精加工之前进行。

 \textcircled{2}时效处理是指在毛坯制造和切削加工的过程中，都会有内应力残留在工件内，为了消除它对加工精度的影响，需要进行的热处理。对于大而结构复杂的铸件，或者精度要求很高的非铸件类工件，需在粗加工前、后各安排一次人工时效。对于一般铸件，只需在粗加工前或后进行一次人工时效。对于要求不高的零件，为了减少工件的往返搬运，有时仅在毛坯铸造以后安排一次时效处理。

 \textcircled{3}最终热处理是指为提高零件表层硬度和强度而进行的热处理，如淬火、渗氮等，一般安排在工艺过程的后期。淬火一般安排在切削加工之后、磨削之前，渗氮则安排在粗磨和精磨之间。应注意在氮化之前要进行调质。

热处理工序在加工工序中的安排如图 16-25 所示。

\begin{figure}[H]
    \centering
    % Placeholder for Figure 16-25
    \includegraphics[page=70, viewport={253.0 18.8 524.0 135.3}, clip]{12-07572-001FigRevised.pdf}
    \caption{热处理工序在加工工序中的安排}
\end{figure}

4）检验工序的安排

检验工序是保证产品质量的有效措施之一，是工艺过程中不可缺少的内容。除了加工过程中操作者的自检外，在下列情况下还应安排检验工序：粗加工阶段之后；关键工序前后；特种检验（如磁力探伤、密封性试验、动平衡试验等）之前；从一个车间转到另一车间加工之前；全部加工结束之后。

5）其他辅助工序的安排

 \textcircled{1}零件的表面处理，如电镀、发蓝、油漆等，一般均安排在工艺过程的最后。但有些大型铸件的内腔不加工面，常在加工之前先涂防锈油漆等。

 \textcircled{2}去毛刺、倒棱边、去磁、清洗等，应适当穿插在工艺过程中进行。这些辅助工序不能忽视，否则会影响装配工作，妨碍机器的正常运行。

5. 工艺文件的编制

前面分析的全部内容，表面看起来十分复杂，如果将它们系统化、条理化，以图表或文字的形式写成工艺文件，便会十分清晰。工艺文件的种类和形式多种多样，其繁简程度也有很大不同，要视生产类型而定，通常有如下几种。

（1）机械加工工艺过程卡片

机械加工工艺过程卡片用于单件、小批生产，格式如表 16-5 所示，它的主要作用是概略地说明机械加工的工艺路线。实际生产中，工艺过程卡片内容的繁简程度也不一样，最简单的只列出各工序的名称和顺序，较详细的则附有主要工序的加工简图等。

（2）机械加工工序卡片

大批、大量生产中，要求工艺文件更加完整和详细，每个零件的各加工工序都要有工序卡片。它是针对某一工序编制的，要画出该工序的工序图，以表示本工序完成后工件的形状、尺寸及其技术要求，还要表示出工件的装夹方式、刀具的形状及其位置等，如表 16-6 所示。

（3）（标准零件或典型零件）工艺过程卡片

（标准零件或典型零件）工艺过程卡片主要用于标准零件和典型零件的成批生产，它比机械加工工艺过程卡片详细，比工序卡片简单且较灵活，是介于两者之间的一种格式（表 16-7）。该卡片既要说明工艺路线，又要说明各工序的主要内容。


总之，机械加工工艺过程包括的内容很广，它的制订是一件十分灵活的工作，必须根据生产实际和现场条件进行综合分析，才能编制出既保证质量，又成本低、效率高、较为合理的工艺规程。

\section{典型零件的工艺过程}
常见的典型零件有三类，即轴类零件、盘套类零件和机架箱体类零件。本节的工艺过程主要围绕上述三类零件进行。

\subsection{轴类零件}
轴类零件是一种常见的典型零件，按其结构特点可以分为简单轴（如光轴）、阶梯轴、空心轴（如车床主轴）和异形轴（如曲轴）等。其主要表面为外圆面、轴肩和端面，某些轴类零件还有内圆面和键槽、退刀槽、螺纹等其他表面。外圆面主要用于安装轴承和轮系（包括带轮、齿轮、链轮、凸轮、槽轮等）。轴肩的作用是使上述零件在轴上轴向定位。轴类零件通过轴上安装的零件起支承、传递运动和扭矩的作用。其加工方法主要有车削、磨削、钻削和铣削等。

现以如图 16-26 所示的传动轴为例，编制其单件、小批生产的加工工艺过程。

\begin{figure}[H]
    \centering
    % Placeholder for Figure 16-26
    \includegraphics[page=71, viewport={90.5 587.8 452.5 750.8}, clip]{12-07572-001FigRevised.pdf}
    \caption{传动轴}
\end{figure}

1. 技术要求分析

本零件的轴颈 $\phi 24\text{h}6$ 和 $\phi 16\text{h}6$ 分别装在箱体的两个孔中，是孔的装配对象，轴通过螺纹 M10 和孔 $\phi 10$ 紧固在箱体上。$\phi 16\text{h}6$ 相对于 $\phi 24\text{h}6$ 有 $0.02\text{ mm}$ 的圆跳动公差要求。轴颈 $\phi 20\text{h}6$ 处是用来安装滚动轴承的，轴上装有齿轮，轴是支承齿轮的。轴中间对称地加工出相距 $22\text{ mm}$ 的两个平行平面，这是为了将轴安装在箱体上时，为采用扳手调整而设计的工艺结构。该轴的材料为 45 钢，调质，硬度为 $235\text{ HBW}$。

2. 工艺分析

如前所述，轴颈 $\phi 24\text{h}6$ 和 $\phi 16\text{h}6$ 处用来装在箱体中，轴颈 $\phi 20\text{h}6$ 处用来装滚动轴承，所以这三个外圆面都是配合表面，其精度要求较高，表面粗糙度 $Ra$ 值要求较小，且两端轴颈对轴线有径向圆跳动要求。因此，这三个表面是该轴的重要加工面。此外，虽然螺纹 M10 未注精度要求，但根据径向圆跳动的要求，说明螺孔的轴线与 $\phi 16\text{h}6$ 的轴线有同轴度要求，而且螺纹的精度（如螺纹中径等各项指标）必须在规定的公差范围内。这两个表面在轴的两端，一般情况下难以在一次安装中全部完成。所以，加工螺纹孔时应特别注意选用精基准定位。

3. 基准选择

（1）以圆钢外圆面为粗基准，粗车端面并钻中心孔。

（2）为保证各外圆面的位置精度，以轴两端的中心孔为定位精基准，这样满足了基准重合和基准同一原则。

（3）调质后，以外圆面定位，精车两端面并修整中心孔。

（4）以修整的两中心孔作为半精车和磨削的定位精基准，这样满足了互为基准原则。

4. 工艺过程

单件、小批生产如图 16-26 所示传动轴的机械加工工艺过程详见表16-8。

\begin{longtable}{|c|c|m{6cm}|m{6cm}|c|}
    \caption{单件、小批生产传动轴的机械加工工艺过程}\\
        \hline
        \textbf{工序号} & \textbf{工序名称} & \textbf{工序内容} & \textbf{加工简图} & \textbf{设备} \\
        \hline
        5 & 准备 & 45 圆钢下料 $\phi 30 \times 150$ & & 锯床 \\
        \hline
        10 & 粗车 & 
       
        \textcircled{1}粗车一面，钻中心孔
        
        \textcircled{2}粗车另一端面，至长 145，钻中心孔
        
        \textcircled{3}粗车一端外圆，分别至 $\phi 22.5 \times 36$, $\phi 18.5 \times 42$
        
        \textcircled{4}粗车另一端外圆至 $\phi 26.5$
        
        &\includegraphics[page=74, viewport={13.5 126.8 176.5 306.3}, clip]{12-07572-001FigRevised.pdf}& 卧式车床 \\
        \hline
        15 & 热处理 & 调质 235，硬度为 HBW && \\
        \hline
        20 & 半精车 & \begin{tabular}[t]{p{4.2cm}}\textcircled{1}精车 $\phi 18.5$ 端面,修整中心孔\\\textcircled{2}精车另一端面,至长 143,钻 M10\\螺纹底孔 $\phi 8.5 \times 25$,孔口车 $60^\circ$ 锥面\\\textcircled{3}半精车一端外圆至 $\phi 24.4_{0}^{+0.1}$\\\textcircled{4}半精车另一端外圆至 $\phi 16.4_{0}^{+0.1}$,\\$\phi 20.4_{0}^{+0.1} \times (36 \pm 1)$ 并保证 $\phi 24.4_{0}^{+0.1} \times 66$\\\textcircled{5}切槽至 $2 \times 0.5$\end{tabular} & \includegraphics[page=74, viewport={190.0 125.3 350.5 317.3}, clip]{12-07572-001FigRevised.pdf} & \makecell{卧式\\车床} \\
        \hline
        25 & 铣削 & \begin{tabular}[t]{p{4.2cm}}按加工简图所注尺寸铣扁,保证尺寸\\22,并去毛刺\end{tabular} & \includegraphics[page=74, viewport={362.0 129.3 532.5 190.8}, clip]{12-07572-001FigRevised.pdf}& \makecell{立式\\铣床} \\
        \hline
        30 & 钻孔 & \begin{tabular}[t]{p{4.2cm}}按加工简图所注尺寸钻 $\phi 10$、$\phi 3.5$ 孔\\深 3,两孔成形\end{tabular} & \includegraphics[page=74, viewport={104.5 20.3 276.0 96.8}, clip]{12-07572-001FigRevised.pdf} & \makecell{立式\\钻床} \\
        \hline
        35 & 磨削 & \begin{tabular}[t]{p{4.2cm}}磨各外圆面,表面粗糙度 $Ra$ 值为\\$0.8\ \mu\text{m}$;靠磨端面,表面粗糙度 $Ra$ 值\\为 $3.2\ \mu\text{m}$\end{tabular} & \includegraphics[page=74, viewport={291.5 18.3 431.0 78.8}, clip]{12-07572-001FigRevised.pdf} & \makecell{外圆\\磨床} \\
        \hline
        40 & 钳工 & 攻 M10 螺纹,去毛刺(切勿伤及表面) & \includegraphics[page=75, viewport={31.5 273.8 189.0 333.3}, clip]{12-07572-001FigRevised.pdf} & \\
        \hline
        45 & 检验 & 按图样检验 & & \\
        \hline
\end{longtable}

\subsection{盘套类零件}

盘套类零件主要由外圆面、内圆面、端面和沟槽组成,其特征是径向尺寸大于轴向尺寸,如联轴器、法兰盘等。一般在法兰盘的底板上设计出均匀分布的、用于连接的孔。根据需要还可设计出螺纹、销孔等结构。现以常用的如图 16-27 所示的法兰端盖为例,说明盘套类零件的机械加工工艺过程。

1. 技术要求分析

本零件的底板为 $80_{-1}^{0}\ \text{mm} \times 80_{-1}^{0}\ \text{mm}$ 的正方形,它的周边不需要加工,其精度直接由铸造保证。底板上有四个均匀分布的通孔 $\phi 9$,其作用是将法兰盘与其他零件相连接,外圆面 $\phi 60\text{d}11$ 是与其他零件相配合的基孔制的轴,内圆面 $\phi 47\text{J}8$ 是与其他零件相配合的基轴制的孔。它们的表面粗糙度 $Ra$ 值均为 $3.2\ \mu\text{m}$。本零件的精度要求较低,可以采用一般加工工艺完成。零件材料为 HT100。

2. 工艺分析

外圆面与正方形底板的相对位置可由铸造时采用整模造型保证,这样不会产生大的偏差。由于本零件精度要求较低,只要选择好定位基准,则只需采用粗车$\to$半精车,即可完成车削加工。因此,可以采用铸造$\to$车削$\to$划线钻孔$\to$检验的工艺路线。


3. 基准选择

(1) 以 $\phi 60\text{d}11$ 的外圆面为粗基准,加工底板的底平面。

(2) 以底板加工好的底平面和不需加工的侧面为精基准(实质上是以 $\phi 60\text{d}11$ 的轴线为精基准),即以底平面定位、四爪单动卡盘夹紧的方式将工序集中,在一次安装中把所有需要车削加工的表面加工出来,这样符合了基准同一原则。

(3) 以 $\phi 60\text{d}11$ 的外圆面的轴线为基准划线,找正孔 $4 \times \phi 9$、$2 \times \phi 2$ 的中心位置,即可钻出上述各小孔。

4. 工艺过程

单件、小批生产法兰端盖的机械加工工艺过程详见表 16-9。

\begin{figure}[H]
    \centering
    \includegraphics[page=71, viewport={16.0 340.3 330.0 570.3}, clip]{12-07572-001FigRevised.pdf}
    \caption{法兰端盖}
\end{figure}

\begin{longtable}[H]{|c|c|m{4.5cm}|m{4.5cm}|c|}

\caption{单件、小批生产法兰端盖的机械加工工艺过程}\\


\hline
\makecell{工序\\号} & \makecell{工序\\名称} & \makecell{工序内容} & \makecell{加工简图} & \makecell{设备} \\
\hline
5 & 铸造 & 铸造毛坯,尺寸如右图所示,清理铸件 &\includegraphics[page=75, viewport={203.5 272.8 310.5 406.3}, clip]{12-07572-001FigRevised.pdf} & \\



\hline

10 & 车削 & \begin{tabular}[t]{m{4.2cm}}\textcircled{1}车 $80 \times 80$ 底平面,保证总长尺寸 26\\\textcircled{2}车 $\phi 60$ 端面,保证尺寸 $23_{-0.5}^{0}$\\\textcircled{3}车 $\phi 60\text{d}11$ 及 $80 \times 80$ 底板的上端\\面,保证尺寸 $15_{0}^{+0.3}$\\\textcircled{4}钻 $\phi 20$ 通孔\\\textcircled{5}镗 $\phi 20$ 孔至 $\phi 22_{0}^{+0.5}$\\\textcircled{6}镗 $\phi 22_{0}^{+0.5}$ 至 $\phi 40_{0}^{+0.5}$,保证尺寸 3\\\textcircled{7}镗 $\phi 47\text{J}8$,保证 $15.5_{0}^{+0.21}$\\\textcircled{8}倒角 $1 \times 45^\circ$\end{tabular} & \includegraphics[page=75, viewport={318.0 274.8 423.0 729.8}, clip]{12-07572-001FigRevised.pdf} & \makecell{卧式\\车床} \\
\hline
15 & 钳工 & 按图样要求划 $4 \times \phi 9$ 及 $2 \times \phi 2$ 孔的加工线 & & 平台 \\
\hline


20 & 钻孔 & 根据划线找正、安装,钻 $4 \times \phi 9$ 及 $2 \times \phi 2$ & \includegraphics[page=75, viewport={435.0 273.3 510.0 393.8}, clip]{12-07572-001FigRevised.pdf} & \makecell{立式\\钻床} \\
\hline
25 & 检验 & 按图样要求,检测零件 & & \\
\hline

\end{longtable}
\subsection{机架箱体类零件}

机架箱体类零件是机器的基础零件,用以支承和装配轴系零件,并使各零件之间保证正确的位置关系,以满足机器的工作性能要求。因此,机架箱体类零件的加工质量对机器的质量影响很大。现以如图 16-28 所示的零件加工过程为例,介绍一般机架箱体类零件的工艺过程。

1. 技术要求分析

(1) 箱座底面与对合面的平行度在 $1000\ \text{mm}$ 长度内不大于 $0.5\ \text{mm}$。

(2) 对合面加工后,其表面不能有条纹、划痕及毛刺。对合面对合间隙不大于 $0.03\ \text{mm}$。

(3) 3 个主要孔(轴承孔)的轴线必须保持在对合面内,其偏差不大于 $\pm 0.2\ \text{mm}$。

(4) 主要孔的距离误差应保持在 $\pm 0.03 \sim \pm 0.05\ \text{mm}$ 的范围内。

(5) 主要孔的尺寸公差等级为 IT6,其圆度与圆柱度误差不超过其孔径公差的 1/2。

(6) 加工后,箱体内部需要清理。

(7) 工件材料为 HT150,毛坯为铸件(去应力退火)。

\begin{figure}[H]
	\centering
	\includegraphics[page=71, viewport={350.5 342.3 528.5 440.3}, clip]{12-07572-001FigRevised.pdf}
	\caption{减速箱}
\end{figure}

2. 工艺分析

减速箱主要有以下加工表面。

(1) 箱座的底平面和对合面、箱盖的对合面和顶部方孔的端面,可采用龙门铣床或龙门刨床加工。

(2) 3 个轴承孔及孔内环槽,可采用坐标镗床镗孔。

3. 基准选择

(1) 粗基准的选择

为了保证对合面的加工精度和表面完整性,选择对合面法兰的不加工面为粗基准加工对合面。

(2) 精基准的选择

箱座的对合面与底面互为精基准,箱盖的对合面与顶面互为精基准。

4. 工艺过程

大批、大量生产减速箱的机械加工工艺过程见表 16-10。

\begin{longtable}[H]{|c|c|m{4.5cm}|m{8cm}|c|}

\caption{大批、大量生产减速箱的机械加工工艺过程}
\\

\hline
\makecell{工序\\号} & \makecell{工序\\名称} & \makecell{工序内容} & \makecell{加工简图} & \makecell{加工\\设备} \\
\hline
5 & 铸造 & & & \\
\hline
10 & 热处理 & 去应力退火 & & \\
\hline
15 & 刨削 & 粗刨对合面 &\includegraphics[page=75, viewport={42.5 118.8 262.0 172.8}, clip]{12-07572-001FigRevised.pdf} & \makecell{龙门\\刨床} \\
\hline
20 & 刨削 & \begin{tabular}[t]{p{4.2cm}}\textcircled{1}粗、精刨箱座的底面及\\两侧面\\\textcircled{2}粗、精刨箱盖的方孔端\\面及两侧面\end{tabular} & \includegraphics[page=75, viewport={283.5 114.8 501.5 241.3}, clip]{12-07572-001FigRevised.pdf} & \makecell{龙门\\刨床} \\
\hline
25 & 刨削 & 精刨对合面 & \includegraphics[page=76, viewport={81.0 660.3 299.0 722.8}, clip]{12-07572-001FigRevised.pdf} & \makecell{龙门\\刨床} \\
\hline
30 & 钻削 & \begin{tabular}[t]{p{4.2cm}}\textcircled{1}钻连接孔\\\textcircled{2}钻螺纹孔\\\textcircled{3} 钻销孔\end{tabular} & & \makecell{摇臂\\钻床} \\
\hline
35 & 钳工 & \begin{tabular}[t]{p{4.2cm}}\textcircled{1}攻螺纹孔\\\textcircled{2}铰销孔\\\textcircled{3} 连接箱体\end{tabular} & & \\
\hline




40 & 镗孔 & \begin{tabular}[t]{m{4.2cm}}\textcircled{1}粗镗 3 个主要孔\\\textcircled{2}半精镗 3 个主要孔\\\textcircled{3}精镗 3 个主要孔\\\textcircled{4}精细镗 3 个主要孔\end{tabular} & \includegraphics[page=76, viewport={315.0 661.3 462.5 744.3}, clip]{12-07572-001FigRevised.pdf} & 镗床 \\
\hline
45 & 终检 & 按图样检验 & & \\
\hline
\end{longtable}
\subsection{飞机起落架减振筒零件}

飞机起落架减振筒(图 16-29)是吸收飞机着陆时撞击能量的零件。可使作用于飞机机身上的撞击载荷足够小,以保证飞机着陆和滑行安全。该零件工作环境恶劣,受高温、高压、气体腐蚀及机械磨损等作用,其内表面易形成麻点、气孔、拉伤、局部损伤等,是报废率较高的零件,故加工精度要求很高。

\begin{figure}[H]
    \centering
    \includegraphics[page=71, viewport={90.5 37.8 447.5 324.3}, clip]{12-07572-001FigRevised.pdf}
    \caption{某飞机起落架减振筒零件图}
\end{figure}

减振筒全长为 $716\ \text{mm}$,外径为 $\phi 88_{-0.3}^{-0.1}\ \text{mm}$,内径为 $\phi 82_{0}^{+0.054}\ \text{mm}$,表面粗糙度 $Ra$ 值为 $0.4\ \mu\text{m}$,壁厚为 $3\ \text{mm}$。其内孔的深径比达 8.73,属深孔;其壁厚与内径比达 $1:27$,属薄壁件。减振筒集深孔、薄壁、复杂型腔等特点于一身,且要求其加工精度高、表面粗糙度值小。另外,减振筒的内孔精加工(精车、磨削)都在焊接底盖后的半封闭状态下进行,故其属难加工件。

减振筒零件加工所用设备主要有 CW6140 车床、M115A 磨床,需进行相应的工装设计。辅助加工设备有 G72 号锯机、Z32K 钻床、X62W 铣床等。

减振筒筒体毛坯选用尺寸为 $\phi 100\ \text{mm} \times \phi 68\ \text{mm} \times 720\ \text{mm}$ 的 30CrMnSiA 热轧管,减振筒底毛坯选用 $\phi 85\ \text{mm} \times 25\ \text{mm}$ 的 30CrMnSiA 棒料。

由于筒形件内孔较长,刀具车削时间长,刀片很容易磨损,故刀头材料选择综合力学性能好的钨钴钛类 M30 硬质合金刀具。

综上分析,减振筒的加工工艺过程见表 16-11。

\begin{longtable}{|c|c|m{2cm}|m{10cm}|m{1cm}|}
\caption{飞机起落架减振筒加工工艺过程}\\
\hline
\makecell{工序号} & \makecell{工序\\名称} & \makecell{工序内容} & \makecell{加工简图} & \makecell{加工\\设备} \\
\hline
1 & 下料 & 毛坯(热轧毛坯管料) & & \\
\hline
5 & 车削 & 车端面锪 $60^\circ$ 角 & \includegraphics[page=73, viewport={13.0 640.3 267.0 706.8}, clip]{12-07572-001FigRevised.pdf} & 车床 \\
\hline
10 & 车削 & 粗车外圆至 $\phi 92$,长为 590 &\includegraphics[page=73, viewport={269.5 639.3 530.0 729.3}, clip]{12-07572-001FigRevised.pdf} & 车床 \\
\hline
15 & 车削 & 调头,车端面,保证全长 694.5,锪 $60^\circ$ 角 & \includegraphics[page=73, viewport={135.0 537.3 404.0 607.8}, clip]{12-07572-001FigRevised.pdf} & 车床 \\
\hline
20 & 车削 & 粗车外圆至 $\phi 92$,长为 100 & \includegraphics[page=73, viewport={133.5 423.8 409.0 501.8}, clip]{12-07572-001FigRevised.pdf} & 车床 \\
\hline
25 & 镗孔 & 调头，采用特制镗孔装置分粗镗、半精镗通孔直径至 $\phi 72$，然后扩孔至 $\phi 78.5$，深度为 684.5 & \includegraphics[page=73, viewport={133.5 314.8 410.5 388.8}, clip]{12-07572-001FigRevised.pdf}& 车床 \\
\hline
30 & 车削 & 调头，车退刀槽；再车外圆至 $\phi 86_{+0.023}^{+0.045}$，长 100、倒角 $2 \times 45^\circ$ &\includegraphics[page=73, viewport={130.5 151.3 415.0 282.3}, clip]{12-07572-001FigRevised.pdf} & 车床 \\
\hline
35 & 镗孔 & 调头，镗孔至 $\phi 84$，长 29 &\includegraphics[page=73, viewport={123.0 36.8 412.5 114.8}, clip]{12-07572-001FigRevised.pdf} & \\
\hline
40 & 焊接 & 调头，使用 30CrMnSiA 焊条，采用直流反接法，将加工好的底盖焊接在工件上 &\includegraphics[page=74, viewport={131.5 680.3 407.5 746.3}, clip]{12-07572-001FigRevised.pdf} & 焊接台 \\
\hline
45 & 热处理 & 淬火 + 回火，抗拉强度 $R_m = 1300 \pm 100 \text{ MPa}$ & & 井式炉 \\
\hline
50 & 车削 & 半精车全部外圆后调头，将内孔扩至 $\phi 81.5$，长 684.5；车内螺纹 M85$\times$1.5，保证螺纹长 29 & \includegraphics[page=74, viewport={132.0 577.3 416.0 663.3}, clip]{12-07572-001FigRevised.pdf}& 车床 \\
\hline
55 & 车削 & 车外圆至 $\phi 88_{-0.3}^{-0.1}$；其余达到要求尺寸 & \includegraphics[page=74, viewport={134.5 457.3 411.5 544.8}, clip]{12-07572-001FigRevised.pdf}& 车床 \\
\hline
60 & 磨削 & 粗、精磨内径，达到 $\phi 82_{0}^{+0.054}$，表面粗糙度 $Ra$ 值为 $0.4 \text{ \textmu m}$ 的技术要求 &\includegraphics[page=74, viewport={129.0 346.3 413.0 422.8}, clip]{12-07572-001FigRevised.pdf} & 内孔磨削工装 \\
\hline
65 & 电镀 & 除内表面及直径为 $\phi 86_{+0.023}^{+0.045}$ 的表面外，所有的表面镀锌 & & 电镀池 \\
\hline
70 & 检验 & 填充压力为 $5690 \sim 6080 \text{ kPa}$ 的氮气，放置 $8 \sim 12 \text{ h}$，要求无明显压降即为合格 & & \\
\hline
\end{longtable}

\subsection{航空发动机压气机转子叶片}

压气机转子叶片(图 16-30)位于发动机前端，工作温度比较低，其结构形式看似较为简单，但其横截面为变曲面；压气机转子叶片尺寸变化范围较大(叶片长度为 30 mm 到几百毫米)，毛坯余量状态不同。因此，机械加工工艺过程各不相同，一般分为两类。普通模锻毛坯主要加工工艺过程为：建立加工基准(包括加工进气和排气边、粗加工榫头等)、粗加工叶型、半精加工叶型、精加工榫头、精加工叶身、加工叶尖(叶冠)；精锻毛坯主要加工工艺过程为：建立加工基准、精加工榫头、加工叶尖(叶冠)。

普通模锻毛坯，由于叶片材料、叶片形状和模锻工艺不同，其型面余量也不同，一般铝叶片毛坯的余量比较小，钢叶片毛坯余量比较大。通常把型面余量 0.5 mm 左右的毛坯称小余量毛坯，余量大于 1 mm 的毛坯称为大余量毛坯。大余量毛坯与小余量毛坯机械加工工艺过程基本相同，只是相应增加型面加工工作量。

普通模锻毛坯压气机转子叶片的主要加工工艺过程见表 16-12。

\begin{figure}[H]
\centering
\includegraphics[page=72, viewport={53.5 612.8 253.0 736.8}, clip]{12-07572-001FigRevised.pdf}
\caption{压气机转子叶片简图}
\end{figure}

\begin{longtable}{|m{2cm}|m{2cm}|m{3cm}|m{5cm}|m{2cm}|}
	\caption{普通模锻毛坯压气机转子叶片的主要机械加工工艺过程}\\
	\hline
	工序号 & 工序名称 & 工序内容 & 加工简图 & 加工设备 \\
	\hline
	0 & 模锻 & 普通模锻，在 $K-K$ 截面压模压线作为基准线 &\includegraphics[page=76, viewport={64.5 503.8 218.0 599.8}, clip]{12-07572-001FigRevised.pdf} & 模锻锤 \\
	\hline
	5 & 打磨 & 初步打磨各表面 &\includegraphics[page=76, viewport={235.0 503.8 377.5 596.3}, clip]{12-07572-001FigRevised.pdf} & 砂轮机 \\
	\hline
	10 & 钻孔 & 钻孔 &\includegraphics[page=76, viewport={393.5 505.8 479.5 631.8}, clip]{12-07572-001FigRevised.pdf} & 专用钻床，翻转式钻具 \\
	\hline
	15 & 加工进气边、排气边 & 铣削各型面、各截面进气边位置及弦宽至标注尺寸 &\includegraphics[page=76, viewport={393.5 505.8 479.5 631.8}, clip]{12-07572-001FigRevised.pdf} & 卧式铣床、立式铣床、液压仿形铣床 \\
	\hline
	20 & 粗加工榫头两端 & 铣削尺寸 $a, b$ 至粗加工尺寸 & \includegraphics[page=76, viewport={192.0 343.8 311.0 472.3}, clip]{12-07572-001FigRevised.pdf}& 龙门铣床、组合刀铣床、拉床 \\
	\hline
	25 & 粗加工榫头侧面、底面 & 铣削、拉削尺寸 $a, b, c$ 至标注尺寸，以加工面为基准，检测 $a$ 尺寸是否合格 &\includegraphics[page=76, viewport={335.5 343.8 487.0 462.3}, clip]{12-07572-001FigRevised.pdf} & 立铣床、组合刀铣床、龙门铣床、拉床 \\
	\hline
	30 & 粗加工型面 & 电解、电脉冲加工，检测型面 &\includegraphics[page=76, viewport={77.0 115.3 220.0 215.8}, clip]{12-07572-001FigRevised.pdf} & 电解加工机床、电脉冲加工机床 \\
	\hline
	35 & 半精加工型面 & 抛光、砂带磨、精密电解加工，保证型面进气边位置 & 同上图 & 抛光机、砂带磨床、精密电解加工机床 \\
	\hline
	40 & 精加工榫头 & 拉削、成形磨削、成形铣削，定位榫头检测叶型，保证各尺寸在误差内 &\includegraphics[page=76, viewport={238.0 116.3 344.0 215.8}, clip]{12-07572-001FigRevised.pdf} & 拉床、成形磨床、成形铣床 \\
	\hline
	45 & 粗加工缘板内侧 & 铣削、电加工、磨削各表面，直至与发动机气道空间形状相符 & \includegraphics[page=76, viewport={361.5 116.8 463.5 181.8}, clip]{12-07572-001FigRevised.pdf}& 铣床、电加工机床、磨床 \\
	\hline
	50 & 加工榫头两端头 & 铣削、拉削，保证 $a, b$ 尺寸 &\includegraphics[page=77, viewport={67.5 653.3 187.5 746.3}, clip]{12-07572-001FigRevised.pdf} & 卧式铣床、龙门铣床、拉床 \\
	\hline
	55 & 加工叶尖 & 粗铣 $R$ 尺寸，铣削、线切割，保证 $a$ 尺寸 & \includegraphics[page=77, viewport={199.5 647.8 332.0 733.8}, clip]{12-07572-001FigRevised.pdf}& 切断床、卧式铣床、仿形铣床、线切割机床 \\
	\hline
	60 & 精加工叶型 & 抛光型面，达到设计要求 &\includegraphics[page=77, viewport={354.0 648.8 483.0 706.8}, clip]{12-07572-001FigRevised.pdf} & 抛光机 \\
	\hline
	65 & 成品检验 & 检验各尺寸是否达到设计要求 & & 检测装置 \\
	\hline
	70 & 测频、测重、探伤、表面处理、光饰处理、强化处理 & 按设计要求进行相应步骤 & & 相应测试和表面处理设备 \\
	\hline
\end{longtable}

\section*{复习思考题}


16-1 何谓生产过程、工艺过程、机械加工工艺过程？

16-2 什么是工序？划分工序的依据是什么？

16-3 什么是生产纲领，生成纲领通常依据什么确定？

16-4 什么是生产类型？生产类型有哪几种？汽车、金属切削机床、大型轧钢机的生产各属于哪种生产类型？各有何特征？

16-5 简述按工序集中原则、工序分散原则组织工艺过程的工艺特征，并说明各适用于什么场合。

16-6 图 16-31 所示的小轴，生产数量为 30 件，毛坯为 $\phi 32 \times 104$ 的圆钢料，若用以下两种方案加工：

(1) 先整批车出大端的端面和外圆，随后仍在该车床上整批车出小端的端面和外圆；
(2) 在一台车床上逐件进行加工，即每个工件车好大端后，立即掉头车小端。

\begin{figure}[H]
\centering
\includegraphics[page=72, viewport={268.5 613.8 492.0 747.8}, clip]{12-07572-001FigRevised.pdf}
\caption{题 16-6 图}
\end{figure}

试问这两种方案分别是几道工序？哪种方案较好？为什么？

16-7 什么是工件的夹紧？

16-8 什么是不完全定位？简述不完全定位和欠定位的区别。

16-9 什么是过定位？试举例说明。

16-10 何谓工件的六点定位原理？工件的定位方式有几种？

16-11 什么是夹具？按使用范围的不同，夹具分为哪几类？各适用于什么场合？

16-12 一般夹具由哪几个组成部分？各起什么作用？

16-13 试分析图 16-32 所示各种安装方法工件的定位情况，各限制了工件的哪些自由度？属于哪种定位？

(1) 三爪自定心卡盘夹持工件，工件较长且紧贴垫铁；
(2) 双顶尖、拨盘、卡箍装夹工件；
(3) 平口钳装夹六面体工件，加工工件的顶平面，工件底面悬空，用划针盘按顶平面加工线找正。

\begin{figure}[H]
\centering
\includegraphics[page=72, viewport={56.5 481.3 486.0 587.3}, clip]{12-07572-001FigRevised.pdf}
\caption{题 16-13 图}
\end{figure}

16-14 何谓基准？根据作用的不同，基准分为哪几种？

16-15 何谓粗基准和精基准？如何选择？

16-16 切削加工工序安排的原则是什么？

16-17 常用的工艺文件有哪几种？各适用于什么场合？

16-18 加工轴类零件时，常以什么作为统一的精基准？为什么？

16-19 如何保证盘套类零件外圆面、内孔及端面的位置精度？

16-20 安排机架箱体类零件的加工工艺时，为什么一般要依据“先面后孔”的原则？

16-21 试分别拟订图 16-33 和图 16-34 所示零件在单件、小批生产中的工艺过程。

\begin{figure}[H]
\centering
\includegraphics[page=72, viewport={98.0 274.8 442.5 467.3}, clip]{12-07572-001FigRevised.pdf}
\caption{题 16-21 图(一)}
\end{figure}
未注倒角C1
外侧表面淬火40~45 HRC
材料：45钢
数量：10件

\begin{figure}[H]
\centering
\includegraphics[page=72, viewport={130.5 22.8 411.5 240.8}, clip]{12-07572-001FigRevised.pdf}
\caption{题 16-21 图(二)}
\end{figure}

未注倒角C1
数量：10件
材料：45钢


\chapter{现代制造技术及发展趋势}

随着计算机技术、信息技术，尤其是网络技术的飞速发展，在全球范围内，一个更加激烈的竞争环境已经形成，消费者的价值观念也在发生结构性的变化，导致产品的制造模式由过去单调的大批量生产转变为多品种、小批量的个性化生产方式。在以消费者为导向的时代，如何对市场环境急剧变化做出快速反应，在时间(T)、质量(Q)、成本(C)、服务(S)和环境(E)目标下，有效地进行生产，提供高质量产品和优质服务是制造业面临的挑战。

\section{超精密加工}

机械制造技术的发展主要体现在提高制造生产效率和产品精度两个方面。从提高产品的制造生产效率方面来看，提高产品生产的自动化程度一直是制造业发展的一个主要方向。从计算机数控发展到计算机集成制造系统，其应用范围也在不断延伸。从提高产品的精度方面来看，从精密加工发展到超精密加工，也一直是制造业主要发展的方向之一。根据不同的应用领域，产品的精度从微米到亚微米，乃至纳米，其应用范围也在日益扩大，并且在高技术应用领域、国防工业、生命科学以及民用工业中都获得了非常广泛的应用，如激光核聚变系统、超大规模集成电路、高密度磁盘、精密雷达、导弹火控系统、惯性导航陀螺、精密机床、精密仪器、录像机磁头、复印机磁鼓、煤气灶转阀等零件都是采用超精密加工技术制造的。

超精密加工是指采用普通精密加工手段无法达到的高精度和高表面质量的加工。就其加工零件的精密量级来说，可以获得亚微米乃至毫微米级的形状和尺寸，且能获得纳米级的表面粗糙度。超精密加工方法主要包括超精密切削（如超精密车削和超精密铣削等）、超精密磨削、超精密研磨、超精密特种加工等。

\subsection{超精密切削}

与一般的普通精密切削相比，超精密切削采用的刀具为金刚石刀具。超精密切削之所以能获得高的切削精度，一方面是由于金刚石刀具与非铁合金亲和力比较小，其硬度、耐磨性以及导热性相对比较好，并且可以刃磨得非常锋利；另一方面，在超精密切削中，采用了高精度的空气轴承、气浮导轨、定位检测元件等零部件，采取恒温、防振以及隔振等措施，从而使其所加工工件的表面粗糙度 $Ra$ 值可以达到 $0.01\text{ }\mu\text{m}$ 以下，形状精度可达 $0.2\text{ }\mu\text{m}$。因此，超精密切削技术在航空航天、光学仪器等领域的应用越来越广泛，正朝着更高精度的方向发展。

\subsection{超精密磨削}

超精密磨削是在一般精密磨削基础上发展起来的一种新技术。在超精密磨削中，其加工材料主要是玻璃、陶瓷等硬脆材料。超精密磨削可以省去传统磨削中的抛光工序，在简化工艺的同时，不但能获得精确的几何形状和尺寸，还可以获得镜面级的表面粗糙度。为了获得这一要求，在超精密磨削中，除了需要采取合理的工艺流程提高加工精度之外，还要求加工设备具有高精度、高刚度以及高阻尼特征等特性，消除各种动态误差的影响，并具有高精度检测手段和补偿手段，以获取高的磨削精度。

\subsection{超精密研磨}

超精密研磨包括机械研磨、化学机械研磨、浮动研磨、弹性发射加工以及磁力研磨等加工方法。超精密研磨加工出的球面跳动可达 $0.025\text{ }\mu\text{m}$，表面粗糙度 $Ra$ 值可达 $0.003\text{ }\mu\text{m}$。最高精度的超精密研磨可加工出平面度为 $0.01\text{ }\mu\text{m}/200\text{ mm}$ 的零件。超精密研磨的关键条件是精密的温度控制、加工过程无振动、洁净的环境以及细小而均匀的研磨剂。此外，高精度的检测方法也必不可少。

\subsection{超精密特种加工}

超精密特种加工技术被国际上公认为是 21 世纪最有前途的技术。它是指采用电能、热能、光能、电化学能、化学能、声能及特殊机械能等能量去除或增加材料的加工方法。其应用对象主要有难加工材料（如钛合金、耐热不锈钢、高强钢、复合材料、工程陶瓷、金刚石、红宝石、硬化玻璃等高硬度、高韧性、高强度、高熔点材料）、难加工零件（如复杂零件三维型腔、型孔、群孔和窄缝等）、低刚度零件（如薄壁零件、弹性元件等零件）。这种加工方法主要包括激光加工技术、电子束加工技术、离子束及等离子加工技术、电加工技术等。这里仅做简要介绍。

1. 激光加工技术

激光发生器将高能量密度的激光进一步聚焦后照射到工件表面，光能被吸收瞬时转化为热能，根据能量密度的高低，可实现打孔、精密切割、加工精微防伪标志等。随着激光加工设备和工艺的迅速发展，已出现了 $100\text{ kW}$ 以上的大功率激光器、千瓦级的高光束固体激光器，同时还配有光导纤维可进行多工位、远距离工作。由于激光加工设备功率大、自动化程度高，已普遍采用计算机数字控制、多坐标联动，并装有激光功率监控、自动聚焦、工业电视显示等辅助系统。目前，激光制孔的最小孔径已达到 $0.002\text{ mm}$，激光切割薄材的速度可达 $15\text{ m/min}$，切缝的宽度仅为 $0.1\sim1\text{ mm}$。激光表面强化、表面重熔、合金化、非晶化处理技术应用越来越广，激光微细加工在电子、生物、医疗工程方面的应用已成为无可替代的特种加工技术。

2. 电子束加工技术

电子枪在真空中从阴极向正极不断发射负电子，负电子从电子枪的阴极向正极移动过程中不断地加速，聚焦成极细的、能量密度极高的电子束流。当高速运动的电子撞击到工件表面时，其动能转化为热能，使材料熔化、气化，并从真空中被抽走。控制电子束的强弱和偏转方向，配合工作台 $x$、$y$ 方向的数控位移，可实现打孔、成形切割、刻蚀、光刻曝光等工艺。集成电路制造中广泛采用波长比可见光短得多的电子束光刻曝光，可达到 $0.25\text{ }\mu\text{m}$ 的线条图形分辨率。电子束加工技术日趋成熟，已广泛用于运载火箭、航天飞机等主承力构件大型结构的组合焊接以及飞机梁、框、起落架部件、发动机整体转子、机匣、功率轴等重要结构件和核动力装置压力容器的制造。电子束加工采用计算机数字控制、多坐标联动，自动化程度高。电子束焊接已成功地应用在特种材料、异种材料、空间复杂曲线、变截面等方面焊接。

3. 离子束加工技术

它是指在真空中将离子源产生的离子加速、聚焦使之撞击工件表面，从而实现材料微量去除和表面改性的加工方法。由于离子是带正电荷且质量比电子大数千万倍，故加速以后可以获得更大的动能。离子束加工技术是靠微观的机械撞击能量而不是靠动能转化为热能来加工的，可用于表面刻蚀、超净清洗，实现原子、分子级的切削加工。

4. 微细电火花加工

它是指在绝缘的工作液中通过工具电极和工件间脉冲火花放电产生的瞬时局部高温来熔化和气化材料，从而实现材料微量去除的加工方法。微细电火花加工过程中工具与工件间没有宏观的切削力，只要精密地控制单个脉冲放电能量并配合精密微量进给，就可实现极微细的金属材料的去除，可加工微细轴、孔、窄缝、平面以及曲面等。电火花加工气膜孔采用多通道、纳秒级超高频脉冲电源和多电极同时加工的专用设备，加工效率 $2\sim3\text{ s/孔}$，表面粗糙度 $Ra$ 值为 $0.4\text{ }\mu\text{m}$。高档电火花成形及线切割已能提供微米级加工精度，可加工 $3\text{ }\mu\text{m}$ 的微细轴和 $5\text{ }\mu\text{m}$ 的孔。

5. 微细电解加工

它是基于电化学阳极溶解原理的一种特种精密加工方法，其工作过程为：在导电的工作液中将水分解为氢离子和氢氧根离子，工件作为阳极，其表面的金属原子成为金属正离子溶入电解液而被逐层地电解下来，随后与电解液中的氢氧根离子发生反应形成金属氢氧化物沉淀，而工件阴极并不损耗。微细电解加工过程中工具与工件间也不存在宏观的切削力，只要精细地控制电流密度和电解部位，就可实现纳米级精度的电解加工，而且表面不会产生加工应力。常用于镜面抛光、精密减薄以及一些需要无应力加工的场合。微细电解加工应用较广，除叶片和整体叶轮外，已扩大到机匣、盘环零件和深小孔加工。微细电解加工可加工出高精度金属反射镜面。目前，电解加工机床的最大电流容量已达到 $50\text{ }000\text{ A}$，并已实现计算机数字控制和多参数自适应控制。

6. 复合加工

它是指采用几种不同能量形式、几种不同的工艺方法，互相取长补短、复合作用的加工技术，例如电解研磨、超声电解加工、超声电解研磨、超声电火花、超声切削加工等，可比单一加工方法更有效，适用范围更广。

\section{纳米加工技术}

正如制造技术在当今各领域所起的重要作用一样，纳米加工技术在纳米技术的各领域中也起着关键作用。纳米加工技术包含机械加工、化学腐蚀、高能束加工以及利用扫描隧道显微镜在工件表面上的电场电工（STM）加工等许多方法。关于纳米加工技术目前还没有一个统一的定义。广义而言，$100\text{ nm}$ 以下的材料的加工或表面粗糙度为纳米级的加工均可称为纳米加工；狭义的纳米加工技术则是指零件加工的尺寸精度、形状精度以及表面粗糙度均为纳米级。通过以下加工技术可以实现纳米级加工。

\subsection{纳米级机械加工技术}

纳米级机械加工方法有单晶金刚石和立方氮化硼（CBN）制作的单点刀具的超精密切削、单晶金刚石和立方氮化硼（CBN）磨料制作的磨具的超精密多点磨料加工，以及研磨、抛光、弹性发射加工等自由磨料加工或机械化学复合加工等。

目前，利用单点金刚石超精密切削加工已在实验室得到了 $3\text{ nm}$ 的表面精度，利用可延性磨削技术也实现了纳米级磨削，而通过弹性发射加工等工艺则可以实现亚纳米级的材料去除，得到埃级的表面粗糙度。

\subsection{高能束加工技术}

高能束加工是利用能量密度很高的激光束、电子束或离子束等去除工件材料的特种加工方法，主要包括离子束加工、电子束加工和光束加工等。此外，电解射流加工、电火花加工、电化学加工、分子束外延、物理和化学气相沉积等也属于高能束加工。

离子束溅射去除加工、沉积加工和表面处理，以及离子束辅助蚀刻也是纳米级加工的研究开发方向。与固体工具切削加工相比，离子束加工的位置和加工速率难以确定，为取得纳米级的加工精度，需要亚纳米级检测系统与加工位置的闭环调节系统。电子束加工以热能的形式去除穿透层表面的原子，可以进行刻蚀、光刻曝光、焊接、微米和纳米级钻削和铣削加工等。1999 年初，$0.18\text{ }\mu\text{m}$ 工艺的深紫外线（DUV）光刻机已相继投放市场。在 $0.1\text{ }\mu\text{m}$ 之后用于替代光学光刻的所谓下一代光刻技术（NGL）主要包括极紫外光、X 射线、电子束以及离子束光刻。下面就各种光刻技术进展情况做简要介绍。

1. 光学光刻

光学光刻通过光学系统以投影方法将掩模上的大规模集成电路器件的结构图形“刻”在涂有光刻胶的硅片上。光学光刻所能获得的最小特征尺寸与光学光刻系统所能获得的分辨率直接相关，而减小光源的波长是提高分辨率的最有效途径。因此，开发新型短波长光源光刻机一直是国际上的研究热点。目前，商品化光刻机的光源波长已经从过去的汞灯光源紫外光波段进入到深紫外波段（DUV），如用于 $0.25\text{ }\mu\text{m}$ 技术的 KrF 准分子激光（波长为 $248\text{ nm}$）和用于 $0.18\text{ }\mu\text{m}$ 技术的 ArF 准分子激光（波长为 $193\text{ nm}$）。除此之外，利用光的干涉特性，采用各种波前技术优化工艺参数也是提高光刻分辨率的重要手段。这些技术是运用电磁理论结合光刻实际对曝光成像进行深入的分析所取得的突破，其中有移相掩模、离轴照明技术、邻近效应校正等。运用这些技术，可在目前的技术水平上获得更高分辨率的光刻图形。如 1999 年初 Canon 公司推出的 FPA-1000ASI 扫描步进机，该机的光源为 $193\text{ nm}$ ArF，通过采用波前技术，可在 $300\text{ mm}$ 硅片上实现 $0.13\text{ }\mu\text{m}$ 光刻线宽。目前，基于同一波长的 $193\text{ nm}$ ArF 浸没式光刻技术，通过多重图案化、光学邻近校正（optical proximity correction, OPC）及先进工艺整合等持续优化，已能稳定生产 $7\text{ nm}$ 及更小等效线宽的芯片。

光学光刻技术包括光刻机、掩模、光刻胶等一系列技术，涉及光、机、电、物理、化学、材料等多个研究领域。目前，科学家正在探索更短波长的 $\text{F}_2$ 激光（波长为 $157\text{ nm}$）光刻技术。由于大量的光吸收，获得用于光刻系统的新型光学及掩模衬底材料是该波段技术的主要困难。

2. 极紫外光刻

极紫外光刻（EUVL）用波长为 $10\sim14\text{ nm}$ 的极紫外光作为光源。虽然该技术最初被称为软 X 射线光刻，但实际上更类似于光学光刻。所不同的是由于极紫外光在材料中的强烈吸收，其光学系统必须采用反射形式。

3. X 射线光刻

X 射线光刻（X-ray lithography, XRL）光源波长约为 $1\text{ nm}$。由于易于实现高分辨率曝光，XRL 技术自 20 世纪 70 年代被发明以来，就受到人们广泛的重视，欧洲、美国、日本和中国等拥有同步辐射装置的国家相继开展了有关研究。XRL 技术是所有下一代光刻技术中最为成熟的技术，其面临的主要困难是获得具有良好机械物理特性的掩模衬底。近年来，掩模技术研究取得了较大进展。SiC 是最合适的衬底材料。由于与 XRL 技术相关问题的研究已经比较深入，加之光学光刻技术的发展和其他光刻技术的新突破，XRL 技术不再是未来唯一的候选技术。最近，美国对 XRL 技术的投入有所减小。尽管如此，XRL 技术仍然是不可忽视的候选技术之一。

4. 电子束光刻

电子束光刻（electron beam lithography, EBL）采用高能电子束对光刻胶进行曝光从而获得结构图形的加工方法。由于其德布罗意波长为 $0.004\text{ nm}$ 左右，电子束光刻不受衍射极限的影响，可获得接近原子尺度的分辨率。电子束光刻可以获得极高的分辨率并能直接产生图形，不但在超大规模集成电路（very large scale integration, VLSI）制作中已成为不可缺少的掩模制备工具，也是加工用于特殊目的的器件和结构的主要方法。目前的电子束曝光机的分辨率已达 $0.1\text{ }\mu\text{m}$ 以下。电子束光刻的主要缺点是生产效率较低，为每小时 $5\sim10$ 个圆片，远小于目前光学光刻的每小时 $50\sim100$ 个圆片的水平。美国朗讯公司开发的角度限制散射投影电子束光刻（scattering with angular limitation projection electron lithography, SCALPEL）技术令人瞩目。该技术如同光学光刻那样对掩模图形进行缩小投影，并采用特殊滤波技术去除掩模吸收体产生的散射电子，从而在保证分辨率条件下提高产出效率。应该指出，无论未来光刻采用何种技术，EBL 技术都将是集成电路研究与生产不可缺少的基础设施。

5. 离子束光刻

离子束光刻（ion beam lithography, IBL）采用液态原子或固态原子电离后形成的离子通过电磁场加速及电磁透镜的聚焦或准直后对光刻胶进行曝光，其原理与电子束光刻类似，但德布罗意波长更短（小于 $0.000\text{ }1\text{ nm}$），且具有无邻近效应小、曝光场大等优点。离子束光刻主要包括聚焦离子束光刻（focused ion beam lithography, FIBL）、离子投影光刻（ion projection lithography, IPL）等。其中，FIBL 技术发展最早，实验研究中已获得 $10\text{ nm}$ 的分辨率。该技术由于效率低，很难在生产中作为曝光工具得到应用，目前主要用做超大规模集成电路中的掩模修补工具和特殊器件的修整。

由于 FIBL 技术的缺点，人们发展了具有较高曝光效率的 IPL 技术。欧洲和美国联合了大量企业、大学和研究机构，开展了一个名为 MEDEA 的合作项目，用于解决设备和掩模等方面的问题，进行可行性验证，目前已取得了不少进展。

\subsection{光刻电模铸造技术}

光刻电模铸造工艺是由深层同步辐射 X 射线光刻、电铸成形、塑铸成形等技术组合而成的综合性技术，其最基本和最核心的工艺是深层同步辐射光刻，而电铸和塑铸工艺是光刻电模铸造产品实用化的关键。与传统的半导体工艺相比，光刻电模铸造技术具有许多独特的优点。

（1）用材广泛，可以是金属及其合金、陶瓷、聚合物、玻璃等。

（2）可以制作高度达数百微米至一千微米，高度比大于 200 的三维立体微结构。

（3）横向尺寸可以小到 $0.5\text{ }\mu\text{m}$，加工精度可达 $0.1\text{ }\mu\text{m}$。

（4）可实现大批量复制、生产，且成本低。

用光刻电模铸造技术可以制作各种微器件、微装置，已研制成功或正在研制的光刻电模铸造产品有微传感器、微电机、微机械零件、集成光学和微光学元件、微波元件、真空电子元件、微型医疗器械、纳米技术元件及系统等。光刻电模铸造产品的应用涉及面广泛，如加工技术、测量技术、自动化技术、汽车及交通技术、电力及能源技术、航空及航天技术、纺织技术、精密工程及光学、微电子学、生物医学、环境科学和化学工程等。

\subsection{扫描隧道显微技术}

宾宁（G. Binning）和罗雷尔（H. Rohrer）发明的扫描隧道显微镜（scanning tunneling microscope, STM），不但使人们可以以单个原子的分辨率观测物体的表面结构，而且也为以单个原子为单位的纳米级加工提供了理想途径。应用扫描隧道显微技术可以进行原子级操作、装配和改型。扫描隧道显微镜将非常尖锐的金属针（探针）接近试件表面至 $1\text{ nm}$ 左右，施加电压时隧道电流产生，隧道电流每隔 $0.1\text{ nm}$ 变化一个数量级。采用电流保持一定的工作方式，对试件表面进行扫描，即可分辨出表面结构。一般隧道电流通过探针尖端的一个原子，因而其横向分辨率为原子级。扫描隧道显微加工技术不仅可以进行单个原子的去除、添加和移动，而且可以进行扫描隧道显微光刻、探针尖电子束感应的沉淀和腐蚀等。

\section{超高速切削加工技术}

高速切削加工是指在高于常规切削速度几倍、十几倍条件下所进行的切削加工。目前国际上对高速切削加工定义尚不统一，主要有高切削速度切削（high cutting speed machining）、高主轴转速切削（high spindle speed machining）、高进给速度切削（high feed speed machining）、高进给高速切削（high speed and feed machining）和高生产率切削（high production machining）等。针对不同材料，高速切削的范围差别很大，如对一些有色金属和易加工钢，高速和超高速切削速度已达到每分钟几千米，而对高硬度材料和高温合金等难加工材料，切削速度只有每分钟几百米。

超高速切削主要有以下特点。

（1）切削过程散热条件好。切屑形态从带状演变为片状或碎屑，切削热的 $95\%$ 将被切屑带走，加工表面受热时间短，切削温度低，热影响区和热影响程度都较小，有利于获得低损伤的表面和保持良好的表面物理性能及力学性能。
（2）可获得较好的表面完整性。高速切削时，在保证生产率效率的同时，可采用较小的进给量，从而降低了加工表面的粗糙度。
（3）刀具使用寿命高。在高速切削过程中，刀具磨损量的增长速度低于切削效率的提高速度，因此刀具的寿命相对有较大的提高。据有关资料介绍，高速切削的刀具寿命比传统切削方法提高了 $70\%$ 左右。

超高速切削已广泛应用于航空、航天和汽车等领域，取得了巨大的经济效益。在模具制造行业，德国、日本、美国等有 $30\%\sim50\%$ 的模具公司用超高速切削加工技术加工电火花加工电极、淬硬模具型腔、塑料和铝合金模型等，加工效率高，质量好，减少了后续手工打磨和抛光工序。在航空航天领域，飞行器的骨架与机翼、桁架均为铝合金整体薄壁构件，加工时均需切除大量的金属，毛坯的切除量甚至达到 $90\%$，采用超高速切削加工技术，加工时间可缩短至原来的几分之一。在汽车制造行业，汽车发动机的铝合金和铸铁缸体广泛采用超高速切削加工技术制造，大大地提高了效率，降低了生产成本。此外，超高速切削加工技术还广泛用于快速成形、光学精密零件和仪器仪表等加工领域。

超高速机床是实现超高速切削的前提条件和关键设备。超高速切削对机床应包括适应于超高速运转的主轴部件及其驱动系统，快速反应的数控伺服系统和进给部件，高压大流量喷射冷却系统和很好的刚度（静刚度、动刚度、热刚度）特性的机床支承件等。

\section{柔性制造系统}

柔性制造系统（flexible manufacturing system, FMS）起源于 20 世纪 70 年代，它是利用系统工程学原理和成组技术解决制造业中的多品种、小批量生产，并使其达到整体最优的自动化加工系统。柔性制造系统通常是通过局域网把数控机床（加工中心）、坐标测量机、物料输送装置、对刀仪、立体刀库、工件装卸站、机器人等设备连接起来，在计算机及控制软件的控制下形成一个加工系统。自 1967 年由英国莫林斯（Molins）公司提出以后，至今已发展了近 60 年。柔性制造技术是管理技术和制造技术的有机集合，是以数控技术为核心，以计算机技术、信息技术、检测技术、质量控制技术与生产管理技术相结合的先进制造技术，因而被世界各国所重视，并在发达国家的制造业中得到了广泛的应用。柔性制造技术的应用，既解决了近百年来中、小批量和大批量、多品种加工自动化的问题，也很好地适应了产品不断迅速更新的需求，即解决了 20 世纪 60 年代以来一直以提高切削用量、减少切削时间和缩短辅助时间为手段来达到提高生产率、但效果不理想的问题。柔性制造技术具有较高的柔性和通用性，转产快、准备时间短，设备利用率高，可实现无人看管 $24\text{ h}$ 连续工作，加工质量高且稳定，生产成本低。

柔性制造技术的这种高效、灵活的特性，使其成为实施敏捷制造、并行工程、精益制造和智能制造的基础，其应用日益广泛，已成为整个机械制造领域的核心技术。

\subsection{柔性制造系统的组成}

一个柔性制造系统主要由加工系统、物料系统、计算机控制系统及系统软件等几部分组成。加工系统采用的设备由待加工工件的类别决定，主要有加工中心、车削中心或数控车床、数控铣床、数控磨床及数控齿轮加工机床等，用以自动地完成多种工序的加工；物料系统用以实现工件及工装的自动供给和装卸，以及完成工序间的自动传送、调运和存贮工作，包括各种传送带、自动导引小车、工业机器人及专用起吊运送机等；计算机控制系统用以处理柔性制造系统的各种信息，输出控制数控机床和物料系统等自动操作所需的信息，通常采用三级（设备级、工作站级、单元级）分布式计算机控制系统，其中单元级控制系统（单元控制器）是柔性制造系统的核心；系统软件用以确保柔性制造系统有效地适应中、小批量和大批量、多品种生产的管理、控制及优化工作，包括设计规划软件、生产过程分析软件、生产过程调度软件、系统管理和监控软件。

\subsection{柔性制造系统的分类}
按规模大小,柔性制造系统可分为如下4类。

1.柔性制造单元
    
    柔性制造单元(flexible manufacturing cell,FMC)的问世并在生产中使用要比柔性制造系统晚6$\sim$8年,它由1至2台加工中心、工业机器人、数控机床及物料运送存贮设备构成,具有适应加工大批量、多品种产品的灵活性。柔性制造单元可视为一个规模最小的柔性制造系统,是柔性制造系统向廉价化及小型化方向发展的一种产物,其特点是实现单机柔性化及自动化,已进入普及应用阶段。
    
2.柔性制造系统
    
    柔性制造系统通常包括4台或更多台全自动数控机床(加工中心与车削中心等),由集中的控制系统及物料系统连接起来,可在不停机的情况下实现大批量、多品种和中、小批量产品的加工及管理。
    
3.柔性制造线
    
    柔性制造线(fleible manufacturing line,FML)是处于单一或少品种、大批量非柔性自动线与中、小批量和大批量、多品种柔性制造系统之间的生产线。其加工设备可以是通用的加工中心、数控机床,亦可采用专用机床或数控专用机床。对物料系统柔性的要求低于柔性制造系统,但生产率更高。它是以离散型生产中的柔性制造系统和连续生产过程中的集散控制系统(distributed control system,DCS)为代表,其特点是实现生产线柔性化及自动化,其技术已日臻成熟,已进入实用化阶段。
    
4.柔性制造工厂
    
    柔性制造工厂(flexible manufactuing factory,FMF)是将多条柔性制造系统连接起来,配以自动化立体仓库,用计算机系统进行联系,采用从订货、设计、加工、装配、检验、运送至发货的完整柔性制造系统。它计算机集成制系统的理念,实现生产系统柔性化及自动化,进而实现全厂范围的生产管理、产品加工及物料贮运等全流程柔性化、智能化运行。柔性制造工厂是自动化生产的最高水平,反映出世界上最先进的自动化应用技术。


\subsection{柔性制造技术的发展}
柔性制造技术是实现未来工厂的新概念模式和新发展趋势,是决定制造企业未来发展前途的具有战略意义的举措。目前,反映工厂整体水平的柔性制造系统是第一代柔性制造系统,今后此种状况仍有可能持续下去。日本从1991年开始实施的智能制造系统(intelligent manufacturing system,IMS)国际性开发项目,可归为第一代柔性制造系统。而真正完善的第二代FMS正加速落地,以数字孪生、AI驱动的动态调度、人机协作机器人及云端供应链协同为核心,逐步实现“智造”与“管理”的柔性集成。届时,智能化机械与人之间将相互融合、柔性地全面协调从接受订单至生产、销售这一企业生产经营的全部活动。

柔性制造系统在20世纪80年代中期获得了迅猛发展,成为生产自动化的热点。一方面,单项技术如数控加工中心、工业机器人、CAD/CAM、资源管理及高技术的发展,为系统集成提供了关键技术基础;另一方面,世界市场发生了重大变化,由过去传统、相对稳定的市场,发展为动态多变的市场。为了在市场中求生存、求发展,提高企业对市场需求的应变能力,人们开始探索新的生产方法和经营模式。近年来,柔性制造作为一种现代化工业生产的科学“哲理”和工厂自动化的先进模式已为国际上所公认。柔性制造系统是在自动化技术、信息技术及制造技术的基础上,将以往企业中相互独立的工程设计、生产制造及经营管理等过程,在计算机及其软件的支持下,构成一个覆盖整个企业的完整而有机的系统,以实现全局动态最优化,总体高效益、高柔性,进而赢得竞争全胜的智能制造系统。柔性制造作为当今世界制造自动化技术发展的前沿科技,为未来机械制造工厂提供了一幅宏伟的蓝图,将成为21世纪机械制造业的主要生产模式。

柔性制造技术已经在各工业发达国家得到广泛的生产应用,数控机床、加工中心、柔性制造线、柔性制造工厂应用日益广泛。计算机集成制造将使生产的设计、制造管理、供销、财务都用计算机统一管理,实现工厂的全盘计算机管理的自动化。20世纪80年代,计算机集成制造系统已成为机械制造的发展方向,但建立计算机集成制造系统往往需要大量的硬件投入和较长的开发时间。但是,由于计算机应用技术的飞速发展,过分强调统一管理,无法适应多变的生产。因此,现在计算机集成制造系统的概念也必须不断更新,要有更多的灵活性能才能适应迅速多变的生产。CAD/CAM一体化技术的发展应用,大大缩短了产品的研制开发周期,同时也促进了设计思想的变化。面向制造的设计(design for manufacturing,DFM)的思想现已被更多的人所接受,从而在保证产品性能要求的前提下大大地减少了制造成本。

\section{虚拟制造技术}
\subsection{虚拟制造的概念}
虚拟制造(virtual manufacturing,VM)是利用计算机仿真与虚拟现实技术,在高性能计算机及高速网络的支持下,采用群组协同工作,通过模型来模拟和预估产品功能、性能及可加工性等各方面可能存在的问题,实现产品制造本质过程的一种先进制造技术。它包括产品的设计、工艺规划、加工制造、性能分析、质量检验,并进行过程管理与控制,以增强制造过程各级的决策与控制能力。

\subsection{虚拟制造的基本原理}
虚拟制造系统(virtual manufacturing system,VMS)可实现产品开发和制造过程各环节的集成。它根据制造企业的制造策略,在信息集成和功能集成的基础上,通过对整个产品开发过程的建模、管理、控制和协调,使企业资源、技术、人员得到合理的组织和配置,在产品的整个生命周期内,实现制造企业策略、企业经营、工程设计和生产活动的集成,以及并行处理虚拟环境下多学科小组的协同工作,以适应市场和用户需求的不断变化,以最快的速度向市场和用户提供优质低价产品。

虚拟制造系统能够在产品开发的各个阶段,根据产品的要求,在虚拟环境下对虚拟产品原型的结构、功能、性能与加工、装配过程进行仿真,并依据产品评价体系的方法、规范和指标,设计修改和优化产品制造过程。由于以上过程都是在虚拟环境下针对虚拟产品原型进行的,所以可以大大缩短产品的开发时间。同时,虚拟制造系统能够在产品开发的早期阶段及时发现可能存在的问题,并在实际制造之前予以解决;开发进程的加快,使设计$\to$评价$\to$修改的设计方式,从串行工作模式下的局部循环过渡到基于并行工作模式的整体循环,进而有利于实现对多个解决方案的比较和优化选择。

虚拟制造系统通过多层次、分阶段、多目标的决策支持,不仅可以提高产品开发过程中的决策和控制能力,而且在人与系统之间、设计人员之间、现实世界与虚拟的计算机制造环境之间、不同的产品开发过程之间、不同的知识来源之间等建立起连接的纽带,从而在产品开发的不同阶段,依据不同的评价体系、评价指标及规范,在产品制造过程中实现从时间、质量、成本等多个方面的综合评价,加强企业的决策和控制能力,以提高运行效率。

\subsection{虚拟制造系统研究的关键技术}
虚拟制造技术和虚拟制造系统的理论基础与体系尚未完全形成,正处在研究探索阶段。研究工作所面对的是多种学科的交叉、多种高新技术的融合,需要研究和尚待解决的关键技术仍然很多。主要包括以下几个方面。

(1)虚拟制造系统的理论体系,以及基于分布式并行处理环境下的虚拟制造系统的开放式体系结构(VMS-OSA)的研究。

(2)虚拟环境下的产品主模型技术。主模型是一个核心,能以此为中心通向设计、制造、生产管理等各个环节并为其提供服务。主模型有统一的数据结构和分布式数据管理系统,它所建立的产品模型是虚拟产品模型,具有代表对象所具备的各种性能和特征。虚拟产品在它投入生产以前就已存在,它具有明显的可视性,能并行处理设计分析、加工制造、生产组织与调度等各种生产环节所面临的诸多问题,并能在供销之间建立信息系统。

(3)虚拟环境下分布式并行处理的分布式智能协同求解技术与系统,以及虚拟环境下系统全局最优决策理论与技术。

(4)应用虚拟现实技术实现虚拟环境及虚拟制造过程中的人机协同求解模型。

(5)现实制造系统与虚拟制造系统之间的映射,虚拟设备、虚拟传感器、虚拟单元、虚拟生产线、虚拟车间及虚拟工厂(公司)的建立,以及各种虚拟设备的重用性和重组性技术。

(6)虚拟制造系统集成开发平台的体系结构、构件库及用户开发环境。

(7)虚拟公司的组织、调度及控制策略与技术。

(8)综合可视化技术(包括计算机可视化技术、虚拟现实技术、多媒体技术和仿真技术)的研究与应用。目前,已开展基于真实感的虚拟产品的装配仿真、生产制造过程及生产调度仿真、数控加工过程仿真等技术研究。


\section{计算机辅助设计与制造技术}
\subsection{计算机辅助设计和计算机辅助制造的基本概念}
计算机辅助设计和计算机辅助制造(CAD/CAM)技术是一项综合性的、技术复杂的系统工程,涉及许多学科领域,如计算机科学和工程,计算数学、几何造型、计算机图形显示、数据结构和数据库、仿真、数控、机器人和人工智能学科与技术,以及与产品设计和制造有关的专业知识等。它是产品设计人员和产品制造人员在计算机系统的辅助下,根据产品的设计和制造程序进行设计和制造的一项新技术,是传统技术与计算机技术的有机结合。同时,CAD/CAM技术本身具有自己的特点和发展规律,随着电子科学技术的发展,新的技术也在不断地涌现,需要人们不断地去探索和研究。

目前,随着CAD/CAM技术的使用日益广泛,其优点越来越突显。早在1985年,美国信息制造业专家W. H. Slatterback曾经预言,从1985年到2000年,美国的制造业面临的变化将比20世纪前75年的变化要大得多,其根本原因是CAD/CAM技术的普遍应用。目前,CAD/CAM技术不仅已广泛用于航空航天、电子和机械制造等工程和产品生产领域,而且发展到服装、装饰、家具和制鞋等应用领域。另外,随着网络化、信息化制造的不断发展,作为计算机集成制造系统的技术基础,CAD/CAM技术会随着网络化的发展而飞速发展。

CAD/CAM技术的普及和应用不仅对传统的制造产业提出新的挑战,对新兴产业的发展、劳动生产率的提高、材料消耗的降低、国际竞争能力的增强具有重要作用,而且也成为衡量一个国家科学技术现代化和工业现代化水平的重要标志之一。

从产品制造的过程来看,产品的制造过程往往要经过图样设计或三维模型、工艺设计才能进行加工。所以,CAD/CAM技术又可细分为CAD/CAPP/CAM,其中计算机辅助工艺规划(CAPP)是连接CAD、CAM的桥梁。

1.计算机辅助设计
    
    计算机辅助设计(computer aided design,CAD)是指工程技术人员在人和计算机组成的系统中以计算机为辅助工具,通过人-机交互操作方式进行产品设计构思和论证、产品总体设计、技术设计、零部件设计,有关零件的强度、刚度、热、电、磁的分析计算和零件加工信息(工程图样或数控加工信息等)的输出,以及技术文档和有关技术报告的编制等,从而达到提高产品设计质量、缩短产品开发周期、降低产品成本的目的。CAD系统的主要功能包括草图设计、零件设计、装配设计、复杂曲面设计、工程图样绘制、工程分析、真实感及渲染和数据交换等。
    
2.计算机辅助工艺规划
    
    计算机辅助工艺规划(computer aided process planning,CAPP)是指在人和计算机组成的系统中,根据产品设计阶段给出的信息,人-机交互地或自动地确定产品加工方法和工艺过程。在CAD/CAM集成环境中,工艺设计人员通常可以根据CAD过程提供的信息和CAM系统的功能,进行零部件加工工艺路线的控制和加工状况的仿真,以生成控制零件加工过程的信息。CAPP的基本功能主要包括毛坯设计,加工方法选择,工艺路线制订,工序、工步设计,刀具、夹具设计等。其中,工序设计又包含机床和刀具选择、切削用量选择、加工余量分配、材料消耗定额的计算以及工时定额计算等。
    
3.计算机辅助制造
    
    在机械制造业中,利用计算机通过各种数值控制机床和设备,自动完成离散产品的加工、装配、检测和包装等制造过程,称为计算机辅助制造(computer aided manufacture,CAM)。
    
    CAM有广义和狭义两种定义。广义CAM一般是指利用计算机辅助完成从生产准备到产品制造整个过程的活动,包括工艺过程设计、工装设计、数控自动编程、生产作业计划、生产控制和质量控制等。狭义CAM通常是指数控程序编制,包括刀具路径规划、刀位文件生成、刀具轨迹仿真及数控代码生成等。


计算机技术的引进,大大地促进了设计和制造能力的提高。这种能力的提高,不但体现在工作效率和工作质量方面,更体现在先进的计算机技术对传统工作方式的促进和变革方面。但要指出,CAD/CAM技术不能代替人们的设计和制造行为,而只是实现这些行为的先进手段和工具。而人们的设计和制造行为,则由专业技术人员的创造能力和工作经验,以及现代设计方法等提供的科学思维方法和实施办法来确定。

\subsection{CAD/CAM 研究的关键技术}
CAD/CAM系统体系结构可分为基础层、支撑层和应用层三个层次。基础层由计算机、外围设备和系统软件组成。其中,系统软件包括各种支撑软件、系统开发和维护的工具软件。支撑层包括CAD/CAM支撑、产品数据管理、图形显示等软件。随着互联网/企业内部网(Intranet/Internet)的广泛使用,分布式协同设计和制造CAD/CAM环境正成为支撑层中的一个重要组成部分。应用层则是针对不同的应用领域需求而开发的各种CAD/CAM应用系统。

1.集成化技术
    
    众所周知,产品的设计过程作为一项创造性的活动,随着自然科学、技术科学、环境科学和人文科学的发展已变成一项综合技术,将系统观点和信息观点引入制造业,出现了计算机集成制造系统的概念。CAD/CAM技术经过长期的发展,其各个单元技术(如CAD、CAPP、CAM、PDM,ERP等)已日渐成熟,在各自的领域起到愈来愈重要的作用。例如,CAD系统主要用于构造产品模型;CAM系统主要用于输出数控指令和数控加工;CAPP系统用于进行工艺的自动设计;产品数据管理(product data management,PDM)用于管理产品数据信息;ERP(enterprise resource planning,企业资源计划)系统用于管理和控制制造资源的分配和使用。但由于这些独立的分系统之间无法实现各系统间信息的自动传递和交换,因此许多工作不得不在各分系统中重复来做,如在CAPP中需要建立产品特征模型,在CAM系统中需要再次建立产品的模型,而普通的CAD模型主要用于图样的生成和产品仿真。
    
    所谓集成通常是指以统一产品数据模型及工程数据库为基础,在系统之间及系统内部实现信息传递、响应、分析及反馈,从而使系统及各模块之间实现有机整合。
    
    随着CAD及相关技术的不断深入,集成的概念也不断深化,目前对集成的认识是以信息集成为基础的多集成的概念。实现多集成的目的,是在TQCSE(T—time,Q—quality,C—cost,S—service,E—environment)目标下,寻求全局最优决策,实现可持续发展的策略。其中,集成的内容主要包括信息集成、智能集成、过程集成、工作机制集成、资源集成、技术集成和人机环境集成等。
    
2.智能化技术
    
    智能制造系统就是将人工智能融入制造系统的各个环节中,通过模拟专家的智能活动取代或延伸制造环境中应由专家完成的那部分活动。在智能制造系统中,系统具有部分人类专家的“智能”。如系统能自动监视自身的运动状态;能够自动调整自身参数来适应外部环境,使自己始终运行在最佳状态下。智能制造系统的研究和应用主要取决于人工智能技术的发展。
    
3.网络化技术
    
    网络化技术包括硬件载体与软件支撑以及各种通信协议及制造自动化协议、通信接口、系统操作控制策略等,是实现各种制造系统自动化的基础。特别是在20世纪90年代以后,互联网/企业内部网的发展,为异地、协同设计的研究和应用提供了平台,CAD/CAM技术也朝着网络化的方向发展。目前,这方面的研究内容主要集中在以下几方面。

(1)因特网/企业内部网络环境下的异地协同设计平台的建立。

(2)并行协同工作原理和实施技术(包括协同问题求解、合作运行机制和管理控制)。

(3)协同工作环境下的产品建模问题。

(4)基于网络的企业制造资源管理。


    

4.可视化技术
    
    虚拟现实技术、多媒体技术及计算机仿真技术可实现产品设计制造过程中的几何仿真、物理仿真、制造过程仿真及工作过程仿真,采用多种介质来存储、表达、处理多种信息,融文字、语音、图像、动画于一体,给人一种真实感及临境感。可视化技术比较典型的应用包括虚拟制造、虚拟现实等,具体体现在以下几个方面。

(1)科学计算结果数据的数字及图形动态显示。

(2)产品及其零部件的几何仿真及装配过程仿真。

(3)产品性能的物理及力学仿真。

(4)产品工作过程仿真,使其具有临境感及可驾驭感。


总之,制造过程的自动化程度是制造技术先进性的主要标志之一,也是21世纪现代制造技术中的一个最活跃的环节。制造自动化的发展将以其柔性化、集成化、敏捷化、智能化、全球化的特征来满足市场快速变化的要求。我国制造自动化的发展是以立足国情、瞄准世界先进水平、提高竞争力为前提,采用人机结合的适度自动化技术,将自动化程度较高的设备(如数控机床、工业机器人)和自动化程度较低的设备有效地组织起来,在此基础上实现以人为中心、以计算机为重要工具,构建具有柔性化、智能化、集成化、快速响应和快速重组的制造自动化系统。显然,制造自动化技术是我国必须大力发展的重要技术领域。

\section{先进制造技术发展趋势}
随着先进机电设备、互联网和人工智能的应用,今天的现代机械制造业已不再是传统意义上的机械加工,而已成为集机械、电子、光学、信息科学、材料学、生物学、管理科学等为一体的综合技术。近年来,先进制造技术正向着高精度、高柔性、绿色化、智能化方向快速发展。

精密制造技术(包括精密加工、超精密加工、微细加工、超微细加工及微纳加工等)的出现,使加工精度由微米级向纳米级发展。普通级数控机床的加工精度已由10 $\mu$m提高到5 $\mu$m,精密级加工中心则从3$\sim$5 $\mu$m提高到1$\sim$1.5 $\mu$m,而超精密加工精度已稳定进入纳米级。

得益于数控技术的快速发展,传统数控系统将逐渐发展为开放式柔性系统,可针对最终用户需求,通过改变、增加或剪裁数控系统的功能,将生产准备、零件加工、产品装配、检测等各个自动化子系统有机地综合集成起来,形成系列化的面向不同品种、不同加工要求的高柔性生产系统。

2011年,德国政府推出了“工业4.0”概念,旨在将传统生产的“集中式控制”模式转变为“分布式智能”模式,以建立一个高度灵活的个性化和数字化的产品与服务的智能制造产业,从而消除传统的行业界限。我国先后出台《中国制造2025》(2015年)、《“十四五”智能制造发展规划》(2021年)等制造强国建设纲领性文件,持续推动我国由制造大国向制造强国迈进。

面对日趋严格的环境及资源约束和实现“双碳”目标的紧迫性,先进制造技术应加快绿色化转型,以提升我国制造业的竞争力和在全球产业链中的位置。目前,以精密成形技术、无切削液加工、增材制造等为代表的绿色制造技术,正在提升能源和原材料的利用效率,降低制造成本与环境影响,推动制造业走向全流程绿色化。

人工智能(artificial intelligence,AI)作为引领新一轮科技与产业变革的核心驱动力,正在深度融入先进制造技术体系与价值链中,推动制造业由自动化、数字化、网络化向智能化阶段跃迁。人工智能与制造业融合,可扩大、延伸和甚至部分取代人类专家在制造过程中的脑力劳动,其应用已从单点辅助性工具,演进为支撑制造系统实现自主感知、决策优化与协同运行的核心支柱。人工智能的广泛应用已驱动研发设计、生产及管理体系的全链条重塑,为发展新质生产力提供了关键支撑。

\section*{复习思考题}


17-1 举例说明超精密加工的重要应用领域。

17-2 简述超精密磨削加工的特点。

17-3 超精密研磨有哪几种方法?各种方法所能达到的精度是多少?

17-4 简述纳米加工技术的特点和应用前景。

17-5 简述光刻技术的种类及特点。

17-6 高速切削是如何定义的?它有哪些优点?

17-7 简述柔性制造系统的分类和特点。

17-8 简述虚拟制造系统的基本概念与原理。

17-9 试解释CAD/CAM的基本概念。

17-10 简述CAD/CAM发展的关键技术。

17-11 计算机辅助工艺规划(CAPP)的基本功能包括哪些方面?

17-12 简述现代工业发展绿色制造技术的主要意义。

17-13 试举例说明人工智能在制造技术上的主要应用。


%\section*{参考文献}
\begin{thebibliography}{99}
    \bibitem{1} 齐乐华.工程材料及成形工艺基础[M].2版.西安:西北工业大学出版社,2020.
    \bibitem{2} 杨方.机械加工工艺基础[M].2版.西安:西北工业大学出版社,2020.
    \bibitem{3} 傅水根,马二恩,张学政.机械制造工艺基础[M].北京:清华大学出版社,1998.
    \bibitem{4} BLACK J T, KOHSER R A. Materials processes in manufacturing [M]. 11th Edition. New York: John Willey \& Sons, Inc., 2011.
    \bibitem{5} KALPAKJIAN S, SCHMID S. Manufacturing processes for engineering materials [M]. 6th Edition. Hoboken: Pearson, 2016.
    \bibitem{6} 徐自立,陈慧敏,吴修德.工程材料[M].武汉:华中科技大学出版社,2012.
    \bibitem{7} 贺毅,向军,胡志华.工程材料[M].2版.西安:西安交通大学出版社,2015.
    \bibitem{8} 周美玲,谢建新,朱宝泉.材料工程基础[M].北京:北京工业大学出版社,2001.
    \bibitem{9} 杜丽娟.工程材料成形技术基础[M].北京:电子工业出版社,2003.
    \bibitem{10} 吕广庶,张远明.工程材料及成形技术基础[M].2版.北京:高等教育出版社,2011.
    \bibitem{11} 梁耀能.工程材料及加工工程[M].北京:机械工业出版社,2001.
    \bibitem{12} 任家烈,吴爱萍.先进材料的连接[M].北京:机械工业出版社,2000.
    \bibitem{13} 刘金合.高能密度焊[M].西安:西北工业大学出版社,1995.
    \bibitem{14} 严绍华.材料成形工艺基础[M].北京:清华大学出版社,2001.
    \bibitem{15} 盛晓敏,邓朝晖.先进制造技术[M].北京:机械工业出版社,2011.
    \bibitem{16} 林建榕.机械制造基础[M].上海:上海交通大学出版社,2000.
    \bibitem{17} 张建华.精密与特种加工技术[M].北京:机械工业出版社,2003.
    \bibitem{18} 王丽英.机械制造技术[M].2版.北京:中国计量出版社,2009.
    \bibitem{19} 黄卫东,等.激光立体成形[M].西安:西北工业大学出版社,2007.
    \bibitem{20} 张文文.TC18钛合金飞机起落架支柱锻造成形研究[D].秦皇岛:燕山大学,2014.
    \bibitem{21} 朱金华.飞机起落架减震筒的加工工艺研究[J].制造技术与机床,2017(2):98-102.
    \bibitem{22} 郭文有,航空发动机叶片机械加工工艺[M].北京:国防工业出版社,1994.
\end{thebibliography}

郑重声明

高等教育出版社依法对本书享有专有出版权。任何未经许可的复制、销售行为均违反《中华人民共和国著作权法》，其行为人将承担相应的民事责任和行政责任；构成犯罪的，将被依法追究刑事责任。为了维护市场秩序，保护读者的合法权益，避免读者误用盗版书造成不良后果，我社将配合行政执法部门和司法机关对违法犯罪的单位和个人进行严厉打击。社会各界人士如发现上述侵权行为，希望及时举报，本社将奖励举报有功人员。

反盗版举报电话 (010)58581999 58582371
反盗版举报邮箱 dd@hep.com.cn
通信地址 北京市西城区德外大街4号
高等教育出版社知识产权与法律事务部
邮政编码 100120

防伪查询说明

用户购书后刮开封底防伪涂层，使用手机微信等软件扫描二维码，会跳转至防伪查询网页，获得所购图书详细信息。

防伪客服电话

(010)58582300
\end{document}